\documentclass[11pt,a4paper]{article}
\usepackage[T1]{fontenc}
\usepackage{lmodern,microtype}
\usepackage[margin=27mm]{geometry}
\usepackage{amsmath,amssymb,amsthm,mathtools,booktabs,array,mathrsfs}
\PassOptionsToPackage{hyphens}{url}
\usepackage[hidelinks]{hyperref}
\emergencystretch=2em
\newtheorem{theorem}{Theorem}[section]
\newtheorem{proposition}[theorem]{Proposition}
\newtheorem{lemma}[theorem]{Lemma}
\newtheorem{corollary}[theorem]{Corollary}
\theoremstyle{definition}
\newtheorem{definition}[theorem]{Definition}
\theoremstyle{remark}
\newtheorem{remark}[theorem]{Remark}
\newcommand{\GT}{\widehat{\mathrm{GT}}}
\newcommand{\GTz}{\widehat{\mathrm{GT}}_0}
\newcommand{\Zh}{\widehat{\mathbf Z}}
\newcommand{\Z}{\mathbf Z}
\newcommand{\Q}{\mathbf Q}
\newcommand{\F}{\mathbf F}
\newcommand{\GQ}{G_{\Q}}
\newcommand{\ab}{\mathrm{ab}}
\newcommand{\cts}{\mathrm{cts}}
\newcommand{\cl}{\overline}
\DeclareMathOperator{\Hom}{Hom}
\DeclareMathOperator{\Gal}{Gal}
\DeclareMathOperator{\Aut}{Aut}
\DeclareMathOperator{\Out}{Out}
\DeclareMathOperator{\Br}{Br}
\DeclareMathOperator{\SL}{SL}
\DeclareMathOperator{\PSL}{PSL}
\DeclareMathOperator{\Ad}{Ad}
\DeclareMathOperator{\res}{res}
\DeclareMathOperator{\infla}{inf}
\newcommand{\workbenchurl}{https://github.com/KokunoYumeto/k-theory-and-galois-groups-workbench}
\newcommand{\recordurl}{https://github.com/KokunoYumeto/mathematics-workbenches}
\title{The arithmetic quotient of the profinite Grothendieck--Teichm\"uller group}
\author{}
\date{25 November 2026}
\begin{document}
\maketitle
\begin{abstract}
Ihara extended the Kummer cocycles to the profinite Grothendieck--Teichm\"uller group, Enriquez proved their distribution relations, and Nakamura and Schneps related the cocycle of two to the Teichm\"uller tower. Starting from the established normalized--Galois decomposition, we identify the arithmetic quotient and the relative abelian obstruction to equality with the Galois group. We compute the entire invariant Galois character target at every prime by explicit corrected cup products. At prime two, finite cyclic-cover homology lifts every target character to full GT, proving that the quadratic restriction map is surjective with an explicit linear section. The quadratic character space and its arithmetic restriction image both have countably infinite dimension.

Using Willwacher's characteristic-two free generators, we prove that the full group and its Galois subgroup have the same dyadic free-group action, with image the free pro-two product of a group of order two and $\mathbf Z_2$. Every integral leading Lie lattice is attained arithmetically. Simultaneously removing this action, all rational Kummer coordinates, the entire coefficient-two cubic CM image and all the chosen corrected cup-product characters preserves the full relative abelian quotient. For its maximal pro-two quotient one can further remove every single-prime word action and both odd-primary CM images. An exact real-unit square-class criterion then determines further arithmetic CM images and permits their simultaneous removal, while preserving the relative obstruction and original character kernel. The local determinant is governed by the cyclotomic relative class number. Global Frobenius elements fill the local deficits at degrees twenty-nine and one hundred thirteen, including the finite signs; at degree one hundred sixty-three an explicit unit identity proves an arithmetic index of four. The latter full image has two possibilities, separated by a fixed finite-word test; the larger one would give a quotient of order twelve. An odd-inertia argument eliminates the conjugation obstruction for every collection of CM scalar kernels, including the infinite collection. Their full and arithmetic images are exact products even when individual images remain unknown. Extending an intrinsic invariant relative character is therefore equivalent to the vanishing of its relative central extension class. The quadratic restriction image is completely determined; its relative kernel remains open.

The corrected adelic modular frames have exactly their joint Galois image. We determine its overlap with rational radicals, prove the normalized harmonic and cubic comparisons, and compute the exact joint odd CM image. At every odd coefficient prime, all arithmetic target characters survive the common kernel of the rational, single-prime, CM and finite-abelian-cover homology realizations; these realizations therefore do not provide the missing odd-prime lifts. We also construct explicit relative quadratic characters whose vanishing is undecided. The mathematics and text were produced by GPT-6 Astra Pro (ChatGPT Pro mode), with independent same-model re-proofs and targeted checks; no human has checked them.
\end{abstract}

\section{The question and the results}
The comparison between the absolute Galois group $\GQ$ and Drinfeld's profinite Grothendieck--Teichm\"uller group $\GT$ is a central question in Grothendieck's programme. Drinfeld introduced the group by the two-cycle, three-cycle and pentagon equations; Ihara developed the arithmetic characters of its Galois action. We use the standard tangential-base-point inclusion $\GQ\subset\GT$. All groups, homomorphisms and cohomology in this paper are profinite or continuous unless another meaning is stated.

Write $G=\GT$. Let $\lambda:G\to\Zh^\times$ be the multiplier, and let
$\rho_2:\GT\to\Zh(1)$ be the cocycle extending the action on the compatible positive roots of $2$. Thus
\[
 \rho_2(gh)=\rho_2(g)+\lambda(g)\rho_2(h).
\]
Put
\[
 K_2=\Q(\mu_\infty,2^{1/\infty}),\qquad
 N=\ker\lambda\cap\ker\rho_2,\qquad H=\GQ,\qquad K=G_{K_2}.
\]
The notation $N$ denotes the normalized subgroup also written $\Lambda_0$. It lies in the group $\Lambda$ of Hatcher, Lochak and Schneps, whose formulas act on the pure Teichm\"uller tower. The input used here is the decomposition
\begin{equation}\label{eq:decomposition}
 \GT=NH.
\end{equation}
It is Theorem 2 of \cite{H17}; Lemma 7 there gives the cocycle law and evenness of $\rho_2$ when $\lambda=1$. The same note proves relation (IV), relation $(\mathrm{III}')$ and the resulting identification of the Nakamura--Schneps subgroup with $\GT$ in Theorems 1 and 20 and Corollary 21. These statements are cited background, with their stated AI checks and no human verification. The gentle group $\GTz$ is defined by the two- and three-cycle equations. Invertibility and a unit multiplier follow from these equations by \cite[Theorems A and E]{Monoid}; no extra invertibility test is imposed below.

For a prime $p$, write
\[
 X_p=H^1(N,\F_p)^{\GT},\qquad
 Y_p=H^1(K,\F_p)^H.
\]
Coefficients have trivial action. The restriction in the question lands in $Y_p$, although its codomain can of course be written $H^1(K,\F_p)$.

\begin{theorem}\label{thm:main-structural}
Assume \eqref{eq:decomposition} and the preceding normalization of $\rho_2$. Then $N$ is closed and normal, $N\cap H=K$, and
\[
 \GT/N\simeq A:=\Gal(K_2/\Q).
\]
In the affine coordinates $(u,t)$, with product
$(u,t)(v,s)=(uv,t+us)$, this quotient is
\[
 A=\{(u,t)\in\Zh^\times\ltimes\Zh:
 t\equiv\varepsilon_8(u)\pmod2\},
\]
where $\varepsilon_8(u)=0$ for $u\equiv1,7\pmod8$ and $1$ for $u\equiv3,5\pmod8$.
There is a canonical topological isomorphism
\[
 \GT^{\ab}\simeq\Zh^\times\times B,\qquad
 B=N/\overline{[\GT,N]K}.
\]
For every prime $p$, restriction identifies
\[
 \Hom(B,\F_p)\simeq\ker(X_p\longrightarrow Y_p).
\]
Consequently the relative kernel vanishes for every prime if and only if the multiplier identifies $\GT^{\ab}$ with $\Zh^\times$.
\end{theorem}

In words, the multiplier and the cocycle of two describe precisely this Galois quotient. All extra abelian characters belong to $B$. This is the precise form of Lochak and Schneps's Question 1.2 \cite[\S1, Question 1.2]{LS2006}. It is weaker than deciding whether $\GT=\GQ$: a group can have trivial relative abelianization and still have a nontrivial relative nonabelian quotient.

Hoshi--Mochizuki--Tsujimura give a group-theoretic construction of the
conjugacy class of the Galois subgroup through their arithmetic Belyi
compatibility and algebraic ampleness conditions
\cite[Theorem~C(ii), p.~5; Definition~5.12, pp.~87--90]{HMT2025}.
Their theorem identifies the subgroup by those additional conditions.
The present question asks for its relative abelian obstruction inside
full $\GT$.

\begin{theorem}\label{thm:main-dimension}
With the same notation,
\[
 \dim_{\F_2} X_2=\dim_{\F_2}\operatorname{im}(X_2\longrightarrow Y_2)=\aleph_0.
\]
For every positive integer $d$, Ihara's integral cocycle $\rho_d^I$ is even on the multiplier-one subgroup of $\GTz$. For every odd prime $q$,
\[
 \psi_q=\frac{\rho_q^I|_N}{2}\pmod2
\]
is an invariant continuous character of $N$. Its restriction is the character of $K_2(q^{1/4})/K_2$, and these restrictions are linearly independent as $q$ varies. For every prime $p$, including odd $p$, $\dim_{\F_p}Y_p=\aleph_0$.
\end{theorem}

Thus the answer to the dimension question is infinite even if $\GT=\GQ$. It is the restriction, rather than the existence of invariant characters, that tests equality. In particular this theorem supplies no nonzero character of $\GT$ trivial on $\GQ$.

The entire arithmetic target can be described explicitly. Put
\(U=\Zh^\times\), let \(a,b\) be the quadratic characters of
\(\Q(\sqrt2)\) and \(\Q(i)\), and for each odd prime \(p\) put
\[
 p^*=(-1)^{(p-1)/2}p,\qquad
 c_p=\epsilon_{p^*},\qquad
 \varepsilon_p=\frac{1-(2/p)}2,\qquad
 \ell_p(d)=\frac{1-(d/p)}2.
\]
Here the last two quantities lie in \(\F_2\), and \(p\nmid d\)
in the last definition. Let \(\delta_2=b^2\). For odd \(p\),
let \(\delta_p\in H^2(U,\F_2)\) be the inflation of the central
cyclic extension
\[
 C_{2^{v_2(p-1)+1}}\longrightarrow C_{2^{v_2(p-1)}}
\]
from the maximal two-primary quotient of \(\Z_p^\times\).

\begin{theorem}\label{thm:main-full-quadratic-target}
With \(A=G/N=H/(H\cap N)\), transgression is an isomorphism
\[
 \operatorname{tr}:Y_2\xrightarrow{\ \sim\ }
 \ker\bigl(H^2(A,\F_2)\longrightarrow\operatorname{Br}(\Q)[2]\bigr).
\]
There is a direct-sum decomposition
\[
 Y_2=\F_2\eta_3
 \oplus\bigoplus_{p\text{ odd}}\F_2\psi_p
 \oplus\bigoplus_{p\text{ odd}}\F_2\kappa_{2,p}
 \oplus\bigoplus_{\substack{p<q\\p,q\text{ odd}}}\F_2\kappa_{p,q},
\]
where \(\eta_3=m_3/8\bmod2\), with
\(m_3(g)=[S^2T]f_g\) in the expansion \(x=1+S,y=1+T\),
is the cubic character and
the additional basis characters are uniquely specified by
\begin{align*}
 \operatorname{tr}(\kappa_{2,p})
 &=\operatorname{Inf}_U^A
       \bigl(ac_p+\varepsilon_p(\delta_2+\delta_p)\bigr),\\
 \operatorname{tr}(\kappa_{p,q})
 &=\operatorname{Inf}_U^A
       \bigl(c_pc_q+\ell_p(q^*)\delta_p+\ell_q(p^*)\delta_q\bigr).
\end{align*}
All products are cup products, and every character is a finite sum.
The character \(\kappa_{5,29}\) is realized by the central quadratic
layer of
\[
 \Q\bigl(\sqrt5,\sqrt{29},\sqrt{7+2\sqrt5}\bigr)
       /\Q(\sqrt5,\sqrt{29}),
\]
restricted to \(G_{K_2}\). It is nonzero and independent of
\(\eta_3\) and the \(\psi_p\).
\end{theorem}

Theorem~\ref{thm:fr-Y2-complete} proves the decomposition by
computing the affine cohomology and cancelling each cup product's
local Brauer invariants. The span of \(\eta_3\) and the
\(\psi_p\) is therefore a proper subspace of \(Y_2\), with quotient
of countably infinite dimension. The next two constructions lift all these basis directions to
\(X_2\), giving the complete restriction image.

For odd coefficients the translation coordinate disappears from the
cohomology, and the target has a uniform description. Fix an odd prime
\(p\), put \(\mathcal S_p=\{\ell:\ell\text{ prime},\ \ell\equiv1\pmod p\}\),
and choose cyclic coordinates in
\[
 U^{(p)}=\Z_p\times\prod_{\ell\in\mathcal S_p}C_{p^{v_p(\ell-1)}}.
\]
Write \(\gamma_0,\gamma_\ell\) for their reductions to \(\F_p\),
and \(\delta_\ell\) for the central cyclic extension class obtained
by increasing the exponent of the \(\ell\)-factor by one, with
kernel coordinate \(x\mapsto p^{v_p(\ell-1)}x\).

\begin{theorem}\label{thm:main-full-odd-target}
For every odd prime \(p\), inflation
\(H^r(U,\F_p)\to H^r(A,\F_p)\) is an isomorphism in every degree.
The arithmetic images of the \(\delta_\ell\) form a basis of
\(H^2(H,\F_p)\). Put
\(a_\ell=(\ell-1)/p^{v_p(\ell-1)}\bmod p\), and define
\begin{align*}
 D_{\ell m}&=a_\ell^{-1}\gamma_m(\ell)\delta_\ell
                  -a_m^{-1}\gamma_\ell(m)\delta_m\quad(\ell<m),\\
 D_{0\ell}&=-a_\ell^{-1}\gamma_0(\ell)\delta_\ell.
\end{align*}
Here Frobenius acts by the positive power map.
Transgression identifies \(Y_p\) with
\[
 \ker\bigl(H^2(A,\F_p)\longrightarrow H^2(H,\F_p)\bigr),
\]
and its inverse images of
\(\gamma_i\smile\gamma_j-D_{ij}\), for
\(i<j\) in \(\{0\}\cup\mathcal S_p\), form a basis of \(Y_p\).
Every character is a finite sum of these basis elements.
\end{theorem}

Theorem~\ref{thm:odd-full-galois-target} proves this by taking the
cyclotomic eigenspace of the global Brauer exact sequence over
\(\Q(\mu_p)\). Only primes congruent to one modulo \(p\) contribute,
and each cup product needs a correction at at most its two positive
indices. This determines the codomain of the odd-prime restriction
problem; its image and kernel for full \(G\) remain open.


The corrected cup products indexed by two odd primes all lift to full
GT. The construction uses finite homology of cyclic covers and keeps
track of the entire central extension, including its cyclic terms.
For a multiplier-one element \(g\), let \(B_g(X,Y)\) denote the
normalized scalar on the completed Alexander commutator module,
as defined in Section~\ref{sec:all-odd-cross-cup-lifts}.

\begin{theorem}\label{thm:main-odd-pair-lifts}
For every ordered pair of distinct odd primes \(\ell,q\), there is
an invariant continuous homomorphism \(\xi_{\ell,q}:M\to\F_2\),
where \(M=\ker\lambda\), whose restriction to \(N\) belongs to
\(X_2\) and satisfies
\[
                    \xi_{\ell,q}|_K=\kappa_{\ell,q}.
\]
Here \(\kappa_{\ell,q}=\kappa_{q,\ell}\) denotes the basis
character indexed by the unordered pair. Its transgression is
\[
 c_\ell c_q+\ell_\ell(q^*)\delta_\ell
                         +\ell_q(\ell^*)\delta_q.
\]
In particular this corrected class inflates to zero in
\(H^2(G,\F_2)\).

Explicitly, put \(d=\operatorname{ord}_\ell(q)\), and choose a
primitive \(\ell\)-th root \(\zeta\in\F_{q^d}\).
If \(d=2s\), choose representatives \(a_i\) of
\((\Z/\ell)^\times/\langle q\rangle\); then
\[
 (-1)^{\xi_{\ell,q}(g)}
   =\prod_i B_g(\zeta^{a_i},\zeta^{a_i})^{(q^s+1)/2}.
\]
If \(d\) is odd, choose representatives of
\((\Z/\ell)^\times/\langle q,-1\rangle\); then
\[
 (-1)^{\xi_{\ell,q}(g)}
   =\prod_i B_g(\zeta^{a_i},\zeta^{a_i})^{(q^d-1)/2}.
\]
The values are signs independent of these choices.
The same construction works for any continuous positive peripheral
action on \(\widehat F(x,y)\) of the form
\(x\mapsto x^\lambda\), \(y\mapsto f^{-1}y^\lambda f\),
with \(f\in\widehat F(x,y)'\) and
\((xy)^{-1}\mapsto\text{a conjugate of }(xy)^{-\lambda}\).

The two ordered constructions have the same arithmetic restriction.
Their sum on \(N\) extends uniquely to a continuous character
\(\varepsilon_{\ell,q}:G\to\F_2\) vanishing on \(H\).
Whether these relative characters vanish is not determined here.
\end{theorem}

Theorem~\ref{thm:all-odd-cross-cup-lifts} proves the general
construction. The pair \((5,29)\) gives the smallest worked example
used below: its two cyclic correction terms are zero, and the
character detects the classical dihedral field in
Theorem~\ref{thm:main-full-quadratic-target}.
The arithmetic relation between Fermat homology and Anderson's beta
function is classical; Otsubo gives its precise cocycle comparison
in \cite[\S7, p.~637, (7.1)--(7.2)]{Otsubo2016}.
Here the finite symplectic normalizer computes the central extension
for the full peripheral action.


Integral homology at coefficient two supplies the remaining family.
Its character sees homology modulo eight, where the square of a lift
of multiplier minus one records the necessary cyclic correction.

\begin{theorem}\label{thm:main-complete-quadratic}
For every odd prime \(p\), there is a continuous invariant
character \(\xi_{2,p}:M\to\F_2\), determined by compact homology
modulo eight of the degree-\(p\) cyclic cover with monodromies
\(1,1,-2\), such that \(\nu_p=\xi_{2,p}|_N\in X_2\) and
\[
 \operatorname{tr}(\nu_p)
     =\operatorname{Inf}_U^A
          \bigl(ac_p+\varepsilon_p(\delta_2+\delta_p)\bigr),
 \qquad \nu_p|_K=\kappa_{2,p}.
\]
The character is defined for every positive peripheral action with
the hypotheses of Theorem~\ref{thm:main-odd-pair-lifts}.
Consequently the original quadratic restriction is surjective, with
an explicit linear section:
\[
 0\longrightarrow\operatorname{Hom}_{\rm cts}(B,\F_2)
 \longrightarrow X_2\xrightarrow{\ \mathrm{res}\ }Y_2
 \longrightarrow0.
\]
One section sends the basis elements \(\eta_3,\psi_p,\kappa_{2,p},
\kappa_{\ell,q}\) to \(\eta_3,\psi_p,\nu_p,\xi_{\ell,q}\),
respectively, with \(\ell<q\). The companion character of the same
integral homology construction satisfies
\[
 \widetilde\psi_p=\rho_p^I/2\pmod2
                 \quad\text{on }M\text{ for every odd }p.
\]
At \(p=3\), the lift \(\nu_3\)
equals the cubic CM lift \(\tau+\psi_3\)
defined below; use this latter section \(s\) in the simultaneous
normalization, and put \(\mathcal S=s(Y_2)\). Then
\[
               X_2=\operatorname{Hom}_{\rm cts}(B,\F_2)
                                      \oplus\mathcal S.
\]
The two inflation kernels on \(H^2(A,\F_2)\), towards
\(H^2(G,\F_2)\) and towards \(\operatorname{Br}(\Q)[2]\),
are equal.
\end{theorem}

Every arithmetic invariant quadratic character therefore extends.
The unresolved part of the original character problem is exactly
\(\operatorname{Hom}_{\rm cts}(B,\F_2)\): surjectivity does not
prove that this kernel vanishes. The complete proofs are in
Sections~\ref{sec:coefficient-two-cross-cups} and
\ref{sec:companion-kummer-equality} and
\ref{sec:complete-quadratic-restriction}.


The normalization can be made substantially stronger. Let
\[
 L_{\mathrm{rad}}=\Q(\mu_\infty,q^{1/n}:q\text{ prime},\ n\geq1),
 \qquad
 N_{\mathrm{rad}}=N\cap\bigcap_{q\text{ odd prime}}\ker\rho_q^I.
\]
Here all radicals are the coherent positive roots.

\begin{theorem}\label{thm:main-rational-normalization}
With the preceding assumptions and notation, the joint map
\[
 N\longrightarrow\prod_{q\text{ odd prime}}2\Zh,
 \qquad n\longmapsto(\rho_q^I(n))_q
\]
is surjective already on $K$. Its kernel $N_{\mathrm{rad}}$ is
a proper closed normal subgroup, and
\[
 \GT=N_{\mathrm{rad}}\GQ,\qquad
 N_{\mathrm{rad}}\cap\GQ=G_{L_{\mathrm{rad}}},\qquad
 \GT/N_{\mathrm{rad}}\simeq\Gal(L_{\mathrm{rad}}/\Q).
\]
The latter quotient consists of the affine coordinates
$(u,(t_q)_{q\text{ prime}})$ satisfying
$t_q\equiv\epsilon_q(u)\pmod2$, with multiplication
$(u,t)(v,s)=(uv,t+us)$; here $\epsilon_q$ is the character of
$\Q(\sqrt q)$. The relative abelian quotient is unchanged:
\[
 B\simeq N_{\mathrm{rad}}/
       \overline{[\GT,N_{\mathrm{rad}}]G_{L_{\mathrm{rad}}}}.
\]
The invariant quadratic characters of $N$ that factor through all
rational Kummer coordinates are exactly the finite linear combinations
of the $\psi_q$ for odd primes $q$, and restriction is injective on
this entire subspace. No nonzero invariant character with values in
$\F_p$, for odd $p$, factors through these coordinates.
\end{theorem}

This removes all rational-radical data from the open problem at once.
An extra character would have to be nonzero on the smaller subgroup
$N_{\mathrm{rad}}$. The theorem does not show that its relative
abelianization is zero. Enriquez's distribution theorem is essential
for describing the whole rational Kummer subspace: it proves
$\rho_{ab}^I=\rho_a^I+\rho_b^I$ integrally on $\GT$
\cite[Theorem~8.1 and Remark~8.6]{Enriquez2007}.

A further normalization removes the entire dyadic free-group action.
Put $M=\ker\lambda$, $D=N_{\mathrm{rad}}$, and let $F^nM$ consist
of the elements whose word $f_g-1$ has no pro-$2$ Magnus term
below degree $n$. Define
\[
 R_D=\{g\in D:f_g=1\text{ in the free pro-$2$ group}\}.
\]

\begin{theorem}\label{thm:main-all-degree-dyadic}
For every $n\ge3$, the full and arithmetic leading images are
\[
 \operatorname{in}_n(F^nM)
 =\operatorname{in}_n(F^nM\cap H)
 =\mathfrak{grt}_1(\Q_2)_n
       \cap\operatorname{Lie}_{\Z_2}(S,T)_n.
\]
Arithmetic odd-degree elements freely generate this integral Lie
algebra, and their reductions are Willwacher's characteristic-two
generators. The full group $\GT$ and $H$ have the same image in
$\operatorname{Aut}(\widehat F_2^{(2)})$. The subgroup $R_D$ is
closed and normal, and
\[
 \GT=R_DH,\qquad
 B\simeq R_D/\overline{[\GT,R_D](R_D\cap H)}.
\]
Put $m_3(g)=[S^2T]f_g$ and $\eta_3=m_3/8\pmod2$ on $N$.
Then $\eta_3\in X_2$, and
\[
 \ker\bigl(X_2\longrightarrow H^1(R_D,\F_2)\bigr)
 =\left(\bigoplus_{q\text{ odd prime}}\F_2\psi_q\right)
                         \oplus\F_2\eta_3.
\]
Arithmetic restriction is injective on this kernel.
Let $P$ be the image of $M$ in the dyadic action. It is the free
pro-$2$ group on a convergent family with one arithmetic generator
in each odd degree at least three. Its integral mass
$\mu=\rho_{2,2}^I/2:P\to\Z_2$ is surjective. The images of $N$
and $D$ are both $P_0=\ker\mu$, and
\[
 H^1(P,\F_2)^G=\F_2(\mu\bmod2),\qquad
 H^1(P_0,\F_2)^G=\F_2\eta_3.
\]
\end{theorem}

Thus the relative character problem survives entirely on elements
invisible to dyadic free-group monodromy and to every rational
Kummer cocycle. It is not settled by the equality of the dyadic
action images. Section~\ref{sec:fr-uniform-dyadic-lie} gives the
uniform proof; the finite layers below provide explicit coefficients,
arithmetic witnesses and the precise extension thresholds of their
quadratic characters.

The group acting on the free pro-two group has an explicit presentation.

\begin{theorem}\label{thm:main-full-dyadic-image}
Under the preceding hypotheses, the full dyadic action image is
\[
 \operatorname{im}\bigl(G\to\operatorname{Aut}(\widehat F_2^{(2)})\bigr)
       \cong C_2*_{\mathrm{pro}\text{-}2}\mathbf Z_2.
\]
The factors can be generated by complex conjugation and an arithmetic
element of dyadic multiplier \(-3\). Its cyclotomic map is its
abelianization, with kernel the free pro-two group \(P=M/R_M\).
\end{theorem}

This determines the complete dyadic realization. It does not identify
its full profinite source with the arithmetic subgroup.
Section~\ref{sec:full-dyadic-image} proves the presentation.
The same abstract free product describes the maximal pro-two Galois
extension of \(\mathbf Q\) unramified outside two and infinity,
by Neumann's Theorem~3.7 \cite[p.~244]{Neumann1977}.
The theorem above determines the image of the specified fundamental-group
representation using its arithmetic leading lattices.

For the quadratic problem there is a still smaller domain. Define
\[
 P_D=\{g\in D:f_g=1\text{ in every free pro-$p$ completion}\},
 \qquad T_D=P_D\cap\ker d_{(3)},
\]
where $d_{(3)}$ is the order-three radical--CM defect defined below.

\begin{theorem}\label{thm:main-all-prime-quadratic-domain}
Both subgroups are closed and normal. Their quotients in $R_D$
are pro-odd, and
\[
 \begin{aligned}
 B_{\max2}&\simeq
 \left(P_D/\overline{[G,P_D](P_D\cap H)}\right)_{\max2}\\
 &\simeq
 \left(T_D/\overline{[G,T_D](T_D\cap H)}\right)_{\max2}.
 \end{aligned}
\]
Restriction from $X_2$ to either subgroup has the same rational
Kummer-plus-cubic kernel as restriction to $R_D$.
The indices $[R_D:P_D]$ and
$[R_D\cap\ker d_{(3)}:T_D]$ are infinite.
No factorization $G=P_DH$ or $G=T_DH$ is asserted.
\end{theorem}

The odd-primary harmonic CM image gives a further strict reduction.
Let \(W:D\to N_4(4)\) be the corrected harmonic CM homomorphism
and let \(W_{\rm odd}\) denote its projection to the maximal pro-odd
factor of its norm-one target.

\begin{theorem}\label{thm:main-odd-cm-domain}
The subgroup \(U_D=T_D\cap\ker W_{\rm odd}\) is closed and normal,
and
\[
 B_{\max2}\cong
 \left(U_D/\overline{[G,U_D](U_D\cap H)}\right)_{\max2}.
\]
Restriction from \(X_2\) to \(U_D\) has the same kernel as restriction
to \(R_D\). Moreover
\[
 W_{\rm odd}(T_D\cap H)=N_4(4)_{\rm odd},\qquad
 T_D=U_D(T_D\cap H),
\]
and \(U_D\) has infinite index in \(T_D\), already on the arithmetic
intersection. No equality \(G=U_DH\) is asserted.
\end{theorem}

The witness is a commutator of arithmetic elements with disjoint
primary orders. It is invisible in every single-prime fundamental-group
action, but has a nonzero harmonic CM value. Thus the CM restriction
removes additional arithmetic directions from the residual domain.

The odd Eisenstein image can be removed as well. If \(T_{\rm odd}\)
is its pro-odd projection, put \(V_D=U_D\cap\ker T_{\rm odd}\).
Then \(V_D\) is closed and normal,
\[
 B_{\max2}\cong
 \left(V_D/\overline{[G,V_D](V_D\cap H)}\right)_{\max2},
 \qquad [U_D:V_D]=\infty.
\]
The original quadratic restriction kernel is unchanged. The infinite
index is witnessed arithmetically by the thirteen-primary cubic CM
value in Theorem~\ref{thm:joint-odd-cm-gate}. The entire joint odd
image can also be determined. Write
\[
 \mathcal N_{3,\ell}=
 \{z\in\mathbf Z_\ell[\omega]^\times:z\bar z=1\},
 \qquad \omega^2+\omega+1=0,
\]
and let its subscript \({\rm odd}\) mean the maximal pro-odd factor.

\begin{theorem}\label{thm:main-exact-joint-odd-cm}
Put
\[
 \mathcal C_3=\prod_{\ell\ge5}
           (\mathcal N_{3,\ell,{\rm odd}})^3,
 \qquad \mathcal A_4=N_4(4)_{\rm odd},
\]
with identity components at coefficient primes two and three in
\(\mathcal C_3\). Then
\[
 (T_{\rm odd},W_{\rm odd})(T_D)
 =(T_{\rm odd},W_{\rm odd})(T_D\cap H)
 =\mathcal C_3\times\mathcal A_4.
\]
The cubic image on \(U_D\) is exactly \(\mathcal C_3\), already
attained on \(U_D\cap H\). Hence
\[
 T_D/V_D\simeq\mathcal C_3\times\mathcal A_4,
 \qquad U_D/V_D\simeq\mathcal C_3,
\]
and
\[
 T_D=V_D(T_D\cap H),\qquad U_D=V_D(U_D\cap H).
\]
In addition, the full coefficient-three cubic frame is the identity
on \(P_D\).
\end{theorem}

The proof separates the two CM factors by arithmetic commutators.
The identity at three uses continuous specialization of the trivial
pro-three beta series at the third roots of unity. These internal
factorizations do not imply \(G=V_DH\); the prime-two relative
comparison above still uses the coprime restriction theorem.

The entire coefficient-two cubic image and the chosen lifts of all
corrected cup products can be removed simultaneously. Put
\[
 \mathcal O_2=\Z_2[\omega],\quad \omega^2+\omega+1=0,
 \quad \delta=1+2\omega,
 \quad \mathcal U=\{z\in1+4\mathcal O_2:z\bar z=1\}.
\]
Let \(\mathfrak t_2=T_{g,2}\) be the coefficient-two projection
of the corrected cubic CM homomorphism on \(M\).
Let \(\mathcal I\) consist of the odd-prime pairs \(\ell<q\),
and let \(\mathcal J\) consist of the odd primes other than three.

\begin{theorem}\label{thm:main-simultaneous-quadratic}
The integral logarithm defines an invariant character
\[
 \tau=\frac{\log\mathfrak t_2}{4\delta}\pmod2\in H^1(M,\F_2)^G,
 \qquad \tau|_K=\kappa_{2,3}+\psi_3.
\]
The joint map
\[
 (\mathfrak t_2,(\xi_{\ell,q})_{\mathcal I},(\nu_p)_{\mathcal J}):R_D
       \longrightarrow\mathcal U\times\prod_{\mathcal I}\F_2
                                   \times\prod_{\mathcal J}\F_2
\]
is surjective already on \(R_D\cap H\), and also on \(V_D\cap H\).
Define
\[
 \begin{split}
 R_{\rm quad}&=R_D\cap\ker\mathfrak t_2
          \cap\bigcap_{\mathcal I}\ker\xi_{\ell,q}
          \cap\bigcap_{\mathcal J}\ker\nu_p,\\
 V_{\rm quad}&=V_D\cap R_{\rm quad}.
 \end{split}
\]
These subgroups are closed and normal in \(G\), and
\[
 \begin{gathered}
 R_D/R_{\rm quad}\simeq V_D/V_{\rm quad}
   \simeq\mathcal U\times\prod_{\mathcal I}\F_2
                                   \times\prod_{\mathcal J}\F_2,\\
 G=R_{\rm quad}H,\qquad
 B\simeq R_{\rm quad}/
           \overline{[G,R_{\rm quad}](R_{\rm quad}\cap H)},\\
 B_{\max2}\simeq
 \left(V_{\rm quad}/
           \overline{[G,V_{\rm quad}](V_{\rm quad}\cap H)}\right)_{\max2}.
 \end{gathered}
\]
The two quotient descriptions also hold on arithmetic intersections.
Restriction from \(X_2\) to either residual subgroup has kernel
exactly \(\mathcal S=s(Y_2)\). Its image is precisely the space
of invariant quadratic characters on that subgroup vanishing on
its arithmetic intersection, canonically identified with
\(\operatorname{Hom}_{\rm cts}(B,\F_2)\).
\end{theorem}

This isolates the relative kernel after removing the complete known
arithmetic quadratic image. The subgroup \(V_{\rm quad}\) has
infinite index in \(V_D\), already arithmetically. The theorem
preserves the relative obstruction and does not compute its value.
The proofs are in Sections~\ref{sec:fr-cubic-coefficient-two},
\ref{sec:simultaneous-quadratic-normalization} and
\ref{sec:complete-quadratic-restriction}.


Prime-degree CM covers give a stronger simultaneous arithmetic
normalization. The distinction between a local reflex-norm image and
the full arithmetic image is essential here.
For each prime \(p\ge5\), put \(h=(p-1)/2\),
\(\mathcal O_p=\mathbb Z[\zeta_p]\otimes\mathbb Z_2\), choose a primitive
residue \(g\), and let \(s=[g]\), \(c=s^h\). Define
\[
 T_p=\{u\in1+2\mathcal O_p:u\,c(u)=1\},\qquad
 P_p=(s-1)T_p.
\]
The algebra \(\mathcal O_p\) includes every factor at two.
The compact-homology scalar of \(y^p=x(1-x)\), restricted to
\(R_{\rm quad}\), will be written \(h_p\); it is the scalar
\(\mathfrak h_p\) used in the earlier cyclic-cover construction.
Define \(Q_p\) by
\[
 \Phi_p^*=\sum_{a=1}^h[a]^{-1},\qquad
 (s-1)\Phi_p^*=(c-1)Q_p,\qquad
 Q_p\in\mathbb Z[s]/(s^h+1).
\]
Its augmentation is \(1\) or \(-1\). Let
\(M_p=(\mathcal O_p/2\mathcal O_p)^c\), with cyclic basis
\(\beta_k=\overline\zeta_p^{\,g^k}+\overline\zeta_p^{-g^k}\),
\(0\le k<h\).

The first condition on \(p\) is that the natural map
\[
 T_p/Q_pT_p\longrightarrow M_p/\overline Q_pM_p,\qquad
 u\longmapsto(u-1)/2\bmod2,
\]
is an isomorphism. Under this condition inclusion identifies
\(P_p/Q_pP_p\) with the same finite elementary abelian group
\(\mathscr C_p\). Write \(\pi_p:M_p\to\mathscr C_p\), and let
\(\mathscr E_p\) be the full unit group of
\(\mathbb Q(\zeta_p+\zeta_p^{-1})\).
The remaining condition is injectivity of
\[
 \mathscr K_p:\mathscr C_p^\vee
       \longrightarrow\mathscr E_p/(\{\pm1\}\mathscr E_p^2),
 \qquad
 \ell\longmapsto
 \left[\prod_{k=0}^{h-1}
   (\zeta_p^{g^k}+\zeta_p^{-g^k})^{\ell(\pi_p\beta_k)}\right].
\]
Let \(\Pi_{\rm global}\) be the set of primes satisfying both
conditions. Square roots in the target may be arbitrary real units.

\begin{theorem}\label{thm:main-cyclotomic-cm-normalization}
For every prime satisfying the first-layer condition, the exact
arithmetic image of \(h_p\) on \(R_{\rm quad}\cap H\) or
\(V_{\rm quad}\cap H\) is the preimage of
\((\ker\mathscr K_p)^\perp\) in \(P_p\). This remains true after
imposing any collection of the other-prime scalar kernels.
For \(p\in\Pi_{\rm global}\), the joint images of
\(R_{\rm quad}\), \(V_{\rm quad}\), and both arithmetic
intersections are all \(\prod_pP_p\), with all finite signs included.
Put
\[
 R_{\rm global}=R_{\rm quad}\cap
          \bigcap_{p\in\Pi_{\rm global}}\ker h_p,\qquad
 V_{\rm global}=V_{\rm quad}\cap R_{\rm global}.
\]
Then
\[
 \begin{gathered}
 G=R_{\rm global}H,\qquad
 R_{\rm global}/V_{\rm global}\simeq R_{\rm quad}/V_{\rm quad},\\
 B\simeq R_{\rm global}/
       \overline{[G,R_{\rm global}](R_{\rm global}\cap H)},\\
 B_{\max2}\simeq
 \left(V_{\rm global}/
       \overline{[G,V_{\rm global}](V_{\rm global}\cap H)}\right)_{\max2}.
 \end{gathered}
\]
The quotient in the first line is pro-odd, and its arithmetic image
is preserved. Both original quadratic restriction kernels are exactly
\(\mathcal S\), with image every invariant quadratic character
vanishing on the arithmetic intersection.

The set \(\Pi_{\rm global}\) contains every prime with odd cyclotomic
class number, and also \(29\) and \(113\); it does not contain \(163\).
At \(29\) and \(113\) the local index is eight but the full arithmetic
image is \(P_p\). At \(163\) the arithmetic image is exactly
\(Q_{163}P_{163}\), of index four.
\end{theorem}

The determinant test is the classical relative class-number formula
in an integral CM-type normalization. The global criterion uses
quadratic Kummer theory of real units and Otsubo's character-labelled
Jacobi-sum comparison. Frobenius at \(59\) and \(227\) fills the local
deficits at \(29\) and \(113\); at \(163\), an explicit negative-square
identity proves that the deficit persists arithmetically.
These arguments, with every finite sign factor, appear in
Sections~\ref{sec:cm-type-unit-criterion}--\ref{sec:simultaneous-global-cm}.
Every added scalar kernel has infinite index.

There is also a restriction statement after all prime scalar kernels,
even when their full arithmetic images are not known to coincide.
For a closed normal subgroup \(L\subseteq N\), put
\[
 B(L)=L/\overline{[G,L](L\cap H)}.
\]

\begin{theorem}\label{thm:main-all-cm-restriction}
Let
\[
 R_{\rm all}=R_{\rm quad}\cap\bigcap_{p\ge5}\ker h_p,\qquad
 V_{\rm all}=V_{\rm quad}\cap R_{\rm all}.
\]
Restriction from \(X_2\) to either group has kernel exactly
\(\mathcal S\). There are canonical maps
\[
 B(V_{\rm all})_{\max2}\simeq B(R_{\rm all})_{\max2}
                         \twoheadrightarrow B_{\max2}.
\]
Moreover \(R_{\rm all}/V_{\rm all}\simeq
R_{\rm quad}/V_{\rm quad}\), preserving the arithmetic image.
The last arrow is asserted to be surjective.
\end{theorem}

Thus every original relative quadratic character remains detectable
on the all-prime kernel. Identifying its intrinsic relative quotient
with \(B_{\max2}\) requires extending those intrinsic characters.
The next theorem gives the exact quadratic extension condition.

\begin{theorem}\label{thm:main-all-cm-central-obstruction}
Let \(\mathcal I\) be any set of primes at least five and put
\[
 R_0=R_{\rm quad},\qquad
 C_p=h_p(R_0),\qquad D_p=h_p(R_0\cap H),\qquad
 J=R_0\cap\bigcap_{p\in\mathcal I}\ker h_p.
\]
The joint full and arithmetic images are exactly
\[
 C=\prod_{p\in\mathcal I}C_p,\qquad
 D_0=\prod_{p\in\mathcal I}D_p.
\]
For the discrete conjugation module
\(U=\operatorname{Hom}_{\rm cont}(C/D_0,\mathbb F_2)\),
\[
 H^j_{\rm cont}(G/R_0,U)=0\qquad(j\ge0).
\]
Let \(\mathscr X(J)\) be the space of continuous \(G\)-invariant
quadratic characters of \(J\) that vanish on \(J\cap H\).
Restriction identifies
\[
 \operatorname{Hom}_{\rm cont}(B,\mathbb F_2)
 \simeq\ker\!\left(
 o_2:\mathscr X(J)\longrightarrow H^2_{\rm cont}(C,D_0;\mathbb F_2)
 \right).
\]
Here \(o_2\) is the relative central extension class, with the
splitting over \(D_0\) specified by the arithmetic subgroup.
Every character in this kernel has a unique extension to \(G\)
which vanishes on \(H\).
\end{theorem}

The proof is Proposition~\ref{prop:all-cm-exact-products},
Corollary~\ref{cor:all-cm-conjugation-vanishing}, and
Theorem~\ref{thm:all-cm-central-extension}, together with the
relative quotient description of \(B\).
The statement includes the infinite collection of all prime degrees.
It does not require a first-layer test or determine the individual
full images \(C_p\). The unknown relative central classes are now
the sole obstruction to extending intrinsic quadratic characters.

At degree \(163\), the possible central obstruction has a particularly
small target.

\begin{theorem}\label{thm:main-163-odd-inertia}
Let \(R=R_{\rm global}\), \(P=P_{163}\), \(A=Q_{163}P\), and
\(U=\operatorname{Hom}_{\rm cts}(P/A,\mathbb F_2)\), with its
cyclotomic conjugation action. Then
\[
 H^j_{\rm cont}(G/R,U)=0\qquad(j\ge0).
\]
If \(h_{163}(R)=P\) and \(J=\ker(h_{163}:R\to P)\), a continuous
\(G\)-invariant quadratic character of \(J\) that is zero on
\(J\cap H\) extends to \(G\), zero on \(H\), exactly when its
relative central extension class is zero. The extension is unique.
The possible invariant relative class is the one-dimensional
quaternion line in \(H^2(P,A;\mathbb F_2)\).
\end{theorem}

The proof is Theorem~\ref{thm:163-concrete-odd-inertia} and
Corollary~\ref{cor:163-single-extension-obstruction}.
The full scalar image is either its arithmetic image \(A\) or \(P\).
The first case gives a further arithmetic normalization. The second
gives a quotient \(A_4\) whose arithmetic image has order three.
The theorem eliminates the conjugation obstruction in the latter
case; the relative central class remains to be determined.
A fixed finite-word detector expresses the full-image alternative.
Neither alternative is selected here, and
\(\operatorname{Hom}(B,\mathbb F_2)\) remains undetermined.



For every coefficient prime there is a corresponding restriction theorem.

\begin{theorem}\label{thm:main-primary-residual}
For a prime \(r\), let \(\Phi_p\) denote the free pro-\(p\) action and put
\[
 E_r=R_D\cap\bigcap_{p\ne r}\ker\Phi_p.
\]
This is a closed normal subgroup of infinite index in \(R_D\), and
\[
 B_{\max r}\cong
 \left(E_r/\overline{[G,E_r](E_r\cap H)}\right)_{\max r}.
\]
For odd \(r\), restriction from \(X_r\) to \(H^1(E_r,\mathbf F_r)\)
is injective. For \(r=2\), one has \(E_2=P_D\), and the kernel is
the rational Kummer span together with \(\mathbf F_2\eta_3\).
\end{theorem}

This restriction loses no original invariant characters at an odd
coefficient prime. Their values and their arithmetic restrictions
remain to be determined.

The smaller primary domains preserve the relevant relative quotient,
but their arithmetic character spaces are still nontrivial. The
following theorem gives a limitation of the combined realizations.

\begin{theorem}\label{thm:main-odd-realization-obstruction}
Let \(B_g(X,Y)\) be the full adelic Alexander multiplier, and put
\[
 J_B=D\cap\bigcap_r\ker\Phi_r\cap\ker T\cap\ker W\cap\ker B.
\]
Here \(T,W\) are the complete corrected cubic and harmonic CM
homomorphisms, including every coefficient prime. For every odd
prime \(p\), restriction
\[
             Y_p\longrightarrow H^1(J_B\cap H,\F_p)
\]
is injective. Its nonzero values can be witnessed by arithmetic
elements acting trivially on the integral first homology of every
finite abelian cover of the three-punctured projective line, at
every coefficient prime.
Consequently, if \(\chi\in X_p\) vanishes on \(J_B\), then
\(\chi|_K=0\). In particular any original character that factors
through the joint realizations defining \(J_B\) has zero arithmetic
restriction.
\end{theorem}

Theorems~\ref{thm:odd-target-survives-all-standard-kernels}
and~\ref{cor:odd-target-invisible-abelian-covers} prove this by
mixed arithmetic inertia and the alternating cup-product pairing.
The conclusion limits these realization methods; it does not show
that the odd arithmetic target characters cannot extend to full GT.

OpenAI proved the all-weight rational bound in the
Deligne--Drinfeld conjecture \cite[Theorem~1.1]{OpenAI}.
Willwacher's proof establishes the stronger characteristic-two
freeness theorem with explicit generators
\cite[Theorems~1.1 and~1.3, p.~2; \S4.3, p.~9]{Willwacher}.
That characteristic-two theorem is the Lie-algebra input here.
The arithmetic step constructs its generators by cyclotomic
conjugation, corrects their lower integral terms by arithmetic
commutators, and then uses compactness to identify the full dyadic
images. Markarian's restricted tangent reconstruction concerns
initial terms of solutions for his covering filtration
\cite[Theorem~4.8 and \S5, p.~10]{Markarian}; the arithmetic
realization in the theorem above uses the ordinary integral word
filtration and explicit Galois elements.

The normalization also controls every modular congruence representation
and identifies the two individual elliptic-frame signs.

\begin{theorem}\label{thm:main-congruence-frames}
Let $g=(1,f)\in\GT$ satisfy $\rho_2(g)=0$, and put
\[
 U=\begin{pmatrix}1&1\\0&1\end{pmatrix},\qquad
 V=\begin{pmatrix}1&0\\-1&1\end{pmatrix}.
\]
In the signed adelic braid representation,
\[
 f(U^2,V^2)=I\quad\text{in }\operatorname{SL}_2(\Zh).
\]
For the signed cubic and harmonic frame bits $a_3(g),a_2(g)$
of \cite[Lemmas 12--15]{H17}, with the positive signed lifts
specified in Section~\ref{sec:global-frame-character}, one has
\[
 a_2(g)=a_3(g)=\rho_3^I(g)/2\pmod2.
\]
In particular both signs vanish on $N_{\mathrm{rad}}$.
\end{theorem}

The first assertion says that every normalized element induces the
identity on every signed modular congruence quotient. The second says
that these two particular frame coordinates are already rational
Kummer data. Neither assertion identifies the action on the full
profinite modular group with its congruence action. The proofs are
Theorem~\ref{thm:individual-frame-kummer} and
Proposition~\ref{prop:normalized-congruence-triviality}.

The uniform dyadic theorem includes the following explicit finite layers.
They distinguish an arithmetic normalization from an invariant character
that extends to the original group. Put \(D_0=D\). In increasing order
\(j=3,5,7,8,9,10\), define \(D_j\) as the kernel of the coordinate
in the following table on its indicated domain. The odd coordinates are
\(m_j(g)=[S^{j-1}T]f_g\); the even ones are
\(m_8(g)=[S^4TST^2]f_g\) and \(m_{10}(g)=[S^5TST^3]f_g\).

\begin{theorem}\label{thm:main-finite-dyadic-examples}
Each coordinate is additive on its domain, has conjugation weight
\(\lambda_2^j\), and has the displayed image both on that domain
and on its intersection with \(H\).
\begin{center}
\begin{tabular}{@{}cccc@{}}
\toprule
Weight \(j\)&Domain&Exact image&First possible degree on \(D_j\)\\
\midrule
3&\(D\)&\(8\Z_2\)&5\\
5&\(D_3\)&\(32\Z_2\)&7\\
7&\(D_5\)&\(128\Z_2\)&8\\
8&\(D_7\)&\(\Z_2\)&9\\
9&\(D_8\)&\(512\Z_2\)&10\\
10&\(D_9\)&\(\Z_2\)&11\\
\bottomrule
\end{tabular}
\end{center}
Every \(D_j\) is proper in its preceding domain, closed and normal
in \(G\), and
\[
 G=D_jH,\qquad
 B\simeq D_j/\overline{[G,D_j](D_j\cap H)}.
\]
For every displayed \(j\), the kernel of restriction from \(X_2\)
to \(H^1(D_j,\F_2)\) is
\[
 \Sigma_0:=\F_2\eta_3\oplus
       \bigoplus_{q\text{ odd prime}}\F_2\psi_q.
\]
Moreover \(m_3(N)\subseteq8\Z_2\), and the joint map
\[
 N\longrightarrow
 \left(\prod_{q\text{ odd prime}}2\Zh\right)\times8\Z_2,
 \qquad n\longmapsto((\rho_q^I(n))_q,m_3(n))
\]
is surjective already on \(K\).

Divide each coordinate by the positive generator of its displayed
image and reduce modulo two, obtaining \(\eta_j\). These characters
have nonzero arithmetic restrictions. Their extension thresholds are
as follows: \(\eta_3\) extends to \(N\); \(\eta_5\) on \(D_3\)
does not extend invariantly to \(D\); \(\eta_7\) on \(D_5\)
does not extend invariantly to \(D_3\); \(\eta_8\) extends to
\(D_5\), but not to \(D_3\); \(\eta_9\) extends to \(D_7\),
but not to \(D_5\); and \(\eta_{10}\) extends to \(D_8\), but
not to \(D_7\). Each failure also excludes an invariant extension
to any larger preceding domain.

The joint map \((m_8,m_9):D_7\to\Z_2\times512\Z_2\)
is surjective already on \(D_7\cap H\), with kernel \(D_9\).
No nonzero combination of \(\eta_7\) and the extension of
\(\eta_8\) on \(D_5\) extends invariantly to \(D_3\), and
\[
 \ker\bigl(H^1(D_3,\F_2)^G\longrightarrow H^1(D_8,\F_2)\bigr)
 =\F_2\eta_5.
\]
\end{theorem}

Soul\'e's characters and their extensions to GT are classical;
Marin obtains the degree-three and degree-five extensions from the
Burau representation \cite[\S6.3, p.~780]{Marin2008}.
Sections~\ref{sec:cubic-magnus-normalization} through
\ref{sec:weight-ten-normalization} give the exact integral images,
arithmetic witnesses and extension obstructions in the displayed
normalizations. The mixed coordinates in even degree show why
odd depth-one coefficients alone do not describe the full dyadic action.

The odd metabelian coordinates admit a uniform continuation. For
$g\in\ker\lambda$, put $a_j(g)=[S^jT]f_g$ and
\[
 c_{2r+1}(g)=[u^{2r}]\sum_{j\ge1}a_j(g)(e^u-1)^j
 \qquad(r\ge1).
\]
These rational two-adic coordinates are the logarithmic boundary
coefficients of the dyadic beta multiplier.

\begin{theorem}\label{thm:main-uniform-dyadic-normalization}
Put $C_0=D=N_{\mathrm{rad}}$ and
\[
 C_r=D\cap\bigcap_{j=1}^r\ker c_{2j+1},\qquad
 C_\infty=\bigcap_{r\ge0}C_r.
\]
For every $r\ge1$, the coefficient
$m_{2r+1}=[S^{2r}T]f$ is additive on $C_{r-1}$, has conjugation
weight $\lambda_2^{2r+1}$, and satisfies
\[
 m_{2r+1}(C_{r-1})=m_{2r+1}(C_{r-1}\cap H)
                     =2^{2r+1}\Z_2,\qquad
 \ker(m_{2r+1}|_{C_{r-1}})=C_r.
\]
All these subgroups are closed and normal, and
\[
 G=C_rH=C_\infty H,\qquad
 B\simeq C_\infty/
       \overline{[G,C_\infty](C_\infty\cap H)}.
\]
The subgroup $C_\infty$ is exactly the kernel in $D$ of the dyadic
beta multiplier. The kernel of restriction from $X_2$ to
$H^1(C_\infty,\F_2)$ is the rational Kummer span together with
$\F_2\eta_3$.
Finally, $E=C_\infty\cap D_{10}$ has infinite index in $D_{10}$,
and
\[
 G=EH,\qquad B\simeq E/\overline{[G,E](E\cap H)}.
\]
Its elements have no pro-$2$ Magnus term below degree eleven and
have trivial dyadic beta multiplier. The kernel of restriction
from $X_2$ to $H^1(E,\F_2)$ is again exactly the rational Kummer
span together with $\F_2\eta_3$.
The mixed coordinate $b_{11}=[S^6TST^3]f$ on $E$ is additive,
has conjugation weight $\lambda_2^{11}$, and has full and arithmetic
image $\Z_2$. Its kernel $E'$ is a proper closed normal subgroup
whose words start in degree at least twelve,
$G=E'H$, and
\[
 B\simeq E'/\overline{[G,E'](E'\cap H)}.
\]
The kernel of restriction from $X_2$ to $H^1(E',\F_2)$ is the
same rational Kummer span together with $\F_2\eta_3$. More precisely,
if $F^n$ denotes the subgroup whose pro-$2$ Magnus words start
in degree $n$, then $E'=C_\infty\cap F^{12}$.
More generally put $E_n=C_\infty\cap F^n$ for $n=12,13,14,15$.
The following leading layers have exact arithmetic images:
\[
 E_{12}/E_{13}\simeq\Z_2^2,\qquad
 E_{13}/E_{14}\simeq\Z_2^2,\qquad
 E_{14}/E_{15}\simeq\Z_2^3.
\]
For each $n=12,13,14,15$, one has
\[
 G=E_nH,\qquad
 B\simeq E_n/\overline{[G,E_n](E_n\cap H)}.
\]
The kernel of restriction from $X_2$ to $E_n$ is still the
rational Kummer span together with $\F_2\eta_3$.
\end{theorem}

Section~\ref{sec:uniform-dyadic-normalization} proves this by combining
integer-valued polynomials with arithmetic cyclotomic differences.
Compactness passes from the finite stages to their intersection.
This identifies a smaller group on which the original relative
character question is decided; the character question itself remains
open. Ichimura and Kaneko characterize the arithmetic beta image at
odd primes in terms of Vandiver's conjecture
\cite[Section~1.2, pp.~314--315, Theorem~1]{IchimuraKaneko1989}.
Their theorem assumes an odd prime. The dyadic reduction here uses
the normalized duplication identity of Furusho and gives its exact
successive coefficient lattices directly.

There is a second arithmetic normalization using the two elliptic
points of the modular curve.  Let $G_1=\ker\lambda\subset\GT$.
Use the signed braid representation $q$, the unique pure frames
$A_g,\beta_g$ at the cubic and harmonic elliptic points, and the
normalized upper-triangular intertwiner $D_g$ defined in
Section~\ref{sec:adelic-cm-frames}.  Put
\[
 T_g=q(A_g)^{-1}D_g,\qquad W_g=q(\beta_g)^{-1}D_g
                       \quad(g\in G_1).
\]

\begin{theorem}\label{thm:main-cm-normalization}
The corrected pair $(T,W)$ has the same full adelic image on $G_1$
and on $G_{\Q^{\ab}}$.  Its kernel $C_{\mathrm{CM}}$ in $G_1$
is closed and normal in $\GT$.  If
$K_3=\Q(\sqrt{-3})$ and $K_4=\Q(i)$, then
\[
 \GT=C_{\mathrm{CM}}\GQ,\qquad
 C_{\mathrm{CM}}\cap\GQ=G_{K_3^{\ab}K_4^{\ab}},\qquad
 \GT/C_{\mathrm{CM}}\simeq
       \Gal(K_3^{\ab}K_4^{\ab}/\Q).
\]
Every continuous abelian character of $\GT$ which is trivial on
$\GQ$ and whose restriction to $G_1$ factors through $(T,W)$ is zero.
\end{theorem}

The complete adelic frame pair therefore contributes no further
relative abelian character.  The proof identifies its Galois image
by complex multiplication and the exact signed Tate-module comparison
on the Legendre family.  The two CM division fields intersect in
$\Q^{\ab}(3^{1/4})$.  The common field of their compositum and the rational-radical field is
\[
 L_{\mathrm{rad}}\cap K_3^{\ab}K_4^{\ab}
 =\Q^{\ab}\bigl(q^{1/12}:q\text{ prime},\,2^{1/8},3^{1/8}\bigr).
\]
Theorem~\ref{thm:radical-cm-intersection} proves this formula by an
explicit derived-subgroup calculation and division points.  Combining
the two normalizations is therefore equivalent to a specified crossed
cocycle with values in
$(\Z/12)^2\times\prod_{q\ne2,3}\Z/6$ being zero.
The equivalence is Proposition~\ref{prop:radical-cm-joint-gate}.
The entire coordinate belonging to $2$ vanishes: on every multiplier-one
gentle element the cubic and eighth-root actions of the frames both agree
with $\rho_2/2$, combined modulo twelve
(Corollary~\ref{cor:complete-two-compatibility}).
The remaining compatibility conditions have explicit finite forms.
On $D=N_{\mathrm{rad}}$, put
\[
 W_{g,3}\bmod3=(-I)^{w_3(g)},\qquad
 T_{g,2}\bmod8=(1+4\tau_3(g))I.
\]
Proposition~\ref{prop:lr-exceptional-three-character} proves that
$\chi_3=w_3+\tau_3$ is $\GT$-invariant and trivial on $D\cap H$;
it therefore defines an element of $\Hom(B,\F_2)$. Its vanishing
is exactly the exceptional eighth-root-of-three compatibility.
For each prime $\ell\ge5$, the other tests are
\[
 B_{g,\ell}(\omega,\omega)^{(\ell-\chi_{-3}(\ell))/3}=1,
 \qquad
 W_{g,\ell}^{(\ell-\chi_{-4}(\ell))/2}=1
\]
in the corresponding residue norm-one tori, for every $g\in D$.
Here $\omega$ is a primitive cube root, and
$\chi_{-3},\chi_{-4}$ are the two quadratic Dirichlet characters.
The wild cubic condition at three already holds.
Section~\ref{sec:lr-relative-gate} proves these equivalences, including
the integral cubic-frame comparison. The character $\chi_3$ is zero
by the two comparison theorems below. An explicit homomorphism from a
192-element group modulo eight to $\mathrm{SL}_2(\mathbf F_3)$
compares this character with two finite harmonic and cubic frame
defects (Proposition~\ref{prop:pair-of-gates}). Each defect is $\GT$-invariant and extends uniquely to a
quadratic character of $\GT$ trivial on $H$, and the two
extensions sum to $\chi_3$
(Proposition~\ref{prop:two-relative-extensions}).
The harmonic defect is zero on full GT by the following stronger
nonabelian identity.

\begin{theorem}[Harmonic comparison after Furusho]
\label{thm:main-harmonic-dyadic}
For every $g=(1,f_g)\in\GT$ with $\rho_2(g)=0$, the pure
harmonic frame satisfies
\[
 \beta_g(u^2,v^2)^{-1}=f_g(u^2,uvu)
 \quad\text{in }\widehat B_3.
\]
Consequently $w_3=\eta$ on $D$, the relative character
$\epsilon_h$ is zero, and $\chi_3=\epsilon_c=\tau_3+\eta$.
\end{theorem}

Section~\ref{sec:harmonic-dyadic-comparison} proves the identity
from Furusho's full dyadic word identities, the established
relation $(\mathrm{IV})$, and Enriquez's theorem on $(K_2)$. 

\begin{theorem}[Cubic comparison from genus-zero rigidity]
\label{thm:main-cubic-normalized}
For every $g=(1,f_g)\in\GT$ with $\rho_3^I(g)=0$, the pure
cubic frame satisfies
\[
 A_g(u^2,v^2)^{-1}=f_g(u^2,uv)
 \quad\text{in }\widehat B_3.
\]
Consequently $\epsilon_c=0$ and $\chi_3=0$ on full GT.
The exceptional eighth-root-of-three compatibility is therefore
satisfied. Joint radical--CM compatibility is equivalent to the
remaining tame cubic and harmonic tests displayed above.
\end{theorem}

Section~\ref{sec:normalized-cubic-identity} proves the comparison.
Boggi--Zalesskii's finite-subgroup and centralizer theorems,
Asada's genus-zero faithfulness and
Hoshi--Minamide--Mochizuki's description of the automorphisms
of the five-point group make the cyclic-cover quotient available
with its full profinite topology. The exact $(K_3)$ relation
identifies its GT parameter. Forgetting the two branch points
then gives the cubic comparison with all inner corrections
accounted for. The tame tests are not decided here.

A different finite homology representation gives a quadratic character
in the remaining relative kernel: it uses the ten character lines
$(\xi,\omega)$ with $\xi$ a nontrivial eleventh root of unity and
$\omega$ a primitive cube root, in characteristic seventeen, in the homology of the genus-ten curve
$u^{11}+v^3=1$.
Corollary~\ref{cor:concrete-radical-anderson-character} gives its
determinant formula. Its vanishing on the Galois subgroup is proved;
whether it is nonzero on $\GT$ remains open.

The harmonic residue tests have a direct interpretation in the
full adelic Alexander module. Let
\[
 B_g^{\rm ad}(X,Y)=1+(X-1)(Y-1)R_g^{\rm ad}(X,Y)
\]
be its normalized unit for $g\in N$. Evaluation at $(i,-1)$
is a continuous finite-character evaluation at each coefficient
prime, and does not require a formal power-series substitution
at an odd prime.

\begin{theorem}\label{thm:main-harmonic-beta-residue}
There is one integral identification of the harmonic quadratic order
with $\widehat{\Z[i]}$, valid for every $g\in N$ and every coefficient
prime, under which $W_g=B_g^{\rm ad}(i,-1)$. Reversing the choice of
primitive fourth root conjugates all these values simultaneously.
For each odd prime $\ell$, reduce this value
modulo $\ell$ and define
\[
 (-1)^{\vartheta_{\ell,4}(g)}
   =B_{g,\ell}^{\rm ad}(i,-1)^{(\ell-\chi_{-4}(\ell))/2}.
\]
Then $\vartheta_{\ell,4}\in X_2$, and
\[
 \Xi_{\ell,4}|_N=\vartheta_{\ell,4}+\rho_\ell^I/2\pmod2
\]
extends uniquely to a continuous quadratic character of $\GT$
that is zero on $H$. On $D$, its vanishing is exactly the tame
harmonic square condition. No vanishing or nonvanishing of these
characters on full $\GT$ is concluded.
\end{theorem}

Section~\ref{sec:harmonic-beta-square-lattice} proves the integral
comparison through a specified four-division construction. The
arithmetic value of the sign follows from the full Anderson
multiplication formula. Its logarithmic derivative alone does
not determine that constant sign. The order-three part of the
relative CM defect can separately be killed in the domain when
computing the maximal pro-$2$ relative abelian quotient; this uses
coprimality and makes no claim that the cubic defect vanishes.

The fixed finite-word detector for the degree-$163$ alternative has
the following complete form. Use $R=R_{\rm global}$, $P=P_{163}$,
$A=Q_{163}P$, and the scalar $h_{163}:R\to P$ defined above. Put
\[
 \mathcal O=\Z_2[\zeta_{163}],\qquad
 s(\zeta_{163})=\zeta_{163}^2,\qquad
 W=\ker\bigl(\operatorname{Tr}_{\F_8/\F_2}\bigr)
       \subset\mathcal O/2\mathcal O.
\]
Here $\F_8$ is the subfield of eight elements in
$\mathcal O/2\mathcal O$; $W$ is its additive trace-zero subspace,
not a field of order four. For $z\in P$ define
\[
 \Psi(z)=\sum_{j=0}^{26}s^{3j}
       \left((z-1)/2\bmod2\mathcal O\right).
\]
Let $\mathcal O_4=\mathcal O/4\mathcal O$ and define
\[
 \mathscr D=(\mathcal O_4)_{\rm add}\rtimes\langle t\rangle,
 \qquad t^{163}=1,\qquad tvt^{-1}=\zeta_{163}v.
\]
With $d=\zeta_{163}-1$, a unit, put
\[
 X=(0,t),\qquad Y=(d^{-1},t),\qquad n=2^{81}+1.
\]
For $0\le j\le26$, let $a_j$ be the least positive residue of
$2^{3j}$ modulo $163$.

\begin{theorem}[The degree-$163$ finite-word detector]
\label{thm:main-163-finite-word}
The homomorphism $\Psi:P\to W$ has kernel $A$. The elements $X,Y$
generate $\mathscr D$, giving one finite quotient of $\widehat F_2$
of order $163\cdot4^{162}$. For every $r\in R$, use the conjugator
$f_{r^n}$ of the group power $r^n$, with conventions
\[
 \phi_f(x)=x,\qquad \phi_f(y)=f^{-1}yf,\qquad
 [x,y]=xyx^{-1}y^{-1}.
\]
Then
\[
 \prod_{j=0}^{26}f_{r^n}(X^{a_j},Y^{a_j})
      =\bigl(2d^{-2}\Psi(h_{163}(r)),1\bigr)
      \quad\text{in }\mathscr D.
\]
In particular, this word is the identity if and only if
$h_{163}(r)\in A$. Its values lie in $2d^{-2}W$, an additive
subgroup of order four.
\end{theorem}

Proposition~\ref{prop:163-finite-word-detector} proves the formula.
The detector does not select either possible full scalar image.

\subsection*{Credit and scope}
The Kummer cocycles used here are Ihara's, and their full distribution relations are Enriquez's Theorem 8.1. The Galois description of rational radicals is classical Kummer theory in the explicit form of Moree and Stevenhagen's Theorem 2.5. The group-theoretic obstruction is the continuous Hochschild--Serre transgression. The structural use of normalized elements starts from \cite[Theorem 2]{H17}. The paper writes out the quotient, restriction, cyclic-cover parity and lattice arguments in the form needed for the character question. The kernel of restriction remains unresolved, as does the invariant character space for odd primes.

There are substantial earlier constructions of nonabelian representations
and quotients. Bleher, Chinburg and Lubotzky extend generalized-Jacobian
representations to finite-index subgroups of
$\operatorname{Aut}(\widehat F_2)$, including irreducible elliptic
examples \cite[Theorems 1.1--1.3]{BCL2022}. Bortnovskyi, Dolgushev,
Holikov and Pashkovskyi construct finite nonabelian quotients of full
$\GT$ from dihedral groups of two-power order, with surjective
arithmetic image. Their inverse limit is also arithmetically
surjective and is an index-two subgroup of
$\Z_2\rtimes\Z_2^\times$
\cite[Theorem 5.3, Corollary 5.4, Theorem 5.7 and Proposition 5.8]{BDHP2026}.
These specified representations and quotients do not identify the
relative central classes considered below; no such identification
is assumed here.

Basualdo Bonatto and Robertson give a modular-operad model of the
higher-genus tower and a faithful action of the Nakamura--Schneps
subgroup. Their object-fixing profinite endomorphisms are determined
by genus at most one \cite[Corollary~4.4 and Theorem~4.9]{BBR2026}.
The background equality of that subgroup with $\GT$ makes this
a full-GT action. The same source constructs its faithful action on
the associated tower of profinite homotopy one-types
\cite[Corollary~6.13 and Proposition~6.14]{BBR2026}.
These are statements about that specified model and its structure maps.

Lochak's Theorem 5.10 identifies the inertia-preserving automorphisms of generic procongruence mapping class groups with his group of compatible pairs \cite[Definition 5.4, Theorem 5.10]{Lochak2025}. Relation (IV) identifies every such pair with a diagonal pair, and the resulting group is $\GT$ \cite[Corollary 22]{H17}. This answers his closing questions about these groups. It does not by itself identify separately defined tower actions, or prove that preservation of nonseparating twist inertia implies preservation of every twist inertia.

\section{Proof overview}
The arithmetic targets are kernels of continuous transgression. At
prime two the affine Kummer translation contributes one extra cubic
line. The remaining classes come from cyclotomic cup products, whose
local Brauer invariants are cancelled by uniquely determined cyclic
extension classes. At odd primes the translation contributes no
cohomology: a central cyclotomic involution acts as minus one.
The cyclotomic eigenspace of the global Brauer sequence then gives
the explicit two-place corrections.

To lift the odd-prime-pair quadratic classes, take the compact homology
of the cyclic degree-\(\ell\) cover modulo \(q\). The peripheral word
makes its pairing a similitude of multiplier \(\lambda_q\), including
for nonarithmetic elements. The normalizer of its deck transformation
has an invariant torus sign. The commutator of a deck-multiplier lift
with a similitude lift gives the mixed cup coefficient; the powers
of the two lifts give exactly the two cyclic corrections. An explicit
finite representation therefore splits the corrected central extension
on full GT. Arithmetic transgression identifies its restriction.

The corrected cubic frame at coefficient two has a logarithmic
coordinate in \(\Z_2\). Its first bit is the sign on
\((2+\sqrt3)^{1/4}\); the degree-sixteen Galois extension computes
its corrected cup class. Exact arithmetic matching removes the whole
logarithmic coordinate. The full arithmetic target basis then proves
that all odd-pair characters remain independent on the resulting
arithmetic subgroup. Write \(R^*,V^*\) for the intermediate groups
defined immediately before
Theorem~\ref{thm:simultaneous-quadratic-normalization}: they are
obtained from \(R_D,V_D\) by killing \(\mathfrak t_2\) and all the
chosen odd-prime-pair lifts \(\xi_{\ell,q}\).
Compactness turns the independence into joint surjectivity onto their
infinite product, giving \(G=R^*H\).
The passage from \(R^*\) to \(V^*\) uses a pro-odd quotient and
preserves precisely the maximal pro-two relative obstruction.

At coefficient two, the same cover must be used integrally. Its
self-dual torus factors have Frobenius coinvariants equal to the four
squareclasses of \(\Z_2^\times\). The character of \(\sqrt2\)
on those coinvariants yields the missing corrected cup. A trace of the
integral logarithm computes its value on the square of a lift of
multiplier minus one; this is the necessary dyadic cyclic correction.
The resulting character already depends on homology modulo eight.
All arithmetic quadratic basis directions now lift. Applying the same
joint matching argument removes the last family, giving
\(R_{\rm quad}\) and \(V_{\rm quad}\). Every original character
on these residual groups is an arithmetic-trivial relative character.

For prime-degree cyclic covers, Milne's adelic CM theorem identifies
the local reflex norm. Its determinant is governed by Kummer's relative
class-number formula. A finite sign term and a Smith-normal-form test
decide whether the residual quotient is its natural first principal-unit
layer. Otsubo's Jacobi-sum character formula then identifies the dual of
the arithmetic image with square classes of explicit real units.
Mixed inertia commutators preserve every preceding normalization and
isolate one prime coordinate. Their axes and compactness give
simultaneous matching on \(\Pi_{\rm global}\).

At degrees twenty-nine and one hundred thirteen, one explicitly counted
Frobenius value and its conjugates fill the residual quotient of order
eight. At degree one hundred sixty-three, three cyclotomic-unit products
are negative squares, giving the exact smaller arithmetic image.
The integral pair \(P,QP\) determines a relative cohomology group whose
invariant central-extension space is a quaternion line. More generally,
the integral CM-type augmentation identity isolates every coordinate
of both full and arithmetic images, without requiring their equality.
The defining normalizations allow an odd inertia normal closure to be
followed through the complete radical quotient. Coprime cohomology
then kills every coordinate of the conjugation module. Continuous
cochains have finite coordinate support, so the same vanishing holds
for the entire infinite module. Only the relative central obstruction
remains in the exact extension test.
Finally, an Alexander-module evaluation realizes the residual image
test as one word in a fixed finite quotient.

The uniform dyadic theorem has an integral arithmetic construction.
Choose a Galois element with dyadic multiplier $-3$. Its word has
odd quadratic coefficient, whereas its power substitutions have no
quadratic or cubic correction modulo two. Commuting with it therefore
implements Willwacher's recursion for the odd free generators.
An induction removes the intervening even integral term by even
powers of arithmetic Hall commutators. These corrections preserve
the next reduced generator. Saturation of the integral equation
lattice and Nakayama's lemma then make the arithmetic leading
image the entire lattice in every degree.

To keep all rational Kummer coordinates zero simultaneously,
choose the arithmetic odd generators with Kummer vectors collinear
with one fixed vector. Their derived subgroup lies in the relative
commutator subgroup and realizes every beta-trivial leading term.
Nested compactness in this fixed arithmetic subgroup removes all
remaining Magnus coefficients. Thus the residual group is $R_D$,
and every original invariant character kills the arithmetic
corrections. This is why the stronger normalization preserves
both $B$ and the stated restriction kernel.

The arithmetic quotient is obtained by combining the cocycle with the cyclotomic character. A positive radical of $2$ lying in an abelian extension can have denominator at most two, so the Kummer translation subgroup over the cyclotomic field is exactly $2\Zh$. This determines the quotient, including its parity condition.

The extension argument uses the factorization $G=NH$: an invariant character of $N$ extends by zero on $H$ exactly when it vanishes on $N\cap H$. To identify the full abelianization one also uses Kronecker--Weber, which makes the image of $H^{\ab}$ a canonical section of the multiplier.

The dimension argument uses a different feature. A positive peripheral automorphism acts on each cyclic cover of the three-punctured sphere. In its class-two quotient, preservation of the boundary word makes the coefficients of the conjugator symmetric under reversal of the cyclic indices. Passing from a cover of degree $2d$ to one of degree $d$ proves the parity of $\rho_d^I$. For each odd rational prime, adjoining its fourth root to $K_2$ gives a nontrivial quadratic extension, and these extensions are independent. The divided cocycles give the desired lifts. Conjugation acts by the odd scalar $\lambda$, so their reductions modulo two are invariant. Compact metrizability gives the matching countable upper bound. This identifies an infinite subspace on which restriction is injective; it does not exhaust the invariant character space.

The simultaneous normalization uses the same arithmetic idea for all rational primes. A positive radical in an abelian extension of $\Q$ has at most two conjugates, since all its conjugates must be real. Finite-character separation then makes every joint Kummer image exactly the even lattice. Compactness prescribes all coordinates at once, so their common kernel still multiplies with $G_\Q$ to give $\GT$. Enriquez's distribution theorem identifies the cocycles at composite indices with the sums of those at prime indices. This shows that the entire rational Kummer character space restricts injectively; the relative obstruction is unchanged on their common kernel.

The odd coefficients are controlled together by the dyadic beta
multiplier. Its integral boundary series is the transform of a measure on
$\mathbf Z_2$. The two cycle relations give reflection, and Furusho's
normalized duplication identity, using relation (IV) and the exact
$(K_2)$ relation, gives its distribution under multiplication by two.
When the lower odd coefficients vanish, the measure has zero lower even
moments. Integrating an integer-valued polynomial then makes the next
coefficient divisible by exactly the required power $2^{2r+1}$.

To attain this upper bound arithmetically, choose a cyclotomic Galois
element whose Kummer coordinate of two is two and whose odd-prime
coordinates vanish. The boundary parity formula makes all its normalized
odd moments odd. Successive cyclotomic conjugation differences kill
the rational coordinates and then each lower odd coefficient. The exact
valuation of the next value follows from
$v_2(3^{2k}-1)=3+v_2(k)$. Thus every finite odd normalization has its
full image already on the Galois subgroup. Nested compactness gives the
same factorization for the common kernel of all odd coefficients,
which is exactly the trivial dyadic beta multiplier.

The first cases have images $8\mathbf Z_2$, $32\mathbf Z_2$ and
$128\mathbf Z_2$. Low-degree cycle calculations and the leading
pentagon then determine the entire Magnus leading spaces, an assertion
stronger than the vanishing of their one-letter-depth coordinates.
The cubic divided character extends to $N$ and is arithmetically
independent of the rational Kummer characters. The arithmetic
conjugation differences show why the next divided characters fail to
extend to earlier normalized groups.

At weight eight the normalized leading space is a primitive integral
line generated by an Ihara bracket. A commutator of two multiplier-one
Galois elements kills every preceding additive coordinate. Comparing
one element with a cyclotomic conjugate gives an odd value on the new
line, so its image is the full two-adic lattice. The same commutator
also obstructs extending its quadratic reduction to the original group.
The exact weight-seven divisibility allows that reduction to extend to
$D_5$. A cyclotomic difference and a complex-conjugation norm of arithmetic
elements give independent evaluations of this extension and $\eta_7$,
ruling out every nonzero combination as the restriction of an invariant
character on $D_3$.

At weight nine the uniform theorem gives the exact odd lattice.
Correction by the degree-eight commutator preserves the beta multiplier
and puts the arithmetic witness in $D_8$. A finite leading-space
calculation then gives $D_9$. At weight ten a primitive Ihara bracket
has an odd mixed coefficient, attained by a Galois commutator, giving
$D_{10}$. Both even arithmetic corrections have beta multiplier one.
They therefore combine with the infinite odd normalization to give
$E=C_\infty\cap D_{10}$ with the same arithmetic factorization.
One further commutator gives an odd value of a mixed degree-eleven
coordinate. Killing that coordinate gives $E'$. Every invariant
character of the original $N$ kills these correcting commutators,
so the original restriction kernel is unchanged at each step.
In degree twelve the two brackets of weights $(3,9)$ and $(5,7)$
have an odd determinant in two selected Magnus coordinates.
Arithmetic elements realize unit multiples of both columns, while
the same rational bound makes them the whole leading space.
Correcting both coordinates gives $E_{13}=C_\infty\cap F^{13}$.
The next two brackets, of types $(3,10)$ and $(5,8)$, have another
odd coefficient determinant. A nonzero arithmetic depth-one vector
separates the remaining rational direction in degree thirteen.
Removing the two beta-trivial directions gives
$E_{14}=C_\infty\cap F^{14}$.

For the finite harmonic comparison, Furusho's four-variable
dyadic identity specializes to a factorization of $f$ as a word
times the inverse of its interchange. The swap involution has
trivial fixed subgroup in the relevant free profinite group, so
this factor is exactly the pure harmonic frame. Specialization
to the signed braid group gives the identity directly, since the
central argument has total exponent zero. Its finite reduction removes the harmonic comparison
character. For the cubic comparison, the order-three centralizer
on an eleven-point sphere has the completed five-point quotient.
Genus-zero rigidity gives a GT parameter for that quotient.
Collision multitwists determine its multiplier and labelling, and
$(K_3)$ identifies it with the original parameter.
This extends the scope of the normalized cubic relation proved by
Schneps for her special-homomorphism subgroup
\cite[Theorem~3.4 and equation~(3.10), p.~667]{Schneps2006Special}. The other
filling then identifies the pure cubic conjugator. Exponent sums
remove the possible central discrepancy in the signed braid group.
Together the two comparisons remove the exceptional character at three.

For the CM normalization, correcting each pure frame by the same
upper-triangular intertwiner produces homomorphisms into the norm-one
tori of $\Q(\sqrt{-3})$ and $\Q(i)$.  Ordinary paths in the Legendre
family identify these signed integral matrices with the Tate actions
at its two CM fibres.  Class field theory and explicit division points
compute the two images and their common quadratic subfield.  The frame
sign equality on GT then forces its joint image to equal the Galois image.

\section{The arithmetic quotient}
\label{sec:arithmetic-quotient}

We first identify the quotient through which the multiplier and the Kummer
cocycle of~$2$ factor. Put
\[
 G=\widehat{\mathrm{GT}},\quad H=G_{\mathbb Q},\quad
 N=\Lambda _0,\quad
 K_2=\mathbb Q(\mu _\infty,2^{1/\infty}),\quad K=G_{K_2},
 \qquad U=\widehat{\mathbb Z}^{\times}.
\]
The embedding of $H$ in $G$ is understood throughout. We use the decomposition
$G=NH$ of \cite[Theorem~2]{H17}, the cocycle identity
\[
 \rho _2(gg')=\rho _2(g)+\lambda(g)\rho _2(g'),
\]
and the fact that $\rho _2$ restricts to the Kummer cocycle of the positive
coherent roots of~$2$; see \cite[Definition~5.1]{NS2000} and
\cite[Lemma~7]{H17}. Define the additive quadratic character
$\epsilon _8:U\to\mathbb F _2$ by
\[
 \epsilon _8(a)=
 \begin{cases}
 0,&a\equiv 1,7\pmod 8,\\
 1,&a\equiv 3,5\pmod 8.
 \end{cases}
\]

\begin{proposition}\label{prop:arithmetic-affine}
The subgroup $N$ is closed and normal in $G$, and $N\cap H=K$.
The map $(\lambda,\rho _2)$ identifies $G/N$ with
\[
 A=\{(a,r)\in U\ltimes\widehat{\mathbb Z}:
                 r\bmod 2=\epsilon _8(a)\},
 \qquad
 (a,r)(b,s)=(ab,r+as).
\]
This identifies $G/N$ canonically with $\operatorname{Gal}(K_2/\mathbb Q)$.
In particular $\rho _2\bmod 2=\epsilon _8\circ\lambda$ on $G$.
\end{proposition}

\begin{proof}
The cocycle identity makes $(\lambda,\rho _2)$ a continuous homomorphism to
the affine semidirect product. Its kernel is $N$, so $N$ is closed and normal.
On $H$ the simultaneous kernel fixes every root of unity and every coherent
root of~$2$, and is consequently $K$.

We need the precise Kummer image over $\mathbb Q^{\mathrm{ab}}
=\mathbb Q(\mu _\infty)$. It is a closed subgroup $J$ of
$\widehat{\mathbb Z}$. For a positive integer $d$, the inclusion
$J\subset d\widehat{\mathbb Z}$ is equivalent to
$2^{1/d}\in\mathbb Q^{\mathrm{ab}}$. This happens exactly for $d=1,2$.
Indeed $X^d-2$ is Eisenstein; if $d>2$, the field generated by its positive
real root is real and misses a nonreal conjugate, so is not normal over
$\mathbb Q$. Every finite subfield of an abelian Galois extension is normal.
Closed subgroups of $\widehat{\mathbb Z}$ have the form
$\prod_p p^{e_p}\mathbb Z_p$ with $e_p\in\mathbb N\cup\{\infty\}$,
and the inclusions in $d\widehat{\mathbb Z}$ determine all the $e_p$.
Thus $J=2\widehat{\mathbb Z}$.

Given $g\in G$, choose $h\in H$ with $\lambda(h)=\lambda(g)$.
The element $gh^{-1}$ has multiplier one, and
\cite[Lemma~7(c)]{H17} gives $\rho _2(gh^{-1})\in2\widehat{\mathbb Z}$.
The cocycle identity therefore gives
$\rho _2(g)\equiv\rho _2(h)\pmod2$. On $H$ this parity is the action
on $\sqrt2=\zeta _8+\zeta _8^{-1}$, namely $\epsilon _8(\lambda(h))$.
For each fixed multiplier, the Kummer values on $H$ fill the corresponding
coset of $2\widehat{\mathbb Z}$, by the preceding calculation. Thus the
image of $H$ is exactly $A$. The decomposition $G=NH$ shows that $G$ has
the same image. Finally the continuous bijection $H/K\to G/N$ is a
homeomorphism, since its source is compact and its target Hausdorff.
\end{proof}

\begin{proposition}\label{prop:arithmetic-commutators}
Let $T=\{(1,2t):t\in\widehat{\mathbb Z}\}\subset A$.
Then
\[
 \overline{[A,A]}=T,\qquad
 \overline{[A,T]}=\{(1,4t):t\in\widehat{\mathbb Z}\}.
\]
Every finite $p$-group quotient of $A$ is abelian when $p$ is odd.
Every finite quotient of $A$ which is a $2$-group of nilpotency class at
most two has commutator subgroup of order at most two.
\end{proposition}

\begin{proof}
The quotient $A/T=U$ is abelian. Choose $a\in U$ whose $2$-adic
component is~$3$ and whose odd-primary components are~$1$.
Then $(a,1)\in A$, and its commutator with $c=(-1,0)$ is $(1,2)$.
The closed subgroup generated by this translation is $T$, proving the first
equality. Conjugation by $(a,r)$ multiplies a translation coordinate by~$a$.
Thus $[A,T]$ consists of coordinates generated by $(a-1)2t$.
They are all divisible by~$4$ at the prime~$2$, and $a=-1$ produces
$(1,-4)$, whose closed cyclic subgroup is $4\widehat{\mathbb Z}$.
This proves the second equality.

The element $c$ has order two and acts on $T$ by inversion. It maps to the
identity in an odd-primary finite group, so the image of $T$ in such a
group has exponent dividing two and is trivial. The quotient is therefore
abelian. A quotient of nilpotency class at most two kills $[A,T]$, and its
commutator subgroup is consequently a quotient of
$T/\overline{[A,T]}\simeq C_2$.
\end{proof}

\section{Characters and the cyclotomic abelianization}
\label{sec:characters-abelianization}

The decomposition $G=NH$ separates the characters obtained from the multiplier
from the characters that could distinguish $G$ from $H$. All characters
in this section are continuous, and all abelianizations are profinite.

\begin{lemma}\label{lem:character-gluing}
Let $G=NH$ be a profinite group with $N$ closed and normal and $H$ closed,
and put $K=N\cap H$. Let $M$ be a finite abelian group with trivial action.
An invariant character $\psi:N\to M$ and a character $\eta:H\to M$
extend to a character of $G$ if and only if $\psi|_K=\eta|_K$.
The extension is unique and is given by
\[
 \chi(nh)=\psi(n)+\eta(h).
\]
\end{lemma}

\begin{proof}
Necessity follows by restriction. If $nh=n'h'$, then
$n'^{-1}n=h'h^{-1}\in K$, so the proposed two values coincide.
The identity
\[
 (n_1h_1)(n_2h_2)=n_1(h_1n_2h_1^{-1})h_1h_2
\]
and invariance of $\psi$ prove additivity. The multiplication map
$N\times H\to G$ is a continuous surjection from a compact space to a
Hausdorff space and hence is a quotient map. The displayed function on
$N\times H$ is continuous and constant on its fibres, proving continuity
on $G$. Uniqueness follows from $G=NH$.
\end{proof}

In our arithmetic situation define
\[
 C_N=\overline{\langle[g,n],k:g\in G,\ n\in N,\ k\in K\rangle},
 \qquad B=N/C_N.
\]
The group $B$ is abelian, since $C_N$ contains $[N,N]$; it is also the
largest quotient of $N$ on which conjugation by $G$ and the image of $K$
are both trivial.

\begin{theorem}\label{thm:abelian-splitting}
There is a canonical isomorphism of profinite abelian groups
\[
 G^{\mathrm{ab}}\simeq U\times B
\]
under which the multiplier is the first projection. Moreover
\[
 C_N=N\cap\overline{[G,G]},\qquad
 B\simeq\bigl(G/\overline{\langle\!\langle H\rangle\!\rangle}\bigr)^{\mathrm{ab}}.
\]
For every prime $p$, put
\[
 X_p=H^1(N,\mathbb F_p)^G,
 \qquad r_p:X_p\longrightarrow H^1(K,\mathbb F_p).
\]
Restriction gives canonical identifications
\begin{align*}
 \ker r_p&\simeq\operatorname{Hom}_{\mathrm{cts}}(B,\mathbb F_p),\\
 H^1(G,\mathbb F_p)&\simeq
 \lambda^*H^1(U,\mathbb F_p)\oplus\ker r_p.
\end{align*}
Consequently $\ker r_p=0$ if and only if every continuous
$\mathbb F_p$-character of $G$ factors through the multiplier.
The multiplier identifies $G^{\mathrm{ab}}$ with $U$ if and only if
$N\subset\overline{[G,G]}$, and this holds if and only if $\ker r_p=0$
for every prime~$p$.
\end{theorem}

\begin{proof}
Kronecker--Weber identifies $H^{\mathrm{ab}}$ with $U$ through the
cyclotomic character. Since $K$ fixes $\mathbb Q^{\mathrm{ab}}$, it lies
in $\overline{[H,H]}$ and hence in $\overline{[G,G]}$.
The map $N\to G^{\mathrm{ab}}$ therefore factors through $B$.
Conversely the rule
\[
 nh\longmapsto(\lambda(h),n\bmod C_N)
\]
is well-defined: changing the decomposition changes $n$ by an element of
$K$, which is zero in $B$. It is continuous by the compact quotient
argument of Lemma~\ref{lem:character-gluing}. It is a homomorphism because
conjugation by $H$ is trivial on $B$. The maps
$H^{\mathrm{ab}}=U\to G^{\mathrm{ab}}$ and $B\to G^{\mathrm{ab}}$
give the inverse of the resulting homomorphism
$G^{\mathrm{ab}}\to U\times B$: both composites are the identity on the
images of $N$ and $H$, which generate the relevant groups.
This proves the splitting and the formula for $C_N$.
Killing $H$ and then abelianizing imposes exactly the relations defining
$C_N$, which proves the second description of~$B$.

Apply Lemma~\ref{lem:character-gluing} with $\eta=0$ to identify the
restriction kernel. Every character of $H$ is the pullback of a unique
character of $U$ and vanishes on $K$. Subtracting that pullback from a
character of $G$ gives the asserted direct sum, or it follows directly
from the abelian splitting. Finally a nonzero profinite abelian group has
a nonzero finite abelian quotient, and that quotient has a quotient
$C_p$ for some prime $p$. Thus $B=0$ if and only if it has no character
to $\mathbb F_p$ for any prime~$p$.
\end{proof}

The expression ``the abelianization is $U$'' in this theorem specifies
the multiplier map. The kernel assertion does not follow from an
unmarked abstract isomorphism of profinite abelian groups.

\subsection{The transgression describes the image}

The kernel of restriction and the image of restriction involve different
obstructions. The next proposition states both in one exact sequence.
Write
\[
 Y_p=H^1(K,\mathbb F_p)^H.
\]
Restriction $r_p$ lands in $Y_p$, and $A=G/N=H/K$ acts on $K$ through
conjugation by $H$.

\begin{proposition}\label{prop:transgression-comparison}
The transgression
$d_H:Y_p\to H^2(A,\mathbb F_p)$ is injective and identifies $Y_p$ with
\[
 \ker\bigl(H^2(A,\mathbb F_p)\longrightarrow H^2(H,\mathbb F_p)\bigr).
\]
There is an exact sequence
\[
 0\longrightarrow\operatorname{Hom}_{\mathrm{cts}}(B,\mathbb F_p)
 \longrightarrow X_p
 \xrightarrow{\ d_Hr_p\ }
 \ker\bigl(H^2(A,\mathbb F_p)\longrightarrow H^2(G,\mathbb F_p)\bigr)
 \longrightarrow0.
\]
Thus, under the injection $d_H$, the image of restriction is precisely
the inflation kernel for~$G$.
\end{proposition}

\begin{proof}
For completeness, we describe the central-extension form of the
five-term argument. Given a conjugation-invariant character
$\psi:N\to\mathbb F_p$, its pushout extension is
\[
 E_\psi=(\mathbb F_p\times G)/
       \{(-\psi(n),n):n\in N\}.
\]
The subgroup being divided out is closed and normal, by the invariance
of $\psi$. The quotient is a central extension of $A$ by $\mathbb F_p$;
its class is $d_G\psi$, with a consistent choice of the transgression
sign. A splitting of this extension is equivalent to a character of
$G$ extending $\psi$: in one direction subtract the splitting from the
natural lift $g\mapsto[0,g]$, and in the other use the character to
define a splitting independent of the lift from~$A$.

More generally, a central extension of $A$ becomes trivial after
inflation to $G$ if and only if it admits a continuous homomorphism
$G\to E$ lifting $G\to A$. Restricting such a lift to $N$ gives an
invariant character $\psi$, and the map
$[z,g]\mapsto z\,\widetilde g$ identifies its pushout with $E$.
This proves exactness at $X_p$ and identifies the image of $d_G$ with
the indicated inflation kernel. The construction respects sums of
characters and Baer sums of central extensions, so these are linear
maps.

Apply the same argument to $H$ and $K$. A character of $K$ in the
kernel of $d_H$ extends to $H$. Every character of $H$ factors through
$U$ and therefore vanishes on $K$, so $d_H$ is injective. The preceding
lifting argument gives its image. Finally pulling the pushout extension
back from $G/N$ along $H/K$ is the same construction with $\psi|_K$.
Thus $d_G=d_Hr_p$. Theorem~\ref{thm:abelian-splitting} identifies the
kernel of $d_G$, proving the displayed exact sequence.
\end{proof}

\section{The arithmetic invariant characters}
\label{sec:arithmetic-target}

The target allowed by Galois invariance is already infinite dimensional.
Its size does not determine which characters extend across $K\subset N$.

\begin{theorem}\label{thm:arithmetic-target-dimension}
For every prime $p$,
\[
 \dim_{\mathbb F_p}H^1(G_{K_2},\mathbb F_p)^{G_{\mathbb Q}}=\aleph_0.
\]
For $p=2$, an explicit independent family is given, for the odd rational
primes $q$, by the quadratic extensions $K_2(q^{1/4})/K_2$.
Equivalently these are the Kummer classes of $\sqrt q$ in
$K_2^\times/(K_2^\times)^2$.
\end{theorem}

\begin{proof}
For an odd prime $q$, the square of the positive fourth root $q^{1/4}$
belongs to $\mathbb Q^{\mathrm{ab}}\subset K_2$. The extension in the
statement is therefore at most quadratic. It is normal over $\mathbb Q$,
because its other conjugates are $i^j q^{1/4}$ and $i\in K_2$.
Its character, if nonzero, is invariant under $H$, because conjugation
acts trivially on a group of order two.

Suppose that a nonempty finite product of these characters is zero.
Let $a$ be the product of the corresponding distinct odd primes. Then
$a^{1/4}\in K_2$. The extension
$\mathbb Q^{\mathrm{ab}}(a^{1/4})/\mathbb Q^{\mathrm{ab}}$ is a
nontrivial quadratic extension. To see nontriviality, $X^4-a$ is
Eisenstein at any prime dividing $a$, and its positive real root
generates a nonnormal real field, which cannot lie in
$\mathbb Q^{\mathrm{ab}}$. By Proposition~\ref{prop:arithmetic-affine},
$\operatorname{Gal}(K_2/\mathbb Q^{\mathrm{ab}})=2\widehat{\mathbb Z}$
has exactly one subgroup of index two. Its quadratic subfield is
$\mathbb Q^{\mathrm{ab}}(2^{1/4})$. The two quadratic fields are
therefore equal, so the product $a^{1/4}2^{1/4}=(2a)^{1/4}$ lies in
$\mathbb Q^{\mathrm{ab}}$: their common nontrivial automorphism negates
both factors. But $X^4-2a$ is Eisenstein at~$2$, and its positive root
again generates a real nonnormal field. This contradiction proves
nontriviality and independence of the whole quadratic family.

Now let $p$ be odd and let $m\geq1$. Put
$P=\operatorname{UT}_3(\mathbb F_p)^m$, a direct product of upper
unitriangular matrix groups. Its centre and commutator subgroup agree
and are isomorphic to $(C_p)^m$, and its abelianization is
$(C_p)^{2m}$. By \v{S}afarevi\v{c}'s theorem, in the precise form
\cite[Theorem~1, p.~1]{SchmidtWingberg}, there is a finite Galois
extension $E/\mathbb Q$ with group~$P$.
Let $L=E\cap K_2$. Its Galois group over $\mathbb Q$ is a finite
$p$-group quotient of $A$, and is abelian by
Proposition~\ref{prop:arithmetic-commutators}. Hence
$L\subset E^{[P,P]}$. The reverse inclusion follows from
$E^{[P,P]}\subset\mathbb Q^{\mathrm{ab}}\subset K_2$. Thus
\[
 \operatorname{Gal}(EK_2/K_2)=\operatorname{Gal}(E/L)=[P,P]
                  \simeq(C_p)^m.
\]
This subgroup is central in $\operatorname{Gal}(EK_2/\mathbb Q)$,
because $[P,P]$ is central in~$P$. Its characters give $m$ independent
members of $Y_p$.

Finally the algebraic field $\overline{\mathbb Q}$ is countable.
Every cyclic degree-$p$ extension of $K_2$ inside it is generated by
one element, so there are at most countably many such extensions.
Each has finitely many characters. This gives the countable upper
bound, while the preceding arguments give arbitrarily large finite
independent families for every~$p$.
\end{proof}

\section{Infinitely many invariant quadratic characters of the normalized group}
\label{sec:quadratic-dimension}

Ihara's Kummer cocycles give explicit invariant quadratic characters and
determine the dimension of their restriction image. They do not settle
the kernel of restriction. We write
$\rho_d^I$ for these cocycles to distinguish their construction from the
frame cocycle used to define~$N$.

Let $d$ be a positive integer and put
\[
 F=\widehat F_2=\langle x,y\rangle,\qquad
 H_d=\ker(F\longrightarrow\mathbb Z/d\mathbb Z),
 \quad x\longmapsto1,\quad y\longmapsto0.
\]
The Schreier basis of $H_d$ consists of $x^d$ and
$y_j=x^jyx^{-j}$ for $0\leq j<d$. Let
$\ell_d:H_d\to\widehat{\mathbb Z}$ send $y_0$ to~$1$ and the other
basis elements to~$0$. For $g=(\lambda,f)\in\GTz$ define
\[
 \rho_d^I(g)=-\ell_d(f).
\]
This is Ihara's cocycle $-\Psi_d^{(0)}$; the construction and its
identification with the Kummer cocycle of $d$ on $G_{\mathbb Q}$ are
recalled in \cite[\S2.2]{Furusho2003}, following Ihara's theory
\cite{Ihara1999}. We use that arithmetic identification as a published
input. The following direct verification fixes the cocycle convention.

\begin{lemma}\label{lem:ihara-cocycle}
For $g,g'\in\GTz$,
\[
 \rho_d^I(gg')=\rho_d^I(g)+\lambda(g)\rho_d^I(g').
\]
If $n\in N$, then
\[
 \rho_d^I(gng^{-1})=\lambda(g)\rho_d^I(n).
\]
\end{lemma}

\begin{proof}
The automorphism $\Phi_g:x\mapsto x^\lambda$,
$y\mapsto f^{-1}y^\lambda f$ preserves $H_d$. On its abelianization,
\[
 [x^d]\longmapsto\lambda[x^d],\qquad
 [y_j]\longmapsto\lambda[y_{\lambda j\bmod d}].
\]
The two conjugating contributions of $f$ cancel. Replacing an exponent
$\lambda j$ by its residue modulo $d$ only conjugates by a power of
$x^d$, which is also invisible in the abelianization. Since $\lambda$
is a unit modulo $d$, index zero has no preimage other than zero.
Consequently $\ell_d\Phi_g=\lambda\ell_d$.
With the convention $\Phi_{gg'}=\Phi_g\Phi_{g'}$, the second coordinate
is $f_{gg'}=f_g\Phi_g(f_{g'})$. Applying $-\ell_d$ proves the first
formula. Applying it twice to $gng^{-1}$, and using $\lambda(n)=1$,
cancels the terms involving $\rho_d^I(g)$ and proves the second.
\end{proof}

\begin{lemma}\label{lem:cyclic-cover-symmetry}
Let $F=\widehat F_2=\langle x,y\rangle$. Suppose a continuous
automorphism $\phi$ satisfies
\[
 \phi(x)=x,\qquad \phi(y)=f^{-1}yf,\qquad
 f\in\overline{[F,F]},
\]
and $\phi(xy)$ is conjugate to $xy$. For a positive integer $n$, use
the free basis
\[
 X=x^n,\qquad Y_j=x^jyx^{-j}\quad(0\leq j<n)
\]
of $H_n$. If $[f]=a[X]+\sum_{j=0}^{n-1}c_j[Y_j]$ in
$H_n^{\mathrm{ab}}$, then
\[
 a=0,\qquad \sum_{j=0}^{n-1}c_j=0,\qquad
 c_j=c_{-j\bmod n}\quad(0\leq j<n).
\]
\end{lemma}

\begin{proof}
The map $H_n^{\mathrm{ab}}\to F^{\mathrm{ab}}$ takes $[X]$ to
$n[x]$ and every $[Y_j]$ to $[y]$. Hence $na=0$ and $\sum c_j=0$.
Multiplication by $n$ is injective on $\widehat{\mathbb Z}$, so $a=0$.
The automorphism preserves $H_n$, fixes $X$, and sends
\[
 Y_i\longmapsto f_i^{-1}Y_if_i,\qquad f_i=x^ifx^{-i}.
\]
Conjugation by $x$ cyclically permutes the $[Y_j]$ in the abelianization
and fixes $[X]$, so $[f_i]=\sum_jc_{j-i}[Y_j]$, with indices modulo $n$.

The boundary word
\[
 P=Y_0Y_1\cdots Y_{n-1}X=(yx)^n
\]
is a power of a conjugate of $xy$, and therefore
$\phi(P)=tPt^{-1}$ for some $t\in F$. Multiplying $t$ on the right
by a suitable integer power of $yx$ makes $t$ belong to $H_n$ without
changing this equality: $yx$ centralizes $P$ and generates $F/H_n$.

Work in the quotient of $H_n$ of nilpotency class two, and use the
commutator convention $[u,v]=uvu^{-1}v^{-1}$. The commutator pairing
identifies its degree-two group with the exterior square of
$H_n^{\mathrm{ab}}$ over $\widehat{\mathbb Z}$. In particular the
commutators of distinct basis elements are independent. This independence
can also be checked in finite unitriangular $3$-by-$3$ matrix groups:
map one chosen pair of generators to the two elementary generators and
all other generators to the identity, for every finite modulus.
Using the same symbols for degree-one vectors, the boundary identity gives
\[
 \sum_{i=0}^{n-1}[Y_i,[f_i]]
       =\left[[t],X+\sum_{i=0}^{n-1}Y_i\right].
\]
Here each outer bracket is the bilinear commutator pairing; $[f_i]$
and $[t]$ denote degree-one vectors. Write
$[t]=t_XX+\sum_i t_iY_i$. The coefficient of $[X,Y_i]$ on the left
is zero because every $[f_i]$ has zero $X$-coordinate. On the right it
is $t_X-t_i$. Thus every $t_i=t_X$, and the right side is zero.
For $i<j$, the left side has coefficient $c_{j-i}-c_{i-j}$ at
$[Y_i,Y_j]$. Taking $i=0$ gives the asserted symmetry.
\end{proof}

\begin{theorem}\label{thm:all-kummer-parity}
For every positive integer $d$, Ihara's integral Kummer cocycle is even
on the multiplier-one subgroup of the gentle group:
\[
 \rho_d^I\bigl(\ker(\lambda:\widehat{\mathrm{GT}}_0\to U)\bigr)
       \subset2\widehat{\mathbb Z}.
\]
On the whole gentle group it satisfies
\[
 \rho_d^I\bmod2=\epsilon_d\circ\lambda,
\]
where $\epsilon_d:U\to\mathbb F_2$ is the additive quadratic
Dirichlet character of $\mathbb Q(\sqrt d)$.
\end{theorem}

\begin{proof}
The two cycle equations at multiplier one give the positive peripheral
preservation in Lemma~\ref{lem:cyclic-cover-symmetry}. Apply that lemma
with $n=2d$. Pairing indices $j$ and $-j$ gives
\[
 0=\sum_{j=0}^{2d-1}c_j
   =c_0+c_d+2\sum_{j=1}^{d-1}c_j.
\]
Under $H_{2d}\subset H_d$, only $Y_0$ and $Y_d$ among these basis
elements have nonzero $y$-coordinate in $H_d^{\mathrm{ab}}$, and that
coordinate is one for each. Indeed $Y_{j+d}=x^dY_jx^{-d}$.
The element $x^{2d}$ has zero $y$-coordinate as well. Consequently
\[
 \rho_d^I=-\ell_d(f)=-(c_0+c_d)\in2\widehat{\mathbb Z}.
\]
For a general gentle element $g$, choose $h\in G_{\mathbb Q}$ with
$\lambda(h)=\lambda(g)$. The cocycle identity and the first assertion
give $\rho_d^I(g)\equiv\rho_d^I(h)\pmod2$. The latter is the action
of $h$ on $\sqrt d$, and hence is $\epsilon_d(\lambda(g))$.
\end{proof}

\begin{theorem}\label{thm:quadratic-dimension}
The vector space of invariant continuous quadratic characters of the
normalized subgroup has dimension
\[
 \dim_{\mathbb F_2}H^1(\Lambda_0,\mathbb F_2)^{\widehat{\mathrm{GT}}}
 =\aleph_0.
\]
For every positive integer $d$ the formula
\[
 \psi_d^N=\frac{\rho_d^I|_N}{2}\bmod2
\]
defines a continuous invariant quadratic character. For the odd rational
primes $q$, the characters $\psi_q^N$ are linearly independent and
restrict to the characters of $K_2(q^{1/4})/K_2$. Thus
\[
 \dim_{\mathbb F_2}\operatorname{im}r_2=\aleph_0.
\]
More precisely, for distinct odd primes $q_1,\ldots,q_m$, let
$R:N\to\mathbb Z_2^m$ be the vector of the $2$-adic components of
$\rho_{q_1}^I,\ldots,\rho_{q_m}^I$. Then
\[
 R(K)=R(N)=2\mathbb Z_2^m,
\]
and restriction is injective on the $m$-dimensional character space
obtained by pullback from this lattice.
\end{theorem}

\begin{proof}
Theorem~\ref{thm:all-kummer-parity} permits division by two on $N$.
Since $\lambda=1$ there, the result is a continuous homomorphism.
The conjugation formula of Lemma~\ref{lem:ihara-cocycle} and the fact
that every multiplier is odd make its reduction modulo two invariant.
On $K$, the character $\psi_q^N$ is the ordinary fourth-root Kummer
character because $K$ fixes $\sqrt q$. Independence follows from
Theorem~\ref{thm:arithmetic-target-dimension}.

The vector $R$ is a continuous homomorphism. On $K$ it is the vector of ordinary Kummer
characters of the $q_j$. All their values are even, since $K$ fixes
$\sqrt{q_j}$. Modulo four, half their values are precisely the
independent quadratic characters of
Theorem~\ref{thm:arithmetic-target-dimension}. Thus the image of $R(K)$
modulo $4\mathbb Z_2^m$ is all of
$2\mathbb Z_2^m/4\mathbb Z_2^m$.

A closed additive subgroup of $\mathbb Z_2^m$ is a
$\mathbb Z_2$-submodule, because integer multiplication extends by
continuity to $2$-adic scalar multiplication. A closed submodule
$M\subset2\mathbb Z_2^m$ whose reduction modulo four is onto equals
$2\mathbb Z_2^m$. Explicitly, approximate any vector modulo four by
an element of $M$, then correct the error modulo eight by twice an
element of $M$, and continue. The convergent sum belongs to $M$ by
closedness. Apply this to the compact image $R(K)$.

The parity theorem gives $R(N)\subset2\mathbb Z_2^m$, and the reverse
inclusion follows from $K\subset N$. Hence $R(N)=R(K)$.
Conjugation acts on this lattice by the odd scalar $\lambda(g)$, so
all its quadratic characters are invariant. Its character space has
dimension $m$, and restriction to $K$ is injective because $R(K)=R(N)$.

There is also a countable upper bound. The finitely generated free
profinite group $\widehat F_2$ has a countable base, as there are only
countably many finite quotients with specified generator images.
Thus $U\times\widehat F_2$, its closed subspace $G$, and the closed
subgroup $N$ are compact metrizable spaces. Such a space has countably
many clopen subsets: each is a finite union of members of a countable
base. A continuous function to $\mathbb F_2$ is determined by its
zero fibre. Consequently $H^1(N,\mathbb F_2)$, and hence its invariant
subspace, is countable. The lower bounds for arbitrary $m$ prove the
dimension formula.

The same upper bound applies to the restriction image, while the
independent quartic restrictions give its countably infinite dimension.
\end{proof}

The dimension computation identifies the explicit quartic subspace. To determine its complete extent, and to normalize all rational Kummer coordinates simultaneously, we next use the distribution theorem.

\section{The classical distribution relations and their arithmetic consequences}
\label{sec:kummer-distribution}

The multiplicative relations among Ihara's Kummer cocycles already hold
integrally on the full profinite Grothendieck--Teichm\"uller group.
They follow from Enriquez's distribution theorem, which is used here as
published background.

For the cyclic-cover subgroup $H_n\subset\widehat F_2$, define
\[
 d_n:H_n\longrightarrow\widehat F_2,\qquad
 x^n\longmapsto x,\quad y_0\longmapsto y,\quad
 y_j\longmapsto1\quad(0<j<n).
\]
Let $c_n=\operatorname{exp}_y\circ d_n$; thus $c_n=\ell_n$ in
our Schreier coordinates. Enriquez proves that, for every
$g=(\lambda,f)\in\widehat{\mathrm{GT}}$ and every $n\geq1$,
\begin{equation}
 d_n(f)=y^{c_n(f)}f.
 \label{eq:full-kummer-distribution}
\end{equation}
This is equation (50) and Theorem 8.1 of \cite{Enriquez2007}; in
particular his subgroup $\mathrm{GTK}$ is the full profinite group.
The map is the filling map on the cyclic cover, so this is an identity
in the full free profinite group, beyond its abelian or congruence images.

\begin{proposition}\label{prop:kummer-multiplicativity}
For positive integers $a,b$ and every $g\in\widehat{\mathrm{GT}}$,
\[
 \rho_{ab}^I(g)=\rho_a^I(g)+\rho_b^I(g).
\]
Consequently
\[
 \rho_n^I=\sum_{q\mid n}v_q(n)\rho_q^I,
 \qquad
 \frac{\rho_{ab}^I-\rho_a^I-\rho_b^I}{2}\bmod2=0.
\]
The last expression therefore supplies no nonzero character in the
restriction kernel.
\end{proposition}
\begin{proof}
On $H_{ab}$ one has $d_b(H_{ab})\subset H_a$ and
$d_a d_b=d_{ab}$. Indeed $x^{ab}$ maps successively to $x^a,x$;
for $y_j=x^jyx^{-j}$, the first map kills it unless $b\mid j$,
and otherwise sends it to $y_{j/b}$. The second map retains only
$j=0$, for $0\leq j<ab$.
Applying \eqref{eq:full-kummer-distribution} twice now gives
\[
 y^{c_{ab}(f)}f=d_{ab}(f)
   =d_a\bigl(y^{c_b(f)}f\bigr)
   =y^{c_b(f)+c_a(f)}f.
\]
The procyclic subgroup generated by $y$ is isomorphic to
$\widehat{\mathbb Z}$, so $c_{ab}(f)=c_a(f)+c_b(f)$.
This is also the observation in \cite[Remark 8.6]{Enriquez2007}.
Since $\rho_n^I=-c_n(f)$, the asserted formulas follow, with the
prime factorization formula obtained by iteration.
\end{proof}

The cocycle normalization agrees with the one used for the modular
frame: \cite[\S2.2, footnote on p.~100]{Furusho2003} identifies the
Nakamura--Schneps cocycle $\rho_2$ with $-\Psi_2^{(0)}$, and
\cite[Lemma 7(a) and Proposition 11]{H17} identifies that cocycle
with the pure-frame definition. Hence $\rho_2^I=\rho_2$ on
$\widehat{\mathrm{GT}}$. The negative-index Schreier generators in
Furusho's convention give the same coefficient $c_n$: for $0<j<n$,
$x^{-j}yx^j$ is conjugate within $H_n$ to $y_{n-j}$, so both
coefficient maps vanish on it.

\begin{corollary}\label{cor:dyadic-anderson}
Every element of $\widehat{\mathrm{GT}}$ satisfies all dyadic
Anderson multiplication relations $(A_{2^j})$, $j\geq1$, in the
notation of \cite[Definition 2.3]{Furusho2003}.
\end{corollary}
\begin{proof}
Relation (IV) holds on the group by \cite[Theorem 1]{H17}, and
$(K_2)$ holds by \eqref{eq:full-kummer-distribution}.
\cite[Theorem 3.1]{Furusho2003} therefore gives the duplication
relation; its dyadic iterates follow from that paper's Lemma 3.2.
Equivalently, combine
$\mathrm{GTK}=\widehat{\mathrm{GT}}$,
$\mathrm{I}\Gamma=\widehat{\mathrm{GT}}$
\cite[Corollary 21]{H17}, and
\cite[Theorem 3.4]{Furusho2003}.
\end{proof}

These are consequences of the cited distribution and multiplication
theorems. They do not determine the full restriction kernel: they show
that the displayed multiplicative defects vanish and impose the
specified dyadic identities.

\section{Simultaneous normalization of all rational Kummer coordinates}
\label{sec:rational-radical-normalization}

The arithmetic part of every rational Kummer coordinate can be removed
at once. This replaces the normalized subgroup by a proper closed
subgroup without changing the relative abelian obstruction.
The Galois calculation used here is the classical description of radical
extensions; see \cite[\S2, especially Theorem~2.5, author pp.~4--6]{MoreeStevenhagen}.
We give the finite-coordinate argument to make its simultaneous use and
the resulting GT quotient explicit.

\begin{lemma}\label{lem:simultaneous-kummer-image}
For distinct rational primes $q_1,\ldots,q_m$, the joint Kummer map
\[
 G_{\mathbb Q^{\mathrm{ab}}}\longrightarrow\widehat{\mathbb Z}^{\,m},
 \qquad \sigma\longmapsto
       (\kappa_{q_1}(\sigma),\ldots,\kappa_{q_m}(\sigma))
\]
has image exactly $2\widehat{\mathbb Z}^{\,m}$.
Consequently, for distinct odd primes $q_1,\ldots,q_m$, its restriction
to $K=G_{K_2}$ also has image $2\widehat{\mathbb Z}^{\,m}$.
\end{lemma}

\begin{proof}
Denote the first image by $J$. It is a compact, hence closed, subgroup
of $\widehat{\mathbb Z}^{\,m}$. Every continuous character of that
product with values in $\mathbb Q/\mathbb Z$ has the form
\[
 (t_1,\ldots,t_m)\longmapsto\frac1n\sum_{i=1}^m a_it_i\pmod{\mathbb Z}
\]
for integers $n>0$ and $a_i$. Such a character vanishes on $J$ exactly
when the positive real radical
\[
 \alpha=\prod_{i=1}^m q_i^{a_i/n}
\]
belongs to $\mathbb Q^{\mathrm{ab}}$. Indeed the Kummer action on
$\alpha$ multiplies it by the corresponding $n$th root of unity.

If $\alpha\in\mathbb Q^{\mathrm{ab}}$, its field over $\mathbb Q$
is normal and is contained in $\mathbb R$. Every conjugate of
$\alpha$ is a real root of $X^n-\alpha^n$, and is therefore one
of $\alpha,-\alpha$. It follows that $\alpha^2\in\mathbb Q$.
Writing this last rational number in its unique prime factorization
gives $2a_i/n\in\mathbb Z$ for each $i$. Conversely these
integrality conditions make $\alpha$ a square root of a positive
rational number, and hence an element of $\mathbb Q^{\mathrm{ab}}$.
Thus the annihilator of $J$ is precisely the group of characters
of exponent at most two. It is also the annihilator of
$2\widehat{\mathbb Z}^{\,m}$, so the two closed subgroups agree.
For this last implication, a proper closed subgroup of a profinite
abelian group is separated from any point outside it by a character
of a finite quotient.

For the second assertion, apply the first to the list
$2,q_1,\ldots,q_m$. Its image is the full product
$2\widehat{\mathbb Z}^{\,m+1}$. Every point of that product whose
first coordinate is exactly zero has a preimage fixing all roots
of two as well as all roots of unity. Such a preimage belongs to
$K$, proving the assertion with no compatibility condition left over.
\end{proof}

Let $\mathcal P$ be the set of rational primes and put
\[
 L_{\mathrm{rad}}=
 \mathbb Q\bigl(\mu_\infty,q^{1/n}:q\in\mathcal P,
                                      \ n\geq1\bigr).
\]
The roots in this expression are positive real roots; adjoining
$\mu_\infty$ also adjoins all their conjugates.
Define
\[
 N_{\mathrm{rad}}=N\cap
       \bigcap_{\substack{q\in\mathcal P\\q\ne2}}\ker\rho_q^I,
 \qquad K_{\mathrm{rad}}=G_{L_{\mathrm{rad}}}.
\]
The frame cocycle $\rho_2$ and $\rho_2^I$ agree on $G$;
the normalization is recorded in \cite[\S2.2, p.~100, footnote]{Furusho2003}
and \cite[Proposition~11]{H17}. Thus $N_{\mathrm{rad}}$ is the
simultaneous kernel of the multiplier and all rational-prime Kummer
cocycles.

\begin{theorem}\label{thm:rational-radical-normalization}
The joint map
\[
 \mathcal R:N\longrightarrow
       \prod_{\substack{q\in\mathcal P\\q\ne2}}2\widehat{\mathbb Z},
 \qquad n\longmapsto(\rho_q^I(n))_q
\]
is surjective already on $K$. Its kernel is $N_{\mathrm{rad}}$,
and
\[
 N=N_{\mathrm{rad}}K,\qquad
 G=N_{\mathrm{rad}}H,\qquad
 N_{\mathrm{rad}}\cap H=K_{\mathrm{rad}}.
\]
In particular $N_{\mathrm{rad}}$ is a proper closed normal subgroup
of $N$, and
\[
 N/N_{\mathrm{rad}}\simeq
       \prod_{\substack{q\in\mathcal P\\q\ne2}}2\widehat{\mathbb Z}.
\]
There are canonical topological isomorphisms
\begin{align*}
 G/N_{\mathrm{rad}}&\simeq
            \operatorname{Gal}(L_{\mathrm{rad}}/\mathbb Q),\\
 B&\simeq N_{\mathrm{rad}}/
       \overline{[G,N_{\mathrm{rad}}]K_{\mathrm{rad}}}.
\end{align*}
The first quotient is explicitly
\[
 A_{\mathrm{rad}}=
 \left\{(u,(t_q)_{q\in\mathcal P}):
 \begin{array}{l}
 u\in\widehat{\mathbb Z}^{\times},\quad
 t_q\in\widehat{\mathbb Z},\\
 t_q\equiv\epsilon_q(u)\pmod2\quad(q\in\mathcal P)
 \end{array}\right\},
\]
with multiplication $(u,t)(v,s)=(uv,t+us)$, where
$\epsilon_q$ is the quadratic character of $\mathbb Q(\sqrt q)$.
\end{theorem}

\begin{proof}
Each coordinate of $\mathcal R$ is a homomorphism on $N$ and is
even by Theorem~\ref{thm:all-kummer-parity}. The joint map is
continuous for the product topology. By
Lemma~\ref{lem:simultaneous-kummer-image}, every finite-coordinate
projection of $\mathcal R(K)$ is the whole corresponding product.
Hence $\mathcal R(K)$ is dense in the full product. It is compact
and therefore closed, proving
\[
 \mathcal R(K)=\mathcal R(N)=\prod_{q\ne2}2\widehat{\mathbb Z}.
\]
Equivalently, one can prescribe successively larger finite lists
of coordinates and use the finite intersection property in $K$.
This compactness step prescribes the whole profinite value of every
coordinate, not only its residue at one finite level.

Given $n\in N$, choose $k\in K$ with $\mathcal R(k)=\mathcal R(n)$.
Then $nk^{-1}\in N_{\mathrm{rad}}$, proving
$N=N_{\mathrm{rad}}K$ and then $G=N_{\mathrm{rad}}H$.
The conjugation identity for the Kummer cocycles shows that
$N_{\mathrm{rad}}$ is normal in $G$; it is closed by definition.
Its Galois intersection fixes precisely the roots of unity and
the roots of all rational primes, so it is $K_{\mathrm{rad}}$.
The quotient isomorphism $N/N_{\mathrm{rad}}\simeq\prod_{q\ne2}
2\widehat{\mathbb Z}$ follows from compactness. This nontrivial
quotient also proves that $N_{\mathrm{rad}}$ is proper.

The factorization $G=N_{\mathrm{rad}}H$ identifies
$G/N_{\mathrm{rad}}$ with $H/K_{\mathrm{rad}}$.
On the Galois group, each translation coordinate has parity
$\epsilon_q(u)$ because it records the action on $\sqrt q$.
For a fixed multiplier choose one Galois lift. The first assertion
of Lemma~\ref{lem:simultaneous-kummer-image}, followed by compactness
over all finite prime lists, allows an arbitrary even translation
vector to be added to this lift. Thus the image is exactly
$A_{\mathrm{rad}}$, with the displayed cocycle multiplication.

Finally apply the proof of Theorem~\ref{thm:abelian-splitting}
to $G=N_{\mathrm{rad}}H$. The Galois intersection
$K_{\mathrm{rad}}$ lies in $\overline{[H,H]}$, since it fixes
$\mathbb Q^{\mathrm{ab}}$. That proof identifies the new relative
abelian quotient with
$(G/\overline{\langle\!\langle H\rangle\!\rangle})^{\mathrm{ab}}$,
which is also $B$. This gives the asserted canonical isomorphism.
\end{proof}

\begin{corollary}\label{cor:rational-kummer-characters}
The invariant quadratic characters of $N$ that factor through all
rational Kummer coordinates form exactly
\[
 X_2^{\mathrm{Kum}}=
 \bigoplus_{\substack{q\in\mathcal P\\q\ne2}}
       \mathbb F_2\,\psi_q^N.
\]
Restriction to $K$ is injective on this subspace. For an odd prime
$p$, no nonzero invariant $\mathbb F_p$-character of $N$ factors
through the rational Kummer coordinates.
For every prime $p$, the unresolved relative kernel is equivalently
\[
 \operatorname{Hom}(B,\mathbb F_p)=
 \left\{\chi\in H^1(N_{\mathrm{rad}},\mathbb F_p)^G:
                       \chi|_{K_{\mathrm{rad}}}=0\right\}.
\]
\end{corollary}

\begin{proof}
Enriquez's distribution theorem gives
$\rho_{ab}^I=\rho_a^I+\rho_b^I$ on $G$ for positive integers
$a,b$; see \cite[Theorem~8.1 and Remark~8.6, author pp.~30--32]{Enriquez2007}.
Since $\rho_2^I$ is zero on $N$, all rational Kummer coordinates
on $N$ are generated by the odd-prime coordinates.
A continuous character of the product in
Theorem~\ref{thm:rational-radical-normalization} to a finite group
depends on finitely many coordinates: its open kernel contains
a basic product neighborhood of zero. Each factor
$2\widehat{\mathbb Z}$ has a one-dimensional quadratic character
space, generated by division by two followed by reduction modulo
two. This proves the direct-sum formula. Every multiplier is odd,
so every one of these characters is $G$-invariant. Surjectivity
already on $K$ proves injectivity of restriction.

Complex conjugation acts on the product by multiplication by $-1$.
If $p$ is odd and a character $\chi$ is invariant, then
$\chi(t)=\chi(-t)=-\chi(t)$, so it is zero. The final identification
is the character-gluing lemma applied to
$G=N_{\mathrm{rad}}H$.
\end{proof}

\begin{proposition}\label{prop:radical-character-pullback}
For every prime $p$, put
\[
 X_p^{\mathrm{rad}}=H^1(N_{\mathrm{rad}},\mathbb F_p)^G,
 \qquad
 Y_p^{\mathrm{rad}}=H^1(K_{\mathrm{rad}},\mathbb F_p)^H.
\]
The restriction maps identify $X_p$ with the fiber product
\[
 X_p\simeq X_p^{\mathrm{rad}}
                  \mathop{\times}_{Y_p^{\mathrm{rad}}}Y_p.
\]
The kernels of $X_p\to X_p^{\mathrm{rad}}$ and
$Y_p\to Y_p^{\mathrm{rad}}$ are identified by restriction.
They are zero for odd $p$, and for $p=2$ they are respectively
$X_2^{\mathrm{Kum}}$ and its restriction to $K$.
\end{proposition}
\begin{proof}
A character of $N$ restricts to the indicated compatible pair,
and the pair determines it because $N=N_{\mathrm{rad}}K$.
Conversely, let $\alpha\in X_p^{\mathrm{rad}}$ and
$\beta\in Y_p$ have the same restriction to $K_{\mathrm{rad}}$.
Define
\[
 \chi(nk)=\alpha(n)+\beta(k)
       \qquad(n\in N_{\mathrm{rad}},\ k\in K).
\]
If $nk=n'k'$, then $n'^{-1}n=k'k^{-1}$ lies in
$N_{\mathrm{rad}}\cap K=K_{\mathrm{rad}}$; the equality of
the two restrictions proves well-definedness. The invariance of
$\alpha$ under $K$ makes this formula a homomorphism.
It is continuous: the multiplication map from the compact space
$N_{\mathrm{rad}}\times K$ onto the Hausdorff group $N$ is
a quotient map, and its pullback is the continuous sum
$\alpha+\beta$. Conjugation by $H$ preserves the formula,
since $\alpha$ is $G$-invariant and $\beta$ is $H$-invariant.
Conjugation by $N$ preserves any abelian character of $N$.
Since $G=NH$, the resulting character belongs to $X_p$.

The kernel of restriction to $N_{\mathrm{rad}}$ is the invariant
character space of $N/N_{\mathrm{rad}}$; the corresponding kernel
on the Galois side is that of $K/K_{\mathrm{rad}}$.
The joint Kummer map identifies both quotients with
$\prod_{q\ne2}2\widehat{\mathbb Z}$, with the same scalar action
of the multiplier. The preceding corollary computes their
invariant character spaces, proving the final assertions.
\end{proof}

The subgroup $N_{\mathrm{rad}}$ gives a strictly smaller location
for the open obstruction. A character of $G$ which vanishes on
$H$ is zero if and only if its restriction to this subgroup is zero.
The theorem does not determine its relative abelianization.

\section{An integral measure for all odd coefficients}
\label{sec:integral-beta-measure}

The dyadic beta multiplier encodes all the odd one-letter-depth
coefficients in a single integral measure. Its symmetry and distribution
identity will give their exact uniform divisibility.

Put $M=\ker\lambda$ and $N=\ker(\lambda,\rho_2^I)$.
For $g\in M$ write
\[
 B_g(X,Y)=1+(X-1)(Y-1)R_g(X,Y),\qquad
 q_g(X)=(X-1)R_g(X,1)=\sum_{n\ge1}a_n(g)(X-1)^n,
\]
where $a_n(g)=[S^nT]f_g$.  All coefficients lie in $\mathbf Z_2$.
The equality with the single-$T$ Magnus coefficients follows by taking
the metabelian image: the ending-$T$ part of $[x,y]$ is $S$, conjugation
by $x$ multiplies it by $X$, and a product of two commutator terms has
at least two occurrences of $T$.

The beta cocycle and its symmetry follow from Ihara's factorization
\cite[Proposition~1.6.1]{Ihara1999}, recalled in
\cite[Section~2.3, p.~101]{Furusho2003}. The boundary reflection is
an elementary consequence of the two cycle relations:
\[
 B(X,Y)B(XY,Z)=B(Y,Z)B(X,YZ),\quad B(X,Y)=B(Y,X),
 \quad q(X^{-1})=q(X).
\]
Here is a direct check of that reflection. The metabelian forms of
relations (I) and (II), at multiplier one, are
\[
 R(X,Y)=R(Y,X),\qquad
 R(X,Y)+(XY)^{-1}R(Y,(XY)^{-1})
       +Y^{-1}R((XY)^{-1},X)=0.
\]
The factors in the second formula are the commutator coefficients
of $[y,(xy)^{-1}]$ and $[(xy)^{-1},x]$, respectively.
Specializing at $Y=1$ and at $X=1$, and using symmetry, gives
$(X-1)R(X,X^{-1})=0$. Since the one-variable power-series ring is
a domain, $R(X,X^{-1})=0$ and $R(X^{-1},1)=-XR(X,1)$.
Multiplication by $X^{-1}-1$ proves $q(X^{-1})=q(X)$.

For $g\in N$, the normalized duplication equation is
\begin{equation}
 B_g(X,X)=B_g(-1,X).\label{eq:measure-duplication}
\end{equation}
The geometric input is \cite[Theorem~3.1 and the displayed
$(A_2^0)$, p.~102]{Furusho2003}: relations (I), (II), (IV) and
$(K_2)$ imply $B(X,X)=X^{-2\Psi_2}B(-1,X)$.
The convention there is $\Psi_2=-\rho_2^I$
\cite[p.~100, footnote]{Furusho2003}.
The full profinite $(K_2)$ relation follows from
\cite[Section~8, Theorem~8.1, equation~(50), author p.~30]{Enriquez2007}.
Relation (IV) on full $G$ is the input cited in Section~1.

\begin{lemma}[Integral boundary measure]
\label{lem:integral-boundary-measure}
For every $g\in N$ there is an integral $\mathbf Z_2$-valued measure
$\nu_{g,+}$ on $1+4\mathbf Z_2$, of total mass zero, such that
\begin{equation}
 \int t^{2k}\,d\nu_{g,+}(t)
  =(1-2^{2k})\tau_{2k+1}(g)\qquad(k\ge1),
 \label{eq:integral-boundary-moments}
\end{equation}
where the normalization is exactly
\[
 c_{2k+1}(g)=[u^{2k}]q_g(e^u),\qquad
 \tau_{2k+1}(g)=\frac{(2k)!}{2}c_{2k+1}(g).
\]
The same formula defines an integral continuous homomorphism
$\tau_{2k+1}:M\to\mathbf Z_2$ with conjugation weight
$\lambda_2^{2k+1}$. Its value modulo two is independent of $k$:
\[
 \tau_{2k+1}(g)\equiv\tau_3(g)=a_2(g)\pmod2
 \qquad(g\in M).
\]
\end{lemma}
\begin{proof}
The isomorphism
$\mathbf Z_2[[\mathbf Z_2]]\simeq\mathbf Z_2[[X-1]]$ gives an integral
measure $\mu_g$ with transform $q_g(X)=\int X^t\,d\mu_g(t)$.
One can verify the isomorphism directly: the ideals
$(2^r,X^{2^s}-1)$ and $(2^r,(X-1)^n)$ give cofinal topologies.
For fixed $r,n$, the coefficients below degree $n$ of $X^{2^s}-1$
are divisible by $2^r$ for large $s$; conversely,
$(X-1)^{2^s}$ is divisible by two modulo $X^{2^s}-1$.
The finite group-ring coefficients give the measure on residue classes.
Continuous functions are integrated by uniform locally constant
approximation, so pointwise divisibility passes to the integral.

Write $D_Y\log B=B^{-1}Y\partial B/\partial Y$; no analytic
logarithm is intended.  Differentiating the beta cocycle in $Z$ at
$Z=1$ gives
\[
 D_Y\log B(X,Y)=q(XY)-q(Y).
\]
Symmetry gives the analogous $X$ derivative.  Differentiating
\eqref{eq:measure-duplication}, with the chain rule on the diagonal,
therefore gives
\begin{equation}
 q(X)+q(-X)=2q(X^2).\label{eq:boundary-doubling-measure}
\end{equation}
Substitution $X\mapsto-X$ is legitimate: $X-1$ becomes
$-2-(X-1)$, and every resulting coefficient converges $2$-adically.

Boundary reflection makes $\mu_g$ invariant under $t\mapsto-t$.
The transforms $q(-X)$ and $q(X^2)$ correspond respectively to
$(-1)^t\mu_g$ and to the pushforward $[2]_*\mu_g$.
Equation~\eqref{eq:boundary-doubling-measure} and the torsion-freeness
of the measure module imply
\[
 \mu_g|_{2\mathbf Z_2}=[2]_*\mu_g.
\]
If $\nu_g=\mu_g|_{\mathbf Z_2^\times}$ and
$\nu_{g,+}=\mu_g|_{1+4\mathbf Z_2}$, it follows that
\[
 2\int t^{2k}\,d\nu_{g,+}
 =\int t^{2k}\,d\nu_g
 =(1-2^{2k})\int t^{2k}\,d\mu_g.
\]
The first equality uses negation symmetry, which exchanges the two
odd residue classes modulo four.  Expansion of $e^{tu}$ shows that
$\tau_{2k+1}=\tfrac12\int t^{2k}\,d\mu_g$, proving
\eqref{eq:integral-boundary-moments}; its factor $1-2^{2k}$ is a unit.
Also $\mu_g(\mathbf Z_2)=q_g(1)=0$, hence
$\nu_g(\mathbf Z_2^\times)=0$ and $\nu_{g,+}(1+4\mathbf Z_2)=0$.

The integral measure $\mu_g$ and its negation symmetry exist on all
of $M$, using only the integral boundary series and reflection.
Modulo two, the even residue classes contribute zero to
$\int t^{2k}\,d\mu_g$, while the two odd residue classes modulo four
have equal mass. The moment is therefore even, so its half is
integral. Moreover $t^{2k}-t^2$ is divisible by four for every
$t\in\mathbf Z_2$. Taking the integral and dividing by two proves
the asserted congruence; the identity $\tau_3=a_2$ follows from
$a_1=0$. Each moment is a finite rational linear combination of
boundary coefficients, hence is continuous.

On $M$ the beta multipliers multiply.  Their conjugation law is
$B_{hgh^{-1}}(X,Y)=B_g(X^{\lambda_2(h)},Y^{\lambda_2(h)})$,
as follows from the semilinear action on the metabelian commutator
module, with the multipliers of $h$ and $h^{-1}$ cancelling.
Thus $q$ is additive on $M$ and
$q_{hgh^{-1}}(X)=\lambda_2(h)q_g(X^{\lambda_2(h)})$.
The coefficient of $u^{2k}$ has the asserted weight.
\end{proof}

\begin{lemma}[Uniform odd-coordinate upper bound]
\label{lem:uniform-odd-coordinate-bound}
If $g\in N$ and $a_2(g)=a_4(g)=\cdots=a_{2r-2}(g)=0$, then
\[
 a_{2r}(g)\in2^{2r+1}\mathbf Z_2.
\]
\end{lemma}
\begin{proof}
Reflection implies that $q(1+S)$ is a power series in
$S^2/(1+S)$.  This follows recursively by cancelling the first
nonzero even-degree term of a reflection-invariant series; an odd
first degree would have its coefficient equal to its negative.
Consequently the assumed lower even coefficients make $q$ start in
$S^{2r}$.  The measure in Lemma~\ref{lem:integral-boundary-measure}
has zero moments in degrees $0,2,\ldots,2r-2$.
For $t=2z+1$ put
\[
 P_r(t)=\prod_{j=0}^{r-1}\bigl(t^2-(2j+1)^2\bigr)
 =2^{2r}(2r)!\binom{z+r}{2r}.
\]
The binomial polynomial is $\mathbf Z_2$-valued, by continuity from
integers.  Hence
\[
 (1-2^{2r})\tau_{2r+1}
 =\int t^{2r}\,d\nu_{g,+}
 =\int P_r(t)\,d\nu_{g,+}
 \in2^{2r}(2r)!\mathbf Z_2.
\]
As $\tau_{2r+1}=\frac{(2r)!}{2}a_{2r}$, cancellation in
$\mathbf Q_2$ gives the asserted divisibility.
\end{proof}


\section{A cubic Magnus normalization beyond rational radicals}
\label{sec:cubic-magnus-normalization}

Put $G=\widehat{\mathrm{GT}}$, let $H\subseteq G$ be the Galois image, and
write $N=\ker(\lambda,\rho_2^I)$.  For a multiplier-one element
$g=(1,f)$, write its pro-$2$ Magnus expansion as
\[
 x=1+S,\qquad y=1+T,\qquad
 m_3(g)=[S^2T]f\in\mathbf Z_2.
\]
Here brackets around a word mean its coefficient, whereas group commutators
are $[a,b]=aba^{-1}b^{-1}$.  The notation $m_3$ separates this cubic
monodromy coordinate from the rational Kummer coordinate $\rho_3^I$.

The integral measure already gives the cubic divisibility. To identify
its arithmetic image and distinguish its quadratic reduction from rational
Kummer characters, we also need the boundary parity congruence below.
The arithmetic input is the simultaneous rational Kummer image
$\prod_q2\widehat{\mathbf Z}$ over $\mathbf Q^{\mathrm{ab}}$,
proved in Theorem~\ref{thm:main-rational-normalization}.

\begin{lemma}[An integral alternating coefficient identity]
\label{lem:integral-alternating-cubic}
Let $f$ be a pro-$2$ commutator word satisfying
\[
 f(y,x)f(x,y)=1,\qquad
 f(z,x)f(y,z)f(x,y)=1,\qquad z=(xy)^{-1}.
\]
Set $a=[S^2T]f$ and $A=[S^2T^2S]f$.  Then
\begin{equation}\label{eq:alternating-cubic-identity}
 [STS]f=-2a,\qquad [ST^2S]f=a,\qquad
 [STSTS]f=-2a-4A.
\end{equation}
In particular, $[STSTS]f\equiv2a\pmod4$.
\end{lemma}

\begin{proof}
The degree-two part of $f$ is a multiple of $[S,T]$.  The degree-two
part of the three-cycle identity is three times that multiple, so it is
zero.  Thus $f\in\gamma_3$ in the pro-$2$ nilpotent quotient relevant
to the calculation.

For completeness, the low-weight collection calculation needed here is
given explicitly.  Let $L_w$ be the standard Lyndon group commutator
associated to a Lyndon word $w$ on $x<y$: $L_x=x$, $L_y=y$, and
$L_w=[L_u,L_v]$ for the standard factorization $w=uv$ with $v$ the
longest proper Lyndon suffix.  Up to weight five, collect $f$ as the
ordered product of these commutators, with coordinates
\[
\begin{gathered}
\begin{array}{c|rrrrr}
 w&xxy&xyy&xxxy&xxyy&xyyy\\
 \hline
 \text{coordinate}&a&d&u&v&w
\end{array}
\\[5pt]
\begin{array}{c|rrrrrr}
 w&xxxxy&xxxyy&xxyxy&xxyyy&xyxyy&xyyyy\\
 \hline
 \text{coordinate}&b&h&i&j&e&k.
\end{array}
\end{gathered}
\]
All coordinates belong to $\mathbf Z_2$.  Since every factor begins in
degree at least three, the product is linear in these coordinates
through degree five.  Write $I=f(y,x)f(x,y)-1$ and
$II=f(z,x)f(y,z)f(x,y)-1$.  Direct substitution of
$z=(1+T)^{-1}(1+S)^{-1}$ gives the following five coefficients:
\[
\begin{array}{c|l}
\text{coefficient}&\text{value}\\ \hline
{[S^2T]I}&a+d\\
{[S^3T]I}&u-w\\
{[S^3T^2]I}&h+j-u+w\\
{[S^3T]II}&a-2d+2u+v+2w\\
{[S^4T]II}&-a+3d-h-j-2u-v-3w.
\end{array}
\]
Each entry is zero.  The first three give $d=-a$, $w=u$, and
$h+j=0$; the last two then give
\[
 3a+4u+v=0,\qquad -4a-5u-v=0.
\]
Consequently $u=w=-a$ and $v=a$.  No division by two has occurred.
The same commutator expansion gives
\begin{equation}\label{eq:cubic-three-small-coefficients}
 [STS]f=-2a,\qquad
 [ST^2S]f=a,\qquad
 [STS^2]f=4a+3u=a.
\end{equation}
This supplies all the low-weight collection identities used below.

Let $C=[ST^2S^2]f$.  The group-ring antipode sends each generator to
its inverse and reverses products.  Since $f-1$ starts in degree three,
$f^{-1}=1-(f-1)$ through degree five.  The coefficient of $S^2T^2S$
in the antipode of $f$ is therefore
\[
 -A=-\bigl(C+[STS^2]f+[ST^2S]f+[STS]f\bigr)=-C.
\]
Thus $C=A$.

Finally, group-likeness of $f$ for
$\Delta(S)=S\otimes1+1\otimes S+S\otimes S$ and the analogous
formula for $T$ gives, on taking the coefficient of
$S\otimes ST^2S$,
\[
 0=2[S^2T^2S]f+[STSTS]f+2[ST^2S^2]f+2[ST^2S]f.
\]
The right-hand side counts the five insertions of the extra $S$ and
the two possible mergers with an existing $S$.  Substitution of
$C=A$ and \eqref{eq:cubic-three-small-coefficients} proves
\eqref{eq:alternating-cubic-identity}.
\end{proof}

\begin{proposition}[Divisibility and the cubic character]
\label{prop:cubic-magnus-eight}
One has
\[
 m_3(N)\subseteq8\mathbf Z_2.
\]
Consequently
\[
 \eta_3:N\longrightarrow\mathbf F_2,
 \qquad \eta_3(g)=m_3(g)/8\pmod2
\]
is an ordinary continuous $G$-invariant quadratic character.
\end{proposition}

\begin{proof}
Lemma~\ref{lem:uniform-odd-coordinate-bound}, with $r=1$, gives
$m_3(N)\subseteq8\mathbf Z_2$ directly.

On $\ker\lambda$, the cubic Magnus coordinate is additive under the
GT product.  Moreover, for $s\in G$ and $n\in\ker\lambda$,
\begin{equation}\label{eq:cubic-scalar-conjugation}
 m_3(sns^{-1})=\lambda_2(s)^3m_3(n).
\end{equation}
Indeed, a multiplier-one word has no part below degree three; the
GT product adds its degree-three terms, and conjugation acts on that
degree by the cube of the degree-one scalar.  Terms from the
conjugating word start in degree two and cannot change this leading
degree-three term.  These statements may equally be read from the
induced tangential automorphism in degree four.  Continuity follows
from the Magnus coefficient map.  Finally, every $2$-adic unit has
odd cube, so division by eight and reduction modulo two turn
\eqref{eq:cubic-scalar-conjugation} into invariance.
\end{proof}

We next give the arithmetic witness, including the small congruence
which detects its parity.

\begin{lemma}[A boundary parity congruence]
\label{lem:rho2-cubic-modfour}
For every multiplier-one gentle element,
\[
 \rho_2^I(g)\equiv-2m_3(g)\pmod4
\]
in the $2$-adic coordinate.
\end{lemma}

\begin{proof}
Use the affine matrix representation
\[
 x\longmapsto D=\begin{pmatrix}-1&0\\0&1\end{pmatrix},\qquad
 y\longmapsto E=\begin{pmatrix}1&1\\0&1\end{pmatrix}.
\]
In the index-two cyclic subgroup, $x^2$ maps to the identity, while
$Y_0=y$ and $Y_1=xyx^{-1}$ have translation coordinates $1$ and $-1$.
Since $f$ has zero total $y$-exponent, its translation is
$2\ell_2(f)=-2\rho_2^I(f)$.

On the Magnus side, $S=D-I=\operatorname{diag}(-2,0)$ and $T=E-I$.
Here $T^2=TS=0$.  Since the first nonzero term of $f$ has degree
three, its translation is
\[
 \sum_{n\ge2}(-2)^n[S^nT]f\equiv4m_3(f)\pmod8.
\]
The series converges $2$-adically.  Comparing translations and
dividing the congruence by $-2$ proves the claim.
\end{proof}

Let
\[
 D=N_{\mathrm{rad}}=\bigcap_q\ker(\rho_q^I)\cap\ker\lambda,
 \qquad
 K_{\mathrm{rad}}=D\cap H=G_L,
\]
where $L=\mathbf Q(\mu_\infty,q^{1/n}:q\text{ prime},n\ge1)$.
The all-prime normalization gives $G=DH$.

\begin{proposition}[An arithmetic value outside the rational Kummer sector]
\label{prop:cubic-arithmetic-witness}
There exists $k\in K_{\mathrm{rad}}$ such that
\[
 m_3(k)\in8\mathbf Z_2^\times.
\]
In particular, $\eta_3|_{K_{\mathrm{rad}}}\ne0$, so $\eta_3$ is
independent of every character generated by rational Kummer coordinates.
Moreover,
\[
 m_3(D)=m_3(K_{\mathrm{rad}})=8\mathbf Z_2.
\]
\end{proposition}

\begin{proof}
Simultaneous rational Kummer surjectivity supplies
$\sigma\in H^{\mathrm{cyc}}=G_{\mathbf Q^{\mathrm{ab}}}$ with
\[
 \rho_2^I(\sigma)=2\in\widehat{\mathbf Z},\qquad
 \rho_q^I(\sigma)=0\quad(q\text{ odd}).
\]
Lemma~\ref{lem:rho2-cubic-modfour} shows that $m_3(\sigma)$ is odd.
Choose $h\in H$ with $\lambda_2(h)=3$, and put
$t=\lambda(h)\in\widehat{\mathbf Z}^{\times}$.  For
$a\in\widehat{\mathbf Z}$, the power $\sigma^a$ means the value of
the continuous map $\widehat{\mathbf Z}\to
\overline{\langle\sigma\rangle}$ extending integer powers.
Define
\[
 k=(h\sigma h^{-1})\sigma^{-t}.
\]
Every rational Kummer coordinate of this element is zero:
on the multiplier-one subgroup the cocycles are homomorphisms,
and conjugation by $h$ multiplies each by $t$, exactly as the
profinite power $\sigma^t$ does.  Thus $k\in K_{\mathrm{rad}}$.
By additivity and \eqref{eq:cubic-scalar-conjugation},
\[
 m_3(k)=(3^3-3)m_3(\sigma)
        =24m_3(\sigma)\in8\mathbf Z_2^\times.
\]
Every rational Kummer character vanishes on $K_{\mathrm{rad}}$, while
$\eta_3(k)=1$, proving the independence assertion.

The image of the compact group $K_{\mathrm{rad}}$ under $m_3$ is a
closed additive subgroup of $\mathbf Z_2$.  It is contained in
$8\mathbf Z_2$ by Proposition~\ref{prop:cubic-magnus-eight}, and it
contains a generator of that closed subgroup.  It is therefore all of
$8\mathbf Z_2$.  The same conclusion follows for $D$ by inclusion.
\end{proof}

\begin{corollary}[A further proper normalization]
\label{cor:cubic-radical-normalization}
Define
\[
 D_3=N_{\mathrm{rad},3}=D\cap\ker m_3,
 \qquad K_{\mathrm{rad},3}=K_{\mathrm{rad}}\cap\ker m_3.
\]
Then $D_3$ is a proper closed normal subgroup of $G$, and
\[
 D=D_3K_{\mathrm{rad}},\qquad
 G=D_3H,\qquad D_3\cap H=K_{\mathrm{rad},3}.
\]
Writing
\[
 B=N/\overline{[G,N]K},\qquad K=N\cap H,
\]
the previously obtained radical normalization refines canonically to
\[
 B\simeq
 D/\overline{[G,D]K_{\mathrm{rad}}}
 \simeq
 D_3/\overline{[G,D_3]K_{\mathrm{rad},3}}.
\]
If $L_3=\overline{\mathbf Q}^{\,K_{\mathrm{rad},3}}$, then
$\operatorname{Gal}(L_3/L)\simeq\mathbf Z_2$.  This defines the
arithmetic extension by the explicit cubic monodromy coordinate; no
radical formula for $L_3$ is asserted.
\end{corollary}

\begin{proof}
Closedness and normality follow from continuity and
\eqref{eq:cubic-scalar-conjugation}.  Properness and the factorization
$D=D_3K_{\mathrm{rad}}$ follow from the equality of the two cubic
images.  The remaining factorization and intersection statements are
immediate from $G=DH$.

For the quotient assertion, one may use the following elementary
group argument.  If a profinite group factors as $G=EH$ with $E$
closed normal, then
\[
 E/\overline{[G,E](E\cap H)}
 \simeq G/\overline{[G,E]H}.
\]
The right-hand quotient is abelian because the image of $E$ is
central and it generates after the image of $H$ is killed.  If also
$E=E'(E\cap H)$ and $G=E'H$, then
$\overline{[G,E]H}=\overline{[G,E']H}$: after quotienting by the latter,
the image of $E'$ is central and $H$ is trivial.  Apply this with
$E=D$, $E'=D_3$, and combine with the earlier radical normalization.
Finally, the kernel of the surjective map
$m_3:G_L\to8\mathbf Z_2$ is exactly $G_{L_3}$.
\end{proof}

\begin{corollary}\label{cor:cubic-radical-joint-image}
The map
\[
 N\longrightarrow
 \left(\prod_{q\text{ odd prime}}2\widehat{\mathbf Z}\right)
       \times8\mathbf Z_2,
 \qquad n\longmapsto((\rho_q^I(n))_q,m_3(n))
\]
has kernel $D_3$ and is surjective; already its restriction to
$K=N\cap H$ is surjective.
Consequently the $G$-invariant continuous $\mathbf F_2$-characters
of $N$ which vanish on $D_3$ form exactly
\[
 \left(\bigoplus_{q\text{ odd prime}}\mathbf F_2\psi_q\right)
                \oplus\mathbf F_2\eta_3.
\]
For an odd prime $p$, no nonzero $G$-invariant
$\mathbf F_p$-character of $N$ vanishes on $D_3$.
\end{corollary}

\begin{proof}
Prescribe any rational Kummer vector and cubic value in the displayed
target.  First choose $h\in K$ with the prescribed Kummer vector.
Its cubic coordinate lies in $8\mathbf Z_2$.  The difference between
that coordinate and the prescribed one can be corrected by an element
of $K_{\mathrm{rad}}$, by
Proposition~\ref{prop:cubic-arithmetic-witness}; this changes none of
the Kummer coordinates.  The kernel is exactly the common kernel $D_3$.
Compactness identifies $N/D_3$ with the displayed product.

Every continuous character of this product into a finite group depends
on finitely many coordinates.  At two, each divided coordinate has a
one-dimensional character space, and conjugation multiplies the Kummer
coordinates by $\lambda$ and the cubic coordinate by $\lambda_2^3$.
Both scalars are odd, so every displayed quadratic character is invariant.
At an odd prime $p$, the pro-$2$ cubic factor has no nonzero character
into $\mathbf F_p$.  A Kummer character is multiplied by $\lambda_p$
under conjugation; choosing a multiplier with $\lambda_p\not\equiv1$
modulo $p$ forces every invariant coefficient to vanish.
\end{proof}

The character $\eta_3$ has nonzero arithmetic restriction.  The result
therefore enlarges the explicitly normalized arithmetic character
sector; it does not construct a nonzero character in the kernel of
restriction to the Galois subgroup.

\section{The complete invariant quadratic arithmetic target}
\label{sec:fr-Y2-target}

Put
\[
 H=G_{\mathbf Q},\qquad
 K_2=\mathbf Q(\mu_\infty,2^{1/\infty}),\qquad
 K=G_{K_2},\qquad
 Y_2=H^1_{\mathrm{cts}}(K,\mathbf F_2)^H.
\]
All cohomology below is continuous, with trivial coefficient module
\(\mathbf F_2\), unless otherwise indicated.
We use the established description
\[
 A=H/K
 =\{(u,t)\in U\ltimes\widehat{\mathbf Z}:
                  t\bmod2=\epsilon_8(u)\},
 \qquad U=\widehat{\mathbf Z}^{\times},
\]
with multiplication
\((u,t)(v,s)=(uv,t+us)\).
Every quadratic character of \(H\) factors through \(U\).
We also use the already established arithmetic characters
\[
 \psi_p=\rho_p^I/2\bmod2\quad(p\text{ odd prime}),\qquad
 \eta_3=m_3/8\bmod2
\]
on \(K\), their \(H\)-invariance, and the cubic identities recalled
in the proof below.

For an odd prime \(p\), set
\[
 p^*=(-1)^{(p-1)/2}p,\qquad
 \varepsilon_p=\frac{1-(2/p)}2\in\mathbf F_2,
 \qquad
 \ell_p(d)=\frac{1-(d/p)}2\in\mathbf F_2
 \quad(p\nmid d).
\]
Here \((d/p)\) is the Legendre symbol.
Let \(b,a,c_p\in H^1(U,\mathbf F_2)\) be the quadratic
characters of \(\mathbf Q(i)\), \(\mathbf Q(\sqrt2)\),
and \(\mathbf Q(\sqrt{p^*})\), respectively.
Thus
\[
 b(u)=\frac{u_2-1}{2}\bmod2,\qquad
 a(u)=\frac{u_2^2-1}{8}\bmod2.
\]

For each odd prime \(p\), let
\(k_p=v_2(p-1)\) and let \(\delta_p\in H^2(U,\mathbf F_2)\)
be the pullback of the unique nonzero class of
\[
 H^2(C_{2^{k_p}},\mathbf F_2)
\]
from the maximal \(2\)-primary cyclic quotient of
\(\mathbf Z_p^\times\).
Equivalently, it is the class of
\(C_{2^{k_p+1}}\to C_{2^{k_p}}\) with central kernel \(C_2\).
This class does not depend on a choice of cyclic generator.
Put \(\delta_2=b^2\).

\begin{theorem}\label{thm:fr-Y2-complete}
There is a direct-sum decomposition
\begin{equation}\label{eq:fr-Y2-basis}
 Y_2=
 \mathbf F_2\eta_3
 \ \oplus\!
 \bigoplus_{p\ {\rm odd}}\mathbf F_2\psi_p
 \ \oplus\!
 \bigoplus_{p\ {\rm odd}}\mathbf F_2\kappa_{2,p}
 \ \oplus\!
 \bigoplus_{\substack{p<q\\p,q\ {\rm odd}}}
                     \mathbf F_2\kappa_{p,q}.
\end{equation}
The additional characters are specified uniquely as follows.
The transgression gives an isomorphism
\begin{equation}\label{eq:fr-Y2-transgression}
 \operatorname{tr}:Y_2
 \xrightarrow{\ \sim\ }
 \ker\!\left(H^2(A,\mathbf F_2)
                   \longrightarrow\operatorname{Br}(\mathbf Q)[2]\right).
\end{equation}
Writing \(\operatorname{Inf}_U^A\) for inflation, one has
\begin{align}
 \operatorname{tr}(\kappa_{2,p})
 &=\operatorname{Inf}_U^A
       \left(a c_p+\varepsilon_p(\delta_2+\delta_p)\right),
                                                        \label{eq:fr-Y2-k2p}\\
 \operatorname{tr}(\kappa_{p,q})
 &=\operatorname{Inf}_U^A
       \left(c_pc_q+\ell_p(q^*)\delta_p
                       +\ell_q(p^*)\delta_q\right).
                                                        \label{eq:fr-Y2-kpq}
\end{align}
All products in these formulas are cup products.
Every sum in \eqref{eq:fr-Y2-basis} is finite.

In particular the quotient of \(Y_2\) by the known span
\(\mathbf F_2\eta_3+\bigoplus_p\mathbf F_2\psi_p\)
has countably infinite dimension.
These are statements about the arithmetic target only.
They do not assert that any additional basis character
extends to an invariant character of the corresponding full-GT subgroup.
\end{theorem}

\subsection{Transgression and the affine quotient}

The five-term sequence for \(1\to K\to H\to A\to1\) begins
\[
 0\longrightarrow H^1(A)
 \longrightarrow H^1(H)
 \longrightarrow Y_2
 \xrightarrow{\operatorname{tr}}H^2(A)
 \longrightarrow H^2(H).
\]
The first map is an isomorphism: every character of \(H\) factors
through \(U\), hence through \(A\), and \(H\to A\) is surjective.
Kummer theory and Hilbert 90 identify
\(H^2(H,\mathbf F_2)\) with \(\operatorname{Br}(\mathbf Q)[2]\).
This proves \eqref{eq:fr-Y2-transgression}.

For precision, the cohomological exact sequences used here are
the continuous versions of the Hochschild--Serre sequences.
The original cochain construction and transgression are
\cite[Chapter~II, pp.~118--119; Chapter~III, Theorem~2 and the
remark on pp.~129--130]{HochschildSerre1953}.
Continuous cochains and passage through finite quotients are
\cite[II, \S4, Proposition~4.2, pp.~87--88]{MilneCFT}.
The Brauer identification is
\cite[IV, Theorem~3.14, p.~137]{MilneCFT}, combined
with the Kummer exact sequence.

The kernel of \(A\to U\) is
\[
 T=2\widehat{\mathbf Z}.
\]
It has \(H^1(T,\mathbf F_2)=\mathbf F_2 z\), where
\(z(2s)=s_2\bmod2\), and \(H^j(T,\mathbf F_2)=0\) for \(j\ge2\).
Indeed its pro-odd factor has vanishing positive cohomology
by finite averaging, while \(\mathbf Z_2\) has cohomological
dimension one. The latter also follows from the length-one
resolution over \(\mathbf F_2[[\mathbf Z_2]]=\mathbf F_2[[X]]\)
given by multiplication by \(X\).
The action of \(U\) on \(H^1(T,\mathbf F_2)\) is trivial.

Use the continuous section \(u\mapsto(u,\widetilde a(u))\),
where \(\widetilde a(u)\) is the integer \(0\) or \(1\)
representing \(a(u)\). Its extension cocycle, evaluated by \(z\), is
\[
 \frac{\widetilde a(u)+u\widetilde a(v)
                    -\widetilde a(uv)}2\bmod2
 =a(u)a(v)+b(u)a(v).
\]
The class \(a\) factors through the torsion-free
\(\mathbf Z_2\) factor of \(\mathbf Z_2^\times\), so \(a^2=0\).
Thus
\begin{equation}\label{eq:fr-Y2-extension-class}
                         d_2(z)=ab.
\end{equation}

Modulo pro-odd factors, the cyclotomic group is
\[
 U^{(2)}\cong C_2\times\mathbf Z_2
                   \times\prod_{p\ {\rm odd}} C_{2^{k_p}}.
\]
The first two factors have the specified dual characters
\(b\) and \(a\). Its first cohomology has basis
\(b,a,c_p\). Its second cohomology has basis
\begin{equation}\label{eq:fr-Y2-U2-basis}
 \delta_2,\ \delta_p,\quad
 ab,\quad bc_p,\quad ac_p,\quad c_pc_q\ (p<q).
\end{equation}
For a finite cyclic \(2\)-group,
\[
 c_p^2=
 \begin{cases}
  \delta_p,&p\equiv3\pmod4,\\
  0,&p\equiv1\pmod4.
 \end{cases}
\]
These assertions follow from the cyclic resolution and the
Kunneth formula over \(\mathbf F_2\). An explicit check of the
square relation uses the carry cocycle for
\(C_{2^{k+1}}\to C_{2^k}\): when \(k=1\) it is the square of
the degree-one class; when \(k\ge2\) that square is the
coboundary of \(j\mapsto\lfloor j/2\rfloor\bmod2\).
Only finitely many factors occur in any continuous cohomology
class. Thus the displayed basis is a direct-sum basis,
even though the group itself is a product.

The two-row spectral sequence for \(T\to A\to U\), together
with \eqref{eq:fr-Y2-extension-class}, gives
\begin{equation}\label{eq:fr-Y2-H2A}
 0\longrightarrow H^2(U)/\mathbf F_2ab
 \xrightarrow{\operatorname{Inf}_U^A}H^2(A)
 \xrightarrow{\operatorname{edge}}\mathbf F_2a
 \longrightarrow0.
\end{equation}
Here the edge target is the kernel of
\[
 H^1(U)\longrightarrow H^3(U),\qquad x\longmapsto xab.
\]
That kernel is precisely \(\mathbf F_2a\).
The products \(ab^2\) and \(abc_p\) are nonzero independent
classes by their distinct Kunneth multidegrees, whereas
\(a^2b=0\). There are no higher differentials: only the rows
of fibre degrees zero and one are nonzero.
The differential as cup product with the extension class
is \cite[Chapter~III, \S6, Theorems~3--4,
pp.~132--133]{HochschildSerre1953}; its sign is immaterial over
\(\mathbf F_2\).

\subsection{The cubic character is exactly the extra edge direction}

Choose \(\sigma\in H\cap\ker\lambda\) with
\(\rho_2^I(\sigma)=2\). Such an element exists by the
given surjection onto the affine group. Put
\(\alpha=m_3(\sigma)\in\mathbf Z_2\).
The established cubic congruence
\(\rho_{2,2}^I\equiv-2m_3\pmod4\) makes \(\alpha\) odd.

The quotient extension determined by \(\operatorname{tr}(\eta_3)\)
splits over \(T\), using the section
\[
             2s\longmapsto \sigma^s\pmod{\ker\eta_3}.
\]
This is a genuine continuous profinite-power section.
For \(h\in H\) with multiplier \(u=\lambda(h)\), its failure
to be equivariant is measured by
\[
 d(h,s)=h\sigma^s h^{-1}\sigma^{-us}\in K.
\]
The full profinite exponent is \(us\); its dyadic component
is used only when evaluating the cubic coordinate.
The Kummer conjugation law proves \(d(h,s)\in K\).
The additive cubic coordinate on the cyclotomic subgroup
and its weight-three conjugation law give
\[
 \eta_3(d(h,s))
 =\frac{u_2^3-u_2}{8}\,\alpha s_2\bmod2
 =a(u)s_2\bmod2.
\]
The last equality holds because \(u_2\) and \(\alpha\) are odd
and \((u_2^2-1)/8\bmod2=a(u)\).
This is the explicit extension-theoretic description of the edge map.
Consequently
\begin{equation}\label{eq:fr-Y2-eta-edge}
       \operatorname{edge}(\operatorname{tr}(\eta_3))=a.
\end{equation}
In particular the cubic character supplies the extra line in
\eqref{eq:fr-Y2-H2A}; it is not a cyclotomic cup class.
Since its transgression inflates to zero on \(H\), subtracting
its appropriate multiple reduces the computation of \(Y_2\)
to the kernel of the Brauer map on \(H^2(U)/\mathbf F_2ab\).

\subsection{The diagonal classes and their exact local invariants}

For every finite rational prime \(v\), the image of \(\delta_v\)
in \(\operatorname{Br}(\mathbf Q)[2]\) is ramified exactly at
\(\{v,\infty\}\). Here ramification means local invariant \(1/2\).

For an odd prime \(p\), the class \(\delta_p\) is the obstruction
to lifting the maximal \(2\)-power cyclotomic character at \(p\)
from \(C_{2^{k_p}}\) to \(C_{2^{k_p+1}}\).
At any finite place other than \(p\), this character is
unramified and lifts through the procyclic unramified quotient.
At \(p\), local reciprocity identifies its restriction to the
torsion units with a primitive character of the cyclic group
of order \(2^{k_p}\). Such a character cannot lift one step:
an element of order at most \(2^{k_p}\) in
\(C_{2^{k_p+1}}\) reduces to an even element.
Thus the local obstruction is nonzero.
At the real place, complex conjugation maps to the unique
element of order two in \(C_{2^{k_p}}\); it likewise has no
order-two lift. This is the other nonzero local obstruction.

For \(\delta_2=b^2\), the same argument uses the quadratic
character of \(\mathbf Q(i)\). It is unramified at all odd
places; at two its value on the torsion unit \(-1\) cannot
lift from \(C_2\) to \(C_4\); the real obstruction is nonzero.
This proves the assertion for all \(v\).

The precise arithmetic inputs are local reciprocity
\cite[I, Theorem~1.1, p.~20]{MilneCFT}, the local
Brauer invariant \cite[III, Theorem~2.1, p.~103]{MilneCFT},
and the global exact sequence
\cite[VIII, Theorem~4.2, p.~241]{MilneCFT}.
The last theorem identifies \(\operatorname{Br}(\mathbf Q)[2]\)
with the finitely supported tuples of local bits of even total sum.
Therefore the classes ramified at \(\{v,\infty\}\), for finite
primes \(v\), form a basis. In particular the diagonal
subspace spanned by the \(\delta_v\) maps isomorphically onto
\(\operatorname{Br}(\mathbf Q)[2]\).

\subsection{The corrected cross cups}

For a quadratic character of a rational number \(d\), cup
product with the character of \(e\) has local Brauer bit
given by the quadratic Hilbert symbol \((d,e)_v\).
This is the cup-product definition and norm criterion in
\cite[III, \S4, pp.~112--115, especially Theorem~4.4]{MilneCFT}.
We compute only the cases needed for the displayed basis.

First, \((-1,2)\) is trivial globally because
\(2=N_{\mathbf Q(i)/\mathbf Q}(1+i)\).
Thus the zero-Brauer cross cup \(ab\) is precisely the class
already killed in passing from \(U\) to \(A\).

For an odd prime \(p\), the cup \(bc_p\) has no finite
invariant except possibly at \(p\), where its bit is
\(\ell_p(-1)\). There is no invariant at two:
\(p^*\equiv1\pmod4\), so its quadratic extension of
\(\mathbf Q_2\) is split or unramified and every unit is a norm.
Hence the corrected class is
\[
 C_{b,p}=bc_p+\ell_p(-1)\delta_p.
\]

The cup \(ac_p\) can have finite invariants only at \(2\) and \(p\).
At \(p\) its bit is \(\varepsilon_p\).
Its real invariant is zero, so the product formula gives
the same bit at two. Its corrected class is
\[
 C_{2,p}=ac_p+\varepsilon_p(\delta_2+\delta_p).
\]

Finally, for distinct odd primes \(p,q\), the cup \(c_pc_q\)
has finite bits
\[
 \ell_p(q^*)\quad\text{at }p,\qquad
 \ell_q(p^*)\quad\text{at }q,
\]
and zero at every other finite place. At two both parameters
are odd and congruent to one modulo four, so the same
split-or-unramified norm observation applies.
The corrected class is
\[
 C_{p,q}=c_pc_q+\ell_p(q^*)\delta_p+\ell_q(p^*)\delta_q.
\]
The real bit is automatically cancelled as well, by the
global sum formula.

For completeness, the odd-prime local calculations just used
are elementary. A quadratic extension defined by a unit is
split or unramified, so a second unit is a norm. A quadratic
extension defined by a parameter of odd valuation is tamely
ramified, and its unit norms have square residues. Thus a
uniformizer against a unit gives the corresponding Legendre
symbol. There are no unlisted finite places.

Adding the uniquely determined diagonal correction to each
cross cup is a triangular change of the basis
\eqref{eq:fr-Y2-U2-basis}. Since the diagonal subspace maps
isomorphically onto the Brauer group, it follows that
\[
 \ker\!\left(H^2(U)\longrightarrow\operatorname{Br}(\mathbf Q)[2]\right)
 =
 \mathbf F_2ab
 \oplus\bigoplus_p\mathbf F_2 C_{b,p}
 \oplus\bigoplus_p\mathbf F_2 C_{2,p}
 \oplus\bigoplus_{p<q}\mathbf F_2 C_{p,q}.
\]
Quotienting by \(ab\), and adjoining the extra cubic line
from \eqref{eq:fr-Y2-eta-edge}, proves the asserted completeness
and independence once the known Kummer coordinates are identified.

\subsection{Identification of the odd-prime Kummer axes}

Let \(\epsilon_p\in H^1(U,\mathbf F_2)\) denote the character
of \(\mathbf Q(\sqrt p)\), with the positive rational prime
as radicand. Then
\[
 \epsilon_p=
 \begin{cases}
 c_p,&p\equiv1\pmod4,\\
 b+c_p,&p\equiv3\pmod4.
 \end{cases}
\]
The mod-four Kummer cocycle of \(p\), whose mod-two value is
\(\epsilon_p\), determines a central quadratic extension over \(A\).
Using its integer-valued parity section, exactly the same carry
calculation as above gives
\[
 \operatorname{tr}(\psi_p)
 =\operatorname{Inf}_U^A(\epsilon_p^2+b\epsilon_p).
\]
If \(p\equiv1\pmod4\), this is \(\operatorname{Inf}(bc_p)\).
If \(p\equiv3\pmod4\), it is
\(\operatorname{Inf}(bc_p+c_p^2)
 =\operatorname{Inf}(bc_p+\delta_p)\).
Both cases give
\[
                  \operatorname{tr}(\psi_p)
                         =\operatorname{Inf}_U^A(C_{b,p}).
\]
This proves every assertion in
Theorem~\ref{thm:fr-Y2-complete}.

\subsection{A completely explicit target outside the known span}

The character \(\kappa_{5,29}\) has a small radical realization.
Put
\[
 \alpha=7+2\sqrt5,\qquad
 L=\mathbf Q(\sqrt5,\sqrt{29},\sqrt\alpha).
\]
Its norm from \(\mathbf Q(\sqrt5)\) is \(29\).
Let \(F=\mathbf Q(\sqrt5)\) and \(E=F(\sqrt{29})\).
Neither \(\alpha\) nor \(\alpha/29\) is a square in \(F\),
because its norm is respectively \(29\) or \(1/29\),
neither a square in \(\mathbf Q\).
An element of \(F\) becomes a square in \(F(\sqrt{29})\)
only if it or its quotient by \(29\) was already a square in \(F\):
expand \((x+y\sqrt{29})^2\) and use \(2xy=0\).
Thus \([L:E]=2\), so \([L:\mathbf Q]=8\).

Write \(c=\sqrt\alpha\), \(d=\sqrt{29}\).
The transformations
\[
 r:\quad \sqrt5\mapsto-\sqrt5,\ d\mapsto d,\ c\mapsto d/c,
 \qquad
 s:\quad \sqrt5\mapsto\sqrt5,\ d\mapsto-d,\ c\mapsto c
\]
are involutions. Their commutator fixes \(E\) and sends
\(c\) to \(-c\). Hence \(L/\mathbf Q\) is dihedral of order
eight, with central quadratic subgroup \(\operatorname{Gal}(L/E)\).

Choose \(u^{(5)},u^{(29)}\in U\) supported at the indicated
coefficient primes, with value \(2\) at that prime and \(1\)
elsewhere. The number \(2\) is a nonsquare modulo both
\(5\) and \(29\). Their dyadic components are one, so
\[
 (u^{(5)},0),\qquad (u^{(29)},0)
\]
are commuting elements of \(A\).
Choose arbitrary lifts \(h_5,h_{29}\in H\).
Then
\[
                         k=[h_5,h_{29}]\in K.
\]
On \(L\), these lifts have the two independent images of
\(r\) and \(s\) modulo the centre. Multiplying a lift by a
central element does not change their commutator.
Consequently \(k\) sends \(c\) to \(-c\).

This proves directly that \(L\not\subset K_2\).
Since \(E\subset\mathbf Q^{\rm ab}\subset K_2\) and
\([L:E]=2\), one has \(L\cap K_2=E\).
The quadratic extension \(K_2L/K_2\) therefore gives a
nonzero \(H\)-invariant character \(\chi_L\in Y_2\),
with \(\chi_L(k)=1\).

Every known \(\psi_p\) vanishes on \(k\).
Indeed both lifts have dyadic multiplier one, so each
dyadic Kummer cocycle is a homomorphism on the subgroup
they generate and is zero on their commutator.
Likewise the cubic coordinate is a homomorphism on
\(\ker\lambda_2\): the dyadic degree-two term vanishes
when \(\lambda_2=1\), and the degree-three term is additive
under the GT product. Hence \(m_3(k)=0\) and
\(\eta_3(k)=0\).
Thus
\[
 \chi_L\notin
 \mathbf F_2\eta_3+\bigoplus_p\mathbf F_2\psi_p.
\]

Finally \((5,29)\) is split, because \(29=N_{F/\mathbf Q}(\alpha)\).
The two coordinate involutions in the dihedral presentation
have square one and nontrivial commutator. Its central
extension class is therefore the cross cup \(c_5c_{29}\).
It follows that \(\chi_L=\kappa_{5,29}\) in the normalization
of \eqref{eq:fr-Y2-kpq}. The norm construction of such dihedral
fields is classical; see \cite[Lemma~5.1 and Corollary~5.2,
author pp.~9--10]{StevenhagenRedei2022}.

The argument proves no lift of \(\chi_L\), or of any other
new target basis vector, to the full-GT invariant character
space. That remains a separate extension problem.


\section{The full arithmetic target at an odd coefficient prime}
\label{sec:odd-galois-target}

Fix an odd prime \(p\).  All cohomology in this section is
continuous cohomology with the trivial discrete coefficient
module \(\mathbf F_p\), unless a twist is displayed.  Put
\[
 K_2=\mathbf Q(\mu_\infty,2^{1/\infty}),\quad
 A=\operatorname{Gal}(K_2/\mathbf Q),\quad
 U=\widehat{\mathbf Z}^{\times},\quad
 Y_p=H^1(G_{K_2},\mathbf F_p)^{G_{\mathbf Q}}.
\]
The established affine description gives
\begin{equation}\label{eq:odd-target-affine-extension}
 1\longrightarrow 2\widehat{\mathbf Z}(1)
       \longrightarrow A\longrightarrow U\longrightarrow1,
\end{equation}
with conjugation on the translation kernel given by
multiplication by the corresponding unit.

Let \(\mathcal S_p\) be the set of rational primes
\(\ell\equiv1\pmod p\).  Put
\(v_\ell=v_p(\ell-1)\), and choose an isomorphism from
the maximal \(p\)-quotient of \(\mathbf Z_\ell^\times\)
to \(\mathbf Z/p^{v_\ell}\mathbf Z\).
Its reduction modulo \(p\) defines \(\gamma_\ell\).
Choose also a generator
\(\gamma_0\in\operatorname{Hom}(1+p\mathbf Z_p,\mathbf F_p)\),
extended to \(\mathbf Z_p^\times\) by zero on its
prime-to-\(p\) torsion.  Regard all these as characters
of \(U\) by the respective coordinate projections.

Define \(\delta_\ell\in H^2(U,\mathbf F_p)\) by pulling
back the central extension
\begin{equation}\label{eq:odd-target-cyclic-extension}
 0\longrightarrow\mathbf F_p
   \xrightarrow{\ a\mapsto p^{v_\ell}a\ }
   \mathbf Z/p^{v_\ell+1}\mathbf Z
   \longrightarrow\mathbf Z/p^{v_\ell}\mathbf Z
   \longrightarrow0.
\end{equation}
This normalization uses the chosen cyclic coordinate.
When \(v_\ell>1\), it must not be confused with applying
the ordinary length-two Bockstein directly to
\(\gamma_\ell:G\to\mathbf F_p\).

\begin{theorem}\label{thm:odd-full-galois-target}
Inflation gives an isomorphism
\[
                         H^2(U,\mathbf F_p)\simeq H^2(A,\mathbf F_p).
\]
In fact the corresponding inflation is an isomorphism in
every degree.  The arithmetic inflation images of the
\(\delta_\ell\), for \(\ell\in\mathcal S_p\), form a
basis of \(H^2(G_{\mathbf Q},\mathbf F_p)\).

Order \(\{0\}\cup\mathcal S_p\) with zero first.  For
each \(i<j\), there is a unique finite linear combination
\(D_{ij}\) of the \(\delta_\ell\) having the same
arithmetic image as \(\gamma_i\smile\gamma_j\).
The transgression is an isomorphism
\[
 Y_p\simeq
 \ker\bigl(H^2(A,\mathbf F_p)
                  \longrightarrow H^2(G_{\mathbf Q},\mathbf F_p)\bigr),
\]
and the inverse images of
\(\gamma_i\smile\gamma_j-D_{ij}\), identified with
classes of \(A\) by inflation, form a basis of \(Y_p\).
Every linear combination here has finite support.
\end{theorem}

There is also an explicit two-place formula for the
correction.  Set
\[
              a_\ell=\frac{\ell-1}{p^{v_\ell}}\bmod p
                                            \in\mathbf F_p^\times.
\]
For distinct positive indices \(\ell<m\) one has
\begin{equation}\label{eq:odd-target-two-place-correction}
 D_{\ell m}=
       a_\ell^{-1}\gamma_m(\ell)\,\delta_\ell
          -a_m^{-1}\gamma_\ell(m)\,\delta_m,
 \qquad
 D_{0\ell}=-a_\ell^{-1}\gamma_0(\ell)\,\delta_\ell.
\end{equation}
The evaluations use \(\ell\in\mathbf Z_m^\times\),
\(m\in\mathbf Z_\ell^\times\), and
\(\ell\in\mathbf Z_p^\times\), respectively.
Frobenius is normalized to act on residue fields by the
positive power map.  Formula
\eqref{eq:odd-target-two-place-correction} corresponds
to the extension convention
\eqref{eq:odd-target-cyclic-extension}.

\subsection{Cohomology of the unit and affine groups}

The maximal pro-\(p\) quotient of \(U\) is
\[
            U^{(p)}\simeq\mathbf Z_p
                    \times\prod_{\ell\in\mathcal S_p}
                                  C_{p^{v_\ell}}.
\]
The kernel has no finite quotient of order divisible by
\(p\), and hence has no positive-degree
\(\mathbf F_p\)-cohomology: for its finite quotients,
averaging makes invariants exact, and one passes to the
continuous direct limit.  Thus \(U\) and \(U^{(p)}\)
have the same cohomology with these coefficients.

The cyclic resolution gives one generator in degrees
one and two for each finite cyclic \(p\)-group.  The
degree-two generator is its nonsplit cyclic-extension
class.  The square of a degree-one class is zero, since
\(p\) is odd.  For \(\mathbf Z_p\), continuous
cohomology has only degrees zero and one: on passing
from \(C_{p^n}\) to \(C_{p^{n+1}}\), the degree-two
extension class inflates to zero because the quotient
map lifts to \(C_{p^{n+1}}\) itself; cyclic periodicity
then kills every degree at least two in the direct
limit.  The cyclic calculation and the continuous
direct-limit assertion are recalled in
\cite[II, Proposition~3.4, p.~81; Proposition~4.2,
pp.~87--88]{MilneCFT}.

Tensoring the cyclic resolutions for a finite product,
or the usual Kunneth formula over \(\mathbf F_p\), now
gives
\begin{align}
 H^1(U,\mathbf F_p)
   &=\bigoplus_{i\in\{0\}\cup\mathcal S_p}\mathbf F_p\gamma_i,
       \label{eq:odd-target-unit-H1}\\
 H^2(U,\mathbf F_p)
   &=\left(\bigoplus_{\ell\in\mathcal S_p}\mathbf F_p\delta_\ell\right)
       \oplus
       \left(\bigoplus_{i<j}\mathbf F_p
                                  (\gamma_i\smile\gamma_j)\right).
       \label{eq:odd-target-unit-H2}
\end{align}
These are direct sums.  A continuous cochain on a compact
profinite group, with finite discrete values, factors
through a finite quotient.  An open subgroup of the
displayed product contains every factor outside a
finite set.  Thus a cohomology class involves only
finitely many cyclic coordinates and a finite quotient
of the \(\mathbf Z_p\) coordinate.  This proves the
finite-support assertion in these formulas.

Let \(T=2\widehat{\mathbf Z}(1)\) in
\eqref{eq:odd-target-affine-extension}.  Its cohomology
with trivial coefficients is
\[
 H^0(T,\mathbf F_p)=\mathbf F_p,\quad
 H^1(T,\mathbf F_p)=\mathbf F_p(-1),\quad
 H^b(T,\mathbf F_p)=0\quad(b\ge2).
\]
The action in the middle expression is the inverse
cyclotomic action, since a unit acts on translations
by multiplication.  The central order-two subgroup
\(\{1,-1\}\subset U\) acts by \(-1\) on this
one-dimensional module.  Its invariants are zero,
and averaging kills its positive cohomology.  The
Hochschild--Serre spectral sequence for that subgroup
therefore gives
\[
                         H^a(U,\mathbf F_p(-1))=0
                                      \quad\hbox{for every }a\ge0.
\]
In the spectral sequence of
\eqref{eq:odd-target-affine-extension}, only the
row \(b=0\) remains.  Its edge map is inflation,
which proves the asserted isomorphism in all degrees.
The spectral-sequence and five-term conventions used
here are those of
\cite[II, Remark~1.35, p.~72; \S4, pp.~87--88]{MilneCFT}.

\subsection{The global second cohomology, including its support}

We derive the required global statement directly from
the Brauer local-global theorem.  Put
\(L=\mathbf Q(\mu_p)\),
\(\Delta=\operatorname{Gal}(L/\mathbf Q)\), and let
\(\chi:\Delta\to\mathbf F_p^\times\) be the
cyclotomic character.  Since \(|\Delta|=p-1\),
restriction gives
\[
 H^2(G_{\mathbf Q},\mathbf F_p)
                     \simeq H^2(G_L,\mathbf F_p)^\Delta.
\]
Choose \(\zeta_p\) to identify the trivial module over
\(L\) with \(\mu_p\).  Kummer theory, Hilbert 90,
and the cohomological interpretation of the Brauer
group identify \(H^2(G_L,\mu_p)\) with
\(\operatorname{Br}(L)[p]\); see
\cite[II, Proposition~1.22; IV, Theorem~3.14,
p.~137]{MilneCFT}.
The identification by \(\zeta_p\) is not
\(\Delta\)-equivariant: it twists the action by
\(\chi\).  Consequently
\begin{equation}\label{eq:odd-target-brauer-eigenspace}
 H^2(G_{\mathbf Q},\mathbf F_p)
              \simeq\operatorname{Br}(L)[p]^{\chi},
\end{equation}
where \(\delta b=\chi(\delta)b\) defines the
eigenspace on the right.  Indeed if \(c\) has trivial
\(\Delta\)-action before changing coefficients,
then applying \(\delta\) to \(\zeta_p^c\) multiplies
its additive exponent by \(\chi(\delta)\).

The Brauer fundamental exact sequence
\cite[VIII, Theorem~4.2, p.~241]{MilneCFT}, together
with the local invariant isomorphisms, gives
\begin{equation}\label{eq:odd-target-brauer-exact}
 0\longrightarrow\operatorname{Br}(L)[p]
       \longrightarrow\bigoplus_{w\ {
m finite}}\mathbf F_p
       \xrightarrow{\ \sum\ }\mathbf F_p\longrightarrow0.
\end{equation}
The field \(L\) is totally imaginary.  The group
\(\Delta\) permutes the local summands and acts
trivially on the final sum.  Taking a prescribed
\(\Delta\)-eigenspace is exact, because \(p\) does
not divide \(|\Delta|\).  The nontrivial character
\(\chi\) has zero eigenspace in the final term.
It follows that the \(\chi\)-eigenspace of
\(\operatorname{Br}(L)[p]\) is the direct sum of
the \(\chi\)-eigenspaces of the local permutation
modules.

For a rational prime \(r\) and \(w\mid r\), that
permutation module is \(\mathbf F_p[\Delta/\Delta_w]\).
Its \(\chi\)-eigenspace has dimension one if
\(\chi|_{\Delta_w}=1\), and zero otherwise.
The character \(\chi\) is faithful, so the former
condition is exactly that \(r\) split completely
in \(L\), namely \(r\in\mathcal S_p\).

The same coefficient-change argument over each local
field shows that \(H^2(G_{\mathbf Q_r},\mathbf F_p)\)
is one-dimensional for \(r\in\mathcal S_p\), and
zero otherwise.  At the prime \(p\), the local
decomposition group is all of \(\Delta\), so it
contributes zero.  Real places contribute zero since
their group has order two.  The local invariant
of a Brauer class is unchanged by a field
automorphism, which is the only local equivariance
used here.  Hence localization itself is an
isomorphism
\begin{equation}\label{eq:odd-target-localization-isomorphism}
 H^2(G_{\mathbf Q},\mathbf F_p)
          \xrightarrow{\ \sim\ }
       \bigoplus_{\ell\in\mathcal S_p}
                            H^2(G_{\mathbf Q_\ell},\mathbf F_p).
\end{equation}
In particular the support of every global class is
finite.  There is no residual global sum relation:
the sum target in \eqref{eq:odd-target-brauer-exact}
has trivial character, whereas the coefficient
module here selects the nontrivial \(\chi\)-component.

\subsection{The cyclic classes give the local basis vectors}

For \(r\ne\ell\), the local restriction of the
cyclotomic \(\ell\)-coordinate is unramified.
The restriction of \(\delta_\ell\) therefore factors
through \(G_{\mathbf Q_r}/I_r\simeq\widehat{\mathbf Z}\),
whose second \(\mathbf F_p\)-cohomology is zero.
Thus its only possible support is at \(\ell\).

At \(\ell\), its restriction is nonzero.  To see
this directly, write \(v=v_\ell\).  Local class
field theory identifies the maximal abelian
pro-\(p\) quotient of \(G_{\mathbf Q_\ell}\) with
\[
                             \mathbf Z_p\times C_{p^v}.
\]
The two factors come from the uniformizer and the
\(p\)-primary torsion in \(\mathbf Z_\ell^\times\).
The cyclotomic \(\ell\)-coordinate restricts to an
isomorphism on that torsion factor.  This is the
usual unit action in local cyclotomic reciprocity
\cite[I, Theorem~1.1, p.~20]{MilneCFT}; it also follows
directly from \cite[p.~380, formula~(2)]{LubinTate1965}.
If the restricted extension class were zero, this
character into \(C_{p^v}\) would lift to a
homomorphism into \(C_{p^{v+1}}\).
But every homomorphism from the torsion factor
\(C_{p^v}\) into \(C_{p^{v+1}}\) has image in
\(pC_{p^{v+1}}\), whose image modulo \(p^v\)
cannot generate \(C_{p^v}\).  This is impossible.
The restriction is therefore nonzero, and is a
basis of its one-dimensional local cohomology.

Together with
\eqref{eq:odd-target-localization-isomorphism},
this proves that the arithmetic images of the
\(\delta_\ell\) form the claimed global basis.

\subsection{The exact two-place correction}

Let \(\eta_\ell:G_{\mathbf Q_\ell}\to\mathbf F_p\)
be the unramified character taking arithmetic
Frobenius to one.  With the normalization
\eqref{eq:odd-target-cyclic-extension}, one has
\begin{equation}\label{eq:odd-target-local-cyclic-cup}
 \operatorname{loc}_\ell(\delta_\ell)
                   =a_\ell\,\gamma_\ell\smile\eta_\ell.
\end{equation}
We check the sign and the exponent.  Since
\(\mu_{p^v}\subset\mathbf Q_\ell\), write
\[
       \chi_p(g)=1+p^v\tau(g)\pmod{p^{v+1}},
       \qquad \tau=a_\ell\eta_\ell.
\]
Choose the lift of \(\gamma_\ell\) to the cyclic
character with values in \(\mathbf Z/p^v\mathbf Z\),
and choose a cochain lift of it to
\(\mathbf Z/p^{v+1}\mathbf Z\).
The connecting cocycle for trivial coefficients is
\(\delta_\ell\).  If instead the middle coefficient
module has cyclotomic action, its differential
acquires the additional term
\(\tau(g)\gamma_\ell(h)\).
The latter connecting homomorphism is zero: under
the root-of-unity identifications its coefficient
sequence is
\[
                  1\longrightarrow\mu_p
                   \longrightarrow\mu_{p^{v+1}}
                     \xrightarrow{\ z\mapsto z^p\ }\mu_{p^v}
                       \longrightarrow1,
\]
and Kummer theory identifies its map on first
cohomology with the surjection
\(\mathbf Q_\ell^\times/(\mathbf Q_\ell^\times)^{p^{v+1}}
\to\mathbf Q_\ell^\times/(\mathbf Q_\ell^\times)^{p^v}\).
Therefore
\[
       \operatorname{loc}_\ell(\delta_\ell)
             =-\tau\smile\gamma_\ell
             =\gamma_\ell\smile\tau,
\]
where the last equality is graded commutativity
at the odd prime \(p\).  This proves
\eqref{eq:odd-target-local-cyclic-cup}.

For distinct \(\ell,m\in\mathcal S_p\), the
character \(\gamma_m\) at \(\ell\) is unramified,
with Frobenius value \(\gamma_m(\ell)\).  Thus
the local contribution of
\(\gamma_\ell\smile\gamma_m\) at \(\ell\)
is \(\gamma_m(\ell)\gamma_\ell\smile\eta_\ell\).
At \(m\) its contribution is
\(-\gamma_\ell(m)\gamma_m\smile\eta_m\).
All other local contributions vanish: both
characters are unramified there, or the local
second cohomology is zero.  The same argument for
\(\gamma_0\smile\gamma_\ell\) leaves only the
contribution \(-\gamma_0(\ell)
\gamma_\ell\smile\eta_\ell\) at \(\ell\).
Using \eqref{eq:odd-target-local-cyclic-cup} and
the localization isomorphism proves
\eqref{eq:odd-target-two-place-correction}.
It proves uniqueness as well, because the cyclic
classes have independent local supports.

\subsection{Transgression and the basis of the target}

Apply the five-term sequence to
\(1\to G_{K_2}\to G_{\mathbf Q}\to A\to1\):
\[
 0\to H^1(A,\mathbf F_p)\to H^1(G_{\mathbf Q},\mathbf F_p)
       \to Y_p\xrightarrow{\ \operatorname{tr}\ }
           H^2(A,\mathbf F_p)\to H^2(G_{\mathbf Q},\mathbf F_p).
\]
Kronecker--Weber identifies
\(G_{\mathbf Q}^{\rm ab}\) with \(U\)
\cite[I, Theorem~4.16, pp.~49--50]{MilneCFT}.
The already proved degree-one inflation isomorphism
therefore makes the first map in this sequence
an isomorphism.  Transgression is consequently
injective and its image is exactly the kernel
of arithmetic inflation in degree two.

By \eqref{eq:odd-target-unit-H2}, every class in
\(H^2(U,\mathbf F_p)\) has a unique finite
decomposition into cyclic classes and pairwise
cups.  The cyclic summand maps isomorphically
onto \(H^2(G_{\mathbf Q},\mathbf F_p)\).
Subtracting \(D_{ij}\) from each cup gives
a basis for the inflation kernel.  Injectivity
of transgression supplies its unique inverse
basis in \(Y_p\), proving the theorem.

This determines the full arithmetic target.
It makes no assertion that a target basis vector
is the restriction of a character of full GT.

\section{All odd coefficients and the dyadic metabelian quotient}
\label{sec:uniform-dyadic-normalization}

The integral measure gives the exact arithmetic image in every odd
weight. Compactness then removes the entire dyadic metabelian multiplier
while preserving the arithmetic factorization.

Use the dyadic beta multiplier $B_g(X,Y)$ and the integral
homomorphisms $\tau_{2r+1}$ of
Lemma~\ref{lem:integral-boundary-measure}. Put
\[
 C_0=D,\qquad
 C_r=D\cap\bigcap_{j=1}^{r}\ker\tau_{2j+1},\qquad
 J(C_r)=C_r\cap H,\qquad C_\infty=\bigcap_{r\ge0}C_r.
\]
Here $C_1=D_3$. The finite leading-term calculations below use
$D_5=C_2$ and $D_7=C_3$. The groups $C_r$ impose only odd
metabelian coordinates; mixed coefficients will be treated separately.

\begin{theorem}\label{thm:uniform-odd-arithmetic-image}
For every integer $r\ge1$, the coefficient
\[
 m_{2r+1}(g)=[S^{2r}T]f_g\qquad(g\in C_{r-1})
\]
is a continuous homomorphism of conjugation weight
$\lambda_2^{2r+1}$, with
\[
 m_{2r+1}(C_{r-1})=m_{2r+1}(J(C_{r-1}))
                 =2^{2r+1}\mathbf Z_2.
\]
Its kernel is $C_r$. In particular each inclusion
$C_r\subset C_{r-1}$ is proper, all these subgroups are closed and
normal in $G$, and
\[
 G=C_rH\quad(r\ge0),\qquad G=C_\infty H.
\]
Moreover
\[
 C_\infty=\{g\in D:B_g(X,Y)=1\},\qquad
 B\simeq C_\infty/
       \overline{[G,C_\infty](C_\infty\cap H)}.
\]
Here $B_g$ is its dyadic power-series component; no assertion about
other components of the adelic beta function is made.
\end{theorem}

\begin{proof}
Boundary reflection makes $q_g(e^u)$ even. Since it has constant
term zero, vanishing of $\tau_3,\ldots,\tau_{2r-1}$ is equivalent
to vanishing of the coefficients of $q_g(1+S)$ below $S^{2r}$.
On $C_{r-1}$ one therefore has
\[
 \tau_{2r+1}=\frac{(2r)!}{2}\,m_{2r+1}.
\]
The asserted additivity and weight law follow, since multiplication
by a nonzero integer on $\mathbf Z_2$ is injective.
Lemma~\ref{lem:uniform-odd-coordinate-bound} gives the upper bound
$m_{2r+1}\in2^{2r+1}\mathbf Z_2$.

Choose the arithmetic elements from the cubic normalization: $\sigma\in H\cap\ker\lambda$ with
$\rho_2^I(\sigma)=2$ and all odd-prime rational Kummer coordinates
zero, and $h\in H$ with $t=\lambda(h)$ and $t_2=3$.
Then $m_3(\sigma)$ is odd. Lemma~\ref{lem:integral-boundary-measure}
gives
\[
 \tau_{2r+1}(\sigma)\equiv m_3(\sigma)\equiv1\pmod2
 \quad\text{for every }r\ge1.
\]
For a positive odd integer $j$ define
\[
 \Delta_j(z)=hzh^{-1}z^{-t^j},\qquad
 v_r=\Delta_{2r-1}\Delta_{2r-3}\cdots\Delta_3\Delta_1(\sigma).
\]
The exponent $t^j$ lies in $\widehat{\mathbf Z}^{\times}$, and the
power is taken in the closed procyclic subgroup generated by $z$.
The operator $\Delta_1$ kills every rational Kummer coordinate,
so $v_1\in D\cap H$. Each later operator kills the indicated
odd-weight homomorphism and preserves its predecessors' kernels.
Thus $v_r\in J(C_{r-1})$.

The exact value of its next coordinate is
\begin{equation}\label{eq:uniform-odd-arithmetic-witness}
 m_{2r+1}(v_r)=
 \frac{2\prod_{j=0}^{r-1}(3^{2r+1}-3^{2j+1})}{(2r)!}
                         \tau_{2r+1}(\sigma).
\end{equation}
For $k\ge1$, one has
$v_2(3^{2k}-1)=3+v_2(k)$.
For odd $k$, factor $9^k-1$ by $9-1$ and note that the remaining
sum is odd. Repeatedly squaring then adds one to the valuation,
because every $9^k+1$ has valuation one. Consequently the numerator
product in \eqref{eq:uniform-odd-arithmetic-witness}, before its
factor two, has valuation
\[
 \sum_{k=1}^{r}(3+v_2(k))=3r+v_2(r!).
\]
The denominator satisfies
$v_2((2r)!)=r+v_2(r!)$, by separating its even and odd factors.
The value in \eqref{eq:uniform-odd-arithmetic-witness} therefore has
valuation exactly $2r+1$. The closed arithmetic image contains all
profinite powers of this element, hence the whole ideal
$2^{2r+1}\mathbf Z_2$. Together with the upper bound, this proves
the two exact image equalities and identifies the kernel.

The weight laws prove normality, and the nonzero arithmetic value
proves properness. Equality of the full and arithmetic images gives
\[
 C_{r-1}=C_rJ(C_{r-1}).
\]
Starting with $G=DH$, induction gives $G=C_rH$ for all $r$.
For any fixed $g\in G$, the sets
\[
 \{a\in H:ga^{-1}\in C_r\}
\]
are nonempty, closed and nested in the compact group $H$.
Their intersection is nonempty. This proves $G=C_\infty H$.

Finally, vanishing of all $\tau_{2r+1}$ is equivalent to $q_g=0$.
The beta differential identity then gives
$B_g^{-1}Y\partial_YB_g=0$. In characteristic zero this makes
$B_g$ independent of $Y$, and $B_g(X,1)=1$ makes it identically
one. The converse follows by differentiation.
The relative-normalization argument used for
Corollary~\ref{cor:cubic-radical-normalization} applies to the
closed normal subgroup $C_\infty$ and the factorization just
proved, giving the displayed canonical quotient for $B$.
\end{proof}

\begin{proposition}\label{prop:uniform-odd-character-kernels}
For $s\ge0$, the kernel of restriction
\[
 H^1(C_s,\mathbf F_2)^G\longrightarrow
 H^1(C_\infty,\mathbf F_2)^G
\]
is the one-dimensional space generated by
\[
 \zeta_s(g)=m_{2s+3}(g)/2^{2s+3}\pmod2.
\]
For the original normalized group $N$, the kernel of restriction
to $C_\infty$ is
\[
 \left(\bigoplus_{q\text{ odd prime}}\mathbf F_2\psi_q\right)
                     \oplus\mathbf F_2\eta_3.
\]
Each of these kernels has injective restriction to the corresponding
arithmetic subgroup.
\end{proposition}

\begin{proof}
The character $\zeta_s$ is defined and invariant by the exact image
and weight law, and it vanishes on $C_{s+1}$. Its arithmetic value
on $v_{s+1}$ is one.
Let $\chi$ be an invariant character on $C_s$ vanishing on
$C_\infty$. Continuity makes $\ker\chi$ open. Compactness and the
nested intersection imply $C_r\subseteq\ker\chi$ for some
$r>s$: otherwise the compact sets $C_r\setminus\ker\chi$ would
have nonempty intersection.

Suppose $r\ge s+2$ and $\chi$ kills $C_r$. Its restriction to
$C_{r-1}$ is zero or the unique nonzero character of the quotient
$C_{r-1}/C_r\simeq\mathbf Z_2$. But
\[
 v_r=h v_{r-1}h^{-1}v_{r-1}^{-t^{2r-1}},
 \qquad v_{r-1}\in C_{r-2}\subseteq C_s.
\]
Invariance and the odd exponent give $\chi(v_r)=0$, whereas the
nonzero quotient character has value one by
\eqref{eq:uniform-odd-arithmetic-witness}. The restriction must
therefore be zero. Descending induction shows that $\chi$ kills
$C_{s+1}$, proving the first assertion.

For a character on $N$, the same compactness and descent argument
works until $C_1=D_3$: every $v_{r-1}$ used with $r\ge2$ lies
in $D\subseteq N$. Thus the kernel on $N$ is exactly its already
computed kernel on restriction to $D_3$, namely the displayed
rational Kummer span and $\eta_3$. The converse is immediate from
$C_\infty\subseteq D_3$. The arithmetic value of $\zeta_s$ and
the previously proved independence of the rational Kummer
characters and $\eta_3$ prove the last assertion.
\end{proof}


\section{The full dyadic beta quotient as a measure algebra}
\label{sec:dyadic-beta-measure-algebra}

The integral measure describes the entire beta quotient, including its
integral lattice and arithmetic image. This also explains the successive
odd-coordinate ideals without treating their images one at a time.
Put
\[
 \mathcal U=1+4\mathbf Z_2,\qquad
 \mathcal A=\mathbf Z_2[[\mathcal U]],\qquad
 \mathfrak a=\ker(\mathcal A\longrightarrow\mathbf Z_2).
\]
We regard the completed group algebra as the integral measures on
$\mathcal U$, with convolution as multiplication and total mass as
augmentation. For $v\in\mathbf Z_2^\times$, let $\langle v\rangle$
be the unique member of $\{v,-v\}$ lying in $\mathcal U$.
The symbol $[t]$ denotes the Dirac measure at $t\in\mathcal U$.

\begin{theorem}\label{thm:full-dyadic-measure-algebra}
There is a continuous additive map
\[
 \nu:M=\ker\lambda\longrightarrow\mathcal A
\]
with the following properties:
\begin{align*}
 \nu_g(\mathcal U)&=\rho_{2,2}^I(g)/2,\\
 \int_{\mathcal U}t^{2k}\,d\nu_g(t)
   &=(1-2^{2k})\tau_{2k+1}(g)\qquad(k\ge1),\\
 \nu_{hgh^{-1}}&=\lambda_2(h)[\langle\lambda_2(h)\rangle]\nu_g.
\end{align*}
Here $\rho_{2,2}^I$ is the two-adic component of $\rho_2^I$.
Its kernel is exactly the kernel of the dyadic beta multiplier, and
\[
 \nu(M\cap H)=\nu(M)=\mathcal A,\qquad
 \nu(D\cap H)=\nu(D)=\nu(N)=\mathfrak a.
\]
Consequently the continuous additive quotients are
\[
 M/\{g\in M:B_g=1\}\simeq\mathcal A,\qquad
 D/C_\infty\simeq\mathfrak a,
\]
with the displayed conjugation action. In particular every dyadic beta
multiplier on either domain is attained on its arithmetic subgroup.
\end{theorem}

\begin{proof}
Use the integral measure $\mu_g$ with transform $q_g(X)$ from
Lemma~\ref{lem:integral-boundary-measure}, which exists on all of $M$.
Set $\nu_g=\mu_g|_{1+4\mathbf Z_2}$. The unnormalized duplication
identity, with the same conventions as
\eqref{eq:measure-duplication}, is
\[
 B_g(X,X)=X^{2\rho_{2,2}^I(g)}B_g(-1,X).
\]
Its logarithmic derivative gives
\[
 q_g(X)+q_g(-X)=2q_g(X^2)-2\rho_{2,2}^I(g).
\]
Thus, as integral measures on the additive group $\mathbf Z_2$,
\begin{equation}\label{eq:unnormalized-measure-doubling}
 \mu_g|_{2\mathbf Z_2}=[2]_*\mu_g-\rho_{2,2}^I(g)\delta_0.
\end{equation}
The total mass of $\mu_g$ is $q_g(1)=0$. Its mass on the odd
integers is therefore $\rho_{2,2}^I(g)$. Negation symmetry exchanges
the two odd residue classes modulo four, giving the claimed mass of
$\nu_g$. Integrating $t^{2k}$ in
\eqref{eq:unnormalized-measure-doubling} removes the term at zero.
The same symmetry and the definition of $\tau_{2k+1}$ give the
stated moment formula.

Additivity follows from additivity of $q_g$ on $M$, including its
measure interpretation. For $v=\lambda_2(h)$, the boundary
conjugation formula is $q_{hgh^{-1}}(X)=v q_g(X^v)$.
It sends $\mu_g$ to $v[v]_*\mu_g$. If $v\equiv1\pmod4$,
restriction to $\mathcal U$ is the pushforward by $v$; if
$v\equiv-1\pmod4$, use negation symmetry to replace this by the
pushforward by $-v$. This proves the asserted action in both cases.
All maps used are continuous in the inverse-limit topologies.

If $\nu_g=0$, its total mass gives $\rho_{2,2}^I(g)=0$.
Symmetry makes the whole odd restriction of $\mu_g$ zero, and
\eqref{eq:unnormalized-measure-doubling} gives
$\mu_g=[2]_*\mu_g$. Iteration shows that $\mu_g$ is supported on
$\bigcap_n2^n\mathbf Z_2=\{0\}$. Its total mass is zero, so it is
the zero measure. Hence $q_g=0$. The beta differential identity
$D_Y\log B_g(X,Y)=q_g(XY)-q_g(Y)$ and $B_g(X,1)=1$ give $B_g=1$.
The converse is immediate. This proves the kernel assertion; it
concerns the dyadic component and does not determine the other
components of $\rho_2^I$.

Choose $\sigma\in H\cap M$ with $\rho_2^I(\sigma)=2$ and all
odd-prime rational Kummer coordinates zero, and $h\in H$ with
$t=\lambda(h)$ and $t_2=3$, as in the earlier arithmetic witnesses.
Write $\gamma=[-3]\in\mathcal A$. The element $-3$ topologically
generates $\mathcal U$: one has
$v_2((-3)^{2^n}-1)=n+2$, so its image generates every quotient
$(1+4\mathbf Z_2)/(1+2^{n+2}\mathbf Z_2)$. Therefore
\[
 \mathcal A\simeq\mathbf Z_2[[Z]],\qquad \gamma=1+Z,
 \qquad \mathfrak a=Z\mathcal A.
\]
The measure $\nu_\sigma$ has augmentation one, so it is a unit of
this power-series ring. The compact additive subgroup
$\nu(H\cap M)$ is a closed $\mathbf Z_2$-submodule, and
conjugation by $h$ acts as $3\gamma$. Multiplication by three is
invertible on a closed $\mathbf Z_2$-submodule. The image therefore
contains every polynomial in $\gamma$ times $\nu_\sigma$, and by
closure every element of $\mathcal A\nu_\sigma=\mathcal A$.
This proves the full arithmetic image.

Finally let $z=h\sigma h^{-1}\sigma^{-t}$. The full profinite
Kummer laws put $z$ in $D\cap H$, and
\[
 \nu_z=3(\gamma-1)\nu_\sigma.
\]
This generates the ideal $\mathfrak a$ as a closed
$\mathcal A$-module. Normality of $D$ and the same conjugation
argument show that $\nu(D\cap H)$ contains this ideal.
Conversely $D\subseteq N$, and the mass formula puts $\nu(N)$
in $\mathfrak a$. All the claimed image equalities follow.
The kernel on $D$ is $C_\infty$ by its beta characterization.
A continuous surjection from a compact group to a Hausdorff group
induces the claimed topological quotient isomorphisms.
\end{proof}

\begin{corollary}\label{cor:unnormalized-odd-images}
Put $T_0=M$ and
\[
 T_r=M\cap\bigcap_{j=1}^r\ker\tau_{2j+1}\qquad(r\ge1).
\]
For every $r\ge1$, the raw coefficient
$m_{2r+1}=[S^{2r}T]f$ is additive on $T_{r-1}$, has conjugation
weight $\lambda_2^{2r+1}$, and has exact full and arithmetic image
\[
 m_{2r+1}(T_{r-1})=m_{2r+1}(T_{r-1}\cap H)
                  =2^{2r-2-v_2(r)}\mathbf Z_2.
\]
Its kernel is $T_r$. More precisely, in the preceding identification
of $\mathcal A$ one has
\[
 \nu(T_r)=\nu(T_r\cap H)
       =\prod_{j=1}^r(\gamma-9^j)\mathcal A,
 \qquad
 \nu(C_r)=\nu(C_r\cap H)
       =\prod_{j=0}^r(\gamma-9^j)\mathcal A.
\]
\end{corollary}

\begin{proof}
The moment $\int t^{2j}\,d\nu(t)$ is the evaluation homomorphism
$\operatorname{ev}_{9^j}:\mathcal A\to\mathbf Z_2$ that sends
$\gamma$ to $9^j$. Its scalar factor $1-2^{2j}$ in the theorem
is a unit. For $\alpha\in2\mathbf Z_2$, ordinary power-series
division gives
\[
 \ker\bigl(\mathbf Z_2[[Z]]\xrightarrow{Z\mapsto\alpha}
                  \mathbf Z_2\bigr)=(Z-\alpha).
\]
Indeed the coefficients of $(F(Z)-F(\alpha))/(Z-\alpha)$ are
convergent sums of integral coefficients times powers of $\alpha$.
Applying this division successively proves that the intersection
of the evaluation kernels at the distinct points $9^j$ is the
product of the displayed linear factors. At each step the
previous factors have nonzero value in $\mathbf Z_2$ at the new
point, so cancellation is valid even though that value need not
be a unit. Theorem~\ref{thm:full-dyadic-measure-algebra} lifts
every measure in either resulting ideal to the corresponding
arithmetic subgroup, proving the two exact ideal formulas.

On $T_{r-1}$, reflection and vanishing of the lower moments make
$q_g(1+S)$ start in degree $2r$, so
\[
 m_{2r+1}=\frac{2}{(2r)!}\tau_{2r+1}.
\]
Evaluation of the ideal for $T_{r-1}$ at $\gamma=9^r$ has image
\[
 \prod_{j=1}^{r-1}(9^r-9^j)\mathbf Z_2.
\]
Its valuation is $3(r-1)+v_2((r-1)!)$, by
$v_2(9^k-1)=3+v_2(k)$. After the factor $2/(2r)!$ the valuation is
\[
 1+3(r-1)+v_2((r-1)!)-r-v_2(r!)=2r-2-v_2(r).
\]
The empty product for $r=1$ gives the same formula. The exact
images and kernel follow. Additivity and the weight law follow
from the corresponding properties of the moments.
\end{proof}

\begin{corollary}\label{cor:beta-quadratic-invariants}
Among continuous invariant quadratic characters which factor through
the dyadic beta multiplier, the space on $M$ is one-dimensional,
generated by $\rho_{2,2}^I/2\pmod2$. On each of $N$ and $D$ it
is one-dimensional, generated by $\eta_3=m_3/8\pmod2$.
\end{corollary}
\begin{proof}
Invariance under the chosen $h$ forces a quadratic character of
$\mathcal A$ to kill $2\mathcal A$ and $(3\gamma-1)\mathcal A$.
Modulo two the latter is $(\gamma-1)\mathcal A$, so the invariant
quotient is $\mathcal A/(2,\gamma-1)\simeq\mathbf F_2$.
Augmentation modulo two is invariant under every multiplier by
the action in the theorem and is nonzero.
For $\mathfrak a=(\gamma-1)\mathcal A$, the analogous quotient
is $\mathfrak a/(2\mathfrak a+\mathfrak a^2)\simeq\mathbf F_2$.
Write $\nu=(\gamma-1)A$. Its second moment is $8A(9)$, whence
\[
 m_3=\tau_3=-\frac{8}{3}A(9),\qquad
 \eta_3=A(1)\pmod2.
\]
This character is invariant and nonzero by the arithmetic image.
The claimed descriptions follow from the exact beta quotients.
\end{proof}

\section{Uniform arithmetic saturation of the dyadic word image}
\label{sec:fr-uniform-dyadic-lie}

The characteristic-two generators admit a uniform arithmetic
construction. The useful multiplier is $-3$ in the $2$-adic
component: its power substitution has no quadratic or cubic
correction modulo two.

Put $M=\ker\lambda$, $H_M=H\cap M$, and
\[
 F^nM=\{g\in M:f_g-1\text{ has no dyadic Magnus term of degree }<n\}.
\]
Every element of $M$ belongs to $F^3M$. Let
\[
 R=\bigcap_{n\ge3}F^nM
   =\ker\bigl(M\longrightarrow
                \operatorname{Aut}(\widehat F_2^{(2)})\bigr),
 \qquad R_D=D\cap R.
\]
The equality in this definition is justified below.

\begin{theorem}\label{thm:fr-uniform-dyadic-saturation}
For every degree $n\ge3$, the full and arithmetic leading images
are the entire integral Grothendieck--Teichm\"uller Lie lattice:
\[
 \operatorname{in}_n(F^nM)
 =\operatorname{in}_n(F^nM\cap H)
 =\mathcal L_n,
 \qquad
 \mathcal L_n=
 \mathfrak{grt}_1(\mathbf Q_2)_n
          \cap\operatorname{Lie}_{\mathbf Z_2}(S,T)_n.
\]
There are arithmetic elements $z_{2r+1}\in H_M\cap F^{2r+1}M$
whose reduced leading terms are precisely the characteristic-two
free generators of Willwacher's Theorem~1.3. Their leading terms
freely generate $\mathcal L=\bigoplus_n\mathcal L_n$ over
$\mathbf Z_2$.

One has
\[
 M=RH_M,\qquad G=RH,\qquad G=R_DH.
\]
In particular the full-GT and arithmetic images in
$\operatorname{Aut}(\widehat F_2^{(2)})$ coincide.
The normalized subgroup $R_D$ is closed and normal, and
\[
 B\simeq R_D/\overline{[G,R_D](R_D\cap H)}.
\]
The kernel of restriction from $X_2=H^1(N,\mathbf F_2)^G$ to
$R_D$ is exactly
\[
 \left(\bigoplus_{q\ {\rm odd\ prime}}\mathbf F_2\psi_q\right)
                   \oplus\mathbf F_2\eta_3.
\]
Moreover $R_D$ is a proper subgroup of infinite index in
$E_{15}=C_\infty\cap F^{15}M$.
\end{theorem}

The proof uses the defining-equation and characteristic-two
freeness results of
\cite[equations (1)--(6), p.~1; \S1.1 and Theorem~1.3, p.~2;
\S4.3, p.~9]{Willwacher}. It does not require an a priori
arithmetic realization of a rational or integral Lie solution.

\begin{lemma}\label{lem:fr-integral-leading-basis}
The finite-rank module $\mathcal L_n$ is saturated in
$\operatorname{Lie}_{\mathbf Z_2}(S,T)_n$, and reduction gives
an injection
\[
 \mathcal L_n/2\mathcal L_n
       \hookrightarrow\mathfrak{grt}_1(\mathbf F_2)_n.
\]
If a collection of elements of $\mathcal L_n$ reduces to a
basis of the space on the right, it is a $\mathbf Z_2$-basis
of $\mathcal L_n$.
\end{lemma}

\begin{proof}
The intersection with a rational subspace is saturated:
if $2v$ belongs to it and $v$ is integral, then $v$ belongs
to the same rational subspace. The integral defining equations,
including cabling, remain valid after reduction modulo two.
Their finite integral matrix description is recalled in
\cite[\S1.1, p.~2]{Willwacher}. Saturation makes the displayed
reduction map injective.

For the final assertion, the proposed reductions are linearly
independent in $\mathcal L_n/2\mathcal L_n$, and their images
span its ambient target. They therefore form a basis of
$\mathcal L_n/2\mathcal L_n$. Nakayama's lemma makes their lifts
generate $\mathcal L_n$; their independence follows by
successive reduction modulo powers of two, or by equality of
the finite ranks.
\end{proof}

The necessary leading-map lemma applies here in every degree.
For completeness, project the full profinite pentagon to the
pro-$2$ and then rational prounipotent completions, as in
\cite[\S3, p.~554, immediately before Corollary~12]{Furusho2010}.
The substitution $x=e^S$, $y=e^T$ preserves the first nonzero
term. That term is primitive and hence is a free Lie polynomial
over $\mathbf Q_2$. Taking the first nonzero degree in the group
pentagon gives the linear pentagon in
$\mathfrak t_4(\mathbf Q_2)$: the degree-one class of a pure
braid generator is $t_{ij}$, products add these classes, and
products of two nonconstant word terms have degree at least
$2n>n$. Duality and the hexagon likewise give antisymmetry
and the three-cycle identity. Cabling is redundant in
characteristic zero
\cite[p.~1, after equation~(6)]{Willwacher}.
Thus the leading term belongs to
$\mathfrak{grt}_1(\mathbf Q_2)_n$.
Its integral associative coefficients are integral Lyndon
coefficients, because the Lyndon expansions are unit triangular.
It consequently belongs to $\mathcal L_n$. This is only a
necessary inclusion; the following construction supplies the
arithmetic realization.

\begin{lemma}\label{lem:fr-minus-three-jet}
Choose $h\in H$ with $t=\lambda(h)$ and $t_2=-3$.
Put $C=[S,T]$. For $z\in F^nM$, $n\ge3$, write
$q=\operatorname{in}_n(z)$ and define, using the full
profinite exponent,
\[
 \Delta_n(z)=hzh^{-1}z^{-t^n}.
\]
Then $\Delta_n(z)\in F^{n+1}M$, and its Magnus series satisfies
\begin{equation}\label{eq:fr-minus-three-jet}
 f_{\Delta_n(z)}
 \equiv1+\{C,\overline q\}
        \pmod{\,2,\ (S,T)^{n+3}}.
\end{equation}
In particular every coefficient through degree $n+1$ is
even, while its degree-$n+2$ term modulo two is the indicated
Ihara bracket.
\end{lemma}

\begin{proof}
We specify the ambient convention to check the jet precisely.
For a pair $(s,f)$ let
\[
 \Phi_{s,f}(x)=x^s,\qquad
 \Phi_{s,f}(y)=f^{-1}y^sf.
\]
Composition gives
\[
 (s,f)(u,g)=(su,\ f\,\Phi_{s,f}(g)).
\]
Consequently, in the ambient group of these automorphisms,
\[
 h=(1,f_h)(t,1).
\]
The power pair $(t,1)$ is an automorphism, with inverse
$(t^{-1},1)$, so this factorization is legitimate even though
its factors need not individually belong to GT.

The degree-two coefficient of $f_h$ is $\tfrac13 C$.
Indeed, put $m=(t_2-1)/2=-2$ in the multiplier hexagon
and compute its logarithm through degree two. The three
$f_h$ terms contribute $3c_2(h)C$, and the intervening
product $z^my^mx^m$ contributes
$-\tfrac12m(m+1)C$. Thus
$c_2(h)=m(m+1)/6=1/3$. In particular its reduction is $C$.

Conjugation by $(t,1)$ replaces the variables of $f_z$ by
$x^t,y^t$. Modulo two,
\[
 (1+S)^{-3}-1=S+O(S^4),\qquad
 (1+T)^{-3}-1=T+O(T^4),
\]
since $\binom{-3}{2}=6$ and $\binom{-3}{3}=-10$ are even.
Therefore this substitution leaves every term of $f_z$
through degree $n+2$ unchanged modulo two, including its
unknown terms in degrees $n+1$ and $n+2$.

For multiplier-one ambient pairs with leading terms $A$ and
$Q$ of degrees $d$ and $n$, the first term of their group
commutator is
\[
 \{A,Q\}=[A,Q]+\partial_AQ-\partial_QA,\qquad
 \partial_A(S)=0,\quad\partial_A(T)=[T,A].
\]
This follows directly by subtracting the two products
$a\,\Phi_a(b)$ and $b\,\Phi_b(a)$: in degree $d+n$ their
difference is $AQ-QA+\partial_AQ-\partial_QA$.
Taking $d=2$ shows that conjugation by $(1,f_h)$ adds exactly
$\{C,\overline q\}$ in degree $n+2$ modulo two. Higher terms
of $f_h$ contribute only in degrees at least $n+3$.

Finally, through degree $n+2$, a profinite power of $z$
has series $1+r_2(f_z-1)$: all product and Ihara
self-substitution corrections begin in degree $2n>n+2$.
Since $t_2^n$ is odd, multiplication by $z^{-t^n}$ cancels
the unchanged $f_z-1$ jet modulo two. This proves
\eqref{eq:fr-minus-three-jet}. Integrally, the degree-$n$
term cancels as well, because conjugation has weight
$t_2^n$ on the leading quotient; hence
$\Delta_n(z)\in F^{n+1}M$.
\end{proof}

\begin{proposition}\label{prop:fr-all-odd-arithmetic-generators}
There exist elements $z_{2r+1}\in H_M\cap F^{2r+1}M$
whose leading reductions are the generators
\[
 s_3=[S+T,[S,T]],\qquad
 s_{2r+3}=\{C,s_{2r+1}\}
\]
of $\mathfrak{grt}_1(\mathbf F_2)$.
In each degree the leading terms of the corresponding
arithmetic Hall commutator words form a
$\mathbf Z_2$-basis of $\mathcal L_n$.
\end{proposition}

\begin{proof}
Take the established arithmetic $\sigma\in H_M$ with
$a=m_3(\sigma)$ odd, and put $z_3=\sigma$.
Its leading reduction is $s_3$. Theorem~1.3 of
\cite[p.~2]{Willwacher} says that the $s_{2r+1}$ freely
generate the characteristic-two Lie algebra under the
Ihara bracket.

Suppose arithmetic generators have been constructed through
the odd degree $n$, and their leading reductions are the
specified generators. Every free Hall word of even weight
$n+1$ is a bracket word in lower odd generators; no new
generator has that even weight. Realize these words by
group commutators in the already constructed $z_j$.
Their leading reductions form the full free Hall basis of
$\mathfrak{grt}_1(\mathbf F_2)_{n+1}$.
Lemma~\ref{lem:fr-integral-leading-basis} makes their leading
terms a $\mathbf Z_2$-basis of $\mathcal L_{n+1}$.
When this characteristic-two space is zero, the integral
space is zero as well.

By Lemma~\ref{lem:fr-minus-three-jet},
$\Delta_n(z_n)$ starts in degree $n+1$ and its degree-$n+1$
term is even. Saturation implies that this term lies in
$2\mathcal L_{n+1}$. Express it in the arithmetic Hall basis
just constructed. All coordinates are even.
Multiply $\Delta_n(z_n)$ by inverse profinite powers of
those Hall words, with these even $2$-adic exponents, to
kill its whole degree-$n+1$ term.
Let the result be $z_{n+2}$.

Each correcting word starts in degree $n+1$. Its even
power is trivial modulo two through degree $n+2$, because
$2(n+1)>n+2$. Products of the corrections have the same
property. Thus the correction does not change the
degree-$n+2$ term in
\eqref{eq:fr-minus-three-jet}. The new arithmetic word
starts in degree $n+2$ and has leading reduction
$\{C,s_n\}=s_{n+2}$.

In odd degree $n+2$, this new generator together with all
Hall bracket words in lower generators reduces to the
complete characteristic-two Hall basis. Apply
Lemma~\ref{lem:fr-integral-leading-basis} once more.
This proves the induction, both in even and odd degrees.
It also proves that the free Lie algebra over $\mathbf Z_2$
on the arithmetic leading terms of the $z_{2r+1}$ maps
isomorphically to $\mathcal L$, degree by degree.
\end{proof}

We record two properties of this construction needed for the
normalized conclusion. Every correcting Hall word in the
induction is a commutator in $M$. Therefore all its rational
Kummer coordinates vanish and its dyadic beta multiplier is
one. If $c_3=1$ and
\[
 c_{n+2}=(t-t^n)c_n\in\widehat{\mathbf Z},
\]
then, for every rational prime including two,
\begin{equation}\label{eq:fr-uniform-kummer-collinearity}
 \rho_q(z_n)=c_n\rho_q(\sigma).
\end{equation}
The same beta-trivial corrections leave every beta moment
unchanged. Hence $z_{2r+1}$ has zero moments in preceding
odd weights, and its depth-one coefficient is
\begin{equation}\label{eq:fr-uniform-primitive-depth}
 [S^{2r}T]f_{z_{2r+1}}
 =\frac{2\prod_{j=1}^{r-1}
       \bigl((-3)^{2r+1}-(-3)^{2j+1}\bigr)}{(2r)!}
                         \tau_{2r+1}(\sigma)\ne0.
\end{equation}
The established all-$M$ moment parity makes the last moment
odd. In particular the depth-one coefficient is nonzero on
each odd integral generator. Its valuation is
\[
 3(r-1)+v_2((r-1)!)+1-v_2((2r)!)
                   =2r-2-v_2(r).
\]
The empty product for $r=1$ has the same interpretation.

\begin{proof}[Proof of Theorem~\ref{thm:fr-uniform-dyadic-saturation}]
Every full leading term belongs to $\mathcal L_n$.
Proposition~\ref{prop:fr-all-odd-arithmetic-generators}
realizes an entire integral basis there by arithmetic
words. Their profinite powers and products realize all
integral linear combinations. This proves the two leading
image equalities.

For any $g\in M$, correct its leading terms successively
by these arithmetic basis words. For each $n$ this gives
an $h_n\in H_M$ with $g h_n^{-1}\in F^nM$.
The sets
\[
 \{h\in H_M:g h^{-1}\in F^nM\}
\]
are nonempty, closed, and nested in the compact group
$H_M$. Their intersection is nonempty. Thus $M=RH_M$.
The cyclotomic multiplier of $H$ is all of
$\widehat{\mathbf Z}^{\times}$, so $G=MH=RH$.

The pro-$2$ Magnus embedding is faithful. Indeed, for every
finite $2$-group quotient $Q$, the augmentation ideal of
$\mathbf F_2[Q]$ is nilpotent. This follows by induction using
a central element $c$ of order two: the ideal generated by $c-1$
has square zero, and quotienting by it gives the group algebra
of $Q/\langle c\rangle$. Substitution of $x-1,y-1$ in this
finite group algebra therefore factors through a finite Magnus
truncation. A word with trivial Magnus series is consequently
trivial in every finite $2$-group quotient, hence in the free
pro-$2$ group. The intersection of the word filtration is thus
exactly the condition that the pro-$2$ word $f_g$ be one. This is also the
kernel of the multiplier-one action on the free pro-$2$
group. Indeed, if that action is the identity, $f_g$
commutes with $y$. In the Magnus algebra the centralizer
of $T$ consists of series in $T$ alone; and $f_g(1,y)=1$
because $f_g$ is in the commutator subgroup. Therefore
$f_g=1$. The converse is immediate. This justifies the
definition of $R$ and the equality of full and arithmetic
dyadic action images. It also shows that $R$ is closed
and normal in $G$.

We now preserve the full rational Kummer normalization.
Let $U$ be the closed subgroup of $H_M$ generated by all
$z_{2r+1}$, and put
\[
 A=\overline{[U,U]}.
\]
Both beta and every rational Kummer coordinate are
homomorphisms on $M$, so $A\subset H\cap C_\infty$.
Moreover \eqref{eq:fr-uniform-kummer-collinearity} gives
$z_n\sigma^{-c_n}\in D$.
In the quotient by $\overline{[G,D]}$, the image of $D$
is central, and all the $z_n$ are central factors times
profinite powers of the same $\sigma$. Their images
therefore commute. By continuity,
\begin{equation}\label{eq:fr-arithmetic-derived-in-relative-commutator}
 A\subset\overline{[G,D]}\subset\overline{[G,N]}.
\end{equation}

In the free integral Lie algebra on the leading
$z_{2r+1}$, the bracket submodule has a Hall basis
consisting of words of generator length at least two.
Every such bracket has zero depth-one coefficient.
In an odd degree the remaining abelianization is
rank one, and its generator has nonzero depth-one
coefficient by
\eqref{eq:fr-uniform-primitive-depth}.
It follows, integrally, that the kernel of the
depth-one functional on $\mathcal L_n$ is exactly
the bracket submodule; in even degrees the whole
space is that submodule.

For an element of $C_\infty$, all depth-one coefficients
vanish because its dyadic beta multiplier is one.
Its leading term can consequently be corrected entirely
by arithmetic Hall bracket words, all lying in $A$.
Repeating this correction degree by degree remains
inside $C_\infty$. For each $c\in C_\infty$, the closed
sets
\[
 \{a\in A:c a^{-1}\in F^nM\}
\]
are nonempty and nested in the compact group $A$.
Their intersection gives
\[
 C_\infty=(C_\infty\cap R)A=R_DA.
\]
The last equality uses the fact that a trivial dyadic
word has trivial dyadic beta multiplier, so
$C_\infty\cap R=D\cap R$.
The established factorization $G=C_\infty H$ now gives
$G=R_DH$. To check the displayed presentation of $B$ explicitly,
put $K=N\cap H$ and $C_R=\overline{[G,R_D]}$.
Then $N=R_DK$. In $G/C_R$, the image of $R_D$ is central
and $G=R_DH$, so the image of $[G,N]$ is contained in the
image of $K$. Hence
\[
 \overline{[G,N]K}=C_RK.
\]
These products are closed because $C_R$ is closed and normal
and $K$ is compact. Since $C_R\subset R_D$,
\[
 R_D\cap C_RK=C_R(R_D\cap H).
\]
Restriction of the quotient map from $N$ therefore yields
the claimed canonical isomorphism for $B$.

Every $\chi\in X_2$ kills $A$ by
\eqref{eq:fr-arithmetic-derived-in-relative-commutator}.
Thus $\chi$ vanishes on $R_D$ if and only if it vanishes
on $C_\infty$. The already computed restriction kernel
for $C_\infty$ is precisely the rational Kummer span
plus $\mathbf F_2\eta_3$, proving the asserted kernel.

Finally, the arithmetic word
$[z_3,[z_3,z_9]]$ has a nonzero primitive leading term
of degree fifteen, since its reduction is a nonzero
free Hall word. It lies in $A\cap F^{15}M\subset E_{15}$.
Its procyclic image modulo $R_D$ is infinite: multiplying
its nonzero leading coefficient by a nonzero $2$-adic
exponent cannot give zero. Hence $R_D$ is proper of
infinite index in $E_{15}$.
\end{proof}

\begin{remark}
The conclusion concerns the dyadic free-group action and
the normalized domain carrying the relative obstruction.
It does not identify $R_D$ with the trivial group or
prove that $B$ has no nonzero quadratic characters.
\end{remark}

\subsection{Arithmetic Kummer coordinates in the dyadic action kernel}

% Arithmetic pro-two-kernel Kummer coordinates.
% Required bibliography key: HainMatsumoto2003.
\begin{proposition}[Arithmetic Kummer coordinates in the pro-two kernel]
\label{prop:arithmetic-kummer-kernel-axes}
Let \(G=\widehat{\mathrm{GT}}\), let \(H\) be its arithmetic subgroup, and put
\[
 M=\ker\lambda,\qquad N=\ker(\lambda,\rho_2^I).
\]
\[
 D=\bigcap_{q\ {\rm prime}}\ker(\rho_q^I|_M),
\]
\[
 R_M=\{g\in M:f_g=1\text{ in the free pro-two group}\}.
\]
For every odd rational prime \(q\), there is an element
\(k_q\in R_M\cap N\cap H\) whose only nonzero rational Kummer
coordinate is
\[
 \rho_{q,2}^I(k_q)=-2.
\]
In particular the tuple of dyadic odd-prime coordinates on the closed
arithmetic subgroup of \(R_M\cap N\) killing every odd-primary
Kummer coordinate has image exactly
\begin{equation}
 \prod_{q\ {\rm odd}}2\mathbf Z_2.                       \label{eq:arithmetic-kernel-axes-1}
\end{equation}
If \(G=DH\), and \(R_D=R_M\cap D\), the natural map
\[
 (R_M\cap N)/R_D\longrightarrow N/D
\]
is an isomorphism. If, in addition,
\(M=R_M(M\cap H)\), then
\[
 D=R_D(D\cap H),\qquad G=R_DH.
\]
\end{proposition}

\begin{proof}
Let \(\varepsilon_2\in\widehat{\mathbf Z}\) have dyadic component
one and every odd-primary component zero.
The tangential action on the pro-two fundamental group of
\(\mathbf P^1-\{0,1,\infty\}\) is unramified at every odd prime:
apply \cite[\S8.5, p.~151]{HainMatsumoto2003} with \(N=1\) and
\(\ell=2\). This is the action into automorphisms, so no inner
automorphism ambiguity is involved.

Choose an inertia element \(\tau_q\) with dyadic Kummer coordinate
one at \(q\). It has coordinate zero at every other rational
prime: these primes are \(q\)-adic units, whose two-power roots
lie in unramified extensions, as do the two-power roots of unity.
Choose complex conjugation \(c\) fixing the compatible positive
real roots of the positive rational primes. Its Kummer coordinates
are zero and its multiplier is \(-1\). Then
\[
 k_q=[c,\tau_q]^{\varepsilon_2}
\]
has full multiplier one and trivial pro-two action.
The crossed Kummer law gives its dyadic coordinate \(-2\) at \(q\)
and zero elsewhere, while the power \(\varepsilon_2\) kills all
odd-primary coordinates. This proves the asserted individual
coordinates.

An arithmetic element of \(M\) fixes every rational square root,
because the quadratic fields are cyclotomic. Its dyadic coordinates
are therefore even. The image in \eqref{eq:arithmetic-kernel-axes-1} is closed by compactness,
contains each coordinate axis and every \(\mathbf Z_2\)-multiple
of that axis, and hence contains the closure of their direct sum.
This proves \eqref{eq:arithmetic-kernel-axes-1}, including its simultaneous assertion.

The normalized pro-two action has pro-two image: its actions on
the finite algebras
\(\mathbf F_2\langle S,T\rangle/(S,T)^d\) are upper unitriangular
and detect its parameter. Consequently
\begin{equation}
 a^{\,1-\varepsilon_2}\in R_M,\qquad
 a^{\varepsilon_2}R_M=aR_M\qquad(a\in M).                \label{eq:arithmetic-kernel-axes-2}
\end{equation}
These powers are taken in the procyclic closure of \(a\).

Suppose \(G=DH\). For \(n\in N\), choose \(a\in N\cap H\)
with \(na^{-1}\in D\). The element \(a^{1-\varepsilon_2}\)
lies in \(R_M\cap N\cap H\) and realizes all its odd-primary
coordinates. Formula \eqref{eq:arithmetic-kernel-axes-1} realizes the remaining dyadic
odd-prime tuple inside the same group, without changing any
other coordinate. Thus \(R_M\cap N\to N/D\) is surjective.
Its kernel is \(R_M\cap D=R_D\), and the compact quotient theorem
gives the stated isomorphism.

Finally take \(d\in D\), and suppose
\(M=R_M(M\cap H)\). Choose \(a\in M\cap H\) with the same
pro-two parameter as \(d\). The dyadic beta measure, and hence
its mass \(\rho_{2,2}^I/2\), depends only on that parameter;
thus \(\rho_{2,2}^I(a)=0\). Replace \(a\) by
\(a^{\varepsilon_2}\), preserving its parameter by \eqref{eq:arithmetic-kernel-axes-2}
and killing all its odd-primary Kummer coordinates.
Use \eqref{eq:arithmetic-kernel-axes-1} to cancel its even dyadic coordinates at every odd
prime by an element \(k\in R_M\cap N\cap H\).
Then \(a^{\varepsilon_2}k^{-1}\in D\cap H\) still matches \(d\).
This gives \(D=R_D(D\cap H)\), and \(G=DH\) gives \(G=R_DH\).
\end{proof}

% Requires bibliography key Willwacher2026 and the established uniform
% leading-lattice, dyadic-measure, and arithmetic kernel-axis results.
\begin{theorem}[The normalized pro-two image]
\label{thm:normalized-protwo-free}
Put \(M=\ker\lambda\) and
\[
 R_M=\{g\in M:f_g=1\text{ in the free pro-two group}\},
 \qquad P=M/R_M.
\]
Then \(P\) is free pro-two on a countable family converging to
the identity. This family can be chosen arithmetically.

More precisely, choose the square-corrected arithmetic family
\(a_3=\sigma,a_5,a_7,\ldots\), whose leading reductions are the
free generators of
\(\mathfrak{grt}_1(\mathbf F_2)\), and for which
\(a_{j+2}\) is \([h_*,a_j]\) followed by inverse-square corrections,
where \(\lambda_2(h_*)=-3\).
If \(F_r\) is the free pro-two group on
\(z_3,\ldots,z_{2r+1}\), with transition maps killing the last
generator, then
\begin{equation}
 \varprojlim_r F_r \ \xrightarrow{\ \sim\ }\ P,
 \qquad z_j\longmapsto a_jR_M.                          \label{eq:normalized-protwo-1}
\end{equation}
In particular \(H^2_{\rm cts}(P,\mathbf F_2)=0\).
\end{theorem}

\begin{proof}
Write \(P_j=\operatorname{in}_j(a_j)\).
The integral leading-basis theorem gives an isomorphism
\begin{equation}
 \operatorname{Lie}_{\mathbf Z_2}(Z_3,Z_5,\ldots)
 \longrightarrow\bigoplus_n L_n,\qquad Z_j\longmapsto P_j,
                                                               \label{eq:normalized-protwo-2}
\end{equation}
where the letters have weights \(j\) and the target bracket is
Ihara's bracket.
For injectivity, divide a nonzero homogeneous polynomial by the
largest common power of two in its Lie-basis coefficients.
Its nonzero reduction remains nonzero after evaluation, by
\cite[Theorem~1.3]{Willwacher}.
Surjectivity is the integral leading-basis assertion.

Let \(\mathcal P=\varprojlim F_r\).
It has the following universal property: a sequence \(b_j\to1\)
in any pro-two group \(Q\) determines a unique continuous map
\(\mathcal P\to Q\) sending \(z_j\) to \(b_j\).
Indeed, in any finite quotient of \(Q\), all sufficiently late
\(b_j\) are one; hence the required map factors through a
finite-rank \(F_r\). These maps are compatible over the finite
quotients. The generators topologically generate \(\mathcal P\),
by compactness and their generation in every \(F_r\).

Use the weighted integral Magnus expansion \(z_j=1+Z_j\), with
\(\operatorname{wt}(Z_j)=j\). It is faithful.
For completeness, a finite-rank pro-two Magnus expansion is
already faithful modulo two: a nonidentity element survives in
a finite two-group \(Q\), and substituting \(Z_j\mapsto q_j-1\)
in \(\mathbf F_2[Q]\) detects it.
The augmentation ideal of this group algebra is nilpotent.
Inductively, a central element of order two gives a square-zero
ideal, and the quotient is the group algebra of a smaller
two-group. Thus the power-series substitution is legitimate.
Faithfulness in \(\mathcal P\) follows by projection to a
finite-rank \(F_r\).

For \(w\ne1\) in \(\mathcal P\), let \(Q_w\ne0\) be its first
nonzero weighted Magnus term, of weight \(n\).
It is an ordinary integral Lie polynomial.
In fact
\(\Delta(Z_j)=Z_j\otimes1+1\otimes Z_j+Z_j\otimes Z_j\);
the last term has larger weight, so the lowest term of a
group-like element is primitive for the ordinary coproduct.
Over \(\mathbf Q_2\), decomposition by word length and the
Dynkin operator give Lie membership.
The leading coefficients one in the Lyndon basis then give
integral Lie membership.

The images \(a_jR_M\) converge to one, since their parameters
start in degree \(j\). They therefore define the map \eqref{eq:normalized-protwo-1}.
Its image is closed and contains all integral leading-basis
corrections, so successive approximation and compactness make
it surjective.

To prove injectivity, act on the original Magnus algebra and put
\(A_j=\Phi_{a_j}-1\).
This operator raises degree by at least \(j\), and its component
of degree shift \(j\) is
\[
 D_{P_j}(S)=0,\qquad D_{P_j}(T)=[T,P_j].
\]
Substitution \(Z_j\mapsto A_j\) into the weighted completed
algebra is well-defined and agrees with the group
action. Its first possible nonzero term on \(w\) is
\[
 Q_w(D_{P_3},D_{P_5},\ldots)
    =D_{\,Q_w(P_3,P_5,\ldots)}.
\]
Here the equality follows from
\([D_A,D_B]=D_{\{A,B\}}\).
The polynomial inside the last derivation is nonzero by \eqref{eq:normalized-protwo-2}.
Its derivation is nonzero as well: a homogeneous polynomial
commuting with \(T\) is a pure power of \(T\), whereas these Lie
polynomials vanish when \(S=0\).
Thus \(w\) acts nontrivially, proving \eqref{eq:normalized-protwo-1}.

Finally, every extension of \(P\) by a finite two-group splits.
In the extension group choose an open subgroup meeting the
finite kernel trivially. It maps homeomorphically to an open
subgroup of \(P\).
The late generators have unique lifts there converging to one;
choose arbitrary lifts of the finitely many remaining
generators. The universal property of \eqref{eq:normalized-protwo-1} gives a section.
In particular every central extension by \(\mathbf F_2\)
splits, proving the cohomological assertion.
\end{proof}

\begin{corollary}[Frattini action and characters of the pro-two image]
\label{cor:protwo-image-invariant-characters}
For the family in Theorem~\ref{thm:normalized-protwo-free}, put
\(V=P/\overline{P^2[P,P]}\) and \(T=h_*-1\).
Then
\begin{equation}
 V\simeq\mathbf F_2[[T]],\qquad
 [a_{3+2r}]=T^r[a_3],\qquad h_*=1+T.                   \label{eq:normalized-protwo-3}
\end{equation}
The dyadic mass
\(\mu=\rho_{2,2}^I/2:P\to\mathbf Z_2\) is surjective.
For \(P_0=\ker\mu\),
\begin{equation}
 N/(N\cap R_M)\simeq D/R_D\simeq P_0,                  \label{eq:normalized-protwo-4}
\end{equation}
and
\begin{equation}
 H^1(P,\mathbf F_2)^G=\mathbf F_2(\mu\bmod2),
 \qquad
 H^1(P_0,\mathbf F_2)^G=\mathbf F_2\eta_3.             \label{eq:normalized-protwo-5}
\end{equation}
Thus every invariant quadratic character of \(M,N,\) or \(D\)
that factors through its full pro-two parameter image already
factors through the dyadic beta quotient.
\end{corollary}

\begin{proof}
The Frattini quotient of the inverse limit in \eqref{eq:normalized-protwo-1} is the product
of its generator copies of \(\mathbf F_2\).
Indeed its closed Frattini subgroup maps onto the Frattini
subgroup in every finite-rank quotient, and compactness
identifies it with the intersection of their inverse images.
The inverse-square corrections disappear in this quotient,
so
\([a_{j+2}]=(h_*-1)[a_j]\), proving \eqref{eq:normalized-protwo-3}.
An invariant continuous functional is consequently determined
by its value at \([a_3]\).
The mass modulo two is invariant and takes value one there,
proving the first assertion of \eqref{eq:normalized-protwo-5}.

The image of \(N\) is \(P_0\): for a lift \(m\) of a point of
\(P_0\), the number \(z=\rho_2^I(m)/2\) has zero dyadic
component. The pro-two image of \(\sigma^z\) is therefore one,
and \(m\sigma^{-z}\in N\) gives a lift.
The image of \(D\) is also \(P_0\).
Choose an arithmetic lift, take its dyadic-primary power, and
cancel its even dyadic odd-prime Kummer coordinates by the
arithmetic kernel axes. Its coordinate at two is already zero.
This proves \eqref{eq:normalized-protwo-4}.

Let \(\ell_r\) be the continuous coordinate functional dual to
\([a_{3+2r}]\). Then
\(H^1(P,\mathbf F_2)=\bigoplus_{r\ge0}\mathbf F_2\ell_r\),
with \(\ell_0=\mu\bmod2\).
The procyclic closure of \(a_3\) splits \(P\to\mathbf Z_2\).
Restriction thus gives the \(G\)-equivariant exact sequence
\[
 0\longrightarrow\mathbf F_2\ell_0
 \longrightarrow H^1(P,\mathbf F_2)
 \longrightarrow H^1(P_0,\mathbf F_2)^P
 \longrightarrow0.
\]
Any \(G\)-invariant character of \(P_0\) is \(P\)-invariant.
For \(\ell=\sum c_r\ell_r\), invariance of its class modulo
\(\ell_0\) under \(h_*\) requires
\[
 \ell\circ(h_*-1)=\sum_{r\ge1}c_r\ell_{r-1}
                   \in\mathbf F_2\ell_0.
\]
It follows that \(c_r=0\) for \(r\ge2\), so this invariant
space has dimension at most one.
The measure theorem gives the nonzero invariant character
\(\eta_3=m_3/8\bmod2\) on \(P_0\), already nonzero on
\(D\cap H\). This proves the remaining assertion of \eqref{eq:normalized-protwo-5}.
\end{proof}

\begin{remark}
The same finite-extension splitting argument gives the exact
sequence
\[
 0\longrightarrow H^1(P,\mathbf F_2)
 \longrightarrow H^1(M,\mathbf F_2)
 \longrightarrow H^1(R_M,\mathbf F_2)^M
 \longrightarrow0.
\]
After taking \(G/M\)-invariants, a possible obstruction remains
in \(H^1_{\rm cts}(G/M,H^1(P,\mathbf F_2))\).
These results do not prove that the relative abelian quotient
of \(R_D\) is zero.
\end{remark}

\begin{proposition}[Full cyclotomic action on the Frattini module]
\label{prop:full-cyclotomic-frattini-action}
Retain \(P=M/R_M\), \(V=P/\Phi(P)\), and the coordinate
\(V=\mathbf F_2[[T]]\), where \(T=h_*-1\) and
\(\lambda_2(h_*)=-3\).
Put
\[
 U=G/M=\widehat{\mathbf Z}^{\,\times},\qquad
 \Gamma=1+4\mathbf Z_2,\qquad
 U_0=\{\pm1\}\times\prod_{q\ {\rm odd}}\mathbf Z_q^\times.
\]
Thus \(U=\Gamma\times U_0\).
For \(u\in U\), let \(\langle u_2\rangle\) be the unique member
of \(\{u_2,-u_2\}\) in \(\Gamma\).
The complete action on \(V\) is
\begin{equation}
 u\cdot F(T)
   =(1+T)^{\,\log\langle u_2\rangle/\log(-3)}F(T).        \label{eq:cyclotomic-frattini-1}
\end{equation}
In particular \(U_0\) acts trivially.
\end{proposition}

\begin{proof}
The group \(M\) acts on \(P\) by inner automorphisms and hence
trivially on \(V\). The action factors through the abelian group
\(U\), so every action operator commutes with \(h_*\), and hence
with \(T\). By continuity it is an endomorphism of the cyclic
completed module \(\mathbf F_2[[T]]\).
An automorphism is therefore multiplication by a unit.
This gives a continuous homomorphism
\[
 U\longrightarrow\mathbf F_2[[T]]^\times;
\]
continuity follows also by applying the operator to the cyclic
vector \(1=[a_3]\).

The target is a torsion-free pro-two group.
Indeed it is \(1+T\mathbf F_2[[T]]\), and its quotients modulo
powers of \(T\) are finite two-groups.
If \(a^{2^n}=1\), then \((a-1)^{2^n}=0\); the power-series
domain forces \(a=1\).
Consequently the finite dyadic sign factor acts trivially.
For odd \(q\), the open pro-\(q\) subgroup of
\(\mathbf Z_q^\times\) has trivial image in a pro-two group,
and the remaining finite image is trivial by torsion-freeness.
Finite-support products are dense in the odd-local product,
so the whole of \(U_0\) acts trivially.

It remains to determine the action of \(\Gamma\).
Its topological generator \(-3\) acts as \(1+T\) by the exact
Frattini recursion. This gives \eqref{eq:cyclotomic-frattini-1}.
The exponent lies in \(\mathbf Z_2\), since \(-3\)
topologically generates \(1+4\mathbf Z_2\).
\end{proof}

\begin{proposition}[The residual extension obstruction]
\label{prop:explicit-residual-extension-obstruction}
Let \(W=H^1(P,\mathbf F_2)\), with its discrete continuous
\(U\)-action. There is a natural isomorphism
\begin{equation}
 H^1_{\rm cts}(U,W)
   \cong\operatorname{Hom}_{\rm cts}(U_0,\mathbf F_2).    \label{eq:cyclotomic-frattini-2}
\end{equation}
In particular the latter is a direct sum of one copy of
\(\mathbf F_2\) for the dyadic sign factor and one for each
odd rational prime.
There is an exact segment
\begin{equation}
 0\longrightarrow\mathbf F_2(\mu\bmod2)
 \longrightarrow H^1(M,\mathbf F_2)^G
 \longrightarrow H^1(R_M,\mathbf F_2)^G
 \longrightarrow
       \operatorname{Hom}_{\rm cts}(U_0,\mathbf F_2),     \label{eq:cyclotomic-frattini-3}
\end{equation}
where \(\mu=\rho_{2,2}^I/2\).
No surjectivity of the last map is asserted.
\end{proposition}

\begin{proof}
Let \(\ell_r\) be dual to \(T^r\), so
\(W=\bigoplus_{r\ge0}\mathbf F_2\ell_r\).
For the pullback action of the chosen generator,
\[
 \ell_r\circ(h_*-1)=\ell_{r-1}\quad(r\ge1),\qquad
 \ell_0\circ(h_*-1)=0.
\]
The inverse-pullback convention has the same image and kernel.
Thus generator minus one acts surjectively on \(W\), and
\(W^\Gamma=\mathbf F_2\ell_0\).

A continuous cocycle on the procyclic group \(\Gamma\) is
determined by its value on a topological generator, and every
value in \(W\) is permitted.
Indeed a vector has finite orbit of two-power size; summing
around that orbit twice gives zero, so its cocycle factors
through a sufficiently large finite cyclic quotient.
Coboundaries are the image of generator minus one.
It follows that \(H^1_{\rm cts}(\Gamma,W)=0\).

Given a \(U\)-cocycle \(z\), subtract a coboundary so that its
restriction to \(\Gamma\) is zero.
Since \(U=\Gamma\times U_0\), comparison of the cocycle identities
for \(\gamma u=u\gamma\) gives \(z(u)\in W^\Gamma\) for
\(u\in U_0\).
The action of \(U_0\) on \(W\) is trivial by the preceding
proposition. Hence this restriction is a continuous
homomorphism \(U_0\to\mathbf F_2\ell_0\).
Conversely every such homomorphism, extended by zero on
\(\Gamma\), is a \(U\)-cocycle.
Changing the normalization by a \(\Gamma\)-invariant
coboundary has no effect on \(U_0\). This proves \eqref{eq:cyclotomic-frattini-2}.

For odd \(q\), the pro-\(q\) unit subgroup contributes no
quadratic character, and the cyclic residue-unit group
contributes one. Continuity permits only finitely many
nonzero components, proving the direct-sum description.

Finally use the \(U\)-equivariant exact sequence
\[
 0\longrightarrow W
 \longrightarrow H^1(M,\mathbf F_2)
 \longrightarrow H^1(R_M,\mathbf F_2)^M
 \longrightarrow0
\]
proved by freeness of \(P\).
Taking invariants and applying \eqref{eq:cyclotomic-frattini-2}, with
\(W^U=\mathbf F_2\ell_0=\mathbf F_2(\mu\bmod2)\),
gives \eqref{eq:cyclotomic-frattini-3}.
\end{proof}

\begin{remark}
Concretely, an invariant character of \(R_M\) first extends to
\(M\). Its conjugation defect is a \(U\)-cocycle in \(W\).
An adjustment by a character of \(P\) removes the
\(\Gamma\)-defect. The remaining defect is a finite-support
quadratic character of \(U_0\), multiplied by the mass
functional. This describes the obstruction to an invariant
extension; it does not make that obstruction or the relative
abelian quotient of \(R_D\) vanish.
\end{remark}

\section{The full dyadic image}
\label{sec:full-dyadic-image}

The normalized image and its cyclotomic action determine the
whole group acting on the free pro-two fundamental group.
Write
\[
 Q_2=\operatorname{im}\bigl(G\longrightarrow
                 \operatorname{Aut}(\widehat F_2^{(2)})\bigr).
\]
The symbol \(Q_2\) denotes a group, not the field \(\mathbf Q_2\).

\begin{theorem}\label{thm:full-dyadic-free-product}
The full dyadic image is the arithmetic image, and
\[
                  Q_2\cong C_2*_{\mathrm{pro}\text{-}2}\mathbf Z_2.
\]
One may take complex conjugation as the generator of \(C_2\),
and any arithmetic element with dyadic multiplier \(-3\) as
the generator of the second factor.
Under this presentation the multiplier is the abelianization
map
\[
 Q_2\longrightarrow\mathbf Z_2^\times
        =\{\pm1\}\times(1+4\mathbf Z_2),
 \qquad c\longmapsto-1,\qquad h\longmapsto-3.
\]
Its kernel is the normalized image \(P=M/R_M\).
\end{theorem}

\subsection{Two arithmetic generators}

\begin{proof}[Generation and the cyclotomic kernel]
Every dyadic action is the identity on the mod-two Frattini
quotient of the free group. The automorphisms with this property
form a pro-two group, by the finite characteristic filtration
argument used for the prime word actions. Thus \(Q_2\) is pro-two.
Uniform arithmetic saturation gives equality of its full and
arithmetic images.

We first check that its cyclotomic kernel is exactly \(P\).
The tangential arithmetic action is unramified at every odd
prime \cite[\S8.5, p.~151]{HainMatsumoto2003}.
For an odd prime \(q\), inertia at \(q\) therefore acts trivially
on \(\widehat F_2^{(2)}\), whereas its cyclotomic image is the
whole \(q\)-local factor \(\mathbf Z_q^\times\) and is trivial
at every other prime. The closed subgroup generated by these
inertia groups has trivial dyadic action. Its full multiplier
image contains the finite-support product of the odd-local
unit groups, and compactness makes that image the whole product.
Consequently an arithmetic element with dyadic multiplier one
can have all its other multiplier coordinates cancelled without
changing its dyadic action. Its resulting lift belongs to
\(M\cap H\). This proves the kernel assertion.

Let \(c\) be complex conjugation, using the positive tangential
base point. Then \(c^2=1\), \(f_c=1\), and \(\rho_2(c)=0\).
Choose \(h\in H\) with \(\lambda_2(h)=-3\) and set
\(a_3=[c,h]\). The cocycle law gives
\[
 \rho_{2,2}(a_3)=-2\rho_{2,2}(h),\qquad
 \mu(a_3)=-\rho_{2,2}(h).
\]
The parity theorem says that \(\rho_{2,2}(h)\) is odd, since
\(-3\equiv5\pmod8\). Hence the cubic leading coefficient of
\(a_3\) is odd. Its reduction is the prescribed weight-three
free Lie generator.

Apply the uniform generator construction with this seed and
this same \(h\). Every constructed generator belongs to the
closed subgroup generated by \(c,h\): the construction uses
conjugation by \(h\), powers, and corrections made from earlier
generators. Its leading terms give the complete integral basis
in every weight. Successive correction and compactness therefore
put all of \(P\) in that closed subgroup.
The multipliers \(-1,-3\) generate \(\mathbf Z_2^\times\), so
\(c,h\) generate \(Q_2\).

For these generators the construction also gives
\[
 [a_{3+2j}]=(h-1)^j[a_3]\quad\hbox{in }P/\Phi(P).
\]
Indeed the conjugation difference in odd weight \(n\) differs
from \([h,a_n]\) by a power with even dyadic exponent, and
the further corrections are squares. All disappear in the
Frattini quotient. The arithmetic leading-basis theorem and
the normalized freeness theorem apply to this choice of seed.
\end{proof}

\subsection{A free kernel with a convergent basis}

We give the pro-two group calculation needed to detect relations.
Let
\[
 F=C_2*_{\mathrm{pro}\text{-}2}\mathbf Z_2
   =\langle s,t\mid s^2=1\rangle_{\mathrm{pro}\text{-}2},
 \qquad
 \pi:F\longrightarrow C_2\times\mathbf Z_2
\]
be the map sending \(s\) and \(t\) to the two factors, and put
\(K_F=\ker\pi\).

\begin{lemma}\label{lem:free-product-cyclotomic-kernel}
The group \(K_F\) is free pro-two on the convergent sequence
\[
 b_0=[s,t],\qquad b_{j+1}=[t,b_j]\quad(j\ge0).
\]
\end{lemma}

\begin{proof}
The kernel of \(F\to C_2\) is free pro-two on
\(t,k=sts\). This follows directly from universal properties:
if \(A\) is free pro-two on \(t,k\), the involution interchanging
these generators defines \(A\rtimes C_2\); the maps in both
directions given by \(s,t,k=sts\) are inverse.
Replacing \(k\) by \(b_0=kt^{-1}\) is a free change of basis.
Thus \(K_F\) is the kernel of the map from the free pro-two
group on \(t,b_0\) to \(\mathbf Z_2\) that sends \(t\) to one
and \(b_0\) to zero.

Let \(\mathcal P\) be free pro-two on the sequence
\(z_0,z_1,\ldots\) converging to one, concretely the inverse
limit of the free pro-two groups on its finite initial segments
with retractions killing the last generator. Define
\[
                   \alpha(z_j)=z_{j+1}z_j.
\]
On every finite initial-segment quotient this is an invertible
triangular change of free basis: solve for the generators
starting from the last one. These automorphisms and their
inverses commute with the retractions. They therefore define
an automorphism of \(\mathcal P\).
Its action on each finite-rank Frattini quotient is unipotent.
The kernel of the action of the automorphism group of a
finite-rank free pro-two group on its Frattini quotient is
pro-two. Hence the closed cyclic group generated by each
finite-rank image of \(\alpha\) is pro-two. Passing to the
inverse limit shows that powers of \(\alpha\) define a
continuous \(\mathbf Z_2\)-action on \(\mathcal P\).

Form \(\mathcal P\rtimes_\alpha\mathbf Z_2\), writing \(t\)
for the second-factor generator. It is generated by \(t,z_0\),
because \([t,z_j]=z_{j+1}\). The free pro-two group on
\(t,b_0\) maps onto it by \(b_0\mapsto z_0\).
Conversely, in that free group the recursively defined
\(b_j\) converge to one: in every finite two-group quotient,
iterated commutators eventually vanish by nilpotence.
The universal property gives a continuous map
\(\mathcal P\to A\), \(z_j\mapsto b_j\).
The recursion makes it equivariant for conjugation by \(t\),
and therefore for the full \(\mathbf Z_2\)-action by continuity.
It extends to a map of semidirect products into \(A\).
The two maps fix the generators \(t,b_0\), respectively
\(t,z_0\), and so are inverse.
The kernel of the exponent map is thus precisely
\(\mathcal P\), proving the lemma.
\end{proof}

\subsection{There are no further relations}

\begin{proof}[Completion of Theorem~\ref{thm:full-dyadic-free-product}]
Generation gives a surjection
\[
                   F\longrightarrow Q_2,
                   \qquad s\longmapsto c,\quad t\longmapsto h.
\]
It commutes with the cyclotomic quotient maps, so restricts
to a surjection \(K_F\to P\).
The preceding lemma and the exact Frattini recursion identify
its map on Frattini quotients with the identity on the
products of copies of \(\mathbf F_2\):
\[
             [b_j]\longmapsto[a_{3+2j}].
\]

We verify carefully that this forces injectivity here.
The surjection \(K_F\to P\) has a continuous section, since
\(P\) is free on a convergent basis. To construct one, choose
a descending countable neighborhood basis of open subgroups
of \(K_F\). Their images are open in \(P\). Late basis
elements of \(P\) lie in each such image. They can therefore
be lifted so that the chosen lifts converge to one; choose
arbitrary lifts for the finitely many initial exceptions.
The free universal property produces a section.
Thus \(K_F=J\rtimes P\), where \(J\) is the kernel.

If \(J\ne1\), it has a nonzero continuous character to
\(\mathbf F_2\), obtained from a nontrivial finite two-group
quotient. The orbit of this character under \(P\) is finite:
continuity and compactness of \(J\) give an open stabilizer.
Its span is a nonzero finite-dimensional \(\mathbf F_2\)-module
for a finite two-group. Such a module has a nonzero invariant
vector. For example, induct on the order of the group, first
taking the nonzero kernel of \(z-1\) for a central element
\(z\) of order two and then taking invariants under the
quotient group.
The resulting invariant character of \(J\) extends to
\(J\rtimes P\) by setting it zero on \(P\).
This contradicts the isomorphism of Frattini quotients, which
says that every character of \(K_F\) is pulled back from \(P\)
and hence vanishes on \(J\). Therefore \(J=1\).

The map \(F\to Q_2\) is injective on the kernel of its
cyclotomic quotient and is an isomorphism on that quotient.
It is consequently an isomorphism. Finally,
\(F^{\mathrm{ab}}=C_2\times\mathbf Z_2\), by the universal
property of the free product. This proves the abelianization
assertion as well.
\end{proof}

This theorem describes one fundamental-group realization.
It does not identify the full profinite group \(G\) with
its arithmetic subgroup or determine the residual quotient
supported on \(R_D\).

\section{An exact weight-five normalization}
\label{sec:weight-five-normalization}

Let $G=\widehat{\mathrm{GT}}$, let $H$ be its Galois image, and write
the radical and cubic normalizations as
\[
 D=N_{\mathrm{rad}},\qquad
 D_3=D\cap\ker m_3,\qquad J_3=D_3\cap H.
\]
The established factorizations give $G=D_3H$.  All Magnus coefficients
in this section are $2$-adic.  For $g=(1,f)\in D_3$, the pro-$2$
Magnus expansion of $f$ starts in degree five: the degree-three and
degree-four Hall coordinates are respectively
$(m_3,-m_3)$ and $(-m_3,m_3,-m_3)$ by the preceding low-weight
calculation.  Define
\[
 m_5(g)=[S^4T]f,\qquad x=1+S,\quad y=1+T.
\]

\begin{theorem}\label{thm:weight-five-exact-image}
The exact integral images are
\[
 m_5(D_3)=m_5(J_3)=32\mathbf Z_2.
\]
In particular,
\[
 \eta_5:D_3\longrightarrow\mathbf F_2,
 \qquad \eta_5(g)=m_5(g)/32\pmod2,
\]
is a continuous $G$-invariant character with nonzero arithmetic
restriction.  The kernel of $m_5$,
\[
 D_5=D_3\cap\ker m_5,
 \qquad J_5=D_5\cap H=J_3\cap\ker m_5,
\]
is a proper closed normal subgroup satisfying
\[
 D_3=D_5J_3,\qquad G=D_5H.
\]
The relative abelian quotient refines canonically to
\[
 B\simeq D_3/\overline{[G,D_3]J_3}
       \simeq D_5/\overline{[G,D_5]J_5}.
\]
\end{theorem}

The exact image is the case $r=2$ of
Theorem~\ref{thm:uniform-odd-arithmetic-image}. We give its explicit
arithmetic value and the resulting extension obstruction.

\begin{proposition}\label{prop:weight-five-upper-bound}
For $g\in D_3$, one has $m_5(g)\in32\mathbf Z_2$.
\end{proposition}

\begin{proof}
Since $D_3=C_1$, apply Theorem~\ref{thm:uniform-odd-arithmetic-image}
with $r=2$. Equivalently, the integral measure has zero mass and zero
second moment, so the integer-valued polynomial of degree four in
Lemma~\ref{lem:uniform-odd-coordinate-bound} gives the factor $2^5$.
\end{proof}

\subsection{An arithmetic generator of the exact image}

The arithmetic lower bound uses a weight-five homomorphism on
the entire multiplier-one subgroup.  For such an element, write
\[
 a_3=[S^2T]f,\qquad b_5=[S^4T]f,
 \qquad \tau_5(f)=12b_5-11a_3\in\mathbf Z_2.
\]

\begin{lemma}\label{lem:tau-five}
The map $\tau_5$ is a continuous homomorphism on $\ker\lambda$,
and
\[
 \tau_5(hzh^{-1})=\lambda_2(h)^5\tau_5(z)
 \qquad(h\in G,\ z\in\ker\lambda).
\]
On $D_3$, it is equal to $12m_5$.
\end{lemma}

\begin{proof}
Let $R(X,Y)$ be the Alexander coefficient of $f$ on $[x,y]$,
and let
\[
 B_f(X,Y)=1+(X-1)(Y-1)R(X,Y).
\]
The action on the rank-one Alexander module is semilinear.
Consequently, on the multiplier-one subgroup the cofactors
multiply, and conjugation by multiplier $t$ sends $B_f(X,Y)$
to $B_f(X^t,Y^t)$.  It follows that
\[
 [u^4v]\log B_f(e^u,e^v)
\]
is an additive weight-five $\mathbf Q_2$-valued homomorphism.
Only terms linear in $Y-1$ contribute to this coefficient, so
the logarithm can be omitted for this extraction.  The boundary
Magnus coefficients and the earlier identity $[S^3T]f=-a_3$ give
\[
 (X-1)R(X,1)
 =a_3(X-1)^2-a_3(X-1)^3+b_5(X-1)^4+O((X-1)^5).
\]
Now
\[
 [u^4](e^u-1)^2=\frac7{12},\quad
 [u^4](e^u-1)^3=\frac32,\quad
 [u^4](e^u-1)^4=1.
\]
Thus the extracted homomorphism is $b_5-11a_3/12$.
Multiplication by $12$ proves the asserted integral formula
and its exact transformation rule.  The last assertion uses
$a_3=0$ on $D_3$.
\end{proof}

\begin{proposition}\label{prop:weight-five-arithmetic-generator}
There exists $w\in J_3$ with $m_5(w)\in32\mathbf Z_2^\times$.
\end{proposition}

\begin{proof}
Simultaneous rational Kummer surjectivity over
$\mathbf Q^{\mathrm{ab}}$ supplies $\sigma\in H\cap\ker\lambda$
with
\[
 \rho_2^I(\sigma)=2\in\widehat{\mathbf Z},\qquad
 \rho_q^I(\sigma)=0\quad(q\text{ odd}).
\]
The established boundary congruence
$\rho_2^I\equiv-2m_3\pmod4$ makes $a_3(\sigma)$ odd.
Therefore $\tau_5(\sigma)=12b_5(\sigma)-11a_3(\sigma)$ is odd.

Choose $h\in H$ with $\lambda_2(h)=3$ and put
$t=\lambda(h)\in\widehat{\mathbf Z}^{\times}$.  For
$c\in\widehat{\mathbf Z}$ and $z\in H\cap\ker\lambda$, define
\[
 \mathcal D_c(z)=hzh^{-1}z^{-c}.
\]
The power $z^{-c}$ is the standard continuous profinite power
in the procyclic closure of $z$.

Every rational Kummer coordinate of $\mathcal D_t(\sigma)$
is zero, because its conjugation weight is $t$.  Set
\[
 w=\mathcal D_{t^3}\bigl(\mathcal D_t(\sigma)\bigr).
\]
All rational Kummer coordinates remain zero.  Its cubic
coordinate also vanishes, since conjugation on that coordinate
has weight $t_2^3$.  Hence $w\in J_3$.

Lemma~\ref{lem:tau-five} and the continuity of powers give the
exact equality
\[
 \tau_5(w)=(3^5-3^3)(3^5-3)\tau_5(\sigma)
          =51840\tau_5(\sigma).
\]
This equality holds for the exact homomorphism $\tau_5$ with the
stated conjugation weight.  Since $w\in D_3$,
\[
 m_5(w)=\frac{51840}{12}\tau_5(\sigma)
       =4320\tau_5(\sigma)
       =32\cdot135\,\tau_5(\sigma),
\]
which is a generator of $32\mathbf Z_2$.
\end{proof}

\begin{proof}[Proof of Theorem~\ref{thm:weight-five-exact-image}]
On $D_3$, $m_5$ is a continuous additive homomorphism, with
conjugation weight $\lambda_2^5$; this follows either from its
leading Magnus degree or from Lemma~\ref{lem:tau-five} and
the torsion-freeness of $\mathbf Z_2$.  The upper bound and
the arithmetic generator show that its compact images on
both $D_3$ and $J_3$ are exactly $32\mathbf Z_2$.  Dividing
by $32$ and reducing modulo two gives a $G$-invariant
character, since every $2$-adic unit has odd fifth power.

Equality of the images gives $D_3=D_5J_3$.  Normality follows
from the scalar conjugation rule, properness from the
arithmetic generator, and $G=D_5H$ follows from $G=D_3H$.
For the quotient assertion, if $G=EH$ with $E$ closed normal,
then
\[
 E/\overline{[G,E](E\cap H)}
 \simeq G/\overline{[G,E]H}.
\]
If $E=E'(E\cap H)$ and $G=E'H$, the closures
$\overline{[G,E]H}$ and $\overline{[G,E']H}$ coincide: in
the quotient by the latter, $H$ is trivial and the image
of $E'$ is central.  Apply this with $E=D_3$ and $E'=D_5$.
\end{proof}

If $J_3=G_{L_3}$ and $L_5=\overline{\mathbf Q}^{\,J_5}$,
then $\operatorname{Gal}(L_5/L_3)\simeq\mathbf Z_2$.
The character $\eta_5$ has nonzero arithmetic restriction on $D_3$.
The same witness shows precisely why it does not give another invariant
character on the original normalized subgroup.

\begin{corollary}\label{cor:weight-five-nonextension}
The character $\eta_5$ does not extend to a $G$-invariant continuous
$\mathbf F_2$-character of $D=N_{\mathrm{rad}}$, and hence does not
extend to one of $N$.  Nevertheless $D_5$ is a strictly smaller
normalization with the same relative abelian quotient.  The kernel of
restriction from $X_2$ to $H^1(D_5,\mathbf F_2)$ is exactly
\[
 \bigoplus_{q\text{ odd prime}}\mathbf F_2\psi_q
       \oplus\mathbf F_2\eta_3.
\]
The kernel of restriction from $H^1(D,\mathbf F_2)^G$ to
$H^1(D_5,\mathbf F_2)$ is exactly $\mathbf F_2\eta_3|_D$.
For every odd prime $p$, restriction from $H^1(N,\mathbf F_p)^G$
to $H^1(D_5,\mathbf F_p)$ is injective.
\end{corollary}

\begin{proof}
Use the elements $h,\sigma,t$ of
Proposition~\ref{prop:weight-five-arithmetic-generator}, and put
$z=\mathcal D_t(\sigma)\in D\cap H$.  Its witness is
\[
 w=hzh^{-1}z^{-t^3}=[h,z]z^{1-t^3}\in D_3\cap H,
 \qquad \eta_5(w)=1.
\]
Since $t$ is a profinite unit, $1-t^3\in2\widehat{\mathbf Z}$.
The continuous profinite power $z^{(1-t^3)/2}$ is defined, so
$w$ belongs to the closed subgroup generated by $[G,D]$ and the
squares in $D$.  Every invariant quadratic character of $D$ kills
that subgroup.  It therefore takes the value zero at $w$, which is
incompatible with extending $\eta_5$.  Restricting an invariant
character of $N$ to $D$ proves the second nonextension assertion.

If $\chi\in X_2$ vanishes on $D_5$, its restriction to $D_3$
factors through $D_3/D_5\simeq\mathbf Z_2$.  That restriction is
either zero or $\eta_5$.  The latter possibility has just been
excluded, so $\chi$ vanishes on $D_3$.  Conversely every character
vanishing on $D_3$ vanishes on $D_5$.  Corollary~\ref{cor:cubic-radical-joint-image}
identifies this common kernel as the displayed direct sum.
The same argument on $D$, using $D/D_3\simeq8\mathbf Z_2$,
gives the one-dimensional kernel there.

For odd $p$, a character of $D_3$ vanishing on $D_5$ must be zero,
because $D_3/D_5$ is pro-$2$.  A character of $N$ vanishing on
$D_3$ factors through the joint rational Kummer and cubic quotient.
The cubic factor is pro-$2$, and conjugation by multiplier $-1$
acts by inversion on every rational Kummer coordinate.  Invariance
therefore kills the remaining factor in $\mathbf F_p$.
\end{proof}

\section{The exact weight-seven normalization}
\label{sec:weight-seven-partial}

Let $G=\widehat{\mathrm{GT}}$, let $H$ be its Galois image, and use
the established normalization
\[
 D_5=N_{\mathrm{rad}}\cap\ker m_3\cap\ker m_5,
 \qquad J_5=D_5\cap H,
 \qquad G=D_5H.
\]
All Magnus coefficients below are $2$-adic, with $x=1+S$ and $y=1+T$.
We prove that the first possible nonzero degree on $D_5$ is seven and
write $m_7(g)=[S^6T]f$ for $g=(1,f)\in D_5$.

\begin{theorem}\label{thm:weight-seven-bounds}
Every word belonging to $D_5$ has zero Magnus components below degree
seven. The map $m_7:D_5\to\mathbf Z_2$ is a continuous homomorphism and
\[
 m_7(hgh^{-1})=\lambda_2(h)^7m_7(g)
 \qquad(h\in G,\ g\in D_5).
\]
Its arithmetic and full images are exactly
\begin{equation}\label{eq:weight-seven-image-interval}
 m_7(J_5)=m_7(D_5)=128\mathbf Z_2.
\end{equation}
The corresponding degree-seven Hall images both equal
$8\mathbf Z_2\mathbf L$, where $\mathbf L$ is the primitive
integral vector specified below.
\end{theorem}

\subsection{The first possible leading term}

Use standard Lyndon brackets on $x<y$, ordered first by degree and
then lexicographically, and let $h_w$ denote the corresponding integral
Hall coordinate. The standard factorization $w=uv$ uses the longest
proper Lyndon suffix, and the group word is $L_w=[L_u,L_v]$.

\begin{lemma}\label{lem:weight-seven-leading-spaces}
Suppose that $f\in\gamma_5$ satisfies the two-cycle and three-cycle
relations with multiplier one. Put $a=[S^4T]f$. Its Hall coordinates
in degrees five and six are
\begin{equation}\label{eq:weight-seven-lower-vector}
 a\left(1,-2,\tfrac32,2,\tfrac12,-1,
 -2,\tfrac72,-\tfrac32,-4,-\tfrac12,\tfrac32,
 \tfrac72,\tfrac32,-2\right).
\end{equation}
In particular, $a=0$ implies that $f$ starts in degree seven.

If $f$ starts in degree seven and satisfies the same two relations,
its degree-seven Hall vector has the form $a\mathbf v+b\mathbf w$,
where $a=[S^6T]f$, $b$ is the sixth degree-seven Hall coordinate, and
\begin{align}
\mathbf v={}&\left(1,-3,5,5,-2,0,-12,-5,-2,\tfrac23,
                  -5,5,8,3,-\tfrac{11}3,4,3,-1\right),
                  \label{eq:weight-seven-v}\\
\mathbf w={}&\left(0,0,0,0,0,1,1,0,0,\tfrac13,
                  -1,-1,-1,0,-\tfrac13,0,0,0\right).
                  \label{eq:weight-seven-w}
\end{align}
These are necessary leading-term conditions; they assert no lifting
of an arbitrary truncated solution to a profinite automorphism.
\end{lemma}

\begin{proof}
The finite identities in Section~\ref{subsec:weight-seven-leading-certificate}
give \eqref{eq:weight-seven-lower-vector} and
\eqref{eq:weight-seven-v}--\eqref{eq:weight-seven-w} directly from
coefficients of the two defining relations through degree nine.
They are exact rational identities between explicitly specified
integral coefficient forms. Consequently they apply to the
$2$-adic Hall coordinates of a word.

On $D_5$, the preceding cubic calculation first gives $f\in\gamma_5$,
and $m_5=0$. Equation~\eqref{eq:weight-seven-lower-vector} therefore
makes every degree-five and degree-six coordinate zero. This proves
the asserted absence of lower terms on $D_5$.
\end{proof}

\subsection{The additional leading pentagon relation}

The original pure-braid pentagon has the form
\[
 f(x_{12},x_{23}x_{24})f(x_{13}x_{23},x_{34})
 =f(x_{23},x_{34})f(x_{12}x_{13},x_{24}x_{34})f(x_{12},x_{23});
\]
see \cite[Remark~1.2, p.~98]{Furusho2003}.
For a word starting in degree seven, its degree-seven term $p$
therefore satisfies
\begin{equation}\label{eq:weight-seven-leading-pentagon}
\begin{split}
 p(t_{12},t_{23}+t_{24})+p(t_{13}+t_{23},t_{34})
 -p(t_{23},t_{34})\hspace{12mm}\\
 -p(t_{12}+t_{13},t_{24}+t_{34})-p(t_{12},t_{23})=0.
\end{split}
\end{equation}

\begin{lemma}\label{lem:weight-seven-integral-line}
For an element of $D_5$, the parameters in
\eqref{eq:weight-seven-v}--\eqref{eq:weight-seven-w} satisfy
\[
 16b=19a.
\]
Consequently $a\in16\mathbf Z_2$, and the integral leading line is
generated by the primitive vector
\begin{equation}\label{eq:weight-seven-primitive-vector}
\mathbf L=(16,-48,80,80,-32,19,-173,-80,-32,
           17,-99,61,109,48,-65,64,48,-16).
\end{equation}
The leading Hall vector of an element is $(a/16)\mathbf L$.
Moreover, its leading alternating Magnus coefficient is
$[STSTSTS]f=-2a$.
\end{lemma}

\begin{proof}
Evaluate \eqref{eq:weight-seven-leading-pentagon} in the infinitesimal
Artin action on the free Lie algebra on $X_1,\ldots,X_4$. Write
\[
 D_{ij}(X_i)=[X_i,X_j],\qquad
 D_{ij}(X_j)=[X_j,X_i],\qquad D_{ij}(X_k)=0\quad(k\ne i,j).
\]
The relations $[D_{ij},D_{kl}]=0$ for four distinct indices and
$[D_{ij},D_{ik}+D_{jk}]=0$ follow from Jacobi. A common change of sign
in this infinitesimal action does not change the required vanishing
in degree seven. No faithfulness assertion is needed.

Take the coefficient of $X_1^4X_2X_3X_2X_4$ in the image of $X_1$.
For $p=a\mathbf v+b\mathbf w$, the five terms, before attaching
their respective signs $+,+,-,-,-$, give
\[
\begin{array}{c|rr}
\text{term}&\text{coefficient of }a&\text{coefficient of }b\\ \hline
1&46/3&-16/3\\
2&-6&0\\
3&0&0\\
4&3&0\\
5&0&0.
\end{array}
\]
This table follows by the displayed formulas for $D_{ij}$ and the
recursion $[D,E](X_i)=D(E(X_i))-E(D(X_i))$, applied to the standard
Lyndon brackets. The pentagon thus gives $(19a-16b)/3=0$.

Both $a$ and $b$ are integral Hall coordinates, and $19$ is odd.
Hence $a\in16\mathbf Z_2$. Substituting $b=19a/16$ into
\eqref{eq:weight-seven-v}--\eqref{eq:weight-seven-w} gives
\eqref{eq:weight-seven-primitive-vector}; its coordinate $19$ proves
primitivity. Direct bracket expansion on the preceding plane gives
$[STSTSTS]p=(32/3)(a-b)$, which becomes $-2a$.
\end{proof}

\subsection{The exact beta coordinate and an arithmetic value}

For a multiplier-one element, put $a_j=[S^jT]f$ and define
\begin{equation}\label{eq:weight-seven-tau}
 \tau_7(g)=360a_6+900a_5+780a_4-239a_2\in\mathbf Z_2.
\end{equation}

\begin{lemma}\label{lem:weight-seven-beta}
The map $\tau_7$ is a continuous homomorphism on $\ker\lambda$, with
conjugation weight $\lambda_2^7$. On $D_5$ it equals $360m_7$.
\end{lemma}

\begin{proof}
Use the normalized Alexander beta cofactor
\[
 B_g(X,Y)=1+(X-1)(Y-1)R_g(X,Y),\qquad [f]=R_g(X,Y)[x,y].
\]
The action on the rank-one Alexander module has scalar
$S_\lambda(X)S_\lambda(Y)B_g(X,Y)$, where
$S_\lambda(X)=(X^\lambda-1)/(X-1)$. The $S_\lambda$ factors satisfy
the same semilinear multiplicative cocycle identity as that scalar;
cancelling them gives the corresponding identity for $B_g$.
Thus
\[
 \vartheta_7(g)=[u^6v]\log B_g(e^u,e^v)
\]
is an additive $\mathbf Q_2$-valued homomorphism on $\ker\lambda$,
with conjugation weight $\lambda_2^7$.

Since $B_g(X,1)=1$, differentiation at $v=0$ gives
\[
 \left.\partial_v\log B_g(e^u,e^v)\right|_{v=0}
   =(e^u-1)R_g(e^u,1)
   =\sum_{j\ge2}a_j(e^u-1)^j.
\]
The last equality follows from the affine specialization
$x=\operatorname{diag}(1+S,1)$, $y=I+TE_{12}$: the commutator has
translation $ST$, conjugation by $x$ multiplies it by $1+S$, and
conjugation by $y$ acts trivially on that translation module.
The coefficients of $u^6$ in $(e^u-1)^j$, for $j=2,\ldots,6$, are
\[
 \frac{31}{360},\quad\frac34,\quad\frac{13}{6},\quad\frac52,
 \quad1.
\]
Using the earlier identity $a_3=-a_2$ therefore gives
\[
 \vartheta_7=a_6+\frac52a_5+\frac{13}{6}a_4
                         -\frac{239}{360}a_2.
\]
Multiplication by $360$ proves the integral formula
\eqref{eq:weight-seven-tau}, continuity and the exact weight law.
On $D_5$, all the terms except $a_6$ vanish by
Lemma~\ref{lem:weight-seven-leading-spaces}.
\end{proof}

\begin{proposition}\label{prop:weight-seven-arithmetic-value}
There is $z_5\in J_5$ with
\[
 m_7(z_5)\in128\mathbf Z_2^\times.
\]
More explicitly, it can be chosen so that
$m_7(z_5)=128\cdot199017\,\tau_7(\sigma)$ for an arithmetic
$\sigma$ with $\tau_7(\sigma)$ odd.
\end{proposition}

\begin{proof}
Simultaneous rational Kummer surjectivity supplies
$\sigma\in H\cap\ker\lambda$ with $\rho_2^I(\sigma)=2$ in
$\widehat{\mathbf Z}$ and every odd-prime rational Kummer coordinate
zero. The established boundary congruence
$\rho_2^I\equiv-2m_3\pmod4$ makes $a_2(\sigma)$ odd.
Equation~\eqref{eq:weight-seven-tau} consequently makes
$\tau_7(\sigma)$ odd.

Choose $h\in H$ with $\lambda_2(h)=3$ and put
$t=\lambda(h)\in\widehat{\mathbf Z}^{\times}$. For fixed $h$, set
\[
 \mathcal D_c(z)=hzh^{-1}z^{-c},\qquad
 z_1=\mathcal D_t(\sigma),\quad
 z_3=\mathcal D_{t^3}(z_1),\quad
 z_5=\mathcal D_{t^5}(z_3).
\]
The powers are the continuous profinite powers in the procyclic
closure of the element. The first operation kills every rational
Kummer coordinate. The second kills $m_3$ and places $z_3$ in $J_3$.
The third kills $m_5$ by its established weight-five transformation
law on $D_3$. All preceding vanishings are preserved, so $z_5\in J_5$.

The exact homomorphism and weight law of $\tau_7$ give
\[
 \tau_7(z_5)=(3^7-3)(3^7-3^3)(3^7-3^5)\tau_7(\sigma).
\]
Since $\tau_7=360m_7$ on $D_5$, division by $360$ gives
\[
 m_7(z_5)=\frac{2184\cdot2160\cdot1944}{360}\tau_7(\sigma)
         =128\cdot199017\,\tau_7(\sigma).
\]
The last two factors are odd. The compact arithmetic image
therefore contains $128\mathbf Z_2$; the uniform odd theorem
gives the reverse inclusion.
\end{proof}

\begin{proof}[Proof of Theorem~\ref{thm:weight-seven-bounds}]
The lower-degree vanishing and integral leading line were proved in
Lemmas~\ref{lem:weight-seven-leading-spaces} and
\ref{lem:weight-seven-integral-line}. The beta-coordinate identity
gives additivity and the scalar conjugation rule for $m_7$; division
by the nonzero integer $360$ is legitimate in the torsion-free group
$\mathbf Z_2$.

The arithmetic value and closedness show that both images contain
$128\mathbf Z_2$. Since $D_5=C_2$,
Theorem~\ref{thm:uniform-odd-arithmetic-image} with $r=3$ gives
the reverse inclusion. Since the leading vector is
$(m_7/16)\mathbf L$, both Hall images are exactly
$8\mathbf Z_2\mathbf L$.
\end{proof}

\subsection{Finite identities for the leading spaces}
\label{subsec:weight-seven-leading-certificate}

Here are finite recursions and explicit matrix identities that verify
Lemma~\ref{lem:weight-seven-leading-spaces}. They specify the
coefficient calculation without a dependence on an implementation.
Work in noncommutative integer polynomials truncated after degree
nine. Put $X=1+S$, $Y=1+T$, $Z=Y^{-1}X^{-1}$ and
$(1+A)^{-1}=\sum_{j=0}^{9}(-A)^j$. Form the standard Lyndon group
words by the preceding commutator recursion.

For $d=5$ or $7$, let the columns be the Lyndon words of lengths
$d,\ldots,9$, ordered first by length and then lexicographically.
There are $119$ columns for $d=5$ and $104$ for $d=7$. Set
\[
 F_d=1+\sum_{d\le |w|\le9}h_w(L_w-1).
\]
All coefficients through degree nine are linear in the $h_w$,
because products of two nonconstant factors start in degree at
least ten. Define the rows of the integer coefficient matrix $A_d$
by
\[
 I_v=[v]\bigl(F_d(Y,X)+F_d(X,Y)-2\bigr),
\]
\[
 II_v=[v]\bigl(F_d(Z,X)+F_d(Y,Z)+F_d(X,Y)-3\bigr).
\]
The rows vanish on the coordinates of any word satisfying
the two-cycle and three-cycle relations.

Encode a word $v$ of length $n$ by
$c(v)=2^n+\operatorname{bin}(v)$, where $S$ is the binary digit $0$
and $T$ is $1$. The row codes for $I_v$ and $II_v$ are respectively
$c(v)$ and $1024+c(v)$. Let $R_5,R_7$ be the row lists displayed
below, let $C_5=(1,\ldots,96)$ and $C_7=(1,\ldots,82)$, and put
$B_d=A_d[R_d,C_d]$.

Let $\mathbf u$ be the fifteen-entry vector in
\eqref{eq:weight-seven-lower-vector}. In the $119$-dimensional
row space let $T_5$ have the fourteen rows
$2(e_j-u_je_1)$, $2\le j\le15$. In the $104$-dimensional row
space let $T_7$ have the sixteen rows
$3(e_j-v_je_1-w_je_6)$ for $1\le j\le18$, $j\ne1,6$.
Both matrices have integral entries. The finite identities are
\begin{equation}\label{eq:weight-seven-finite-matrix-identity}
 \det B_5=-6603882696,\qquad \det B_7=42661728,
 \qquad
 T_d[:,C_d]B_d^{-1}A_d[R_d,:]=T_d\quad(d=5,7).
\end{equation}
Every entry in these identities is specified by the integer
recursions and the row lists. The inverses are taken over
$\mathbf Q$, and the identities are exact, not congruences.
Applying them to $A_dh=0$ gives $T_dh=0$ and hence the two leading
vector formulas. This is a proof of necessary vanishings and
leading relations; it supplies no lifting assertion for the
remaining coordinates.

% The explicit row lists are appended immediately below.

\[
\begin{split}
R_{5}=\bigl( &33,35,37,65,67,69,77,129,131,133,135,137,\\
 &139,141,147,149,257,259,261,263,265,267,269,275,\\
 &277,281,283,285,293,301,513,515,517,519,521,523,\\
 &525,527,529,531,533,535,537,539,541,547,549,551,\\
 &553,555,557,563,565,569,587,589,595,597,1059,1089,\\
 &1093,1099,1153,1155,1157,1161,1163,1165,1281,1283,1285,1287,\\
 &1289,1291,1299,1301,1317,1325,1537,1539,1541,1545,1547,1549,\\
 &1553,1555,1557,1559,1561,1563,1571,1573,1577,1579,1581,1589\bigr).
\end{split}
\]

\[
\begin{split}
R_{7}=\bigl( &129,131,133,135,137,139,141,147,149,257,259,261,\\
 &263,265,267,269,275,277,281,283,285,293,301,513,\\
 &515,517,519,521,523,525,527,529,531,533,535,537,\\
 &539,541,547,549,551,553,555,557,563,565,569,587,\\
 &589,595,597,1155,1157,1161,1281,1283,1285,1287,1289,1291,\\
 &1299,1301,1317,1325,1537,1539,1541,1545,1547,1549,1553,1555,\\
 &1557,1559,1561,1563,1571,1573,1577,1579,1581,1589\bigr).
\end{split}
\]

\subsection{The exact weight-seven image}
\label{subsec:weight-seven-exact}

The uniform odd theorem supplies the integral image; the leading-space
calculation identifies the corresponding full degree-seven kernel.

\begin{theorem}\label{thm:weight-seven-exact-image}
One has
\[
 m_7(D_5)=m_7(J_5)=128\mathbf Z_2.
\]
Consequently, for $D_7=D_5\cap\ker m_7$ and $J_7=D_7\cap H$,
\[
 D_5=D_7J_5,\qquad G=D_7H,\qquad
 B\simeq D_7/\overline{[G,D_7]J_7}.
\]
The subgroup $D_7$ is proper, closed and normal in $G$, and
\[
 \eta_7:D_5\longrightarrow\mathbf F_2,
 \qquad \eta_7(g)=m_7(g)/128\pmod2
\]
is a continuous $G$-invariant character with nonzero arithmetic
restriction.
\end{theorem}

\begin{proof}
Here $D_5=C_2$. Theorem~\ref{thm:uniform-odd-arithmetic-image}
with $r=3$ gives the two exact images and the factorization
$D_5=D_7J_5$, hence $G=D_7H$. The scalar conjugation law makes
$D_7$ closed and normal and makes $m_7/128$ invariant modulo two.
The explicit arithmetic value in
Proposition~\ref{prop:weight-seven-arithmetic-value} proves properness
and nonvanishing on $J_5$. The canonical relative quotient is the
centralization argument of Corollary~\ref{cor:cubic-radical-normalization}.
Finally, Lemma~\ref{lem:weight-seven-integral-line} shows that killing
$m_7$ kills the entire degree-seven leading term.
\end{proof}

\begin{corollary}\label{cor:weight-seven-nonextension}
The character $\eta_7$ has no $G$-invariant continuous extension to
$D_3$, $N_{\mathrm{rad}}$ or $N$.
Moreover, writing $\psi_q=\rho_q^I/2\pmod2$ for the divided
odd-prime Kummer characters on $N$, one still has
\[
 \ker\bigl[H^1(N,\mathbf F_2)^G\longrightarrow
                         H^1(D_7,\mathbf F_2)\bigr]
 =\left(\bigoplus_{q\text{ odd prime}}\mathbf F_2\psi_q\right)
       \oplus\mathbf F_2\eta_3.
\]
On $D_3$, the corresponding kernel is exactly $\mathbf F_2\eta_5$.
\end{corollary}

\begin{proof}
Use the arithmetic elements constructed above, for which
\[
 z_3\in J_3,\qquad
 z_5=h z_3h^{-1}z_3^{-t^5}\in J_5,
 \qquad \eta_7(z_5)=1.
\]
For every continuous character $\chi$ to $\mathbf F_2$ and every
profinite exponent $c$, continuity of powers gives
$\chi(z^c)=(c\bmod2)\chi(z)$. If $\chi$ is $G$-invariant on
$D_3$, then
\[
 \chi(z_5)=\chi(z_3)-(t_2^5\bmod2)\chi(z_3)=0,
\]
because $t_2$ is odd. Thus its restriction to $D_5$ cannot be
$\eta_7$. Restricting an invariant character from either larger
group to $D_3$ proves the other nonextension assertions.

Now suppose that a $G$-invariant character of $N$ vanishes on
$D_7$. Its restriction to $D_5$ factors through
$D_5/D_7=128\mathbf Z_2$, so is zero or $\eta_7$.
Nonextension excludes the latter, and hence it vanishes on $D_5$.
Its restriction to $D_3$ is similarly zero or $\eta_5$.
The weight-five arithmetic word
$z_3=h z_1h^{-1}z_1^{-t^3}$, with $z_1\in N_{\mathrm{rad}}$,
satisfies $\eta_5(z_3)=1$, whereas the same odd-power calculation
makes every invariant character of $N$ zero on $z_3$.
Consequently the character vanishes on $D_3$.
The established cubic joint quotient
\[
 N/D_3\simeq
 \left(\prod_{q\text{ odd prime}}2\widehat{\mathbf Z}\right)
       \times8\mathbf Z_2
\]
now gives exactly the displayed direct sum. Conversely every one
of those characters already vanishes on $D_3$ and hence on $D_7$.
For a character defined on $D_3$, the first nonextension argument
again forces vanishing on $D_5$; the quotient
$D_3/D_5=32\mathbf Z_2$ then gives precisely $\mathbf F_2\eta_5$.
\end{proof}

If $J_5=G_{L_5}$ and $L_7=\overline{\mathbf Q}^{\,J_7}$, the exact
image also gives $\operatorname{Gal}(L_7/L_5)\simeq\mathbf Z_2$.
This specifies the extension by the weight-seven monodromy
coordinate; no expression by rational radicals is asserted.

\section{An exact weight-eight commutator normalization}
\label{sec:weight-eight-normalization}

Put $M=\ker\lambda$, and write
\[
 \begin{gathered}
 D=N_{\mathrm{rad}},\qquad D_3=D\cap\ker m_3,
 \qquad D_5=D_3\cap\ker m_5,\\
 D_7=D_5\cap\ker m_7,\qquad J_7=D_7\cap H.
\end{gathered}
\]
The preceding normalizations give $G=D_7H$. All Magnus expansions
below are in the free pro-$2$ group, with $x=1+S$ and $y=1+T$.
For $g=(1,f)\in D_7$, define the mixed coefficient
\[
 m_8(g)=[S^4TST^2]f.
\]
The depth-one coefficient $[S^7T]f$ is zero; the first available
degree-eight direction is an Ihara bracket.

\begin{theorem}\label{thm:weight-eight-exact-image}
Every word belonging to $D_7$ starts in degree eight. Its homogeneous
part of degree eight lies in a primitive rank-one integral Lie module,
and
\[
 m_8(D_7)=m_8(J_7)=\mathbf Z_2.
\]
Set $D_8=D_7\cap\ker m_8$ and $J_8=D_8\cap H$. Then $D_8$ is a
proper closed normal subgroup of $G$,
\[
 D_7=D_8J_7,\qquad G=D_8H,\qquad
 B\simeq D_8/\overline{[G,D_8]J_8}.
\]
Every word belonging to $D_8$ starts in degree at least nine.
The continuous character $\eta_8=m_8\bmod2$ is $G$-invariant on
$D_7$ and has nonzero arithmetic restriction. It has no
$G$-invariant extension to $D$ or to $N=\ker(\lambda,\rho_2)$.
\end{theorem}

We give the integral leading calculation and an arithmetic
commutator. No existence claim is drawn from a truncated system of
necessary equations.

\subsection{The integral leading line}

For a Lyndon word $w$ in $S<T$, let $L_w$ be its standard Lie
bracket, using the longest proper Lyndon suffix. Write
\[
 \psi_3=L_{SST}-L_{STT},
\]
and, in the degree-five Lyndon order, put
\[
 \psi_5=L_{SSSST}-2L_{SSSTT}
       +\tfrac32 L_{SSTST}+2L_{SSTTT}
       +\tfrac12 L_{STSTT}-L_{STTTT}.
\]
For a Lie element $P$, define the derivation
$\partial_P(S)=0$, $\partial_P(T)=[T,P]$, and use the Ihara bracket
\[
 \{P,Q\}=[P,Q]+\partial_P(Q)-\partial_Q(P).
\]
Let
\[
 \mathcal L_8=2\{\psi_3,\psi_5\}.
\]
Its integral Lyndon coordinates are as follows:
\[
\begin{array}{cr@{\qquad}cr@{\qquad}cr}
 w&[L_w]\mathcal L_8&w&[L_w]\mathcal L_8&w&[L_w]\mathcal L_8\\ \hline
 SSSSSSST&0&SSSTSTTT&-6&SSTTSTST&21\\
 SSSSSSTT&0&SSSTTSST&4&SSTTSTTT&3\\
 SSSSSTST&-4&SSSTTSTT&-3&SSTTTSTT&-15\\
 SSSSSTTT&0&SSSTTTST&-6&SSTTTTST&-3\\
 SSSSTSST&-2&SSSTTTTT&0&SSTTTTTT&0\\
 SSSSTSTT&7&SSTSSTST&-5&STSTSTTT&8\\
 SSSSTTST&17&SSTSSTTT&3&STSTTSTT&-11\\
 SSSSTTTT&0&SSTSTSTT&3&STSTTTTT&-4\\
 SSSTSSTT&-1&SSTSTTST&6&STTSTTTT&-10\\
 SSSTSTST&-12&SSTSTTTT&7&STTTTTTT&0
\end{array}
\]
In particular $\mathcal L_8$ is primitive. Direct associative
expansion gives
\begin{equation}\label{eq:weight-eight-mixed-coefficient}
 [S^4TST^2]\mathcal L_8=7,
 \qquad [S^7T]\mathcal L_8=0.
\end{equation}
The bracket formula and the displayed finite table specify this
identity without a choice of a group lift of either rational Lie
normalization.

\begin{lemma}\label{lem:weight-eight-leading-line}
Suppose $f\in\gamma_8$ satisfies the two-cycle and three-cycle
relations
\[
 f(y,x)f(x,y)=1,\qquad
 f(z,x)f(y,z)f(x,y)=1,\qquad z=(xy)^{-1}.
\]
Then its homogeneous part of degree eight is
\[
 \operatorname{in}_8(f)=b\mathcal L_8\quad\text{for a unique }b\in\mathbf Z_2.
\]
If $h_w$ denotes a degree-eight Lyndon coordinate, then
$b=-h_{SSSTSSTT}$ and $m_8(f)=7b$.
\end{lemma}

\begin{proof}
The following degree-nine certificate proves the necessary rational
line. Its coordinate $-1$ then proves the stated integral assertion.
Collect $f$ modulo $\gamma_{10}$ in the standard Lyndon group
commutators $C_w$, with integral coordinates $h_w$. Explicitly, for the
standard factorization $w=uv$, put
\[
 C_S=x,\qquad C_T=y,\qquad C_w=C_uC_vC_u^{-1}C_v^{-1}.
\]
Set
\[
 F=1+\sum_{8\le |w|\le9}h_w(C_w-1).
\]
There are $30+56=86$ columns, in order of weight and then lexicographic
order. The expression for $F$ is exact through degree nine, since
products of two nonconstant factors have degree at least sixteen.
All the following polynomial operations are truncated after degree
nine. Put $X=1+S$, $Y=1+T$, $Z=Y^{-1}X^{-1}$, and form
\[
 I_v=[v]\bigl(F(Y,X)+F(X,Y)-2\bigr),\qquad
 II_v=[v]\bigl(F(Z,X)+F(Y,Z)+F(X,Y)-3\bigr).
\]
These are integral linear forms on the $86$ Hall coordinates and
vanish for the given word.

Encode $I_v$ by $2^{|v|}+\operatorname{bin}(v)$ and $II_v$ by
$1024+2^{|v|}+\operatorname{bin}(v)$, with $S=0$, $T=1$. Let $A_R$
be the $66$ by $86$ matrix consisting of the following rows, in the
written order:
\[
\begin{aligned}
R=( &257,259,261,263,265,267,269,275,277,281,283,285,293,301,\\
    &513,515,517,519,521,523,525,527,529,531,533,535,537,539,\\
    &541,547,549,551,553,555,557,563,565,569,587,589,595,597,\\
    &1281,1287,1291,1299,1301,1317,1537,1539,1541,1545,\\
    &1547,1549,1553,1555,1557,1559,1561,1563,1571,1573,\\
    &1577,1579,1581,1589).
\end{aligned}
\]
Let $C$ be its first $66$ columns and let $A_{R,C}$ denote the
corresponding square matrix. If $\ell_j$ is the $j$th entry of the
thirty-coordinate table above, ordered lexicographically by its word,
define $T$ to have the $29$ rows
\[
 T_j=e_j+\ell_j e_9\quad(1\le j\le30,\ j\ne9),
\]
where the $e_j$ are rows of length $86$. The finite identities are
\begin{equation}\label{eq:weight-eight-leading-certificate}
 \det A_{R,C}=-12312=-2^3\,3^4\,19,
 \qquad
 T_C A_{R,C}^{-1}A_R=T.
\end{equation}
Here every matrix entry is determined by the displayed group
commutator recursion and the finite geometric expansion
$(1+A)^{-1}=\sum_{j=0}^{9}(-A)^j$. Thus
\eqref{eq:weight-eight-leading-certificate} is an explicit finite
rational matrix identity, not a rank assumption. Multiplying
$A_Rh=0$ by $T_C A_{R,C}^{-1}$ gives
$h_j=-\ell_jh_9$ for all $j\le30$. This proves the claimed line
over $\mathbf Q_2$. Since $h_9\in\mathbf Z_2$ and $\ell_9=-1$,
its scalar $b=-h_9$ is integral. The mixed coefficient is then
$7b$ by \eqref{eq:weight-eight-mixed-coefficient}.
\end{proof}

The earlier leading degree-seven calculation shows that a word in
$D_5$ has its degree-seven part determined by $m_7$. Therefore every
word in $D_7$ starts in degree eight and satisfies the preceding
lemma. In particular its degree-eight part vanishes exactly when
$m_8$ vanishes. This upper-lattice argument uses only the two-cycle
and three-cycle identities once the preceding normalization has
been made; it uses no new pentagon coefficient.

\subsection{A primitive arithmetic commutator}

Recall the three homomorphisms on all of $M$, together with their
scalar conjugation laws:
\[
 \begin{array}{c|c|c}
  \text{homomorphism}&\text{formula}&\text{conjugation weight}\\ \hline
  m_3 &[S^2T]f&\lambda_2^3\\
  \tau_5&12[S^4T]f-11m_3&\lambda_2^5\\
  \tau_7&360[S^6T]f+900[S^5T]f+780[S^4T]f-239m_3
            &\lambda_2^7.
 \end{array}
\]
They are the coefficients of the logarithmic Alexander beta
cocycle: $m_3=c_3$, $\tau_5=12c_5$, and $\tau_7=360c_7$, where
$c_j=[u^{j-1}v]\log B(e^u,e^v)$. Thus these formulas apply before
any radical normalization. On $D_3$ one has $\tau_5=12m_5$, and
on $D_5$ one has $\tau_7=360m_7$.

\begin{lemma}\label{lem:weight-eight-arithmetic-commutator}
One has $\overline{[M,M]}\subseteq D_7$. There are
$\sigma\in H\cap M$ and $h\in H$ such that, with
$s=h\sigma h^{-1}$,
\[
 w_8=[\sigma,s]\in J_7,
 \qquad m_8(w_8)=63\,m_3(\sigma)\,\tau_5(\sigma)
                 \in\mathbf Z_2^\times.
\]
Moreover $w_8=[\sigma,k]$ for some $k\in D\cap H$.
\end{lemma}

\begin{proof}
Every rational-prime Kummer cocycle is a homomorphism on $M$.
Consequently $[M,M]\subseteq D$. The homomorphism $m_3$ makes
these commutators belong to $D_3$; the homomorphism $\tau_5$, which
there equals $12m_5$, makes them belong to $D_5$; and $\tau_7=360m_7$
on $D_5$ makes them belong to $D_7$. All the kernels are closed,
which proves the first assertion.

The exact arithmetic Kummer image supplies
$\sigma\in G_{\mathbf Q^{\mathrm{ab}}}\subset H\cap M$ with
\[
 \rho_2(\sigma)=2\in\widehat{\mathbf Z},\qquad
 \rho_q(\sigma)=0\quad(q\text{ odd prime}).
\]
Write $a=m_3(\sigma)$. The established congruence
$\rho_2\equiv-2m_3\pmod4$ makes $a$ odd, and the formula for
$\tau_5$ makes $\tau_5(\sigma)$ odd as well.
Choose $h\in H$ with $\lambda_2(h)=3$, and put $t=\lambda(h)$.
Profinite powers are defined by the continuous map
$\widehat{\mathbf Z}\to\overline{\langle\sigma\rangle}$ sending
$1$ to $\sigma$. Hence both
\[
 k=s\sigma^{-t}\quad\text{and}\quad r=s\sigma^{-27}
\]
are elements of $H\cap M$. Conjugation has weight $\lambda$
on every Kummer homomorphism, so $k\in D\cap H$. Since powers of
$\sigma$ commute with $\sigma$,
\[
 [\sigma,k]=[\sigma,s]=[\sigma,r]=w_8.
\]
The weight-three law gives $m_3(r)=(3^3-27)a=0$. The low-degree
cycle calculation, whose degree-three and degree-four Hall
coordinates are $(m_3,-m_3)$ and $(-m_3,m_3,-m_3)$, therefore shows
that $r$ starts in degree five. Its leading part is a multiple of
$\psi_5$, by the preceding integral degree-five calculation. The
weight-five law gives its multiplier explicitly:
\[
 [S^4T]f_r=\frac{3^5-27}{12}\tau_5(\sigma)
       =18\tau_5(\sigma).
\]
The leading part of $\sigma$ is $a\psi_3$.

For completeness, with the normalized action
$\Phi_f(x)=x$, $\Phi_f(y)=f^{-1}yf$, multiplication of parameters
is $f\star g=f\Phi_f(g)$. Its leading Lie bracket is
$[P,Q]+\partial_P(Q)-\partial_Q(P)$, since the first-order change
in $T$ is $[T,P]$. Thus the leading part of this group
commutator is
\[
 \operatorname{in}_8(w_8)
   =18a\tau_5(\sigma)\{\psi_3,\psi_5\}
   =9a\tau_5(\sigma)\mathcal L_8.
\]
Taking the coefficient in
\eqref{eq:weight-eight-mixed-coefficient} gives the claimed odd
value $63a\tau_5(\sigma)$.
\end{proof}

\begin{proof}[Proof of Theorem~\ref{thm:weight-eight-exact-image}]
The leading coefficient $m_8$ is an additive continuous map on
$D_7$. This follows either by multiplying the filtered normalized
automorphisms or by noting that substitutions affect a leading
word only above degree eight. Conjugation acts on it by
\[
 m_8(gng^{-1})=\lambda_2(g)^8m_8(n)
 \quad(g\in G,\ n\in D_7).
\]
Indeed the associated graded action on the two free generators is
scalar $\lambda_2$, and conjugation of a tangential derivation of
degree eight has weight $\lambda_2^8$.

The image of $J_7$ is a compact additive subgroup of $\mathbf Z_2$
containing the unit given by the arithmetic commutator. It is
therefore all of $\mathbf Z_2$; the same holds for $D_7$.
The scalar law proves that $D_8$ is normal and that $\eta_8$ is
invariant. The unit proves properness and arithmetic nonvanishing.
Equality of the arithmetic and full images gives $D_7=D_8J_7$, hence
$G=D_8H$. The leading-line lemma shows that words in $D_8$ start
in degree at least nine.

The relative abelian quotient is unchanged by this replacement.
Explicitly, a continuous $G$-invariant homomorphism from $D_8$ to
a profinite abelian group, trivial on $J_8$, extends uniquely to
$G=D_8H$ by sending $dh$ to the value at $d$. The intersection
condition proves well-definedness, and the compact surjection
$D_8\times H\to G$ proves continuity. This is exactly the universal
property of $D_8/\overline{[G,D_8]J_8}$ and of the preceding quotient
$B$.

Finally let $\chi:D\to\mathbf F_2$ be $G$-invariant. The arithmetic
word satisfies $w_8=[\sigma,k]$ with $k\in D$, so
\[
 \chi(w_8)=\chi(\sigma k\sigma^{-1})-\chi(k)=0,
 \qquad \eta_8(w_8)=1.
\]
Thus $\chi$ cannot extend $\eta_8$. An extension to $N$ would
restrict to one on $D$, so it is impossible as well.
\end{proof}

\begin{corollary}\label{cor:weight-eight-original-character-kernel}
Let $E$ be any closed $G$-stable subgroup containing $D$ and contained
in $M$, and let $A$ be a finite abelian group with trivial $G$-action.
Then
\[
 \{\chi\in\operatorname{Hom}_{\mathrm{cts}}(E,A)^G:
                         \chi(D_8)=0\}
 =
 \{\chi\in\operatorname{Hom}_{\mathrm{cts}}(E,A)^G:
                         \chi(D_7)=0\}.
\]
In particular,
\[
 \ker\!\left(H^1(N,\mathbf F_2)^G
                 \longrightarrow H^1(D_8,\mathbf F_2)\right)
 =\left(\bigoplus_{q\text{ odd prime}}\mathbf F_2\psi_q\right)
       \oplus\mathbf F_2\eta_3.
\]
\end{corollary}

\begin{proof}
Every invariant homomorphism on $E$ kills $w_8=[\sigma,k]$, since
$k\in D\subseteq E$. If it also kills $D_8$, its restriction to
$D_7$ factors through $D_7/D_8\simeq\mathbf Z_2$. The image of
$w_8$ topologically generates this quotient, so the restriction is
zero. The reverse inclusion is immediate. Apply this assertion
to $E=N$, $A=\mathbf F_2$, and use the previously proved kernel
formula for restriction to $D_7$.
\end{proof}

\subsection{The extension threshold of the eighth character}
\label{subsec:eta-eight-extension-threshold}

Complex conjugation gives an arithmetic obstruction to extending the
new character through the preceding normalizations.

\begin{proposition}\label{prop:eta-eight-no-extension-d3}
The character $\eta_8$ on $D_7$ extends to the invariant character
\[
 \alpha_8:D_5\longrightarrow\mathbf F_2,\qquad
 \alpha_8(1,f)=[S^4TST^2]f\pmod2,
\]
but has no $G$-invariant extension to $D_3$. More precisely, no
nonzero linear combination of $\eta_7$ and $\alpha_8$ on $D_5$
extends to a $G$-invariant character of $D_3$.
\end{proposition}


\begin{proof}
The exact weight-seven theorem puts the degree-seven part of every
$D_5$ word in $8\mathbf Z_2\mathcal L_7$, where $\mathcal L_7$
is the primitive vector of \eqref{eq:weight-seven-primitive-vector}.
Thus its Magnus expansion over $\mathbf F_2$ starts in degree eight.
The first nonzero term is additive under the normalized GT product,
and conjugation acts by $\lambda_2^8$, which is one modulo two.
This follows by substituting the linear terms in the GT automorphism:
higher terms raise degree, and a homogeneous ambiguity commuting with
$T$ would be a pure power of $T$, absent from a commutator word.
Hence $\alpha_8$ is an invariant character and restricts to $\eta_8$.

On $M=\ker\lambda$ use the additive coordinates
\[
 a=m_3,\qquad T_5=12[S^4T]f-11a,\qquad
 T_7=360[S^6T]f-1020[S^4T]f+661a.
\]
Their conjugation weights are $\lambda_2^3,\lambda_2^5,\lambda_2^7$.
Together with the rational Kummer homomorphisms they imply
\[
 [M,M]\subset D_7,\qquad [M,D_7]\subset D_8.
\]
For the first inclusion all four families vanish on commutators;
then $T_5=12m_5$ and $T_7=360m_7$, and $\mathbf Z_2$ has no torsion.
For the second inclusion use the weight-eight law on $D_7$ and
$\lambda(M)=1$. Therefore $M/D_8$ has nilpotency class at most two,
and its commutator subgroup is contained in $D_7/D_8\simeq\mathbf Z_2$.
The previously proved commutator formula is
\begin{equation}\label{eq:eta-eight-commutator-pairing}
 m_8([z,z'])=\frac7{24}\bigl(a(z)T_5(z')-a(z')T_5(z)\bigr)
 \qquad(z,z'\in M).
\end{equation}

Let $c\in H$ be complex conjugation, and choose $h\in H$ with
$\lambda_2(h)=3$. Put $t=\lambda(h)$ and $\sigma=[h,c]$.
Since $c^2=1$, $\lambda(c)=-1$ and $\rho_2(c)=0$, one has
\[
 c\sigma c^{-1}=\sigma^{-1},\qquad
 \rho_2(\sigma)=2\rho_2(h).
\]
The established parity law makes $\rho_2(h)$ odd at two, and
$\rho_2(\sigma)\equiv-2m_3(\sigma)\pmod4$ makes
\[
 A=a(\sigma),\qquad B=T_5(\sigma),\qquad C=T_7(\sigma)
\]
all odd. The last two assertions also follow directly from their
formulas above. The odd-prime Kummer coordinates of $\sigma$ need
not vanish.

For $e\in\widehat{\mathbf Z}$ put
$D_e(z)=hzh^{-1}z^{-e}$, and define
\[
 z_1=D_t(\sigma),\qquad z_3=D_{t^3}(z_1),\qquad
 z_5=D_{t^5}(z_3).
\]
The weight laws give $z_1\in D\cap H$, $z_3\in D_3\cap H$ and
$z_5\in D_5\cap H$. Moreover
\[
 m_7(z_5)=\frac{(3^7-3)(3^7-3^3)(3^7-3^5)}{360}C
        =128\cdot199017\,C,
\]
so $\eta_7(z_5)=1$. Every invariant quadratic character on $D_3$
kills $z_5$, since conjugation fixes its value on $z_3$ and
$t^5$ is odd at two.

A second arithmetic element is $v=cz_3c^{-1}z_3$.
Every rational Kummer coordinate and each of $a,T_5,T_7$ changes
sign under $c$, so $v\in D_7$. Every invariant character on $D_3$
also kills $v$. We compute $m_8(v)$ in the class-two quotient.
Put $X=\sigma$, $Y=h\sigma h^{-1}$, $Z=h^2\sigma h^{-2}$.
The identity $chc^{-1}=X^{-1}h$ gives
\[
 cXc^{-1}=X^{-1},\quad
 cYc^{-1}=X^{-1}Y^{-1}X,\quad
 cZc^{-1}=X^{-1}Y^{-1}Z^{-1}YX.
\]
For profinite exponents $p,r$, set $W=ZY^{-p}(YX^{-p})^{-r}$.
Class-two collection gives
\begin{equation}\label{eq:eta-eight-conjugation-norm}
 (cWc^{-1})W=[X,Y]^{pr(p-1)-p-r}
                  [X,Z]^{1-pr}[Y,Z]^{1+p+r}.
\end{equation}
Here is a direct verification. Encode the normal form
$X^aY^bZ^c[X,Y]^d[X,Z]^e[Y,Z]^f$ by $(a,b,c,d,e,f)$.
Multiplication is
\[
 (a,b,c,d,e,f)(A,B,C,D,E,F)
 =(a+A,b+B,c+C,d+D-Ab,e+E-Ac,f+F-Bc),
\]
and its $n$th power has coordinates
\[
 (na,nb,nc,nd-\tbinom n2ab,
             ne-\tbinom n2ac,nf-\tbinom n2bc).
\]
The normal form of $W$ is
\[
 (pr,-p-r,1,\ p^2r+pr(r-1)/2,-pr,p+r).
\]
Conjugation by $c$ sends a coordinate tuple to
$(-a,-b,-c,d+b,e+c,f+c)$. Multiplying these two tuples gives
\eqref{eq:eta-eight-conjugation-norm}. Integer exponent identities
extend to profinite exponents by continuity, including the integral
binomial polynomial $n(n-1)/2$.

For $W=z_3$, the exponents are $p=t,r=t^3$. Only after applying
$m_8$ to the central commutator group do their two-adic values become
$3,27$. Thus the exponents in \eqref{eq:eta-eight-conjugation-norm}
are $132,-80,31$ in that coordinate. Formula
\eqref{eq:eta-eight-commutator-pairing} gives
\[
 m_8([X,Y])=63AB,\quad m_8([X,Z])=17010AB,\quad
 m_8([Y,Z])=413343AB.
\]
Consequently
\[
 m_8(v)=(132\cdot63-80\cdot17010+31\cdot413343)AB
       =11461149\,AB
\]
is odd. Hence $(\eta_7(v),\alpha_8(v))=(0,1)$.

Suppose an invariant character on $D_3$ restricts to
$\mu\eta_7+\nu\alpha_8$ on $D_5$. Evaluation at $v$ forces
$\nu=0$; evaluation at $z_5$ then forces $\mu=0$.
The value $\alpha_8(z_5)$ is unnecessary. Finally, an extension
of $\eta_8$ to $D_3$ would restrict on $D_5$ to
$\alpha_8+\epsilon\eta_7$, because $D_5/D_7\simeq128\mathbf Z_2$.
Both possibilities have been excluded.
\end{proof}

\begin{corollary}\label{cor:d3-character-kernel-d8}
One has
\[
 \ker\!\left(H^1(D_3,\mathbf F_2)^G
                    \longrightarrow H^1(D_8,\mathbf F_2)\right)
       =\mathbf F_2\eta_5.
\]
\end{corollary}
\begin{proof}
Let $\chi$ be an invariant character on $D_3$ that vanishes on $D_8$.
Since $D_7/D_8\simeq\mathbf Z_2$ through $m_8$, its restriction
to $D_7$ is zero or $\eta_8$. Subtracting the corresponding multiple
of $\alpha_8$ from $\chi|_{D_5}$ makes it vanish on $D_7$.
It then factors through $D_5/D_7\simeq128\mathbf Z_2$, so
$\chi|_{D_5}=\mu\eta_7+\nu\alpha_8$.
The proposition forces both coefficients to be zero.
Thus $\chi$ factors through $D_3/D_5\simeq32\mathbf Z_2$ and is
a multiple of $\eta_5$. Conversely $\eta_5$ is invariant on $D_3$
and vanishes on $D_5$, hence on $D_8$.
\end{proof}


\section{The exact weight-nine arithmetic normalization}
\label{sec:weight-nine-bounds}

The next odd coefficient gives another exact arithmetic quotient.
Use $G=D_8H$, $J_7=D_7\cap H$ and $J_8=D_8\cap H$.
For $g=(1,f)\in D_7$, put
\[
 m_9(g)=[S^8T]f.
\]
We prove that its image on each of $D_7,D_8,J_7,J_8$ is
$512\mathbf Z_2$. Killing this coefficient on $D_8$ therefore
preserves the relative abelian obstruction and leaves only words
starting in degree ten. The uniform odd theorem supplies the integral bound. We correct an
arithmetic value by the degree-eight commutator and then determine the
whole leading line.

\subsection{An additive coefficient and an arithmetic value}

For $g\in M=\ker\lambda$, write $a_j=[S^jT]f$. The logarithmic
Alexander beta coefficient
\[
 c_9(g)=[u^8v]\log B_g(e^u,e^v)
\]
is a continuous homomorphism on $M$ and has conjugation weight
$\lambda_2^9$. The boundary formula for the beta cocycle gives
\[
 c_9=[u^8]\sum_{j\ge1}a_j(e^u-1)^j.
\]
Using $a_1=0$, $a_3=-a_2$, and the finite exponential expansion gives
\begin{equation}\label{eq:weight-nine-beta-formula}
 \begin{split}
 \tau_9:=20160c_9={}&20160a_8+70560a_7+95760a_6\\
                  &+63000a_5+20412a_4-2771a_2.
 \end{split}
\end{equation}
This is an integral homomorphism on all of $M$. Explicitly, the
coefficients of $a_2,a_4,a_5,a_6,a_7,a_8$ in $c_9$ are
\[
 -\frac{2771}{20160},\quad\frac{81}{80},\quad
 \frac{25}{8},\quad\frac{19}{4},\quad\frac72,\quad1.
\]
These numbers can also be checked from
$[u^8](e^u-1)^j=j!\genfrac{\{}{\}}{0pt}{}{8}{j}/8!$.
On $D_7$ the lower depth-one terms vanish: the possible degree-eight
bracket has no depth-one component. Consequently $c_9=m_9$ already
on $D_7$, and in particular on $D_8$.

\begin{proposition}\label{prop:weight-nine-arithmetic-value}
There is $w_9\in J_8$ such that
\[
 m_9(w_9)=512\cdot12746612385\,\tau_9(\sigma),
 \qquad \tau_9(\sigma)\in\mathbf Z_2^\times,
\]
where $\sigma$ is the arithmetic element used in the preceding
normalizations. Thus $512\mathbf Z_2\subseteq m_9(J_8)$.
\end{proposition}

\begin{proof}
Choose $\sigma\in H\cap M$ with $\rho_2(\sigma)=2$ and all the
odd-prime Kummer coordinates zero, and choose $h\in H$ with
$\lambda_2(h)=3$. Put $t=\lambda(h)$ and
\[
 \Delta_j(z)=hzh^{-1}z^{-t^j},\qquad
 z_1=\Delta_1(\sigma),\quad z_3=\Delta_3(z_1),\quad
 z_5=\Delta_5(z_3),\quad z_7=\Delta_7(z_5).
\]
All powers are the continuous profinite powers. The previously
proved weight-one, three, five and seven conjugation laws show
successively that $z_1\in D\cap H$, $z_3\in J_3$, $z_5\in J_5$,
and $z_7\in J_7$. By the weight-nine law,
\[
 c_9(z_7)=\frac{\prod_{j\in\{1,3,5,7\}}(3^9-3^j)}{20160}
                          \tau_9(\sigma)
        =512\cdot12746612385\,\tau_9(\sigma).
\]
The integer multiplying $512$ is odd; the full numerator divided
by $20160$ is $2^9 3^{14}\cdot5\cdot13\cdot41$.
Also $m_3(\sigma)$ is odd, and all other coefficients in
\eqref{eq:weight-nine-beta-formula} are even. Thus
$\tau_9(\sigma)$ is odd.

Let $w_8\in J_7$ be the primitive arithmetic commutator of the
preceding section. Its $m_8$ value is a unit. Choose
$u\in\widehat{\mathbf Z}$ whose $2$-adic component is
$m_8(z_7)/m_8(w_8)$, and set
\[
 w_9=z_7w_8^{-u}.
\]
Then $w_9\in J_8$. Since $w_8$ is a commutator of two elements of
$M$ and $c_9$ is a homomorphism on $M$, one has $c_9(w_8)=0$.
Therefore $m_9(w_9)=c_9(w_9)=c_9(z_7)$, as claimed. Compactness
of the arithmetic image supplies the entire ideal
$512\mathbf Z_2$.
\end{proof}

\subsection{The full leading line}

Use the standard Lyndon Lie and group commutators, with $S<T$ and
the longest proper Lyndon suffix. The degree-nine leading Lie
module of multiplier-one GT words that start in degree nine is contained
in $\mathbf Z_2\mathcal L_9$, where the primitive vector
$\mathcal L_9$ has the following coordinates:

{\small
\[
\begin{array}{cr@{\quad}cr@{\quad}cr}
w&[L_w]\mathcal L_9&w&[L_w]\mathcal L_9&w&[L_w]\mathcal L_9\\ \hline
SSSSSSSST&144&SSSTSTSST&-1118&SSTSTTTTT&-3686\\
SSSSSSSTT&-576&SSSTSTSTT&-5303&SSTTSSTTT&-1013\\
SSSSSSTST&1512&SSSTSTTST&-2951&SSTTSTSTT&2059\\
SSSSSSTTT&1344&SSSTSTTTT&4353&SSTTSTTST&5340\\
SSSSSTSST&-1440&SSSTTSSTT&-682&SSTTSTTTT&-3340\\
SSSSSTSTT&-346&SSSTTSTST&-11276&SSTTTSTST&4388\\
SSSSSTTST&-6322&SSSTTSTTT&479&SSTTTSTTT&1868\\
SSSSSTTTT&-2016&SSSTTTSST&-1762&SSTTTTSTT&3558\\
SSSSTSSST&360&SSSTTTSTT&-8109&SSTTTTTST&2434\\
SSSSTSSTT&-1598&SSSTTTTST&-6663&SSTTTTTTT&576\\
SSSSTSTST&3214&SSSTTTTTT&-1344&STSTSTSTT&3299\\
SSSSTSTTT&-2337&SSTSSTSTT&-89&STSTSTTTT&-5446\\
SSSSTTSST&2722&SSTSSTTST&-5496&STSTTSTTT&-1768\\
SSSSTTSTT&4724&SSTSSTTTT&-1949&STSTTTSTT&6452\\
SSSSTTTST&8823&SSTSTSSTT&291&STSTTTTTT&1368\\
SSSSTTTTT&2016&SSTSTSTST&1041&STTSTTTTT&2088\\
SSSTSSSTT&600&SSTSTSTTT&7472&STTTSTTTT&936\\
SSSTSSTST&2148&SSTSTTSTT&-4352&STTTTTTTT&-144\\
SSSTSSTTT&3259&SSTSTTTST&-2416&&\\
\end{array}
\]
}

Its first coordinate is $144$, and its nineteenth coordinate is
$3259$, an odd integer. Hence, if the leading vector is
$A\mathcal L_9$, then $A\in\mathbf Z_2$ and $m_9=144A$.

Here is complete finite data for the necessary leading line.
Collect group commutators of weights nine and ten, giving $155$
columns. Form $I_v,II_v$ as in the preceding section, now through
degree ten. Encode them by $2^{|v|}+\operatorname{bin}(v)$ and
$2048+2^{|v|}+\operatorname{bin}(v)$, respectively.
For the pentagon rows, let $d_{ij}$ be the derivation on the four
free letters $X_1,\ldots,X_4$ defined by
\[
 d_{ij}(X_i)=[X_i,X_j],\qquad
 d_{ij}(X_j)=[X_j,X_i],\qquad d_{ij}(X_k)=0\ (k\ne i,j).
\]
Use the five pairs and signs
\[
 (d_{12},d_{23}+d_{24},+),\quad
 (d_{13}+d_{23},d_{34},+),\quad
 (d_{23},d_{34},-),\quad
 (d_{12}+d_{13},d_{24}+d_{34},-),\quad
 (d_{12},d_{23},-).
\]
For a homogeneous Lie polynomial $P$, let $\mathcal P(P)$ be the
signed sum of its five evaluations on these pairs. This is the
leading Artin image of the pentagon, up to the common sign $(-1)^9$;
therefore it vanishes for every GT word under consideration.
Let $P_j$ be the coefficient on $X_1$ at the following five words,
respectively:
\[
 1111112324,\quad1111123324,\quad1111124324,\quad
 1111232324,\quad1111232424,
\]
where a digit $i$ denotes $X_i$. Give $P_j$ the row code $4096+j$.
Their entries on the degree-ten Hall columns are zero. The
coefficients on the degree-nine columns are determined by
$L_w(\mathcal U,\mathcal V)(X_1)$, with Lie bracket interpreted as
the commutator of derivations.

Let $A_R$ consist of the following $121$ rows, in the written order:
\[
\begin{aligned}
R=(&513,515,517,519,521,523,525,527,529,531,533,535,537,539,\\
&541,547,549,551,553,555,557,563,565,569,587,589,595,597,\\
&1025,1027,1029,1031,1033,1035,1037,1039,1041,1043,1045,1047,1049,1051,\\
&1053,1059,1061,1063,1065,1067,1069,1073,1075,1077,1079,1081,1083,1085,\\
&1093,1095,1097,1099,1101,1107,1109,1113,1115,1117,1125,1129,1131,1133,\\
&1139,1171,1173,1197,1203,1237,2563,2565,2569,2577,2583,2587,2603,2605,\\
&2613,3073,3075,3077,3079,3081,3083,3085,3089,3091,3093,3095,3097,3099,\\
&3101,3107,3109,3113,3115,3117,3119,3121,3123,3141,3145,3147,3149,3157,\\
&3159,3161,3163,3221,3243,3253,4097,4098,4101).
\end{aligned}
\]
Let $C$ be its first $121$ columns, and write
$q_j=\ell_j/144$, where $\ell_j$ is the $j$th coordinate of
$\mathcal L_9$ in lexicographic order. Define $T$ by its $55$ rows
$e_j-q_je_1$ for $2\le j\le56$. The exact finite identities are
\[
 \det A_{R,C}=22883210311680,\qquad
 T_C A_{R,C}^{-1}A_R=T.
\]
All entries are specified by the finite group-commutator and
geometric-series recursions through degree ten, and the displayed
derivations. Thus the identities give $h_j=q_jh_1$ for every
leading vector of a GT word. In particular it is a rational multiple
of $\mathcal L_9$, and its nineteenth coordinate makes the multiple
integral. This proves the asserted upper lattice without claiming
that arbitrary solutions of the necessary equations lift.

\subsection{The exact weight-nine normalization}
\label{subsec:weight-nine-exact}

The integral boundary measure forces the ninth coefficient into
its arithmetic lattice.

\begin{theorem}\label{thm:weight-nine-exact-image}
One has
\[
 m_9(D_8)=m_9(J_8)=512\mathbf Z_2.
\]
The same equality holds with $D_8,J_8$ replaced by $D_7,J_7$.
For $D_9=D_8\cap\ker m_9$ and $J_9=D_9\cap H$,
\[
 D_8=D_9J_8,\qquad G=D_9H,\qquad
 B\simeq D_9/\overline{[G,D_9]J_9}.
\]
The subgroup $D_9$ is proper, closed and normal in $G$, and every
word belonging to it starts in degree at least ten. Moreover,
\[
 (m_8,m_9):D_7\longrightarrow\mathbf Z_2\times512\mathbf Z_2
\]
is surjective already on $J_7$ and has kernel $D_9$.
\end{theorem}

\begin{proof}
The arithmetic word in Proposition~\ref{prop:weight-nine-arithmetic-value}
gives a value of exact valuation nine on $J_8$.
For the upper bound, use Lemma~\ref{lem:uniform-odd-coordinate-bound}
with $r=4$. On $D_7$ the coefficients $a_2,a_4,a_6$ vanish, so
$m_9=a_8\in512\mathbf Z_2$. Concretely, the polynomial
\[
 (t^2-1)(t^2-9)(t^2-25)(t^2-49)
 =t^8-84t^6+1974t^4-12916t^2+11025
\]
is divisible by $2^{15}$ on odd $t$. Its lower moments vanish,
and its eighth moment is $(1-2^8)\tau_9$, with
$\tau_9=20160m_9$ and $v_2(20160)=6$. This is exactly the required
valuation-nine bound. Compactness and the arithmetic witness
therefore prove equality of the two images.

The coefficient $m_9=c_9$ is already an additive, weight-nine
homomorphism on $D_7$. Given any element of $D_7$, multiply it by
a suitable profinite power of the primitive arithmetic
weight-eight commutator $w_8$ to make $m_8$ zero. This does not
change $c_9$, because every logarithmic beta homomorphism vanishes
on $w_8$. Consequently the same exact image equality holds on
$D_7$ and $J_7$. The two arithmetic words $w_8,w_9$ have respective
coordinate values
\[
 (m_8(w_8),0),\qquad(0,m_9(w_9)),
\]
where the first is a unit and the second has exact valuation
nine. They prove the asserted joint surjectivity.

The scalar conjugation laws imply normality of $D_9$. The
arithmetic value proves properness, and equality of the two images
gives $D_8=D_9J_8$, hence $G=D_9H$. The same compact gluing and
centralization argument as in the preceding normalizations gives
the displayed description of $B$. The leading-line calculation
shows that $m_9=0$ kills the whole degree-nine term, so the words
of $D_9$ start in degree at least ten.
\end{proof}

\subsection{The ninth character and its extension obstruction}

\begin{proposition}\label{prop:eta-nine-nonextension}
The character
\[
 \eta_9:D_7\longrightarrow\mathbf F_2,\qquad
 \eta_9(g)=m_9(g)/512\pmod2
\]
is continuous and $G$-invariant. Its restriction to $D_8$ has
nonzero arithmetic restriction and has no $G$-invariant extension
to $D_5$.
\end{proposition}

\begin{proof}
Integrality on $D_7$ and invariance follow from the exact image
and the weight-nine law. Recall the arithmetic construction
\[
 z_7=h z_5h^{-1}z_5^{-t^7}\in J_7,
 \qquad z_5\in J_5,
 \qquad w_9=z_7w_8^{-u}\in J_8,
\]
where $u_2=m_8(z_7)/m_8(w_8)$ and $m_8(w_8)$ is a unit.
The invariant character $\alpha_8$ on $D_5$ vanishes on $z_7$,
since $t$ is odd. Its restriction to $D_7$ is $m_8\bmod2$, so
$u_2$ is even.

Let $\chi:D_5\to\mathbf F_2$ be $G$-invariant. The same odd-power
argument gives $\chi(z_7)=0$. Also
$\chi(w_8^{-u})=-u_2\chi(w_8)=0$. Hence $\chi(w_9)=0$, whereas
$\eta_9(w_9)=1$ by its exact arithmetic valuation. This rules out
an invariant extension from $D_8$ to $D_5$.
\end{proof}

\begin{corollary}\label{cor:weight-nine-character-kernels}
The invariant quadratic-character kernels on restriction to $D_9$
are as follows:
\[
\begin{array}{c|l}
 \text{domain}&\text{kernel}\\ \hline
 D_8&\mathbf F_2\eta_9\\
 D_7&\mathbf F_2\eta_8\oplus\mathbf F_2\eta_9\\
 D_5&\mathbf F_2\eta_7\oplus\mathbf F_2\alpha_8\\
 D_3&\mathbf F_2\eta_5\\
 N&\displaystyle\left(\bigoplus_{q\text{ odd prime}}\mathbf F_2\psi_q\right)
          \oplus\mathbf F_2\eta_3.
\end{array}
\]
\end{corollary}

\begin{proof}
The first two rows follow from the exact scalar and joint quotient
maps. A character on $D_5$ killed on $D_9$ restricts on $D_8$ to
zero or $\eta_9$; the latter is impossible by the preceding
proposition. Its restriction therefore already kills $D_8$,
whose kernel on $D_5$ is the displayed span of $\eta_7,\alpha_8$.
No nonzero combination of those two characters extends to $D_3$,
so the $D_3$ kernel is unchanged. Similarly any invariant character
on $N$ killed on $D_9$ must kill $D_8$, since otherwise its
restriction to $D_5$ would extend $\eta_9$. The preceding kernel
formula for $N$ then gives the last row.
\end{proof}


\section{An exact weight-ten commutator normalization}
\label{sec:weight-ten-normalization}

Use $G=D_9H$ and $J_9=D_9\cap H$. For $g=(1,f)\in D_9$, define
\[
 m_{10}(g)=[S^5TST^3]f.
\]
This mixed coefficient detects the next even-degree bracket.

\begin{theorem}\label{thm:weight-ten-exact-image}
The full and arithmetic images are
\[
 m_{10}(D_9)=m_{10}(J_9)=\mathbf Z_2.
\]
For $D_{10}=D_9\cap\ker m_{10}$ and $J_{10}=D_{10}\cap H$,
\[
 D_9=D_{10}J_9,\qquad G=D_{10}H,\qquad
 B\simeq D_{10}/\overline{[G,D_{10}]J_{10}}.
\]
The subgroup $D_{10}$ is proper, closed and normal, and its words
start in degree at least eleven.
\end{theorem}

\subsection{The primitive leading bracket}

Let $L_w$ denote the standard Lyndon Lie brackets in $S<T$, and set
\[
 \psi_3=L_{SST}-L_{STT},\qquad
 \mathcal L_{10}=\{\psi_3,\mathcal L_7\},
\]
where $\mathcal L_7$ is the primitive degree-seven vector displayed
above. As before,
\[
 \{P,Q\}=[P,Q]+\partial_P(Q)-\partial_Q(P),\qquad
 \partial_P(S)=0,\quad\partial_P(T)=[T,P].
\]
All coefficients of $\mathcal L_{10}$ are integral. Two exact
coefficient calculations are
\begin{equation}\label{eq:weight-ten-unit-coefficient}
 [L_{SSSSSTSTTT}]\mathcal L_{10}=-141,
 \qquad [S^5TST^3]\mathcal L_{10}=-141.
\end{equation}
In particular it is primitive over $\mathbf Z_2$.

\begin{lemma}\label{lem:weight-ten-leading-line}
For every multiplier-one GT word starting in degree ten, the
homogeneous term of degree ten is $A\mathcal L_{10}$ with
$A\in\mathbf Z_2$. It vanishes exactly when $m_{10}$ vanishes.
\end{lemma}

\begin{proof}
The following explicit finite certificate proves the necessary
rational line. Collect the $99+186=285$ standard Lyndon group
commutators of weights ten and eleven. Through degree eleven
write
\[
 F=1+\sum_{10\le|w|\le11}h_w(C_w-1).
\]
Products of two nonconstant factors start in degree twenty, so
this expression is exact at the indicated precision. Define the
cycle forms
\[
 I_v=[v]\bigl(F(Y,X)+F(X,Y)-2\bigr),\qquad
 II_v=[v]\bigl(F(Z,X)+F(Y,Z)+F(X,Y)-3\bigr),
\]
where $X=1+S$, $Y=1+T$ and $Z=Y^{-1}X^{-1}$. Encode the rows by
$2^{|v|}+\operatorname{bin}(v)$ and
$4096+2^{|v|}+\operatorname{bin}(v)$, respectively.

Adjoin one leading Artin pentagon row $P$, with code $8192$.
It is the coefficient of the output word
\[
 11123234324
\]
on $X_1$. Explicitly, use the infinitesimal derivations
$d_{ij}(X_i)=[X_i,X_j]$, $d_{ij}(X_j)=[X_j,X_i]$, with the other
letters fixed. Evaluate the degree-ten Lie polynomial on the
five pairs
\[
 (d_{12},d_{23}+d_{24},+),\quad
 (d_{13}+d_{23},d_{34},+),\quad
 (d_{23},d_{34},-),\quad
 (d_{12}+d_{13},d_{24}+d_{34},-),\quad
 (d_{12},d_{23},-),
\]
and take the indicated coefficient of the signed sum. This is a
necessary coefficient of the full pentagon: a word starting in
degree ten contributes first in Artin output degree eleven, so
only the leading derivations enter. The common sign relating the
ordinary pure Artin action to the displayed $d_{ij}$ is raised to
the even power ten. The row is zero on all degree-eleven columns.

Let $A_R$ be the following $222$ rows, in the written order:
\[
\begin{aligned}
R=(&1025,1027,1029,1031,1033,1035,1037,1039,1041,1043,1045,1047,1049,1051,\\
&1053,1059,1061,1063,1065,1067,1069,1073,1075,1077,1079,1081,1083,1085,\\
&1093,1095,1097,1099,1101,1107,1109,1113,1115,1117,1125,1129,1131,1133,\\
&1139,1171,1173,1197,1203,1237,2049,2051,2053,2055,2057,2059,2061,2063,\\
&2065,2067,2069,2071,2073,2075,2077,2079,2081,2083,2085,2087,2089,2091,\\
&2093,2095,2097,2099,2101,2103,2105,2107,2109,2115,2117,2119,2121,2123,\\
&2125,2127,2129,2131,2133,2135,2137,2139,2141,2147,2149,2151,2153,2155,\\
&2157,2161,2163,2165,2167,2169,2185,2187,2189,2195,2197,2199,2201,2203,\\
&2205,2211,2213,2215,2217,2219,2221,2227,2229,2233,2247,2249,2251,2253,\\
&2259,2261,2265,2267,2277,2281,2341,2347,2349,2355,2357,2379,2381,2387,\\
&2389,5121,5127,5131,5133,5139,5141,5145,5155,5157,5161,5167,5169,5189,\\
&5193,5207,5211,5291,5301,6145,6147,6149,6151,6153,6155,6157,6159,6161,\\
&6163,6165,6167,6169,6171,6177,6179,6181,6183,6185,6187,6189,6191,6193,\\
&6195,6197,6199,6211,6213,6215,6217,6219,6221,6225,6227,6229,6231,6233,\\
&6235,6237,6245,6249,6251,6253,6281,6283,6285,6291,6293,6295,6297,6299,\\
&6307,6309,6313,6315,6317,6325,6345,6347,6357,6437,6443,8192).
\end{aligned}
\]
Let $C$ be the first $222$ columns. Write $\ell_j$ for the
$99$ degree-ten Lyndon coordinates of the explicitly defined
bracket $\mathcal L_{10}$. Its twelfth coordinate is $-141$.
Define the $98$ target rows, of length $285$, by
\[
 T_j=e_j-\frac{\ell_j}{\ell_{12}}e_{12}
       \qquad(1\le j\le99,\ j\ne12).
\]
The exact finite identities are
\[
 \det A_{R,C}=-63256059346747392\equiv37\pmod{101},
 \qquad T_C A_{R,C}^{-1}A_R=T.
\]
Every entry is specified by standard group commutators, geometric
inversion through degree eleven, and the displayed derivation
recursion. Thus the identity forces
$h_j=(\ell_j/\ell_{12})h_{12}$ for every GT word. Its rational
leading scalar is $A=h_{12}/(-141)$, which belongs to
$\mathbf Z_2$ because $141$ is odd. Equation
\eqref{eq:weight-ten-unit-coefficient} gives $m_{10}=-141A$.
This proves the assertion without inferring a lift from the
necessary equations.
\end{proof}

\subsection{An arithmetic unit}

\begin{lemma}\label{lem:weight-ten-arithmetic-unit}
There is an arithmetic commutator $w_{10}\in J_9$ such that
\[
 m_{10}(w_{10})=-141\cdot729\,m_3(\sigma)\tau_7(\sigma)
                                  \in\mathbf Z_2^\times.
\]
Moreover $w_{10}=[\sigma,k]$ for some $k\in D\cap H$.
\end{lemma}

\begin{proof}
Use the established arithmetic element $\sigma\in H\cap M$ with
$\rho_2(\sigma)=2$ and every odd-prime Kummer coordinate zero.
Both $a=m_3(\sigma)$ and $\tau_7(\sigma)$ are odd. Choose
$h\in H$ with $\lambda_2(h)=3$, set $t=\lambda(h)$, and put
\[
 r_3=h\sigma h^{-1}\sigma^{-27},\qquad
 r_5=h r_3h^{-1}r_3^{-243}.
\]
These are multiplier-one arithmetic elements. The first has
$c_3=m_3=0$, so its parameter starts in degree five. The second
has $c_3=c_5=0$, so its parameter starts in degree seven. These
vanishings use the cycle calculations before imposing any radical
normalization.

The weight-seven beta law gives
\[
 c_7(r_5)=\frac{(3^7-3^5)(3^7-3^3)}{360}\tau_7(\sigma)
        =16\cdot729\tau_7(\sigma).
\]
For a word starting in degree seven, $c_7=m_7$. The degree-seven
leading relation, which holds for every multiplier-one GT word
starting in that degree, therefore gives
\[
 \operatorname{in}_7(f_{r_5})=729\tau_7(\sigma)\mathcal L_7.
\]
No normalized weight-seven divisibility claim is being applied to
$r_5$. Its parameter is allowed to have a primitive leading term.
The leading parameter of $\sigma$ is $a\psi_3$. Hence the
commutator
\[
 w_{10}=[\sigma,r_5]
\]
starts in degree ten and has leading term
\[
 729a\tau_7(\sigma)\{\psi_3,\mathcal L_7\}
       =729a\tau_7(\sigma)\mathcal L_{10}.
\]
Every Kummer homomorphism on $M$ vanishes on this commutator. Its
lower Magnus coefficients also vanish, so $w_{10}\in J_9$.
The unit assertion follows from
\eqref{eq:weight-ten-unit-coefficient}.

For the last assertion, the all-prime conjugation law gives
\[
 \rho_q(r_5)=(t-243)(t-27)\rho_q(\sigma)
\]
for every rational prime $q$. Put $c=(t-243)(t-27)$ and
$k=r_5\sigma^{-c}$. The continuous profinite power is defined for
this $c\in\widehat{\mathbf Z}$, and all Kummer coordinates of $k$
are zero. Thus $k\in D\cap H$. Since powers of $\sigma$ commute
with $\sigma$, one has $[\sigma,k]=[\sigma,r_5]=w_{10}$.
\end{proof}

\begin{proof}[Proof of Theorem~\ref{thm:weight-ten-exact-image}]
The leading coefficient $m_{10}$ is an additive continuous
homomorphism on $D_9$, and conjugation acts by
$\lambda_2^{10}$. The arithmetic image contains the unit furnished
by the preceding lemma. Being compact, it is all of
$\mathbf Z_2$, as is the full image. Equality gives
$D_9=D_{10}J_9$ and $G=D_{10}H$; the usual compact gluing argument
identifies the relative quotient $B$ as displayed. The scalar law
proves normality and the unit proves properness. The leading-line
lemma shows that the words of $D_{10}$ start in degree at least
eleven.
\end{proof}

\subsection{The tenth character}

\begin{proposition}\label{prop:eta-ten-extension-kernels}
The invariant character $\eta_{10}=m_{10}\pmod2$ on $D_9$ extends
invariantly to $D_8$ by
\[
 \alpha_{10}(1,f)=[S^5TST^3]f\pmod2,
\]
but has no invariant extension to $D$ or $N$. The original kernel is
\[
 \ker\!\left(H^1(N,\mathbf F_2)^G
                  \longrightarrow H^1(D_{10},\mathbf F_2)\right)
 =\left(\bigoplus_{q\text{ odd prime}}\mathbf F_2\psi_q\right)
                  \oplus\mathbf F_2\eta_3.
\]
The joint raw-coefficient map $(m_9,m_{10})$ on $D_8$ has image
$512\mathbf Z_2\times\mathbf Z_2$ already on $J_8$, with kernel
$D_{10}$. Its invariant quadratic-character kernel on restriction
to $D_{10}$ is $\mathbf F_2\eta_9\oplus\mathbf F_2\alpha_{10}$.
\end{proposition}
\begin{proof}
Indeed the exact weight-nine theorem puts every degree-nine leading
vector on $D_8$ in $32\mathbf Z_2\mathcal L_9$. Modulo two, such
words therefore start in degree at least ten, making the displayed
coefficient additive and invariant under the odd multiplier.

In the other direction, $\eta_{10}$ has no invariant extension to
$D$ or $N$. The arithmetic witness is $w_{10}=[\sigma,k]$ with
$k\in D$, so every such invariant character kills it, whereas
$\eta_{10}(w_{10})=1$. Consequently the original character kernel
is still
\[
 \ker\!\left(H^1(N,\mathbf F_2)^G
                  \longrightarrow H^1(D_{10},\mathbf F_2)\right)
 =\left(\bigoplus_{q\text{ odd prime}}\mathbf F_2\psi_q\right)
                  \oplus\mathbf F_2\eta_3.
\]
To see the last assertion, restrict a character in this kernel to
$D_9$. It is zero or $\eta_{10}$, since $D_9/D_{10}\simeq\mathbf Z_2$;
the second case cannot extend to $N$. The previously proved kernel
on $D_9$ then applies.

Finally, extend $m_{10}$ to $D_8$ by the same raw coefficient.
The map $(m_9,m_{10})$ is additive there, since all
nonconstant product and substitution corrections start above
degree ten. The arithmetic witnesses $w_9,w_{10}$ show that its
image is $512\mathbf Z_2\times\mathbf Z_2$, with kernel $D_{10}$.
Thus the invariant quadratic-character kernel on $D_8$ is
$\mathbf F_2\eta_9\oplus\mathbf F_2\alpha_{10}$.

\end{proof}

\subsection{A group-theoretic extension test in weight ten}
\label{subsec:eta-ten-extension-threshold}

\begin{proposition}\label{prop:eta-ten-no-extension-d7}
The character $\eta_{10}$ on $D_9$ extends invariantly to $D_8$
by $\alpha_{10}$, but has no $G$-invariant extension to $D_7$.
Consequently it has no such extension to $D_5$, $D_3$, $D$ or $N$.
\end{proposition}

\begin{proof}
The extension to $D_8$ was proved above. We give an obstruction
using only class-two commutator identities and the previously
established leading terms.

For $r\ge3$, let $\mathcal F^r$ be the subgroup of
$M=\ker\lambda$ whose pro-$2$ Magnus parameter starts in degree
at least $r$. The filtered normalized group law gives
\[
 [\mathcal F^r,\mathcal F^s]\subseteq\mathcal F^{r+s}.
\]
One has $M=\mathcal F^3$ and
$[M,M]\subseteq D_7\subseteq\mathcal F^8$. Thus
$[M,[M,M]]\subseteq\mathcal F^{11}$, and the quotient
$M/\mathcal F^{11}$ has nilpotence class at most two.
Moreover,
\[
 [\mathcal F^5,\mathcal F^5]\subseteq\mathcal F^{11}.
\]
Indeed its possible degree-ten leading term is the Ihara bracket
of two multiples of $\psi_5$, and this is zero. Vanishing over
$\mathbf Q_2$ implies vanishing in the integral leading module,
which is torsion-free.

Keep the arithmetic $\sigma,h$ of the weight-ten construction, and
put
\[
 a_0=\sigma,\qquad a_1=h\sigma h^{-1},\qquad
 a_2=h^2\sigma h^{-2},\qquad c_{ij}=[a_i,a_j].
\]
The weight-three law shows that
\[
 r=a_1a_0^{-27},\qquad q=a_2a_0^{-729}
\]
belong to $\mathcal F^5$. Therefore $[r,q]=1$ in
$M/\mathcal F^{11}$. All commutators are central in this quotient,
and expanding this identity gives
\[
 1=[a_1a_0^{-27},a_2a_0^{-729}]
   =c_{12}c_{01}^{729}c_{02}^{-27}.
\]
Hence $c_{12}=c_{02}^{27}c_{01}^{-729}$ there.

The arithmetic weight-eight and weight-ten witnesses are
\[
 w_8=c_{01},\qquad
 w_{10}=[a_0,(a_2a_1^{-27})(a_1a_0^{-27})^{-243}].
\]
The same class-two calculation gives
\[
 w_{10}=c_{02}c_{01}^{-270}\pmod{\mathcal F^{11}}.
\]
Consequently the word
\[
 z=h w_8h^{-1}w_8^{-6561}=c_{12}c_{01}^{-6561}
\]
satisfies
\[
 z=c_{02}^{27}c_{01}^{-7290}
   =w_{10}^{27}\pmod{\mathcal F^{11}}.
\]
Since $z\in J_7$ and its parameter agrees through degree ten with
that of $w_{10}^{27}$, one has $z\in J_9$ and
\[
 m_{10}(z)=27m_{10}(w_{10}),\qquad\eta_{10}(z)=1.
\]
On the other hand, every $G$-invariant character
$\chi:D_7\to\mathbf F_2$ satisfies
\[
 \chi(z)=\chi(h w_8h^{-1})-6561\chi(w_8)=0,
\]
because $6561$ is odd. Thus such a character cannot restrict to
$\eta_{10}$ on $D_9$. Any extension to one of the larger listed
subgroups would restrict to an extension on $D_7$, so those are
also impossible.
\end{proof}

\begin{corollary}\label{cor:weight-ten-character-kernels}
The invariant quadratic-character kernels on restriction to
$D_{10}$ are
\[
\begin{array}{c|l}
 \text{domain}&\text{kernel}\\ \hline
 D_9&\mathbf F_2\eta_{10}\\
 D_8&\mathbf F_2\eta_9\oplus\mathbf F_2\alpha_{10}\\
 D_7&\mathbf F_2\eta_8\oplus\mathbf F_2\eta_9\\
 D_5&\mathbf F_2\eta_7\oplus\mathbf F_2\alpha_8\\
 D_3&\mathbf F_2\eta_5\\
 N&\displaystyle\left(\bigoplus_{q\text{ odd prime}}\mathbf F_2\psi_q\right)
                    \oplus\mathbf F_2\eta_3.
\end{array}
\]
\end{corollary}

\begin{proof}
The first two rows follow from the exact scalar and joint maps.
For any larger listed domain, a character killed on $D_{10}$
restricts on $D_9$ to zero or $\eta_{10}$. The latter is impossible
by the preceding proposition. Such a character therefore already
kills $D_9$, and the previously proved weight-nine kernel table
applies.
\end{proof}


\section{Combining the odd and mixed normalizations}
\label{sec:dyadic-metabelian-intersection}

The arithmetic commutators in degrees eight and ten have trivial beta
multiplier. They therefore let us combine the infinite odd normalization
with the finite leading-term normalization.

\begin{corollary}\label{cor:dyadic-metabelian-deep-gate}
Put
\[
 E=C_\infty\cap D_{10},\qquad J_E=E\cap H.
\]
Then $E$ is a proper closed normal subgroup of $D_{10}$ and of
infinite index there, and
\[
 G=EH,\qquad B\simeq E/\overline{[G,E]J_E}.
\]
Every pro-$2$ Magnus word on $E$ starts in degree at least eleven,
and its dyadic beta multiplier is identically one.
The kernel of restriction from $H^1(N,\mathbf F_2)^G$ to
$H^1(E,\mathbf F_2)$ is exactly
\[
 \left(\bigoplus_{q\text{ odd prime}}\mathbf F_2\psi_q\right)
                   \oplus\mathbf F_2\eta_3.
\]
\end{corollary}

\begin{proof}
We first verify the factorization without assuming any joint-image
statement. Since $C_\infty\subseteq D_7$, the coordinate $m_8$
is defined there. The arithmetic word $w_8$ of the eighth
normalization is a commutator of multiplier-one elements. The
beta multiplier is a homomorphism on $M=\ker\lambda$, so
$B_{w_8}=1$. Hence $w_8\in C_\infty\cap H$, and its $m_8$
value is a unit. Correcting by its powers gives
\[
 C_\infty=(C_\infty\cap D_8)(C_\infty\cap H).
\]
Since $m_9=0$ on $C_\infty$, the first factor lies in $D_9$.
The arithmetic word $w_{10}$ of the tenth normalization is also
a commutator of multiplier-one elements, so belongs to
$C_\infty\cap J_9$ and has an odd $m_{10}$ value. Correcting
by its powers gives
\[
 C_\infty\cap D_8=E(C_\infty\cap D_8\cap H).
\]
Together with $G=C_\infty H$, these equations prove $G=EH$.
Normality and the relative quotient follow as before.

To prove strictness, take the word $v_5\in C_4\cap H$ in
\eqref{eq:uniform-odd-arithmetic-witness}. It has
$m_{11}(v_5)\in2^{11}\mathbf Z_2^{\times}$.
First correct its eighth coordinate by a power of $w_8$.
The result lies in $C_4\cap D_8\subseteq D_9$.
Then correct its tenth coordinate by a power of $w_{10}$.
The result $v$ belongs to $D_{10}\cap H$.
Both correcting words have beta multiplier one, so they leave
$\tau_{11}$, and therefore the nonzero eleventh depth-one
coefficient, unchanged. Thus $v\notin E$.
In fact $D_{10}\subseteq C_4$, so the uniform bound and the
corrected arithmetic word give the exact images
\[
 m_{11}(D_{10})=m_{11}(D_{10}\cap H)=2^{11}\mathbf Z_2.
\]
This additive coordinate vanishes on $E$, proving infinite index. The final assertions follow
from $E\subseteq D_{10}$ and its definition.

For the last kernel assertion, every invariant character $\chi$ on
$N$ kills both correcting words. Indeed,
\[
 w_8=[\sigma,h\sigma h^{-1}]=[\sigma,k_8],\qquad
 k_8=(h\sigma h^{-1})\sigma^{-t}\in D\subseteq N;
\]
the rational Kummer coordinates of $k_8$ cancel by their exact
weight-one law. Invariance under $\sigma\in G$ gives
$\chi(w_8)=0$. The tenth construction gives
$w_{10}=[\sigma,k_{10}]$ for $k_{10}\in D$ and hence
$\chi(w_{10})=0$ as well. The two correction factorizations in
this proof now imply that $\chi$ vanishes on $E$ if and only if
it vanishes on $C_\infty$. Proposition~\ref{prop:uniform-odd-character-kernels}
gives the asserted kernel.
\end{proof}

\section{A primitive mixed coordinate in degree eleven}
\label{sec:eleven-mixed-mixed}

We remove a primitive mixed coefficient while preserving the arithmetic
factorization. Use the groups $G$, $H$, $N$, $D$, and $D_{10}$ defined above.
Write $M=\ker\lambda$ and
\[
 C_\infty=\{g\in D:B_g(X,Y)=1\},\qquad
 E=C_\infty\cap D_{10},\qquad J(E)=E\cap H.
\]
Here $B_g$ is the dyadic beta multiplier. The preceding normalizations
give $G=EH$, and every Magnus series $f_g$ with $g\in E$ has no
nonconstant term below degree eleven. We use these established facts
and the exact degree-eight leading lattice below. The argument uses one coordinate of the degree-eleven leading space.

\begin{theorem}\label{thm:eleven-mixed-mixed-reduction}
The coordinate
\[
 b_{11}(g)=[S^6TST^3]f_g\qquad(g\in E)
\]
is a continuous additive homomorphism satisfying
\[
 b_{11}(aga^{-1})=\lambda_2(a)^{11}b_{11}(g)
 \quad(a\in G,\ g\in E).
\]
Its full and arithmetic images are both saturated:
\[
 b_{11}(E)=b_{11}(J(E))=\mathbf Z_2.
\]
Consequently
\[
 E'=\ker(b_{11}|_E)
\]
is a proper closed normal subgroup of $G$, and
\[
 E/E'\simeq\mathbf Z_2,\qquad G=E'H,
 \qquad
 B\simeq E'/\overline{[G,E'](E'\cap H)}.
\]
For $X_2=H^1(N,\mathbf F_2)^G$, its restriction kernel remains
\[
 \ker\bigl(X_2\longrightarrow H^1(E',\mathbf F_2)\bigr)
 =\left(\bigoplus_{q\text{ odd prime}}\mathbf F_2\psi_q\right)
       \oplus\mathbf F_2\eta_3.
\]
The invariant character $b_{11}\bmod2$ on $E$ is nonzero on $J(E)$
and has no $G$-invariant extension to $N$.
\end{theorem}

\begin{proof}
We first give the integral coefficient calculation that supplies an
arithmetic unit. For Lie polynomials use
\[
 P_3=[S,[S,T]]+[T,[S,T]],\qquad
 \{A,B\}=[A,B]+\partial_A B-\partial_B A,
\]
where $\partial_A(S)=0$ and $\partial_A(T)=[T,A]$.
This is the Ihara bracket in the convention of the GT product.
Let $L_8$ be the integer Lie polynomial for the already established
degree-eight leading line, normalized so that
\[
 [S^4TST^2]L_8=7.
\]
For specificity it is completely specified by its coordinates in the
lexicographically ordered Lyndon basis of weight eight, with $S<T$:
\begin{align*}
(&0,0,-4,0,-2,7,17,0,-1,-12,-6,4,-3,-6,0,\\
 &-5,3,3,6,7,21,3,-15,-3,0,8,-11,-4,-10,0).
\end{align*}
The Lyndon bracket of a nonletter word is $[L_u,L_v]$, where $v$
is its longest proper Lyndon suffix and the word is $uv$.
This finite recursion gives, in particular,
\begin{equation}\label{eq:eleven-mixed-four-coefficients}
 [S^4TST^2]L_8=7,\qquad [S^5TST]L_8=-4,\qquad
 [S^5T^3]L_8=[S^6T^2]L_8=0.
\end{equation}
An element $z\in D_7$ therefore has degree-eight term
$(m_8(z)/7)L_8$.

Put $\omega=S^6TST^3$. Expanding
\[
 P_3=S^2T-2STS+TS^2+2TST-T^2S-ST^2
\]
shows that $[\omega]P_3L_8=[\omega]L_8P_3=0$: the required
prefix or suffix of length three is absent from $P_3$. Moreover
\[
 [T,P_3]=-2TSTS+T^2S^2+3T^2ST-T^3S-3TST^2
             -S^2T^2+2STST+ST^3.
\]
Replacing one $T$ in a word of $L_8$ by this expression gives
\[
 [\omega]\partial_{P_3}L_8
     =2[S^5T^3]L_8-2[S^6T^2]L_8=0.
\]
For $\partial_{L_8}P_3$, the only surviving contributions are
$-S^2L_8T$ from $S^2T$ and $+SL_8T^2$ from $-ST^2$.
Thus the final, negative derivation term contributes
\[
 -[\omega]\partial_{L_8}P_3
    =[S^4TST^2]L_8-[S^5TST]L_8=7+4=11.
\]
The four Ihara summands have coefficients $(0,0,0,11)$, proving
\begin{equation}\label{eq:eleven-mixed-bracket-unit}
 [S^6TST^3]\{P_3,L_8\}=11.
\end{equation}

Choose the arithmetic elements from the earlier exact
normalizations: $\sigma\in H\cap M$ with
$\rho_2^I(\sigma)=2$ and all odd-prime rational Kummer coordinates
zero, and $h\in H$ with $t=\lambda(h)$ and $t_2=3$.
Let
\[
 a=m_3(\sigma),\qquad u=\tau_5(\sigma).
\]
The established integral moment parity gives $a,u\in\mathbf Z_2^\times$.
The arithmetic word
\[
 w_8=[\sigma,h\sigma h^{-1}]
\]
belongs to $D_7\cap H$ and satisfies the exact earlier formula
\begin{equation}\label{eq:eleven-mixed-old-eight}
 m_8(w_8)=63au.
\end{equation}
Recall that all rational Kummer coordinates and the dyadic beta
multiplier are homomorphisms on $M$, the latter into a commutative
group. In particular $B_{w_8}=1$.

Define the new arithmetic commutator
\[
 w_{11}=[\sigma,w_8].
\]
The first nonconstant terms of $f_\sigma$ and $f_{w_8}$ have
degrees three and eight. The filtered GT commutator has no term
below degree eleven, and its term in degree eleven is their
Ihara bracket. Its rational Kummer coordinates and beta multiplier
are trivial, since both inputs lie in $M$. Hence $w_{11}\in E\cap H$.
Equations \eqref{eq:eleven-mixed-bracket-unit} and
\eqref{eq:eleven-mixed-old-eight} now give the exact value
\begin{equation}\label{eq:eleven-mixed-arithmetic-unit}
 b_{11}(w_{11})
  =\frac{11}{7}a\,m_8(w_8)
  =99a^2u\in\mathbf Z_2^\times.
\end{equation}
There are no contributions from higher terms of either input to
degree eleven.

On a subgroup whose words start in degree eleven, the degree-eleven
part of the GT product is the sum of the two leading parts.
Conjugation by an arbitrary element of $G$ multiplies this
homogeneous part by $\lambda_2^{11}$: on the associated graded
algebra both generators are multiplied by $\lambda_2$, and all
conjugating corrections have higher degree. This proves the
claimed homomorphism and weight laws. Coefficient extraction is
continuous. The closed arithmetic image contains the image of the
procyclic group generated by $w_{11}$, which is all of $\mathbf Z_2$
by \eqref{eq:eleven-mixed-arithmetic-unit}. This proves both image
equalities and the normality and properness of $E'$.

For each $e\in E$, choose $r\in\widehat{\mathbf Z}$ whose
2-adic component is $b_{11}(e)/b_{11}(w_{11})$. The legitimate
profinite power $w_{11}^r$ lies in $J(E)$ and has the same
coordinate as $e$. Thus $e w_{11}^{-r}\in E'$, proving
$E=E'J(E)$ and hence $G=E'H$.

For completeness, let $K=N\cap H$ and
$L=\overline{[G,E']}$. Since $N=E'K$, the image of $E'$ is
central in $G/L$, while $[H,K]\subseteq K$. Consequently
\[
 \overline{[G,N]K}=LK.
\]
The product is closed because both factors are compact and $L$
is normal. Also $E'\cap LK=L(E'\cap H)$. The inclusion of $E'$
therefore induces the stated canonical presentation of $B$.

Finally let $\chi\in X_2$. Since $w_8\in N$ and $\sigma\in G$,
invariance gives
\[
 \chi(w_{11})=\chi(\sigma w_8\sigma^{-1})-\chi(w_8)=0.
\]
If $\chi$ vanishes on $E'$, the same correction by $w_{11}^r$
shows that it vanishes on $E$. The converse is immediate, so the
restriction kernels to $E$ and $E'$ coincide. The previously proved
kernel for $E$ is exactly the displayed rational Kummer span plus
$\mathbf F_2\eta_3$. On the other hand $b_{11}\bmod2$ is invariant
and takes value one on $w_{11}$, proving both its arithmetic
nonvanishing and its nonextension to $N$.
\end{proof}

\begin{remark}\label{rem:eleven-mixed-scope}
The character $b_{11}\bmod2$ is nonzero on the arithmetic subgroup,
so it gives no nonzero character of the relative quotient $B$.
The next proposition shows that this coordinate detects the entire
degree-eleven leading term on $E$.
\end{remark}

\subsection{The entire beta-trivial leading term in degree eleven}
\label{subsec:eleven-completeness}

Willwacher's rational dimension bound identifies the complete leading
term removed by Theorem~\ref{thm:eleven-mixed-mixed-reduction}. The
arithmetic unit then gives its exact integral lattice.

For $g\in\ker\lambda$ whose dyadic Magnus word has no term below
degree $n$, write $\operatorname{in}_n(g)$ for its degree-$n$
homogeneous term. Write $F^n$ for this closed word filtration.

\begin{lemma}\label{lem:full-gt-leading-term}
If $g$ is a full-GT element in $F^n$, with $n\ge3$, then
\[
 \operatorname{in}_n(g)\in
 \mathfrak{grt}_1(\mathbf Q_2)_n
       \cap\mathbf Z_2\langle S,T\rangle_n.
\]
Furthermore
\[
 \mathfrak{grt}_1(\mathbf Q_2)_n
  =\mathfrak{grt}_1(\mathbf Q)_n\otimes_{\mathbf Q}\mathbf Q_2.
\]
\end{lemma}

\begin{proof}
Project the profinite word and its full pentagon relation to the
pro-$2$ completions, then to the rational prounipotent completions.
The pro-$\ell$ to prounipotent passage is explicitly recalled in
Furusho's proof of the pro-$\ell$ corollary
\cite[\S3, p.~554, immediately before Corollary~12]{Furusho2010}.
Set $x=e^S$, $y=e^T$ in the free prounipotent group. The leading
term is unchanged by this substitution, because
$e^S-1=S+O(S^2)$ and similarly for $T$. Since the image of $f_g$
is group-like, its first nonzero term is a primitive element of
the free associative Hopf algebra in characteristic zero, hence
a homogeneous free Lie polynomial.

Taking the first nonzero degree in the full pentagon gives
\begin{align*}
 q(t_{12},t_{23}+t_{24})+q(t_{13}+t_{23},t_{34})
 ={}&q(t_{23},t_{34})\\
   &+q(t_{12}+t_{13},t_{24}+t_{34})+q(t_{12},t_{23}),
\end{align*}
where $q=\operatorname{in}_n(g)$ and the equality is in the
Drinfeld--Kohno Lie algebra $\mathfrak t_4(\mathbf Q_2)$.
Indeed, the degree-one class of a pure braid generator is
$t_{ij}$, a product of generators has the sum of these classes,
and a product of two nonconstant $f_g$ terms has degree at least
$2n>n$. This is the associated graded form of the usual
prounipotent pure braid realization. Equivalently it follows by
the associator transport used in the proof of
\cite[proof of Theorem~2, p.~554]{Furusho2010}: that transport preserves
linear terms, so it preserves the displayed leading equation.

Duality and the multiplier-one hexagon similarly give the
antisymmetry and three-cycle relations. In characteristic zero
the cabling relation is redundant, as recalled with Drinfeld's
Proposition~5.7 in \cite[p.~1, after equation~(6)]{Willwacher}.
Thus $q$ satisfies precisely the defining equations (1)--(4)
of that source. One may also use the pentagon implication
\cite[Theorem~3, pp.~547--548]{Furusho2010}, together with the
same cabling redundancy. The hypotheses on degree exclude all
linear and quadratic exceptional terms. This proves the first
assertion. It is a necessary leading-term assertion only; no
surjectivity from profinite GT to a prounipotent group is used.

In fixed degree the free Lie and infinitesimal braid Lie algebras
have finite bases over $\mathbf Q$, and all defining equations
have rational, indeed integral, coefficients. This integral
matrix description is noted explicitly in
\cite[\S1.1, p.~2]{Willwacher}. Choose these bases and write the
finite defining map as a matrix $A_n$ over $\mathbf Q$.
Its rank is unchanged after extension to $\mathbf Q_2$, since a
minor is nonzero over one characteristic-zero field if and only
if it is nonzero over the other. Equivalently, tensoring the
kernel sequence with the field extension is exact. Hence
$\ker(A_n\otimes\mathbf Q_2)=(\ker A_n)\otimes\mathbf Q_2$,
which is the asserted base-change identity.
\end{proof}

\begin{proposition}\label{prop:eleven-complete-leading-line}
Let $L_8$ and $P_3$ have the normalizations in
Theorem~\ref{thm:eleven-mixed-mixed-reduction}, and put
\[
 V_{11}=\frac1{11}\{P_3,L_8\}.
\]
Then $V_{11}$ is an integral $2$-adic Lie polynomial with
$[S^6TST^3]V_{11}=1$, and
\[
 \operatorname{in}_{11}(g)=b_{11}(g)V_{11}\quad(g\in E).
\]
In particular
\[
 \operatorname{in}_{11}(E)
 =\operatorname{in}_{11}(E\cap H)=\mathbf Z_2V_{11},
 \qquad E'=E\cap F^{12}=C_\infty\cap F^{12}.
\]
Thus every dyadic word belonging to $E'$ starts in degree twelve.
All arithmetic factorizations, relative quotient identifications,
and character-kernel conclusions of the strict reduction remain
valid for this deeper subgroup.
\end{proposition}

\begin{proof}
Willwacher's dimension bound
\cite[Theorem~1.1, p.~2]{Willwacher} gives
\[
 \dim_{\mathbf Q_2}\mathfrak{grt}_1(\mathbf Q_2)_{11}
 =\dim_{\mathbf Q}\mathfrak{grt}_1(\mathbf Q)_{11}\le d_{11}=2.
\]
Here the equality of dimensions is
Lemma~\ref{lem:full-gt-leading-term}. To check the last integer
directly, the free Lie algebra on generators of odd weights
$3,5,7,\ldots$ has only two possibilities in weight eleven:
the generator of weight eleven, and the multidegree with two
weight-three letters and one weight-five letter. The latter
has one-dimensional span, represented by
$[\sigma_3,[\sigma_3,\sigma_5]]$; brackets of two generators
have even weight and brackets of at least four have weight at
least twelve. We need no injectivity theorem for free generators.

We supply the two independent directions using arithmetic
elements. The preceding strict reduction already supplies
$w_{11}\in E\cap H$ with
\[
 [S^{10}T]\operatorname{in}_{11}(w_{11})=0,\qquad
 b_{11}(w_{11})=99a^2u\in\mathbf Z_2^\times.
\]
The first equality follows because its beta multiplier is one.

For the second direction, let
\[
 v_5=\Delta_9\Delta_7\Delta_5\Delta_3\Delta_1(\sigma),
 \qquad \Delta_j(z)=hzh^{-1}z^{-t^j},\quad t_2=3.
\]
The uniform odd normalization gives $v_5\in C_4\cap H$ and
\[
 m_{11}(v_5)=
 \frac{2\prod_{j=0}^{4}(3^{11}-3^{2j+1})}{10!}
                     \tau_{11}(\sigma)
       \in2^{11}\mathbf Z_2^\times.
\]
Indeed the numerator product has valuation $18$, the extra
factor two adds one, and $v_2(10!)=8$; the moment is odd.
Here $C_4$ kills the depth-one weights $3,5,7,9$ but need not
kill the mixed terms of weights eight and ten.

Use the established beta-trivial arithmetic units $w_8$ and
$w_{10}$, with $m_8(w_8)$ and $m_{10}(w_{10})$ odd. First multiply $v_5$ by
a suitable profinite power of $w_8$ to kill $m_8$. The result
is in $D_8\cap C_4\cap H$, hence in $D_9$, since its
weight-nine depth-one coordinate is still zero. Next multiply
by a suitable power of $w_{10}$ to kill $m_{10}$. This gives
$z_{11}\in D_{10}\cap H$. Both correcting words have beta
multiplier one, so all beta coefficients of $z_{11}$ equal
those of $v_5$. In particular
\[
 [S^{10}T]\operatorname{in}_{11}(z_{11})
       =m_{11}(z_{11})=m_{11}(v_5)\ne0.
\]
The established degree-ten normalization ensures that
$z_{11}$ has no term below degree eleven.

By Lemma~\ref{lem:full-gt-leading-term}, both leading polynomials
belong to $\mathfrak{grt}_1(\mathbf Q_2)_{11}$. They are linearly
independent, as detected by $[S^{10}T]$ and $b_{11}$. This proves
that the dimension is exactly two and that the kernel of the
depth-one functional $[S^{10}T]$ is the one-dimensional line
spanned by $\operatorname{in}_{11}(w_{11})$.

Every $g\in E$ lies in that kernel, since $B_g=1$. Its leading
term is therefore a scalar multiple of $V_{11}$. The exact
coefficient identity
$[S^6TST^3]\{P_3,L_8\}=11$ identifies that scalar with
$b_{11}(g)$. Division by eleven is harmless in $\mathbf Z_2$,
so $V_{11}$ itself is integral. The arithmetic unit from the
strict reduction realizes all scalars in $\mathbf Z_2$ and
establishes both leading-image equalities. This is the integral
saturation step; it does not follow from rational dimension
alone.

It follows that $b_{11}(g)=0$ on $E$ if and only if its entire
degree-eleven term vanishes. This proves $E'=E\cap F^{12}$.
Finally $C_\infty\cap F^{12}\subset D_{10}$, because all the
earlier mixed and depth-one coordinates vanish when the word
starts in degree twelve. The last displayed equality follows.
\end{proof}


\section{Two saturated arithmetic directions in degree twelve}
\label{sec:twelve-normalization}

Write $E_{12}=C_\infty\cap F^{12}=E'$ for the group just obtained,
where $F^n$ denotes the dyadic word filtration. The next entire
leading layer can also be removed arithmetically.

\begin{theorem}\label{thm:twelve-normalization}
On $E_{12}$ put
\[
 \begin{split}
 b_{12,1}(g)&=[S^6TST^4]f_g,\\
 b_{12,2}(g)&=[S^5TST^2ST^2]f_g,
 \end{split}
 \qquad b_{12}=(b_{12,1},b_{12,2}).
\]
This is a continuous additive map of scalar conjugation weight
$\lambda_2^{12}$, and
\[
 b_{12}(E_{12})=b_{12}(E_{12}\cap H)=\mathbf Z_2^2.
\]
Its kernel is exactly
\[
 E_{13}=C_\infty\cap F^{13}.
\]
Consequently $E_{13}$ is a proper closed normal subgroup, with
\[
 E_{12}/E_{13}\simeq\mathbf Z_2^2,\qquad G=E_{13}H,
 \qquad
 B\simeq E_{13}/\overline{[G,E_{13}](E_{13}\cap H)}.
\]
The kernel of restriction of $X_2=H^1(N,\mathbf F_2)^G$ to
$E_{13}$ is again the odd-prime rational Kummer span plus
$\mathbf F_2\eta_3$. No nonzero linear combination of
$b_{12,1}\bmod2$ and $b_{12,2}\bmod2$ on $E_{12}$ extends to
a $G$-invariant character on $N$.
\end{theorem}

\begin{proof}
We first specify the two integral Lie directions and an exact
coefficient certificate. In lexicographically ordered Lyndon
bases, with the longest-proper-Lyndon-suffix bracketing and $S<T$,
let $L_5,L_7,L_9$ have the following integer coordinates:
\begin{align*}
L_5:\quad & (2,-4,3,4,1,-2),\\
L_7:\quad & (16,-48,80,80,-32,19,-173,-80,-32,17,-99,61,\\
          &\hspace{12mm}109,48,-65,64,48,-16),\\
L_9:\quad & (144,-576,1512,1344,-1440,-346,-6322,-2016,\\
 &360,-1598,3214,-2337,2722,4724,8823,2016,\\
 &600,2148,3259,-1118,-5303,-2951,4353,-682,\\
 &-11276,479,-1762,-8109,-6663,-1344,-89,-5496,\\
 &-1949,291,1041,7472,-4352,-2416,-3686,-1013,\\
 &2059,5340,-3340,4388,1868,3558,2434,576,\\
 &3299,-5446,-1768,6452,1368,2088,936,-144).
\end{align*}
These are the already established low-weight leading vectors,
normalized by
\[
 [S^4T]L_5=2,\qquad [S^6T]L_7=16,\qquad [S^8T]L_9=144.
\]
For any full-GT word starting in the indicated degree, its leading
term lies on the corresponding rational line. The earlier finite
leading calculations give this fact directly. It also follows from
Lemma~\ref{lem:full-gt-leading-term} and Willwacher's upper bound,
since $d_5=d_7=d_9=1$ and the known nonzero arithmetic leading
vectors supply the lower bounds. No rational Kummer normalization
is required for this statement about a word that already starts
in the specified degree.

Set
\[
 V=\{P_3,L_9\},\qquad W=\{L_5,L_7\},\qquad
 \omega_1=S^6TST^4,\quad\omega_2=S^5TST^2ST^2.
\]
Here the bracket is the same Ihara bracket used above. The exact
coefficient calculation is
\begin{equation}\label{eq:twelve-unit-minor}
 \begin{pmatrix}
 [\omega_1]V&[\omega_1]W\\
 [\omega_2]V&[\omega_2]W
 \end{pmatrix}
 =\begin{pmatrix}3385&46\\-4694&-1397\end{pmatrix},
 \qquad \det=-4512921\in\mathbf Z_2^\times.
\end{equation}
For a short reproducible certificate, the contributions of
$AB,-BA,\partial_A B,-\partial_B A$ are respectively
\[
\begin{array}{c|r|r|r|r|r}
(A,B;\omega)&AB&-BA&\partial_A B&-\partial_B A&\text{sum}\\\hline
(P_3,L_9;\omega_1)&0&0&-6720&10105&3385\\
(P_3,L_9;\omega_2)&0&-4378&12412&-12728&-4694\\
(L_5,L_7;\omega_1)&0&32&-2208&2222&46\\
(L_5,L_7;\omega_2)&0&-1152&1887&-2132&-1397
\end{array}
\]
All entries follow from the displayed Lyndon vectors by ordinary
concatenation and the finite coefficient rule
\[
 [\omega]\partial_A B=
 \sum_{B\ni w}\sum_{i:w_i=T}
       [w]B\,[\omega]\bigl(w_{<i}(TA-AT)w_{>i}\bigr).
\]
This specifies the entire integer calculation independently of
any software. In particular the two columns in
\eqref{eq:twelve-unit-minor} reduce to the identity matrix
modulo two.

We now realize odd multiples of both columns arithmetically.
Use $\sigma,h,t$ as in the previous theorem, with $t_2=3$, and
write $\tau_j=\tau_j(\sigma)$ and $a=m_3(\sigma)$; all of
$a,\tau_5,\tau_7,\tau_9$ are odd. Put
\[
 u_5=\Delta_3(\sigma),\qquad
 u_7=\Delta_5(u_5),\qquad
 \widetilde u_9=\Delta_7(u_7),\qquad
 \Delta_j(z)=hzh^{-1}z^{-t^j}.
\]
Every exponent $t^j$ is the full profinite multiplier power.
The operator $\Delta_j$ kills the homogeneous leading term of
degree $j$. There is no full-GT leading term in degrees four or
six: Lemma~\ref{lem:full-gt-leading-term} and the bound $d_4=d_6=0$
give this immediately. Thus $u_5\in F^5$, $u_7\in F^7$, and
$\widetilde u_9\in F^8$.

The degree-eight leading space is the single line detected by
$m_8$, and $w_8$ has an odd $m_8$ value. Choose
$r\in\widehat{\mathbf Z}$ with
$r_2=m_8(\widetilde u_9)/m_8(w_8)$ and put
\[
 u_9=\widetilde u_9 w_8^{-r}.
\]
Then $u_9\in F^9\cap H$. Since $B_{w_8}=1$, this correction
does not change any beta moment.

The additive weight laws for the moments and the triangular
depth-one formulas give the exact leading terms
\begin{equation}\label{eq:twelve-primitive-inputs}
 \begin{split}
 \operatorname{in}_5(u_5)&=9\tau_5 L_5,\\
 \operatorname{in}_7(u_7)&=729\tau_7 L_7,\\
 \operatorname{in}_9(u_9)&=2302911\tau_9 L_9.
 \end{split}
\end{equation}
Indeed, after the preceding moments have been killed, the raw
depth-one coordinate in degree $2k+1$ is
$2\tau_{2k+1}/(2k)!$. Therefore the three raw coordinates are
\[
 \begin{split}
 m_5(u_5)&=(3^5-3^3)\tau_5/12=18\tau_5,\\
 m_7(u_7)&=(3^7-3^5)(3^7-3^3)\tau_7/360=11664\tau_7,\\
 m_9(u_9)&=(3^9-3^7)(3^9-3^5)(3^9-3^3)\tau_9/20160\\
         &=331619184\tau_9.
 \end{split}
\]
Dividing these by the displayed depth-one coefficients $2,16,144$
gives \eqref{eq:twelve-primitive-inputs}. Only odd scalars
multiply the primitive integer vectors. This construction does
not divide an arithmetic group element by two.

Define
\[
 z_{3,9}=[\sigma,u_9],\qquad z_{5,7}=[u_5,u_7].
\]
These lie in $H\cap C_\infty\cap F^{12}=H\cap E_{12}$:
their leading degrees are at least $3+9$ and $5+7$, and every
rational Kummer coordinate and the beta multiplier vanish on
commutators in $M$. Their leading terms are
\[
 \operatorname{in}_{12}(z_{3,9})=2302911a\tau_9 V,
 \qquad
 \operatorname{in}_{12}(z_{5,7})=6561\tau_5\tau_7 W.
\]
Thus both columns of \eqref{eq:twelve-unit-minor} are realized
up to odd arithmetic factors.

Willwacher's Theorem~1.1 and the base-change lemma give
$\dim\mathfrak{grt}_1(\mathbf Q_2)_{12}\le d_{12}=2$.
To count $d_{12}$ directly, the only nonzero bracket types are
$(3,9)$ and $(5,7)$: four weight-three letters span zero in a
free Lie algebra on that single letter, and all other generator
lengths or weights are excluded. The two arithmetic leading
vectors just constructed are independent by the odd determinant,
so they span the full rational leading space in degree twelve.

Every leading polynomial $q$ of an element of $E_{12}$ is
therefore $cV+dW$ with $c,d\in\mathbf Q_2$. Its two selected
Magnus coefficients are integral. The inverse of the matrix in
\eqref{eq:twelve-unit-minor} is integral over $\mathbf Z_2$,
so $c,d\in\mathbf Z_2$. Conversely, products of profinite powers
of $z_{3,9},z_{5,7}$ realize every such pair, since their scalar
factors are odd. Hence the entire integral leading lattice, and
its arithmetic image, is $\mathbf Z_2V+\mathbf Z_2W$.

Leading coefficients on $F^{12}$ add under the GT product and
have scalar conjugation weight $\lambda_2^{12}$. The unit minor
shows that vanishing of $b_{12}$ is equivalent to vanishing of
the whole degree-twelve term. This proves the asserted kernel
and image statements. Arithmetic correction by the two words
gives $E_{12}=E_{13}(E_{12}\cap H)$, whence $G=E_{13}H$.
The same closed normal subgroup argument used in the strict
degree-eleven reduction identifies the relative quotient $B$.

Finally the two correcting words are killed by every original
$N$-invariant character. To see this with all coefficient primes
preserved, put
\[
 c=(t-t^7)(t-t^5)(t-t^3)\in\widehat{\mathbf Z}.
\]
All rational Kummer coordinates satisfy
$\rho_q(u_9)=c\rho_q(\sigma)$; the degree-eight correction has
zero rational Kummer coordinates. Thus $u_9\sigma^{-c}\in D$
and
\[
 z_{3,9}=[\sigma,u_9\sigma^{-c}].
\]
Likewise $\rho_q(u_7)=(t-t^5)\rho_q(u_5)$ for every $q$, so
$u_7u_5^{-(t-t^5)}\in D$ and
\[
 z_{5,7}=[u_5,u_7u_5^{-(t-t^5)}].
\]
Both are commutators $[g,n]$ with $g\in G$ and $n\in D\subset N$.
Every character in $X_2$ therefore vanishes on both. Correcting
by their profinite powers shows that its restriction to $E_{13}$
vanishes if and only if its restriction to $E_{12}$ vanishes.
This proves the unchanged kernel. The same two evaluations,
whose matrix is invertible modulo two, exclude every nonzero
linear combination of the two quotient characters as an
invariant restriction from $N$.
\end{proof}

\begin{remark}
Both new quotient directions are already arithmetic. Their
removal preserves the relative obstruction; it does not exhibit
a nonzero element of $\operatorname{Hom}(B,\mathbf F_2)$.
The proof uses a rational dimension upper bound together with
an explicit odd integral minor. Neither ingredient substitutes
for the other.
\end{remark}

\section{The beta-trivial leading plane in degree thirteen}
\label{sec:thirteen-normalization}

Continue with $E_{13}=C_\infty\cap F^{13}$ and the same dyadic
word filtration. We use the canonical degree-ten normalization
\[
 L_{10}=\{P_3,L_7\},\qquad [S^5TST^3]L_{10}=-141.
\]

\begin{theorem}\label{thm:thirteen-normalization}
The coordinate pair
\[
 b_{13}(g)=
 \bigl([S^7TST^4]f_g,\ [S^6TST^2ST^2]f_g\bigr)
 \qquad(g\in E_{13})
\]
is continuous and additive, with scalar conjugation weight
$\lambda_2^{13}$. Its full and arithmetic images are both
$\mathbf Z_2^2$, and its kernel is
\[
 E_{14}=C_\infty\cap F^{14}.
\]
In particular $E_{13}/E_{14}\simeq\mathbf Z_2^2$,
$G=E_{14}H$, and
\[
 B\simeq E_{14}/\overline{[G,E_{14}](E_{14}\cap H)}.
\]
The kernel of restriction of $X_2=H^1(N,\mathbf F_2)^G$ to
$E_{14}$ is still the odd-prime rational Kummer span plus
$\mathbf F_2\eta_3$. No nonzero linear combination of the two
coordinates modulo two on $E_{13}$ extends invariantly to $N$.
\end{theorem}

\begin{proof}
Put
\[
 V=\{P_3,L_{10}\},\qquad W=\{L_5,L_8\},\qquad
 \omega_1=S^7TST^4,\quad \omega_2=S^6TST^2ST^2.
\]
The polynomials $P_3,L_5,L_7,L_8$ have the explicit Lyndon
definitions in the preceding sections, and the displayed
definition of $L_{10}$ therefore specifies $V,W$ entirely by
finite integer bracket recursions. The required minor is
\begin{equation}\label{eq:thirteen-unit-minor}
 \begin{pmatrix}[\omega_1]V&[\omega_1]W\\
                 [\omega_2]V&[\omega_2]W\end{pmatrix}
 =\begin{pmatrix}-285&-56\\168&125\end{pmatrix},
 \qquad \det=-26217\in\mathbf Z_2^\times.
\end{equation}
In the order $AB,-BA,\partial_A B,-\partial_B A$, the four
contributions are
\[
\begin{array}{c|r|r|r|r|r}
(A,B;\omega)&AB&-BA&\partial_A B&-\partial_B A&\text{sum}\\\hline
(P_3,L_{10};\omega_1)&0&0&0&-285&-285\\
(P_3,L_{10};\omega_2)&0&144&-194&218&168\\
(L_5,L_8;\omega_1)&0&0&0&-56&-56\\
(L_5,L_8;\omega_2)&0&0&43&82&125
\end{array}
\]
The Lyndon recursion and the derivation coefficient rule printed
for degree twelve give an exact finite certificate for each
entry. Both selected columns again reduce to the standard basis
modulo two.

Use the arithmetic elements of the degree-twelve construction:
$\sigma$ has leading term $aP_3$,
$u_5=\Delta_3(\sigma)$ has leading term $9\tau_5L_5$, and
$u_7=\Delta_5(u_5)$ has leading term $729\tau_7L_7$, where
$a,\tau_5,\tau_7$ are odd. We have
\[
 \operatorname{in}_8(w_8)=9a\tau_5L_8,\qquad
 w_{10}=[\sigma,u_7],\qquad
 \operatorname{in}_{10}(w_{10})=729a\tau_7L_{10}.
\]
The last word is beta-trivial and has odd mixed degree-ten
coordinate because $-141$ is odd. Define
\[
 z_{3,10}=[\sigma,w_{10}],\qquad z_{5,8}=[u_5,w_8].
\]
Both are in $H\cap C_\infty\cap F^{13}=H\cap E_{13}$.
Their rational Kummer coordinates and beta multipliers vanish
on these commutators in $M$, and their leading terms are
\begin{equation}\label{eq:thirteen-arithmetic-columns}
 \operatorname{in}_{13}(z_{3,10})=729a^2\tau_7V,\qquad
 \operatorname{in}_{13}(z_{5,8})=81a\tau_5^2W.
\end{equation}
Both scalars are odd. This already supplies two independent
arithmetic beta-zero directions with an odd integral minor.

We verify that these are the entire beta-zero leading plane.
By the leading-map and base-change lemma and
\cite[Theorem~1.1, p.~2]{Willwacher},
\[
 \dim_{\mathbf Q_2}\mathfrak{grt}_1(\mathbf Q_2)_{13}
       \le d_{13}=3.
\]
The count consists of the single generator of weight thirteen
and the bracket multidegrees $(3,3,7)$ and $(3,5,5)$, each of
the latter having dimension one. Other generator lengths or
weights give no contribution.

There is a genuine arithmetic leading vector on which the
depth-one functional $[S^{12}T]$ is nonzero. Start with the
uniform odd witness
\[
 v_6=\Delta_{11}\Delta_9\Delta_7\Delta_5\Delta_3
                  \Delta_1(\sigma)\in C_5\cap H.
\]
Its moments of weights three through eleven vanish and
\[
 m_{13}(v_6)=
 \frac{2\prod_{j=0}^{5}(3^{13}-3^{2j+1})}{12!}
                        \tau_{13}(\sigma)
            \in 2^{13}\mathbf Z_2^\times.
\]
For completeness the numerator product has valuation
$3\cdot6+v_2(6!)=22$, its factor two adds one, and
$v_2(12!)=10$; the moment is odd.

Correct $v_6$ successively by beta-trivial arithmetic words.
First use $w_8$ to remove degree eight. Degree nine then vanishes
because its rational leading space is one-dimensional and its
depth-one coordinate is zero. Use $w_{10}$ to remove degree ten.
At degree eleven the depth-one coordinate is zero, so the
complete beta-zero leading-line calculation permits correction
by $w_{11}$. Finally use the two degree-twelve arithmetic words
to remove the complete degree-twelve term. Each correction is
a product with legitimate profinite powers and remains inside
$H$. The leading-lattice arguments used here apply to any full-GT
word that starts in the indicated degree: for degree eleven they
require only that its degree-eleven depth-one coefficient vanish,
not that all its higher beta coefficients vanish.

The result is an arithmetic word $z_{13}^0\in F^{13}\cap H$
whose beta multiplier is still that of $v_6$. In particular
\[
 [S^{12}T]\operatorname{in}_{13}(z_{13}^0)
       =m_{13}(v_6)\ne0.
\]
Together with the two independent beta-zero vectors
in \eqref{eq:thirteen-arithmetic-columns}, this gives three
independent vectors in $\mathfrak{grt}_1(\mathbf Q_2)_{13}$.
The rational dimension is therefore exactly three, and the
kernel of its depth-one functional is precisely
$\mathbf Q_2V+\mathbf Q_2W$.

For $g\in E_{13}$ its leading term belongs to that plane,
because $B_g=1$. The selected coefficients are integral, and
the inverse of the matrix in \eqref{eq:thirteen-unit-minor}
is over $\mathbf Z_2$. Thus the leading term belongs to
$\mathbf Z_2V+\mathbf Z_2W$. Conversely the two arithmetic
words in \eqref{eq:thirteen-arithmetic-columns}, by
profinite powers, realize every integral pair. This proves
saturation of both the full and arithmetic leading images.

Leading coordinates on $F^{13}$ are additive and transform
by the scalar $\lambda_2^{13}$. Their simultaneous vanishing
removes the entire leading term, giving the asserted kernel
$E_{14}$. Arithmetic correction now gives
$E_{13}=E_{14}(E_{13}\cap H)$, hence $G=E_{14}H$ and the same
canonical relative quotient.

Finally $w_{10},w_8\in N$, so every $G$-invariant character
on $N$ kills $[\sigma,w_{10}]$ and $[u_5,w_8]$ directly.
Correcting by those two arithmetic words shows that its
restriction kernels on $E_{13}$ and $E_{14}$ agree.
This gives the stated original kernel. Their independent
evaluations modulo two also show that no nonzero linear
combination of the two quotient characters can extend
invariantly to $N$.
\end{proof}

\section{The arithmetic leading lattice in degree fourteen}
\label{sec:degree-fourteen-lattice}
\subsection{An exact coefficient certificate for the pure degree-eleven direction}
\label{subsec:fr-pure-eleven-certificate}

The next normalization uses a specific integral vector in the
degree-eleven rational leading space. This subsection gives all
finite data needed to identify it; it does not rely on assuming
that an arbitrary infinitesimal solution lifts to profinite GT.

Use the lexicographic Lyndon basis, with $S<T$ and the longest
proper Lyndon suffix bracketing. Let $N=(N_0,\ldots,N_{185})$ be
the integer vector whose entries are listed below, row by row,
and let $L_{11}$ be the corresponding Lie polynomial of degree
eleven:
\[
\begin{array}{rrrrrr}
8448 & -42240 & 152064 & 126720 & -236544 & 183632 \\
-635824 & -253440 & 147840 & -337248 & -37328 & -1013760 \\
1632032 & -259888 & 1144704 & 354816 & -25344 & 132496 \\
-126672 & 2133568 & -1739824 & 1611760 & 7131552 & 1700885 \\
31568 & -2423936 & 3414176 & 1545536 & -2644688 & 195408 \\
-1665643 & -354816 & -59136 & -713392 & 194976 & 145696 \\
4456096 & -2655024 & -2873338 & 428624 & -1785296 & -2847760 \\
-1650127 & 283648 & -2631426 & -10156516 & -2055701 & -246992 \\
3755168 & -813519 & 5663888 & -8344398 & -11894626 & -2047685 \\
-116048 & 2904305 & -3569317 & 2634032 & 1206347 & 2650091 \\
1319275 & 253440 & -53360 & 318784 & 782288 & -238768 \\
-4135792 & -2459504 & -4508616 & -267936 & -883885 & 6968697 \\
973418 & -163168 & -313456 & 5564885 & -16000 & -2331033 \\
6357753 & -2439271 & 4203469 & 16940303 & 673697 & -188229 \\
6915145 & 3851869 & 1774080 & -411952 & -1684752 & -457131 \\
-5215638 & 71769 & -6124789 & 11170493 & 1866801 & 1075125 \\
4395714 & 6237830 & 3319616 & -2154437 & 718591 & -2801034 \\
6313939 & 13062190 & -859952 & 943076 & 1130196 & -7260032 \\
42853 & -1695184 & -375936 & -126720 & -103312 & 4819387 \\
2892779 & 11269886 & -130547 & -2565834 & 2461246 & -120507 \\
1218785 & -2554734 & 137392 & 3648115 & -2401133 & -1955433 \\
-1373195 & -1404391 & 2248757 & -7660926 & -1349408 & 2951056 \\
1650701 & 1642820 & 1621284 & 1329456 & 1895999 & -4035212 \\
-7424912 & -6067296 & 130704 & -817232 & 1755411 & 1346142 \\
1712064 & -4708057 & -5711589 & -6562098 & 5671072 & -7032843 \\
5243088 & 2583680 & 1328720 & -2287408 & 1452176 & -3228615 \\
1958704 & -2843056 & 1873664 & -2198032 & -3771728 & -51344 \\
1149616 & -372592 & 1042688 & -163840 & 10672 & 42240 \\
2783575 & -402752 & -1588080 & 1981456 & -300592 & 6888560 \\
-2296000 & -830880 & 1240240 & 3721840 & -699120 & 143616 \\
1631792 & 950752 & 321024 & 375936 & 181632 & -8448 \\
\end{array}
\]
In particular $N_0=8448$. Let
$M=(M_0,\ldots,M_{185})$ be the integral Lyndon vector of
$\{P_3,L_8\}$, already specified by the finite bracket recursion.
Its depth-one coefficient is zero and its $b_{11}$ coefficient
is eleven.

\begin{lemma}\label{lem:fr-pure-eleven-certificate}
For every full-GT element whose dyadic word starts in degree
eleven, with degree-eleven leading term $q$, one has
\[
 q=\frac{[S^{10}T]q}{8448}L_{11}
       +\frac{[S^6TST^3]q}{11}\{P_3,L_8\}.
\]
In particular $[S^{10}T]L_{11}=8448$,
$[S^6TST^3]L_{11}=0$, and $L_{11}$ is primitive as an
integral $2$-adic coefficient vector.
\end{lemma}

\begin{proof}
Here is an entirely finite necessary-equation certificate.
In logarithmic variables write
$\log f=q+r+O(13)$, with $q$ homogeneous of degree eleven
and $r$ homogeneous of degree twelve. The associated Lyndon
coordinates form $521=186+335$ columns, ordered first by degree
and then lexicographically, with indices $0,\ldots,520$.
Define
\[
 \begin{split}
 I(a)&=a(S,T)+a(T,S),\\
 C(a)&=a(S,T)+a(-S-T,S)+a(T,-S-T),\\
 J(a)&=[\varepsilon]\bigl(a(-S-T-\varepsilon[S,T],S)
                 +a(T,-S-T-\varepsilon[S,T])\bigr).
 \end{split}
\]
The full duality and multiplier-one cycle imply
\[
 I(q)=0,\qquad I(r)=0,\qquad C(q)=0,\qquad
                         J(q)+2C(r)=0.
\]
Indeed $\log((xy)^{-1})=-S-T-\tfrac12[S,T]+O(3)$,
and products of two nonconstant terms of $f$ start in degree
twenty-two. Thus the last equation is twice the degree-twelve
cycle coefficient, with no quadratic-in-$f$ contribution.

Form the integer matrix $A_0$ by extracting every nonzero
associative word coefficient of these four expressions, in
the displayed order, with words lexicographic within each
degree. This gives $12092$ rows and $521$ columns.
To these append seven necessary leading-pentagon rows.
For free Lie generators $X_1,\ldots,X_4$, define
\[
 \delta_{ij}(X_i)=[X_j,X_i],\qquad
 \delta_{ij}(X_j)=[X_i,X_j],\qquad
 \delta_{ij}(X_k)=0\quad(k\notin\{i,j\}).
\]
These derivations satisfy the infinitesimal braid relations.
For a Lie polynomial $a$ write
\[
\begin{split}
\mathcal P(a)={}&a(\delta_{12},\delta_{23}+\delta_{24})
 +a(\delta_{13}+\delta_{23},\delta_{34})
 -a(\delta_{23},\delta_{34})\\
 &-a(\delta_{12}+\delta_{13},\delta_{24}+\delta_{34})
 -a(\delta_{12},\delta_{23}).
\end{split}
\]
Take the coefficients of the following words in
$\mathcal P(q)(X_1)$, where a digit denotes its corresponding
free generator:
\[
\begin{gathered}
111111123244,\quad111112323244,\quad111123232344,\quad
111123234234,\\
111232342344,\quad112323423244,\quad121232323444.
\end{gathered}
\]
They have zero entries in the degree-twelve columns. These are
genuine necessary equations because the leading term of a
full-GT word satisfies the infinitesimal pentagon; the
derivation representation only evaluates that equation.
This fully defines a $12099$ by $521$ integer matrix $A$.

The following minor and identity are an exact integer/rational
certificate. Put
\[
 C=\{0,1,\ldots,408\}\setminus\{180,185\}.
\]
The row-index set $R$, using zero-based indices of the matrix
just defined, is the following list read row by row:
{\scriptsize
\[
\begin{array}{rrrrrrrrrrrr}
0 & 2 & 4 & 6 & 8 & 10 & 12 & 14 & 16 & 18 & 20 & 22 \\
24 & 26 & 28 & 30 & 32 & 34 & 36 & 38 & 40 & 42 & 44 & 46 \\
48 & 50 & 52 & 54 & 56 & 58 & 60 & 66 & 68 & 70 & 72 & 74 \\
76 & 78 & 80 & 82 & 84 & 86 & 88 & 90 & 92 & 98 & 100 & 102 \\
104 & 106 & 108 & 112 & 114 & 116 & 118 & 120 & 136 & 138 & 140 & 146 \\
148 & 150 & 152 & 154 & 156 & 162 & 164 & 166 & 168 & 170 & 172 & 178 \\
180 & 184 & 198 & 200 & 202 & 204 & 210 & 212 & 216 & 218 & 228 & 232 \\
292 & 298 & 300 & 306 & 308 & 330 & 332 & 338 & 340 & 2046 & 2048 & 2050 \\
2052 & 2054 & 2056 & 2058 & 2060 & 2062 & 2064 & 2066 & 2068 & 2070 & 2072 & 2074 \\
2076 & 2078 & 2080 & 2082 & 2084 & 2086 & 2088 & 2090 & 2092 & 2094 & 2096 & 2098 \\
2100 & 2102 & 2104 & 2106 & 2111 & 2113 & 2115 & 2117 & 2119 & 2121 & 2123 & 2125 \\
2127 & 2129 & 2131 & 2133 & 2135 & 2137 & 2139 & 2141 & 2143 & 2145 & 2147 & 2149 \\
2151 & 2153 & 2155 & 2157 & 2159 & 2161 & 2163 & 2165 & 2167 & 2175 & 2177 & 2179 \\
2181 & 2183 & 2185 & 2186 & 2188 & 2190 & 2192 & 2194 & 2196 & 2198 & 2204 & 2206 \\
2208 & 2210 & 2212 & 2214 & 2217 & 2219 & 2221 & 2223 & 2225 & 2227 & 2229 & 2237 \\
2239 & 2241 & 2243 & 2245 & 2248 & 2250 & 2252 & 2254 & 2256 & 2258 & 2260 & 2266 \\
2268 & 2270 & 2272 & 2274 & 2276 & 2279 & 2281 & 2312 & 2314 & 2316 & 2318 & 2320 \\
2322 & 2328 & 2330 & 2332 & 2334 & 2336 & 2338 & 2344 & 2346 & 2349 & 2351 & 2363 \\
2365 & 2367 & 2369 & 2375 & 2377 & 2380 & 2382 & 2384 & 2390 & 2392 & 2394 & 2395 \\
2397 & 2399 & 2405 & 2431 & 2437 & 2439 & 2442 & 2444 & 2446 & 2454 & 2457 & 2459 \\
2461 & 2467 & 2469 & 2488 & 2490 & 2500 & 2615 & 2617 & 2623 & 2625 & 2639 & 2645 \\
2647 & 2686 & 2709 & 2715 & 2748 & 2841 & 6016 & 6018 & 6022 & 6028 & 6030 & 6036 \\
6040 & 6046 & 6052 & 6056 & 6058 & 6064 & 6066 & 6084 & 6088 & 6090 & 6096 & 6098 \\
6102 & 6114 & 6118 & 6152 & 6154 & 6160 & 6162 & 6166 & 6176 & 6178 & 6182 & 6214 \\
6306 & 8060 & 8062 & 8064 & 8066 & 8068 & 8070 & 8072 & 8076 & 8078 & 8080 & 8082 \\
8084 & 8086 & 8088 & 8092 & 8094 & 8096 & 8098 & 8100 & 8102 & 8104 & 8106 & 8108 \\
8110 & 8112 & 8114 & 8116 & 8118 & 8126 & 8128 & 8130 & 8132 & 8134 & 8136 & 8138 \\
8140 & 8142 & 8144 & 8146 & 8148 & 8150 & 8152 & 8155 & 8157 & 8159 & 8161 & 8163 \\
8165 & 8167 & 8169 & 8171 & 8175 & 8191 & 8195 & 8197 & 8199 & 8202 & 8204 & 8206 \\
8208 & 8210 & 8212 & 8214 & 8220 & 8222 & 8226 & 8228 & 8230 & 8234 & 8236 & 8238 \\
8240 & 8242 & 8244 & 8246 & 8254 & 8256 & 8258 & 8260 & 8262 & 8266 & 8268 & 8330 \\
8332 & 8334 & 8336 & 8338 & 8346 & 8348 & 8352 & 8354 & 8356 & 8364 & 8384 & 8386 \\
8388 & 8396 & 8400 & 8402 & 8404 & 8410 & 8412 & 8415 & 8417 & 8419 & 8451 & 8638 \\
8648 & 8710 & 8734 & 8774 & 12092 & 12093 & 12094 & 12095 & 12096 & 12097 & 12098 &  \\
\end{array}
\]}
It has $407$ entries. For $B=A_{R,C}$ the determinant is
\[
 \begin{aligned}
 \det B={}&570662177953320090405597978960021340179106\cdot10^{36}\\
 &+875700910966765341012729477988352000.
 \end{aligned}
\]
The $b_{11}$ functional on the degree-eleven Lyndon coordinates
is
\[
 \ell=e_{11}-4e_7,
\]
extended by zero on the degree-twelve columns. Define the
$186$ by $521$ integer matrix $T$ by its rows
\[
 T_j=92928e_j-11N_j e_0-8448M_j\ell
       \qquad(0\le j<186).
\]
Direct integer matrix multiplication gives the certificate
\begin{equation}\label{eq:fr-pure-eleven-row-certificate}
 T_{*,C}\operatorname{adj}(B)A_{R,*}=(\det B)\,T.
\end{equation}
All entries here are determined by the printed vectors,
bracket recursions, word lists, and row/column sets. Thus
\eqref{eq:fr-pure-eleven-row-certificate} is a finite exact
coefficient identity, with no floating-point or modular
inference.

For the coefficient vector $(q,r)$ arising from a full-GT element, one has
$A(q,r)=0$. Since $\det B\ne0$, the certificate gives
$T(q,r)=0$. Dividing its $j$th equality by $92928=8448\cdot11$
gives precisely the stated leading-term formula.
The displayed vector has $N_{11}-4N_7=0$ and includes odd
entries, so its two claimed normalizations and primitivity
also follow directly.
\end{proof}

\subsection{The complete leading layer in degree fourteen}
\label{subsec:fr-fourteen-normalization}

Use the pure integral $L_{11}$ of
Lemma~\ref{lem:fr-pure-eleven-certificate}, normalized by
$[S^{10}T]L_{11}=8448$ and $[S^6TST^3]L_{11}=0$.
We continue to write $E_n=C_\infty\cap F^n$ for the dyadic
word filtration.

\begin{theorem}\label{thm:fr-fourteen-normalization}
On $E_{14}$ the coordinate map
\[
 b_{14}(g)=\bigl(
 [S^8TST^4]f_g,\ 
 [S^7TST^5]f_g,\ 
 [S^6TST^2ST^3]f_g
 \bigr)
\]
is continuous and additive with scalar conjugation weight
$\lambda_2^{14}$, and
\[
 b_{14}(E_{14})=b_{14}(E_{14}\cap H)=\mathbf Z_2^3.
\]
Its kernel is $E_{15}=C_\infty\cap F^{15}$. Consequently
\[
 E_{14}/E_{15}\simeq\mathbf Z_2^3,\qquad
 G=E_{15}H,\qquad
 B\simeq E_{15}/\overline{[G,E_{15}](E_{15}\cap H)}.
\]
The original $X_2=H^1(N,\mathbf F_2)^G$ restriction kernel
to $E_{15}$ remains the odd-prime rational Kummer span plus
$\mathbf F_2\eta_3$. No nonzero linear combination of the three
coordinates modulo two on $E_{14}$ extends invariantly to $N$.
\end{theorem}

\begin{proof}
First construct the required primitive arithmetic pure
degree-eleven input without dividing a group element.
Starting with the arithmetic $u_9$ constructed in degree
twelve, form $\Delta_9(u_9)$. It starts in degree ten.
Use a profinite power of the beta-trivial arithmetic $w_{10}$
to remove its degree-ten term, obtaining a word starting in
degree eleven. Then use a profinite power of $w_{11}$ to make
its coefficient $b_{11}=[S^6TST^3]$ zero. Denote the result
by $u_{11}\in H\cap F^{11}$. The two corrections preserve
every beta moment and every rational Kummer coordinate.

The lower positive even moments already vanish. The exact
weight laws therefore give
\[
 m_{11}(u_{11})=
 \frac{(3^{11}-3^9)(3^{11}-3^7)(3^{11}-3^5)(3^{11}-3^3)}
             {1814400}\,\tau_{11}(\sigma)
 =475764757947648\,\tau_{11}(\sigma).
\]
The finite coefficient certificate now identifies its entire
leading term:
\begin{equation}\label{eq:fr-fourteen-pure-eleven-input}
 \operatorname{in}_{11}(u_{11})=\epsilon_{11}L_{11},
 \qquad
 \epsilon_{11}=\frac{619485361911}{11}\tau_{11}(\sigma)
                  \in\mathbf Z_2^\times.
\end{equation}
The denominator eleven is a unit in $\mathbf Z_2$; the only
operations on group elements were the stated profinite powers.

Define three degree-fourteen integer Lie polynomials
\[
 V_1=\{P_3,L_{11}\},\qquad
 V_2=\{L_5,L_9\},\qquad
 V_3=\{P_3,\{P_3,L_8\}\}.
\]
At the three words
\[
 \omega_1=S^8TST^4,\qquad
 \omega_2=S^7TST^5,\qquad
 \omega_3=S^6TST^2ST^3
\]
their coefficient matrix is
\begin{equation}\label{eq:fr-fourteen-unit-minor}
 \bigl([\omega_i]V_j\bigr)_{i,j}
 =
 \begin{pmatrix}
 -563792&3526&15\\
 1143317&494&-24\\
 -4406053&62987&71
 \end{pmatrix},
 \qquad \det=327441420801\in\mathbf Z_2^\times.
\end{equation}
The following table supplies the four Ihara contributions
$AB,-BA,\partial_A B,-\partial_B A$ to every entry:
\[
\begin{array}{c|r|r|r|r}
(A,B;\omega)&AB&-BA&\partial_A B&-\partial_B A\\\hline
(P_3,L_{11};\omega_1)&0&0&-760320&196528\\
(P_3,L_{11};\omega_2)&0&0&1216512&-73195\\
(P_3,L_{11};\omega_3)&0&0&-3419688&-986365\\
(L_5,L_9;\omega_1)&0&288&-38688&41926\\
(L_5,L_9;\omega_2)&0&0&57696&-57202\\
(L_5,L_9;\omega_3)&0&-21312&4003&80296\\
(P_3,\{P_3,L_8\};\omega_1)&0&0&0&15\\
(P_3,\{P_3,L_8\};\omega_2)&0&0&0&-24\\
(P_3,\{P_3,L_8\};\omega_3)&0&0&28&43
\end{array}
\]
The explicit Lyndon vectors and derivation coefficient rule
already printed, together with the complete finite certificate
for $L_{11}$, specify this calculation exactly.

The three corresponding genuine arithmetic words are
\[
 z_1=[\sigma,u_{11}],\qquad
 z_2=[u_5,u_9],\qquad
 z_3=[\sigma,w_{11}].
\]
They belong to $E_{14}\cap H$: their leading degrees are at
least fourteen, and commutators in $M$ have trivial beta and
rational Kummer coordinates. Their leading terms are
\[
 \operatorname{in}_{14}(z_1)=a\epsilon_{11}V_1,\qquad
 \operatorname{in}_{14}(z_2)=20726199\,\tau_5\tau_9V_2,
 \qquad
 \operatorname{in}_{14}(z_3)=9a^3\tau_5V_3.
\]
Here $a=m_3(\sigma)$ and $\tau_j=\tau_j(\sigma)$; every
displayed scalar is odd. The second scalar is
$9\cdot2302911$, and the last uses
$\operatorname{in}_{11}(w_{11})=9a^2\tau_5\{P_3,L_8\}$.

By Lemma~\ref{lem:full-gt-leading-term} and
\cite[Theorem~1.1, p.~2]{Willwacher}, the full rational leading
space in degree fourteen has dimension at most $d_{14}=3$.
Indeed its free Lie types are $(3,11)$, $(5,9)$, and
$(3,3,3,5)$; the apparent type $(7,7)$ vanishes. The last
listed type has one-dimensional span, represented by the
triple adjoint of the weight-three generator on the
weight-five generator.
The three genuine arithmetic leading vectors above are
independent by \eqref{eq:fr-fourteen-unit-minor}. They
therefore span the entire rational leading space.

For an integral leading polynomial in that rational space,
its three selected Magnus coefficients are integral.
The inverse of the matrix
\eqref{eq:fr-fourteen-unit-minor} is over $\mathbf Z_2$,
so its coefficients on $V_1,V_2,V_3$ all belong to
$\mathbf Z_2$. Conversely the three arithmetic words, with
their odd scalar factors, realize every integral triple
by profinite powers and products. This proves exact
integral saturation of both images, not merely equality
after tensoring with $\mathbf Q_2$.

Leading terms on $F^{14}$ add and transform by
$\lambda_2^{14}$. The three selected coefficients vanish
if and only if the entire degree-fourteen term vanishes.
Thus the kernel is $E_{15}$ and the coordinate map has
the stated properties. Arithmetic correction gives
$E_{14}=E_{15}(E_{14}\cap H)$, hence $G=E_{15}H$ and
the unchanged canonical relative quotient.

Finally all three correcting words are killed by each
$\chi\in X_2$. For the first, put
\[
 c_{3,11}=(t-t^9)(t-t^7)(t-t^5)(t-t^3)
                    \in\widehat{\mathbf Z}.
\]
The full profinite weight-one Kummer laws and the
beta-trivial corrections give
$\rho_q(u_{11})=c_{3,11}\rho_q(\sigma)$ for every
rational prime, including two. Hence
$u_{11}\sigma^{-c_{3,11}}\in D$ and
\[
 z_1=[\sigma,u_{11}\sigma^{-c_{3,11}}].
\]
For the second put $c_{5,9}=(t-t^7)(t-t^5)$.
Then $\rho_q(u_9)=c_{5,9}\rho_q(u_5)$ for all such
coordinates, so $u_9u_5^{-c_{5,9}}\in D$ and
\[
 z_2=[u_5,u_9u_5^{-c_{5,9}}].
\]
For the third, $w_{11}\in N$ already. Each is consequently
a commutator of an element of $G$ with an element of $N$.
Invariance gives $\chi(z_i)=0$ for all three. Arithmetic
correction preserves the restriction kernel from
$E_{14}$ to $E_{15}$, yielding the displayed original
kernel and excluding all nonzero invariant extensions
of the three quotient characters.
\end{proof}

\section{A concrete quadratic character from the eleventh multiplication formula}
\label{sec:odd-anderson-characters}

This section constructs the explicit continuous quadratic character
$\Theta_{17,11}$, which vanishes on the arithmetic subgroup.  Its
nonvanishing on the full
profinite Grothendieck--Teichm\"uller group is not established.  A separate
formal example shows why beta factorization, the unitary relation, and
the dyadic multiplication formulas do not by themselves eliminate these
characters.  The example is not asserted to lift to a pentagon element.

Let $G_0$ be the profinite gentle group and let
$G=\widehat{\mathrm{GT}}\subset G_0$.  Write
$F=\widehat F(x,y)$ and $z=(xy)^{-1}$.  We use the convention
\[
 \phi_g(x)=x^{\lambda_g},\qquad
 \phi_g(y)=f_g^{-1}y^{\lambda_g}f_g,
 \qquad f_g\in F'.
\]
Let $B_g(X,Y)$ be the normalized Alexander coefficient, so that the
semilinear coefficient on $F'/F''$ is
\[
 A_g(X,Y)=S_{\lambda_g}(X)S_{\lambda_g}(Y)B_g(X,Y),
 \qquad S_a(T)=\frac{T^a-1}{T-1}.
\]
Composition gives
\begin{equation}
 B_{gh}(X,Y)=B_g(X,Y)
                 B_h(X^{\lambda_g},Y^{\lambda_g}).
 \label{eq:odd-beta-cocycle}
\end{equation}
These are the conventions of \cite[\S2.3, p.~101]{Furusho2003}.

\subsection{The boundary characters in the unitary relation}

\begin{lemma}\label{lem:all-finite-unitarity}
Suppose that $\lambda_g=1$.  Let $p$ be an odd prime and let
$a,b\in\overline{\mathbb F}_p^\times$ have finite orders prime to $p$.
Then
\[
 B_g(a,b)B_g(a^{-1},b^{-1})=1.
\]
In fact $B_g(a,b)=1$ if $a=1$, $b=1$, or $ab=1$.
\end{lemma}

\begin{proof}
Choose $m$ prime to $p$ divisible by the orders of $a$ and $b$, and
put $U_m=\ker(F\to(\mathbb Z/m)^2)$.  Since $f_g\in F'$, the
automorphism $\phi_g$ induces the identity on $F/\gamma_3F$.  Choose
$c\in F$ such that $\phi_g(z)=c^{-1}zc$, and write
$[c]=(u,v)\in F^{\mathrm{ab}}=\widehat{\mathbb Z}^2$.  In the
class-two quotient the commutator subgroup is freely generated over
$\widehat{\mathbb Z}$ by $[x,y]$.  The equality $c^{-1}zc=z$ in this
quotient implies $u=v$.

The cover with fundamental group $U_m$ has deck group $(\mathbb Z/m)^2$.
Its boundary subrepresentation in $H_1(U_m,\mathbb F_p)$ is the sum
of the three permutation representations induced from the cyclic
subgroups generated by the images of $x$, $y$, and $z$, with one
trivial relation among the boundary classes.  On the boundaries above
$x$ the conjugator is $1$; above $y$ it is $f_g$, whose deck image is
$1$; above $z$ it is $c$, whose deck image is diagonal.  Thus $\phi_g$
fixes the boundary representation pointwise.  More explicitly, a
character $(a,b)$ occurring above $z$ satisfies $ab=1$, and its
conjugator scalar is $a^u b^v=(ab)^u=1$.  The corresponding assertions
above $x$ and $y$ are immediate.  Every nontrivial deck character
occurs in $H_1(U_m,\overline{\mathbb F}_p)$ with multiplicity one.
The coefficient $B_g(a,b)$ is its scalar action.  This proves the
last assertion for all nontrivial boundary characters.  At the trivial
character $B_g(1,1)=1$, since the class-two commutator is fixed.

For completeness, the complement of the boundary subrepresentation
carries the usual nondegenerate alternating form.  One can see its
preservation without a topological realization of $\phi_g$.  A surface
presentation of $U_m$ has the form
\[
 \prod_{j=1}^{(m-1)(m-2)/2}[a_j,b_j]
       \prod_{i=1}^{3m}c_i=1.
\]
In the class-two exponent-$p$ quotient, project the degree-two
commutators to the exterior square of homology modulo the boundary
classes.  The relation yields the bivector
$\sum_j[a_j]\wedge[b_j]$.  Conjugating a boundary word changes its
degree-two coordinate by a wedge containing its boundary class, so
this change vanishes after projection.  Positive boundary preservation
therefore fixes the bivector and the associated alternating form.

The deck action preserves the form, so the character spaces at
$(a,b)$ and $(a^{-1},b^{-1})$ pair perfectly whenever
$a,b,ab\ne1$.  Preservation of the form makes the product of their
scalar actions equal to one.  The boundary calculation proves the
same identity in the remaining cases.
\end{proof}

\subsection{An explicit family of arithmetic-trivial characters}

\begin{theorem}\label{thm:odd-anderson-character}
Let $p\equiv5\pmod{12}$ be prime, and let $n\geq1$ satisfy
$\gcd(n,3p)=1$.  Choose a primitive cube root
$\omega\in\mathbb F_{p^2}$.  The expression
\begin{equation}
 E_{p,n}(g)=
 \frac{\displaystyle\prod_{k=0}^{n-1}B_g(\omega,\omega^k)}
      {\displaystyle\prod_{\xi^n=1}B_g(\xi,\omega)}
 \label{eq:odd-anderson-defect}
\end{equation}
belongs to the norm-one subgroup of $\mathbb F_{p^2}^\times$.
Here all evaluations are reduced modulo $p$, and the denominator
ranges over all $n$th roots in $\overline{\mathbb F}_p$.
The rule
\begin{equation}
 (-1)^{\Theta_{p,n}(g)}=E_{p,n}(g)^{(p+1)/2}
 \label{eq:odd-anderson-sign}
\end{equation}
defines a continuous homomorphism $G_0\to\mathbb F_2$, independent
of the choice between $\omega$ and $\omega^2$.  It vanishes on
$G_{\mathbb Q}$.

In particular its restriction to $G$ is an explicit candidate for a
nonzero element of
$\operatorname{Hom}_{\mathrm{cts}}(G/\langle\!\langle
G_{\mathbb Q}\rangle\!\rangle,\mathbb F_2)$.
The theorem does not assert that this restriction is nonzero.
\end{theorem}

\begin{proof}
The evaluations are nonzero because they are coefficients of
invertible semilinear actions on the appropriate character spaces.
The $p^2$-power Frobenius fixes $\omega$ and permutes the $n$th
roots.  Hence the expression lies in $\mathbb F_{p^2}^\times$.
Equation~\eqref{eq:odd-beta-cocycle} and reindexing the denominator
give
\begin{equation}
 E_{p,n}(gh)=E_{p,n}(g)\,
             \sigma_{\lambda_g}\bigl(E_{p,n}(h)\bigr),
 \label{eq:odd-anderson-crossed}
\end{equation}
where $\sigma_\lambda$ is the identity for $\lambda\equiv1\pmod3$
and the $p$-power Frobenius for $\lambda\equiv-1\pmod3$.
Indeed $\lambda$ permutes all $n$th roots, and its only remaining
effect in the displayed expression is its action on $\omega$.

If $\lambda_g=1$, Lemma~\ref{lem:all-finite-unitarity} gives
$E_{p,n}(g)^p=E_{p,n}(g)^{-1}$.  In the denominator one reindexes
the $p$th powers of the $n$th roots as their inverses; in the numerator
$\omega^p=\omega^{-1}$.  For arbitrary $g$, choose
$h\in G_{\mathbb Q}$ with $\lambda_h=\lambda_g$ and write $g=uh$,
where $\lambda_u=1$.  The arithmetic Anderson multiplication formula
\cite[Corollary~8.6.3]{Anderson1989}, in the beta form stated explicitly
in \cite[Remark~2.4, p.~101]{Furusho2003}, gives
\begin{equation}
 E_{p,n}(h)=\omega^{n\Psi_n(h)}\in\mu_3.
 \label{eq:odd-anderson-arithmetic}
\end{equation}
Its exact sign convention is immaterial here; in the convention above
$\Psi_n=-\rho_n$.  Since $p\equiv-1\pmod3$, the group $\mu_3$ also
has norm one.  Equation~\eqref{eq:odd-anderson-crossed} now proves
that $E_{p,n}(g)$ has norm one.

The norm-one subgroup is cyclic of order $p+1$, so
\eqref{eq:odd-anderson-sign} is a sign.  Frobenius acts on that
subgroup by inversion and therefore fixes its quadratic sign.
Consequently \eqref{eq:odd-anderson-crossed} makes $\Theta_{p,n}$
an ordinary homomorphism.  It is continuous because it uses finitely
many finite character evaluations.  Since $3$ divides $(p+1)/2$,
\eqref{eq:odd-anderson-arithmetic} makes this homomorphism zero on
$G_{\mathbb Q}$.  Replacing $\omega$ by $\omega^2$ applies Frobenius
and leaves the sign unchanged.
\end{proof}

For a pentagon element, Ihara's beta factorization gives
$B_g(s,t)=\Gamma_g(s)\Gamma_g(t)/\Gamma_g(s+t)$ in additive
notation.  Consequently
\begin{equation}
 E_{p,n}(g)=
 \frac{\displaystyle\prod_{c\in\frac1n\mathbb Z/\mathbb Z}
                    \Gamma_g(\tfrac13+c)}
      {\displaystyle\Gamma_g(\tfrac n3)
          \prod_{c\in\frac1n\mathbb Z/\mathbb Z}\Gamma_g(c)}.
 \label{eq:odd-anderson-gamma}
\end{equation}
This expression is independent of the character ambiguity in
$\Gamma_g$.  If the Anderson relation $(A_n)$ holds for $g$, then
$\Theta_{p,n}(g)=0$.  Thus the known dyadic multiplication formulas
eliminate the cases $n=2^a$.  Enriquez's theorem that
$\mathrm{GTK}=\widehat{\mathrm{GT}}$ also eliminates every proposed
character obtained from a multiplicative defect of the integral
Kummer coordinates: his Remark~8.6 gives
$\rho_{ab}=\rho_a+\rho_b$ before any reduction or division.
The exact reference is \cite[Theorem~8.1 and Remark~8.6]{Enriquez2007}.
It does not assert the equality
$\mathrm{GTA}=\widehat{\mathrm{GT}}$; the relation $(K_n)$ is recorded
as implying the logarithmic derivative of $(A_n)$ in
\cite[p.~102]{Furusho2003}.

\begin{corollary}\label{cor:concrete-radical-anderson-character}
Let $H=G_{\mathbb Q}$ and let $N_{\mathrm{rad}}$ be the simultaneous
kernel in $G$ of the multiplier and the rational-prime Kummer cocycles.
Put
\[
 L_{\mathrm{rad}}=\mathbb Q(\mu_\infty,q^{1/m}:
                          q\text{ prime},\ m\geq1),
 \qquad K_{\mathrm{rad}}=G_{L_{\mathrm{rad}}}.
\]
Assume the simultaneous normalization
$G=N_{\mathrm{rad}}H$, with
$N_{\mathrm{rad}}\cap H=K_{\mathrm{rad}}$, and identify the relative
abelian quotient as
\[
 B=N_{\mathrm{rad}}/
       \overline{[G,N_{\mathrm{rad}}]K_{\mathrm{rad}}}.
\]
For a primitive cube root $\omega\in\mathbb F_{289}$ the formula
\begin{equation}
 (-1)^{\psi_{17,11}(u)}=
    \left(\prod_{\substack{\xi^{11}=1\\\xi\ne1}}
                        B_u(\xi,\omega)\right)^9,
 \qquad u\in N_{\mathrm{rad}},
 \label{eq:concrete-radical-anderson-sign}
\end{equation}
defines a continuous $G$-invariant quadratic character vanishing on
$K_{\mathrm{rad}}$.  It therefore gives a precisely specified element
of $\operatorname{Hom}_{\mathrm{cts}}(B,\mathbb F_2)$.
Its unique extension to $G$ vanishing on $H$ is $\Theta_{17,11}$.
In particular $\psi_{17,11}$ is nonzero if and only if
$\Theta_{17,11}|_G$ is nonzero; neither nonvanishing assertion is
proved here.
\end{corollary}

\begin{proof}
Take $p=17$, $n=11$ in Theorem~\ref{thm:odd-anderson-character}.
For $u\in N_{\mathrm{rad}}$ the multiplier is one.
In the numerator of \eqref{eq:odd-anderson-defect}, four values of
$k$ are $1$ modulo $3$, four are $0$ modulo $3$, and three are $2$
modulo $3$.  Lemma~\ref{lem:all-finite-unitarity} gives
\[
 \prod_{k=0}^{10}B_u(\omega,\omega^k)=B_u(\omega,\omega)^4,
 \qquad B_u(\omega,\omega)^{18}=1.
\]
Its ninth-power contribution is therefore one.  The denominator's
$\xi=1$ term is also one.  Since inversion fixes a sign, the
restriction of \eqref{eq:odd-anderson-sign} is exactly
\eqref{eq:concrete-radical-anderson-sign}.

Being a restriction of an ordinary character of $G$, this map is
$G$-invariant; being zero on $H$, it is zero on $K_{\mathrm{rad}}$.
In fact the arithmetic formula gives $E_{17,11}(u)=1$ on
$K_{\mathrm{rad}}$, since $\Psi_{11}(u)=0$ there.
It consequently factors through the displayed quotient $B$.
The factorization $G=N_{\mathrm{rad}}H$ proves both uniqueness
of the extension and the equivalence of the two nonvanishing claims.
\end{proof}

There is also a finite-dimensional interpretation of the product in
\eqref{eq:concrete-radical-anderson-sign}.  Take the degree-$33$ cover
with deck group $C_{11}\times C_3$, given algebraically by
$u^{11}+v^3=1$ and the map $t=u^{11}$.  There are three boundary
components above $t=0$, eleven above $t=1$, and one above infinity.
Euler characteristic therefore gives genus ten for its compactification.
Its homology has the ten character lines $(\xi,\omega)$ with
$\xi$ a nontrivial $11$th root, together with their inverse lines.
The individual evaluations may lie in $\mathbb F_{17^{10}}$;
$\mu_{11}$ is not contained in $\mathbb F_{289}$.
The sum descends to a ten-dimensional space over $\mathbb F_{289}$,
since the $17^2$-power Frobenius permutes these ten characters.  It is
the whole $\omega$-eigenspace of the cubic deck action on the compact
homology.  The product is the determinant of an element of
$N_{\mathrm{rad}}$ on this space, and
\eqref{eq:concrete-radical-anderson-sign} is its quadratic norm-one
sign.  Thus the proposed test involves one finite homology
representation of a genus-ten curve.

\subsection{A formal sign ambiguity which survives duplication}

\begin{proposition}\label{prop:formal-odd-sign}
There exists a Frobenius-equivariant sign-valued formal gamma function
at the prime $17$ whose beta coboundary satisfies beta associativity,
symmetry, the unitary relation, and the duplication formula, but for
which the expression in \eqref{eq:odd-anderson-gamma} with $n=11$
equals $-1$.  It may be taken constant in the pro-$17$ directions,
so all its logarithmic derivatives vanish.

This is a non-implication statement for the listed formal identities.
It does not provide a solution of the full nonabelian relations
$(K_n)$ or of the pentagon equation.
\end{proposition}

\begin{proof}
Let $T=(\mathbb Q/\mathbb Z)_{17'}$ be the subgroup of elements of
order prime to $17$.  In $(\mathbb Z/33)^\times$ put
\[
 S=\langle2\rangle
   =\{1,2,4,8,16,32,31,29,25,17\}.
\]
Define $\gamma:T\to\mathbb F_2$ by
\[
 \gamma(s)=
 \begin{cases}
  1,&s=a/33\text{ for }a\in S,\\
  0,&\text{otherwise},
 \end{cases}
 \qquad \Gamma(s)=(-1)^{\gamma(s)}.
\]
For an arbitrary rational torsion argument $s$, the definition means
$\gamma(s)=\gamma(s_{17'})$, where $s_{17'}$ is its prime-to-$17$
component.  Thus the value is constant in the $17$-primary variable;
it is not declared to be $1$ merely because the denominator is
divisible by $17$.  In particular $\gamma$ is zero whenever the
prime-to-$17$ component has even order.
The set $S$ is stable under $2$, $-1=2^5$, and $17=2^9$.
Thus $\Gamma(-s)=\Gamma(s)$ and $\Gamma(17s)=\Gamma(s)$.

We check the duplication relation
\begin{equation}
 \gamma(s)+\gamma(s+\tfrac12)=\gamma(2s)+\gamma(\tfrac12).
 \label{eq:formal-sign-duplication}
\end{equation}
The last term is zero.  If $s$ has odd order, then $s+1/2$ has
even order and the assertion is the invariance under doubling on
odd torsion.  If the $2$-primary component of $s$ has order exactly
two, put $v=s+1/2$.  Then $v$ has odd order, $2s=2v$, and the same
invariance proves the identity.  If the $2$-primary component has
order at least four, all three possibly nonzero arguments on the
left and right have even order, so all terms vanish.  This proves
\eqref{eq:formal-sign-duplication}, or equivalently
\[
 \Gamma(s)\Gamma(s+\tfrac12)=\Gamma(2s)\Gamma(\tfrac12).
\]
Iteration gives all dyadic multiplication identities with zero Kummer
correction.

Put $B(s,t)=\Gamma(s)\Gamma(t)/\Gamma(s+t)$.  It is symmetric,
satisfies beta associativity, is one on the axes and on $s+t=0$,
and satisfies $B(s,t)B(-s,-t)=1$.  These assertions follow directly
from the coboundary definition and the fact that all values are signs.
For the prime-to-$17$ spectral components, the Frobenius invariance
allows the assignment to descend to $\mathbb Z_{17}$ coefficients.
Make it constant on each pro-$17$ formal neighborhood.  Equivalently,
use the product decomposition of
$\mathbb Z_{17}[[\widehat{\mathbb Z}]]$ into the unramified formal
power-series rings indexed by Frobenius orbits of prime-to-$17$
characters.  The specified constants define a unit in this completed
group algebra.  The beta coboundary is a corresponding unit in the
two-variable algebra, and its logarithmic derivatives are zero.

Finally $\Gamma(11/3)=\Gamma(2/3)=1$, and $\Gamma(c)=1$ for
every $c\in\frac1{11}\mathbb Z/\mathbb Z$.  The numerators of the
points $1/3+k/11$, expressed with denominator $33$, are $11+3k$.
Their intersection with $S$ is
\[
 \{2,8,17,29,32\},
\]
which has five elements.  Therefore the numerator of
\eqref{eq:odd-anderson-gamma} is $(-1)^5=-1$ and its denominator
is $1$.  At $p=17$ the quadratic exponent is $(17+1)/2=9$, so
the associated formal quadratic sign is also $-1$.
\end{proof}

The remaining question is a full-group realization or elimination
problem.  Neither the finite sign computation nor the fact that
logarithmic derivatives forget this sign permits one to infer an
element of $\widehat{\mathrm{GT}}$ with nonzero
$\Theta_{17,11}$.  Conversely, eliminating these candidates requires
an additional consequence of the full relations, rather than only
the four formal identities in Proposition~\ref{prop:formal-odd-sign}.

\subsection{Two limitations of a direct eleventh-Kummer comparison}

Put $D=\ker\lambda\cap\bigcap_q\ker\rho_q^I$ in the full
profinite Grothendieck--Teichm\"uller group.  By \cite[Theorem~8.1 and Remark~8.6]{Enriquez2007}, the exact
Kummer relation, with zero Kummer coordinate on $D$, is
\begin{equation}\label{eq:eleventh-local-system-K}
d_n(f_g)=f_g\qquad(g\in D).
\end{equation}
Here $H_n\leq\widehat F(x,y)$ is the kernel of $x\mapsto1,y\mapsto0$
in $\mathbb Z/n$, freely generated by
\[
 X=x^n,\qquad Y_j=x^{-j}yx^j\quad(0\leq j<n),
\]
and $d_n:H_n\to\widehat F(x,y)$ sends $X\mapsto x$, $Y_0\mapsto y$
and $Y_j\mapsto1$ for $j\ne0$.  We apply the full nonabelian relation \eqref{eq:eleventh-local-system-K}.

\begin{lemma}\label{lem:kn-different-local-systems}
For $n=11$ and a primitive cube root $\omega$ in
$\mathbb F_{17^2}$, consider characters over $\overline{\mathbb F}_{17}$.
The character of $H_{11}$ obtained by pulling
back $x\mapsto1,y\mapsto\omega$ along $d_{11}$ is different from
the restriction of every character
$x\mapsto\xi,y\mapsto\omega$ with $\xi^{11}=1$.
Nor is either kind a deck translate of the other.
\end{lemma}
\begin{proof}
The pullback along $d_{11}$ has values
\[
 X\mapsto1,\qquad Y_0\mapsto\omega,\qquad
 Y_j\mapsto1\ (j\ne0).
\]
Each character of the second kind has values
\[
 X\mapsto\xi^{11}=1,\qquad Y_j\mapsto\omega
 \quad\hbox{for every }j.
\]
The deck generator cyclically permutes the conjugacy classes of
the $Y_j$.  It therefore permutes the eleven characters having a
single nontrivial $Y_j$-value and fixes the character whose eleven
$Y_j$-values are all $\omega$.  These two orbits are disjoint.
\end{proof}

Consequently the direct operation ``apply \eqref{eq:eleventh-local-system-K}, then pass to the
pulled-back rank-one local system'' does not identify the product
\begin{equation}\label{eq:eleventh-local-system-P}
P_g=\prod_{\substack{\xi^{11}=1\\\xi\ne1}}B_g(\xi,\omega)
\end{equation}
with the determinant on that pulled-back local system.  This
observation is a limitation of this proposed argument, not a claim
that the full relation \eqref{eq:eleventh-local-system-K} cannot constrain \eqref{eq:eleventh-local-system-P} by another route.

Here is a separate, explicit limitation of a pairing-only argument.
Let $C$ be the smooth compact curve with affine equation
$u^{11}+v^3=1$, and let
$A=C_{11}\times C_3$ be its deck group over the $t$-line,
$t=u^{11}$.  We use the geometric finite homology description:
over an algebraic closure of $\mathbb F_{17}$, its $20$-dimensional
homology has exactly the character lines
\begin{equation}\label{eq:eleventh-local-system-H}
(\xi^a,\omega^b),\qquad
 a\in(\mathbb Z/11)^\times,\quad b\in\{1,2\}.
\end{equation}
Every line is one-dimensional.  The alternating intersection form
pairs a line only with its inverse-character line.
For completeness, the character description follows from the
two-loop cellular complex of the three-punctured sphere.  On a
nontrivial character the boundary from the two one-cells to the
one zero-cell has rank one, so the open cover has a single
homology line.  Compactification kills precisely the peripheral
lines whose character is trivial on one of the three inertia
groups.  The inertia values here are $\xi^a$, $\omega^b$ and
$(\xi^a\omega^b)^{-1}$.  Since eleven and three are coprime,
all three are nontrivial exactly for the pairs in \eqref{eq:eleventh-local-system-H}.
The trivial character contributes no compact homology.

\begin{proposition}\label{prop:kn-linear-sign-model}
There is an $\mathbb F_{17}$-linear automorphism of the space \eqref{eq:eleventh-local-system-H}
that commutes with the whole deck group $A$, preserves its
alternating intersection form, and has determinant $-1$ on the
ten-dimensional $\omega$-eigenspace over $\mathbb F_{17^2}$.
Its quadratic norm-one determinant sign is therefore $-1$.
This is a linear model, not a GT element or a solution of the
pentagon or of \eqref{eq:eleventh-local-system-K}.
\end{proposition}
\begin{proof}
The square of Frobenius acts on the nonzero exponents modulo
eleven by multiplication by $17^2\equiv3$.  Its two orbits are
\[
 A_0=\{1,3,9,5,4\},\qquad
 A_1=\{2,6,7,10,8\}=-A_0.
\]
Ordinary Frobenius acts on the pairs $(a,b)$ in \eqref{eq:eleventh-local-system-H} as
\[
 (a,b)\longmapsto(6a,-b).
\]
It thus preserves the set
\[
 \mathcal O=(A_0\times\{1\})\cup(A_1\times\{2\}).
\]
Indeed $6A_0=A_1$ and $6A_1=A_0$.
The set $\mathcal O$ is also stable under inversion
$(a,b)\mapsto(-a,-b)$, as is its complement.

Define the automorphism to be $-1$ on the ten character lines in
$\mathcal O$ and $+1$ on the ten remaining lines.  Frobenius
stability makes it descend to $\mathbb F_{17}$, and its scalar
action on each character line makes it commute with $A$.
The intersection form pairs only mutually inverse characters:
if $v,w$ have characters $\chi,\psi$, invariance gives
$\langle v,w\rangle=\chi(a)\psi(a)\langle v,w\rangle$ for every
deck element $a$, which forces $\psi=\chi^{-1}$ whenever the
pairing is nonzero.  Our automorphism has equal signs on these
paired lines, so it preserves that form.

There are exactly five negative lines in the $\omega$-eigenspace,
namely those indexed by $A_0$.  Its determinant there is
$(-1)^5=-1$.  This value belongs to the norm-one group of
$\mathbb F_{17^2}^\times$, and raising it to
$(17+1)/2=9$ leaves $-1$, as asserted.
\end{proof}

The exact remaining question is whether the full nonabelian
relations, including \eqref{eq:eleventh-local-system-K} and the pentagon, exclude this sign in
the GT action.  Neither Lemma~\ref{lem:kn-different-local-systems}
nor Proposition~\ref{prop:kn-linear-sign-model} answers it.
In particular, no assertion about nonvanishing of
$\Theta_{17,11}|_{\widehat{\mathrm{GT}}}$ follows from these two statements.


\section{The comparison with Galois and the gentle group}
The factorization also makes the nonabelian equality question and the
role of gentle characters precise.

\begin{proposition}\label{prop:normalized-comparison}
There is a natural homeomorphism of right-coset spaces
\[
 N/K\longrightarrow G/H,\qquad nK\longmapsto nH.
\]
In particular $G=H$ if and only if $N=K$.
Put $N_0=\ker(\lambda:\GTz\to U)\cap\ker\rho_2$.
Then $\GTz=N_0H$. If a continuous homomorphism from $\GTz$
to any topological group is trivial on $H$, its restriction to $G$
is trivial if and only if its restriction to $N$ is trivial.
\end{proposition}
\begin{proof}
Surjectivity of the coset map is $G=NH$. Two elements of $N$ have
the same right $H$-coset precisely when their quotient lies in
$N\cap H=K$. The map is continuous; compactness of the source and
the Hausdorff property of the target make it a homeomorphism.
There is a single coset exactly when $N=K$, proving the equivalence.

For $g\in\GTz$, choose $h\in H$ with $\lambda(h)=\lambda(g)$.
By \cite[Lemma~7(c)]{H17}, the value $\rho_2(gh^{-1})$ is even;
that lemma uses only the three-cycle equation. The Galois Kummer
image over $\Q^{\mathrm{ab}}$ is $2\Zh$, as proved in
Proposition~\ref{prop:arithmetic-affine}. Choose
$k\in G_{\Q^{\mathrm{ab}}}$ with
$\rho_2(k)=\rho_2(gh^{-1})$. The cocycle law gives
$gh^{-1}k^{-1}\in N_0$, proving the gentle factorization.
Finally $G=NH$ shows directly that a homomorphism killing $H$
kills $G$ exactly when it kills $N$.
\end{proof}

The equivalence is about equality of the two full profinite groups.
The quotient $B$ records its abelian consequence. In particular a
gentle character trivial on Galois proves strictness of $G\supset H$
only if it is nonzero on $N$; a character vanishing on the pentagon
group does not provide such a witness.

\section{Comparison of the two higher-genus actions}
The explicit action on the groups with zero or one marked point agrees,
as an outer action, with the action constructed from quilts. The following
calculation gives the change of representative.

\begin{proposition}\label{prop:ns-lns-comparison}
Use the compatible standard markings of Nakamura--Schneps
\cite[Theorem 11.2]{NS2000} and Lochak--Nakamura--Schneps
\cite[Theorem 1.1]{LNS}. For $F=(\lambda,f)\in\widehat{\operatorname{GT}}$,
write $\rho=\rho_2(F)$ and
\[
 w_1=1,\qquad w_j=(a_1\cdots a_{j-1})^j\quad(2\leq j\leq2g),
 \qquad
 K_F=\prod_{j=2}^{2g}w_j^{\,(-1)^{j+1}4\rho}.
\]
Let $A_F^{\mathrm{NS2000}}$ denote the standard quilt representative and
$A_F^{\mathrm{LNS}}$ the representative given by the explicit chain
formulas. On each of $\widehat\Gamma_{g,0}$ and
$\widehat\Gamma_{g,1}$ one has
\[
 A_F^{\mathrm{NS2000}}=\operatorname{Ad}(K_F)\circ A_F^{\mathrm{LNS}}.
\]
In particular the two outer actions are equal. If $\rho_2(F)=0$, these
standard representatives themselves are equal.
\end{proposition}
\begin{proof}
The equality of the Nakamura--Schneps subgroup with
$\widehat{\operatorname{GT}}$ is used as established background
\cite[Corollary 21]{H17}; it ensures that their formulas apply to every
$F$ in the statement.

The elements $a_1,\ldots,a_{2g}$ satisfy the braid relations. The full
braid twist $w_j$ is central in the subgroup generated by
$a_1,\ldots,a_{j-1}$. Hence the $w_j$ commute with one another, and
$w_j$ commutes with $a_i$ whenever $j\ne i$. The usual full-twist
factorization gives
\[
 w_i=y_iw_{i-1},\qquad
 y_i=a_{i-1}\cdots a_1^2\cdots a_{i-1}.
\]
The element $w_{i-1}$ commutes both with $a_i$ and with $y_i$.
If $z$ centralizes $a,b$ and $f$ belongs to the closed commutator
subgroup of a free profinite group on two generators, then
$f(a,bz)=f(a,b)$. Indeed the homomorphism from
$\langle a,b\rangle\times\langle z\rangle$ given by multiplication
sends the second coordinate of a commutator word to the identity;
the identity extends from commutator words by continuity.
Consequently
\[
 f(a_i^2,w_i)=f(a_i^2,y_i),\qquad
 f(w_i,a_i^2)=f(y_i,a_i^2).
\]

Set $c_i=f(a_i^2,y_i)$. All $w_j$ with $j\ne i$ commute with
$c_i$ and with $a_i$. The explicit chain action is
$A_F^{\mathrm{LNS}}(a_i)=c_i a_i^\lambda c_i^{-1}$.
Conjugating by $K_F$ therefore gives
\[
 w_i^{\,(-1)^{i+1}4\rho}
 f(a_i^2,w_i)a_i^\lambda f(w_i,a_i^2)
 w_i^{-(-1)^{i+1}4\rho},
\]
which is the standard quilt formula. For $a_1$ both maps give
$a_1^\lambda$, and $K_F$ commutes with $a_1$.

For $g\geq2$, the remaining generator is $d=d_2$, a curve of the
standard pants decomposition. In the same marking,
\[
 w_2=a_1^2,\qquad
 w_{2i}=d_i d_{-i}\quad(i\geq2),\qquad
 w_{2i-1}^{2}=e_{i-1}\quad(i\geq2).
\]
All curves appearing on the right belong to that pants decomposition.
Since $4\rho$ is divisible by two, these identities express $K_F$ as
a product of powers of its mutually commuting pants twists. Thus
$K_F$ commutes with $d$. Both representatives take $d$ to $d^\lambda$.
Equality on the displayed generating set proves the assertion; in
genus one the chain generators already suffice. The same identities
survive the relation filling the marked point, so the argument also
applies to $\widehat\Gamma_{g,0}$.
\end{proof}

This is an equality on the overlapping families of full profinite
mapping class groups. A comparison with another completed-tower
construction additionally requires that construction's action formulas
or a stated uniqueness theorem for its objects and morphisms. The
existence of actions with the same genus-zero restriction alone does
not identify their chosen representatives on every object of a tower.

\section{The scope of inertia-preserving identifications}
For a generic surface $S$ of type $(g,n)$, namely
$g>0$, $3g-3+n>2$, and $(g,n)\ne(2,0)$, we use the established
identification
\[
 \operatorname{Out}^{*}(\check\Gamma(S))
       \cong\widehat{\operatorname{GT}}.
\]
Here the star requires preservation of the twist inertia subgroups.
This is \cite[Corollary 22]{H17}, using Lochak's compatible-pair
classification \cite[Definition 5.4, Proposition--Definition 5.6,
and Theorem 5.10(ii)]{Lochak2025}. It is a cited result about the
procongruence completion. Passing from nonseparating inertia to all
inertia is a further assertion.

Boggi states this further assertion as
$\operatorname{Aut}^{I_0}(\check\Gamma(S))=
\operatorname{Aut}^{I}(\check\Gamma(S))$ in
\cite[Theorem 3.12 and Corollary 3.15]{Boggi}.
The proof proceeds through the nonseparating curve complex and requires
recognition of the multicurves with connected complement, as used in
\cite[Remark 3.11]{Boggi}. Commutation of individual twists recognizes
disjointness. It does not by itself recognize connectedness of the
complement: a bounding pair and a pair with connected complement both
consist of disjoint nonseparating curves. Thus the latter recognition
step requires its own justification.

\begin{proposition}\label{prop:nonunique-stabilizer-extension}
Let $\sigma=\{\gamma_1,\ldots,\gamma_{g-1}\}$ be a cut system with
connected complement, and let $H$ be its oriented stabilizer in the
full profinite pure mapping class group. Let $H_i$ be the oriented
stabilizer of $\sigma\setminus\{\gamma_i\}$. The identity
of $H$ has more than one inertia-preserving extension to $H_i$.
\end{proposition}
\begin{proof}
The twist $\tau_{\gamma_i}$ is central in $H$, since every element
of the oriented stabilizer fixes its twist. Thus the identity and
$\operatorname{Ad}(\tau_{\gamma_i})$ restrict to the same automorphism
of $H$. Both are inertia-preserving automorphisms of $H_i$.
Choose a simple curve $\beta_i$ meeting $\gamma_i$ once and disjoint
from the other curves of the cut system. Its twist belongs to $H_i$.
On the symplectic plane generated by these two curves, their twist
matrices are
\[
 U=\begin{pmatrix}1&1\\0&1\end{pmatrix},\qquad
 V=\begin{pmatrix}1&0\\-1&1\end{pmatrix}.
\]
They do not commute, already after reduction modulo three. Thus
$\tau_{\gamma_i}$ is not central in $H_i$, and its inner automorphism
is not the identity. This finite symplectic quotient also shows that
the distinction survives profinite completion.
\end{proof}

Accordingly, uniqueness as a genuine automorphism in
\cite[Lemma 6.16]{Boggi} fails as written. This counterexample does
not refute a statement of uniqueness up to inner automorphisms.
There is a separate group identification issue in the proof of
\cite[Proposition 6.9]{Boggi}: the quotient
$B_{2g+1}/\langle z^2\rangle$ obtained by filling a distinguished
Weierstrass boundary point preserves that point. It therefore gives
the corresponding stabilizer in the closed hyperelliptic group,
whose index is $2g+2$, rather than the whole closed hyperelliptic
group. Neither observation is a counterexample to the earlier
nonseparating-inertia theorem; they concern distinct proof assertions.

\section{A global character from the two elliptic frames}
\label{sec:global-frame-character}

We define the frame signs, extend their sum to every multiplier, and
fix its arithmetic normalization. The construction uses the normalized
frame vanishing theorem of \cite[Theorem~3]{H17}. The extension and
the conjugation invariance are proved here.

\subsection{The frames and their normalization}

Let $\widehat B_3=\langle u,v\mid uvu=vuv\rangle$ be the profinite
braid group and put $c=(uv)^3$. Its central quotient is
\[
 \Gamma=\widehat B_3/\overline{\langle c\rangle}
      =\langle s,r\mid s^2=r^3=1\rangle,
 \qquad s=uvu,\quad r=uv,\quad \tau=sr=u^{-1}.
\]
The free pure subgroup $F=\ker(\Gamma\to S_3)$ has generators
$x=u^2=\tau^{-2}$ and $y=v^2=sxs$. We also use the subgroup
$\langle u^2,v^2\rangle\subset\widehat B_3$, which maps isomorphically
to $F$. A gentle element $g=(\lambda_g,f_g)$ acts on the braid
group by
\[
 \Phi_g(u)=u^{\lambda_g},\qquad
 \Phi_g(v)=f_g(u^2,v^2)^{-1}v^{\lambda_g}f_g(u^2,v^2),
 \qquad \Phi_g(c)=c^{\lambda_g}.
\]
These are automorphisms and induce the identity on the permutation
quotient \cite[Section 3, before Lemma~5]{H17}.

The finite-subgroup and centralizer theorem for the free profinite
product gives unique pure frames
\[
 \Phi_g(s)=\beta_gs\beta_g^{-1}\quad(g\in\GTz),
 \qquad
 \Phi_h(r)=A_hrA_h^{-1}\quad(\lambda_h=1),
 \qquad \beta_g,A_h\in F.
\]
These are the frames of \cite[Lemma~5]{H17}. For the cubic frame,
the multiplier-one condition fixes the conjugacy class of $r$ in
$\Gamma^{\mathrm{ab}}=C_6$, distinguishing it from $r^{-1}$.
The frames vary continuously: their graphs are closed in compact
products, and uniqueness makes the corresponding projections
homeomorphisms. Their composition identities are
\begin{equation}
 \beta_{gh}=\Phi_g(\beta_h)\beta_g,
 \qquad A_{hk}=\Phi_h(A_k)A_h
       \quad(\lambda_h=\lambda_k=1).
 \label{eq:frame-cocycles}
\end{equation}

Let $\theta$ interchange $x,y$. Comparison of the two expressions
for $\Phi_g(y)$ gives
$f_g=\theta(\beta_g)\beta_g^{-1}$. Indeed their possible discrepancy
centralizes $y^{\lambda_g}$; abelianization makes its exponent zero.
Thus the two coordinates of $[\beta_g^{-1}]\in F^{\mathrm{ab}}$
are equal. Define $\rho_2(g)$ by
\[
 [\beta_g^{-1}]=(\rho_2(g),\rho_2(g)).
\]
It is a continuous cocycle,
\[
 \rho_2(gh)=\rho_2(g)+\lambda_g\rho_2(h).
\]
Its restriction to $G_{\mathbb Q}$ is the Kummer cocycle of the
coherent positive roots of two \cite[Lemma~7]{H17}.
Write
\[
 \mathcal G=\GTz,\qquad
 N_0=\{g\in\mathcal G:\lambda_g=1,\ \rho_2(g)=0\},
 \qquad H=G_{\mathbb Q}.
\]
The cocycle makes $N_0$ closed and normal.

We will use the arithmetic decomposition
\begin{equation}
 \mathcal G=N_0H,\qquad
 N_0\cap H=G_{K_2},\qquad
 K_2=\mathbb Q(\mu_\infty,2^{1/\infty}).
 \label{eq:gentle-arithmetic-normalization}
\end{equation}
This is Proposition~\ref{prop:normalized-comparison}; its intersection formula follows from the two defining Galois actions.

\subsection{The signed matrices at level twelve}

Put
\[
 U=\begin{pmatrix}1&1\\0&1\end{pmatrix},\qquad
 L=\begin{pmatrix}1&0\\-1&1\end{pmatrix},\qquad
 S=ULU,\qquad R_+=UL,\qquad R=-R_+.
\]
The assignments $u\mapsto U$, $v\mapsto L$ define the signed
representation $\widetilde q_\ell:\widehat B_3\to
\operatorname{SL}_2(\mathbb Z_\ell)$ for every prime $\ell$.
Its restriction to $F$ is denoted $q_\ell$. Reducing the signed
representation modulo twelve gives $q_{12}$ on $F$ and a projective
representation $q$ of $\Gamma$.
The signed pure image is
\[
 E_{12}=V_4\times\operatorname{SL}_2(\mathbb F_3),
 \qquad V_4=\{I,U^2,L^2,U^2L^2\}\pmod4.
\]
The matrix $-I$ is not in $V_4$, so this image embeds in the
projective group modulo twelve. Its projective centralizers of
$R_+$ and $S$ are respectively
\begin{equation}
 \{I\}\times\bigl(\langle-I\rangle\times\langle R\rangle\bigr),
 \qquad
 \{I,U^2L^2\}\times\langle S\rangle.
 \label{eq:finite-frame-centralizers}
\end{equation}
These follow directly by multiplying the indicated matrices; they
are the finite groups in \cite[Lemma~12]{H17}.

\begin{lemma}\label{lem:frame-all-multipliers-congruence}
For every $g\in\mathcal G$ and every prime $\ell$ there is a unique
matrix
\[
 D_{g,\ell}=\begin{pmatrix}\lambda_{g,\ell}&b_{g,\ell}\\0&1\end{pmatrix}
\]
such that
$\widetilde q_\ell\Phi_g=\operatorname{Ad}(D_{g,\ell})\widetilde q_\ell$.
It depends continuously on $g$, and $D_{gh,\ell}=D_{g,\ell}D_{h,\ell}$.
For $h\in N_0$, its reduction modulo twelve is the identity.
\end{lemma}

\begin{proof}
Put $a=\lambda_{g,\ell}$ and
$M=\widetilde q_\ell(\Phi_g(v))=
\left(\begin{smallmatrix}\alpha&b\\d&2-\alpha\end{smallmatrix}\right)$.
It is conjugate to $L^a$, so its trace is two and its determinant
is one. The lower-left entry of the braid relation
$U^aMU^a=MU^aM$ gives $d(1+ad)=0$.
If $d=0$, the domain property of $\mathbb Z_\ell$, together with
the trace and determinant, makes $M$ upper unipotent. The images
of both braid generators would commute, contradicting surjectivity
onto $\operatorname{SL}_2(\mathbb Z_\ell)$. Hence $d=-a^{-1}$
and $b=a(\alpha-1)^2$. Direct multiplication gives
\[
 D_{g,\ell}=
  \begin{pmatrix}a&a(1-\alpha)\\0&1\end{pmatrix}.
\]
This proves existence and continuity. A matrix commuting with both
$U,L$ is scalar; the specified lower-right entry gives uniqueness.
Composition then gives the product identity.

For $\lambda_h=1$, write the reduction modulo twelve as $U^{t_h}$.
Modulo four the conjugator $f_h\in F'$ maps trivially because
$V_4$ is abelian, so $t_h=0\pmod4$. Modulo three,
$U^{-t_h}q_3(\beta_h)$ lies in the quaternion subgroup of
$\operatorname{SL}_2(\mathbb F_3)$: its projective class centralizes
the involution represented by $S$. In the quotient of order three,
the classes of $U$ and $L$ agree, whereas
$[\beta_h]=(-\rho_2(h),-\rho_2(h))$ gives class
$-4\rho_2(h)=-\rho_2(h)$. Thus
$t_h=-\rho_2(h)\pmod3$. This is the shear formula of
\cite[Lemma~13]{H17}; in particular $D_h=I\pmod{12}$ on $N_0$.
\end{proof}

Define the closed subgroup
\[
 H_0=\{h\in\mathcal G:\lambda_h=1,\ t_h=0\pmod{12}\}.
\]
It contains $N_0$. For $h\in H_0$, the two frames lie in
\eqref{eq:finite-frame-centralizers}, so define
\begin{equation}
 q_3(A_h)\in(-I)^{a_3(h)}\langle R\rangle,
 \qquad
 q_3(\beta_h)\in S^{a_2(h)}\langle-I\rangle,
 \qquad a_2(h),a_3(h)\in\mathbb F_2.
 \label{eq:frame-sign-definition}
\end{equation}
These are exactly the signs of \cite[Definition~14]{H17}.
The cocycles \eqref{eq:frame-cocycles}, the identity action on
$E_{12}$, and the abelian groups in
\eqref{eq:finite-frame-centralizers} make $a_2,a_3$ continuous
characters on $H_0$. Put $\delta=a_2+a_3$.

The sign convention can also be recorded by a single unit. Let
$\xi(h)\in(\mathbb Z/12\mathbb Z)^\times$ be determined by
\[
 \xi(h)\equiv(-1)^{a_2(h)}\pmod4,\qquad
 \xi(h)\equiv(-1)^{a_3(h)}\pmod3.
\]
For the quadratic Dirichlet character $\chi_{12}$ the complete table is
\[
 \begin{array}{c|rrrr}
  t&1&5&7&11\\ \hline
  \chi_{12}(t)&1&-1&-1&1
 \end{array}
 \qquad (-1)^{\delta(h)}=\chi_{12}(\xi(h)).
\]
All elements of $H_0$ have multiplier one. The unit $\xi$ records
the two frame signs separately from that multiplier.

\subsection{Conjugation by elements of every multiplier}

\begin{proposition}\label{prop:frame-sign-conjugation}
The restrictions of $a_2,a_3$, and hence $\delta$, to $N_0$ are
invariant under conjugation by $\mathcal G$.
\end{proposition}

\begin{proof}
Let $h\in N_0$, $g\in\mathcal G$, and $\eta=ghg^{-1}$.
Normality gives $\eta\in N_0$. All matrix calculations in this
proof are modulo three. Write $\overline D_g$ for the reduction
of $D_{g,3}$.

Applying $\Phi_\eta\Phi_g=\Phi_g\Phi_h$ to $s$ and comparing
unique pure frames gives
\[
 \beta_\eta=\Phi_\eta(\beta_g)^{-1}\Phi_g(\beta_h)\beta_g.
\]
Since $\overline D_\eta=I$,
\[
 q_3(\beta_\eta)=q_3(\beta_g)^{-1}\overline D_g
      q_3(\beta_h)\overline D_g^{-1}q_3(\beta_g).
\]
Thus the traces of their squares agree. The defining groups give
\[
 \operatorname{tr}(q_3(\beta_h)^2)=2(-1)^{a_2(h)},
\]
which proves invariance of $a_2$.

Put $\epsilon\in\{1,-1\}$ equal to $\lambda_g$ modulo three.
The torsion theorem gives
$\Phi_g(r)=\gamma_g r^\epsilon\gamma_g^{-1}$.
The exponent follows from the action on $\Gamma^{\mathrm{ab}}=C_6$.
Since $\Phi_g$ induces the identity on $S_3$, multiplying
$\gamma_g$ by a power of $r$ permits the form
$\gamma_g=B_gs^e$, with $B_g\in F$ and $e=0$ or one.
Here $e=0$ when $\epsilon=1$ and $e=1$ when $\epsilon=-1$.
Comparison of the two actions on $r$ gives
\begin{equation}
 A_\eta=\Phi_\eta(\gamma_g)^{-1}\Phi_g(A_h)\gamma_g.
 \label{eq:cubic-frame-conjugation}
\end{equation}
The right side is pure and conjugates $r$ to $\Phi_\eta(r)$,
so equality follows from uniqueness of the pure frame.

We check this equality in the signed representation. Lift
$\gamma_g$ to $\widetilde\gamma_g=B_g(uvu)^e$ in
$\widehat B_3$. Conjugation by $uvu$ interchanges $x,y$ and
normalizes the embedded free subgroup. Since $\lambda_\eta=1$,
\[
 \Phi_\eta(uvu)=\beta_\eta(uvu)\beta_\eta^{-1}.
\]
Both sides project to the same modular element and have square $c$.
A central discrepancy would have square one, and the central
procyclic group $\overline{\langle c\rangle}\simeq\widehat{\mathbb Z}$
is torsion free. This proves the equality.
It follows that the lifted right side of
\eqref{eq:cubic-frame-conjugation} belongs to the embedded free
subgroup: the two factors involving $uvu$ cancel through its
normalizing action. Its projection equals $A_\eta$, and the free
subgroup injects into the central quotient, so it is the signed
lift of $A_\eta$ itself.

Applying the signed matrix representation therefore gives
\[
 q_3(A_\eta)=
 \widetilde q_3(\widetilde\gamma_g)^{-1}\overline D_g
 q_3(A_h)\overline D_g^{-1}
 \widetilde q_3(\widetilde\gamma_g).
\]
The trace is unchanged. Every power of $R$ has trace two in
$\mathbb F_3$, so
$\operatorname{tr}(q_3(A_h))=2(-1)^{a_3(h)}$.
This proves invariance of $a_3$ as well.
\end{proof}

\begin{theorem}[Arithmetic normalization of the frame character]
\label{thm:global-frame-character}
There is a unique continuous character
\[
 \Delta:\mathcal G\longrightarrow\mathbb F_2
\]
which restricts to $\delta$ on $N_0$ and is zero on $G_{\mathbb Q}$.
For any $g\in\mathcal G$, choose $\sigma\in G_{\mathbb Q}$
with $(\lambda,\rho_2)(\sigma)=(\lambda,\rho_2)(g)$. Then
\[
 \Delta(g)=\delta(g\sigma^{-1}).
\]
This character vanishes on the full pentagon group.
\end{theorem}

\begin{proof}
By \cite[Theorem~3]{H17}, $\delta$ vanishes on the normalized
pentagon subgroup. In particular it vanishes on
$N_0\cap H=G_{K_2}$. Using
\eqref{eq:gentle-arithmetic-normalization}, define
$\Delta(n\sigma)=\delta(n)$ for $n\in N_0$, $\sigma\in H$.
Two decompositions differ by an element of $N_0\cap H$, so this
is well defined. Proposition~\ref{prop:frame-sign-conjugation}
proves its homomorphism law. The multiplication map
$N_0\times H\to\mathcal G$ is a quotient map, since it is a
continuous surjection from a compact space to a Hausdorff space.
Thus the character is continuous. The two specified restrictions
and the decomposition prove uniqueness, as well as the displayed
formula. Finally \cite[Theorem~2]{H17} decomposes the pentagon
group into its normalized subgroup and $G_{\mathbb Q}$, both of
which have zero image. Hence $\Delta$ vanishes on $\GT$.
\end{proof}

\section{The individual frame signs and the Kummer cocycle of three}
\label{sec:harmonic-kummer-comparison}

The two normalized frame signs can be identified separately on the
pentagon group. The comparison uses Anderson duplication, already a
consequence of the cited distribution and braid identities.

\begin{theorem}\label{thm:individual-frame-kummer}
For every $g\in N=\ker(\lambda,\rho_2)\subset\widehat{\mathrm{GT}}$,
\[
 a_2(g)=a_3(g)=\frac{\rho_3^I(g)}2\pmod2.
\]
In particular the character
$a_2+\rho_3^I/2\bmod2$ is zero on $N$. Both frame signs vanish
on the simultaneous radical-normalized subgroup $N_{\mathrm{rad}}$.
\end{theorem}
\begin{proof}
Let $g=(1,f)$ and put $F=\widehat F_2=\langle x,y\rangle$.
Write $c=[x,y]=xyx^{-1}y^{-1}$. The module $F'/F''$ is free of
rank one over
\[
 \mathcal A=\widehat{\mathbb Z}[[F^{\mathrm{ab}}]],
\]
with generator $[c]$; see \cite[\S2.3, p.~101]{Furusho2003}.
Let $X,Y$ denote the two conjugation operators and write
$[f]=A_f(X,Y)[c]$. The two-cycle identity gives
\begin{equation}
 A_f(X,Y)=A_f(Y,X),
 \label{eq:harmonic-A-symmetry}
\end{equation}
because interchange of $x,y$ sends $[c]$ to $-[c]$.
The action $x\mapsto x$, $y\mapsto f^{-1}yf$ gives
\[
 [\Phi_g(c)]=[c]+(X-1)(Y-1)[f].
\]
Thus its multiplier on this module is
\begin{equation}
 B_g(X,Y)=1+(X-1)(Y-1)A_f(X,Y).
 \label{eq:harmonic-beta-coefficient}
\end{equation}
For multiplier one this is also the adelic beta function, since
the normalizing cyclotomic factors in \cite[\S2.3]{Furusho2003}
are equal to one.

Work in $\mathcal O=\mathbb Z_2[\omega]$, where
$\omega^2+\omega+1=0$. Let $\nu_g$ be the triangle-group lattice
unit of \cite[Lemma 15]{H17}. Its relation to the beta multiplier is
\begin{equation}
 \nu_g=B_g(\omega,\omega),\qquad
 \nu_g\equiv(-1)^{a_2(g)}\pmod{4\mathcal O}.
 \label{eq:harmonic-triangle-beta}
\end{equation}
For the first equality, in the affine triangle quotient the two pure
generators have multiplier $\omega$ and $[x,y]$ is the translation
by $-3$. The action on translations is multiplication by $\nu_g$.
Applying the module action to $[x,y]$ therefore gives
$-3\nu_g=-3B_g(\omega,\omega)$, and $3$ is a unit in
$\mathcal O$. The second equality is the cited lattice calculation.

Since $\rho_2^I(g)=\rho_2(g)=0$, the Anderson duplication identity
at the torsion parameter $s=1/3$ gives
\begin{equation}
 B_g(\omega,\omega)=B_g(-1,\omega).
 \label{eq:harmonic-duplication}
\end{equation}
Indeed in the notation of \cite[Theorem 3.1]{Furusho2003} the
quotient of these two values is the third root of unity with exponent
$\Psi_2^{(0)}(g)=-\rho_2^I(g)$, which is zero here. The identity
is in the full adelic beta function; its evaluation in
$\mathcal O$ is therefore legitimate.

It remains to compute the right side modulo four. In the degree-three
cyclic-cover subgroup with basis $x^3,y_0,y_1,y_2$, write the
abelianized image of $f$ as
$c_0y_0+c_1y_1+c_2y_2$. The $x^3$ coefficient is zero.
The cyclic-cover symmetry of
Lemma~\ref{lem:cyclic-cover-symmetry} gives $c_1=c_2$;
abelianizing in $F$ gives $c_0+c_1+c_2=0$.
Consequently, for $b=\ell_3(f)/2=-\rho_3^I(g)/2$,
\[
 (c_0,c_1,c_2)=(2b,-b,-b).
\]
Map this cover homology to $\mathcal O$ by
$y_j\mapsto\omega^j$, $x^3\mapsto0$.
Conjugation by $x$ acts as multiplication by $\omega$ and
conjugation by $y$ acts trivially. The commutator $c$ maps to
$\omega-1$. Applying the map to $[f]=A_f[c]$ therefore gives
\begin{equation}
 (\omega-1)A_f(\omega,1)
   =2b-b\omega-b\omega^2=3b.
 \label{eq:harmonic-cover-coefficient}
\end{equation}

The evaluations at $(-1,\omega)$ and $(1,\omega)$ agree
modulo $2\mathcal O$. By \eqref{eq:harmonic-A-symmetry},
$A_f(1,\omega)=A_f(\omega,1)$. Using
\eqref{eq:harmonic-beta-coefficient} and
\eqref{eq:harmonic-cover-coefficient}, we obtain
\[
 \begin{split}
 \nu_g
 &=B_g(-1,\omega)\\
 &\equiv1-2(\omega-1)A_f(\omega,1)\\
 &=1-6b\equiv1+2b\pmod{4\mathcal O}.
 \end{split}
\]
Together with \eqref{eq:harmonic-triangle-beta}, this proves
$a_2(g)=b=\rho_3^I(g)/2\pmod2$.
The established vanishing $a_2+a_3=0$ on $N$
\cite[Theorem 3]{H17} gives the assertion for $a_3$.
Every rational-prime Kummer cocycle is zero on $N_{\mathrm{rad}}$,
so both signs vanish there.
\end{proof}

The comparison identifies these particular finite-frame coordinates.
It does not determine the residual quotient $B$ or its quotient $B/2B$.

\section{The full adelic image of the two elliptic frames}
\label{sec:adelic-cm-frames}

The two corrected elliptic frames detect no further relative abelian
characters on the profinite Grothendieck--Teichm\"uller group.  The
reason is an equality of full adelic images: the multiplier-one
subgroup and the cyclotomic Galois subgroup have the same image in the
product of the two CM tori.  We prove the integral comparison, identify
the one quadratic entanglement between the division fields, and obtain
a second normalization of the main character problem.

\subsection{Corrected frames and their integral targets}

Write $\mathcal G=\widehat{\mathrm{GT}}_0$ for the gentle group,
$G=\widehat{\mathrm{GT}}$, $H=G_{\mathbb Q}$, and
\[
 M_0=\ker(\lambda:\mathcal G\longrightarrow
                       \widehat{\mathbb Z}^{\times}),
 \qquad H^{\mathrm{cyc}}=G_{\mathbb Q^{\mathrm{ab}}}.
\]
We recall the normalization of the frames.  In the profinite braid
group put
\[
 \widehat B_3=\langle u,v\mid uvu=vuv\rangle,\qquad
 c=(uv)^3,\qquad s=uvu,\quad r=uv,
\]
\[
 \Gamma=\widehat B_3/\overline{\langle c\rangle}
             =\langle s,r\mid s^2=r^3=1\rangle.
\]
The free pure subgroup is
$F=\ker(\Gamma\to S_3)=\widehat F(x,y)$, where
$x=u^2$ and $y=v^2$.  The subgroup generated by these two elements
in $\widehat B_3$ identifies with $F$.  For $g\in\mathcal G$ its
braid action is
\[
 \Phi_g(u)=u^{\lambda_g},\qquad
 \Phi_g(v)=f_g^{-1}v^{\lambda_g}f_g,\qquad
 \Phi_g(c)=c^{\lambda_g},\qquad f_g\in F'.
\]
The unique pure frames satisfy
\[
 \Phi_g(s)=\beta_gs\beta_g^{-1},\qquad
 \Phi_h(r)=A_hrA_h^{-1}\quad(h\in M_0),
 \qquad \beta_g,A_h\in F.
\]
Their existence, uniqueness, and cocycles are as in
\cite[Lemmas~5 and~7]{H17}; in particular
\begin{equation}
 \beta_{gh}=\Phi_g(\beta_h)\beta_g,\qquad
 A_{gh}=\Phi_g(A_h)A_g\quad(g,h\in M_0),\qquad
 [\beta_g^{-1}]=(\rho_2(g),\rho_2(g)).
 \label{eq:adelic-frame-cocycles}
\end{equation}
The parity result gives $\rho_2(M_0)\subset2\widehat{\mathbb Z}$.

Set
\[
 U=\begin{pmatrix}1&1\\0&1\end{pmatrix},\quad
 L=\begin{pmatrix}1&0\\-1&1\end{pmatrix},\quad
 S=ULU=\begin{pmatrix}0&1\\-1&0\end{pmatrix},\quad
 R=-UL=\begin{pmatrix}0&-1\\1&-1\end{pmatrix}.
\]
Let $\widetilde q:\widehat B_3\to
\operatorname{SL}_2(\widehat{\mathbb Z})$ send $u,v$ to $U,L$,
and let $q$ be its restriction to $F$.  The intertwiner proved above
has the form
\begin{equation}
 D_g=\begin{pmatrix}\lambda_g&b_g\\0&1\end{pmatrix},\qquad
 \widetilde q\Phi_g=\operatorname{Ad}(D_g)\widetilde q,\qquad
 D_{gh}=D_gD_h.
 \label{eq:adelic-frame-intertwiner}
\end{equation}
Here all entries are adelic.  Its existence can also be seen directly,
prime by prime: if
$\widetilde q(\Phi_g(v))=\left(\begin{smallmatrix}
 a&b\\d&2-a\end{smallmatrix}\right)$, the braid equation gives
$d(1+\lambda_gd)=0$.  The zero possibility would make the two
braid images commute, contrary to their surjectivity onto
$\operatorname{SL}_2(\mathbb Z_\ell)$.  Thus
$d=-\lambda_g^{-1}$, $b=\lambda_g(a-1)^2$, and
$D_g=\left(\begin{smallmatrix}\lambda_g&\lambda_g(1-a)\\0&1
\end{smallmatrix}\right)$.  Its normalization makes it unique.

\begin{lemma}\label{lem:adelic-explicit-intertwiner}
For every \(g\in G\), the normalized intertwiner is
\[
 D_g=\begin{pmatrix}\lambda_g&8\rho_2(g)\\0&1\end{pmatrix}.
\]
In particular \(D_g=I\) whenever \(\lambda_g=1\) and
\(\rho_2(g)=0\).
\end{lemma}

\begin{proof}
We evaluate the established relation~\((\mathrm{IV})\) to determine
its upper-right entry, including its sign at two. Write
\(F_g=q(f_g)\) and \(b=b_g\), \(\lambda=\lambda_g\).
Conjugation by \(S\) exchanges \(U^2,L^2\), so duality gives
\(F_gSF_g=S\). The determinant-one identity
\(F_gSF_g^{\mathsf t}=S\) then makes \(F_g\) symmetric.
The intertwining relation at \(v\) is
\[
 (F_gD_g)L=L^\lambda(F_gD_g).
\]
Solving it with symmetry and determinant one gives
\[
 F_g=\frac{e}{\lambda}
 \begin{pmatrix}1&-b\\-b&\lambda^2+b^2\end{pmatrix},
 \qquad e\in\widehat\Z^\times,\quad e^2=1.
\]
Here \(e\) may have independent signs at different primes;
no global-sign assumption is made.

In the finite adeles let \(J=\operatorname{diag}(1,2)\).
Since \(JU^2J^{-1}=U\) and \(JL^2J^{-1}=L^4\),
relation~\((\mathrm{IV})\), already established for full GT
\cite[Theorem~1]{H17}, gives
\[
 JF_gJ^{-1}=L^{8\rho_2(g)}F_gU^{4\rho_2(g)}.
\]
The central factor has image
\((-I)^{-2\rho_2(g)}=I\). Comparing upper-right entries and
cancelling \(e\) yields
\[
 -\frac{b}{2\lambda}
       =\frac{4\rho_2(g)-b}{\lambda},
 \qquad b=8\rho_2(g).
\]
Both sides are integral, so the equality holds in
\(\widehat\Z\), including its component at two.
\end{proof}

Put $O_3=\mathbb Z[\omega]$, where $\omega^2+\omega+1=0$,
and $O_4=\mathbb Z[i]$.  Embed these rings in matrices by
$\omega\mapsto R$ and $i\mapsto S$.  A bar denotes their quadratic
involution.  Define
\begin{equation}
 N_j=\{a\in(O_j\otimes\widehat{\mathbb Z})^{\times}:
                                      a\bar a=1\},\qquad
 N_j(4)=\{a\in N_j:a_2\equiv1\pmod4\}.
 \label{eq:adelic-cm-targets}
\end{equation}
The subscript $2$ in this formula means the component at the prime two.

\begin{lemma}\label{lem:adelic-corrected-frames}
The formulas
\begin{equation}
 T_g=q(A_g)^{-1}D_g,\qquad
 W_g=q(\beta_g)^{-1}D_g\qquad(g\in M_0)
 \label{eq:adelic-corrected-frames}
\end{equation}
define continuous homomorphisms
$T:M_0\to N_3(4)$ and $W:M_0\to N_4(4)$.
\end{lemma}

\begin{proof}
The centralizations are signed equalities.  Indeed the two lifts of
$\Phi_g(s)=\beta_gs\beta_g^{-1}$ in $\widehat B_3$ differ by
$c^e$.  Squaring, and using $s^2=c$ and $\Phi_g(c)=c$, gives
$2e=0$ in $\widehat{\mathbb Z}$, hence $e=0$.  Cubing the
corresponding identity for $r$, with $r^3=c$, similarly gives
$3e=0$.  Thus $T_g$ commutes with $R$ and $W_g$ commutes with $S$
as integral matrices, without a scalar ambiguity.

The centralizer rings of $R$ and $S$ in
$\operatorname{Mat}_2(\widehat{\mathbb Z})$ are respectively
$\widehat{\mathbb Z}[R]$ and $\widehat{\mathbb Z}[S]$.
For example $aI+bR$ has determinant $a^2-ab+b^2$, and
$aI+bS$ has determinant $a^2+b^2$.  Since $\det D_g=1$ on
$M_0$, this places $T_g,W_g$ in $N_3,N_4$.
Using \eqref{eq:adelic-frame-cocycles} and
\eqref{eq:adelic-frame-intertwiner} gives, with the order fixed,
\[
 q(A_{gh})^{-1}D_{gh}
 =q(A_g)^{-1}D_gq(A_h)^{-1}D_h=T_gT_h;
\]
the same calculation applies to $W$.

Modulo four the image of $F$ is the abelian four-group
\[
 V_4=\{I,U^2,L^2,U^2L^2\}.
\]
Consequently $q_4(f_g)=1$.  The braid covariance at $v$ makes
$D_g\bmod4$ commute with $L$, which forces the upper-right entry
of $D_g$ to be zero modulo four.  Since $q_4(A_g)\in V_4$
commutes with $R$ and $C_{V_4}(R)=\{I\}$, one has
$T_g\equiv I\pmod4$.  Finally the two abelian coordinates of
$\beta_g$ are $-\rho_2(g)$, both even.  The map $q_4$ factors
through their reductions modulo two, so $q_4(\beta_g)=I$ and
$W_g\equiv I\pmod4$ as well.  All maps used are continuous.
\end{proof}

\begin{lemma}\label{lem:adelic-kernel-normality}
The joint kernel of $T,W$ in $M_0$ is normal in $\mathcal G$.
\end{lemma}

\begin{proof}
The formula $q(\beta_g)^{-1}D_g$ defines a homomorphism to
$\operatorname{GL}_2(\widehat{\mathbb Z})$ for every $g\in\mathcal G$:
the calculation in the preceding proof uses the beta cocycle,
which is valid for every multiplier.  Thus its kernel is normal.

For the cubic frame, fix $g\in\mathcal G$, $h\in M_0$, and
$\eta=ghg^{-1}$.  Put $\epsilon\in\{1,-1\}$ equal to
$\lambda_g$ modulo three.  The torsion theorem gives a conjugator
$\gamma_g$ with
$\Phi_g(r)=\gamma_g r^\epsilon\gamma_g^{-1}$.
Because $\Phi_g$ induces the identity on $S_3$, this conjugator
may be chosen as $\gamma_g=B_gs^e$ with $B_g\in F$ and
$e\in\{0,1\}$.  Comparing
$\Phi_\eta\Phi_g(r)$ with $\Phi_g\Phi_h(r)$ and using the
uniqueness of the pure frame gives
\[
 A_\eta=\Phi_\eta(\gamma_g)^{-1}\Phi_g(A_h)\gamma_g.
\]
Lift $\gamma_g$ to $\widetilde\gamma_g=B_g(uvu)^e$ in
$\widehat B_3$.  Since $\lambda_\eta=1$, the equality
$\Phi_\eta(uvu)=\beta_\eta(uvu)\beta_\eta^{-1}$ is signed,
by the squaring argument above.  Also $uvu$ normalizes the embedded
free group.  Consequently the displayed expression, formed with
these braid lifts, belongs to that free group and is its unique
lift of $A_\eta$.  Writing
$Q_g=\widetilde q(\widetilde\gamma_g)$, covariance therefore yields
\[
 q(A_\eta)=D_\eta Q_g^{-1}D_\eta^{-1}
                 D_gq(A_h)D_g^{-1}Q_g.
\]
As $D_\eta=D_gD_hD_g^{-1}$, inversion and multiplication by
$D_\eta$ give
\begin{equation}
 T_\eta=(Q_g^{-1}D_g)\,T_h\,(Q_g^{-1}D_g)^{-1}.
 \label{eq:adelic-cubic-kernel-conjugation}
\end{equation}
Thus the kernel of $T$ in $M_0$ is also stable under every
conjugation by $\mathcal G$.  Intersecting the two kernels proves
the assertion.
\end{proof}

There are two parity maps on the target.  Let
$\pi_3=\sqrt{-3}$, an associate of $1-\omega$.  For $t\in N_3(4)$
and $w\in N_4(4)$ define
\begin{equation}
 (-1)^{\epsilon_3(t)}=t_3\bmod\pi_3\in\mathbb F_3^{\times},\qquad
 w_3\bmod3=S^a\quad\Longrightarrow\quad
 \epsilon_4(w)=a\bmod2.
 \label{eq:adelic-parities}
\end{equation}
The latter is well defined because the norm-one group in
$(O_4/3)^{\times}=\mathbb F_9^{\times}$ is cyclic of order four.
Both maps are surjective.  Write
\begin{equation}
 \mathcal E=\{(t,w)\in N_3(4)\times N_4(4):
                       \epsilon_3(t)=\epsilon_4(w)\}.
 \label{eq:adelic-equal-parity}
\end{equation}

\subsection{The CM division fields and the normalized images}

Let $K_j=O_j\otimes\mathbb Q$, and consider
\[
 E_3:Y^2=X^3+1,\qquad E_4:Y^2=X^3-X.
\]
Their adelic Tate representations are denoted $\rho_3,\rho_4$.
Over $K_j$ they take values in $(O_j\otimes\widehat{\mathbb Z})^
\times$.  We use the following precise CM input: the torsion
reciprocity theorem identifies the action modulo geometric
automorphisms with the ray class action
\cite[Theorem~12.27, printed p.~124]{MilneMF110}.
Since both rings have class number one, the ray class exact sequence
identifies the inverse limit with
$(O_j\otimes\widehat{\mathbb Z})^\times/O_j^\times$.
The action on the Tate module is abelian over $K_j$
\cite[Remark~12.12 and Proposition~12.13, printed p.~119]{MilneMF110}.

Put $\alpha=3^{1/4}$ and $r_0=(3/4)^{1/4}$, with positive real
roots, and define on $H^{\mathrm{cyc}}$
\[
 \kappa(\sigma)=\frac{\sigma(r_0)}{r_0}
               =\frac{\sigma(\alpha)}{\alpha}\in\{1,-1\}.
\]
The equality follows from $r_0=\alpha/\sqrt2$.

\begin{proposition}\label{prop:adelic-cm-images}
The full division fields are
\[
 \mathbb Q(E_j[\mathrm{tors}])=K_j^{\mathrm{ab}}.
\]
On $H^{\mathrm{cyc}}$, one has
\begin{equation}
 \rho_4(H^{\mathrm{cyc}})=N_4(4),\qquad
 \rho_3^{\circ}(H^{\mathrm{cyc}})=N_3(4),\qquad
 \rho_3^{\circ}=\kappa\rho_3.
 \label{eq:adelic-arithmetic-projections}
\end{equation}
The representations $\rho_3^{\circ}$ and $\rho_3$ have the same
kernel on $H^{\mathrm{cyc}}$.
\end{proposition}

\begin{proof}
First determine the images over the CM fields.  Put $\pi_4=1+i$.
The eight points of $E_4[\pi_4^3]$ are rational over $K_4$:
besides the four points of $E_4[2]$, they have coordinates
\[
 (i,\pm(1-i)),\qquad(-i,\pm(1+i)).
\]
For instance, with $[i](X,Y)=(-X,iY)$, the addition law gives
$(1+i)(i,1-i)=(1,0)$, then $(1+i)(1,0)=(0,0)$, and then zero.
The other three points give the same conclusion.  The eight points
exhaust the kernel because the endomorphism has degree eight.
The four roots of unity map bijectively to
$(O_4/\pi_4^3)^{\times}$.  Hence
\[
 C_4=\{u\in\widehat O_4^{\times}:u\equiv1\pmod{\pi_4^3}\}
       \longrightarrow\widehat O_4^{\times}/\mu_4
\]
is an isomorphism.  Rationality puts $\rho_4(G_{K_4})$ inside
$C_4$, and CM reciprocity makes its image modulo $\mu_4$ full.
Thus $\rho_4(G_{K_4})=C_4$.

For $E_3$, the two-torsion abscissae are
$-1,-\omega,-\omega^2$, and
$E_3[1-\omega]=\{O,(0,1),(0,-1)\}$.  These coprime kernels
make $E_3[2\pi_3]$ rational over $K_3$.
The reduction of $\mu_6$ is an isomorphism onto
\[
 (O_3/(2\pi_3))^{\times}
              \simeq\mathbb F_4^{\times}\times\mathbb F_3^{\times}.
\]
The same argument gives
\[
 \rho_3(G_{K_3})=C_3
       =\{u\in\widehat O_3^{\times}:u\equiv1\pmod{2\pi_3}\}.
\]

The CM theorem implies that the action on torsion points modulo
geometric automorphisms already cuts out $K_j^{\mathrm{ab}}$.
Since the full torsion action is
abelian over $K_j$, the full division field over $K_j$ is the same
field.  Moreover $K_3$ occurs among the two-torsion coordinates,
and $K_4$ occurs by the Weil pairing on four-torsion.  This proves
the displayed division-field identities over $\mathbb Q$.

Complex conjugation acts on $C_j$ by $u\mapsto\bar u$ and has
order two.  Thus the closed derived subgroup of the full Galois
image is $\{u/\bar u:u\in C_j\}$; this is already closed, as
the image of a compact group.  Kronecker--Weber identifies
$H^{\mathrm{cyc}}$ with the closed derived subgroup of $H$.
Consequently
\begin{equation}
 \rho_j(H^{\mathrm{cyc}})=\{u/\bar u:u\in C_j\}.
 \label{eq:adelic-ratio-image}
\end{equation}

At split primes the ratio map on units is
$(a,b)\mapsto(a/b,b/a)$ and is onto the local norm-one group.
At an unramified quadratic prime, Hilbert~90 supplies a preimage;
multiplication by a rational uniformizer power makes it a unit.
At the inert prime two for $K_3$, a norm-one unit $v\equiv1\pmod2$
has a unit preimage whose residue is fixed by the residue-field
involution.  Its residue therefore belongs to
$\mathbb F_2^\times=\{1\}$, so the preimage is also one modulo two.
At the ramified prime three, a Hilbert--90 preimage of a norm-one
unit $v\equiv1\pmod{\pi_3}$ has even $\pi_3$-valuation:
an odd valuation contributes the residue
$\pi_3/\bar\pi_3=-1$.  Multiplying by a rational power of three
and then by a sign makes it a unit congruent to one modulo $\pi_3$.
The converses follow by reduction.  Thus
\begin{equation}
 \rho_3(H^{\mathrm{cyc}})
 =\{v\in N_3:v_2\equiv1\pmod2,\quad
                     v_3\equiv1\pmod{\pi_3}\}=:M_3.
 \label{eq:adelic-unnormalized-cubic-image}
\end{equation}

For $K_4$ at two, a unit $u=a+bi\equiv1\pmod{\pi_4^3}$ has
$b$ even, so $u/\bar u\equiv1\pmod4$.  Conversely if
$v\bar v=1$ and $v\equiv1\pmod4$, then
$b_0=(1+v)/2$ is a unit congruent to one modulo two and
$b_0/\bar b_0=v$.  The two unit classes modulo $\pi_4^3$ that
are one modulo two are $1$ and $-1$.  Hence one of $b_0,-b_0$
belongs to $1+\pi_4^3O_{4,2}$.  Together with the other local
calculations, this proves $\rho_4(H^{\mathrm{cyc}})=N_4(4)$.

We now determine the cubic normalization without an unspecified
four-division comparison.  The norm-one units modulo four in
$O_3/4$ that are one modulo two are exactly $1,-1$:
write a unit as $1+2a$, and use
$\operatorname{Tr}_{\mathbb F_4/\mathbb F_2}(a)=0$.
Let $k=\mathbb Q^{\mathrm{ab}}$.  If
$Y^2=\prod_{i=1}^3(X-e_i)$ has all its roots in $k$, the halves
of $(e_i,0)$ are obtained from
\[
 a^2=e_i-e_j,\quad b^2=e_i-e_k,\qquad
 (X,Y)=(e_i+ab,\,ab(a+b)),
\]
with the possible signs.  Conversely the ratios $Y/(X-e_i)$
for the choices $(a,b)$ and $(-a,b)$ recover $a+b$ and $b-a$,
and hence $a,b$.  Thus its four-division field over $k$ is
generated by the square roots of the root differences.
For $E_3$ every nonzero root difference is a root of unity times
$\sqrt{-3}$.  Since $k$ contains every root of unity,
\begin{equation}
 k(E_3[4])=k(3^{1/4}).
 \label{eq:adelic-cubic-four-division}
\end{equation}
This extension is nontrivial: the real field
$\mathbb Q(3^{1/4})$ misses the nonreal conjugates of its
Eisenstein polynomial $X^4-3$, so it cannot lie in an abelian
extension of $\mathbb Q$.
It follows that the modulo-four sign of $\rho_3(\sigma)$ is
exactly $\kappa(\sigma)$.

Multiplication by this sign gives an isomorphism
$M_3\to N_3(4)$.  Its inverse sends $w\in N_3(4)$ to
$\varepsilon(w_3)w$, where
$\varepsilon(w_3)\in\{1,-1\}$ is its residue modulo $\pi_3$.
Indeed this makes the three-component one modulo $\pi_3$, while
the modulo-four sign at two becomes $\varepsilon(w_3)$.
The two maps are inverse homomorphisms.  This proves both the
normalized cubic image and the equality of kernels.
\end{proof}

\begin{lemma}\label{lem:adelic-cm-entanglement}
The intersection of the two full division fields is
\begin{equation}
 K_3^{\mathrm{ab}}\cap K_4^{\mathrm{ab}}
       =\mathbb Q^{\mathrm{ab}}(3^{1/4}).
 \label{eq:adelic-cm-intersection}
\end{equation}
In the representations of Proposition~\ref{prop:adelic-cm-images},
the joint image of $H^{\mathrm{cyc}}$ is exactly $\mathcal E$.
\end{lemma}

\begin{proof}
Put $L=K_3^{\mathrm{ab}}\cap K_4^{\mathrm{ab}}$.
It contains $\mathbb Q^{\mathrm{ab}}$ and is Galois over
$\mathbb Q$.  In $J=\operatorname{Gal}(L/\mathbb Q)$, the two
distinct index-two subgroups
$J_3=\operatorname{Gal}(L/K_3)$ and
$J_4=\operatorname{Gal}(L/K_4)$ are abelian.
Their intersection $Z$ is central, since $J_3,J_4$ generate $J$.
Choose $a\in J_3\setminus Z$ and $b\in J_4\setminus Z$.
Then $J=\langle Z,a,b\rangle$, $a^2,b^2\in Z$, and
$J'$ is generated by the central commutator $[a,b]$.
The identity $[a^2,b]=[a,b]^2=1$ makes $|J'|\leq2$.
Since $L$ contains the maximal abelian extension of $\mathbb Q$,
this proves $[L:\mathbb Q^{\mathrm{ab}}]\leq2$.

The splitting field $\mathbb Q(i,\alpha)$ of $X^4-3$ has dihedral
group of order eight.  Its subgroup over $\mathbb Q(i)$ is cyclic
of order four; its subgroup over $\mathbb Q(\sqrt{-3})$ is a
four-group.  It is therefore contained in both $K_j^{\mathrm{ab}}$.
The nontrivial extension in
\eqref{eq:adelic-cubic-four-division} gives the reverse bound and
proves \eqref{eq:adelic-cm-intersection}.

It remains to identify the common quadratic quotient in the stated
matrix coordinates.  By
\eqref{eq:adelic-unnormalized-cubic-image},
$\rho_3(\sigma)$ is one modulo $\pi_3$; hence
\[
 (-1)^{\epsilon_3(\rho_3^{\circ}(\sigma))}
                      =\kappa(\sigma).
\]
For $E_4$ its three-division polynomial is
$3X^4-6X^2-1$.  Choose a root $x$ with
$x^2=1+2\sqrt3/3$.  Then
\[
 \frac{x^2}{\sqrt3}=\frac{2+\sqrt3}{3}
     =\left(\frac{\sqrt6+\sqrt2}{2\sqrt3}\right)^2.
\]
Thus $3^{1/4}\in k(E_4[3])$.  The latter extension over $k$ is
cyclic of order four, since the reduction of $N_4(4)$ modulo
three is $\langle S\rangle$.  Its unique quadratic subextension
is consequently $k(3^{1/4})$, and
\[
 (-1)^{\epsilon_4(\rho_4(\sigma))}=\kappa(\sigma).
\]
The joint image of two Galois representations is the fibre product
of their separate images over the Galois group of the intersection
of their fields.  The separate images are full by
Proposition~\ref{prop:adelic-cm-images}; the intersection has order
two and the two quotient maps have just been computed.  The fibre
product is exactly \eqref{eq:adelic-equal-parity}.
\end{proof}

\subsection{A signed comparison on the Legendre family}

\begin{proposition}\label{prop:adelic-signed-path-comparison}
Choose the Tate-module bases by transport along paths to the
harmonic and equianharmonic fibres of the Legendre family.  For
$\sigma\in H^{\mathrm{cyc}}$ one then has the exact integral
identities
\begin{equation}
 W_\sigma=\rho_4(\sigma),\qquad
 T_\sigma=\rho_3^{\circ}(\sigma)
                       =\kappa(\sigma)\rho_3(\sigma).
 \label{eq:adelic-signed-path-comparison}
\end{equation}
In particular $(T,W)(H^{\mathrm{cyc}})=\mathcal E$.
\end{proposition}

\begin{proof}
Let $Y=\mathbb P^1\setminus\{0,1,\infty\}$ and consider the
Legendre family
\[
 \mathcal L_t:\quad Y_1^2=X_1(X_1-1)(X_1-t)
\]
with tangential base point at $t=0$ towards $1$.
Choose the integral vanishing-cycle basis in the peripheral
convention fixed by $x,y$.  The two local monodromies are
$U^2,L^2$: a loop about either collision of two ordered branch
points lifts to the square of the corresponding Dehn twist, and
the two vanishing cycles have intersection of absolute value one.
This identifies the geometric Tate-module representation with $q$
on $F$, at every prime.

At zero the special cubic is $Y_1^2=X_1^2(X_1-1)$, whose two
branches at the node have slopes $i,-i$.  Over
$\mathbb Q(i)((t))$ the reduction is split multiplicative.
Tate uniformization supplies a primitive cyclotomic line and a
free rank-one quotient in every integral Tate module; the only
additional character over $\mathbb Q((t))$ is the splitting
character of $\mathbb Q(i)$.
For an explicit uniformization over an arbitrary complete
nonarchimedean field, including the Galois compatibility and the
split multiplicative criterion, see
\cite[\S2, Theorem~2.2 and pp.~3--4]{PapikianTate}.
The integral filtration follows directly: in
$\overline K^{\times}/Q^{\mathbb Z}$, the $n$-torsion has kernel
$\mu_n$ and quotient $\mathbb Z/n$ under
$[z]\mapsto m\bmod n$ for $z^n=Q^m$.
The characteristic is zero, so this argument applies to every
prime, including two.

Both characters on the two graded pieces are trivial on
$H^{\mathrm{cyc}}$.  The tangential arithmetic matrix
$V_\sigma$ is therefore upper unipotent with both diagonal
entries one.  Functoriality gives
$q\Phi_\sigma=\operatorname{Ad}(V_\sigma)q$ on $F$.
Equation~\eqref{eq:adelic-frame-intertwiner} gives the same
covariance for $D_\sigma$.
An upper unipotent matrix commuting with $U^2$ and $L^2$ is the
identity: the second commutation gives twice its upper-right
entry equal to zero in each $\mathbb Z_\ell$.
Consequently
\begin{equation}
 V_\sigma=D_\sigma.
 \label{eq:adelic-tangential-matrix}
\end{equation}
This uniqueness uses the primitive integral filtration, so it
does not leave a sign undetermined at the prime two.

The geometric fundamental group of the quotient orbifold
$[Y/S_3]$ is $\Gamma$, with pure subgroup $F$.
The arithmetic action induced by the image of the chosen
tangential base point restricts to the canonical Ihara action
on $F$.  It is exactly $\Phi_\sigma$ on $\Gamma$.
Indeed two automorphisms of $\Gamma$ agreeing on $F$ differ,
on each group element, by an element of $C_\Gamma(F)$, and this
centralizer is trivial.  For an explicit argument, if $z$ centralizes
$F$, then $z^6\in F$ and $z^6$ lies in the centre of the nonabelian
free profinite group $F$, which is trivial.  Thus $z$ has finite
order.  A nontrivial torsion element of $\Gamma$ is conjugate into
one of its two finite factors, and its centralizer lies in that
factor \cite{HerfortRibes1985}.  Such a finite centralizer cannot
contain $F$, so $z=1$.

The transposition $t\mapsto1-t$ fixes $p_4=1/2$.
Choose an ordinary geometric path $\xi_4$ in $Y$ from the
tangential base point to $p_4$ so that its transported orbifold
isotropy is $s$.  Such a choice exists: every order-two subgroup
is conjugate into the order-two factor, and a conjugator preserving
the specified transposition in $S_3$ belongs to $F\langle s\rangle$.
It can therefore be replaced by a pure conjugator.
For the order-three isotropy use one of
$p_3=(1\pm\sqrt{-3})/2$ and choose an ordinary path $\xi_3$
whose transported positive generator is $r$.
The two points give the two possible pure conjugacy classes;
choose the one corresponding to $r$.  The same conjugator
argument, with centralizer $\langle r\rangle$ in $S_3$, then
makes the path adjustment pure.  Both chosen points and their
coordinate-permutation isotropies are fixed by $H^{\mathrm{cyc}}$.

Paths are composed as maps of fibre functors.  If $\xi$ is either
chosen path, set
\[
 c_\sigma=(\sigma\xi)^{-1}\xi\in F.
\]
For the isotropy $s_p$ at its endpoint, the identity
$s=\xi^{-1}s_p\xi$ gives
\[
 \Phi_\sigma(s)=(\sigma\xi)^{-1}s_p(\sigma\xi)
                        =c_\sigma s c_\sigma^{-1}.
\]
The analogous formula holds for $r$.  Uniqueness of the pure
frames yields $c_\sigma=\beta_\sigma$ or $A_\sigma$.

Now use the Tate sheaf on $Y$, where the Legendre family itself
is defined.  If $\rho_p(\sigma)$ is its matrix at the ordinary
point $p$, functoriality of parallel transport gives
\[
 q(\sigma\xi)=\rho_p(\sigma)q(\xi)D_\sigma^{-1}.
\]
In the transported basis put
$M_p(\sigma)=q(\xi)^{-1}\rho_p(\sigma)q(\xi)$.
Taking the inverse in the preceding formula and multiplying by
$q(\xi)$ proves
\begin{equation}
 q(c_\sigma)=D_\sigma M_p(\sigma)^{-1},\qquad
 q(c_\sigma)^{-1}D_\sigma=M_p(\sigma).
 \label{eq:adelic-path-order}
\end{equation}
The signed representation here is the one already fixed on $F$.
The calculation applies on every finite torsion sheaf, with
compatible transition maps.  It therefore proves an integral
adelic identity.  No descent of the signed Tate sheaf to
$[Y/S_3]$ is needed.

It remains to identify the two ordinary fibres.  At $p_4=1/2$,
the substitution $X_1=X+1/2$ gives
$Y_1^2=X^3-X/4$.  This is the quadratic twist of $E_4$ by
$1/2$, whose square root lies in $\mathbb Q^{\mathrm{ab}}$.
Its representation on $H^{\mathrm{cyc}}$ is therefore $\rho_4$.
At $p_3=(1+\sqrt{-3})/2$, substituting
$X_1=X+(1+p_3)/3$ gives
\[
 Y_1^2=X^3+\frac{\sqrt{-3}}9
             =X^3+\left(\frac1{\sqrt{-3}}\right)^3.
\]
This is the quadratic twist of $E_3$ by
$d=1/\sqrt{-3}$.  The quotient of a square root of $d$ by $r_0$
lies in $\mathbb Q^{\mathrm{ab}}$; its square is $-2i/3$, and
one may take $\sqrt{2/3}\,\zeta_8^{-1}$.
Its twisting character on $H^{\mathrm{cyc}}$ is consequently
$\kappa$.  The conjugate equianharmonic point changes $d$ by
$-1$ and gives the same character.

The transported CM algebras are the centralizers of $S$ and $R$
already computed in Lemma~\ref{lem:adelic-corrected-frames}.
Indeed Proposition~\ref{prop:adelic-cm-images} makes each local
arithmetic image open in its one-dimensional CM torus, so its
commutant over $\mathbb Q_\ell$ is the corresponding quadratic
algebra.  The matrices $S$ and $R$, which commute with that image,
generate its full quadratic integer ring.
At a prime, interchanging the two CM embeddings changes a unit
to its conjugate; it preserves the norm-one groups, their level-four
subgroups, and both parity maps.  Thus the CM image calculations
apply in these transported integral bases.
Equation~\eqref{eq:adelic-path-order} now gives
\eqref{eq:adelic-signed-path-comparison}, and
Lemma~\ref{lem:adelic-cm-entanglement} identifies the joint image.
\end{proof}

\subsection{Eliminating the entire adelic frame contribution}

Recall the global quadratic character $\Delta:\mathcal G\to
\mathbb F_2$ of Section~\ref{sec:global-frame-character}.
It is zero on $H$ and on $G$; on
$N_0=\ker(\lambda,\rho_2)$ it equals $a_3+a_2$.
The vanishing on $G$ uses the normalized frame theorem
\cite[Theorem~3]{H17}.  We require no stronger arithmetic input.

\begin{theorem}\label{thm:adelic-frame-elimination}
Put $G_1=G\cap M_0$.  Then
\begin{equation}
 (T,W)(G_1)=(T,W)(H^{\mathrm{cyc}})=\mathcal E.
 \label{eq:adelic-gt-arithmetic-image}
\end{equation}
In particular, the closed normal subgroup
\begin{equation}
 C_{\mathrm{CM}}=
 \ker\bigl((T,W):G_1\longrightarrow N_3(4)\times N_4(4)\bigr)
 \label{eq:adelic-cm-kernel}
\end{equation}
satisfies
\begin{equation}
 G=C_{\mathrm{CM}}H,\qquad
 C_{\mathrm{CM}}\cap H=
 G_{K_3^{\mathrm{ab}}K_4^{\mathrm{ab}}}.
 \label{eq:adelic-cm-normalization}
\end{equation}
The inclusion of $H$ induces the canonical isomorphism
\begin{equation}
 G/C_{\mathrm{CM}}\simeq
 \operatorname{Gal}(K_3^{\mathrm{ab}}K_4^{\mathrm{ab}}/\mathbb Q).
 \label{eq:adelic-cm-arithmetic-quotient}
\end{equation}
Every continuous abelian character of $G$ that is zero on $H$
and whose restriction to $G_1$ factors through $(T,W)$ is zero.
\end{theorem}

\begin{proof}
On $N_0$ the matrix $D_g$ is the identity modulo twelve.
The reductions of $q(A_g)$ and $q(\beta_g)$ therefore give
\[
 \epsilon_3(T_g)=a_3(g),\qquad
 \epsilon_4(W_g)=a_2(g).
\]
For the cubic group the first equality is reduction modulo
$\pi_3$: $R$ reduces to one and $-I$ to minus one.
For the harmonic group the second equality is exponent parity
in $\langle S\rangle$.  Inverting a frame does not change either
parity.  Thus
$(\epsilon_3\circ T)+ (\epsilon_4\circ W)=\Delta$ on $N_0$.

The arithmetic normalization by the Kummer cocycle of two gives
$M_0=N_0H^{\mathrm{cyc}}$: its value on $M_0$ is even, and its
image on $H^{\mathrm{cyc}}$ is exactly
$2\widehat{\mathbb Z}$.  On $H^{\mathrm{cyc}}$ both the sum of
the two parities and $\Delta$ are zero, by
Proposition~\ref{prop:adelic-signed-path-comparison} and the
arithmetic vanishing of $\Delta$.
All these maps are homomorphisms on $M_0$, so
\begin{equation}
 (\epsilon_3\circ T)+ (\epsilon_4\circ W)=\Delta
                          \quad\hbox{on }M_0.
 \label{eq:adelic-parity-delta}
\end{equation}
Since $\Delta|_G=0$, the image of $G_1$ lies in $\mathcal E$.
It contains the image of $H^{\mathrm{cyc}}$, which is already
all of $\mathcal E$.  This proves
\eqref{eq:adelic-gt-arithmetic-image}.

For $g\in G$, first choose $\sigma\in H$ with
$\lambda_\sigma=\lambda_g$.  Then $g\sigma^{-1}\in G_1$.
By \eqref{eq:adelic-gt-arithmetic-image}, choose
$h\in H^{\mathrm{cyc}}$ with the same corrected frame pair.
It follows that $g\sigma^{-1}h^{-1}\in C_{\mathrm{CM}}$, proving
$G=C_{\mathrm{CM}}H$.
The intersection formula follows from the exact signed comparison,
the equality of the normalized and unnormalized cubic kernels,
and the two full division-field identities.
Normality follows from Lemma~\ref{lem:adelic-kernel-normality}.
The compact surjection $H\to G/C_{\mathrm{CM}}$ has exactly the
kernel in \eqref{eq:adelic-cm-normalization}, proving
\eqref{eq:adelic-cm-arithmetic-quotient}.

Finally let $\chi$ be a character as in the statement.
Its restriction to $G_1$ is a character of the image
$\mathcal E$, and it is zero there because
$H^{\mathrm{cyc}}$ maps onto that image.  Thus $\chi|_{G_1}=0$.
Matching a multiplier by an element of $H$ then gives $\chi=0$
on $G$.
\end{proof}

In particular the same relative abelian quotient can be computed as
\[
 B\simeq C_{\mathrm{CM}}\big/
       \overline{[G,C_{\mathrm{CM}}]\,
                 G_{K_3^{\mathrm{ab}}K_4^{\mathrm{ab}}}}.
\]
For any prime $p$, its continuous $\mathbb F_p$-characters are the
$G$-invariant characters of $C_{\mathrm{CM}}$ that vanish on the
displayed Galois subgroup.  This follows from the general extension
criterion for a decomposition $G=NH$ with $N$ closed and normal.

This normalization concerns the full pair of CM frames.  Combining
it with simultaneous rational-radical normalization would require
the joint arithmetic image of the CM and radical representations;
the two separate surjectivity statements do not imply that joint
compatibility.

\section{The exact intersection of the radical and CM fields}
\label{sec:radical-cm-entanglement}

Put $k=\mathbb Q^{\mathrm{ab}}=\mathbb Q(\mu_\infty)$ and
\[
 K_3=\mathbb Q(\sqrt{-3}),\qquad K_4=\mathbb Q(i),\qquad
 F_{12}=K_3K_4=\mathbb Q(\zeta_{12}).
\]
Let $\mathcal P$ be the set of positive rational primes, and set
\[
 L_{\mathrm{rad}}=k(q^{1/n}:q\in\mathcal P,\ n\geq1),\qquad
 L_{\mathrm{CM}}=K_3^{\mathrm{ab}}K_4^{\mathrm{ab}}.
\]
All unqualified radicals are the positive real ones.  The presence
of all roots of unity makes these fields independent of that choice.

\begin{theorem}\label{thm:radical-cm-intersection}
One has
\begin{align}
 L_{\mathrm{rad}}\cap L_{\mathrm{CM}}
 &=L_{\mathrm{rad}}\cap F_{12}^{\mathrm{ab}}\notag\\
 &=k\bigl(q^{1/12}:q\in\mathcal P,\,
                         2^{1/8},3^{1/8}\bigr)=:I.
 \label{eq:radical-cm-intersection}
\end{align}
In particular
\begin{equation}
 \operatorname{Gal}(I/k)\simeq
 (\mathbb Z/12)^2\times
       \prod_{q\in\mathcal P\setminus\{2,3\}}\mathbb Z/6.
 \label{eq:radical-cm-common-group}
\end{equation}
The first two factors are indexed by $2,3$.  In these coordinates
conjugation by an element of $G_{\mathbb Q}$ is multiplication by
its cyclotomic character, reduced modulo the indicated orders.
\end{theorem}

\begin{proof}
We first compute the entire maximal subextension of $L_{\mathrm{rad}}$
that is abelian over $F_{12}$.  Recall the simultaneous radical
description
\[
 A_{\mathrm{rad}}=\operatorname{Gal}(L_{\mathrm{rad}}/\mathbb Q)
 =
 \left\{(u,(t_q)_q):
 \begin{array}{l}
 u\in\widehat{\mathbb Z}^{\times},\
 t_q\in\widehat{\mathbb Z},\\
 t_q\bmod2=\epsilon_q(u)\quad(q\in\mathcal P)
 \end{array}\right\},
\]
with multiplication $(u,t)(v,s)=(uv,t+us)$.
Here $\epsilon_q$ is the quadratic character of $\mathbb Q(\sqrt q)$,
written additively.  The translation subgroup over $k$ is
$T=\prod_q2\widehat{\mathbb Z}$.
This is the classical radical Galois description
\cite[\S2, Theorem~2.5]{MoreeStevenhagen}; the passage from finite
sets of primes to this product follows by compactness.
It also follows directly from the simultaneous Kummer-image lemma:
a positive rational radical lying in an abelian extension of
$\mathbb Q$ has its square in $\mathbb Q$, which is exactly the
displayed parity condition.

Write
\[
 U_{12}=\{u\in\widehat{\mathbb Z}^{\times}:u\equiv1\pmod{12}\},
 \qquad A_{12}=\operatorname{Gal}(L_{\mathrm{rad}}/F_{12}).
\]
Thus $A_{12}$ is the inverse image of $U_{12}$ in $A_{\mathrm{rad}}$.
Let $C_{12}=\overline{[A_{12},A_{12}]}$.  Since the multiplier quotient is
abelian, $C_{12}\subset T$.  For the convention
$[a,b]=aba^{-1}b^{-1}$, direct multiplication gives
\begin{equation}
 [(u,t),(v,s)]=
 \bigl(1,(1-v)t+(u-1)s\bigr).
 \label{eq:radical-cm-affine-commutator}
\end{equation}
We claim the exact simultaneous equality
\begin{equation}
 C_{12}=(24\widehat{\mathbb Z})_{\{2,3\}}\times
             \prod_{q\ne2,3}12\widehat{\mathbb Z}.
 \label{eq:radical-cm-derived-subgroup}
\end{equation}
The subscripts $\{2,3\}$ here index the rational-prime coordinates,
not their local components.

We make the distinction between these coordinates and the local
prime $\ell$ explicit in the proof.  Every closed subgroup of the
profinite abelian group $T$ is stable under
$\widehat{\mathbb Z}$-multiplication: approximate an element of
$\widehat{\mathbb Z}$ by integers and use closedness.
The idempotents of $\widehat{\mathbb Z}=\prod_\ell\mathbb Z_\ell$
therefore project $C_{12}$ onto its local components inside $C_{12}$ itself.
Consequently it suffices to determine each simultaneous local
subgroup $C_{12,\ell}$; there can be no additional entanglement between
distinct local primes.

Commuting a lift of $u\in U_{12}$ with a translation gives
$(u-1)T$.  At $\ell\geq5$ one may choose $u_\ell=2$ and all
other local components one.  Hence
$C_{12,\ell}=\prod_q\mathbb Z_\ell$.
At $\ell=3$, choosing $u_3=4$ gives
$\prod_q3\mathbb Z_3\subset C_{12,3}$.
Both multipliers in \eqref{eq:radical-cm-affine-commutator} are one
modulo three, proving the reverse inclusion.  Thus
$C_{12,3}=\prod_q3\mathbb Z_3$.
At $\ell=2$, choosing $u_2=5$ and using translations in
$\prod_q2\mathbb Z_2$ gives
\[
 \prod_q8\mathbb Z_2\subset C_{12,2}.
\]

It remains to determine the image modulo eight simultaneously
in every coordinate.  For $u\in U_{12}$ put
$a(u)=(u_2-1)/4\bmod2$.  Since $u_2\equiv1\pmod4$,
\[
 a(u)=\epsilon_2(u).
\]
Also $\epsilon_3(u)=0$, since $\sqrt3\in F_{12}$.
Reducing \eqref{eq:radical-cm-affine-commutator} modulo eight gives
in coordinate $q$
\begin{equation}
 4\bigl(\epsilon_2(u)\epsilon_q(v)
               -\epsilon_2(v)\epsilon_q(u)\bigr).
 \label{eq:radical-cm-mod-eight}
\end{equation}
It is zero for $q=2,3$, and is a multiple of four for every
other $q$.  This proves the necessary simultaneous upper bound.

For the reverse inclusion, the characters
$\epsilon_q|_{U_{12}}$ with $q\ne3$ are jointly independent.
Indeed a finite product would be trivial on $U_{12}$ precisely if
$\mathbb Q(\sqrt a)\subset F_{12}$, where $a$ is the corresponding
positive squarefree product of primes.  The three quadratic
subfields of $F_{12}$ are
$\mathbb Q(i)$, $\mathbb Q(\sqrt{-3})$, and $\mathbb Q(\sqrt3)$;
among positive squareclasses only $a=3$ occurs.  Excluding $3$
therefore gives independence.  Every finite evaluation map onto
the corresponding product of $\mathbb F_2$ is surjective.
Compactness then gives surjectivity onto the entire product.

Fix $u\in U_{12}$ with $\epsilon_2(u)=1$.  For an arbitrary family
$(b_q)_{q\ne2,3}\in\prod_{q\ne2,3}\mathbb F_2$, choose
$v\in U_{12}$ with
\[
 \epsilon_2(v)=0,\qquad \epsilon_q(v)=b_q\quad(q\ne2,3).
\]
Choose any lifts of $u,v$ to $A_{12}$.  Formula
\eqref{eq:radical-cm-mod-eight} now gives exactly the family
$4b_q$ in all the other coordinates, with zero at $2,3$.
This proves surjectivity onto the whole required product modulo
eight, not merely onto its individual projections.  Together
with $\prod_q8\mathbb Z_2\subset C_{12,2}$, it proves
\[
 C_{12,2}=(8\mathbb Z_2)_{\{2,3\}}\times
                       \prod_{q\ne2,3}4\mathbb Z_2.
\]
Combining the local equalities proves
\eqref{eq:radical-cm-derived-subgroup}.

The fixed field of $C_{12}$ is $L_{\mathrm{rad}}\cap F_{12}^{\mathrm{ab}}$.
In the radical coordinates, the subgroup
\eqref{eq:radical-cm-derived-subgroup} fixes precisely the
$24$th roots of $2,3$ and the $12$th roots of all other primes.
Since a $24$th root is generated by its eighth and twelfth powers
as root orders, this is the field $I$ in
\eqref{eq:radical-cm-intersection}.  More explicitly,
$q^{1/24}=(q^{1/8})/(q^{1/12})$.
The quotient of $2\widehat{\mathbb Z}$ by
$24\widehat{\mathbb Z}$ is $\mathbb Z/12$, and by
$12\widehat{\mathbb Z}$ is $\mathbb Z/6$.  This proves
\eqref{eq:radical-cm-common-group} and its conjugation action.

Each $K_j^{\mathrm{ab}}/F_{12}$ is abelian, so
$L_{\mathrm{CM}}\subset F_{12}^{\mathrm{ab}}$.  We have therefore
proved $L_{\mathrm{rad}}\cap L_{\mathrm{CM}}\subset I$.
For the reverse inclusion, all rational cube roots lie in
$K_3^{\mathrm{ab}}$, since $K_3$ contains $\mu_3$.
All rational fourth roots lie in $K_4^{\mathrm{ab}}$, since
$K_4$ contains $\mu_4$.

Let $\alpha=2^{1/8}$.  The field $\mathbb Q(i,\alpha)$ contains
\[
 \zeta_8=\frac{1+i}{\alpha^4},
\]
so it is the splitting field of $X^8-2$.  An automorphism fixing
$i$ sends $\alpha$ to $\zeta_8^a\alpha$ and sends $\zeta_8$ to
$\zeta_8^{1+4a}$.  Composition on the index $a\bmod8$ is
$a*b=a+b+4ab$, which is commutative.  Thus this extension is
abelian over $K_4$, and $2^{1/8}\in K_4^{\mathrm{ab}}$.

For the remaining radical put
$v=(2+\sqrt3)^{1/4}$ and $M=\mathbb Q(i,v)$.
The field $M$ contains $F_{12}$ and all conjugates
$i^a v^{\pm1}$, so it is Galois over $\mathbb Q$.
Every automorphism of $M$ fixing $\sqrt{-3}=i\sqrt3$ has the form
\[
 i\longmapsto i^\epsilon,\qquad
 v\longmapsto i^a v^\epsilon,
 \qquad a\in\mathbb Z/4,\quad\epsilon\in\{1,-1\}.
\]
Indeed the sign on $\sqrt3=v^4-2$ must equal the sign on $i$.
The product of two such automorphisms has exponent
$a\eta+\epsilon b$ and sign $\epsilon\eta$; this rule is
symmetric in $(a,\epsilon)$ and $(b,\eta)$.
Therefore $M/K_3$ is abelian and $v\in K_3^{\mathrm{ab}}$.

On the CM elliptic curve $E_4:Y^2=X^3-X$, the three-division
polynomial is $3X^4-6X^2-1$.  Choose a three-torsion point $(x,y)$
with $x^2=1+2\sqrt3/3$.  The curve equation gives
\[
 y^4=x^2(x^2-1)^2
       =\frac49\sqrt3(2+\sqrt3).
\]
All its torsion coordinates are abelian over $K_4$, since the
endomorphism by $i$ is defined there and the Tate-module action
commutes with that rank-one CM action
\cite[Proposition~12.13, printed p.~119]{MilneMF110}.
Thus $y\in K_4^{\mathrm{ab}}$.
As $\sqrt{2/3}\in k$, the element
\[
 z=\frac{y}{\sqrt{2/3}\,v}\in L_{\mathrm{CM}}
\]
satisfies $z^4=\sqrt3$ and hence $z^8=3$.
All eighth roots of unity already lie in $k$, so
$3^{1/8}\in L_{\mathrm{CM}}$.
This supplies every generator of $I$ and completes the proof.
\end{proof}

\subsection{The remaining compatibility condition on GT}

The arithmetic intersection theorem by itself does not combine
the two separate GT normalizations.  Here is the exact additional
condition it leaves.

Suppose the separate arithmetic quotient identifications are fixed:
\[
 G/N_{\mathrm{rad}}\simeq
        A_{\mathrm{rad}}=\operatorname{Gal}(L_{\mathrm{rad}}/\mathbb Q),
 \qquad
 G/C_{\mathrm{CM}}\simeq
        A_{\mathrm{CM}}=\operatorname{Gal}(L_{\mathrm{CM}}/\mathbb Q),
\]
with both identifications induced by the inclusion
$H=G_{\mathbb Q}\subset G$.
Restricting these quotient maps to $I$ gives two homomorphisms
$a,b:G\to A_I=\operatorname{Gal}(I/\mathbb Q)$, which agree on
$H$ and induce the same multiplier.  Their difference
\begin{equation}
 d(g)=a(g)b(g)^{-1}\in C=\operatorname{Gal}(I/k)
 \label{eq:radical-cm-compatibility-defect}
\end{equation}
is a continuous crossed homomorphism:
\[
 d(gh)=d(g)+\lambda_gd(h),\qquad d|_H=0.
\]
The addition uses the coordinates
\eqref{eq:radical-cm-common-group}.
The formula follows by writing $a(g)=d(g)b(g)$ and using the
scalar cyclotomic conjugation on $C$.

\begin{proposition}\label{prop:radical-cm-joint-gate}
The following are equivalent:
\begin{enumerate}
\item $G=(N_{\mathrm{rad}}\cap C_{\mathrm{CM}})H$;
\item the joint radical and CM image of $G$ equals that of $H$;
\item $a=b$, or equivalently $d=0$.
\end{enumerate}
Thus the remaining compatibility has exponent dividing twelve
over $k$: it consists of one cubic coordinate for every rational
prime, one quadratic coordinate for each prime other than $2,3$,
and one coordinate modulo four for each of $2,3$.
\end{proposition}

\begin{proof}
The first two assertions are equivalent by taking the common
kernel of the two quotient maps.
The joint Galois image is exactly
$A_{\mathrm{rad}}\times_{A_I}A_{\mathrm{CM}}$,
by the intersection theorem and the Galois correspondence.
The joint image of $G$ contains it, since $H\subset G$.
It is contained in it precisely when its two restrictions to
$A_I$ coincide.  This proves the equivalences.  The final list
is the primary decomposition of
\eqref{eq:radical-cm-common-group}.
\end{proof}

For example, on $G_1=\ker\lambda$, write the coordinate of the
CM restriction to $I$ as $\eta_q(g)$ in $\mathbb Z/n_q$, where
$n_2=n_3=12$ and $n_q=6$ for $q\ne2,3$.
Normalize it arithmetically by
$\eta_q(\sigma)=\kappa_q(\sigma)/2\bmod n_q$ for
$\sigma\in G_{\mathbb Q^{\mathrm{ab}}}$.
Then the required coordinate identities are precisely
\[
 \frac{\rho_q(g)}2=\eta_q(g)\pmod{n_q}
 \qquad(q\in\mathcal P,\ g\in G_1).
\]
They compare the rational Kummer cocycles with finite quotients
of the CM frames.  Separate surjectivity of the two quotient maps
does not prove these identities.

\section{Compatibility of the complete common coordinate of two}
\label{sec:eighth-two-compatibility}

Both the eighth-root and the cubic coordinates of $2$ in the
radical--CM intersection are compatible with the corrected CM
frames.  These assertions hold on the multiplier-one gentle group.
We first identify the harmonic torsion level by an exact conductor
computation and prove its congruence directly from the cycle
equations.  An abelianized calculation in the modular group then
gives the cubic coordinate.

\subsection{The relevant Gaussian ray fields}

Put $k=\mathbb Q^{\mathrm{ab}}$, $K_4=\mathbb Q(i)$,
$\mathfrak p=(1+i)$, and $E_4:Y^2=X^3-X$.

\begin{lemma}\label{lem:eighth-two-conductor}
For $n=2,4,8$, the field
$L_n=\mathbb Q(i,2^{1/n})$ is cyclic of degree $n$ over $K_4$.
Its conductor is respectively
$\mathfrak p^4,\mathfrak p^6,\mathfrak p^8$.
In particular
\[
 L_4\subset K_4(E_4[8]),\qquad
 L_8\subset K_4(E_4[16]).
\]
\end{lemma}

\begin{proof}
Write $n=2^r$, $\alpha=2^{1/n}$, and $F_n=\mathbb Q(\alpha)$.
For $n=2,4$, cyclicity follows because $K_4$ contains the required
roots of unity.  For $n=8$ the field $L_8$ contains
$\zeta_8=(1+i)/\alpha^4$, and hence is the splitting field of
$X^8-2$.  Automorphisms fixing $i$ are indexed by
$a\in\mathbb Z/8$, with $\alpha\mapsto\zeta_8^a\alpha$.
Their composition is
$a*b=a+b+4ab$ modulo eight.  The element indexed by $1$ has
order eight, so the group is cyclic.  The degrees are $n$:
$X^n-2$ is Eisenstein over $\mathbb Q$ and $F_n$ is real,
so $[F_n(i):K_4]=n$.

We compute the discriminants without assuming a ray conductor.
The polynomial $X^n-2$ is Eisenstein at two, and $\alpha$ is a
uniformizer of its totally ramified local extension.  Therefore
the local integer ring is $\mathbb Z_2[\alpha]$.
The polynomial discriminant is supported only at two, so
$\mathcal O_{F_n}=\mathbb Z[\alpha]$ globally, and
\begin{equation}
 |\operatorname{disc}(F_n/\mathbb Q)|
       =n^n2^{n-1}=2^{rn+n-1}.
 \label{eq:eighth-two-pure-discriminant}
\end{equation}

In the completion of $F_n$ at its unique prime over two, use
the valuation with $v(\alpha)=1$.  Define
\[
 a_2=1+\alpha,\qquad
 a_4=1+\alpha^2+\alpha^3,\qquad
 a_8=1+\alpha^4+\alpha^6+\alpha^7,
\]
where in each expression $\alpha^n=2$ for its indicated $n$.
Direct expansion gives
\[
 a_n^2+1\equiv\alpha^{2n-1}\pmod{\alpha^{2n}},
 \qquad v(a_n^2+1)=2n-1.
\]
The element
$b_n=(i+a_n)/\alpha^{n-1}$ satisfies
\begin{equation}
 b_n^2-a_n\alpha\,b_n+
              \frac{a_n^2+1}{\alpha^{2n-2}}=0.
 \label{eq:eighth-two-eisenstein}
\end{equation}
Both nonleading coefficients have valuation one, so this is an
Eisenstein quadratic polynomial.  It generates the relative local
integer ring.  Its discriminant has valuation two: the square
of its linear coefficient has valuation two, whereas four times
the constant coefficient has valuation $2n+1$.
Thus the relative discriminant of $L_n/F_n$ has exponent two
at the unique prime over two and has no odd prime factor.
The discriminant tower formula and
\eqref{eq:eighth-two-pure-discriminant} give
\[
 |\operatorname{disc}(L_n/\mathbb Q)|
       =2^{\,2(rn+n-1)+2}=2^{\,2n(r+1)}.
\]
Since $|\operatorname{disc}(K_4/\mathbb Q)|=4$ and
$N_{K_4/\mathbb Q}(\mathfrak p)=2$, a second application of
the tower formula yields
\begin{equation}
 \operatorname{disc}(L_n/K_4)=\mathfrak p^{\,2nr}.
 \label{eq:eighth-two-gaussian-discriminant}
\end{equation}

Let $c_2,c_4,c_8$ be the conductor exponents of faithful characters
of the three cyclic extensions.  Characters of the same order
have the same kernel and hence the same conductor.  The cyclic
tower $L_2\subset L_4\subset L_8$ and the conductor--discriminant
formula \cite[V, Theorem~3.27 and Remark~3.28, p.~165]{MilneCFT}
give
\[
 c_2=4,\qquad c_2+2c_4=16,\qquad
 c_2+2c_4+4c_8=48.
\]
Consequently $(c_2,c_4,c_8)=(4,6,8)$.
The conductors of the extensions themselves are the same as
those of their faithful characters.

The torsion reciprocity theorem for $E_4$ identifies its
torsion points modulo geometric automorphisms with the
corresponding Gaussian ray class fields
\cite[Theorem~12.27, printed p.~124]{MilneMF110}.
Thus the ray class field modulo $(8)=\mathfrak p^6$ is contained
in $K_4(E_4[8])$, and that modulo $(16)=\mathfrak p^8$ is
contained in $K_4(E_4[16])$.  The conductors just computed
give the two asserted inclusions.
\end{proof}

\begin{proposition}\label{prop:eighth-two-division-fields}
Over $k=\mathbb Q^{\mathrm{ab}}$ one has
\[
 k(E_4[8])=k(2^{1/4}),\qquad
 k(E_4[16])=k(2^{1/8}).
\]
\end{proposition}

\begin{proof}
The arithmetic CM image on $G_k$ is
$N_4(4)=\{w\in\mathbb Z_2[i]^\times:
w\bar w=1,\ w\equiv1\pmod4\}$ at the prime two.
This is the image calculated in the adelic frame comparison.
Write $w=a+bi$, where $a=1+4s$ and $b=4t$.
The norm equation gives
$s+2(s^2+t^2)=0$.  In particular $s$ is even and
$s\equiv2t\pmod4$, so
\begin{equation}
 w\equiv1+8t+4ti\pmod{16}.
 \label{eq:eighth-two-norm-reduction}
\end{equation}
Every class $t\bmod4$ occurs: the norm-one unit
$(1+2ti)/(1-2ti)$ has the required imaginary coefficient modulo
sixteen.  Thus the image modulo sixteen is cyclic of order four,
generated by $9+4i$, and its image modulo eight has order two.
It follows that
$[k(E_4[16]):k]=4$ and $[k(E_4[8]):k]=2$.
The two inclusions from Lemma~\ref{lem:eighth-two-conductor}
become inclusions of extensions of these same degrees, since
the radical Kummer image over $k$ is $2\widehat{\mathbb Z}$.
They are therefore equalities.
\end{proof}

\subsection{The cycle equations fix the harmonic frame modulo sixteen}

Use the signed matrices
\[
 U=\begin{pmatrix}1&1\\0&1\end{pmatrix},\quad
 L=\begin{pmatrix}1&0\\-1&1\end{pmatrix},\quad S=ULU,
 \qquad q(x)=U^2,\quad q(y)=L^2.
\]
Write $\gamma_3F=\overline{[F,[F,F]]}$ for the third lower
central subgroup of $F=\widehat F(x,y)$.

\begin{lemma}\label{lem:eighth-two-class-two-image}
Modulo sixteen the commutator subgroup of $q(F)$ is the central
cyclic group of order four generated by
\[
 C=[U^2,L^2]
   =\begin{pmatrix}13&8\\8&5\end{pmatrix}\pmod{16}.
\]
In particular $q_{16}(\gamma_3F)=1$.
\end{lemma}

\begin{proof}
Direct multiplication gives $C^2=9I$, $C^4=I$, and shows that
$C$ commutes with both $U^2$ and $L^2$ modulo sixteen.
In a group generated by two elements, the commutator subgroup
is the normal closure of their commutator.  Here that normal
closure is precisely the central group $\langle C\rangle$.
\end{proof}

\begin{lemma}\label{lem:eighth-two-normalized-gentle}
If $h=(1,f)$ is gentle and $\rho_2(h)=0$, then
\[
 \beta_h\in\gamma_3F,\qquad
 q_{16}(\beta_h)=D_h\bmod16=W_h\bmod16=I.
\]
\end{lemma}

\begin{proof}
In $F/\gamma_3F$, write $f=[x,y]^a$.
At multiplier one, the three-cycle equation makes $3a=0$,
because its three commutators have the same central image.
The commutator coordinate is in the torsion-free group
$\widehat{\mathbb Z}$, so $a=0$ and $f\in\gamma_3F$.
The equation $\rho_2(h)=0$ gives $\beta_h\in F'$.
Write its image modulo $\gamma_3F$ as $[x,y]^d$.
The interchange $\theta(x)=y$, $\theta(y)=x$ inverts this
central commutator, and the frame identity
$f=\theta(\beta_h)\beta_h^{-1}$ gives $2d=0$.
Therefore $d=0$, proving $\beta_h\in\gamma_3F$.

Lemma~\ref{lem:eighth-two-class-two-image} gives
$q_{16}(\beta_h)=I$.  The corrected frame
$W_h=q(\beta_h)^{-1}D_h$ centralizes $S$ as a signed matrix.
Consequently the upper-unipotent matrix $D_h\bmod16$ centralizes
$S$.  Comparing the two diagonal entries in $D_hS=SD_h$
forces its shear to be zero modulo sixteen.  Thus both $D_h$
and $W_h$ are the identity modulo sixteen.
\end{proof}

\begin{theorem}\label{thm:eighth-two-compatibility}
Let $M_0=\ker\lambda$ in the gentle group.  Define the continuous
homomorphism
$\nu_2:N_4(4)\to\mathbb Z/4$ by its arithmetic normalization:
\[
 \nu_2(\rho_4(\sigma))
       =\frac{\kappa_2(\sigma)}2\pmod4
       \qquad(\sigma\in G_{\mathbb Q^{\mathrm{ab}}}).
\]
This is well defined and has kernel the subgroup congruent to
one modulo sixteen.  For every $h\in M_0$,
\begin{equation}
 \nu_2(W_h)=\frac{\rho_2(h)}2\pmod4.
 \label{eq:eighth-two-frame-kummer}
\end{equation}
In particular the exceptional eighth-root coordinate of $2$
in the radical--CM compatibility defect is identically zero
on full GT.
\end{theorem}

\begin{proof}
Proposition~\ref{prop:eighth-two-division-fields} identifies
the quotient of the arithmetic harmonic image modulo sixteen
with $\operatorname{Gal}(k(2^{1/8})/k)$.
The Kummer action on the eighth root has even exponent on $G_k$,
so division by two identifies this group with $\mathbb Z/4$.
This proves the stated definition and kernel of $\nu_2$.

The normalization of the gentle group gives
$M_0=N_0G_k$, where $N_0=\ker(\lambda,\rho_2)$.
Explicitly, $\rho_2(h)$ is even, and the Kummer map
$G_k\to2\widehat{\mathbb Z}$ is onto; choose $\sigma\in G_k$
with $\kappa_2(\sigma)=\rho_2(h)$ and put $n=h\sigma^{-1}$.
Then $n\in N_0$.
By Lemma~\ref{lem:eighth-two-normalized-gentle},
$W_n\equiv I\pmod{16}$.
The corrected-frame homomorphism and its exact arithmetic
comparison give
\[
 \nu_2(W_h)=\nu_2(W_\sigma)
           =\nu_2(\rho_4(\sigma))
           =\kappa_2(\sigma)/2
           =\rho_2(h)/2\pmod4.
\]
This proves \eqref{eq:eighth-two-frame-kummer}.
It is the required compatibility on the multiplier-one subgroup.
Matching an arbitrary multiplier by a Galois element extends the
vanishing to the corresponding crossed-cocycle component on $G$.
\end{proof}

\subsection{The cubic coordinate of two}

Recall the modular presentation
$\Gamma=\langle s,r\mid s^2=r^3=1\rangle$ and its pure subgroup
$F=\ker(\Gamma\to S_3)$, with
$x=u^2$, $y=v^2$, $s=uvu$, $r=uv$, and $sr=u^{-1}$.
For $h\in M_0$ the positive peripheral automorphism satisfies
\[
 \Phi_h(s)=\beta_hs\beta_h^{-1},\qquad
 \Phi_h(r)=A_hrA_h^{-1},\qquad
 \Phi_h(sr)=sr,
 \qquad A_h,\beta_h\in F .
\]
Its corrected cubic frame is
\[
 T_h=q(A_h)^{-1}D_h,\qquad
 R=-UL=\begin{pmatrix}0&-1\\1&-1\end{pmatrix}.
\]
Thus $T_h$ centralizes $R$ as a signed matrix.

\begin{lemma}\label{lem:cubic-two-pure-frame}
For every $h\in M_0$, one has $A_h\in F'$.
\end{lemma}

\begin{proof}
Write the abelian group $F^{\mathrm{ab}}$ additively.
Expanding $\Phi_h(s)\Phi_h(r)=sr$ in the extension
$\Gamma/F'$ gives
\[
 (1-s)[\beta_h]+s(1-r)[A_h]=0.
\]
The harmonic normalization is
$[\beta_h]=-\rho_2(h)([x]+[y])$.
Conjugation by $s$ interchanges $[x]$ and $[y]$, so the first
term vanishes and $(1-r)[A_h]=0$.
Conjugation by $r$ sends
$x\mapsto y$ and $y\mapsto(xy)^{-1}$; its matrix on
$F^{\mathrm{ab}}=\widehat{\mathbb Z}^2$ is therefore
\[
 \begin{pmatrix}0&-1\\1&-1\end{pmatrix}.
\]
The determinant of $1-r$ is $3$.  Since multiplication by $3$
is injective on $\widehat{\mathbb Z}^2$, the adjugate identity
shows that $1-r$ is injective.  Hence $[A_h]=0$.
\end{proof}

The centralizer of $R$ in $\mathrm{SL}_2(\mathbb F_3)$ is
$\{\pm R^j:j\in\mathbb Z/3\}$.  Define
\[
 j_3:N_3(4)\longrightarrow\mathbb Z/3,\qquad
 j_3(t)=j\quad\hbox{if }t\bmod3=(-I)^eR^j .
\]
This is the quotient that kills the central sign.

\begin{proposition}\label{prop:cubic-two-compatibility}
For every $h\in M_0$,
\[
 j_3(T_h)=\rho_2(h)\pmod3.
\]
On the arithmetic subgroup, this quotient is the Kummer action
on $k(2^{1/3})/k$, with the indicated choice of exponent.
\end{proposition}

\begin{proof}
The abelianization of $\mathrm{SL}_2(\mathbb F_3)$ is cyclic of
order three.  Its quotient map $\ell$ can be normalized by
$\ell(U)=\ell(L)=1$; then
$\ell(S)=0$ and $\ell(R)=2$.
For completeness, its derived subgroup is the quaternion group
generated by $S$ and $USU^{-1}$.
These generators have square $-I$, anticommute, and generate
a normal subgroup of order eight.  Conjugation by $U$ permutes
its three cyclic subgroups of order four.  Commutators with $U$
generate this quaternion group, proving the abelianization
assertion.

Write $D_h=\left(\begin{smallmatrix}1&b_h\\0&1\end{smallmatrix}\right)$.
Since $q(x)=U^2$ and $q(y)=L^2$, the harmonic normalization gives
\[
 \ell(q(\beta_h))=-4\rho_2(h)=-\rho_2(h)\pmod3.
\]
The centralizer of $S$ in $\mathrm{SL}_2(\mathbb F_3)$ is
$\langle S\rangle$, so $\ell(W_h)=0$.  The identity
$W_h=q(\beta_h)^{-1}D_h$ now gives
$b_h=-\rho_2(h)$ modulo three.
By Lemma~\ref{lem:cubic-two-pure-frame},
$\ell(q(A_h))=0$, and consequently
\[
 2j_3(T_h)=\ell(T_h)=\ell(D_h)
       =-\rho_2(h)\pmod3 .
\]
Multiplying by $2$ yields the stated formula.

The arithmetic identification can also be seen directly from
$E_3:Y^2=X^3+1$.  The duplication formula
$X(2P)=9X(P)^4/(4Y(P)^2)-2X(P)$ shows that
its three-division polynomial is
$3X(X^3+4)$.  The nonzero three-torsion coordinates satisfy
either $X=0$, $Y=\pm1$, or
$X^3=-4$, $Y^2=-3$.
Since $k$ contains the roots of unity and $\sqrt{-3}$, this gives
\[
 k(E_3[3])=k(2^{1/3}).
\]
The signed arithmetic comparison is $T_\sigma=\kappa(\sigma)
\rho_3(\sigma)$, where $\kappa(\sigma)\in\{\pm1\}$.
Its reduction modulo three can therefore have an additional
central sign.  Passing to $j_3$ kills that sign, and gives the
cubic field just identified; no equality with the full corrected
modulo-three field is asserted.
On $\sigma\in G_k$, the formula already proved reads
$j_3(T_\sigma)=\kappa_2(\sigma)\pmod3$.
This fixes its normalization as the usual Kummer exponent.
\end{proof}

\begin{corollary}\label{cor:complete-two-compatibility}
There is a continuous homomorphism
\[
 \nu_{2,\mathrm{CM}}:N_3(4)\times N_4(4)
          \longrightarrow\mathbb Z/12
\]
specified by
\[
 \nu_{2,\mathrm{CM}}(t,w)\equiv\nu_2(w)\pmod4,\qquad
 \nu_{2,\mathrm{CM}}(t,w)\equiv2j_3(t)\pmod3 .
\]
For every $h\in M_0$,
\[
 \nu_{2,\mathrm{CM}}(T_h,W_h)=\frac{\rho_2(h)}2\pmod{12}.
\]
Arithmetically this is the common quotient
$\operatorname{Gal}(k(2^{1/24})/k)$, with Kummer exponent
$\kappa_2/2$.
In particular the entire order-twelve coordinate belonging to
$2$ in the radical--CM compatibility defect vanishes on full GT.
\end{corollary}

\begin{proof}
The Chinese remainder theorem defines the homomorphism.
The two preceding results give its asserted value because
$2$ is its own inverse modulo three.
The fields $k(2^{1/8})$ and $k(2^{1/3})$ have degrees four and
three over $k$, and together generate $k(2^{1/24})$.
Their arithmetic Kummer exponents agree with the two displayed
congruences.  Finally, the compatibility defect is a crossed
cocycle that vanishes on the arithmetic subgroup.  Matching
the multiplier of any full-GT element by an arithmetic element
reduces its value to the multiplier-one case proved here.
\end{proof}

\section{Finite tests for the radical--CM compatibility defect}
\label{sec:lr-relative-gate}

Write $G=\widehat{\mathrm{GT}}$, $H=G_{\mathbb Q}$,
$k=\mathbb Q^{\mathrm{ab}}$, and $D=N_{\mathrm{rad}}$.
We use the established individual arithmetic quotients, the exact
radical--CM intersection, the signed CM-frame comparison, the normalized
signed-congruence identity, and the full hexagonal comparison of
\cite[Proposition~18]{H17} as inputs. Thus $G=DH$, and on $D$ the
multiplier and every rational Kummer cocycle vanish. The crossed defect is
\[
 d:G\longrightarrow C=\operatorname{Gal}(I/k),\qquad
 d(gh)=d(g)+\lambda_gd(h),\quad d|_H=0,
\]
with $C\simeq(\mathbb Z/12)^2\times\prod_{q\ne2,3}\mathbb Z/6$.
Nothing below assumes $d=0$, joint factorization, or $B/2B=0$.
Complete prime-two compatibility is a separate input when that
coordinate is removed from the final list of gates.

\subsection{An obstruction to a group-theoretic shortcut}

\begin{proposition}\label{prop:lr-abstract-countermodel}
Separate arithmetic quotient identifications do not imply vanishing
of a cyclotomically crossed defect. Even a trivial ordinary relative
abelian quotient does not force its order-three component to vanish.
\end{proposition}
\begin{proof}
Let $U=C_2$ act by negation on $V=\mathbb F_3$, and put
\[
 G_*=(V\oplus V)\rtimes U,\qquad
 H_*=\{(x,x,u):x\in V,\ u\in U\},\qquad A_*=V\rtimes U.
\]
The two projections $a_*,b_*:G_*\to A_*$ restrict to isomorphisms
on $H_*$. Hence $G_*=(\ker a_*)H_*=(\ker b_*)H_*$. Nevertheless
$d_*(x,y,u)=x-y$ is a nonzero crossed homomorphism, trivial on $H_*$.
Since $[G_*,G_*]=V\oplus V$, the abelianization of $G_*$ is $U$,
already covered by $H_*$. Its ordinary relative abelian quotient
is therefore zero. This is a countermodel to the proposed abstract
inference, not a construction of an element of GT.
\end{proof}

There is an exact cohomological reduction. Put
$G^+=\ker(\lambda\bmod3)$ and $H^+=H\cap G^+$, and choose complex
conjugation $c\in H$. Then $G=G^+\rtimes\langle c\rangle$.
Crossed homomorphisms to $\mathbb F_3(1)$ trivial on $H$ correspond
to homomorphisms
\[
 \chi:G^+\longrightarrow\mathbb F_3,\qquad
 \chi|_{H^+}=0,\qquad \chi(cgc^{-1})=-\chi(g).
\]
Restriction gives one direction; the inverse is $d(gc^e)=\chi(g)$.
Thus the relevant group is the minus eigenspace in the relative
abelianization of an index-two subgroup. Ordinary transfer applies
$1+c$ and kills this eigenspace; it does not show that it is zero.
Complex conjugation is not central in $G$, so the argument for a
central element acting nontrivially on coefficients cannot be used.
Odd conjugation weights all agree with the cyclotomic weight modulo
three and supply no weight mismatch.

By contrast, the two-primary image is the subgroup $d_2(D)$, of
exponent at most four, because $G=DH$. The action on
$d_2(D)/2d_2(D)$ is trivial, so $d_2$ induces ordinary relative
quadratic characters. A nonzero group of exponent four has a nonzero
quotient by twice itself. This observation does not prove that the
resulting characters vanish.

\subsection{The integral cubic-frame comparison}

Put $N=\ker(\lambda,\rho_2)$ and $\nu_g=B_g(\omega,\omega)$,
where $\omega^2+\omega+1=0$. If
$[f_g]=A_{f_g}(X,Y)[x,y]$ in $F'/F''$, then, for multiplier one,
\[
 B_g(X,Y)=1+(X-1)(Y-1)A_{f_g}(X,Y).
\]

\begin{proposition}\label{prop:lr-integral-cubic-beta}
For $g\in N$, an integral identification of the quadratic orders
identifies $T_g$ with either $(-1)^{a_2(g)}\nu_g$ or
$(-1)^{a_2(g)}\overline{\nu_g}$. For $g\in D$ the sign is one.
No uniform orientation is needed for cube membership or for the
scalar modulo-eight bit.
\end{proposition}
\begin{proof}
We specify the groups in the full comparison of
\cite[Proposition~18]{H17}. Put
\[
 \Gamma=B_3/\langle(uv)^3\rangle\simeq C_2*C_3,\qquad
 \Gamma_{1,1}=B_3/\langle(uv)^6\rangle,\qquad
 \Gamma_{1,2}=B_4/\langle(abc)^4\rangle,
\]
where $a,b,c$ are the standard chain generators of $B_4$.
The forgetting map
$\operatorname{forg}:\widehat\Gamma_{1,2}\to
\widehat\Gamma_{1,1}$ sends $(a,b,c)$ to $(u,v,u)$.
Its kernel is the completed punctured-torus group $\widehat\Pi$,
free on
\[
 \xi_1=ca^{-1},\qquad
 \xi_2=abac^{-1}b^{-1}a^{-1}.
\]
On $\widehat\Pi^{\mathrm{ab}}=\widehat{\mathbb Z}^2$,
$a,c$ act by $U$ and $b$ by $L$.

For $g=(1,f_g)$ let $\mathcal A_g$ be the braid action on
$\widehat\Gamma_{1,2}$:
\[
 \mathcal A_g(a)=a,\quad
 \mathcal A_g(b)=f_g(a^2,b^2)^{-1}b f_g(a^2,b^2),\quad
 \mathcal A_g(c)=P_g^{-1}cP_g,\quad
 P_g=f_g(ba^2b,c^2).
\]
Let $\widetilde A_g=A_g(a^2,b^2)$ be the evaluation of the
pure cubic conjugator $A_g\in\widehat F(x,y)$, and put
$\Psi_g=\operatorname{Ad}(\widetilde A_g^{-1})\mathcal A_g$.
Writing $J_3=\operatorname{forg}^{-1}(\langle uv\rangle)$,
the cited proposition gives an isomorphism
$\varphi:\widehat\Gamma\to J_3$, induced by a discrete
isomorphism and carrying $\overline{\Gamma'}$ to
$\widehat\Pi$, and an $h\in J_3$ such that
\[
 \Psi_g\varphi=\operatorname{Ad}(h)\varphi\Phi_g .
\]
Here $\overline{\Gamma'}$ is the closure of the commutator
subgroup of $\Gamma$, whose abelianization is the hexagonal
lattice $\widehat{\mathbb Z}[\omega]$.
We take this identity on the whole integral homology.

Let $M_g$ be the action of $\mathcal A_g$ on
$\widehat\Pi^{\mathrm{ab}}$. Equivariance of the forgetting map
and the established normalized signed-congruence identity imply
that $M_g$ commutes with $\mathrm{SL}_2(\widehat{\mathbb Z})$.
The centralizer of its two elementary unipotents consists of
scalars. Moreover $P_g\in\widehat\Pi$: after forgetting,
its two arguments are $(uv)^3u^{-2}$ and $u^2$, which commute,
and $f_g$ is a commutator word. Since
$ba^2b=(ab)^3a^{-2}$ and $c$ commutes with $a$, the element
$P_g$ commutes with $a$. Thus
$\mathcal A_g(\xi_1)=P_g^{-1}\xi_1P_g$, so $M_g$ fixes this
primitive homology class and $M_g=I$.

The corrected action on homology is consequently
$q(A_g)^{-1}=T_g$, since $D_g=I$ on $N$. The other side is
\[
 R_+^j C\nu_gC^{-1},\qquad j\in\mathbb Z/6,
\]
where $R_+=UL$ and the discrete homology identification $C$ sends the order-six
lattice rotation to $R_+$ or $R_+^{-1}$.
Lemma~15 of \cite{H17} gives
$\nu_g\equiv(-1)^{a_2(g)}\pmod4$.
Since $T_g\equiv I\pmod4$ and $R_+$ has exact order six modulo
four, $j=3a_2(g)\pmod6$. The two embeddings of the quadratic
order differ by its involution, proving the assertion.
On $D$, the established sign identity is
$a_2(g)=\rho_3^I(g)/2\bmod2=0$.

The identification of the lattice scalar with $B_g(\omega,\omega)$
is integral: $[x,y]$ maps to the translation $-3$ in the affine
triangle quotient, so $-3\nu_g=-3B_g(\omega,\omega)$.
Multiplication by three is injective in
$\widehat{\mathbb Z}[\omega]$, including at three, and can be cancelled.
\end{proof}

\subsection{The order-three gate}

\begin{lemma}\label{lem:lr-cubic-cm-quotient}
The maximal exponent-three subextension of $K_3^{\mathrm{ab}}/k$
is $k(q^{1/3}:q\text{ rational prime})$. The order-three part
of $d$ is therefore zero if and only if
$T_g\in N_3(4)^3$ for every $g\in D$.
\end{lemma}
\begin{proof}
Let $A=\operatorname{Gal}(K_3^{\mathrm{ab}}/K_3)$ and let $c$
act by complex conjugation. The full group over $\mathbb Q$ is
$A\rtimes\langle c\rangle$, whose closed derived subgroup is
$(c-1)A=\operatorname{Gal}(K_3^{\mathrm{ab}}/k)$.
On the pro-three part the idempotents $(1\pm c)/2$ split $A$,
and $(c-1)A$ is precisely its minus part.
Every cubic character of this subgroup thus extends to $A$,
with $c$ acting by negation on the cyclic cubic quotient.

Kummer theory writes the corresponding field as
$K_3(\alpha^{1/3})$, where the class of $\alpha$ in
$K_3^\times/(K_3^\times)^3$ is fixed by conjugation: inversion
on the Galois quotient is already accounted for by the action
on $\mu_3$. A fixed class comes from $\mathbb Q$ because
$[N_{K_3/\mathbb Q}(\alpha)]=[\alpha]^2$, and squaring again
recovers $[\alpha]$ in this exponent-three group.
Conversely rational cube roots give such extensions.
Character separation proves the simultaneous field statement.
The normalized cubic torus is an isomorphic model of the Galois
group over $k$; its quadratic sign correction has no effect on
the quotient by cubes. On $D$ the radical restriction is trivial.
\end{proof}

\begin{lemma}\label{lem:lr-wild-three-cube}
For every $g\in D$, the three-adic component of $\nu_g$ is a cube
in the norm-one units of $\mathbb Z_3[\omega]$.
\end{lemma}
\begin{proof}
Put $\pi=\omega-1$. The class-two three-cycle relation gives
$3A_{f_g}(1,1)=0$, hence $A_{f_g}(1,1)=0$.
Symmetry gives a three-adic expansion
\[
 A_{f_g}(1+S,1+T)=a(S+T)+O((S,T)^2).
\]
In the degree-three cyclic-cover homology modulo three one has
$Y=1$, $X^3=1$, and $[x,y]\mapsto(X-1)y_0$.
Since $(X-1)^3=0$ in $\mathbb F_3[C_3]$, the image of $f_g$ is
$a(X-1)^2y_0=a(1+X+X^2)y_0$.
Its $y_0$ coefficient is $a$, proving
$\rho_3^I(g)\equiv-a\pmod3$. This is zero on $D$. Thus
\[
 \nu_g=1+\pi^2A_{f_g}(\omega,\omega)
      =1+2a\pi^3+O(\pi^4)\equiv1\pmod{\pi^4}.
\]
The unit has norm one by
Proposition~\ref{prop:lr-integral-cubic-beta}, and it is one
modulo three by \cite[Lemma~15]{H17}.
The logarithm identifies the norm-one units in $1+\pi^3O_3$
with $3\sqrt{-3}\,\mathbb Z_3$. Congruence modulo $\pi^4$
forces the logarithm into $9\sqrt{-3}\,\mathbb Z_3$, whose
valuation is five. This is three times the previous lattice;
exponentiation proves cube membership. The possible $\mu_3$
factor was already excluded by congruence to one modulo three.
\end{proof}

\begin{theorem}\label{thm:lr-tame-cubic-gate}
The order-three component of $d$ vanishes if and only if, for
all $g\in D$ and primes $\ell\ge5$,
\begin{equation}
 B_{g,\ell}(\omega,\omega)^{(\ell-\chi_{-3}(\ell))/3}=1
 \quad\text{in the residue norm-one torus}.
 \label{eq:lr-tame-cubic-test}
\end{equation}
Use a primitive cube root in $\mathbb F_\ell$ in the split
case and in $\mathbb F_{\ell^2}$ in the inert case.
\end{theorem}
\begin{proof}
Cubing is bijective on the pro-two factor $N_{3,2}(4)$.
The three-component is handled by the preceding lemma.
For $\ell\ge5$, local reduction has a pro-$\ell$ kernel and a
cyclic norm-one image of order $\ell-\chi_{-3}(\ell)$.
Cubing is bijective on the kernel, and its image on the residue
group is exactly the kernel of the displayed power.
Apply Proposition~\ref{prop:lr-integral-cubic-beta} and
Lemma~\ref{lem:lr-cubic-cm-quotient}. Inversion preserves the tests.
For the product, the cube image is the product of the local
cube images, so no simultaneous condition has been discarded.
\end{proof}

The Anderson relation $(A_\ell)$ is sufficient for
\eqref{eq:lr-tame-cubic-test}. In its beta form
\cite[Remark~2.4]{Furusho2003}, specialize at $s=1/3$ and
reduce in characteristic $\ell$. The $\ell$th roots in the
denominator all reduce to one, where $B_g(1,\omega)=1$.
The numerator contains $B_g(\omega,\omega)$ exactly
$(\ell-\chi_{-3}(\ell))/3$ times; all other factors are boundary
values equal to one. The arithmetic right side is one on $D$.
This implication does not assert $(A_\ell)$ for all of GT.

\subsection{A formal obstruction to removing the tame tests}

\begin{proposition}\label{prop:lr-formal-beta-obstruction}
Beta associativity, unitarity, all dyadic Anderson identities,
vanishing integral Kummer coordinates, and the logarithmic
derivatives of the Anderson identities do not, as formal integral
Alexander identities, imply all tests
\eqref{eq:lr-tame-cubic-test}.
\end{proposition}
\begin{proof}
Fix $\ell\ge5$. By the full affine radical image choose
$\sigma\in H^{\mathrm{cyc}}$ whose sole nonzero radical
translation is $\rho_\ell^I(\sigma)=2e_3$, with
$e_3\in\widehat{\mathbb Z}$ the coefficient-prime-three
idempotent. In particular $\rho_2^I(\sigma)=\rho_3^I(\sigma)=0$
and every $\ell$-adic Kummer coordinate is zero.
The arithmetic $(A_\ell)$ relation at $s=1/3$ gives
\[
 B_{\sigma,\ell}(\omega,\omega)^{
          (\ell-\chi_{-3}(\ell))/3}=\omega^{-2\ell}\ne1
 \quad\bmod\ell.
\]
Take the coefficient-$\ell$ component of $A_{f_\sigma}$ and
set every other coefficient-prime component to zero.
Its beta unit has component $B_{\sigma,\ell}$ at $\ell$ and
one elsewhere. This is integral in the completed group algebra.
Associativity, unitarity, and the dyadic identities hold
componentwise, because they hold for $\sigma$ and the identity.
Every cyclic-cover coefficient is a
$\widehat{\mathbb Z}$-linear functional on $F'/F''$.
Thus every Kummer coordinate of the patched data is zero.
The logarithmic derivatives here are taken in the coefficient-$p$
formal variable near each torsion specialization, for every
coefficient prime $p$. At $p=\ell$, write the variable as
$\zeta(1+S)$, with $\zeta$ of order prime to $\ell$.
The arithmetic Anderson discrepancy for $\sigma$ has principal
factor $(1+S)^{n\rho_{n,\ell}^I(\sigma)}=1$; its remaining
factor is constant torsion on that specialization.
Its logarithmic derivative is therefore zero. At every other
coefficient prime the patched unit is one. This proves all the
stated derivative identities while leaving the displayed
nontrivial constant $\omega^{-2\ell}$ intact.
The displayed tame test still fails.

This constructs formal Alexander data, not a lift satisfying the
full profinite pentagon. It also does not assert the full
nonabelian $(K_n)$ identities for a lift; these contain information
beyond the numerical Kummer and logarithmic-derivative consequences.
The example isolates a missing constraint in the proposed formal
inference without supplying a counterexample inside GT.
\end{proof}

\subsection{The quartic coordinates}

\begin{lemma}\label{lem:lr-quartic-cm-quotient}
The maximal elementary-two subextension of
$K_4^{\mathrm{ab}}/k$ is
\[k(q^{1/4}:q\text{ rational prime}).\]
\end{lemma}
\begin{proof}
Use CM reciprocity in the primary-unit realization
\[
 \operatorname{Gal}(K_4^{\mathrm{ab}}/K_4)
 \simeq\{u\in\widehat O_4^\times:
                    u_2\equiv1\pmod{(1+i)^3}\},
\]
whose commutator image is $N_4(4)$.
The local two-component is isomorphic by logarithm to
$4i\mathbb Z_2$ and has one quadratic character.
Each odd component has one quadratic character, since its
residue norm-one group is cyclic of order
$q-\chi_{-4}(q)$ and its principal kernel is pro-$q$.

The character of $2^{1/4}$ is supported only at two and is
nontrivial over $k$, so it spans the first character space.
For odd $q$, $K_4(q^{1/4})/K_4$ has conductor supported at
two and $q$, and its tame character at $q$ has order four.
On norm-one units it restricts to the unique nontrivial
quadratic character: writing such a unit as $u/\bar u$,
complex conjugation inverts the Kummer character, so its value
is the square of the order-four tame character of $u$.
The local ratio map is surjective.
Thus the restriction is the tame $q$-character plus a possible
two-component character. Conductor support and the primary-unit
form of reciprocity exclude every other local component.
These characters, together with the character of $2^{1/4}$,
span the continuous quadratic character space of the product.
Character separation proves the result. The spanning argument
uses local support, not a comparison of infinite dimensions.
\end{proof}

On $N$ one has $W_{g,2}\equiv I\pmod8$. Indeed
$\beta_g\in F'$ and $D_g=I$, and modulo eight the commutator
subgroup generated by $U^2,L^2$ is $\{I,5I\}$.
The reduction $5I$ cannot lift to a norm-one element $a+bi$:
if $a\equiv5$ and $b\equiv0\pmod8$, then
$a^2+b^2\equiv9\pmod{16}$. This observation is subsumed by
the separate stronger prime-two compatibility theorem.
Consequently the local two-component of $W_g$ is a square.
At three, the established sign identity on $D$ places
$W_{g,3}\bmod3$ in $\{I,-I\}$, the square subgroup of the
cyclic order-four residue torus.
Lemma~\ref{lem:lr-quartic-cm-quotient} now gives the exact
simultaneous reduction: all rational quartic compatibilities
hold if and only if
\begin{equation}
 W_{g,\ell}^{(\ell-\chi_{-4}(\ell))/2}=1
 \quad\text{in the residue norm-one torus}
 \qquad(g\in D,\ \ell\ge5).
 \label{eq:lr-tame-quartic-test}
\end{equation}
Proposition~\ref{prop:tame-harmonic-word-gate} expresses these
tests directly as finite word evaluations by the normalized
harmonic identity. Their vanishing remains open.

\subsection{The exceptional eighth root of three}

Put $E_3:y^2=x^3+1$, $\alpha=3^{1/4}$, and
$v=(2+\sqrt3)^{1/4}$, using positive roots.

\begin{lemma}\label{lem:lr-cubic-eight-field}
One has $k(E_3[8])=k(\alpha,v)$.
The scalar automorphism $[-1]$ of this division field fixes $v$.
On $H^{\mathrm{cyc}}$, if
\[
 T_{\sigma,2}\equiv(1+4\tau(\sigma))I\pmod8,
 \qquad \tau(\sigma)\in\mathbb F_2,
\]
then $\sigma(v)/v=(-1)^{\tau(\sigma)}$.
\end{lemma}
\begin{proof}
Write $e_j=-\omega^j$, for $j=0,1,2$.
The abscissae of the non-two-torsion points in $E_3[4]$ are
$\omega^j(-1+\sqrt3)$ and $\omega^j(-1-\sqrt3)$.
The established four-division calculation is
$k(E_3[4])=k(\alpha)$.
To halve $P=(x_P,y_P)$ one adjoins square roots of $x_P-e_j$
with product of the prescribed sign $-y_P$.
Conversely these square roots belong to the field of a half
$Q=(x_Q,y_Q)$, because duplication gives
\[
 x_P-e_j=
 \left(\frac{(x_Q-e_j)^2-(e_j-e_{j+1})(e_j-e_{j+2})}
 {2y_Q}\right)^2.
\]
Indices are modulo three. This supplies both directions of
the square-root criterion inside the full division field.

For $x_P=-1+\sqrt3$, set $A=\sqrt3-1+\omega$.
Since $v^{-2}=\sqrt{2-\sqrt3}$, direct multiplication gives
\[
 A=\zeta_{24}^{\,5}\sqrt3\,v^{-2},\qquad
 \sqrt A=\zeta_{48}^{\,5}\alpha/v
 \quad\text{up to sign}.
\]
Hence $v\in k(E_3[8])$.
Conversely every square root needed to halve the displayed
four-torsion abscissae is a root of unity times
$\alpha$, $\alpha/v$, or $\alpha v$.
For the minus choice the relevant identity is
\[
 -\sqrt3-1+\omega=\zeta_{24}^{\,11}\sqrt3\,v^2.
\]
Conjugating and multiplying by roots of unity gives all cases.
Thus the halving formulas place all of $E_3[8]$ in $k(\alpha,v)$.

The two quadratic extensions $k(\alpha)$ and $k(v)$ are
nontrivial and distinct. The positive numbers $\alpha,v,v/\alpha$
each have a nonreal conjugate, whereas any real element of
$k$ generates a totally real abelian field.
For $v/\alpha$, the automorphism of
$\mathbb Q(i,\sqrt2,\sqrt3)$ fixing $\sqrt3$ and changing
the sign of $\sqrt2$ negates $(v/\alpha)^2$.
The same criterion, or the defining polynomials, handles
$\alpha$ and $v$. Hence $[k(\alpha,v):k]=4$.

Under $[-1]$, the four-torsion ordinate changes sign,
so $\alpha$ changes sign. The recovered square root in
the duplication formula also changes sign when $Q$ is
replaced by $-Q$. Their quotient $v$ is therefore fixed.
The unnormalized cubic norm-one image modulo eight,
one modulo two, consists of the scalar units $1,3,5,7$.
The scalar $-1=7$ fixes $v$, and the other nontrivial
automorphism fixing $E_3[4]$, the scalar $5$, changes its sign.
The normalization from the Tate representation to $T$
multiplies by a scalar sign, which fixes $v$.
Its image modulo eight is $\{I,5I\}$, proving the assertion.
\end{proof}

For $g\in D$ define $w_3(g),\tau_3(g)\in\mathbb F_2$ by
\[
 W_{g,3}\bmod3=(-I)^{w_3(g)},\qquad
 T_{g,2}\bmod8=(1+4\tau_3(g))I.
\]
The first uses the known vanishing of the individual frame
parity; the second follows from norm one and congruence one
modulo four. Both are homomorphisms on $D$.

\begin{proposition}\label{prop:lr-exceptional-three-character}
The character $\chi_3=w_3+\tau_3$ on $D$ is $G$-invariant
and zero on $D\cap H$. It extends uniquely to an ordinary
continuous quadratic character on $G$ that is zero on $H$.
Its vanishing is exactly the remaining compatibility for
$3^{1/8}$, once the known quartic coordinate of three is imposed.
Moreover
\[
 \tau_3(g)=\frac{B_{g,2}(\omega,\omega)-1}{4}\pmod2,
\]
where the right side is scalar, in
$\mathbb F_2\subset\mathbb F_4$.
\end{proposition}
\begin{proof}
The harmonic frame extends to a matrix representation on the
whole gentle group, so conjugation fixes its scalar classes
$\{I,-I\}$ modulo three. The full cubic-frame conjugation
formula fixes the scalar classes $\{I,5I\}$ modulo eight.
Thus both bits are $G$-invariant.

Use the arithmetic witness in the intersection theorem:
an ordinate $y$ of a three-torsion point of
$E_4:y^2=x^3-x$ satisfies
\[
 y^4=\frac49\sqrt3(2+\sqrt3),\qquad
 z=\frac{y}{\sqrt{2/3}\,v},\qquad z^8=3.
\]
For $\sigma\in D\cap H$, the harmonic residue frame sends
$y$ to $(-1)^{w_3(\sigma)}y$, while
Lemma~\ref{lem:lr-cubic-eight-field} gives the sign
$(-1)^{\tau_3(\sigma)}$ on $v$.
Since $\sigma$ fixes every rational radical, it fixes $z$,
so the sum of the two bits is zero.

The extension is $\chi_3(dh)=\chi_3(d)$ for $d\in D,h\in H$.
Invariance and vanishing on the intersection prove
well-definedness and multiplicativity. Continuity follows
from the compact surjection $D\times H\to G$.
On $D$, the radical coordinate and the quartic coordinate of
three are zero, and the sign on the eighth root is exactly
this sum, proving its interpretation as the remaining gate.

Finally the integral cubic-frame comparison identifies $T_g$
with $B_g(\omega,\omega)$ or its conjugate.
Modulo eight a norm-one unit that is one modulo four is
the scalar $1$ or $5$. Conjugation changes neither, proving
the beta formula. Its vanishing on all of $G$ is proved in
Section~\ref{sec:normalized-cubic-identity}.
\end{proof}

\subsection{Status}

After the separate prime-two compatibility theorem, joint
compatibility is reduced to the tame cubic tests
\eqref{eq:lr-tame-cubic-test}, the tame harmonic quadratic
tests \eqref{eq:lr-tame-quartic-test}, and the single relative
character $\chi_3$. The wild three-adic cubic part is eliminated.
The two subsequent frame comparisons also give $\chi_3=0$,
so only the tame tests remain.
The exceptional eighth root of three is represented by a
concrete finite combination of cubic eight-division and
harmonic three-division frames. The individual quotients,
transfer, and the formal beta identities considered here do
not force the remaining tests to vanish. Full joint
factorization remains unproved.

\section{A finite interchange map for the residual character at three}

This section compares the exceptional character $\chi_3$ with two
finite specializations of the harmonic and equianharmonic formulas of
Nakamura and Tsunogai.  An explicit homomorphism between finite matrix
groups places both comparisons in the same quaternion central sign.
Furusho's dyadic identities will identify the harmonic comparison
on full GT in Section~\ref{sec:harmonic-dyadic-comparison}. The
cubic comparison follows in Section~\ref{sec:normalized-cubic-identity}.

Write $G=\widehat{\mathrm{GT}}$ and let
$D=\ker\lambda\cap\bigcap_q\ker\rho_q^I$.  The pure harmonic and
cubic conjugators are denoted $\beta_g,A_g\in\widehat F(x,y)$,
respectively, where $x=u^2,y=v^2$ in the modular group.
We use
\[
 U=\begin{pmatrix}1&1\\0&1\end{pmatrix},\quad
 L=\begin{pmatrix}1&0\\-1&1\end{pmatrix},\quad
 S=ULU,\quad R=-UL.
\]
On $D$ the established signed normalization is $D_g=I$, and the
established finite frame targets give
\[
 q_3(\beta_g)^{-1}=(-I)^{w(g)},\qquad
 q_8(A_g)^{-1}=5^{\tau(g)}I,\qquad
 \chi_3(g)=w(g)+\tau(g).
\]
Here $q_n$ is the signed braid representation reduced modulo $n$.
All signs in this section are in $\mathbb F_2$.  In particular there is
no extra minus sign in the definition of the harmonic frame.

\begin{proposition}[Finite interchange map]\label{prop:finite-interchange}
Put $X=U^2\bmod8$ and
$P=\langle X,R\bmod8\rangle\subset\mathrm{SL}_2(\mathbb Z/8)$.
There is a surjection
\[
 \kappa:P\longrightarrow\mathrm{SL}_2(\mathbb F_3)
\]
such that
\[
 \kappa(X)=-S\bmod3,\qquad
 \kappa(R)=U\bmod3,\qquad
 \kappa(5I)=-I.
\]
The group $P$ has order $192$ and the kernel has order eight.
For every profinite commutator word $f\in\widehat F(x,y)'$,
\[
 \kappa\bigl(f(X,R)\bigr)=f(S,U)=f(U,S)^{-1}
\]
if $f$ satisfies the two-cycle relation.
\end{proposition}
\begin{proof}
Put $Y=L^2\bmod8$, $C=5I$ and $Z=-I$.  Direct matrix
multiplication gives
\[
 [X,Y]=C,\qquad X^4=Y^4=C^2=Z^2=I,
\]
and $X^2,Y^2,C,Z$ are central in $N=\langle X,Y,Z\rangle$.
Consequently every element of $N$ has the form
$X^aY^bC^eZ^d$ with $0\le a,b<4$ and $e,d\in\{0,1\}$.
These 64 forms are distinct.  Indeed, if one equals $I$, its two
off-diagonal entries give $2a=2b=0\bmod8$, so $a=b=0$;
then $5^e7^d=1\bmod8$ gives $e=d=0$.
This also proves that the displayed class-two presentation is a
presentation of $N$.

Conjugation by $R$ has the following action:
\[
 X\longmapsto Y,\qquad
 Y\longmapsto Z(XY)^{-1},\qquad C\longmapsto C,\qquad
 Z\longmapsto Z.
\]
Since $R^3=I$ and its reduction modulo two has order three,
whereas every element of $N$ is one modulo two, one has
$P=N\rtimes\langle R\rangle$, of order $64\cdot3$.
Here $Y=RXR^{-1}$ and
$Z=(RYR^{-1})XY$, so this semidirect product is generated by
$X,R$ as asserted.

In $\mathrm{SL}_2(\mathbb F_3)$ put
\[
 i=-S=\begin{pmatrix}0&2\\1&0\end{pmatrix},\qquad
 t=U,\qquad
 j=tit^{-1}=\begin{pmatrix}1&1\\1&2\end{pmatrix}.
\]
Then $i^2=j^2=[i,j]=-I$, $t^3=I$, and
$tjt^{-1}=ij$.  Thus $i,j$ generate the quaternion group and
$t$ acts on its three noncentral cyclic subgroups by a three-cycle.
The assignment
\[
 X\mapsto i,\quad Y\mapsto j,\quad C\mapsto-I,\quad Z\mapsto-I
\]
respects the presentation of $N$.  It also respects conjugation by
$R$, because
$(-I)(ij)^{-1}=ij$.  It therefore extends with $R\mapsto t$
to a homomorphism $P\to Q_8\rtimes C_3$
$=\mathrm{SL}_2(\mathbb F_3)$.  It is surjective, and its kernel
has order $192/24=8$.

Evaluation of a profinite word commutes with a homomorphism to a
finite group.  Multiplying either argument of a commutator word
by a central element does not change its evaluation, since the
two exponent sums are zero.  Hence
$\kappa(f(X,R))=f(-S,U)=f(S,U)$, and the two-cycle relation gives
the final equality.
\end{proof}

\begin{lemma}[A common finite evaluation]\label{lem:common-sign}
For every $g\in D$ the matrix $f_g(U,S)\bmod3$ is central.
Consequently
\[
 f_g(U,S)=(-I)^{\eta(g)}\quad\text{in }\mathrm{SL}_2(\mathbb F_3)
\]
defines a continuous homomorphism $\eta:D\to\mathbb F_2$.
\end{lemma}
\begin{proof}
Write $[f_g]=A_f(X,Y)[x,y]$ in the completed Alexander module.
The two-cycle relation gives $A_f(X,Y)=A_f(Y,X)$.
We recall precisely the cyclic-cover input already established for
full GT.  In the homology of the subgroup
$\ker(\widehat F(x,y)\to C_3)$ with $x\mapsto1,y\mapsto0$,
use the basis $x^3,y_j=x^jyx^{-j}$, $0\le j<3$.
The three $y_j$ coefficients of $f_g$ are
$(2b,-b,-b)$, where $b=-\rho_3^I(g)/2$.
Sending $y_j$ to $\omega^j$ gives
\[
 (\omega-1)A_f(\omega,1)=3b.
\]
On $D$, $b=0$.  Since $\omega-1$ is a unit in
$\mathbb Z_2[\omega]$, it follows that
$A_f(\omega,1)=A_f(1,\omega)=0$, also after reduction modulo two.

Now project $\mathrm{SL}_2(\mathbb F_3)$ to
$A_4=V\rtimes C_3$, where $V=Q_8/\{\pm I\}$ is identified
with the additive group of $\mathbb F_4$.
The image of $S$ belongs to $V$ and the image of $U$ acts on
$V$ by a primitive cube root $\omega$ or $\omega^{-1}$.
Conjugation by $S$ is trivial on $V$.  Since this quotient is
metabelian, the image of $f_g(S,U)$ is
$A_f(1,\omega)[S,U]$, with either choice of the root; it is zero.
Thus $f_g(S,U)$, and hence $f_g(U,S)$ by the two-cycle relation,
lies in $\{\pm I\}$.

For $g\in D$, the free-group automorphism
$x\mapsto x$, $y\mapsto f_g^{-1}yf_g$ induces the identity under
the evaluation $x\mapsto U,y\mapsto S$, because the evaluated
conjugating word is central.  Applying this observation to the
composition law $f_{gh}=f_g\phi_g(f_h)$ proves that the evaluated
signs multiply.  Continuity follows from evaluation in the finite
group $\mathrm{SL}_2(\mathbb F_3)$.
\end{proof}

\begin{proposition}[An exact pair of finite comparison gates]
\label{prop:pair-of-gates}
For $g\in D$ define
\[
 \epsilon_h(g)=w(g)+\eta(g),\qquad
 \epsilon_c(g)=\tau(g)+\eta(g).
\]
These are continuous homomorphisms on $D$, and
\[
 \boxed{\quad \chi_3=\epsilon_h+\epsilon_c.\quad}
\]
More explicitly, the two signs are the following finite defects:
\[
 \begin{split}
 (-I)^{\epsilon_h(g)}
   &=q_3(\beta_g)^{-1}f_g(U,S)^{-1},\\
 (-I)^{\epsilon_c(g)}
   &=\kappa\!\left(q_8(A_g)^{-1}f_g(X,R)^{-1}\right).
 \end{split}
\]
In particular the two finite comparisons
\begin{equation}
 q_3(\beta_g)^{-1}=f_g(U,S),\qquad
 \kappa(q_8(A_g)^{-1})=\kappa(f_g(X,R))
 \label{eq:finite-harmonic-cubic-comparisons}
\end{equation}
would imply $\chi_3(g)=0$.  They are sufficient conditions;
vanishing of $\chi_3$ requires only equality of their two defects.
Both defects vanish on $D\cap G_{\mathbb Q}$.
\end{proposition}
\begin{proof}
The established frame homomorphisms make $w,\tau$ homomorphisms
on $D$, and Lemma~\ref{lem:common-sign} does the same for $\eta$.
The first finite formula follows directly from their definitions.
For the second, Proposition~\ref{prop:finite-interchange} gives
\[
 \kappa(q_8(A_g)^{-1})=(-I)^{\tau(g)},\qquad
 \kappa(f_g(X,R))=f_g(U,S)^{-1}=(-I)^{\eta(g)}.
\]
This also proves the displayed sum and the sufficient conditions.

For the arithmetic assertion use exactly the scope of
\cite[Theorem~2, pp.~199--200]{NT2002}.  Its $(\mathrm{GF}_0)$ and
$(\mathrm{HF}_1)$ formulas are stated for Galois elements.
When $\lambda=1$ and $\rho_2=\rho_3=0$, they reduce to
\[
 \beta_g^{-1}=f_g(u,s),\qquad A_g^{-1}=f_g(u^2,r).
\]
Here its harmonic parameter is $\beta^{-1}$.  Its cubic
parameter is $A^{-1}$, since $\operatorname{Ad}(r)$ cyclically
permutes $x,y,(xy)^{-1}$ and
$f_g=\operatorname{Ad}(r)(A_g)A_g^{-1}$.
These identifications also follow from the uniqueness of the
pure conjugators.  The equalities lift without a central error
from the modular quotient to the signed braid group: both
sides have degree zero in the braid abelianization, whereas
the central generator has degree six.
Under the signed representation $q(r)=UL=-R$; the central
minus sign disappears in a commutator word.
Reduction modulo three and eight therefore gives
\eqref{eq:finite-harmonic-cubic-comparisons} on the arithmetic
subgroup.  No use of those two formulas on all of GT is made.
\end{proof}

\subsection{Two relative quadratic characters}

Continue with the same notation:
$G=\widehat{\mathrm{GT}}$, $H=G_{\mathbb Q}$,
$D=\ker\lambda\cap\bigcap_q\ker\rho_q^I$, and
$f_g(U,S)=(-I)^{\eta(g)}$ for $g\in D$.
The two frame signs are $w,\tau$, and
$\epsilon_h=w+\eta$, $\epsilon_c=\tau+\eta$.
We use the previously established facts $G=DH$, the invariance
of $D$, and the signed frame conjugation laws.  The latter make
$w$ and $\tau$ invariant under $G$: on $D$ the finite frames are
scalar matrices, which every conjugation fixes.

\begin{lemma}\label{lem:eta-quotient-invariant}
The kernel of the evaluation
\[
 e:\widehat F(x,y)\longrightarrow Q=\mathrm{SL}_2(\mathbb F_3),
 \qquad e(x)=U,\quad e(y)=S,
\]
is invariant under every gentle automorphism.  Consequently $\eta$
is invariant under conjugation by all of $G$.
\end{lemma}
\begin{proof}
We first verify the precise finite presentation:
\begin{equation}\label{eq:finite-quotient-presentation}
Q=\langle x,y\mid x^3=y^4=(xy)^3=1,\ [x,y^2]=1\rangle.
\end{equation}
In the presented group, $z=y^2$ is central and has order at most
two.  Modulo $z$ one has
$\langle x,y\mid x^3=y^2=(xy)^3=1\rangle$, of order at most
twelve.  For an explicit bound, put
$b_j=x^jyx^{-j}$ for $j=0,1,2$.  They are involutions, their
product in cyclic order is one by $(xy)^3=1$, and hence two
of them commute and generate a group of order at most four.
This group is normalized by $x$, which cyclically permutes the
$b_j$.  Thus the whole quotient has order at most $4\cdot3$.
The group in \eqref{eq:finite-quotient-presentation} has order at most 24.  The matrices $U,S$ satisfy
all four relators.  They generate $\mathrm{SL}_2(\mathbb F_3)$:
$-S$ and its conjugate by $U$ generate the quaternion subgroup,
and $U$ has order three and permutes its three noncentral cyclic
subgroups.  This gives 24 distinct matrices, proving \eqref{eq:finite-quotient-presentation}.

Let $g=(\lambda,f)$ be a gentle element and use the convention
\[
 \phi_g(x)=x^\lambda,\qquad
 \phi_g(y)=f^{-1}y^\lambda f.
\]
The peripheral three-cycle condition says that $\phi_g(xy)$
is conjugate to $(xy)^\lambda$.  Therefore the images under
$e\phi_g$ satisfy the first three relators of \eqref{eq:finite-quotient-presentation}.
Moreover
\[
 e\phi_g(y^2)=e(f)^{-1}(S^2)^\lambda e(f)=-I,
\]
since $S^2=-I$ is central and $\lambda$ is odd.
This proves the fourth relator as well.  There is consequently
an induced endomorphism $\alpha_g$ of $Q$ with
$e\phi_g=\alpha_ge$.  Applying the same argument to $g^{-1}$
shows that $\alpha_g$ is an automorphism.
In particular every $\alpha_g$ fixes the central involution
$-I$, directly also visible in the displayed formula.

Put $s(g)=e(f_g)$.  The composition convention is
$f_{gh}=f_g\phi_g(f_h)$, so
\[
 s(gh)=s(g)\alpha_g(s(h)).
\]
For $d\in D$, Lemma~\ref{lem:common-sign} proves
$s(d)\in\{\pm I\}$.  Since $\lambda_d=1$, the formulas on
the two generators then give $\alpha_d=\mathrm{id}_Q$.
For $g\in G$ and $d\in D$, it follows that
\[
 \begin{split}
 s(gdg^{-1})
 &=s(g)\alpha_g(s(d))\alpha_{gd}(s(g^{-1}))\\
 &=s(g)s(d)\alpha_g(s(g^{-1}))=s(d).
 \end{split}
\]
The last equality uses centrality of $s(d)$ and
$s(g)\alpha_g(s(g^{-1}))=1$.
This proves the conjugation invariance of $\eta$.
\end{proof}

\begin{proposition}\label{prop:two-relative-extensions}
Each of $\epsilon_h,\epsilon_c$ has a unique extension to a
continuous homomorphism $G\to\mathbb F_2$ that is zero on $H$.
Their sum is the established residual relative character $\chi_3$.
The harmonic extension is proved to vanish in
Section~\ref{sec:harmonic-dyadic-comparison}. The cubic extension
and $\chi_3$ vanish by Section~\ref{sec:normalized-cubic-identity}.
\end{proposition}
\begin{proof}
By Lemma~\ref{lem:eta-quotient-invariant} and the scalar frame
conjugation laws, both homomorphisms on $D$ are invariant under
$G$.  Proposition~\ref{prop:pair-of-gates}, using only the Galois
scope of the published $(\mathrm{GF}_0)$ and $(\mathrm{HF}_1)$
formulas, proves that both are zero on $D\cap H$.

For either homomorphism $\epsilon$, define
\[
 \widetilde\epsilon(dh)=\epsilon(d),\qquad d\in D,\ h\in H.
\]
If $dh=d'h'$, then $(d')^{-1}d\in D\cap H$, so this is
well-defined.  Conjugation invariance and normality of $D$ give
\[
 \widetilde\epsilon(d_1h_1d_2h_2)
 =\epsilon(d_1)+\epsilon(h_1d_2h_1^{-1})
 =\epsilon(d_1)+\epsilon(d_2).
\]
Thus it is a homomorphism, zero on $H$.
The continuous surjection $D\times H\to G$ is a quotient map,
because its source is compact and its target is Hausdorff.
Its pullback is the continuous function $(d,h)\mapsto\epsilon(d)$,
which proves continuity.  The decomposition $G=DH$ proves
uniqueness.  The two extensions sum to $w+\tau$ on $D$ and
are zero on $H$, hence their sum is $\chi_3$ by the same uniqueness.
\end{proof}

\subsection{A limitation of a low-degree peripheral calculation}

The following completely explicit test records why checking the
modular order-three relation through degree four does not remove
the central ambiguity.  It is not a counterexample to that relation
in the full profinite group and is not a GT counterexample.

Let $\Pi=\langle a,b\rangle$ be the discrete free group, with
automorphisms
\[
 \theta(a)=a^{-1},\quad\theta(b)=b^{-1},\qquad
 \alpha(a)=a^{-1}b,\quad\alpha(b)=a^{-1}.
\]
These are the actions of the order-two and order-three modular
generators on the free subgroup $\Gamma'$ for the basis
$a=srsr^{-1}$, $b=sr^2sr^{-2}$.
For a candidate conjugator $\beta$, fixing $sr$ makes the
image of $r$ equal to $dr$, where
$d=\beta\theta(\beta)^{-1}$.  Its cube is one exactly when
$d\alpha(d)\alpha^2(d)=1$.

\begin{proposition}[A truncated peripheral counterexample]
\label{prop:truncated-counterexample}
In the mod-two Magnus ring
$\mathcal R=\mathbb F_2\langle A,B\rangle/(\text{degree}\ge5)$,
send $a$ to $1+A$ and $b$ to $1+B$.
The word $\beta=a^3bab^3$ has exponent sums $(4,4)$ and
satisfies
\[
 d\alpha(d)\alpha^2(d)=1\quad\text{in }\mathcal R.
\]
Nevertheless its image under $a\mapsto i,b\mapsto j$ in the
quaternion group is $-1$.  Thus this truncated norm equation
and the two exponent sums modulo four do not determine the
quaternion central sign.
\end{proposition}
\begin{proof}
The exponent sums are immediate, and
$i^3jij^3=(ij)^2=-1$.
Here are direct expansions verifying the norm equation.  Write
$d=1+D$, $\alpha(d)=1+E$, $\alpha^2(d)=1+F$ in $\mathcal R$.
Substitute $a^{-1}=1+A+A^2+A^3+A^4$ and the analogous
expression for $b^{-1}$.  One obtains
\[
\begin{split}
D={}&AAB+ABB+BAA+BBA\\
 &+AAAB+AABA+AABB+ABAA+ABBA+ABBB\\
 &+BAAA+BAAB+BABB+BBAA+BBAB+BBBA,\\
E={}&ABB+BBA+AABB+ABAB+ABBB+BAAB\\
 &+BABA+BABB+BBAB+BBBA,\\
F={}&AAB+BAA+AAAB+AABA+ABAA+ABAB\\
 &+ABBA+BAAA+BABA+BBAA.
\end{split}
\]
Every monomial in $D+E+F$ occurs twice, so the sum is zero.
Each of $D,E,F$ has degree at least three, hence all their
pairwise products vanish in $\mathcal R$.  Therefore
$(1+D)(1+E)(1+F)=1$.
The trivial word has the same exponent sums modulo four
and passes the same norm equation, but has quaternion image
$+1$.  This proves the claimed limitation.
\end{proof}

The displayed expansion proves only the limitation of this truncated
norm equation.  It verifies neither the full pentagon nor the full
nonabelian distribution relation, and does not produce a profinite
GT element.

\subsection{Two-dimensional specializations of relation $(\mathrm{III}^\prime)$}

Write
\[
 B_4=\langle a,b,c\mid aba=bab,\ bcb=cbc,\ ac=ca\rangle.
\]
In the notation of \cite[\S10, Theorem~20]{H17}, the two inputs
to the harmonic auxiliary word $g_g=\beta_g^{-1}$ are
\[
 x_{45}=(ab)^3,\qquad x_{51}=(bc)^3.
\]
On the subgroup $D$, $g_g\in\widehat F(x,y)'$.

\begin{lemma}\label{lem:iii-prime-two-dimensional}
Let $k$ be any field and let $\rho:B_4\to\mathrm{GL}_2(k)$ be
a homomorphism.  Then $\rho((ab)^3)$ and $\rho((bc)^3)$ commute.
For a finite field $k$ and $g\in D$, it follows that
\[
 g_g\bigl(\rho((ab)^3),\rho((bc)^3)\bigr)=1.
\]
Thus such a direct matrix specialization of $(\mathrm{III}')$
does not retain its harmonic auxiliary-word term.
\end{lemma}
\begin{proof}
The elementary braid identity says that $(ab)^3$ is central in
$\langle a,b\rangle$, and $(bc)^3$ is central in
$\langle b,c\rangle$.  Both displayed images therefore belong
to the centralizer of $B=\rho(b)$.

If $B$ is not scalar, its centralizer in $M_2(k)$ is the
commutative algebra $k[B]$.  To see this directly, choose a vector
$v$ such that $v,Bv$ are a basis.  Such a vector exists for every
nonscalar two-by-two matrix.  An operator commuting with $B$ is
determined by its value on $v$, and the operators $\alpha I+\beta B$
give every possible such value.  In particular the two centralizer
elements commute.

If $B=tI$, the braid relation with $A=\rho(a)$ gives
$tA^2=t^2A$, and invertibility implies $A=tI$.  The other braid
relation similarly gives $\rho(c)=tI$.  All images then commute.

Over a finite field the representation extends continuously to
the profinite completion.  Evaluation of a commutator proword on
two commuting elements is one.  Applying this to $g_g$ proves the
second assertion.
\end{proof}

This lemma only excludes this direct way of retaining the
$g_g$ term.  It does not assert that $(\mathrm{III}')$ has no
consequences for the finite defects, nor that higher-dimensional
representations or comparisons of covering groups cannot provide
the required equality.  No vanishing or nonvanishing statement for
the two defects follows from it.


\subsection{A genus-two quotient carrying the cubic comparison sign}

Use the finite-comparison notation
$U=\left(\begin{smallmatrix}1&1\\0&1\end{smallmatrix}\right)$,
$L=\left(\begin{smallmatrix}1&0\\-1&1\end{smallmatrix}\right)$,
$R=-UL$, $X=U^2$, $Y=L^2$, $C=5I$, and $Z=-I$ modulo eight.
Write $G_8=\mathrm{SL}_2(\mathbb Z/8)$ and
$P=\langle X,R\rangle$, of order $192$.
The already proved presentation has
\[
 N=\langle X,Y,C,Z\rangle,\quad |N|=64,\quad
 P=N\rtimes\langle R\rangle,
\]
with the unique forms $X^aY^bC^eZ^d$,
$0\leq a,b<4$, $0\leq e,d<2$, and
\[
 [X,Y]=C,\qquad RXR^{-1}=Y,\qquad
 RYR^{-1}=Z(XY)^{-1}.
\]
The elements $C,Z,X^2,Y^2$ are central in $N$.

In the target put
\[
 S=\begin{pmatrix}0&1\\-1&0\end{pmatrix},\quad
 T=\begin{pmatrix}1&1\\0&1\end{pmatrix},\quad
 A=I-S=\begin{pmatrix}1&2\\1&1\end{pmatrix}
 \quad\text{over }\mathbb F_3.
\]

\begin{proposition}\label{prop:genus-two-finite-quotient}
There is a surjection
\[
 \ell:G_8\longrightarrow\mathrm{GL}_2(\mathbb F_3),\qquad
 U\longmapsto A,\quad R\longmapsto T,
\]
whose kernel is the order-eight subgroup
\[
 K=\langle Z,X^2C,Y^2C\rangle.
\]
It kills $Z$, so the signed congruence representation induces
\[
 \bar\ell:\Gamma=C_2*C_3\longrightarrow
                 \mathrm{GL}_2(\mathbb F_3).
\]
For $s=uvu$, $r=uv$ and $u=r^{-1}s$, their images have orders
$2,3,8$, respectively.  On the subgroup $P$, this map agrees
with the previous interchange map on all commutator prowords in
$X,R$, although its value on $X$ is $S$ instead of $-S$.
\end{proposition}
\begin{proof}
The presentation of $P$ defines a character
$\chi:P\to\{\pm1\}$ by
\[
 \chi(X)=\chi(Y)=\chi(Z)=-1,\qquad
 \chi(C)=\chi(R)=1.
\]
For example the only nontrivial conjugation check is
$\chi(Y)=\chi(Z(XY)^{-1})=-1$; the remaining relations are
immediate.  Multiply the previously established interchange
homomorphism $\kappa:P\to\mathrm{SL}_2(\mathbb F_3)$ by this
central character.  The resulting surjection $\kappa'$ has
\[
 \kappa'(X)=S,\quad \kappa'(R)=T,\quad
 \kappa'(C)=-I,\quad\kappa'(Z)=I.
\]
Its kernel contains $K$.  The unique normal forms show that
$Z,X^2C,Y^2C$ are three independent commuting involutions, so
$|K|=8$.  Surjectivity and $|P|=192$ imply that this is the whole
kernel.  A character is trivial on commutators, giving the last
assertion about prowords.

The subgroup $P$ has index two in $G_8$, and $U^2=X\in P$.
The whole group is generated by $P$ and $U$: indeed
$L=U^{-1}ZR$.  Direct multiplication gives
\begin{equation}\label{eq:bolza-1}
UXU^{-1}=X,\qquad URU^{-1}=XYR^2.
\end{equation}
In the target, $A^2=S$, $A$ commutes with $S$, and
\begin{equation}\label{eq:bolza-2}
ATA^{-1}=S(TST^{-1})T^2.
\end{equation}
For clarity, both sides of \eqref{eq:bolza-2} are
$\left(\begin{smallmatrix}2&2\\1&0\end{smallmatrix}\right)$.
Since $X,R$ generate $P$, \eqref{eq:bolza-1}--\eqref{eq:bolza-2} show that conjugation by
$A$ on $\kappa'(P)$ agrees with conjugation by $U$ before
applying $\kappa'$.  Together with $A^2=\kappa'(U^2)$, these
identities give a well-defined homomorphism on the two cosets
$P\sqcup PU$, extending $\kappa'$ and sending $U$ to $A$.
Its image contains $\mathrm{SL}_2(\mathbb F_3)$ and the matrix
$A$ of determinant $-1$, hence is all of $\mathrm{GL}_2(\mathbb F_3)$.
Its kernel has order $384/48=8$ and is $K$.

The signed braid center maps to $Z$, so the map descends to
$\Gamma$.  There $r=uv$ maps to $T$, since $ZR$ maps to $T$;
and $s=ru$ maps to
\[
 TA=\begin{pmatrix}2&0\\1&1\end{pmatrix},
\]
an involution.  Finally $A^2=S$, $S^2=-I$ and $A^4=-I$,
so $A$ has order eight.  This proves the stated orders.
\end{proof}

\begin{corollary}\label{cor:genus-two-sign-cover}
The discrete kernel of $\bar\ell$ is the fundamental group of
a genus-two surface with six punctures.  Its compactification
carries a faithful deck action of $\mathrm{GL}_2(\mathbb F_3)$
of signature $(2,3,8)$.  The central involution is its
hyperelliptic involution.

On $D$ the induced action on this finite deck group is the
identity.  The pure harmonic conjugator maps to $I$, and the
pure cubic conjugator maps to $(-I)^{\tau(g)}$.
\end{corollary}
\begin{proof}
The subgroup has index $48$ in the modular group $\Gamma$.
It is torsion-free because the orders two and three of the
finite factors are preserved.  The cusp stabilizer has image
of order eight, so there are $48/8=6$ cusps.  Since
$\chi(\Gamma)=-1/6$, its Euler characteristic is $-8$.
Thus $2-2g-6=-8$ and $g=2$.  Filling the cusps gives the
signature $(2,3,8)$ action, by the same stabilizer counts.

The central matrix $-I=A^4$ lies in every conjugate of the
cusp stabilizer, so it fixes all six filled cusp points.
For an involution of a genus-two compact surface the
Riemann--Hurwitz formula is $2=4h-4+r$.  Since $r\geq6$,
one has $h=0$ and $r=6$.  This is the hyperelliptic involution.

On $D$ the established normalized signed covariance gives
$q_8\Phi_g=q_8$, hence identity on the deck group after $\ell$.
The established prime-two normalization gives
$q_8(\beta_g)=I$.  Also
$q_8(A_g)^{-1}=5^{\tau(g)}I=C^{\tau(g)}$, and
$\ell(C)=-I$.  Inversion leaves this sign unchanged, yielding
the asserted cubic conjugator value.
\end{proof}

This gives a concrete genus-two setting for the cubic finite
comparison. The equality $\tau=\eta$ requires the cubic branch-filling
comparison proved in Section~\ref{sec:normalized-cubic-identity}.

\paragraph{Identification with the classical Bolza cover.}
The genus-two cover in Proposition~\ref{prop:genus-two-finite-quotient}
is the classical Bolza cover.  Katz--Katz--Schein--Vishne
\cite[\S3 and Corollaries 13.3--13.4]{KKSV2016} describe its
index-two triangle-group inclusion and the corresponding
$\mathrm{SL}_2(\mathbb F_3)$ and $\mathrm{GL}_2(\mathbb F_3)$
quotients.  Their displayed matrices immediately before
Corollary 13.4 send the order-two and order-three generators to
\[
 s\longmapsto\begin{pmatrix}0&1\\1&0\end{pmatrix},\qquad
 r\longmapsto T=\begin{pmatrix}1&1\\0&1\end{pmatrix}.
\]
These are exactly conjugate to our pair:
\[
 T(TA)T^{-1}=\begin{pmatrix}0&1\\1&0\end{pmatrix},\qquad
 TTT^{-1}=T.
\]
Thus the kernels in the triangle group agree.  The compactification is the Bolza
curve $y^2=x^5-x$; the same model, its six hyperelliptic branch
points $0,\infty,\pm1,\pm i$, and its classical automorphism
group are recalled in \cite[\S4, equation (4.1)]{KRRBolza}.
That paper attributes the classification to Bolza's 1887 work.

Here we compute the factorization through
$\mathrm{SL}_2(\mathbb Z/8)$ with kernel
$\langle Z,X^2C,Y^2C\rangle$, together with the values of the
normalized frame conjugators. Section~\ref{sec:normalized-cubic-identity}
then identifies their cubic sign with $\eta$.


\section{The harmonic comparison from Furusho's dyadic identities}
\label{sec:harmonic-dyadic-comparison}

Let $g=(1,f)$ be an element of full profinite $\mathrm{GT}$ with
$\rho_2^I(g)=0$.  We use the established equality between this
Ihara coordinate and the harmonic coordinate $\rho_2$, so the
pure harmonic conjugator $\beta\in\widehat F(a,b)$ has
abelianization zero.  Write $\theta(a)=b$, $\theta(b)=a$.
The fixed convention for this conjugator is
\begin{equation}\label{eq:harmonic-word-H}
f(a,b)=\theta(\beta)\beta^{-1}.
\end{equation}
Relation $(\mathrm{IV})$ holds on full GT by
\cite[Theorem~1]{H17}, and $(K_2)$ holds there by
\cite[Theorem~8.1]{Enriquez2007}. We apply these geometric
identities to the harmonic frame.

\subsection{The precise dyadic word identities}

The kernel of $\widehat F(x,y)\to\mathbb Z/2$ sending
$x\mapsto1$ and $y\mapsto0$ is free on $xyx^{-1},x^2,y$.  Since $f$ has zero
exponent sums, there is a unique word $f^+(W,X,Y)$ such that
\begin{equation}\label{eq:harmonic-word-1}
f(x,y)=f^+(xyx^{-1},x^2,y).
\end{equation}
Furusho's general definition includes the factor
$y^{-2\Psi_2^{(0)}(g)}$ on the left.  Here that factor is one,
because $\Psi_2^{(0)}=-\rho_2^I=0$.
The exact $(K_2)$ relation, in the form
\cite[\S3, equation~(2), p.~102]{Furusho2003}, gives
\begin{equation}\label{eq:harmonic-word-2}
f^+(1,a,b)=f(a,b).
\end{equation}
We use the following two identities in their full nonabelian
domains:
\begin{equation}\label{eq:harmonic-word-3}
f^+(W,X,Y)\,
 f^+(X^{-1}W^{-1}Y^{-1},Y,X)=1,
\end{equation}
and, whenever $WXYZ=1$,
\begin{equation}\label{eq:harmonic-word-4}
f^+(W,X,Y)f^+(Z,W,X)
 f^+(Y,Z,W)f^+(X,Y,Z)=1.
\end{equation}
These are respectively Furusho~\cite[Lemmas 4.1 and 4.2,
pp.~105--106]{Furusho2003}, specialized to
$\lambda=1$ and $\Psi_2^{(0)}=0$.
Their stated hypotheses are $(\mathrm{IV})$, and
$(\mathrm{I}), (\mathrm{II}), (\mathrm{IV})$, respectively.
They are not restricted to Galois elements.  Equation \eqref{eq:harmonic-word-3} is
an identity in the free profinite group on $W,X,Y$; equation
\eqref{eq:harmonic-word-4} is an identity in the free group on three of the four
variables with their indicated product relation.  We are not
using just their Fox derivatives or abelianizations.

\begin{lemma}\label{lem:dyadic-branch-word}
Put $h(a,b)=f^+(a,1,b)$.  Then
\[
 f(a,b)=h(a,b)h(b,a)^{-1}
 \quad\text{in }\widehat F(a,b).
\]
\end{lemma}
\begin{proof}
Work with $abc=1$.  Substitute $(W,X,Y,Z)=(1,b,c,a)$ in
\eqref{eq:harmonic-word-4}, obtaining
\begin{equation}\label{eq:harmonic-word-5}
f^+(1,b,c)f^+(a,1,b)f^+(c,a,1)f^+(b,c,a)=1.
\end{equation}
By \eqref{eq:harmonic-word-2}, the first factor is $f(b,c)$, and the second is
$h(a,b)$.  Applying \eqref{eq:harmonic-word-3} to $(W,X,Y)=(b,1,a)$ gives
\[
 f^+(b,1,a)f^+(b^{-1}a^{-1},a,1)=1.
\]
Since $b^{-1}a^{-1}=c$, the third factor of \eqref{eq:harmonic-word-5} is therefore
$h(b,a)^{-1}$.
Applying \eqref{eq:harmonic-word-3} to $(W,X,Y)=(b,c,a)$ gives
\[
 f^+(b,c,a)f^+(c^{-1}b^{-1}a^{-1},a,c)=1.
\]
Here $c^{-1}b^{-1}a^{-1}=1$, so \eqref{eq:harmonic-word-2} and the two-cycle
identity make the fourth factor of \eqref{eq:harmonic-word-5} equal to
$f(a,c)^{-1}=f(c,a)$.
Consequently
\[
 f(b,c)h(a,b)h(b,a)^{-1}f(c,a)=1,
\]
whence
\[
 h(a,b)h(b,a)^{-1}=f(b,c)^{-1}f(c,a)^{-1}.
\]
The normalized three-cycle identity is
$f(a,b)f(c,a)f(b,c)=1$ for $abc=1$.
It identifies the right side with $f(a,b)$, in precisely
the displayed order.  This proves the lemma.
\end{proof}

\begin{lemma}\label{lem:unique-swap-frame}
The map $\gamma\mapsto\theta(\gamma)\gamma^{-1}$ on
$\widehat F(a,b)$ is injective.  In particular
\[
 \beta(a,b)=h(b,a).
\]
\end{lemma}
\begin{proof}
If $\theta(\gamma_1)\gamma_1^{-1}=
\theta(\gamma_2)\gamma_2^{-1}$, write
$\gamma_2=\gamma_1 k$.  Cancellation gives $\theta(k)=k$.
The semidirect product
$\widehat F(a,b)\rtimes\langle s\rangle$, where $s^2=1$
and $sas=b$, is the free profinite product
$\widehat{\mathbb Z}\amalg C_2$, generated by $a,s$.
The centralizer of $s$ is the factor $\langle s\rangle$,
by the centralizer theorem for a nonidentity factor element
\cite[Theorem~B, p.~155]{HerfortRibes1985}.
Its intersection with $\widehat F(a,b)$ is trivial, so $k=1$.

By Lemma~\ref{lem:dyadic-branch-word}, the word
$\gamma(a,b)=h(b,a)$ satisfies
$\theta(\gamma)\gamma^{-1}=f(a,b)$.
Equation \eqref{eq:harmonic-word-H} and injectivity now give $\gamma=\beta$.
\end{proof}

\begin{proposition}\label{prop:normalized-GF1-from-dyadic}
For every $g$ as above, in the full signed profinite braid group
on generators $u,v$, with $s=uvu$, one has
\[
 \beta_g(u^2,v^2)^{-1}=f_g(u^2,s).
 \tag{GF1$_0$}
\]
\end{proposition}
\begin{proof}
Equation \eqref{eq:harmonic-word-2} shows that the $X$- and $Y$-exponent sums
of $f^+$ vanish, because $f$ is a commutator word.
Equation \eqref{eq:harmonic-word-1} then shows that its $W$-exponent sum also
vanishes. In particular its $X$-exponent is zero.

In the signed braid group, put $c=(uv)^3$.
The braid relation gives $su^2s^{-1}=v^2$ and $s^2=c$,
and $c$ is central. Specializing \eqref{eq:harmonic-word-1} at $(x,y)=(s,u^2)$
therefore gives
\[
 f(s,u^2)=f^+(v^2,c,u^2)=f^+(v^2,1,u^2)
         =h(v^2,u^2)=\beta(u^2,v^2).
\]
The second equality holds because a central argument can
be collected into its total exponent, which here is zero.
This assertion for profinite words follows by continuity
from the corresponding identity for finite words.
The two-cycle identity now gives
$f(u^2,s)=f(s,u^2)^{-1}=\beta(u^2,v^2)^{-1}$,
as required.
\end{proof}

\begin{corollary}\label{cor:harmonic-defect-zero}
On $D=\ker\lambda\cap\bigcap_q\ker\rho_q^I$, the finite
harmonic frame sign satisfies $w=\eta$.  Thus the previously
constructed relative character $\epsilon_h=w+\eta$ is zero
on the whole of full GT, and the remaining exceptional
character is $\chi_3=\epsilon_c=\tau+\eta$.
The cubic comparison in Section~\ref{sec:normalized-cubic-identity}
will also make this character zero.
\end{corollary}
\begin{proof}
Use the signed matrices over $\mathbb F_3$
\[
 U=\begin{pmatrix}1&1\\0&1\end{pmatrix},\qquad
 S=\begin{pmatrix}0&1\\-1&0\end{pmatrix}.
\]
The earlier finite comparison proves
$f_g(U,S)=(-I)^{\eta(g)}$ on $D$.
Let $J=\operatorname{diag}(-1,1)$.  Then
$JUJ^{-1}=U^{-1}=U^2$ and $JSJ^{-1}=-S$.
A commutator proword is unaffected by multiplying either
input by a central element: its exponent sum in that input
is zero.  Hence
\[
 f_g(U^2,S)=f_g(U^2,-S)
          =Jf_g(U,S)J^{-1}=(-I)^{\eta(g)}.
\]
Proposition~\ref{prop:normalized-GF1-from-dyadic} identifies
the left side with $q_3(\beta_g)^{-1}=(-I)^{w(g)}$.
This proves $w=\eta$ and $\epsilon_h|_D=0$.
Its established extension is zero on $G_{\mathbb Q}$, and
$\widehat{\mathrm{GT}}=D G_{\mathbb Q}$, so it is zero
everywhere.  The final equality follows from the already
proved decomposition $\chi_3=\epsilon_h+\epsilon_c$.
\end{proof}

\section{The cubic branch filling and the normalized cubic identity}
\label{sec:normalized-cubic-identity}

We use the standard full genus-zero GT action and the exact
nonabelian $(K_3)$ identity. The rigidity input is
\begin{equation}\label{eq:cubic-normalizer-1}
\operatorname{Out}(\widehat\Gamma_{0,n})
 =\widehat{\mathrm{GT}}\times\mathfrak S_n\qquad(n\ge5),
\end{equation}
as proved by Hoshi--Minamide--Mochizuki
\cite[Definition~2.7 and Corollary~2.8(i), author pp.~31--32]{CubicHMM}.
Here the set of primes is the full set; their Definition~2.7
identifies the resulting GT group with the classical full
profinite GT group. The other topological inputs and their
precise scope are stated where they are used.

Schneps described the cyclic genus-zero quotient construction and its
discrete normalizers in \cite[Section~4.3,
Theorems~4.3.1--4.3.2, author pp.~54--58]{SchnepsSpecialLoci}.
The mixed version and its full profinite topology are checked below.
She introduced special homomorphisms from these cyclic automorphisms
of punctured spheres. Her subgroup $\widehat{\mathrm{GS}}^1_{0,0}$
is defined by two such compatibilities, and its elements satisfy
the normalized cubic relation $(\mathrm{NT}_2)$
\cite[Definition~1.2, p.~658; Theorems~3.1 and~3.4,
pp.~665 and~667]{Schneps2006Special}.
The argument here proves the cubic comparison on full GT with
multiplier one and vanishing $\rho_3^I$.


Throughout, $\Gamma_{0,n}$ denotes the pure mapping class
group of an oriented sphere with $n$ marked points. Brackets
around a set of marked points mean that its permutations are
allowed. A circumflex denotes full profinite completion.
Let $g=(1,f_g)$ belong to full profinite GT. Its standard genus-zero
action \cite[equation~(1.3), p.~657]{Schneps2006Special}
is represented by an automorphism $\alpha_g$
of each completed mixed mapping class group. It induces the
identity on the finite permutation quotient and the identity
on the abelianization of each pure subgroup. We distinguish
this action from its outer class on a pure subgroup.
The labelled forgetting maps are compatible with these actions
up to inner automorphisms by pure elements. One may choose
compatible based representatives for a single forgetting map
and obtain equality; changing such a choice adds a pure inner
automorphism. This is the usual labelled genus-zero action.
Below we keep these
pure inner corrections explicitly. On the four-point group it
has the usual outer class of
\begin{equation}\label{eq:cubic-normalizer-2}
\phi_g(x)=x,\qquad
 \phi_g(y)=f_g(x,y)^{-1}y f_g(x,y).
\end{equation}
The multiplier-one hypothesis is used here.

\subsection{The centralizer and its full completion}

Fix a primitive cube root $\omega$. On $\mathbb P^1$ mark
$0,\infty$ and three free orbits
\[
 D_{i,j}=\omega^j t_i,
 \qquad i\in\{1,2,3\},\quad j\in\mathbb Z/3,
 \qquad t_i\ne0,\quad t_i^3\ne t_k^3\ (i\ne k).
\]
Thus there are eleven marked points. Let $h(z)=\omega z$,
$\mathcal H=\langle h\rangle$, and let $\mathcal G$ be the subgroup of the
mixed mapping class group whose permutations fix $0,\infty$
individually and commute with the displayed permutation of
$h$. Put $\mathcal N=C_{\mathcal G}(\mathcal H)$. The cyclic quotient has coordinate
$w=z^3$, with its two branch points and three further marked
points. We first prove the discrete centralizer sequence
\begin{equation}\label{eq:cubic-normalizer-3}
1\longrightarrow \mathcal H\longrightarrow \mathcal N\longrightarrow
 Q_3\longrightarrow1,
 \qquad Q_3=\Gamma_{0,\{0\},\{\infty\},[3]}.
\end{equation}
The argument works with two or three free orbits. Delete all
marked points from the covering sphere and its quotient.
The resulting cover $S\to S/\mathcal H$ is unbranched, regular and
three-sheeted. Its monodromy character takes the values $1$
and $-1$ on positive loops around the two distinguished
branch labels and zero on the other punctures.

Every mapping class centralizing $h$ has a
representative commuting with $h$. Give $S$ its displayed
complex structure, in which $h$ is conformal, and choose the
extremal quasiconformal representative $n$ of that class.
Existence and uniqueness for a finitely punctured compact
surface follow from Gardiner--Hu, Section~5, Theorem~1,
and Section~8, Theorem~6
\cite[author pp.~14--15 and~22]{CubicGH}; the punctured
hypothesis is specified immediately before their Theorem~1.
Both $nh$ and $hn$ lie in the same homotopy class, and both
are extremal: any improvement, composed with $h^{-1}$ on
the appropriate side, would improve $n$. Uniqueness implies
that they differ by a conformal map homotopic to the identity.
Such a map fixes all punctures, hence at least three points
of the compactified sphere, and is the identity. Thus $nh=hn$.
This spells out in the punctured setting the extremal-map
argument of Birman--Hilden \cite[p.~433]{CubicBH}.

The representative $n$ descends to the punctured quotient.
If two such representatives are isotopic upstairs, their
quotient is fiber-preserving and isotopic to the identity.
Birman--Hilden Theorem~1
\cite[p.~424]{CubicBH}, applied to this unbranched cover,
makes it fiber-isotopic to the identity. Their descended
mapping classes consequently agree. Conversely every element
of $Q_3$ preserves the monodromy character and therefore
lifts; every lift commutes with $h$ and has exactly the
required label permutation. The lifts differ by $\mathcal H$.
Lifting an isotopy downstairs shows that the kernel is $\mathcal H$,
which proves \eqref{eq:cubic-normalizer-3}. This argument proves the required
liftability directly. It does not use Birman--Hilden's later
general lifting assertion in Lemma~5.1 or Theorem~5;
their erratum leaves Theorems~1--4 unaffected
\cite{CubicBHErr}.

We prove that $\mathcal N$ has its full profinite topology in the
ambient mapping class group. Let
$\Phi=\pi_1(S/\mathcal H)$ and $\Pi=\pi_1(S)$, with compatible
basepoints. Thus $\Pi\lhd\Phi$ is the kernel of the above
character $\chi:\Phi\to C_3$, and
$\widehat\Pi$ is open in the free profinite group
$\widehat\Phi$. Put
\[
 \mathcal L=
 \operatorname{Aut}_{\chi}(\widehat\Phi)
                 /\operatorname{Inn}(\widehat\Pi),
\]
where the subscript means preservation of $\chi$ pointwise.
Since $\widehat\Phi$ is finitely generated and
$\operatorname{Inn}(\widehat\Pi)$ is compact, these are
profinite groups. There is an exact sequence
\begin{equation}\label{eq:cubic-normalizer-4a}
1\longrightarrow C_3\longrightarrow\mathcal L
 \longrightarrow\operatorname{Out}_{\chi}(\widehat\Phi)
 \longrightarrow1.
\end{equation}
Indeed, its kernel is
$\operatorname{Inn}(\widehat\Phi)/
\operatorname{Inn}(\widehat\Pi)=\widehat\Phi/\widehat\Pi$;
the centre of a nonabelian free profinite group is trivial.

The lifted mapping classes give a homomorphism
$\mathcal N\to\mathcal L$. Choose a path upstairs from the basepoint
to its image under a commuting representative and project
that path downstairs. The resulting based automorphism of
$\Phi$ is well-defined modulo $\operatorname{Inn}(\Pi)$:
changing the path gives just such an inner automorphism, and
changing the representative is harmless by the preceding
fiber-isotopy argument. This construction is compatible with
composition. Its profinite extension fits with \eqref{eq:cubic-normalizer-3}, \eqref{eq:cubic-normalizer-4a}
and the ordinary outer action of $\widehat Q_3$ on
$\widehat\Phi$.

The completion of \eqref{eq:cubic-normalizer-3} has kernel exactly $\mathcal H$. This finite
subgroup is separated by its permutation of the upstairs
punctures; the universal property of completion identifies
$\widehat{\mathcal N}/\mathcal H$ with $\widehat Q_3$. Asada's faithfulness
theorem makes
$\widehat Q_3\to\operatorname{Out}(\widehat\Phi)$
injective in the pure case
\cite[introduction Theorem, p.~3, parts (i)--(ii)]{CubicAsada}.
The mixed case follows as well: a nontrivial puncture
permutation is visible on peripheral conjugacy classes,
and its permutation kernel is the pure subgroup. The map
$\mathcal H\to C_3$ in \eqref{eq:cubic-normalizer-4a} is an isomorphism, since a deck
transformation is represented by a projected path of the
corresponding monodromy. The two exact sequences therefore
make $\widehat{\mathcal N}\to\mathcal L$ injective.

Restriction $\mathcal L\to
\operatorname{Out}(\widehat\Pi)$ is injective. If a
representative restricts to an inner automorphism, adjust by
that inner automorphism so that it fixes $\widehat\Pi$
pointwise. For any $x\in\widehat\Phi$, comparison of its
conjugation on $\widehat\Pi$ then shows that
$x^{-1}\theta(x)$ centralizes $\widehat\Pi$.
This centralizer is trivial: it intersects $\widehat\Pi$
in the latter's trivial centre, so injects into the finite
group $\widehat\Phi/\widehat\Pi$; a free profinite group
has no nontrivial finite subgroup. Thus $\theta$ is the
identity. We have proved faithfulness of the full completed
$\mathcal N$-action on $\widehat\Pi$. This action factors through
the completion of the ambient mapping class group, so
$\widehat{\mathcal N}\to\widehat{\mathcal G}$ is injective. In particular
\begin{equation}\label{eq:cubic-normalizer-4}
\overline{\mathcal N}=\widehat{\mathcal N},\qquad
 \overline{\mathcal N}/\mathcal H=\widehat Q_3.
\end{equation}
This establishes the full profinite topology on $\mathcal N$.

Finally the Boggi--Zalesskii centralizer theorem gives
$C_{\widehat{\mathcal G}}(\mathcal H)=\overline{\mathcal N}$
\cite[Theorem~A(iii), p.~3; Section~6.1 and
Theorem~6.2(iii), pp.~16--17]{CubicBZ}.
The theorem applies to the present finite-index subgroup:
genus-zero mapping class groups are good and have finite
virtual cohomological type by the cited Section~6.1;
goodness and that finiteness property pass to finite-index
subgroups (by the finite-index cohomology comparison and
Shapiro's lemma). Residual finiteness is inherited as well.
Only finite $3$-subgroups are used.

\subsection{Correction and identification of the parameter}

Let $P=\Gamma_{0,11}$ and let
$\mathcal G_{\mathrm{cyc}}=P \mathcal H$ be the inverse image of the cyclic permutation
subgroup generated by $h$. Since $\alpha_g$ is the identity
on label permutations, $\alpha_g(\mathcal H)\subset\widehat{\mathcal G}_{\mathrm{cyc}}$.
The Boggi--Zalesskii finite-$3$-subgroup theorem
\cite[Theorem~A(i), p.~3; Theorem~6.2(i), p.~17]{CubicBZ}
applied to $\mathcal G_{\mathrm{cyc}}$ conjugates $\alpha_g(\mathcal H)$ into $\mathcal G_{\mathrm{cyc}}$.
The conjugator can be taken pure: write an initial conjugator
as $h^j p$ with $p\in\widehat P$, and absorb $h^j$ by
conjugating the resulting discrete subgroup. The resulting
generator has the same permutation of all eleven labels as
$h$, since the permutation quotient of $\mathcal G_{\mathrm{cyc}}$ is abelian.

Nielsen realization realizes this finite cyclic group by
conformal transformations of a punctured sphere. The required
finite-type, punctured scope is included in Kerckhoff's
extension to finitely generated Fuchsian groups with cusps
\cite[Section~4, pp.~255--256]{CubicKerckhoff}.
Putting its fixed points at $0,\infty$, the generator is
$z\mapsto\omega z$ or $z\mapsto\omega^{-1}z$.
In the first case its labelled configurations form the
connected space
\[
 \{(t_1,t_2,t_3)\in(\mathbb C^\times)^3:
                         t_i^3\ne t_k^3\text{ for }i\ne k\}.
\]
In the second case exchange the two fixed-point labels first.
It follows that there is an $e$ which is pure, or pure
followed by a transposition of $0,\infty$, and which fixes
every free-orbit label, such that
\begin{equation}\label{eq:cubic-normalizer-5}
\alpha'=\operatorname{Ad}(e^{-1})\alpha_g
       \quad\text{fixes }\mathcal H\text{ pointwise}.
\end{equation}
Here the assertion about the labels concerns the permutation
of $e$; its pure factor need not fix points pointwise as a
homeomorphism.

We verify invariance of the mixed groups before taking
outer classes. The permutation of $e$ is either the identity
or the single transposition $(0\ \infty)$; it fixes all nine
free labels. Either permutation centralizes the permutation
subgroup defining $\mathcal G$. Since $\alpha_g$ induces the identity
on permutations, $\alpha'$ preserves $\widehat{\mathcal G}$ and
induces the identity on its permutation quotient. Since it
fixes $\mathcal H$ pointwise, it preserves
$C_{\widehat{\mathcal G}}(\mathcal H)=\widehat{\mathcal N}$. Equation \eqref{eq:cubic-normalizer-4} therefore
gives an automorphism $\delta_3$ of $\widehat Q_3$.
It preserves its pure subgroup
$P_5=\widehat\Gamma_{0,5}$ and induces the identity on
$\widehat Q_3/P_5=\mathfrak S_3$.

Write $\vartheta=[\delta_3|_{P_5}]=(g',\sigma)$ by \eqref{eq:cubic-normalizer-1}.
For each $\pi\in\mathfrak S_3$, choose $z_\pi\in
\widehat Q_3$ over $\pi$. Conjugation by $z_\pi$ on
$P_5$ has outer class $(1,\pi)$ in \eqref{eq:cubic-normalizer-1}. As
$\delta_3(z_\pi)$ has the same permutation, the identity
\[
 \delta_3\operatorname{Ad}(z_\pi)\delta_3^{-1}
       =\operatorname{Ad}(\delta_3(z_\pi))
\]
restricted to $P_5$ gives
$\vartheta(1,\pi)\vartheta^{-1}=(1,\pi)$.
Hence $\sigma\pi\sigma^{-1}=\pi$ for every $\pi$.
The centralizer in $\mathfrak S_5$ of the full symmetric
group on the three orbit labels consists exactly of the
identity and the transposition of the other two labels.
Thus
\begin{equation}\label{eq:cubic-normalizer-6}
\sigma=1\quad\text{or}\quad\sigma=(0\ \infty).
\end{equation}
To forget the third orbit, first restrict $\widehat{\mathcal N}$
to the subgroup $\widehat{\mathcal N}^{\rm orb}$ fixing each free
orbit as a set; its quotient by $\mathcal H$ is $P_5$. This subgroup
is preserved by $\alpha'$, as just proved. Ambient forgetting
of all three labels in the third orbit maps it onto the
two-orbit centralizer $\widehat{\mathcal N}_2^{\rm orb}$.
On quotients by $\mathcal H$ this is the usual surjection
$P_5\to P_4=\widehat\Gamma_{0,4}$, and it is the
identity on $\mathcal H$. These facts prove surjectivity upstairs too.

Here are the inner corrections in this step. Choose standard
GT and permutation representatives on $P_5$. By definition
of an outer class, $\delta_3|_{P_5}$ is their product
followed by $\operatorname{Ad}(a_3)$ with $a_3\in P_5$.
Both possible $\sigma$ fix the deleted orbit label.
The standard representatives therefore preserve the
forgetting kernel, and so does $\delta_3$; the correction
becomes the image of $a_3$ in $P_4$. Thus the
induced automorphism $\delta_2$ has outer class
$(g',\sigma)$ on $P_4$.
On the ambient mixed group the same conclusion follows from
labelled forgetting compatibility. If a choice of based
representatives introduces a pure correction $c$, the
forgotten action is
$\operatorname{Ad}(\bar e^{-1}c)\alpha_g$.
Its effective conjugator $c^{-1}\bar e$ still fixes every
remaining free label and has exactly the same possible
branch transposition as $e$. Thus the label and homology
arguments below are unchanged. This notation also covers
all later forgetting maps: every extra based correction is
pure in its own target.

We make the resulting two-orbit calculation explicit. Scale
the first orbit to $\{1,\omega,\omega^2\}$, write the
second as $\{t,\omega t,\omega^2t\}$, and put $w=t^3$.
Then
\[
 \mathcal N_2/\mathcal H\cong\pi_1(\mathbb P^1_w\setminus\{0,1,\infty\})
       =F(x,y).
\]
The subgroup $\mathcal N_2^0$ fixing the chosen representatives $1,t$
fixes all six free-orbit labels and intersects $\mathcal H$ trivially.
Its image is
\begin{equation}\label{eq:cubic-normalizer-7}
H_3=\ker(F(x,y)\longrightarrow C_3),\qquad x\mapsto1,
 \quad y\mapsto0.
\end{equation}
It has basis $X=x^3$, $y_j=x^jyx^{-j}$, $j=0,1,2$.
The element $y_j$ is the positive simultaneous collision
multitwist pairing the three points of one free orbit with
the three points of the other, with relative offset $j$.
It is pure. Its abelian class $v_j$ in the ambient pure
eight-point group is fixed by $\alpha'$: $\alpha_g$ acts
as the identity on this abelianization and the only possible
permutation of $e$ exchanges $0,\infty$, fixing all the
paired labels.

The three $v_j$ are distinguished integrally by forgetting
maps. Retain $0,\infty$ and just one pair with offset $j$.
The class $v_j$ maps to the positive generator around their
collision in the four-point group, whereas the other two
classes map to zero. In particular it is nonzero and
primitive in this test quotient.

Put $\lambda'=\lambda(g')$. In \eqref{eq:cubic-normalizer-7}, an automorphism of
outer class $(g',\sigma)$ takes the conjugacy class of $y_j$
to that of
\begin{equation}\label{eq:cubic-normalizer-8}
y_{\varepsilon\lambda'j+c}^{\lambda'},
 \qquad \varepsilon=1\text{ if }\sigma=1,
 \quad\varepsilon=-1\text{ otherwise},
\end{equation}
with subscripts modulo three and a common offset $c$.
For the GT factor this follows directly from
$x\mapsto x^{\lambda'}$, $y\mapsto f_{g'}^{-1}y^{\lambda'}f_{g'}$;
the interchange of $0,\infty$ lifts to $t\mapsto1/t$ and
reflects the three collision points, and an inner correction
adds a common offset. Once the subscript in \eqref{eq:cubic-normalizer-8} has been
chosen, its conjugating element is in $\widehat H_3$, hence
is pure upstairs.
The preceding three abelian tests therefore give
\begin{equation}\label{eq:cubic-normalizer-9}
v_j=\lambda'v_{\varepsilon\lambda'j+c}
 \quad(j=0,1,2),\qquad
 \lambda'=1,\quad c=0,\quad\varepsilon=1.
\end{equation}
Thus $\sigma=1$. Moreover the two-orbit quotient
automorphism has the form
\begin{equation}\label{eq:cubic-normalizer-10}
\delta_2=\operatorname{Ad}(a)\phi_{g'},
                  \qquad a\in\widehat H_3,
\end{equation}
because its common offset is now the image of $a$ in $C_3$.

Forget all but $0,\infty,1,t$ upstairs. The resulting map
on \eqref{eq:cubic-normalizer-7} is exactly the standard cyclic-cover filling
\begin{equation}\label{eq:cubic-normalizer-11}
d_3:\widehat H_3\longrightarrow\widehat F(x,y),
 \qquad X\mapsto x,\quad y_0\mapsto y,
 \quad y_1,y_2\mapsto1.
\end{equation}
The exact $(K_3)$ identity says
$d_3(f_{g'})=y^{c_3(g')}f_{g'}$. In the automorphism formula
this extra power commutes with $y$, so
\begin{equation}\label{eq:cubic-normalizer-12}
d_3\phi_{g'}|_{\widehat H_3}=\phi_{g'}d_3.
\end{equation}
For clarity, \eqref{eq:cubic-normalizer-12} is checked on the four free generators:
on $X$ both sides give $x$; on $y_0$ both give
$f_{g'}^{-1}yf_{g'}$; and on $y_j$ for $j=1,2$ use
\[
 \phi_{g'}(y_j)=(x^j f_{g'}^{-1}x^{-j})\,y_j\,
                         (x^j f_{g'}x^{-j}),
\]
whose image under $d_3$ is one. The conjugating factors
belong to $H_3$ because $f_{g'}$ is a commutator.
Equation \eqref{eq:cubic-normalizer-10} now contributes the inner automorphism
$\operatorname{Ad}(d_3(a))$ in the target. Thus
\eqref{eq:cubic-normalizer-10}--\eqref{eq:cubic-normalizer-12} make the induced outer action downstairs
equal to that of $g'$. On the other hand, the ambient
forgetting compatibility and \eqref{eq:cubic-normalizer-5} make it the action of $g$,
possibly composed with the transposition of the chosen
labels $0,\infty$. That permutation is a single transposition on the
four chosen labels: the two free labels are individually
fixed. It is not a double transposition in the Klein four
kernel of the four-point permutation action. Equivalently,
in the coordinate $t$ it is $t\mapsto1/t$, which exchanges
the peripheral classes at zero and infinity. Both GT
factors and all pure inner corrections preserve the three
peripheral classes individually. The branch transposition
is therefore absent. Thus $e$ is pure and $g'=g$.

For completeness, the usual GT embedding into the outer
automorphism group of $\widehat F(x,y)$ is injective in this
comparison. An inner conjugation equating two standard
representatives with the same multiplier must centralize
$x$, so is a power of $x$. Comparing the conjugators of
$y$ and then taking their exponent sums makes that power
and the possible power of $y$ zero; the two commutator
parameters are equal.

\begin{lemma}[The other cubic filling]\label{lem:cubic-branch-filling}
Let $g=(1,f_g)\in\widehat{\mathrm{GT}}$.
Fill the punctures $0,\infty$ in the threefold cyclic cover
$\mathbb P^1_t\setminus\{0,\infty,1,\omega,\omega^2\}$.
The restriction of the standard action $\phi_g$ on the
base group, followed by this filling, induces on
$\pi_1(\mathbb P^1_t\setminus\{1,\omega,\omega^2\})$
the outer class of the same $g$.
\end{lemma}
\begin{proof}
Retain the three points $1,\omega,\omega^2$ of the first
orbit and the chosen point $t$ of the second, forgetting
$0,\infty$ and the other two points of the second orbit.
On $\mathcal N_2^0=H_3$ this is precisely the displayed branch
filling. The parameter comparison above gives
$\delta_2=\operatorname{Ad}(a)\phi_g$ with
$a\in\widehat H_3$, and the correction $e$ is pure.
Both corrections therefore become inner automorphisms in
the target four-point group. Ambient forgetting supplies
the standard outer action of $g$ there, proving the claim.
\end{proof}

\subsection{The exact cubic conjugator}

Let $\Gamma=\widehat{C_2*C_3}$ be the modular quotient of
the signed braid group on $u,v$, and put $r=uv$,
$a=u^2$, $b=v^2$. The subgroup $\mathcal E=\langle r,a\rangle$
has presentation
\[
 \mathcal E=\langle r,a\mid r^3=(ra)^3=1\rangle
       \cong C_3\amalg C_3.
\]
Its index-three subgroup $F=\langle a,b\rangle$ is free;
conjugation by $r$ cyclically permutes $a,b,(ab)^{-1}$.
It is the filled cover in Lemma~\ref{lem:cubic-branch-filling},
with $x\mapsto r$, $y\mapsto a$. These choices fix the
orientation and the order of the three peripheral loops.
The signs can also be seen directly. Send the distinguished
point $t$ of $\{1,\omega,\omega^2,t\}$ to infinity and make an
affine change of coordinates. The three moving points become
\[
 Z_i(t)=\omega^{2i}+t\omega^i\quad(0\leq i<3),
 \qquad Z^3-3tZ-(1+t^3)=0.
\]
A positive loop of $w=t^3$ around zero ends at $\omega t$;
the endpoint configuration is $\omega^2$ times the initial
one. Closing by a positive one-third rotation gives $uv=r$.
At $w=1$, the difference
$Z_1(t)-Z_2(t)=(\omega^2-\omega)(1-t)$ makes a positive
full turn, giving $u^2=a$. A full-rotation ambiguity is central
and disappears in the modular quotient.

\begin{proposition}\label{prop:normalized-HF1-from-normalizer}
If $g=(1,f_g)\in\widehat{\mathrm{GT}}$ and
$\rho_3^I(g)=0$, then its pure cubic conjugator satisfies
\begin{equation}\label{eq:cubic-normalizer-13}
\boxed{\quad A_g(u^2,v^2)^{-1}=f_g(u^2,uv)\quad}
\end{equation}
in the full signed profinite braid group.
\end{proposition}
\begin{proof}
The standard action $\phi_g$ descends to $\mathcal E$: it preserves
the two finite peripheral relations because it sends each
peripheral generator to a conjugate of itself. Denote this
action by $\psi$. Put $H_g=f_g(a,r)$. The two-cycle relation
gives
\begin{equation}\label{eq:cubic-normalizer-14}
\psi(r)=r,\qquad\psi(a)=H_g aH_g^{-1}.
\end{equation}
The braid action $\Phi_g$ preserves $\mathcal E$ and has
$\Phi_g(a)=a$, $\Phi_g(r)=A_g rA_g^{-1}$.
Lemma~\ref{lem:cubic-branch-filling} says that on the open
free subgroup $F$ one has
\begin{equation}\label{eq:cubic-normalizer-15}
\psi|_F=\operatorname{Ad}(k)\Phi_g|_F
                  \quad\text{for some }k\in F.
\end{equation}
The inner correction is in $F$, since the target four-point
group is pure and all corrections in the preceding lemma
were pure. Two automorphisms of $\mathcal E$ agreeing on $F$ agree
everywhere: their quotient fixes $F$ pointwise, and comparison
of conjugation on $F$ shows that its discrepancy at every
element centralizes $F$, whose centralizer is trivial.
Thus \eqref{eq:cubic-normalizer-15} holds on $\mathcal E$.

Applying it to $r$ gives $kA_g\in C_\mathcal E(r)=\langle r\rangle$.
Since $k,A_g\in F$ and $F\cap\langle r\rangle=1$, this
forces $k=A_g^{-1}$. Applying \eqref{eq:cubic-normalizer-15} to $a$ now gives
$A_gH_g\in C_F(a)=\overline{\langle a\rangle}$.
Write $A_gH_g=a^t$ with $t\in\widehat{\mathbb Z}$.

The factor $A_g$ has abelianization zero. This follows from
$f_g=\operatorname{Ad}(r)(A_g)A_g^{-1}$: on $F^{\rm ab}$
the order-three rotation minus the identity has determinant
three, and multiplication by three on $\widehat{\mathbb Z}$
is injective, so that rotation has no fixed vector.
Write the class of $f_g(x,y)$ in the full Schreier
abelianization of $H_3$ as
\[
 (0,\psi_0,\psi_1,\psi_2),\qquad
 \psi_0+\psi_1+\psi_2=0.
\]
The first coordinate and the displayed sum vanish because
$f_g$ is a commutator. Ihara's full positive-root normalization
is $\rho_3^I=-\psi_0$
\cite[Section~5, p.~205]{NT2002}.
Branch filling sends $y_0,y_1,y_2$ to $a,b,(ab)^{-1}$.
The two-cycle relation therefore gives
\[
 [H_g]=(-\psi_0+\psi_2,-\psi_1+\psi_2)=(t,0).
\]
The second coordinate gives $\psi_1=\psi_2$, while
$\rho_3^I(g)=0$ gives $\psi_0=0$. Their sum is zero and
multiplication by two on $\widehat{\mathbb Z}$ is injective;
hence all three coordinates vanish and $t=0$.
This proves \eqref{eq:cubic-normalizer-13} in
the modular quotient.

Finally both sides of \eqref{eq:cubic-normalizer-13} have exponent sum zero in the
signed braid group: $A_g\in F'$ and $f_g$ is a commutator
word. The kernel of the modular quotient is generated
topologically by $(uv)^3$, whose exponent sum is six.
Multiplication by six on $\widehat{\mathbb Z}$ is injective,
so the possible central discrepancy is zero. This proves
the signed identity exactly.
\end{proof}

In the modular quotient the identity also gives the explicit
normalized cubic relation
\[
 f_g(a,b)=f_g(r,b)f_g(r,a)^{-1}.
\]
Indeed $f_g(a,b)=\operatorname{Ad}(r)(A_g)A_g^{-1}$ and
$\operatorname{Ad}(r)(a)=b$. Substituting
$A_g^{-1}=f_g(a,r)$ and using the two-cycle identity proves
the displayed formula. Its inverse is equation~(3.10) of
\cite[p.~667]{Schneps2006Special}; that source derives it on
$\widehat{\mathrm{GS}}^1_{0,0}$ by forgetting two strands.


\begin{corollary}\label{cor:cubic-defect-zero-normalizer}
On $D=\ker\lambda\cap\bigcap_q\ker\rho_q^I$ the finite
cubic comparison defect $\epsilon_c=\tau+\eta$ vanishes.
Together with the harmonic comparison, this gives
$\chi_3=0$ on full GT.
\end{corollary}
\begin{proof}
Reduce \eqref{eq:cubic-normalizer-13} modulo eight under the signed representation.
Writing $X=U^2$ and $R=-UL$, its right side is $f_g(X,R)$:
the central sign in $q(uv)=-R$ has no effect on a commutator
word. The previously proved finite interchange map sends
$q_8(A_g)^{-1}$ to $(-I)^\tau$ and $f_g(X,R)$ to
$f_g(U,S)^{-1}=(-I)^\eta$. Hence $\tau=\eta$.
The relative extension of $\epsilon_c$ is zero on the Galois
subgroup, and $\widehat{\mathrm{GT}}=D G_{\mathbb Q}$,
so it is zero everywhere. The already proved harmonic
equality $w=\eta$ now gives $\chi_3=w+\tau=0$.
\end{proof}

\subsection{The cusp-swap form of the level-three gate}

Write \(\mathcal M=\operatorname{PSL}_2(\mathbf Z)\), with generators
\[
 u=\begin{pmatrix}1&1\\0&1\end{pmatrix},\qquad
 v=\begin{pmatrix}1&0\\-1&1\end{pmatrix},\qquad s=uvu.
\]
All matrix identities in this subsection are in the centreless
group. Put \(K_0=\Gamma_0(3)=\langle u,v^3\rangle\) and
let \(d_3:K_0\to\mathcal M\) be conjugation by
\(\operatorname{diag}(3,1)\). These maps also denote their
extensions to full profinite completions.

\begin{lemma}\label{lem:level3-cusp-swap-gate}
There is an involution \(\vartheta_3\) of \(K_0\) with
\[
 \vartheta_3(u)=v^3,\qquad \vartheta_3(v^3)=u,\qquad
 \vartheta_3=\operatorname{Ad}(s)\circ d_3.
\]
For an element \(g=(1,f_g)\) of full profinite GT with
\(\rho_2(g)=0\),
put \(q=f_g(u^2,v^2)\). Then
\[
 d_3(q)=q\quad\Longleftrightarrow\quad
 \vartheta_3(q)=q^{-1}.
\]
\end{lemma}

\begin{proof}
The rational matrix
\[
 W_3=\begin{pmatrix}0&1\\-3&0\end{pmatrix}
       =s\operatorname{diag}(3,1)
\]
conjugates \(u\) to \(v^3\), and \(v^3\) to \(u\).
Its square is scalar. It therefore normalizes \(K_0\) and
induces the displayed involution. The same identities hold
after full profinite completion, since they are identities
between homomorphisms of finite-index discrete groups.

Normalized congruence triviality puts \(q\) in
\(\overline{K_0}\). Conjugation by \(s\) swaps \(u\) and \(v\);
the two-cycle identity consequently gives
\[
 s q s^{-1}=f_g(v^2,u^2)=q^{-1}.
\]
Combining this with \(\vartheta_3=\operatorname{Ad}(s)d_3\)
proves the equivalence.
\end{proof}

\begin{lemma}\label{lem:level3-principal-image}
Let \(P=\Gamma(3)\), the kernel of
\(\mathcal M\to\operatorname{PSL}_2(\mathbf F_3)\). Then
\[
 d_3(P)=\Gamma^0(9),\qquad
 \vartheta_3(P)=\Gamma_0(9).
\]
In particular \(\vartheta_3\) does not preserve \(P\).
The intersection \(P\cap\Gamma_0(9)\) has index \(36\) in
\(\mathcal M\); its modular curve has genus one and six cusps, of
widths \(3,3,3,9,9,9\).
\end{lemma}

\begin{proof}
A matrix in \(P\) has, up to its common scalar sign, diagonal
entries \(1\) modulo \(3\), and off-diagonal entries divisible
by \(3\). Conjugating by \(\operatorname{diag}(3,1)\)
multiplies its upper-right entry by \(3\) and divides its
lower-left entry by \(3\). Its image therefore has upper-right
entry divisible by \(9\).
Conversely a determinant-one integral matrix with upper-right
entry divisible by \(9\) has both diagonal entries congruent
to the same sign modulo \(3\). Its inverse conjugate belongs
to \(P\). This proves the first equality; conjugation by \(s\)
gives the second. For example \(v^3\in P\) does not belong
to \(\Gamma_0(9)\).

The group \(P\) is torsion-free of index \(12\) in \(\mathcal M\).
It has four cusps, each of width \(3\): the image of \(u\)
has order \(3\) in \(\operatorname{PSL}_2(\mathbf F_3)\).
Its Euler characteristic is \(-12/6=-2\), so its compactified
modular curve has genus zero.

On \(P\), choose the representative congruent to \(I\)
modulo \(3\). Its lower-left entry divided by \(3\), reduced
modulo \(3\), is a surjective homomorphism to
\(\mathbf F_3\). Its kernel is \(P\cap\Gamma_0(9)\);
thus the latter has index three in \(P\).
For a cusp represented by \(p/q\), a positive peripheral
matrix is
\[
 I+3\begin{pmatrix}-pq&p^2\\-q^2&pq\end{pmatrix}.
\]
Its value under this homomorphism is \(-q^2\) modulo \(3\).
The cusp at infinity therefore has three unramified lifts,
of width \(3\). Each of the other three cusps has one lift
of ramification index three, of width \(9\).
There are six cusps in all. Riemann--Hurwitz for this
three-sheeted cover, branched at three points, gives
\[
 2g-2=3(-2)+3(3-1)=0.
\]
Thus \(g=1\), as claimed.
\end{proof}

\begin{lemma}\label{lem:level3-cubic-frame-principal}
For the same normalized element \(g\), let \(A_g\) be the
normalized pure cubic conjugator, so that
\[
 \Phi_g(uv)=A_g(uv)A_g^{-1},\qquad
 A_g\in\widehat F(u^2,v^2)'.
\]
Then \(A_g\in\overline{\Gamma(6)}\subseteq\overline{\Gamma(3)}\).
If also \(\rho_3^I(g)=0\), the normalized cubic identity gives
\[
 f_g(u^2,uv)\in\overline{\Gamma(6)}.
\]
\end{lemma}

\begin{proof}
Normalized congruence triviality makes the image of \(A_g\)
centralize the image of \(uv\) in
\(\operatorname{PSL}_2(\mathbf F_3)\cong A_4\).
The images of \(u^2,v^2\) generate \(A_4\), so the image of
their commutator subgroup is contained in its derived subgroup
\(V_4\). The centralizer of a nontrivial \(3\)-cycle in
\(A_4\) is its cyclic group of order three, whose intersection
with \(V_4\) is trivial. Hence \(A_g\) has trivial image. Since it already lies in the closure
of \(\Gamma(2)\), it lies in the closure of
\(\Gamma(2)\cap\Gamma(3)=\Gamma(6)\).
Intersections of finite-index subgroup closures equal the closures of
their intersections.
The last assertion follows from
\(A_g^{-1}=f_g(u^2,uv)\).
\end{proof}

These reductions leave an exact full-group question:
whether the restricted conjugator satisfies
\(\vartheta_3(q)=q^{-1}\). The principal level-three group
and its image under the cusp swap have genus-zero quotients,
but their common cover has genus one. The genus-zero
rigidity theorem therefore does not apply to that common
cover. No vanishing of the full degeneracy defect is
asserted here.

\begin{lemma}[The level-two marked cover and its regular closure]\label{lem:level3-marked-regular-closure}
Let \(F=\Gamma(2)\) and \(J=F\cap\Gamma_0(3)\).
The inclusion \(J\subset F\) has index four and its
compactified modular curve has genus zero, with cusp
widths \(2,2,2,6,6,6\).
Its regular closure over \(F\) is \(\Gamma(6)\), with
deck group \(A_4\). The latter group has index \(72\) in
\(\operatorname{PSL}_2(\mathbf Z)\), twelve cusps of width
six, and compactified genus one.
\end{lemma}
\begin{proof}
Reduction modulo three maps \(F\) onto
\(\operatorname{PSL}_2(\mathbf F_3)=A_4\): the images of
\(u^2,v^2\) generate this group. The subgroup \(J\) is the
inverse image of a cyclic subgroup of order three, namely
the stabilizer of one point of \(\mathbf P^1(\mathbf F_3)\).
Thus the coset action has degree four. Each of the three
positive peripheral elements of \(F\) maps to a nontrivial
three-cycle. Its cycle structure on the four cosets is
\(1+3\). Each of the three base cusps, of width two,
therefore has one lift of width two and one of width six.
The index of \(J\) in the modular group is \(24\), and its
six cusps give
\[
 2-2g-6=-24/6=-4,
\]
hence \(g=0\).

The core of a point stabilizer in the faithful \(A_4\)
action on four points is trivial. The core of \(J\) in
\(F\) is consequently the kernel of this reduction, which
is \(\Gamma(2)\cap\Gamma(3)=\Gamma(6)\).
Its modular index is \(6\cdot12=72\). The image of \(u\)
in the corresponding modular quotient has order six,
so all its cusp widths are six and there are twelve cusps.
Finally
\[
 2-2g-12=-72/6=-12
\]
gives \(g=1\).
\end{proof}

This calculation concerns the cover with the level-two
marking preserved. Its genus-one regular closure is different
from the genus-zero curve \(X(3)\). It does not preclude a
different quotient or filling argument, but genus-zero
rigidity cannot be applied to this regular closure.

\begin{proposition}[An explicit free-group form of the level-three gate]\label{prop:level3-free-subgroup}
Put \(F=\langle a,b\rangle=\Gamma(2)\), where \(a=u^2,b=v^2\),
and \(J=F\cap\Gamma_0(3)\). Then \(J\) has the free basis
\[
 a,\quad b^3,\quad b^2ab^{-1},\quad ba^2b^{-2},
       \quad baba^{-1}b^{-1}.
\]
The restriction of the level-three degeneracy map is the
homomorphism \(d:J\to F\) determined by
\[
\begin{array}{c|c}
 \text{generator of }J&\text{image under }d\\ \hline
 a&a^3\\
 b^3&b\\
 b^2ab^{-1}&(ba)^3\\
 ba^2b^{-2}&a^{-1}b^{-2}a^{-1}b^{-1}\\
 baba^{-1}b^{-1}&a^{-1}b^3a
\end{array}
\]
These formulas extend to full profinite completions. For
\(g=(1,f_g)\) in full GT with \(\rho_2(g)=0\), the exact
level-three gate is therefore the explicit equality
\[
             \widehat d(f_g(a,b))=f_g(a,b)
                       \quad\text{in }\widehat F.
\]
No equality in this last display is asserted by this proposition.
\end{proposition}
\begin{proof}
The left cosets of \(J\) in \(F\) are identified with the
projective lower rows of matrices modulo three. A Schreier
transversal is \(1,b,b^2,ba\), corresponding to
\([0:1],[1:1],[2:1],[1:0]\). Right multiplication by the
positive free generators gives the table
\[
\begin{array}{c|cc}
 t&\overline{ta}&\overline{tb}\\\hline
 1&1&b\\
 b&ba&b^2\\
 b^2&b&1\\
 ba&b^2&ba
\end{array}
\]
The five nontrivial Schreier words
\(t x\,\overline{tx}^{-1}\) are exactly the displayed basis.

For the signed matrices
\[
 a=\begin{pmatrix}1&2\\0&1\end{pmatrix},\qquad
 b=\begin{pmatrix}1&0\\-2&1\end{pmatrix},
\]
conjugation by \(\operatorname{diag}(3,1)\) sends
\((p,q;r,s)\) to \((p,3q;r/3,s)\).
The five basis words above have lower-left entry divisible
by three. Substitution verifies all five displayed images.
The identification with the modular degeneracy map follows
from its definition as this conjugation. The faithful
embedding of the discrete free group in the modular group
then proves the word identities.

A finite-index subgroup has its full profinite topology.
Both inclusions \(J\to F\) and the map \(d\), whose image
is the finite-index subgroup \(F\cap\Gamma^0(3)\), therefore
extend injectively to completions. Normalized congruence
triviality puts \(f_g(a,b)\) in \(\widehat J\).
The two elements in the last equality belong to
\(\widehat F\), whose inclusion in the completed modular
group is injective. Their equality is thus exactly the
previous full modular gate.
\end{proof}

The permutation action on the four cosets is the primitive
action of \(A_4\) on four points: its point stabilizer has
order three and is maximal. Consequently this particular
degree-four cover has no nontrivial intermediate cover.
It cannot be decomposed as two successive degree-two
cyclic covers. This does not rule out an identity deduced
from a larger system of covers, but the cyclic index-three
\((K_3)\) relation is a different restriction and filling map.

\section{The tame quartic gate as a finite word evaluation}

Put $G=\widehat{\mathrm{GT}}$ and
\[
 D=\ker\lambda\cap\bigcap_{q\ \mathrm{prime}}\ker\rho_q^I.
\]
Write $g=(1,f_g)$ for an element of $D$, and use the signed
matrices
\[
 U=\begin{pmatrix}1&1\\0&1\end{pmatrix},\qquad
 L=\begin{pmatrix}1&0\\-1&1\end{pmatrix},\qquad
 S=ULU=\begin{pmatrix}0&1\\-1&0\end{pmatrix}.
\]
For an odd prime $\ell$, a subscript $\ell$ on these matrices
will mean reduction in $\mathrm{SL}_2(\mathbb F_\ell)$.
Let
\[
 \chi_{-4}(\ell)=\left(\frac{-1}{\ell}\right),\qquad
 m_\ell=\frac{\ell-\chi_{-4}(\ell)}2
 =\begin{cases}(\ell-1)/2,&\ell\equiv1\pmod4,\\
               (\ell+1)/2,&\ell\equiv3\pmod4.
   \end{cases}
\]

\begin{proposition}\label{prop:tame-harmonic-word-gate}
For every $g\in D$ the corrected harmonic frame satisfies
\begin{equation}\label{eq:tame-harmonic-1}
W_g=f_g(U^2,S)\qquad\text{in }\mathrm{SL}_2(\widehat{\mathbb Z}).
\end{equation}
Consequently its residue belongs to the cyclic group
\[
 \mathcal T_\ell
 =\{aI+bS_\ell:a,b\in\mathbb F_\ell,\ a^2+b^2=1\},
 \qquad |\mathcal T_\ell|=\ell-\chi_{-4}(\ell).
\]
The simultaneous compatibility of the radical and CM actions
on every rational quartic root is equivalent to the following
family of finite word identities:
\begin{equation}\label{eq:tame-harmonic-2}
\boxed{\quad
 f_g(U_\ell^2,S_\ell)^{m_\ell}=I
 \qquad(g\in D,\ \ell\ge5\text{ prime}).\quad}
\end{equation}
For a fixed $g$, the same equivalence holds pointwise. This is
a reformulation of the remaining tame quartic gate; it does
not assert the identities in \eqref{eq:tame-harmonic-2}.
\end{proposition}
\begin{proof}
The harmonic frame is $W_g=q(\beta_g)^{-1}D_g$.
On $D$, the established signed normalization gives $D_g=I$.
Proposition~\ref{prop:normalized-GF1-from-dyadic} gives
\[
 \beta_g(u^2,v^2)^{-1}=f_g(u^2,uvu)
\]
in the signed profinite braid group. Applying its continuous
signed representation $q(u)=U$, $q(v)=L$ proves \eqref{eq:tame-harmonic-1}, and
evaluation of a profinite word commutes with reduction modulo
$\ell$.

The frame belongs to the integral norm-one torus of
$\mathbb Z[i]$, represented by $i\mapsto S$.
Reduction gives the displayed group $\mathcal T_\ell$.
When $\ell\equiv1\pmod4$, its two eigenspaces identify it
with $\mathbb F_\ell^\times$; when $\ell\equiv3\pmod4$,
it is the kernel of the norm
$\mathbb F_{\ell^2}^\times\to\mathbb F_\ell^\times$.
Thus it is cyclic of the stated even order in both cases.
Its square subgroup is exactly the kernel of the power
$z\mapsto z^{m_\ell}$. The kernel of reduction from the
integral local torus is pro-$\ell$, so squaring is an
automorphism on that kernel. Hence a local frame component
at odd $\ell$ is a square if and only if its residue passes
the displayed power test.

Lemma~\ref{lem:lr-quartic-cm-quotient} identifies the
maximal elementary-two subextension of
$\mathbb Q(i)^{\mathrm{ab}}/\mathbb Q^{\mathrm{ab}}$
with
\[
 \mathbb Q^{\mathrm{ab}}(q^{1/4}:q\text{ rational prime}).
\]
Accordingly the action on all these quartic roots is trivial
if and only if $W_g$ is a square in the normalized adelic norm-one
group $N_4(4)$. The radical action is already trivial for $g\in D$.
At two the proved prime-two compatibility places $W_{g,2}$
in the local square subgroup. At three the established
harmonic residue is $W_g\bmod3\in\{I,-I\}$, which is
the square subgroup of the residue torus of order four;
the pro-three reduction kernel contributes no obstruction.
Thus precisely the primes $\ell\ge5$ remain.

Finally, squaring in the product defining $N_4(4)$ has image
equal to the product of their square subgroups: local square
roots may be chosen independently, and their tuple is an
adelic element of the required product. This retains the
simultaneous condition, rather than only checking a separate
list of necessary conditions. Substitution of \eqref{eq:tame-harmonic-1} now proves
both the pointwise and the universal equivalences in \eqref{eq:tame-harmonic-2}.
\end{proof}

\section{A square-lattice comparison for the harmonic frame}
\label{sec:harmonic-beta-square-lattice}

Put \(N=\ker(\lambda,\rho_2)\) in full profinite GT.
We use the full mixed genus-one action, its forgetting compatibility,
and its two low-genus descriptions from
\cite[Lemmas~16--17]{H17}. We also use the normalized harmonic identity
on all of \(N\), Proposition~\ref{prop:normalized-GF1-from-dyadic}, and
the signed congruence identity established above.
No compatibility with an arbitrary finite cover is assumed.
Here \(B_g^{\mathrm{ad}}\) denotes the normalized full adelic
Alexander unit in the completed group algebra
\(\widehat{\mathbf Z}[[\widehat{\mathbf Z}^{\,2}]]\):
\[
 B_g^{\mathrm{ad}}(X,Y)
       =1+(X-1)(Y-1)R_g^{\mathrm{ad}}(X,Y).
\]
Its value at \((i,-1)\) is defined by the continuous finite-character
map to \(\mathbf Z_\ell[i]\) at each coefficient prime \(\ell\).
At two, its augmentation branch also admits the usual dyadic formal
power-series description; the finite-character definition is used
at every odd prime. Throughout this section only, abbreviate
\(B_g^{\mathrm{ad}}\) by \(B_g\). Write
\[
 O_4=\mathbf Z[i],\qquad
 \nu_{4,g}=B_g(i,-1).
\]

\begin{theorem}\label{thm:harmonic-beta-square-lattice}
There is one integral identification of the harmonic quadratic order
with \(\widehat O_4\), valid for every \(g\in N\) and every coefficient
prime, under which the corrected harmonic frame \(W_g\) is
\(\nu_{4,g}\). Reversing the choice of primitive fourth root conjugates
all these values simultaneously.
Consequently, on \(D=N_{\rm rad}\), the tame harmonic square
test at an odd prime \(\ell\) is equivalent to
\[
 B_{g,\ell}(i,-1)^{(\ell-\chi_{-4}(\ell))/2}=1
\]
in the residue norm-one group. Either choice of the primitive
fourth root gives the same test.
\end{theorem}

\subsection{The finite translation group and the quotient action}

Let \(T=\mathbf R^2/\mathbf Z^2\) and mark all sixteen points
of \(T[4]\), with their individual labels. In the full mixed
mapping class group \(\Gamma=\Gamma_{1,[16]}\), the translations
form \(H=T[4]\simeq(\mathbf Z/4)^2\).
Let \(\alpha=\alpha_P\) be the mixed action of \(g\).
The finite-\(2\)-subgroup conjugacy theorem
\cite[Theorem~A(i); Theorem~6.2(i)]{CubicBZ}
puts \(\alpha(H)\) into a discrete finite subgroup by a pure
profinite conjugation. To obtain purity, apply the theorem
in the inverse image of the translation permutation group
\(H\); its finite quotient is abelian, so the discrete
translation part of an initial conjugator can be absorbed
without changing the labelled generators. The discrete group obtained in this
way has the same permutation on the sixteen labels as \(H\).
The homology covariance of \(\alpha\) makes its action on
closed-torus homology trivial. Nielsen realization therefore
realizes it by translations on a complex torus \(T'\)
\cite[Section~4]{CubicKerckhoff}, simply
transitive on the marked points. After choosing an origin,
the labelled translation group is an isomorphism
\(H\to T'[4]\).

We check the orientation of the marking that identifies these
two translation actions. After an oriented linear identification
of the tori, the labelled isomorphism is a matrix in
\(\mathrm{GL}_2(\mathbf Z/4)\). It lifts to
\(\mathrm{GL}_2(\mathbf Z)\): matrices of determinant one lift
from \(\mathrm{SL}_2(\mathbf Z/4)\), and the other determinant
class is represented by a reflection. Let its orientation
be \(\varepsilon\in\{1,-1\}\).
We spell out the point-pushing input that determines this sign.
Put \(\Pi=\pi_1(T\setminus T[4],*)\). Add a seventeenth marked
point and use the forgetting-compatible action on its stabilizer.
Its kernel is \(\widehat\Pi\), so this action supplies
\(\eta\in\operatorname{Aut}(\widehat\Pi)\) such that the
outer action of \(\alpha(h)\) is that of
\(\eta h\eta^{-1}\) for every \(h\in H\).
This is precisely the point-pushing construction and positive
puncture-loop assertion of \cite[Lemma~16(b),(c)]{H17}.
Let \(e\) be the initial pure profinite correction and let
\(r\) be the discrete labelled marking of orientation
\(\varepsilon\) above, so that
\(\alpha(h)=erhr^{-1}e^{-1}\).
The outer automorphism induced by \(\eta^{-1}er\) centralizes
the outer deck action of \(H\) on \(\widehat\Pi\), and it
carries each puncture loop to an \(\varepsilon\)-th power
up to \(\widehat\Pi\)-conjugacy. The signs here follow from
the positive loop assertion for \(\eta\), purity of \(e\),
and the orientation of \(r\).

Let
\[
 F=\pi_1\bigl((T\setminus T[4])/H\bigr),\qquad
 \widehat\Pi=\ker(\widehat F\longrightarrow H),
\]
where \(F\) is free of rank two and the quotient is
abelianization modulo four. Conjugation identifies
\(\widehat F\) with the full preimage of the outer deck
subgroup \(H\) under
\[
 \operatorname{Aut}(\widehat\Pi)
        \longrightarrow\operatorname{Out}(\widehat\Pi).
\]
Indeed its kernel is \(\operatorname{Inn}(\widehat\Pi)\),
its image on the outer quotient is exactly \(H\), and
its conjugation map is injective since the centralizer of
an open nonabelian free subgroup is trivial.
Choose an automorphism representing the outer class of
\(\eta^{-1}er\). Conjugation by it preserves this full
preimage and therefore defines
\(\Theta\in\operatorname{Aut}(\widehat F)\) acting as
that representative on \(\widehat\Pi\) and as the identity
on \(H\).
A positive boundary loop of the once-punctured quotient
torus is a puncture loop in \(\Pi\). Hence \(\Theta\)
sends it to a conjugate of its \(\varepsilon\)-th power.
In the class-two quotient it is the commutator of a free
basis, so \(\det\Theta^{\mathrm{ab}}=\varepsilon\).
Since \(\Theta\) acts identically modulo four, this
number is one modulo four. Thus \(\varepsilon=1\).
The resulting correction is pure and orientation preserving,
and we obtain
\begin{equation}
                  \alpha'|_H=\mathrm{id}_H.                 \label{eq:harmonic-beta-q1}
\end{equation}
This is the point-pushing and determinant argument of
\cite[Proposition~18, Step~1]{H17}, with the characteristic
four-division cover specified explicitly. It uses no
lifting-surjectivity assertion for a cyclic isogeny.

The discrete normalizer of \(H\) is
\(H\rtimes\mathrm{SL}_2(\mathbf Z)\). Its subgroup fixing
the origin is \(N_0=\mathrm{SL}_2(\mathbf Z)\), and forgetting
the other fifteen points identifies its completion with
\(\widehat\Gamma_{1,1}\). The finite-subgroup normalizer
theorem \cite[Theorem~A(iii); Theorem~6.2(iii)]{CubicBZ}
and \eqref{eq:harmonic-beta-q1} show that \(\alpha'\) preserves this completed
subgroup. Forgetting compatibility gives
\begin{equation}
     \alpha'|_{\widehat N_0}
        =\operatorname{Ad}(c_2)\Phi_g
              \quad(c_2\in\widehat\Gamma_{1,1}).             \label{eq:harmonic-beta-q2}
\end{equation}
Its central involution \(\iota=-1\) is fixed.
Moreover \(c_2\) is pure modulo two. Indeed \eqref{eq:harmonic-beta-q1} makes the
action on \(N_0\) trivial modulo four, and normalized
congruence triviality does the same for \(\Phi_g\).
Thus \(c_2\bmod2\) centralizes
\(\mathrm{SL}_2(\mathbf F_2)\), whose centre is trivial.

Take the centralizer of \(\iota\) in the subgroup fixing the
labelled origin. Its quotient by \(\langle\iota\rangle\) is
the mixed mapping class group \(\mathcal D\) of a sphere
with ten marked points: four branch points \(b_0,b_1,b_2,b_3\),
the images of \(T[2]\), and six other points, the images of
the pairs \(\{t,-t\}\subset T[4]\setminus T[2]\).
The group fixes \(b_0\) and permits permutations within the
other three branch labels and within the six pair labels.
The centralizer has commuting representatives: take the
unique extremal representative in a fixed punctured conformal
structure for which \(\iota\) is conformal. Composing on
either side with \(\iota\) gives extremal representatives
of the same class, so they agree. Existence and uniqueness
in this finite-type setting are the inputs
\cite[Sections~5 and~8, Theorems~1 and~6]{CubicGH}.
Every such mapping class lifts: after deleting all ten
points, the monodromy of the double cover is one at each
branch point and zero at the other six. The lifts differ
by \(\iota\). The unbranched fibre-isotopy theorem
\cite[Theorem~1]{CubicBH} makes descent independent of the
chosen commuting representative. This proves the discrete
quotient sequence using the unchanged part of that theorem.

The sequence has the required full profinite topology.
The upstairs and downstairs punctured-surface groups are
an open pair of nonabelian free groups. The completed
downstairs mixed mapping class group acts faithfully on the
completed downstairs surface group: Asada's theorem
\cite[introduction Theorem, (i)--(ii)]{CubicAsada} gives
the pure assertion, and a nontrivial permutation is visible
on peripheral conjugacy classes. Restriction to the open
upstairs group is faithful after quotienting by the deck
group, because an automorphism fixing an open free subgroup
pointwise is the identity. The deck involution itself is
detected by its permutation of the twelve nonfixed labels.
It follows that the completed centralizer embeds in the
completed ambient group. The finite-\(2\)-subgroup
centralizer theorem then identifies it with the full
profinite centralizer of \(\iota\).
Consequently \(\alpha'\) induces an automorphism
\(\delta_{10}\) of \(\widehat{\mathcal D}\).

On the pure ten-point group, genus-zero rigidity
\cite[Definition~2.7 and Corollary~2.8(i)]{CubicHMM} gives
\[
 [\delta_{10}]=(g',\sigma)
 \quad\text{in }\operatorname{Out}(\widehat\Gamma_{0,10})
       =\widehat{\mathrm{GT}}\times\mathfrak S_{10}.
\]
All sixteen labels were marked, and the correction in \eqref{eq:harmonic-beta-q1}
was pure. Thus the induced action on
\(\mathfrak S_3\times\mathfrak S_6\) is the identity.
The centralizer of this permutation subgroup in
\(\mathfrak S_{10}\) is trivial, so \(\sigma=1\).
To identify \(g'\), retain the four branch points and one
pair label. Upstairs this is forgetting all points except
the corresponding pair \(\{t,-t\}\); quotienting by
\(\iota\) gives the five-point description of
\(\Gamma_{1,2}\). The low-genus comparison in
\cite[Lemma~17(b)]{H17} identifies the action there with
the standard braid action of \(g\), up to an inner
correction. On the pure five-point group its class is also
that of \(g'\). Faithfulness of the GT factor gives
\begin{equation}
                         g'=g.                             \label{eq:harmonic-beta-q3}
\end{equation}
Every labelled forgetting map used below therefore carries
\(\delta_{10}\) to the standard action of \(g\), up to a
pure inner automorphism. On a mixed subgroup this follows
from the pure assertion: two automorphisms that agree on
the open pure group agree, since its centralizer is trivial.

\subsection{The harmonic four-division construction}

Fix a nonzero \(t_0\in T[2]\), and restrict \(N_0\) to its
stabilizer. Modulo its central involution this is the modular
subgroup \(R_0\simeq C_2*\mathbf Z\) of index three.
In the ten-point quotient retain \(b_0\) and the following
four branch points for a new double cover: the images of
the two nonzero two-torsion points other than \(t_0\), and
the two pairs \(\{t,-t\}\) with \(2t=t_0\).
The stabilizer of \(t_0\) preserves this set and its two
pairs, so restriction followed by this forgetting map is
defined. The two points above \(b_0\) in the new double cover
are individually marked in this order.

Here is the entire algebraic identification of the resulting
map. Write an elliptic curve with its chosen two-torsion
point as
\[
 E:\ y^2=x(x^2+Ax+B),\qquad t_0=(0,0).
\]
The other two-torsion abscissae are the roots of
\(x^2+Ax+B\). The halves of \(t_0\) have abscissae
\(\pm\sqrt B\), as follows directly from the duplication
formula. Hence the new double cover is
\begin{equation}
 E':\ z^2=(x^2+Ax+B)(x^2-B)
              =x^4+Ax^3-ABx-B^2.                           \label{eq:harmonic-beta-q4}
\end{equation}
Put
\[
 X=2z+2x^2+Ax,\qquad
 Y=2Xx-\frac{A^2}2x+\frac A2X-AB.
\]
Expansion gives
\begin{equation}
 E':\quad Y^2=X^3-B(A^2-4B)X.                               \label{eq:harmonic-beta-q5}
\end{equation}
One point above infinity becomes the origin. The other is
\begin{equation}
 P=\left(\frac{A^2}4,\frac{A(8B-A^2)}8\right).               \label{eq:harmonic-beta-q6}
\end{equation}
Thus the new curve has constant \(j=1728\).
The two infinity branches are globally ordered: they are the
sections \(z/x^2=+1\) and \(z/x^2=-1\) of the monic quartic
\eqref{eq:harmonic-beta-q4}. Under every Weierstrass coordinate change preserving
the origin and \(t_0\), the coordinates scale as
\(x\mapsto c^2x\), \(y\mapsto c^3y\), and
\(z\mapsto c^4z\). The ratio \(z/x^2\) is unchanged.
Thus the two sections cannot be interchanged by monodromy;
the first gives \(O\), the second gives \(P\), and the
construction defines a map to the group of two individually
marked points \(\Gamma_{1,2}\).

For \(E_4:y^2=x^3-x\), its quotient by \(\mu_4\) has
coordinate \(w=x^2\), with orbifold orders four at \(0\),
two at \(1\), and a puncture at infinity. With
\(t=A^2/B\), equations \eqref{eq:harmonic-beta-q5}--\eqref{eq:harmonic-beta-q6} give
\begin{equation}
                      w=\frac{t^2}{16(t-4)}.                \label{eq:harmonic-beta-q7}
\end{equation}
This is a degree-two orbifold covering: it is branched at
\(t=0\) above the order-four point and at \(t=8\) above the
order-two point. Its two unramified cusps are \(t=4,\infty\).
The source therefore has one order-two point and two cusps,
as required for \(R_0\).
It follows, without using a congruence computation to
identify a profinite map, that the construction induces a
discrete injection
\[
 \varphi_4:R_0\hookrightarrow G_4,\qquad
 G_4=q^{-1}(\langle s\rangle)\subset\Gamma_{1,2},
 \quad s=uvu,
\]
whose image has index two. Here \(q\) forgets the second
point and \(s\) has order four in \(\Gamma_{1,1}\).
The group \(G_4\) is the orbifold group of
\((E_4\setminus\{O\})/\mu_4\), hence is \(C_4*C_2\).
These inclusions retain the full profinite topology. Indeed the
point-pushing kernel \(\ker q\) has the full induced topology and
index four in \(G_4\). The quotient \(G_4/\ker q=\langle s\rangle\)
is detected injectively by signed reduction modulo four. Given a
finite-index subgroup of \(G_4\), choose an ambient open subgroup
whose intersection with \(\ker q\) is contained in it, and intersect
with this congruence kernel. The resulting intersection with
\(G_4\) lies in the prescribed subgroup. This proves the assertion
for \(G_4\); its index-two subgroup \(\varphi_4(R_0)\) inherits it.

Choose the conjugate of \(R_0\) generated by \(x=u^2\) and
the order-two image of \(s\). The cusp \(x\) has width two
and is \(t=\infty\); the other cusp \(t=4\) has width one.
The rotation character induced by \(\varphi_4\) is
\begin{equation}
                        x\longmapsto i,\qquad
                        s\longmapsto-1,                    \label{eq:harmonic-beta-q8}
\end{equation}
or its complex conjugate. This also follows from the
ramification in \eqref{eq:harmonic-beta-q7}: the order-two point lifts the square
of the order-four stabilizer, while a positive cusp loop
has primitive fourth-root rotation.

\subsection{Comparison on integral homology}

The normalized harmonic identity gives
\[
 \beta_g(u^2,v^2)^{-1}=f_g(u^2,s).
\]
Consequently the restriction of \(\Phi_g\) to \(R_0\) is
exactly the standard action in the generators \eqref{eq:harmonic-beta-q8}:
\begin{equation}
 x\longmapsto x,\qquad
 s\longmapsto f_g(x,s)^{-1}s f_g(x,s).                       \label{eq:harmonic-beta-q9}
\end{equation}
Let \(\mathcal A_g\) be the braid action on
\(\widehat\Gamma_{1,2}\), let
\(\widetilde\beta_g=\beta_g(a^2,b^2)\), and put
\(\Psi_4=\operatorname{Ad}(\widetilde\beta_g^{-1})\mathcal A_g\).
The action on the chain subgroup is \(\Phi_g\), so
\(\Psi_4\) fixes \(aba\), projects to an action fixing \(s\),
and preserves \(G_4\).

Equations \eqref{eq:harmonic-beta-q2}--\eqref{eq:harmonic-beta-q3}, restriction to the stabilizer of \(t_0\),
and the specified forgetting map give
\begin{equation}
 \Psi_4\varphi_4
       =\operatorname{Ad}(h)\varphi_4\Phi_g|_{R_0}            \label{eq:harmonic-beta-q10}
\end{equation}
for an \(h\in\widehat\Gamma_{1,2}\). More explicitly, write the
descended target action as
\(\delta_5=\operatorname{Ad}(c_3)\mathcal A_g\), with \(c_3\) pure.
Restriction and forgetting give
\[
 \delta_5\varphi_4=
 \operatorname{Ad}(\varphi_4(c_2))\varphi_4\Phi_g|_{R_0},
 \qquad
 h=\widetilde\beta_g^{-1}c_3^{-1}\varphi_4(c_2).
\]
The source inner correction \(c_2\) really belongs to the
stabilizer of \(t_0\), because it is pure modulo two; thus
its image under \(\varphi_4\) is defined. The target correction
comes from the pure five-point group, so is also defined in
\(\widehat\Gamma_{1,2}\). These account for all inner
corrections in \eqref{eq:harmonic-beta-q10}.
The rotation character \eqref{eq:harmonic-beta-q8} is unchanged by \eqref{eq:harmonic-beta-q9}.
Projecting \eqref{eq:harmonic-beta-q10} by \(q\) therefore shows that \(q(h)\)
centralizes \(\langle s\rangle\). In the completed modular
group the centralizer of its order-two factor is that factor;
lifting to \(\widehat\Gamma_{1,1}\) gives
\(q(h)\in\langle s\rangle\). Hence \(h\in G_4\).

Let \(K=\ker q=\pi_1(E_4\setminus O)\).
As in the integral cubic comparison, \(\mathcal A_g\) acts
as the identity on \(\widehat K^{\rm ab}\) for \(g\in N\).
For clarity, its action commutes with the full signed
congruence image, hence is scalar, and it fixes the
primitive class of \(ca^{-1}\). The latter assertion follows
from
\[
 \mathcal A_g(c)=P_g^{-1}cP_g,\qquad
 P_g=f_g(ba^2b,c^2)\in\widehat K,\qquad [P_g,a]=1.
\]
Thus the scalar is one. The left side of \eqref{eq:harmonic-beta-q10} consequently
acts on this homology by \(q(\beta_g)^{-1}=W_g\).

On the source, the kernel of \eqref{eq:harmonic-beta-q8} is the fundamental group
of a twice-punctured torus. After filling its two punctures,
its integral homology is a square lattice. Its map to
\(K^{\rm ab}\) is the inclusion of an index-two square
lattice: filling the two unramified cusps in \eqref{eq:harmonic-beta-q7} gives an
unramified double cover of the compact torus.
The source action on this lattice is multiplication by
\(\nu_{4,g}=B_g(i,-1)\). Here is a direct calculation.
The quotient is the Euclidean triangle group obtained from
\(R_0\) by imposing \(x^4=(xs)^4=1\). These are the two
filled positive cusp loops, so the peripheral action \eqref{eq:harmonic-beta-q9}
preserves their normal closure and descends to this quotient.
In the affine square-lattice quotient write
\[
 x=(a,i),\qquad s=(b,-1),\qquad
 [x,s]=(2a+(i-1)b,1)=:(c,1).
\]
The scalar \(c\) is nonzero and generates a finite-index
ideal in the lattice. Since
\([f_g]=R_g(X,Y)[x,s]\), the image of \(f_g(x,s)\) is
\((R_g(i,-1)c,1)\). If the induced affine automorphism
acts on translations by \(\nu\), allow a change of origin
by \(d\) in the rational quadratic algebra at each prime.
Every crossed translation term has this form, since
\(1-i\) is invertible in that algebra. Fixing \(x\) gives
\((\nu-1)a=(i-1)d\); the prescribed conjugation of \(s\)
gives \((\nu-1)b+2d=-2R_g(i,-1)c\).
Eliminating \(d\), and cancelling the nonzero lattice
factor \(c\), yields
\begin{equation}
               \nu=1-2(i-1)R_g(i,-1)=B_g(i,-1).             \label{eq:harmonic-beta-q11}
\end{equation}
Cancellation is valid in every completed integral component;
the ideal inclusion is torsion-free, including at two.

The inner automorphism by \(h\in G_4\) acts on
\(K^{\rm ab}\) by a power \(S^j\) of the order-four rotation.
The integral map between the two square lattices
intertwines their rotations, possibly with complex
conjugation. Equation \eqref{eq:harmonic-beta-q10} gives
\begin{equation}
                     W_g=S^j\nu_{4,g}^{\,\pm},              \label{eq:harmonic-beta-q12}
\end{equation}
where the superscript records the one geometric choice of embedding of
\(\widehat O_4\), fixed for all \(g\) and all coefficient primes.

Finally both nonrotation factors are one modulo four.
This is established for \(W_g\). For the beta value, the
class-two three-cycle relation gives \(R_g(1,1)=0\).
In \(\mathbf Z_2[i]\), with \(\pi=i-1\), the differences
\(i-1\) and \(-1-1\) have valuations one and two.
Thus \(R_g(i,-1)\in\pi\mathbf Z_2[i]\), and \eqref{eq:harmonic-beta-q11} is
one modulo \(\pi^4=(4)\). Since \(S\) has order four
modulo four, \eqref{eq:harmonic-beta-q12} forces \(j=0\).
This proves the asserted integral comparison. Reduction,
the cyclic order of the norm-one residue group, and the
bijectivity of squaring on its pro-\(\ell\) kernel give
the stated equivalent finite test.

This comparison does not prove that the displayed beta
value is always a square on \(D\). It identifies the
harmonic residue with an exact finite beta value; a further
relation is still needed to settle its tame square class.

\subsection{Harmonic residue characters and the exact remaining condition}

Here \(B_g\) denotes the full adelic Alexander unit evaluated by a
finite character at each coefficient prime. Put \(G=\widehat{\mathrm{GT}}\),
\(H=G_{\mathbf Q}\), \(N=\ker(\lambda,\rho_2)\), and
\(D=N_{\rm rad}\). For an odd prime \(\ell\), choose \(i\) in
\(\overline{\mathbf F}_\ell\) with \(i^2=-1\), and put
\[
 h_\ell=\frac{\ell-\chi_{-4}(\ell)}2,\qquad
 (-1)^{\vartheta_{\ell,4}(g)}
       =B_{g,\ell}(i,-1)^{h_\ell}\quad(g\in N).
\]
The value is a sign, independent of the choice of \(i\).
Theorem~\ref{thm:all-kummer-parity} gives the evenness of every
\(\rho_d^I\) on multiplier one, so the divided Kummer coordinate
in the following statement is defined on all of \(N\).

\begin{proposition}\label{prop:harmonic-finite-relative-character}
The map \(\vartheta_{\ell,4}:N\to\mathbf F_2\) is a continuous,
\(G\)-invariant homomorphism. The difference
\[
 \Xi_{\ell,4}(g)=\vartheta_{\ell,4}(g)
                    +\frac{\rho_\ell^I(g)}2\bmod2
                   \qquad(g\in N)
\]
vanishes on \(N\cap H\), and therefore extends uniquely
to a continuous character \(G\to\mathbf F_2\) that is zero
on \(H\).
On \(D\) this character is \(\vartheta_{\ell,4}\).
By Theorem~\ref{thm:harmonic-beta-square-lattice}, its vanishing
there is exactly the tame harmonic square test at \(\ell\).
No nonvanishing of \(\Xi_{\ell,4}\) on full GT is asserted.
\end{proposition}
\begin{proof}
On the multiplier-one group, the beta cocycle is
multiplicative. The unitary identity gives
\[
 B_{g,\ell}(i,-1)B_{g,\ell}(-i,-1)=1.
\]
At an inert prime, Frobenius exchanges these values, so the
first lies in the norm-one group of order \(\ell+1\).
At a split prime the pair identifies the split norm-one
torus with \(\mathbf F_\ell^\times\), of order \(\ell-1\).
Its \(h_\ell\)-th power is consequently a sign and is
unchanged by inversion. This proves the homomorphism and
orientation assertions.
The semilinear beta law gives, for \(a\in G\) and \(g\in N\),
\[
 B_{aga^{-1}}(X,Y)=B_g(X^{\lambda_a},Y^{\lambda_a}).
\]
At \((i,-1)\), the only possible change is \(i\mapsto-i\);
the sign is unchanged. Thus \(\vartheta_{\ell,4}\) is invariant.
The divided Kummer character is invariant for the same
reason: conjugation scales \(\rho_\ell^I|_N\) by the odd
unit \(\lambda_a\).

We check the arithmetic restriction with the full Anderson
formula, not only its logarithmic derivative. In the beta
form of \cite[Remark~2.4, p.~101]{Furusho2003}, put
\(n=\ell\) and \(s=1/4\), and reduce in characteristic
\(\ell\). Every \(\ell\)-th root in the denominator reduces
to one; the denominator is therefore one.
Normalized duplication gives
\[
 B_\sigma(i,i)=B_\sigma(-1,i)=B_\sigma(i,-1)
                     \quad(\sigma\in N\cap H).
\]
In the numerator, the residues \(k=1,2\bmod4\) contribute
this common value and the residues \(k=0,3\bmod4\)
contribute one. Their total number for
\(0\le k\le\ell-1\) is exactly \(h_\ell\):
it is \(2r\) if \(\ell=4r+1\), and \(2r+2\) if
\(\ell=4r+3\). The arithmetic right side is
\[
 i^{-\ell\rho_\ell^I(\sigma)}
       =(-1)^{\rho_\ell^I(\sigma)/2},
\]
since the Kummer coordinate is even on the cyclotomic
kernel. This proves \(\Xi_{\ell,4}|_{N\cap H}=0\).

Finally \(G=NH\). The invariant character extends by
\(\Xi_{\ell,4}(nh)=\Xi_{\ell,4}(n)\). Vanishing on the intersection
proves that this is well-defined, invariance proves
multiplicativity, and compactness of \(N\times H\to G\)
proves continuity. Its uniqueness follows from \(G=NH\).
\end{proof}

There is a finite-cover interpretation of the beta
coordinate used here. The cyclic degree-four cover
\[
                 v^4=t(t-1)^2
\]
has monodromy \(x\mapsto1,\ y\mapsto2\) in
\(\mathbf Z/4\). Its compactification is the elliptic curve
\[
             v^2=z^3-z,\qquad t=z^2,\quad
             z=\frac{v^2}{t-1}.
\]
The \((i,-1)\) eigenspace in its homology modulo \(\ell\)
is one-dimensional, and the restricted GT action
on that eigenspace is \(B_{g,\ell}(i,-1)\).
Indeed the three nontrivial deck characters occur once in
the homology of the punctured cover; the order-two character
is a boundary character, and the two primitive characters
form the closed elliptic homology. The Alexander generator
maps nontrivially to each primitive eigenspace, so its
coefficient is exactly the stated beta value.
This interpretation uses restriction of an automorphism
that preserves the displayed subgroup. It makes no
assertion that every possible scalar lifts to a GT element.

\begin{proposition}\label{prop:harmonic-formal-residue-gap}
Vanishing of the integral rational Kummer coordinates,
beta associativity, unitarity, the dyadic Anderson formulas,
and logarithmic Anderson identities do not by themselves
force every harmonic square test as identities of formal
integral Alexander data.
\end{proposition}
\begin{proof}
Fix an odd prime \(\ell\ge5\). The established simultaneous
arithmetic radical image supplies
\(\sigma\in H^{\rm cyc}\) with
\[
 \rho_\ell^I(\sigma)=2e_2,\qquad
 \rho_q^I(\sigma)=0\quad(q\ne\ell),
\]
where \(e_2\in\widehat{\mathbf Z}\) is the idempotent
supported at the coefficient prime two. In particular
\(\sigma\in N\) and every coefficient-\(\ell\) Kummer
coordinate is zero. The preceding arithmetic calculation
gives
\[
              B_{\sigma,\ell}(i,-1)^{h_\ell}=-1.
\]
Keep only the coefficient-\(\ell\) component of
\(R_\sigma\) and set its other coefficient-prime components
to zero. The resulting beta unit is \(B_{\sigma,\ell}\)
in this component and one in all others.
All integral cyclic-cover coefficient maps are
\(\widehat{\mathbf Z}\)-linear on the Alexander module.
Thus every patched Kummer coordinate is zero.
Associativity, unitarity and normalized dyadic identities
hold componentwise. For an Anderson discrepancy near a
prime-to-\(\ell\) torsion point, its arithmetic factor is
a constant root of unity times the principal factor
\((1+S)^{n\rho_{n,\ell}^I(\sigma)}\).
The latter is one, so its logarithmic derivative vanishes.
At the other coefficient primes the beta unit is one.
The patched data therefore satisfy all the listed formal
conditions while retaining the nontrivial harmonic sign.

This is formal Alexander data. No lift satisfying the full
profinite pentagon, or the full nonabelian cyclic-filling
identities, is constructed. The construction isolates a
missing implication and is not a counterexample in GT.
\end{proof}

\subsection{Finite-cover coordinates for the cubic restriction}

In this subsection \(B_g\) is the full adelic normalized Alexander
unit, evaluated through a finite character at each coefficient prime.
In particular its evaluation at a cube root in characteristic
\(\ell\ne3\) is not obtained by substituting into a dyadic formal
power series. Put \(G=\widehat{\mathrm{GT}}\),
\(H=G_{\mathbf Q}\), \(N=\ker(\lambda,\rho_2)\), and
\(D=N_{\rm rad}\).
For every prime \(\ell\ge5\), choose a primitive cube root
\(\omega\in\overline{\mathbf F}_\ell\), and write
\[
 n_\ell=\frac{\ell-\chi_{-3}(\ell)}3,\qquad
 \omega^{c_{\ell,3}(g)}
       =B_{g,\ell}(\omega,\omega)^{n_\ell}\quad(g\in N).
\]
The exponent \(c_{\ell,3}(g)\) belongs to \(\mathbf F_3\).
The notation \(\mathbf F_3(1)\) below means the additive group
\(\mathbf F_3\) with \(G\)-action through \(\lambda\bmod3\).

\begin{proposition}\label{prop:cubic-finite-relative-cocycle}
The map \(c_{\ell,3}:N\to\mathbf F_3(1)\) is a continuous
homomorphism, independent of the choice of \(\omega\), and satisfies
\[
 c_{\ell,3}(aga^{-1})=\lambda_a c_{\ell,3}(g).
\]
The homomorphism
\[
 \Xi_{\ell,3}(g)=c_{\ell,3}(g)+
                  \ell\,\rho_\ell^I(g)\bmod3
                       \qquad(g\in N)
\]
vanishes on \(N\cap H\). It extends uniquely to a continuous
crossed homomorphism \(G\to\mathbf F_3(1)\) that is zero on \(H\).
On \(D\) its value is \(c_{\ell,3}\).
\end{proposition}
\begin{proof}
The multiplier-one beta law is multiplicative. Its unitary identity
gives
\[
 B_{g,\ell}(\omega,\omega)
 B_{g,\ell}(\omega^{-1},\omega^{-1})=1.
\]
The first value lies in the residue norm-one torus, of order
\(\ell-\chi_{-3}(\ell)=3n_\ell\). Its \(n_\ell\)-th power
therefore belongs to \(\mu_3\). Replacing \(\omega\) by
\(\omega^{-1}\) inverts both the left side of the defining equation
and its chosen generator, leaving the exponent \(c_{\ell,3}\)
unchanged. This proves well-definedness, continuity and additivity.
The semilinear beta law on \(N\) gives
\[
 B_{aga^{-1}}(X,Y)
       =B_g(X^{\lambda_a},Y^{\lambda_a}).
\]
At a cube root this substitutes either \(\omega\) or
\(\omega^{-1}\), and proves the displayed equivariance.
The Kummer coordinate has the same equivariance.

For \(\sigma\in N\cap H\), use the full arithmetic Anderson
formula in the beta form of
\cite[Remark~2.4, p.~101]{Furusho2003}, with \(s=1/3\) and
the multiplication index equal to \(\ell\). Reduce at the
coefficient prime \(\ell\). The denominator is one, since every
\(\ell\)-th root reduces to one and \(B_\sigma(1,\omega)=1\).
In the numerator the factors with \(k=0,2\bmod3\) are one:
these are the boundary values
\(B_\sigma(\omega,1)\) and \(B_\sigma(\omega,\omega^{-1})\).
The factors with \(k=1\bmod3\) equal
\(B_\sigma(\omega,\omega)\). There are \(n_\ell\) such factors
among \(0\le k<\ell\): the count is \(r\) when
\(\ell=3r+1\), and \(r+1\) when \(\ell=3r+2\).
The right side is \(\omega^{-\ell\rho_\ell^I(\sigma)}\).
Consequently
\[
 c_{\ell,3}(\sigma)=-\ell\rho_\ell^I(\sigma)\pmod3,
\]
as required. This step uses the full formula, including its
constant root of unity; its logarithmic derivative is insufficient.

Use \(G=NH\) and define
\(\Xi_{\ell,3}(nh)=\Xi_{\ell,3}(n)\).
Vanishing on \(N\cap H\) makes this well-defined.
For \(n_i\in N,h_i\in H\), equivariance gives
\[
 \Xi_{\ell,3}(n_1h_1n_2h_2)
   =\Xi_{\ell,3}(n_1)+\lambda_{h_1}\Xi_{\ell,3}(n_2),
\]
which is exactly the crossed law, since \(\lambda_{n_1}=1\).
The continuous surjection \(N\times H\to G\) from a compact
space proves continuity of the descended map. The factorization
\(G=NH\) also proves uniqueness.
\end{proof}

The beta scalar here is detected by a finite cover.
The degree-three cover
\[
                 v^3=t(t-1)
\]
has monodromy \(x\mapsto1,\ y\mapsto1\) in
\(\mathbf Z/3\). Its compactification has equation
\[
                Y^2=4v^3+1,\qquad Y=2t-1,
\]
and genus one: it is fully ramified at the three deleted points,
so Riemann--Hurwitz gives \(2g-2=-6+3(3-1)=0\).
The two nontrivial deck-character eigenspaces in its homology
modulo \(\ell\) have dimension one. The punctured cover has
rank-four first homology; its two-dimensional boundary subspace
is deck-trivial because each of its three punctures is fixed.
Thus the two primitive eigenspaces survive compactification.
For \(g\in N\), the standard automorphism of the punctured line
preserves this cyclic subgroup and the boundary classes. On the
Alexander module its commutator generator is multiplied by
\(B_g(X,Y)\). Specializing \(X=Y=\omega\) therefore gives precisely
the scalar \(B_{g,\ell}(\omega,\omega)\) on the corresponding
elliptic eigenspace. This argument only restricts an existing
GT automorphism. It does not assert that every norm-one scalar
lifts to a GT automorphism.

Finally apply the established integral cubic-frame comparison
and the already eliminated local conditions at two and three.
On \(D\), the kernel of the cubic radical--CM defect has the
explicit description
\[
 \ker(d_{(3)}|_D)
       =\bigcap_{\ell\ge5}\ker(c_{\ell,3}|_D).
\]
Indeed the kernel of the \(n_\ell\)-th power on the cyclic residue
torus is its cube subgroup, and cubing is an automorphism on the
pro-\(\ell\) reduction kernel for \(\ell\ne3\).
Consequently, for the established normalization \(G=R_DH\),
\[
 R_D\cap\ker d_{(3)}
       =R_D\cap\bigcap_{\ell\ge5}\ker\Xi_{\ell,3}.
\]
This makes the permitted cubic restriction in the maximal
pro-two relative quotient a concrete intersection of finite-cover
kernels. It gives no nonzero value of a relative cocycle and does
not assert that the intersection is a proper subgroup of \(R_D\).

\section{Primary residual character groups}
\begin{lemma}[Coprime relative extension without an intersection assumption]
\label{lem:coprime-relative-general}
Let \(G=RH\), with \(R\) closed normal and \(H\) closed.
Let \(E\triangleleft G\) be closed and contained in \(R\), and
suppose every finite quotient of \(R/E\) has order prime to \(r\).
For each finite abelian \(r\)-group \(A\), restriction gives a
bijection between continuous characters \(G\to A\) killing \(H\)
and continuous \(G\)-invariant characters \(E\to A\) killing
\(E\cap H\).
\end{lemma}

\begin{proof}
Let \(\chi:E\to A\) be such a character.
The quotient \(L=R/\ker\chi\) has the central finite
\(r\)-subgroup \(Z=E/\ker\chi\cong\chi(E)\), and \(L/Z\)
is pro-\(r'\).
It has a unique pro-\(r'\) complement to \(Z\).

Here is a direct proof of this splitting assertion.
In a finite central extension of a group \(Q\), of order
\(m\) prime to \(r\), by \(Z\), choose a section with normalized
two-cocycle \(c\). Define
\[
 b(q)=m^{-1}\sum_{t\in Q}c(q,t).
\]
The inverse of \(m\) acts on the finite \(r\)-group \(Z\).
Summing the cocycle identity over its last variable gives
\(c(q,t)=b(q)+b(t)-b(qt)\).
Correcting the section by \(-b\) therefore gives a homomorphic
section. Two such sections differ by a homomorphism \(Q\to Z\),
which is zero. The complement is unique.
For \(L\), pass through the cofinal open normal subgroups
\(U\) with \(U\cap Z=1\).
The unique complements in \(L/U\) are compatible and map onto
each other. Their inverse limit gives \(L=Z\times C\).

The \(G\)-action fixes \(Z\) pointwise and preserves its unique
complement. Projection therefore gives a continuous
\(G\)-invariant extension \(\theta:R\to A\).
It is unique, because \(R/E\) has no nonzero character in \(A\).
The restriction of \(\theta\) to \(R\cap H\) factors through
\((R\cap H)/(E\cap H)\), a closed subgroup of \(R/E\), and
hence vanishes.
Now \(\Phi(rh)=\theta(r)\) is well-defined, is a homomorphism
by normality and invariance, and is continuous by compactness.
It kills \(H\), and uniqueness on \(R\) proves uniqueness on
\(G\). The reverse restriction is immediate.
\end{proof}

\begin{theorem}[A prime-specific residual group]
\label{thm:primary-residual-gate}
For a prime \(r\), put
\[
 E_r=R_D\cap\bigcap_{p\ne r}\ker\Phi_p,
 \qquad X_r=H^1(N,\mathbf F_r)^G,
\]
where \(\Phi_p\) is the action on the free pro-\(p\) group.
Then \(E_r\) is closed and normal in \(G\), and
\begin{equation}
 B_{\max r}\cong
 \left(E_r/\overline{[G,E_r](E_r\cap H)}\right)_{\max r}. \label{eq:primary-residual-1}
\end{equation}
For odd \(r\), restriction
\[
 X_r\longrightarrow H^1(E_r,\mathbf F_r)
\]
is injective.
For \(r=2\), its kernel is
\begin{equation}
 \bigoplus_{q\ {\rm odd}}\mathbf F_2\psi_q
                         \oplus\mathbf F_2\eta_3.       \label{eq:primary-residual-2}
\end{equation}
Moreover \(E_r\) has infinite index in \(R_D\).
\end{theorem}

\begin{proof}
The normalized image \(\Phi_p(M)\) is pro-\(p\).
On each finite algebra
\(\mathbf F_p\langle S,T\rangle/(S,T)^d\), its action is
upper unitriangular because it fixes the associated graded.
These finite \(p\)-group actions separate normalized parameters.
Consequently
\[
 R_D/E_r\lhook\joinrel\longrightarrow
                    \prod_{p\ne r}\Phi_p(R_D)
\]
is pro-\(r'\).
Indeed every finite quotient of the closed image factors
through a finite collection of finite \(p\)-group coordinates,
with all \(p\ne r\).
Normality and closedness of \(E_r\) follow from its definition.

Apply Lemma~\ref{lem:coprime-relative-general} with \(R=R_D\),
using the established factorization \(G=R_DH\).
It identifies all finite abelian \(r\)-valued characters of
the two sides of \eqref{eq:primary-residual-1}.
More directly, the natural map from the right side into
\(B_{\max r}\) is injective because these finite characters
separate points and all extend uniquely.
Its cokernel is a pro-\(r\) quotient of \(R_D/E_r\), so is
trivial. Compactness proves \eqref{eq:primary-residual-1}.

A character of \(R_D\) with values in \(\mathbf F_r\) and
trivial on \(E_r\) factors through the pro-\(r'\) quotient
and is zero. Thus
\begin{equation}
 \ker\bigl(X_r\to H^1(E_r,\mathbf F_r)\bigr)
 =
 \ker\bigl(X_r\to H^1(R_D,\mathbf F_r)\bigr).             \label{eq:primary-residual-3}
\end{equation}
For odd \(r\), a character in the latter kernel first vanishes
on \(D\), since \(D/R_D\) is pro-two.
It factors through the rational Kummer quotient \(N/D\).
Complex conjugation acts on this quotient by inversion.
Invariance therefore gives \(\chi=-\chi\), and hence
\(\chi=0\) in odd characteristic.
For \(r=2\), \eqref{eq:primary-residual-3} gives \eqref{eq:primary-residual-2} from the established original
restriction-kernel theorem.

It remains to prove infinite index.
Choose an odd prime \(p\ne r\).
Only one nonzero arithmetic degree-three class at \(p\) is
needed. Its existence follows from
\cite[Remark~8.6, p.~151]{HainMatsumoto2003}; the outer-action
filtration is defined there in \S1, equation~\eqref{eq:primary-residual-2}, p.~120.
Choose an arithmetic representative \(t\in I_p^3H\) whose
degree-three outer class \(v\) is nonzero over \(\mathbf Q_p\).
In particular \(\lambda_p(t)=1\).

Choose \(h\in H\) with \(u=\lambda_p(h)=1+p\), let
\(b=[h,t]\), and let \(\varepsilon_p\in\widehat{\mathbf Z}\)
have \(p\)-component one and all other components zero.
Set
\begin{equation}
 z=b^{\varepsilon_p},\qquad
 w_p=hzh^{-1}z^{-\lambda(h)}.                            \label{eq:primary-residual-4}
\end{equation}
The commutator \(b\) has full multiplier one.
Thus, for every \(\ell\ne p\), its normalized pro-\(\ell\)
image is killed by the power \(\varepsilon_p\).
Both \(z\) and \(w_p\) have trivial \(\ell\)-action for
\(\ell\ne p\).
All rational Kummer coordinates cancel exactly in \eqref{eq:primary-residual-4},
because their conjugation weight on \(M\) is the full
multiplier \(\lambda(h)\). Therefore \(w_p\in E_p\cap H\).

Conjugation by multiplier \(u\) acts on the degree-three
outer class by \(u^3\); this is also immediate from the
linear scaling of a leading derivation.
The class of \eqref{eq:primary-residual-4} is consequently
\begin{equation}
 \operatorname{gr}_{3,p}(w_p)
      =(u^3-u)(u^3-1)v\ne0.                             \label{eq:primary-residual-5}
\end{equation}
Both factors are nonzero, also when \(p=3\).
Working with the outer class suffices to detect a
nontrivial action and avoids any additional signed
parameter comparison.
Its procyclic image has an infinite \(p\)-adic graded
coordinate.
Since \(E_r\subset\ker\Phi_p\), this gives an infinite
image in \(R_D/E_r\), proving infinite index.
\end{proof}

\begin{remark}
The theorem does not assert \(G=E_rH\).
For odd \(r\), it proves injectivity of the indicated
restriction map; it does not compute \(X_r\).
The strictness argument uses one degree-three arithmetic
class and requires no full leading-lattice result at an
odd prime.
\end{remark}

\subsection{Removing an odd relative defect for quadratic questions}

The distinction between ordinary relative characters and the cubic
crossed defect allows the following exact reduction.

\begin{lemma}\label{lem:odd-relative-kernel}
Let \(G=DH\) be a profinite group with \(D\) closed and normal,
and put \(K=D\cap H\). Suppose \(E\) is a closed normal subgroup
of \(G\), with \(K\subset E\subset D\), such that \(D/E\) is
pro-\(3\). For every finite abelian \(2\)-group \(A\), restriction
is a canonical bijection
\[
 \begin{aligned}
 \{\chi\in\operatorname{Hom}_{\rm cts}(G,A):\chi(H)=0\}
 &\longrightarrow\\
 &\{\psi\in\operatorname{Hom}_{\rm cts}(E,A)^G:\psi(K)=0\}.
 \end{aligned}
\]
Consequently the maximal pro-\(2\) quotient of the relative
abelian group can be computed as the maximal pro-\(2\) quotient of
\[
                  E/\overline{[G,E]K}.
\]
This does not assert \(G=EH\) or equality of the full relative
abelian quotients before taking their pro-\(2\) parts.
\end{lemma}
\begin{proof}
Let \(\psi\) be a character on the right and put \(L=\ker\psi\).
It is closed and normal in \(G\). In \(D/L\), the subgroup
\(E/L\) is a finite central \(2\)-group, and its quotient is
pro-\(3\). This group has a unique pro-\(3\) complement.
To see this directly, pass to finite quotients. A finite
central extension of a \(3\)-group by a \(2\)-group is the
direct product of its central \(2\)-subgroup with its Sylow
\(3\)-subgroup: a Sylow \(3\)-subgroup has the full required
order, and all its conjugates are itself because the
central subgroup and that Sylow subgroup generate the group.
The unique Sylow \(3\)-subgroups are compatible under quotient
maps, so their inverse limit is the asserted unique complement.
Hence
\[
                    D/L=(E/L)\times P
\]
with \(P\) pro-\(3\). Projection to \(E/L\), followed by
\(\psi\), extends \(\psi\) uniquely to a continuous
homomorphism \(\widetilde\psi:D\to A\).
Uniqueness follows also because a pro-\(3\) group has no
nontrivial continuous homomorphism to \(A\).

For every \(g\in G\), the conjugate of
\(\widetilde\psi\) is another extension of \(\psi\);
thus it is the same extension. The character on \(D\)
is therefore \(G\)-invariant and zero on \(K\).
It extends to \(G=DH\) by
\(\chi(dh)=\widetilde\psi(d)\).
The intersection and invariance conditions give
well-definedness and multiplicativity, and the compact
surjection \(D\times H\to G\) gives continuity.
Conversely every character on the left restricts as
asserted, and uniqueness on \(D\), followed by \(G=DH\),
proves injectivity. Taking all finite abelian \(2\)-groups
gives the pro-\(2\) assertion.
\end{proof}

Normality of $D$ is necessary. Let $G=C_2\times\mathfrak S_4$,
$D=C_2\times\langle(123)\rangle$, and $E=C_2\times\{1\}$.
In the dihedral subgroup $J=\langle r=(1234),s=(13)\rangle$,
let $\theta:J\to C_2$ have $\theta(r)=0$, $\theta(s)=1$,
and let $H$ be its graph in $G$. Then $G=DH$, $D\cap H=1$,
and $D/E=C_3$, while $E$ is central. The nonzero character of
$E$ cannot extend to a character of $G$ killing $H$:
every character of $G$ is a combination of the first projection
and the sign character, and $\theta$ is independent of the
restriction of the sign to $J$. Thus dropping normality would
invalidate the displayed bijection.

Apply this lemma to the radical--CM defect of
Section~\ref{sec:lr-relative-gate}
\[
 d:G\longrightarrow C,\qquad d|_H=0,\qquad
 d(gh)=d(g)+\lambda_gd(h).
\]
Let \(d_{(3)}\) be its order-three component and put
\[
                  D^{\rm cub}=\ker(d_{(3)}|_D).
\]
On \(D\), the crossed law is an ordinary homomorphism.
Its \(G\)-equivariance makes \(D^{\rm cub}\) normal in \(G\);
it contains \(K=D\cap H\), and \(D/D^{\rm cub}\) has
exponent three. Thus the ordinary relative quadratic
character problem, and more generally the maximal
pro-\(2\) relative abelian quotient, may be computed on
\(D^{\rm cub}\) without proving \(d_{(3)}=0\) on \(G\).
On this smaller domain the cubic radical--CM compatibility
is imposed by definition. The tame harmonic square tests
remain arithmetic-trivial quadratic characters there.
If the cubic defect is nonzero, this is a proper reduction;
if it vanishes, \(D^{\rm cub}=D\).

More precisely, use the established closed normal subgroup
\(R_D\subset D\) with \(G=R_DH\), and put
\[
 S=R_D\cap\ker d_{(3)},\qquad K_R=R_D\cap H.
\]
Then \(K_R\subset S\), the subgroup \(S\) is closed and normal
in \(G\), and \(R_D/S\) has exponent three. The preceding
lemma, applied with \(R_D\) in place of \(D\), gives the
canonical identification
\[
 B_{\mathrm{max}\,2}\ \simeq\
 \left(S/\overline{[G,S]K_R}\right)_{\mathrm{max}\,2},
\]
where the subscript denotes the maximal pro-\(2\) quotient.
For the original invariant-character group
\(X_2(N)=\operatorname{Hom}_{\rm cts}(N,\mathbf F_2)^G\),
restriction to \(S\) has exactly the same kernel as
restriction to \(R_D\). Indeed a character of \(N\) that is
zero on \(S\), when restricted to \(R_D\), factors through
the pro-\(3\) quotient \(R_D/S\) and is therefore zero.
The converse is immediate. Thus this extra cubic-kernel
restriction loses none of the original quadratic characters
that survived the \(R_D\) normalization.
It is a proper restriction precisely when
\(d_{(3)}|_{R_D}\ne0\); no such nonvanishing is asserted.

\subsection{Removing all single-prime word images from the quadratic question}

Here a pro-odd group means a profinite group all of whose finite
continuous quotients have odd order. This term does not require
commutativity.

\begin{lemma}\label{lem:relative-proodd-restriction}
Let \(G=QH\), where \(Q\) is a closed normal subgroup of the
profinite group \(G\) and \(H\) is closed. Let \(E\subset Q\)
be closed and normal in \(G\), and suppose that \(Q/E\) is pro-odd.
For every finite abelian \(2\)-group \(A\), restriction induces
a bijection
\[
 \begin{aligned}
 \{\chi\in\operatorname{Hom}_{\rm cts}(G,A):\chi|_H=0\}
 &\longrightarrow\\
 &\{\psi\in\operatorname{Hom}_{\rm cts}(E,A)^G:
                  \psi|_{E\cap H}=0\}.
 \end{aligned}
\]
All maps on the left are continuous homomorphisms. In particular
there is a canonical isomorphism
\[
 \left(Q/\overline{[G,Q](Q\cap H)}\right)_{\max2}
 \ \simeq\
 \left(E/\overline{[G,E](E\cap H)}\right)_{\max2}.
\]
No containment \(Q\cap H\subset E\) is required.
\end{lemma}
\begin{proof}
Given \(\psi\) on the right, put \(L=\ker\psi\).
Then \(L\) is closed and normal in \(G\), and \(E/L\) is a
finite central \(2\)-subgroup of \(Q/L\), whose quotient is pro-odd.
We first record why a central finite \(2\)-subgroup \(C\) of a
finite group \(J\), with \(J/C\) of odd order, has a unique normal
complement. The transfer \(V:J\to C\) is a homomorphism and
satisfies \(V(c)=c^{[J:C]}\) for \(c\in C\). These assertions
follow directly by multiplying the coset-representative factors;
centrality permits their reordering. Raising \(V\) to an exponent
inverse to \([J:C]\) modulo the exponent of \(C\) gives a
retraction \(J\to C\). Its kernel is an odd normal complement.
Thus \(J=C\times P\). Any other complement is the graph of a
homomorphism from the odd group \(P\) to \(C\), which must be
zero. The complement is unique.

Apply this observation to every finite quotient of \(Q/L\).
The unique odd complements are compatible with all quotient maps.
Taking inverse limits gives the canonical direct product
\[
                    Q/L=(E/L)\times P
\]
with \(P\) pro-odd. Projection to \(E/L\), followed by \(\psi\),
extends \(\psi\) to a continuous homomorphism \(\widetilde\psi:Q\to A\).
Any two extensions differ by a homomorphism \(Q/E\to A\),
which is zero. Uniqueness implies that \(\widetilde\psi\) is
\(G\)-invariant.

The restriction of \(\widetilde\psi\) to \(Q\cap H\) vanishes on
\(E\cap H\). The quotient
\[
 (Q\cap H)/(E\cap H)\ \hookrightarrow\ Q/E
\]
is pro-odd, so this restriction is zero. Hence
\(\chi(qh)=\widetilde\psi(q)\) is well-defined.
Invariance proves that it is a homomorphism, and the compact
surjection \(Q\times H\to G\) proves continuity.
The same pro-odd argument gives uniqueness and injectivity.
Applying the bijection to all finite abelian \(2\)-groups proves
the asserted canonical isomorphism of maximal pro-two quotients.
\end{proof}

Return to \(G=\widehat{\mathrm{GT}}\), \(H=G_{\mathbf Q}\),
\(D=N_{\rm rad}\), and the established normalization
\(G=R_DH\), where \(R_D\) is the kernel on \(D\) of the
free pro-two word action. For every prime \(p\), let
\[
 \Phi_p:G\longrightarrow
       \operatorname{Aut}(\widehat F_2^{(p)})
\]
be the standard action on the free pro-\(p\) quotient, and put
\[
 P_D=D\cap\bigcap_p\ker\Phi_p,\qquad
 T_D=P_D\cap\ker d_{(3)}.
\]
On multiplier one, \(\Phi_p(g)=1\) is equivalent to \(f_g=1\)
in the free pro-\(p\) group. Indeed the normalized action fixes
\(x\), and triviality on \(y\) says that \(f_g\) centralizes
the free generator \(y\). Its procyclic centralizer intersects
the commutator subgroup trivially.

\begin{theorem}\label{thm:all-prime-residual-quadratic-domain}
The subgroups \(P_D,T_D\) are closed and normal in \(G\).
Both \(R_D/P_D\) and \(R_D/T_D\) are pro-odd. Consequently
\[
 \begin{aligned}
 B_{\max2}&\simeq
 \left(P_D/\overline{[G,P_D](P_D\cap H)}\right)_{\max2}\\
 &\simeq
 \left(T_D/\overline{[G,T_D](T_D\cap H)}\right)_{\max2}.
 \end{aligned}
\]
Restriction from \(X_2(N)\) to either \(P_D\) or \(T_D\)
has the same kernel as restriction to \(R_D\), namely
\[
 \left(\bigoplus_{q\ {\rm odd\ prime}}\mathbf F_2\psi_q\right)
                         \oplus\mathbf F_2\eta_3.
\]
These assertions concern the maximal pro-two relative quotient;
they do not assert \(G=P_DH\) or \(G=T_DH\).
\end{theorem}
\begin{proof}
For a finitely generated free pro-\(p\) group, the kernel of
its automorphism action on its Frattini quotient is pro-\(p\).
For completeness, use the lower exponent-\(p\) central series
\(V_1=F\), \(V_{n+1}=V_n^p[V_n,F]\).
The \(V_n\) are characteristic open subgroups forming a
neighborhood basis. At each successive finite quotient,
automorphisms that are the identity one step earlier are
determined by central corrections to the free generators in
\(V_n/V_{n+1}\); their group is a subgroup of a finite
elementary abelian \(p\)-group. Induction, beginning with the
Frattini quotient, and inverse limits prove the assertion.

Every element of \(D\) has multiplier one. Its pro-\(p\)
action is therefore in this pro-\(p\) automorphism kernel.
Since the pro-two action is already trivial on \(R_D\),
the quotient \(R_D/P_D\) embeds as a closed subgroup in a
product of pro-\(p\) groups over odd primes \(p\).
Every finite quotient of that product, and of its closed
subgroups, has odd order. Thus \(R_D/P_D\) is pro-odd.
The quotient \(R_D/(R_D\cap\ker d_{(3)})\) has exponent three.
The diagonal map embeds \(R_D/T_D\) in the product of these
two pro-odd quotients, proving the second assertion.
Normality follows from the characteristic nature of the
free pro-\(p\) quotients and the equivariance of \(d_{(3)}\).

Apply Lemma~\ref{lem:relative-proodd-restriction} with \(Q=R_D\).
For the original character kernel, a character on \(N\) that
vanishes on \(P_D\), or on \(T_D\), restricts to a homomorphism
from the corresponding pro-odd quotient of \(R_D\) to
\(\mathbf F_2\). It is therefore zero on \(R_D\).
The reverse implication is immediate. The known kernel for
restriction to \(R_D\) completes the proof.
\end{proof}

\begin{proposition}\label{prop:all-prime-residual-strictness}
Both \(P_D\) and \(T_D\) have infinite index in \(R_D\).
Moreover \(T_D\) has infinite index in \(R_D\cap\ker d_{(3)}\).
\end{proposition}
\begin{proof}
Fix any odd prime \(p\). We use only the existence of a
nonzero arithmetic degree-three class, not arithmetic
saturation at \(p\): Ihara's nonzero Soul\'e class is recalled
in \cite[Remark~8.6, p.~151]{HainMatsumoto2003}.
Choose an arithmetic element \(\sigma\) with \(\lambda_p(\sigma)=1\)
and nonzero degree-three pro-\(p\) leading term \(q\).
Here the passage from the outer-action statement to the
normalized word has no kernel in this degree. The normalized
word has zero degrees one and two when \(\lambda_p=1\).
Its degree-three term \(q\) induces the derivation
\(S\mapsto0,\ T\mapsto[T,q]\), up to the harmless convention
of reversing this bracket. An inner derivation with this
property would be given by a degree-three Lie element
commuting with \(S\); the centralizer of \(S\) in the free
Lie algebra has no degree-three part. Thus a nonzero outer
degree-three class has \(q\ne0\). Multiplying an integral
representative if necessary realizes the nonzero rational
class by a Galois element.

Choose \(h\in H\) with \(u=\lambda_p(h)=1+p\), and use the
commutator convention \([a,b]=aba^{-1}b^{-1}\). Put
\[
 z=[h,\sigma],\qquad
 d=hzh^{-1}z^{-\lambda(h)}.
\]
The exponent in the second expression is the full adelic
integer \(-\lambda(h)\): the continuous power map
\(\widehat{\mathbf Z}\to\overline{\langle z\rangle}\)
defines it. The multiplier of \(z\) is one at every prime.
The degree-three word coordinate is additive on multiplier
one, and conjugation has weight \(\lambda_p^3\). Hence
\[
 \operatorname{in}_3(z)=(u^3-1)q,\qquad
 \operatorname{in}_3(d)=(u^3-u)(u^3-1)q\ne0
                  \quad\text{over }\mathbf Z_p.
\]
Each rational Kummer coordinate is additive on multiplier
one and has conjugation weight \(\lambda(h)\).
It therefore vanishes on \(d\), so \(d\in D\cap H\).

Let \(\epsilon_p\in\widehat{\mathbf Z}=\prod_r\mathbf Z_r\)
have component one at \(p\) and zero at every other prime.
The element \(w=d^{\epsilon_p}\) is still in \(D\cap H\).
The pro-two image of \(d\) lies in a pro-two automorphism
group, so the zero two-component of \(\epsilon_p\) kills
that image. Thus \(w\in R_D\cap H\).
Its degree-three pro-\(p\) term equals that of \(d\), and is
nonzero. Since this term is an additive coordinate in a
torsion-free \(\mathbf Z_p\)-module, the pro-\(p\) image
of the closed cyclic subgroup generated by \(w\) is infinite.
The subgroup \(P_D\) kills this image, proving infinite index.
Finally \(d_{(3)}\) is zero on \(H\), so \(w\) already lies in
\(R_D\cap\ker d_{(3)}\); no additional cubing is needed.
The same infinite pro-\(p\) image proves the assertion for
\(T_D\).
\end{proof}

\begin{theorem}[A strict gate from the odd harmonic CM image]
\label{thm:odd-harmonic-cm-gate}
Put
\[
 P_D=D\cap\bigcap_p\ker\Phi_p,\qquad
 T_D=P_D\cap\ker d_{(3)},
\]
where \(d_{(3)}\) is the cubic radical--CM defect.
Let \(W_{\rm odd}\) be the projection of the harmonic CM
homomorphism \(W:D\to N_4(4)\) onto the maximal pro-odd factor,
and put
\[
 U_D=T_D\cap\ker W_{\rm odd}.
\]
Then \(U_D\) is closed and normal in \(G\), and
\begin{equation}
 B_{\max2}\cong
 \left(U_D/\overline{[G,U_D](U_D\cap H)}\right)_{\max2}.  \label{eq:odd-cm-1}
\end{equation}
The kernel of the original restriction map is
\begin{equation}
 \ker\bigl(X_2\to H^1(U_D,\mathbf F_2)\bigr)
  =\bigoplus_{q\ {\rm odd}}\mathbf F_2\psi_q
                           \oplus\mathbf F_2\eta_3.     \label{eq:odd-cm-2}
\end{equation}
Moreover
\begin{equation}
 [T_D:U_D]=\infty,\qquad
 [T_D\cap H:U_D\cap H]=\infty.                          \label{eq:odd-cm-3}
\end{equation}
\end{theorem}

\begin{proof}
The pro-odd projection is characteristic, and the conjugation
action on \(W\) is identity or inversion. Thus \(U_D\) is
closed and \(G\)-normal. Its quotient in \(T_D\) is the
pro-odd group \(W_{\rm odd}(T_D)\).

A \(G\)-invariant character of \(U_D\) with values in a finite
abelian two-group extends uniquely to a \(G\)-invariant
character of \(T_D\).
Indeed its pushout is a central extension of a pro-odd group
by a finite two-group. In a finite quotient, if \(c\) is the
central two-cocycle and the odd quotient has order \(m\), the
cochain
\[
 b(q)=m^{-1}\sum_t c(q,t)
\]
satisfies \(c(q,t)=b(q)+b(t)-b(qt)\).
The resulting complement is unique, since an odd-order group
has no nonzero character in a two-group. Passing to the
inverse limit gives the unique extension.
Uniqueness makes it \(G\)-invariant.
If the original character kills \(U_D\cap H\), the extension
kills \(T_D\cap H\), because
\((T_D\cap H)/(U_D\cap H)\) is pro-odd.
Thus the relative maximal pro-two abelian quotients on
\(T_D\) and \(U_D\) agree.
The established reduction on \(T_D\) gives \eqref{eq:odd-cm-1}.

Likewise a character of \(T_D\) with values in \(\mathbf F_2\)
and vanishing on \(U_D\) factors through a pro-odd quotient
and is zero. The original restriction kernels on \(T_D\)
and \(U_D\) therefore agree, giving \eqref{eq:odd-cm-2}.
Neither argument assumes \(G=T_DH\).

We construct an arithmetic witness for \eqref{eq:odd-cm-3}.
The arithmetic CM image theorem and the harmonic comparison
give
\[
 \rho_4(M\cap H)=N_4(4)
\]
for the Tate representation of \(E_4:y^2=x^3-x\), and identify
its odd projection on \(D\cap H\) with \(W_{\rm odd}\).
The primary reciprocity input is
\cite[Theorem~12.27, p.~124]{MilneMF110}.
At the split prime five,
\[
 \mathbf Z[i]\otimes\mathbf Z_5
       \cong\mathbf Z_5\times\mathbf Z_5,
\]
and the norm-one torus contains \(v=(6,6^{-1})\), of
infinite pro-five order.
Place this value at five and the identity at all other
components of \(N_4(4)\). Its two-component satisfies the
level-four condition.
Choose \(m\in M\cap H\) with \(\rho_4(m)=v\), and put
\[
 a=m^{\varepsilon_5},
\]
where \(\varepsilon_5\in\widehat{\mathbf Z}\) retains the
five-component and kills all others.
Then \(a\) is pro-five, has full multiplier one, and
\(\rho_4(a)=v\).

Choose \(g\in H\) with full multiplier
\[
 \lambda_2(g)=-1,\qquad
 \lambda_p(g)=1\quad(p\text{ odd}),
\]
and set \(h=g^{\varepsilon_2}\).
Then \(h\) is pro-two and has the same multiplier tuple.
For every \(p\), the action image of \(\ker\lambda_p\) is
pro-\(p\), by its upper-unitriangular action on finite
mod-\(p\) Magnus truncations.
It follows that \(\Phi_p(h)=1\) for every odd \(p\):
its procyclic image is both pro-two and pro-\(p\).
Similarly \(\Phi_p(a)=1\) for \(p\ne5\).

Consequently the arithmetic commutator
\[
                         w=[h,a]
\]
has trivial \(\Phi_p\)-image for every prime \(p\).
Its full multiplier is one.
Every rational Kummer value of \(a\) is supported at five,
where \(\lambda_5(h)=1\), so the conjugation law gives
\[
 \rho_q^I(w)
       =(\lambda(h)-1)\rho_q^I(a)=0
       \qquad(q\text{ a rational prime}).
\]
Thus \(w\in P_D\cap H\).
The cubic defect vanishes on \(D\cap H\), hence
\(w\in T_D\cap H\).

Finally \(h\) acts nontrivially on \(\mathbf Q(i)\), because
its dyadic multiplier is \(-1\) modulo four.
It conjugates the CM coefficient ring by complex conjugation
and therefore inverts its norm-one torus.
Compute in the Tate representation, so that no value
of \(W\) outside its stated domain is needed:
\[
                    W_{\rm odd}(w)=v^{-1}v^{-1}=v^{-2}.
\]
This element still has infinite pro-five order.
The procyclic closure of \(w\) lies in \(T_D\cap H\) and has
infinite image in the quotient by \(U_D\).
Both assertions in \eqref{eq:odd-cm-3} follow.
\end{proof}

\begin{remark}
The strictness witness simultaneously kills every primary
fundamental-group action and every rational Kummer coordinate.
It requires no joint radical--CM image assertion.
The theorem does not claim \(G=U_DH\) or compute the remaining
relative abelian quotient.
\end{remark}

\begin{corollary}\label{cor:full-odd-harmonic-arithmetic-image}
Let \(A_{4,\rm odd}=N_4(4)_{\rm odd}\). Then
\[
 W_{\rm odd}(T_D\cap H)=W_{\rm odd}(T_D)=A_{4,\rm odd},
 \qquad T_D/U_D\cong A_{4,\rm odd}.
\]
Moreover \(T_D=U_D(T_D\cap H)\).
This factorization is internal to \(T_D\).
\end{corollary}
\begin{proof}
For any \(v\in A_{4,\rm odd}\), the arithmetic CM image gives
\(m\in M\cap H\) with odd CM value \(v\).
Put \(a=m^{1-\varepsilon_2}\). Its procyclic closure is pro-odd,
its full multiplier is one, its odd CM value is \(v\), and
its rational Kummer values have zero dyadic component.
In particular \(\Phi_2(a)=1\).
Use the same arithmetic pro-two element \(h\) as in the theorem.
It has trivial \(\Phi_p\)-image for every odd \(p\).
Thus \(w=[h,a]\) has trivial image at every prime.
The Kummer identity
\[
 \rho_q^I(w)=(\lambda(h)-1)\rho_q^I(a)=0
\]
holds componentwise: the dyadic coordinate of \(a\) is zero,
and the odd multiplier coordinates of \(h\) are one.
Therefore \(w\in P_D\cap H\), and arithmetic vanishing of
\(d_{(3)}\) puts it in \(T_D\cap H\).
The same CM conjugation calculation gives
\(W_{\rm odd}(w)=v^{-2}\).
Raising to \(-2\) is an automorphism of the profinite abelian
pro-odd group \(A_{4,\rm odd}\), so this proves the arithmetic
surjectivity. Compactness gives the asserted quotient.
For \(t\in T_D\), choose \(j\in T_D\cap H\) with the same
odd CM value. Then \(tj^{-1}\in U_D\), proving the factorization.
\end{proof}

\begin{theorem}[A strict joint odd CM gate]
\label{thm:joint-odd-cm-gate}
Let \(T_{\rm odd}\) be the maximal pro-odd projection of the
corrected Eisenstein CM homomorphism, and put
\[
                       V_D=U_D\cap\ker T_{\rm odd}.
\]
Then \(V_D\) is closed and \(G\)-normal, and
\begin{equation}
 B_{\max2}\cong
 \left(V_D/\overline{[G,V_D](V_D\cap H)}\right)_{\max2}.
                                                               \label{eq:joint-odd-cm-1}
\end{equation}
The original restriction kernel is
\begin{equation}
 \ker\bigl(X_2\to H^1(V_D,\mathbf F_2)\bigr)
   =\bigoplus_{q\ {\rm odd}}\mathbf F_2\psi_q
                            \oplus\mathbf F_2\eta_3.    \label{eq:joint-odd-cm-2}
\end{equation}
Moreover
\begin{equation}
 [U_D:V_D]=\infty,\qquad
 [U_D\cap H:V_D\cap H]=\infty.                          \label{eq:joint-odd-cm-3}
\end{equation}
\end{theorem}

\begin{proof}
Semilinear conjugation acts on the cubic norm-one torus by
identity or inversion, and preserves its pro-odd factor.
Thus \(V_D\) is closed and \(G\)-normal, and
\(U_D/V_D=T_{\rm odd}(U_D)\) is pro-odd.
A \(G\)-invariant character of \(V_D\) with values in a finite
abelian two-group extends uniquely to \(U_D\), by the central
coprime averaging argument. If it kills \(V_D\cap H\), its
extension kills \(U_D\cap H\), since the quotient of these
intersections is pro-odd.
Hence the relative maximal pro-two abelian quotients agree,
giving \eqref{eq:joint-odd-cm-1} from the established reduction on \(U_D\).
A character of \(U_D\) with values in \(\mathbf F_2\) and
vanishing on \(V_D\) also vanishes, so the original
restriction kernels agree and give \eqref{eq:joint-odd-cm-2}.
No factorization \(G=U_DH\) is used.

For strictness use the arithmetic image of the corrected
Tate representation of
\[
                         E_3:y^2=x^3+1.
\]
The established cubic CM image is \(N_3(4)\).
At the split prime thirteen,
\[
 \mathbf Z[\zeta_3]\otimes\mathbf Z_{13}
           \cong\mathbf Z_{13}\times\mathbf Z_{13},
\]
choose the norm-one element \(v=(14,14^{-1})\), of infinite
pro-thirteen order, and put the identity at every other
component. The level-four condition at two is satisfied.
Choose \(m\in M\cap H\) with this corrected cubic CM value,
and put \(a=m^{\varepsilon_{13}}\).
Then \(a\) is pro-thirteen and has full multiplier one.
Its cubic Tate value is \(v\): the central
two-primary correction in the corrected representation is
killed by the power \(\varepsilon_{13}\).

Choose \(g\in H\) with full multiplier
\[
 \lambda_3(g)=-1,\qquad
 \lambda_p(g)=1\quad(p\ne3),
\]
and let \(h=g^{\varepsilon_2}\).
It is pro-two with the same multiplier tuple.
For every odd prime \(p\ne3\), its action \(\Phi_p(h)\)
is trivial, since it is both pro-two and contained in the
pro-\(p\) image of \(\ker\lambda_p\).
For \(p\ne13\), the normalized action \(\Phi_p(a)\) is
trivial because \(a\) is pro-thirteen.
Thus the commutator
\[
                              w=[h,a]
\]
has trivial \(\Phi_p\)-image for every prime \(p\):
at thirteen use the first entry, and at all other primes
use the second.
Its rational Kummer coordinates vanish as well, because
those of \(a\) are supported at thirteen, where
\(\lambda_{13}(h)=1\):
\[
 \rho_q^I(w)
       =(\lambda(h)-1)\rho_q^I(a)=0.
\]
Consequently \(w\in P_D\cap H\).
The cubic defect vanishes arithmetically on \(D\), so
\(w\in T_D\cap H\).

The element \(h\) fixes \(\mathbf Q(i)\), because its
dyadic multiplier is one. Its harmonic Tate action
is therefore linear over the CM coefficient ring and
commutes with that of \(a\). The harmonic value of \(w\),
and in particular \(W_{\rm odd}(w)\), is one.
Hence \(w\in U_D\cap H\).

On the other hand, \(h\) conjugates
\(\mathbf Q(\sqrt{-3})\), since its multiplier at three
is minus one. It inverts the cubic norm-one torus.
The Tate commutator has value \(v^{-2}\);
the central two-primary correction disappears under the
odd projection. Thus
\[
                         T_{\rm odd}(w)=v^{-2},
\]
which has infinite pro-thirteen order.
The procyclic closure of \(w\) therefore has infinite
image in \(U_D/V_D\), proving both assertions in \eqref{eq:joint-odd-cm-3}.
\end{proof}

\begin{remark}
This witness deliberately uses a prime distinct from two
and three. It proves a strict joint gate but makes no claim
that the entire odd cubic CM image survives on \(U_D\).
In particular the wild three-primary part requires a
separate analysis. The theorem does not assert \(G=V_DH\).
\end{remark}

\section{The exact joint odd CM image on the primary kernel}
\label{sec:exact-joint-odd-cm-image}

We use the groups
\[
 M=\ker\lambda,\qquad
 P_D=D\cap\bigcap_p\ker\Phi_p,\qquad
 T_D=P_D\cap\ker d_{(3)},
\]
\[
 U_D=T_D\cap\ker W_{\rm odd},\qquad
 V_D=U_D\cap\ker T_{\rm odd}.
\]
Here \(D\) is the joint rational Kummer kernel inside \(M\),
\(\Phi_p\) is the free pro-\(p\) action, and
\(T,W\) are the corrected cubic and harmonic CM homomorphisms.
For a set of primes, its primary idempotent in
\(\widehat{\mathbf Z}\) has component one at those primes
and zero at the others; \(\varepsilon_p\) denotes a single
primary idempotent. All profinite powers below are taken
in the closed procyclic subgroup generated by their argument.

For a coefficient prime \(\ell\), write
\[
 \mathcal N_{3,\ell}
  =\{z\in(\mathbf Z_\ell[\omega])^\times:z\bar z=1\},
 \qquad \omega^2+\omega+1=0,
\]
and let \(\mathcal N_{3,\ell,\rm odd}\) be its maximal
pro-odd factor. Put
\begin{equation}\label{eq:exact-joint-odd-1}
\mathcal A_4=N_4(4)_{\rm odd},\qquad
 \mathcal C_3=
 \prod_{\ell\ge5}\bigl(\mathcal N_{3,\ell,\rm odd}\bigr)^3.
\end{equation}
We regard \(\mathcal C_3\) as a subgroup of \(N_3(4)_{\rm odd}\)
by putting the identity at coefficient primes two and three.
At two the level-four subgroup already has trivial odd factor.
The identity at three is an additional restriction.

\begin{theorem}\label{thm:exact-joint-odd-cm-image}
The full and arithmetic joint images on \(T_D\) agree:
\begin{equation}\label{eq:exact-joint-odd-2}
\begin{split}
 (T_{\rm odd},W_{\rm odd})(T_D)
   &=(T_{\rm odd},W_{\rm odd})(T_D\cap H)\\
   &=\mathcal C_3\times\mathcal A_4.
 \end{split}
\end{equation}
Moreover
\begin{equation}\label{eq:exact-joint-odd-3}
T_{\rm odd}(U_D)=T_{\rm odd}(U_D\cap H)=\mathcal C_3.
\end{equation}
Consequently
\begin{equation}\label{eq:exact-joint-odd-4}
T_D/V_D\cong\mathcal C_3\times\mathcal A_4,\qquad
 U_D/V_D\cong\mathcal C_3,
\end{equation}
and
\begin{equation}\label{eq:exact-joint-odd-5}
T_D=V_D(T_D\cap H),\qquad
 U_D=V_D(U_D\cap H).
\end{equation}
These are internal factorizations of the indicated groups.
\end{theorem}

\subsection{The upper bound, including the wild coefficient prime}

\begin{lemma}\label{lem:wild-three-primary-kernel}
For every \(g\in P_D\), the entire coefficient-three cubic
frame is the identity:
\[
                              T_{g,3}=1.
\]
\end{lemma}

\begin{proof}
The action \(\Phi_3(g)\) is the identity.
Hence its beta multiplier in the pro-three metabelian
quotient is the constant series one in
\(\mathbf Z_3[[X-1,Y-1]]\).
Specialization \(X=Y=\omega\) is continuous: in
\(\mathbf Z_3[\omega]\), the element \(\omega-1\) is
topologically nilpotent. Equivalently the corresponding
character of the abelianization has three-power order,
so it factors through the pro-three abelianization.
Thus the coefficient-three value of \(B_g(\omega,\omega)\)
is one.

Apply the established full integral cubic comparison,
Proposition~\ref{prop:lr-integral-cubic-beta}.
It identifies \(T_g\) with \(B_g(\omega,\omega)\) or its
conjugate on \(D\); the sign is one there.
Since \(P_D\subset D\), the asserted identity follows.
This uses the integral comparison at the ramified coefficient
prime, not a reduction modulo three.
\end{proof}

For \(g\in T_D\), the definition of \(d_{(3)}\), together with
\(\rho_\ell^I(g)=0\), gives \(c_{\ell,3}(g)=0\) for every
\(\ell\ge5\).
The defining exponent
\[
 \omega^{c_{\ell,3}(g)}
 =B_{g,\ell}(\omega,\omega)^{
                  (\ell-\chi_{-3}(\ell))/3}
 \quad\text{after reduction modulo }\ell
\]
and the integral cubic comparison say precisely that
\(T_{g,\ell}\) is a cube in \(\mathcal N_{3,\ell}\).
At these primes cubing is an automorphism on the principal
pro-\(\ell\) subgroup; the only condition is on the
three-primary part of the finite residue torus.
Lemma~\ref{lem:wild-three-primary-kernel} supplies the
coefficient-three identity, and the coefficient-two odd
factor is trivial. Therefore
\begin{equation}\label{eq:exact-joint-odd-6}
(T_{\rm odd},W_{\rm odd})(T_D)
        \subset\mathcal C_3\times\mathcal A_4.
\end{equation}

\subsection{Local arithmetic inertia in the cubic CM representation}

We use Tate representations for the arithmetic
constructions below. On \(D\cap H\) the established
comparison identifies their odd projections with the
corrected frames. The central two-primary correction in
the cubic comparison has no effect on these projections.

\begin{lemma}\label{lem:local-cubic-cm-inertia}
Let \(\ell\ge5\), let
\(e_\ell=v_3(\ell-\chi_{-3}(\ell))\), and let
\(Z_\ell\) be the cyclic three-Sylow subgroup of
\(\mathcal N_{3,\ell}\), of order \(3^{e_\ell}\).
There exists an arithmetic element
\[
 x_\ell\in M\cap H\cap\bigcap_p\ker\Phi_p
\]
whose uncorrected cubic Tate value is supported only at \(\ell\)
and generates \(Z_\ell\).
\end{lemma}

\begin{proof}
Write \(K=\mathbf Q(\omega)\), \(O=\mathbf Z[\omega]\), and
\(\pi=1-\omega\).
The earlier arithmetic CM computation gives
\begin{equation}\label{eq:exact-joint-odd-7}
\rho_3(G_K)=
 \{u\in\widehat O^\times:u\equiv1\pmod{2\pi}\}.
\end{equation}
The reduction of \(\mu_6\) onto
\((O/(2\pi))^\times\) is an isomorphism.
Thus \eqref{eq:exact-joint-odd-7} is the unique set of representatives of
\(\widehat O^\times/\mu_6\).
CM reciprocity
\cite[Theorem~12.27, p.~124]{MilneMF110} identifies the
inertia of a prime \(\mathfrak l\) over \(\ell\) with the
local units supported at \(\mathfrak l\), modulo \(\mu_6\).
Because \(\mathfrak l\nmid6\), those supported units already
satisfy the normalization in \eqref{eq:exact-joint-odd-7}.
Consequently their Tate image is the same supported
unit group, up to inversion and the choice of the CM
embedding. In particular its three-primary finite image
is full.

Choose \(\tau\) in this arithmetic inertia with three-primary
Tate value generating that local three-Sylow group, and
replace it by \(\tau^{\varepsilon_3}\).
If \(\ell\) splits in \(K\), its Tate value is, up to orientation,
\((\zeta,1)\), where \(\zeta\) generates the three-Sylow
subgroup of \(\mathbf Z_\ell^\times\).
If \(\ell\) is inert, its value is a generator of the
three-Sylow subgroup in the unramified quadratic unit group.
This latter subgroup has norm one, since \(3\nmid\ell-1\).

Let \(c\) be complex conjugation and set
\[
                         x_\ell=[c,\tau]^{\varepsilon_3}.
\]
The commutator has full multiplier one.
At every coefficient prime \(p\ne\ell\), its free pro-\(p\)
action is trivial because arithmetic inertia at \(\ell\)
acts trivially there
\cite[\S8.5, p.~151]{HainMatsumoto2003}.
At \(p=\ell\), its normalized image is pro-\(\ell\),
so the subsequent three-primary power kills that image.
Thus \(x_\ell\) has every primary action trivial.

Its cubic Tate value is the ratio \(\bar u/u\) of the
chosen local unit.
At a split prime this is \((\zeta^{-1},\zeta)\).
At an inert prime it is \(u^{-2}\), because the
three-Sylow subgroup has norm one.
In both cases it generates \(Z_\ell\), since two is
prime to three. It is supported only at \(\ell\), and
the power \(\varepsilon_3\) leaves it unchanged.
\end{proof}

\subsection{Realizing every permitted cubic component}

We will repeatedly use two elementary facts.
The image of \(\ker\lambda_p\) under \(\Phi_p\) is pro-\(p\):
it acts trivially on the Frattini quotient of the free
pro-\(p\) group, and the kernel of that finite linear action
on automorphisms is pro-\(p\).
Also, for \(a\in M\), every rational Kummer cocycle satisfies
\[
 \rho_q^I(hah^{-1})=\lambda(h)\rho_q^I(a),\qquad
 \rho_q^I([h,a])=(\lambda(h)-1)\rho_q^I(a).
\]
Thus complementary primary orders kill the corresponding
normalized free-group actions, and the displayed formulas
give exact cancellation of the full adelic Kummer values.

Choose once and for all an arithmetic pro-two element
\(h_3\) whose full multiplier is \(-1\) at three and one
at every other prime. It fixes \(\mathbf Q(i)\) and
conjugates \(\mathbf Q(\sqrt{-3})\).

First consider a prescribed cubic odd value in the product of the Sylow subgroups for odd primes
different from three.
The arithmetic corrected cubic image supplies a
multiplier-one lift \(m\).
Take its primary power \(a=m^\varepsilon\), where
\(\varepsilon\) is one at primes other than two and three
and zero at those two primes.
The uncorrected cubic Tate value of \(a\) is the prescribed
value; the central two-primary correction has disappeared.

The commutator \([h_3,a]\) has every primary action trivial.
At two and three use the primary order of \(a\); at all
other primes use that \(h_3\) is pro-two with local
multiplier one.
Its rational Kummer values vanish, since those of \(a\)
are supported where the multiplier of \(h_3\) is one.
It is arithmetic, so its cubic defect vanishes.
Its harmonic CM value is one, because \(h_3\) fixes
\(\mathbf Q(i)\). Its cubic odd value is the inverse
square of the prescribed value.
Since squaring is an automorphism on an odd group, this
realizes the entire part of \(\mathcal C_3\) of primary
order prime to three inside \(U_D\cap H\).

It remains to realize \(Z_\ell^3\) for each \(\ell\ge5\).
Let \(x_\ell\) be as in
Lemma~\ref{lem:local-cubic-cm-inertia}, with cubic value
\(z_\ell\) generating \(Z_\ell\).
Choose \(h_4\in H\) with multiplier four at three and one
at every other prime, and put
\[
 y_\ell=h_4x_\ell h_4^{-1}x_\ell^{-\lambda(h_4)}.
\]
Every primary action remains trivial.
All rational Kummer values cancel by their multiplier-one
conjugation law, so \(y_\ell\in P_D\cap H\).
The element \(h_4\) fixes the cubic CM field, and \(x_\ell\)
has a three-primary Tate value. Therefore
\[
                       \rho_3(y_\ell)=z_\ell^{-3}.
\]
Now set
\[
                            w_\ell=[h_3,y_\ell].
\]
It still belongs to \(P_D\cap H\), hence to \(T_D\cap H\).
The harmonic value is one because \(h_3\) fixes
\(\mathbf Q(i)\).
The cubic odd value is \(z_\ell^6\), which generates
\(Z_\ell^3\). Thus \(w_\ell\in U_D\cap H\) realizes the
full permitted three-primary subgroup at the single
coefficient prime \(\ell\), with identity at the others.

The image of the compact group \(U_D\cap H\) is closed.
It contains every such finite coordinate axis, their
finite products, and therefore the full product of the
groups \(Z_\ell^3\). Combining this with the preceding
prime-to-three construction gives
\begin{equation}\label{eq:exact-joint-odd-8}
T_{\rm odd}(U_D\cap H)=\mathcal C_3.
\end{equation}
No convergence of an ordered infinite noncommuting product
has been assumed.

\subsection{Separating the harmonic component}

For any \(v\in\mathcal A_4\), the established harmonic
arithmetic image supplies \(m\in M\cap H\) with that odd
harmonic value.
Take \(a=m^{1-\varepsilon_2}\).
Choose the arithmetic pro-two element \(h_2\) with
multiplier \(-1\) at two and one at every odd prime.
The commutator \([h_2,a]\) belongs to \(T_D\cap H\) by
the argument for the full odd harmonic image:
its primary actions and all rational Kummer values vanish.
It has harmonic odd value \(v^{-2}\).

In addition, \(h_2\) fixes the cubic CM field, since its
multiplier at three is one.
Its uncorrected cubic Tate action therefore commutes with that
of \(a\), and its cubic odd commutator value is one.
It follows that the joint arithmetic image contains
\[
                         \{1\}\times\mathcal A_4.
\]
Together with \eqref{eq:exact-joint-odd-8}, it contains
\(\mathcal C_3\times\mathcal A_4\).
The upper bound \eqref{eq:exact-joint-odd-6} proves \eqref{eq:exact-joint-odd-2} and \eqref{eq:exact-joint-odd-3}.

The kernel of the joint map on \(T_D\) is \(V_D\).
Compactness gives \eqref{eq:exact-joint-odd-4}.
For each element of \(T_D\), choose an arithmetic member
of \(T_D\) with the same pair of odd CM values and divide
by it. This proves the first equality in \eqref{eq:exact-joint-odd-5}.
The same argument on \(U_D\), using \eqref{eq:exact-joint-odd-8}, proves the second.

\begin{remark}[Internal factorization and relative coinvariants]
The quotients in \eqref{eq:exact-joint-odd-4} are pro-odd, so the established
maximal pro-two relative quotient and original \(X_2\)
restriction kernel remain unchanged.
The internal factorizations in \eqref{eq:exact-joint-odd-5} do not, by themselves,
identify the full relative coinvariant quotients.

For a precise counterexample, let \(p\) be odd and let
\(G=\langle a,b,c\rangle\) be the order-\(p^3\) Heisenberg
group, with \(c\) central, \([a,b]=c\), and each generator
of order \(p\). Put
\[
 H=\langle b\rangle,\qquad T=\langle b,c\rangle,\qquad
 V=\langle c\rangle.
\]
Both \(T,V\) are \(G\)-normal and \(T=V(T\cap H)\).
Nevertheless
\[
 V/[G,V](V\cap H)\cong C_p,\qquad
 T/[G,T](T\cap H)=1.
\]
The character of \(V\) sending \(c\) to one extends to \(T\)
by sending \(b\) to zero, but this extension is not
\(G\)-invariant: conjugation by \(a\) sends \(b\) to \(bc\).
Thus cross-conjugation of \(T\cap H\) by \(G\) is a real
additional obstruction. No full relative coinvariant
identification is inferred from \eqref{eq:exact-joint-odd-5}.
\end{remark}

\subsection{An arithmetic cross-cup witness with trivial word and CM actions}
\label{sec:cross-cup-residual-inertia}

\begin{proposition}\label{prop:cross-cup-residual-inertia}
Choose pro-two elements \(\tau_5\in I_5\) and
\(\tau_{29}\in I_{29}\) of arithmetic inertia with
\(c_5(\tau_5)=c_{29}(\tau_{29})=1\).
Then
\[
                         w=[\tau_5,\tau_{29}]
\]
belongs to \(V_D\cap H\). In fact every rational Kummer
coordinate, every free pro-\(p\) action, and both
complete corrected CM frames vanish on \(w\).
The character \(\kappa_{5,29}\) of
\[
 \mathbf Q(\sqrt5,\sqrt{29},\sqrt{7+2\sqrt5})
\]
has value one on \(w\).
\end{proposition}

\begin{proof}
Cyclotomic inertia at an odd prime \(\ell\) is supported
only at the coefficient prime \(\ell\), and its projection
onto the quadratic character \(c_\ell\) is surjective.
Taking the pro-two power of a suitable inertia element
therefore supplies the stated \(\tau_\ell\).

Put \(u=\lambda(\tau_5)\) and \(v=\lambda(\tau_{29})\).
Both are one at the coefficient prime two, while \(u\)
is one away from five and \(v\) is one away from twenty-nine.
For any rational Kummer cocycle \(\rho_q\), the crossed
homomorphism identity gives, componentwise,
\begin{equation}
 \rho_q([\tau_5,\tau_{29}])
       =(1-v)\rho_q(\tau_5)+(u-1)\rho_q(\tau_{29}).        \label{eq:cross-cup-inertia-cocycle}
\end{equation}
At two both coefficients vanish.
At five the only possible term is
\((u_5-1)\rho_{q,5}(\tau_{29})\), and this is zero:
on the pro-two cyclic group generated by \(\tau_{29}\),
the fifth multiplier is one, so \(\rho_{q,5}\) is a
continuous homomorphism to the pro-five additive group.
At twenty-nine the symmetric argument applies.
At every other prime both coefficients vanish.
Thus \eqref{eq:cross-cup-inertia-cocycle} is zero in the full adelic target, for every
rational \(q\), including \(q=2\). The commutator also
has full multiplier one. Therefore \(w\in D\cap H\).

Arithmetic inertia at \(\ell\) acts trivially in the
free pro-\(p\) action for \(p\ne\ell\), by
tangential unramifiedness
\cite[\S8.5, p.~151]{HainMatsumoto2003}.
For \(p=5\), the second commutator entry has trivial
action; for \(p=29\), the first does; at every other
prime both do. Hence \(w\in P_D\cap H\).
The cubic defect vanishes on \(D\cap H\), so
\(w\in T_D\cap H\).

Both inertia elements fix \(\mathbf Q(i)\) and
\(\mathbf Q(\sqrt{-3})\), whose finite ramification is
confined to two and three.
Their CM Tate actions therefore belong to the
respective commutative CM unit groups.
Both Tate commutators are the identity.
On \(D\cap H\), the established signed comparisons
identify the corrected frames with these Tate actions;
the central cubic correction is trivial on \(D\cap H\).
Thus \(T(w)=W(w)=1\), in every coefficient component,
and in particular \(w\in V_D\cap H\).

In the dihedral field defining \(\kappa_{5,29}\),
the images of \(\tau_5\) and \(\tau_{29}\) have respective
biquadratic coordinates \((1,0)\) and \((0,1)\):
the other quadratic character is unramified at that
inertia prime.
For its central extension class \(c_5c_{29}\), the
commutator of such two lifts is the central generator,
independently of the choices of lifts.
Thus \(\kappa_{5,29}(w)=1\).
\end{proof}

The proposition proves arithmetic nonvanishing on the
residual subgroup $V_D$ itself. It does not assert
that \(\kappa_{5,29}\) extends to a full \(G\)-invariant
character. That is a separate cohomological lifting
question.

\section{A mixed finite cover lifts a cross-cup arithmetic character}
\label{sec:cross-cup-five-twenty-nine-lift}

Let \(G_0\) be the positive peripheral automorphism group
containing the profinite Grothendieck--Teichm\"uller group
\(G\). Its action on \(F=\widehat F(x,y)\) has the form
\begin{equation}
 \phi_g(x)=x^{\lambda_g},\qquad
 \phi_g(y)=f_g^{-1}y^{\lambda_g}f_g,\qquad
 \phi_g(z)\sim z^{\lambda_g},\qquad z=(xy)^{-1}.
                                                               \label{eq:cross-cup-5-29-1}
\end{equation}
Here conjugacy in the last formula is in \(F\), and the
same multiplier occurs on all three peripheral generators.
Assume \(f_g\in F\,{}^{\prime}\), as in the standard GT action. Put \(M_0=\ker(\lambda:G_0\to\widehat{\mathbf Z}^{\times})\).

\begin{theorem}\label{thm:five-twenty-nine-cross-cup-lift}
There is a continuous \(G_0\)-invariant homomorphism
\begin{equation}
                    \xi_{5,29}:M_0\longrightarrow\mathbf F_2
                                                               \label{eq:cross-cup-5-29-2}
\end{equation}
whose restriction to \(G_K\), for
\(K=\mathbf Q(\mu_\infty,2^{1/\infty})\), is precisely the
central quadratic character \(\kappa_{5,29}\) of
\[
 \mathbf Q\!\left(\sqrt5,\sqrt{29},\sqrt{7+2\sqrt5}\right).
\]
In particular its restriction to
\(N=\ker(\lambda,\rho_2)\subset G\) belongs to
\(X_2=H^1(N,\mathbf F_2)^G\), and its arithmetic restriction
is outside the span of the rational quartic Kummer
characters and \(\eta_3\).

Equivalently, the corrected cross-cup class
\begin{equation}
                         e_{c_5,c_{29}}=c_5\cup c_{29}
                                                               \label{eq:cross-cup-5-29-3}
\end{equation}
has zero inflation to \(H^2(G,\mathbf F_2)\).
The construction uses the homology modulo twenty-nine of
a cyclic degree-five cover, and applies already to \(G_0\).
\end{theorem}

\subsection{The cyclic cover and its complete peripheral pairing}

Let
\begin{equation}
 U_5=\ker\bigl(F\longrightarrow C_5,\quad x\mapsto1,\ y\mapsto1\bigr).
                                                               \label{eq:cross-cup-5-29-4}
\end{equation}
Every \(\phi_g\) preserves \(U_5\), because it multiplies
the indicated abelian quotient by
\(\lambda_g\bmod5\). Put
\begin{equation}
 V=H_1(U_5,\mathbf F_{29})/
             \langle\text{peripheral homology classes}\rangle .
                                                               \label{eq:cross-cup-5-29-5}
\end{equation}
This definition uses the subgroup action in \eqref{eq:cross-cup-5-29-1} directly;
it does not assume equivariance of any additional
covering or filling operation.

\begin{lemma}\label{lem:five-cover-similitude}
The vector space \(V\) has dimension four and carries a
nondegenerate alternating form.
The induced representation
\begin{equation}
                    \rho:G_0\longrightarrow
                                  \operatorname{GSp}(V)
                                                               \label{eq:cross-cup-5-29-6}
\end{equation}
has similitude multiplier \(\lambda_g\bmod29\).
If \(D\) is the deck transformation induced by conjugation
by \(x\), then
\begin{equation}
                  \rho(g)D\rho(g)^{-1}=D^{\lambda_g\bmod5}.
                                                               \label{eq:cross-cup-5-29-7}
\end{equation}
The characteristic polynomial of \(D\) on \(V\) is
\(1+X+X^2+X^3+X^4\).
\end{lemma}

\begin{proof}
Start with the discrete oriented pair of pants.
The degree-five cover specified by \eqref{eq:cross-cup-5-29-4} has Euler
characteristic \(-5\).
Each of \(x,y,z\) maps to a generator of \(C_5\)
(their images are \(1,1,-2\)), so the cover has exactly
three punctures.
Its compactification has genus two.
Consequently \(U_5\) is the profinite completion of a
free group of rank six with a surface presentation
\begin{equation}
 [a_1,b_1][a_2,b_2]c_1c_2c_3=1.                           \label{eq:cross-cup-5-29-8}
\end{equation}
The three peripheral classes have the sole homology
relation \(c_1+c_2+c_3=0\).
Their span has dimension two, and \eqref{eq:cross-cup-5-29-5} has dimension four.

There is a single cusp above each original cusp.
Thus \eqref{eq:cross-cup-5-29-1}, restricted to \(U_5\), sends each \(c_i\) to
a \(U_5\)-conjugate of \(c_i^{\lambda_g}\).
To check the conjugator assertion, a conjugator in \(F\)
may be multiplied by a power of the relevant original
peripheral generator to move it into \(U_5\), since that
peripheral generator maps onto the deck group; this
multiplication does not change its conjugation on \(c_i\).

We prove the multiplier assertion algebraically.
Take the class-two, exponent-twenty-nine quotient of
\(U_5\).
Its commutator space is the exterior square of its
first homology.
Project that exterior square onto \(\bigwedge^2 V\).
Conjugating a peripheral word changes its central
coordinate by a wedge with its peripheral homology
class, which vanishes under this projection.
Taking a \(\lambda_g\)-th power multiplies its projected
central coordinate by \(\lambda_g\bmod29\).
Relation \eqref{eq:cross-cup-5-29-8} therefore implies
\begin{equation}
 \bigwedge\nolimits^2\rho(g)
       (a_1\wedge b_1+a_2\wedge b_2)
   =\lambda_g
       (a_1\wedge b_1+a_2\wedge b_2)
       \quad\text{in }\bigwedge\nolimits^2V.               \label{eq:cross-cup-5-29-9}
\end{equation}
Here the displayed \(a_i,b_i\) mean their homology classes.
For an explicit way to read this calculation, eliminate
\(c_3\) from \eqref{eq:cross-cup-5-29-8}, use
\(a_1,b_1,a_2,b_2,c_1,c_2\) as free generators, and take
class-two logarithms, which are valid because twenty-nine
is odd. The projected central coordinate of \(c_3\) is
the negative of the bivector in \eqref{eq:cross-cup-5-29-9}, while those of
\(c_1,c_2\) are zero. The peripheral conjugation and
power rule gives exactly \eqref{eq:cross-cup-5-29-9}.

That bivector is nondegenerate. Its inverse alternating
form consequently has the same similitude multiplier.
This proves \eqref{eq:cross-cup-5-29-6}, including the full multiplier statement.

Conjugation by \(x\) induces the deck action, and
\(\phi_g(x)=x^{\lambda_g}\) proves \eqref{eq:cross-cup-5-29-7}.
Finally the cellular chain complex of the five-sheeted
cover of the two-generator graph has
\[
 C_1=\mathbf F_{29}[C_5]^2,\qquad C_0=\mathbf F_{29}[C_5].
\]
At every nontrivial deck character \(\zeta^j\), its
boundary map is the nonzero row
\((\zeta^j-1,\zeta^j-1)\); its kernel is one-dimensional.
The trivial character has multiplicity two in homology,
and is exactly the peripheral subspace.
Thus each nontrivial fifth-root character occurs once
in \(V\), proving the characteristic polynomial.
\end{proof}

\subsection{The finite symplectic normalizer}

Write \(q=29\), \(k=\mathbf F_{q^2}\), and
\(\bar z=z^q\). Choose \(\zeta\in k\) of order five.
Since \(q\equiv-1\pmod5\), the four deck characters
split into the two Frobenius pairs
\(\{\zeta,\zeta^{-1}\}\) and
\(\{\zeta^2,\zeta^{-2}\}\).
We may identify \(V\) with \(k\oplus k\), with
\begin{equation}
                         D(z_1,z_2)=(\zeta z_1,\zeta^2z_2).
                                                               \label{eq:cross-cup-5-29-10}
\end{equation}
The two summands are orthogonal and nondegenerate.
On each, a deck-invariant alternating form is
\(\operatorname{Tr}_{k/\mathbf F_q}(d_i z\bar w)\),
with \(d_i\ne0\) and \(\bar d_i=-d_i\).
Rescaling the coordinate by an element of prescribed
norm allows taking \(d_1=d_2=d\).
The full form is then
\begin{equation}
 \langle z,w\rangle
 =\operatorname{Tr}_{k/\mathbf F_q}
       \bigl(d(z_1\bar w_1+z_2\bar w_2)\bigr),
 \qquad \bar d=-d.                                        \label{eq:cross-cup-5-29-11}
\end{equation}

Let
\[
                    P(z_1,z_2)=(z_2,\bar z_1).
\]
Then \(P^4=1\) and \(PDP^{-1}=D^2\).
Every linear map normalizing \(\langle D\rangle\)
is uniquely of the form \((v_1,v_2)P^j\), with
\(v_i\in k^\times\) and \(j\in\mathbf Z/4\).
Let \(E\) be the subgroup of such maps which are
similitudes of \eqref{eq:cross-cup-5-29-11}, and write
\[
 \pi:E\longrightarrow C_4\times\mathbf F_q^\times
\]
for deck exponent \(2^j\) and similitude multiplier.
Its kernel is
\begin{equation}
 T=\{(v_1,v_2):N(v_1)=N(v_2)=1\}
                   \simeq C_{30}\times C_{30}.             \label{eq:cross-cup-5-29-12}
\end{equation}
Define
\begin{equation}
                 (-1)^{\chi(v_1,v_2)}=(v_1v_2)^{15}.
                                                               \label{eq:cross-cup-5-29-13}
\end{equation}
This is a nonzero homomorphism \(T\to\mathbf F_2\).

\begin{lemma}\label{lem:finite-cross-cup-extension}
The subgroup \(\ker\chi\) is normal in \(E\), and
\begin{equation}
 1\longrightarrow C_2\longrightarrow E/\ker\chi
       \longrightarrow C_4\times\mathbf F_{29}^\times
       \longrightarrow1                                 \label{eq:cross-cup-5-29-14}
\end{equation}
is a central extension with class \(x\cup y\).
Here \(x\) is the nonzero quadratic character of \(C_4\),
and \(y\) is the quadratic character of
\(\mathbf F_{29}^\times\).
\end{lemma}

\begin{proof}
Conjugation by the scalar subgroup is trivial on \(T\).
Conjugation by \(P\) sends \((v_1,v_2)\) to
\((v_2,\bar v_1)=(v_2,v_1^{-1})\) on \(T\).
This leaves \eqref{eq:cross-cup-5-29-13} unchanged, because \(v_1^{15}\)
has order at most two.
Thus \(\ker\chi\) is normal and its quotient kernel is
central.

Choose \(\beta\in k^\times\) with \(N\beta=-1\), and set
\[
                          J=(1,\beta)P.
\]
From \eqref{eq:cross-cup-5-29-11}, \(J\) is a similitude of multiplier one:
\(P\) interchanges the two summands and reverses the
sign of the second contribution, which multiplication
by \(\beta\) reverses again.
More explicitly, a map \((v_1,v_2)P\) has multiplier
\(t\) precisely when \(Nv_1=t\) and \(Nv_2=-t\).
Thus \(\pi(J)=(1,1)\).

Choose \(v\in k^\times\) with norm a generator \(t\)
of \(\mathbf F_{29}^\times\), and put
\[
                             L=(v,v).
\]
Then \(\pi(L)=(0,t)\).
These two images generate the entire base of \eqref{eq:cross-cup-5-29-14};
in particular \(\pi\) is surjective.

There are three independent degree-two classes on
\(C_4\times C_{28}\): the two cyclic extension classes
and the mixed cup product \(xy\).
Their coefficients are determined respectively by
the fourth power of a lift of the first generator,
the twenty-eighth power of a lift of the second,
and their commutator.
We compute all three in the central quotient.

First,
\[
                       J^4=(-1,-1).
\]
This need not be the identity
in \(E\).
Nevertheless
\[
                       \chi(J^4)=0,
\]
because the two norm-one quadratic signs cancel.
Next,
\[
                       L^{28}=(v^{28},v^{28})\in T,
 \qquad                 \chi(L^{28})=0.
\]
Finally,
\[
                       [J,L]=(1,\bar v/v)\in T,
\]
and
\begin{equation}
 \begin{split}
 (-1)^{\chi([J,L])}
   &=(\bar v/v)^{15}
     =v^{420}
     =(Nv)^{14}
     =t^{14}
     =-1 .
 \end{split}                                               \label{eq:cross-cup-5-29-15}
\end{equation}
Thus both cyclic coefficients vanish and the mixed
coefficient is one. This proves the claimed exact
extension class.
\end{proof}

\subsection{The full GT lift and its explicit beta formula}

By Lemma~\ref{lem:five-cover-similitude}, the representation
\(\rho\) takes values in \(E\), and
\begin{equation}
                   \pi\rho(g)=(\lambda_g\bmod5,\lambda_g\bmod29).
                                                               \label{eq:cross-cup-5-29-16}
\end{equation}
The identification of \((\mathbf Z/5)^\times\) with
\(C_4\) uses its generator \(2\).
Under \eqref{eq:cross-cup-5-29-16}, the characters \(x,y\) in
Lemma~\ref{lem:finite-cross-cup-extension} pull back
to \(c_5,c_{29}\), respectively.
Thus \(\rho\), followed by \(E\to E/\ker\chi\), is a
full-\(G_0\) lift of the cyclotomic finite quotient
through the central extension of class \(c_5c_{29}\).
This proves its vanishing in \(H^2(G_0,\mathbf F_2)\),
and hence in \(H^2(G,\mathbf F_2)\).

For \(g\in M_0\), both components of \eqref{eq:cross-cup-5-29-16} are trivial,
so \(\rho(g)\in T\).
Define
\[
                          \xi_{5,29}(g)=\chi(\rho(g)).
\]
Normality and centrality in the preceding lemma prove
that this is a continuous \(G_0\)-invariant homomorphism.

The construction also gives a direct beta expression.
Write \(B_g(X,Y)\) for the normalized Alexander coefficient
in the usual convention: when \(\lambda_g=1\), it is the
scalar action on the nontrivial character line of
\(F'/F''\).
Then
\begin{equation}
 (-1)^{\xi_{5,29}(g)}
  =\left(
      B_g(\zeta,\zeta)
      B_g(\zeta^2,\zeta^2)
    \right)^{15}\qquad(g\in M_0),                          \label{eq:cross-cup-5-29-17}
\end{equation}
with coefficients reduced modulo twenty-nine.
The commutator \([x,y]\) projects nontrivially to each
of these character lines: in the graph chain complex
its value at \((\zeta^j,\zeta^j)\) is a nonzero multiple
of the difference of the two edge classes.
Thus the scalar in the homology construction is indeed
the indicated beta value, without a change-of-basis
factor.
The norm-one property follows from \eqref{eq:cross-cup-5-29-11} and multiplier
one.
Replacing \(\zeta\) by another primitive fifth root
permutes the two factors with possible inversions,
which leave their quadratic sign unchanged.
For the Alexander convention, compare
\cite[\S2.3, p.~101]{Furusho2003}; the finite homology
construction above is also an independent definition
of \eqref{eq:cross-cup-5-29-2}.

Apply naturality of transgression first to
\(1\to M_0\to G_0\to\lambda(G_0)\to1\), and then to
\(1\to N\to G\to A\to1\).
The central extension \eqref{eq:cross-cup-5-29-14} gives
\begin{equation}
                         d_G(\xi_{5,29}|_N)=c_5c_{29}
                              \quad\text{in }H^2(A,\mathbf F_2).
                                                               \label{eq:cross-cup-5-29-18}
\end{equation}
Here \eqref{eq:cross-cup-5-29-18} can equivalently be read by pulling back \eqref{eq:cross-cup-5-29-14}
along \(A\to U\to C_4\times\mathbf F_{29}^\times\);
the map from \(G\) to that pullback has central restriction
\(\xi_{5,29}|_N\).

On the arithmetic subgroup, the independently computed
dihedral character \(\kappa_{5,29}\) has the same
transgression \(c_5c_{29}\).
Its Hilbert corrections vanish because
\[
                         N_{\mathbf Q(\sqrt5)/\mathbf Q}
                                  (7+2\sqrt5)=29.
\]
The arithmetic transgression is injective, since \(K\)
contains \(\mathbf Q^{\mathrm{ab}}\).
Therefore
\[
                         \xi_{5,29}|_{G_K}=\kappa_{5,29}.
\]
This proves Theorem~\ref{thm:five-twenty-nine-cross-cup-lift}.

\subsection{A further strict normalization}

Retain the established factorization \(G=R_DH\), the
residual subgroup \(V_D\), and the established
fact that \(R_D/V_D\) is pro-odd.
Define
\begin{equation}
 R_D^\xi=R_D\cap\ker\xi_{5,29},\qquad
 V_D^\xi=V_D\cap\ker\xi_{5,29}.                            \label{eq:cross-cup-5-29-19}
\end{equation}
Restrictions of \(\xi_{5,29}\) to these groups are
understood through their inclusions in \(M_0\).

\begin{corollary}\label{cor:cross-cup-strict-normalization}
The subgroups in \eqref{eq:cross-cup-5-29-19} are closed and \(G\)-normal, and
\begin{equation}
 [R_D:R_D^\xi]=[V_D:V_D^\xi]=2,\qquad
                         G=R_D^\xi H.                     \label{eq:cross-cup-5-29-20}
\end{equation}
The maximal pro-two relative quotient is unchanged
on passage to \(V_D^\xi\):
\begin{equation}
 B_{\max2}\simeq
 \left(
 V_D^\xi/
 \overline{[G,V_D^\xi](V_D^\xi\cap H)}
 \right)_{\max2}.                                         \label{eq:cross-cup-5-29-21}
\end{equation}
If the established kernel of restriction from \(X_2\)
to \(H^1(V_D,\mathbf F_2)\) is the span of the rational
Kummer characters and \(\eta_3\), then the kernel of
restriction to \(H^1(V_D^\xi,\mathbf F_2)\) is exactly
that span plus the independent character
\(\xi_{5,29}|_N\).
No factorization \(G=V_D^\xi H\) is asserted.
\end{corollary}

\begin{proof}
The arithmetic inertia commutator of
Proposition~\ref{prop:cross-cup-residual-inertia} belongs
to \(V_D\cap H\) and has \(\xi_{5,29}\)-value one.
This proves both indices in \eqref{eq:cross-cup-5-29-20}.
For any \(r\in R_D\), multiplying by that arithmetic
element when necessary gives an element of \(R_D^\xi\).
Together with \(G=R_DH\), this proves \(G=R_D^\xi H\).

The structural relative quotient may therefore be
computed using \(R_D^\xi\).
The quotient \(R_D^\xi/V_D^\xi\) embeds into the pro-odd
group \(R_D/V_D\), so the established central coprime
extension argument identifies their maximal pro-two
relative quotients. This proves \eqref{eq:cross-cup-5-29-21}, using a global
factorization at \(R_D^\xi\), not an unproved internal
factorization assertion at \(V_D^\xi\).

Finally, a character in \(X_2\) vanishing on \(V_D^\xi\)
restricts on \(V_D\) to a character of its quotient
\(C_2\).
Thus its restriction is either zero or
\(\xi_{5,29}|_{V_D}\).
Subtracting the latter when necessary gives precisely
the old restriction kernel.
The arithmetic commutator proves independence from
that old kernel.
\end{proof}

\section{Lifting every cross-cup axis indexed by two odd primes}
\label{sec:all-odd-cross-cup-lifts}

Let \(G_0\) act on \(\widehat F(x,y)\) by positive
peripheral automorphisms with common multiplier
\(\lambda\), with \(x\mapsto x^\lambda\),
\(y\mapsto f_g^{-1}y^\lambda f_g\) and \(f_g\in\widehat F(x,y)'\).
Let \(M_0=\ker\lambda\).
For distinct odd primes \(\ell,q\), write
\[
 m=\ell-1,\qquad d=\operatorname{ord}_{\ell}(q),\qquad
 \ell^*=(-1)^{(\ell-1)/2}\ell,\qquad
 q^*=(-1)^{(q-1)/2}q .
\]
Use the cyclotomic quadratic characters \(c_\ell,c_q\)
and the cyclic degree-two classes \(\delta_\ell,\delta_q\).
Write \(h_v(r,s)\) for the additive local Hilbert sign; in
the formulas below \(h_\ell(\ell^*,q^*)=\ell_\ell(q^*)\)
and \(h_q(\ell^*,q^*)=\ell_q(\ell^*)\).

\begin{theorem}\label{thm:all-odd-cross-cup-lifts}
For every distinct pair of odd primes \(\ell,q\), there
is a continuous \(G_0\)-invariant homomorphism
\[
                         \xi_{\ell,q}:M_0\to\mathbf F_2
\]
whose transgression on \(\lambda(G_0)\), pulled back from the
finite cyclotomic quotient, is
\begin{equation}
 e_{c_\ell,c_q}
   =c_\ell c_q+
        h_\ell(\ell^*,q^*)\delta_\ell+
        h_q(\ell^*,q^*)\delta_q.                          \label{eq:all-odd-cross-cup-1}
\end{equation}
Thus every such class has zero inflation to full
profinite GT. The restriction of \(\xi_{\ell,q}\) to
\(G_K\), for \(K=\mathbf Q(\mu_\infty,2^{1/\infty})\),
is exactly the arithmetic basis character
\(\kappa_{\ell,q}\) transgressing to \eqref{eq:all-odd-cross-cup-1}.

In particular all the odd-prime-pair axes of the complete
arithmetic target lie in the image of \(X_2\).
The characters indexed by \(\{a,c_p\}\) are not covered
by this theorem.
\end{theorem}

We give both an explicit beta formula and the complete
finite central extension calculation. The order of the
two primes records which cyclic cover and which
coefficient field are used; no equality between
\(\xi_{\ell,q}\) and \(\xi_{q,\ell}\) on full GT is assumed.

\subsection{The cover and its symplectic representation}

Put
\[
 U_\ell=\ker\bigl(\widehat F(x,y)\to C_\ell,\quad x,y\mapsto1\bigr).
\]
The associated degree-\(\ell\) cover of the oriented
pair of pants has exactly three cusps, since
\(1,1,-2\) are generators of \(C_\ell\).
Its compact genus is \((\ell-1)/2\).
Define \(V\) to be its first homology over \(\mathbf F_q\),
modulo the peripheral homology subspace.
Then
\[
 \dim_{\mathbf F_q}V=m,
\]
and the deck generator \(D\) has each nontrivial
\(\ell\)-th-root character exactly once.

The positive peripheral action gives
\begin{equation}
 \rho_{\ell,q}:G_0\to\operatorname{GSp}(V),\qquad
 \rho_{\ell,q}(g)D\rho_{\ell,q}(g)^{-1}
                  =D^{\lambda_g\bmod\ell},                \label{eq:all-odd-cross-cup-2}
\end{equation}
with similitude multiplier \(\lambda_g\bmod q\).
Here is the algebraic verification of the multiplier,
to include its full scope. The surface presentation is
\[
       \prod_{i=1}^{(\ell-1)/2}[a_i,b_i]\,c_1c_2c_3=1 .
\]
There is one cusp above each original cusp.
A peripheral conjugator in \(\widehat F(x,y)\) can
therefore be moved into \(U_\ell\) by multiplying it by
a suitable power of that original peripheral generator.
Thus each \(c_i\) is sent to a \(U_\ell\)-conjugate of
\(c_i^\lambda\).
In the class-two, exponent-\(q\) quotient, project
commutators to \(\bigwedge^2 V\).
Conjugation changes a peripheral central coordinate
only by a wedge with a peripheral homology class,
which vanishes after projection.
The \(\lambda\)-th power multiplies that projected
coordinate by \(\lambda\bmod q\).
Consequently the symplectic bivector
\(\sum_i a_i\wedge b_i\), and hence its inverse alternating
form, has multiplier \(\lambda\bmod q\).
The deck spectrum follows directly from the chain
complex \(\mathbf F_q[C_\ell]^2\to\mathbf F_q[C_\ell]\):
each nontrivial character has a one-dimensional kernel,
while the two trivial classes are the peripheral subspace.
This proves \eqref{eq:all-odd-cross-cup-2} without an additional covering or filling
equivariance assumption.

Let \(E\) be the normalizer of \(\langle D\rangle\)
in \(\operatorname{GSp}(V)\), and let
\begin{equation}
                 \pi:E\to C_m\times\mathbf F_q^\times
                                                               \label{eq:all-odd-cross-cup-3}
\end{equation}
record the deck exponent and similitude multiplier.
We now describe its kernel \(T\) and an invariant
quadratic character \(\chi:T\to\mathbf F_2\).

\subsection{Even Frobenius order}

Suppose \(d=2s\), and put \(r=m/d\).
The \(r\) Frobenius orbits are self-dual.
Write
\[
 k=\mathbf F_{q^{2s}},\qquad k_0=\mathbf F_{q^s},\qquad
 \bar z=z^{q^s}.
\]
Then \(V=k^r\) as an \(\mathbf F_q\)-space.
Each summand has alternating form
\(\operatorname{Tr}_{k/\mathbf F_q}(d_i z\bar w)\),
where \(\bar d_i=-d_i\).
Rescaling by elements of prescribed \(k/k_0\)-norm
makes all \(d_i\) equal to a common \(d_0\).
Thus
\begin{equation}
 T=\prod_{i=1}^r\ker(N_{k/k_0}:k^\times\to k_0^\times),
 \qquad
 (-1)^{\chi(t)}=\prod_{i=1}^r t_i^{(q^s+1)/2}.            \label{eq:all-odd-cross-cup-4}
\end{equation}
Every factor has even order \(q^s+1\), so \(\chi\ne0\).

Choose a primitive residue \(a\bmod\ell\).
Write \(q=a^{rh}\), with \(\gcd(h,d)=1\), and choose
\(k_1\) modulo \(d\) such that \(a^r=q^{k_1}\).
Then \(\gcd(k_1,d)=1\), so \(k_1\) is odd.
Using orbit representatives \(1,a,\ldots,a^{r-1}\),
a lift of the deck exponent \(a\) cyclically permutes
the \(r\) summands, with a \(q^{k_1}\)-Frobenius twist
on the last passage.

Such a lift can be chosen symplectic.
Indeed, the plain semilinear permutation preserves
the common form on each unwrapped summand.
On the wrapped summand its form coefficient is changed
from \(d_0\) to \(d_0^{q^{-k_1}}\).
Their ratio belongs to \(k_0^\times\).
Multiplying that summand by an element with the
appropriate \(k/k_0\)-norm restores the form, since
the norm map is onto. Denote the resulting symplectic
lift by \(J\).

Choose a generator \(g\) of \(k^\times\), and take
\[
 v=g^{(q^s-1)/(q-1)},\qquad \lambda=Nv.
\]
Then \(\lambda\) generates \(\mathbf F_q^\times\).
The uniform scalar \(L=(v,\ldots,v)\) is a similitude
of multiplier \(\lambda\).
Thus \eqref{eq:all-odd-cross-cup-3} is surjective.

The character \eqref{eq:all-odd-cross-cup-4} is invariant under every permutation
of its factors and every Frobenius automorphism.
Scalar multiplication used in the lift \(J\) commutes
with \(T\). Hence \(\chi\) is \(E\)-invariant.

The scalar correction in \(J\) also commutes with \(L\).
Therefore their commutator is the identity in every
factor except the wrapped one, where it is
\(v^{q^{k_1}-1}\).
This has norm one, and
\begin{equation}
 \begin{split}
 \left(v^{q^{k_1}-1}\right)^{(q^s+1)/2}
 &=g^{\,\frac{q^{2s}-1}{2}
                    \frac{q^{k_1}-1}{q-1}}\\
 &=(-1)^{\,1+q+\cdots+q^{k_1-1}}
 =-1 .
 \end{split}                                               \label{eq:all-odd-cross-cup-5}
\end{equation}
Thus \(\chi([J,L])=1\).
Also
\begin{equation}
                         \chi(L^{q-1})=r\bmod2.            \label{eq:all-odd-cross-cup-6}
\end{equation}

\subsection{Odd Frobenius order}

Suppose \(d\) is odd, and put \(r=m/(2d)\) and
\(k=\mathbf F_{q^d}\).
Frobenius orbits occur in distinct dual pairs.
The representation is a sum of \(r\) hyperbolic
\(k\)-pairs, with form
\[
 \sum_{i=1}^r
 \operatorname{Tr}_{k/\mathbf F_q}(x_i y_i'-y_i x_i').
\]
Rescaling paired coordinates gives this normalization.
The torus and its character are
\begin{equation}
 T=(k^\times)^r,\qquad
 t_i:(x_i,y_i)\mapsto(t_i x_i,t_i^{-1}y_i),\qquad
 (-1)^{\chi(t)}
      =\prod_{i=1}^r t_i^{(q^d-1)/2}.                     \label{eq:all-odd-cross-cup-7}
\end{equation}
Again \(\chi\ne0\).

For a primitive residue \(a\bmod\ell\), the quotient
by \(\langle q,-1\rangle\) has order \(r\), and
\[
                            a^r=-q^{k_1}
\]
for an integer \(k_1\).
Thus the primitive deck multiplier cycles the \(r\)
pairs and on the last passage exchanges the two
dual coordinates with a Frobenius twist.
A sign on one of the two exchanged coordinates
makes this last operation symplectic.
This gives a symplectic lift \(J\).

Take \(L\) on every pair to be
\[
                              (x,y)\mapsto(x,\lambda y),
\]
where \(\lambda\) generates \(\mathbf F_q^\times\).
It is a similitude of multiplier \(\lambda\), proving
surjectivity of \eqref{eq:all-odd-cross-cup-3}.
Permutations, Frobenius twists, and inversion preserve
the sum of quadratic characters in \eqref{eq:all-odd-cross-cup-7}, so \(\chi\)
is invariant under \(E\).

The commutator \([J,L]\) is trivial except on the
wrapped pair, where its torus parameter is
\(\lambda\) or its inverse.
Its quadratic sign in \(k\) is
\begin{equation}
               \lambda^{(q^d-1)/2}=(-1)^d=-1.             \label{eq:all-odd-cross-cup-8}
\end{equation}
Thus the mixed coefficient is again one.
Here \(L^{q-1}=1\), so
\begin{equation}
                              \chi(L^{q-1})=0.             \label{eq:all-odd-cross-cup-9}
\end{equation}

\subsection{The diagonal coefficients and Hilbert corrections}

In both cases \(\ker\chi\) is normal in \(E\), and
\begin{equation}
 1\to C_2\to E/\ker\chi\to C_m\times\mathbf F_q^\times\to1
                                                               \label{eq:all-odd-cross-cup-10}
\end{equation}
is a central extension.
The degree-two cohomology of the base has the three
basis classes \(\delta_\ell,\delta_q,c_\ell c_q\).
Equations \eqref{eq:all-odd-cross-cup-5} and \eqref{eq:all-odd-cross-cup-8} show that the mixed coefficient
is one. Equations \eqref{eq:all-odd-cross-cup-6} and \eqref{eq:all-odd-cross-cup-9} give the coefficient
\(\beta\) of \(\delta_q\):
\begin{equation}
 \beta=
 \begin{cases}
 r\bmod2,&d\text{ even},\\
 0,&d\text{ odd}.
 \end{cases}                                               \label{eq:all-odd-cross-cup-11}
\end{equation}

There is a uniform exact computation of the other
diagonal coefficient. Over an algebraic closure,
\(J\) permutes the \(m\) distinct deck eigenlines
in a single \(m\)-cycle.
Since \(J\) is symplectic, \(\det J=1\).
Since \(m\) is even, the sign of that permutation
is \(-1\). The product of the scalars around the
cycle is therefore \(-1\), and
\begin{equation}
                                J^m=-I.                   \label{eq:all-odd-cross-cup-12}
\end{equation}
In particular one must not replace \(J^m\) by the
identity before applying the central character.
The coefficient \(\alpha\) of \(\delta_\ell\) is
\begin{equation}
 \alpha=\chi(-I)=
 \begin{cases}
 r(q^s+1)/2\bmod2,&d=2s,\\
 r(q^d-1)/2\bmod2,&d\text{ odd}.
 \end{cases}                                               \label{eq:all-odd-cross-cup-13}
\end{equation}
Thus the class of \eqref{eq:all-odd-cross-cup-10} is exactly
\begin{equation}
                     c_\ell c_q+\alpha\delta_\ell+\beta\delta_q.
                                                               \label{eq:all-odd-cross-cup-14}
\end{equation}

Put \(b_\ell=(\ell-1)/2\bmod2\) and
\(b_q=(q-1)/2\bmod2\).
Equations \eqref{eq:all-odd-cross-cup-11}--\eqref{eq:all-odd-cross-cup-13} simplify to
\begin{equation}
 \beta=\frac{1-\left(\frac q\ell\right)}2,\qquad
                         \alpha=\beta+b_\ell b_q.          \label{eq:all-odd-cross-cup-15}
\end{equation}
For the first identity, if \(d\) is odd then \(q\)
is a square modulo \(\ell\).
If \(d=2s\), writing \(q=a^{rh}\) as above gives
\((q/\ell)=(-1)^r\), because \(h\) is odd.
For the second identity, in odd degree
\(b_\ell\equiv r\pmod2\), and
\((q^d-1)/2\equiv b_q\pmod2\).
In even degree, distinguish \(q\equiv1\) or \(3\pmod4\)
and use \(b_\ell\equiv rs\pmod2\).
This gives precisely \eqref{eq:all-odd-cross-cup-15}.

The Hilbert symbol of \((\ell^*,q^*)\) has no finite
ramification outside \(\ell,q\):
both discriminants are odd and one modulo four,
so the symbol at two is also one.
At \(q\), its sign is
\((\ell^*/q)=(q/\ell)\).
At \(\ell\), its sign is
\((q^*/\ell)=(-1)^{b_\ell b_q}(q/\ell)\).
These are the usual odd-prime norm-residue calculations
and quadratic reciprocity
\cite[III, Theorem~4.4, pp.~113--114; Remark~4.7, p.~115]{MilneCFT}.
Thus \(\alpha,\beta\) are exactly the two Hilbert
corrections in \eqref{eq:all-odd-cross-cup-1}.
This identifies the entire extension class directly;
arithmetic Brauer vanishing is not used to infer
vanishing on GT.

\subsection{The resulting characters and their arithmetic restrictions}

Composing \eqref{eq:all-odd-cross-cup-2} with \eqref{eq:all-odd-cross-cup-10} gives a homomorphism from
\(G_0\) lifting its finite cyclotomic quotient
\[
       g\longmapsto
        (\lambda_g\bmod\ell,\lambda_g\bmod q).
\]
Consequently the class \eqref{eq:all-odd-cross-cup-14}, hence \eqref{eq:all-odd-cross-cup-1}, inflates to
zero on \(G_0\).
For \(g\in M_0\), \(\rho_{\ell,q}(g)\in T\), and we set
\[
                       \xi_{\ell,q}(g)=
                                  \chi(\rho_{\ell,q}(g)).
\]
This is continuous, additive, and \(G_0\)-invariant.

Here are explicit formulas in the normalized beta
convention of \cite[\S2.3, p.~101]{Furusho2003}.
Choose a primitive \(\ell\)-th root \(\zeta\) in
\(\mathbf F_{q^d}\).
If \(d=2s\), choose representatives \(a_i\) for
\((\mathbf Z/\ell)^\times/\langle q\rangle\); then
\begin{equation}
 (-1)^{\xi_{\ell,q}(g)}
   =\prod_{i=1}^{r}
       B_g(\zeta^{a_i},\zeta^{a_i})^{(q^s+1)/2}.            \label{eq:all-odd-cross-cup-16}
\end{equation}
Each factor before exponentiation has norm one from
\(\mathbf F_{q^{2s}}\) to \(\mathbf F_{q^s}\).
If \(d\) is odd, choose representatives for
\((\mathbf Z/\ell)^\times/\langle q,-1\rangle\); then
\begin{equation}
 (-1)^{\xi_{\ell,q}(g)}
   =\prod_{i=1}^{r}
       B_g(\zeta^{a_i},\zeta^{a_i})^{(q^d-1)/2}.            \label{eq:all-odd-cross-cup-17}
\end{equation}
In this case the reciprocal deck character has the
inverse scalar.
Changing representatives or the primitive root only
permutes factors and applies Frobenius or inversion,
which preserve the displayed signs.
The commutator class generates each nontrivial
character line of the cyclic cover, so these scalars
are exactly the homology scalars used above.

Naturality of transgression gives \eqref{eq:all-odd-cross-cup-1} on the quotient
\(A=G/N\).
The arithmetic transgression on \(G_K\) is injective,
and its basis character \(\kappa_{\ell,q}\) has that
same class. Hence
\[
                       \xi_{\ell,q}|_{G_K}=\kappa_{\ell,q}.
\]
This completes the proof.

\begin{remark}[Comparing the two ordered constructions]
The characters \(\xi_{\ell,q}|_N\) and
\(\xi_{q,\ell}|_N\) have the same arithmetic restriction,
because their transgression is the symmetric class
\(e_{c_\ell,c_q}\).
Therefore
\[
          \varepsilon_{\ell,q}
             =\left.(\xi_{\ell,q}+\xi_{q,\ell})\right|_N
\]
is \(G\)-invariant and vanishes on \(N\cap H\).
The established factorization \(G=NH\) consequently
gives a unique continuous character \(G\to\mathbf F_2\)
which is \(\varepsilon_{\ell,q}\) on \(N\) and zero on
\(H\): define its value at \(nh\) to be
\(\varepsilon_{\ell,q}(n)\).
Agreement on the intersection makes this well-defined,
and invariance makes it a homomorphism.
This constructs explicit relative characters but does
not prove that any of them is nonzero.
\end{remark}

\section{Removing the entire cubic coefficient-two image}
\label{sec:fr-cubic-coefficient-two}

Use \(M=\ker\lambda\), \(D=N_{\mathrm{rad}}\), and
\[
 R_D=D\cap\ker\Phi_2,\qquad G=R_DH.
\]
The corrected cubic frame \(T\) is the homomorphism of
Lemma~\ref{lem:adelic-corrected-frames}. We distinguish its
\emph{coefficient-two component} from the maximal pro-two factor
of its whole adelic target. Put
\[
 O=\mathbf Z_2[\omega],\qquad
 \delta=1+2\omega,\qquad \delta^2=-3,
\]
\[
 \mathcal U=\{z\in1+4O:z\bar z=1\},\qquad
 \mathfrak t_2:M\longrightarrow\mathcal U,\quad
 \mathfrak t_2(g)=T_{g,2}.
\]
The corrected frame already belongs to \(1+4O\); no additional
projection away from roots of unity is needed in this definition.

\begin{lemma}\label{lem:fr-cubic-two-character}
The logarithm identifies \(\mathcal U\) with
\(4\delta\mathbf Z_2\). The formula
\[
 \tau(g)=\frac{\log\mathfrak t_2(g)}{4\delta}\pmod2
 \qquad(g\in M)
\]
defines a continuous homomorphism \(M\to\mathbf F_2\), invariant
under conjugation by all of \(G\). Equivalently,
\[
 \mathfrak t_2(g)\equiv(1+4\tau(g))I\pmod8.
\]
Its kernel in \(\mathcal U\) is \(\mathcal U^2\). An element of
\(\mathcal U\) with \(\tau=1\) topologically generates \(\mathcal U\).
\end{lemma}

\begin{proof}
Logarithm and exponential are inverse continuous homomorphisms
between \(1+4O\) and \(4O\). The norm condition becomes trace zero:
\(\log(z\bar z)=\operatorname{Tr}\log z\).
Since
\[
 \operatorname{Tr}(a+b\omega)=2a-b,
\]
the trace-zero lattice in \(O\) is precisely
\(\delta\mathbf Z_2\). Hence
\(\log\mathcal U=4\delta\mathbf Z_2\).
This proves the assertions concerning the logarithm, squares, and
topological generators.

Modulo eight, exponential is \(x\mapsto1+x\), and
\(\delta\equiv1\pmod2\). This gives the displayed scalar formula.
The frame is a homomorphism on \(M\), so \(\tau\) is a continuous
homomorphism. The integral conjugation formula
\eqref{eq:adelic-cubic-kernel-conjugation} identifies
\(T_{ghg^{-1},2}\), for \(h\in M\), with an integral similarity
of \(T_{h,2}\). Its reduction modulo eight fixes a scalar matrix.
Thus \(\tau(ghg^{-1})=\tau(h)\) for every \(g\in G\).
The same formula makes the kernel of \(\mathfrak t_2\) in \(M\)
normal in \(G\).
\end{proof}

\begin{lemma}\label{lem:fr-cubic-two-arithmetic-bit}
The restriction of \(\tau\) to \(D\cap H\) is nonzero. There is,
on the other hand, an element \(k_0\in D\cap H\) with
\[
                  \eta_3(k_0)=1,\qquad \tau(k_0)=0.
\]
\end{lemma}

\begin{proof}
Let
\[
 k=\mathbf Q^{\mathrm{ab}},\quad
 \alpha=3^{1/4},\quad
 v=(2+\sqrt3)^{1/4},\quad
 L_{\mathrm{rad}}=k(q^{1/n}:q\text{ prime},\ n\ge1).
\]
The established eight-division calculation
(Lemma~\ref{lem:lr-cubic-eight-field}) gives
\[
 k(E_3[8])=k(\alpha,v),\qquad [k(\alpha,v):k]=4,
\]
and identifies the action on \(v\) with the corrected scalar bit
\(\tau\). In particular \(v\in K_3^{\mathrm{ab}}\), and
\(v\notin k(\alpha)\). Also \(k(v)/k\) is quadratic, since
\[
 v^2=\frac{\sqrt6+\sqrt2}{2}\in k.
\]

We claim that \(v\notin L_{\mathrm{rad}}\).
The exact radical--CM intersection theorem gives
\[
 I=L_{\mathrm{rad}}\cap
       (K_3^{\mathrm{ab}}K_4^{\mathrm{ab}})
   =k(q^{1/12}:q\text{ prime},\,2^{1/8},3^{1/8}).
\]
Its explicit product Galois group shows that its largest elementary
abelian subextension of exponent two is
\[
                         k(q^{1/4}:q\text{ prime}).
\]
This follows coordinate by coordinate from the cyclic factors of
orders twelve and six; every continuous quadratic character uses
only finitely many coordinates. Each rational fourth root is abelian
over \(K_4=\mathbf Q(i)\), so this elementary subextension is contained
in \(K_4^{\mathrm{ab}}\).
If \(v\) belonged to \(L_{\mathrm{rad}}\), its quadratic field over
\(k\) would consequently lie in both CM fields. But
Lemma~\ref{lem:adelic-cm-entanglement} gives
\[
 K_3^{\mathrm{ab}}\cap K_4^{\mathrm{ab}}=k(\alpha),
\]
contrary to \(v\notin k(\alpha)\).
Thus \(L_{\mathrm{rad}}(v)/L_{\mathrm{rad}}\) is quadratic.
An element of \(G_{L_{\mathrm{rad}}}=D\cap H\) moving \(v\)
has \(\tau=1\).

For the second assertion, choose the established arithmetic element
\(\sigma\in M\cap H\) with
\[
 \rho_2^I(\sigma)=2,\qquad
 \rho_q^I(\sigma)=0\quad(q\text{ odd}).
\]
The cubic congruence
\(\rho_2^I\equiv-2m_3\pmod4\) makes \(m_3(\sigma)\) odd.
Choose \(h\in H\) with \(\lambda_2(h)=3\), and write
\(t=\lambda(h)\in\widehat{\mathbf Z}^{\times}\). Set
\[
                 k_0=(h\sigma h^{-1})\sigma^{-t}.
\]
The exponent here is the full profinite exponent \(t\), not the
integer three. Exact weight-one conjugation of every rational
Kummer cocycle gives \(\rho_q^I(k_0)=0\), hence \(k_0\in D\cap H\).
The cubic scalar law gives
\[
 m_3(k_0)=(3^3-3)m_3(\sigma)=24m_3(\sigma),
\]
so \(\eta_3(k_0)=1\).
By Lemma~\ref{lem:fr-cubic-two-character}, \(\tau\) is a
\(G\)-invariant homomorphism on all of \(M\). Since \(t_2\) is odd,
\[
 \tau(k_0)=\tau(\sigma)-(t_2\bmod2)\tau(\sigma)=0.
\]
\end{proof}

\begin{theorem}\label{thm:fr-cubic-two-exact-image}
The full and arithmetic coefficient-two images are
\[
 \boxed{\quad
 \mathfrak t_2(R_D)
   =\mathfrak t_2(R_D\cap H)
   =\mathcal U.
 \quad}
\]
In particular the upper image is the entire level-four norm-one
group, rather than just its square subgroup.
For
\[
                    R'=R_D\cap\ker\mathfrak t_2
\]
one has
\[
 R_D/R'\simeq\mathcal U\simeq\mathbf Z_2,\qquad
 G=R'H,\qquad
 B\simeq R'/\overline{[G,R'](R'\cap H)}.
\]
Thus \(R'\) is a proper closed normal subgroup of infinite index in
\(R_D\), and it preserves the full relative abelian obstruction.
\end{theorem}

\begin{proof}
Write \(J=D\cap H\), \(J_\tau=J\cap\ker\tau\), and
\(P_0=\Phi_2(D)\). The established dyadic normalization gives
\[
 \Phi_2(J)=P_0,\qquad
 H^1(P_0,\mathbf F_2)^G=\mathbf F_2\eta_3
\]
by Corollary~\ref{cor:protwo-image-invariant-characters}.
Let \(L=\Phi_2(J_\tau)\).
It is compact, hence closed, and normal in \(P_0\).
The quotient \(P_0/L\) is a quotient of \(J/J_\tau\); its order
is at most two.

Moreover \(L\) is stable under \(H\), since \(J_\tau\) is.
The action of \(G\) on \(P_0\) factors through its dyadic realization.
The factorization \(G=R_DH\), with \(\Phi_2(R_D)=1\), therefore
makes \(L\) stable under all \(G\).
If \(L\ne P_0\), the nonzero character of its quotient is
\(G\)-invariant and so equals \(\eta_3\).
But the element \(k_0\) of the preceding lemma belongs to
\(J_\tau\) and has \(\eta_3(k_0)=1\), a contradiction.
Consequently
\[
                         \Phi_2(J_\tau)=P_0.
\]

Choose \(r\in J\) with \(\tau(r)=1\).
There is \(y\in J_\tau\) with
\(\Phi_2(y)=\Phi_2(r)\). Then
\[
                         w=ry^{-1}\in R_D\cap H,
 \qquad \tau(w)=1.
\]
Its coefficient-two value topologically generates \(\mathcal U\)
by Lemma~\ref{lem:fr-cubic-two-character}.
The compact procyclic closure of \(w\) therefore maps onto
\(\mathcal U\). This proves arithmetic surjectivity; the full image
is contained in \(\mathcal U\) by definition, proving equality.

The kernel \(R'\) is closed and \(G\)-normal. For every \(r\in R_D\),
choose \(j\in R_D\cap H\) with
\(\mathfrak t_2(j)=\mathfrak t_2(r)\). Then \(rj^{-1}\in R'\).
Thus
\[
 R_D=R'(R_D\cap H),\qquad G=R'H.
\]
The compact quotient theorem gives \(R_D/R'\simeq\mathcal U\).

Finally, for any finite abelian group \(A\), a \(G\)-invariant
character of \(R'\) killing \(R'\cap H\) extends uniquely to
\(G\) by \(\chi(rh)=\chi(r)\). This is well-defined and
continuous by the compact quotient map \(R'\times H\to G\).
It is exactly a character of \(G\) trivial on \(H\), equivalently
a character of the established relative abelian obstruction \(B\).
Character separation over all finite abelian targets gives the
displayed canonical presentation of \(B\).
\end{proof}

\begin{corollary}\label{cor:fr-cubic-two-joint-gate}
Let
\[
 V_D=
 \bigl(D\cap\bigcap_p\ker\Phi_p\cap\ker d_{(3)}\bigr)
       \cap\ker W_{\mathrm{odd}}\cap\ker T_{\mathrm{odd}},
 \qquad
 V'=V_D\cap\ker\mathfrak t_2.
\]
Then
\[
 \mathfrak t_2(V_D)=\mathfrak t_2(V_D\cap H)=\mathcal U,
 \qquad [V_D:V']=[V_D\cap H:V'\cap H]=\infty,
\]
and
\[
 B_{\max2}\simeq
 \left(V'/\overline{[G,V'](V'\cap H)}\right)_{\max2}.
\]
If \(X_2=H^1(N,\mathbf F_2)^G\), its exact restriction kernels
to \(R'\) and to \(V'\) are both
\[
 \boxed{\quad
 \left(\bigoplus_{q\ \mathrm{odd}}\mathbf F_2\psi_q\right)
 \oplus\mathbf F_2\eta_3\oplus\mathbf F_2(\tau|_N).
 \quad}
\]
The last summand is independent of all preceding summands.
\end{corollary}

\begin{proof}
The established reductions show that \(R_D/V_D\) is pro-odd.
Explicitly, \(R_D/P_D\) is pro-odd; the successive cubic-defect,
odd harmonic, and odd cubic quotients are pro-odd as well.
Extensions of pro-odd groups are pro-odd.

For the element \(w\in R_D\cap H\) constructed above, take
\(w_2=w^{\varepsilon_2}\), where
\(\varepsilon_2\in\widehat{\mathbf Z}\) is one at two and
zero at all odd primes. Its image in the pro-odd quotient
\(R_D/V_D\) is one, so \(w_2\in V_D\cap H\).
Its image in the pro-two group \(\mathcal U\) is unchanged.
It therefore still topologically generates \(\mathcal U\), proving
the exact images and both infinite indices.

The quotient
\[
                         R'/V'\hookrightarrow R_D/V_D
\]
is pro-odd. Apply the relative pro-odd restriction lemma to the
already established factorization \(G=R'H\).
This gives the asserted presentation of \(B_{\max2}\).
The factorization \(G=R'H\) is essential to this application;
an internal arithmetic factorization inside \(V_D\) would not
by itself give the conclusion.

For the character kernel, suppose \(\chi\in X_2\) vanishes on
\(R'\). Its restriction to \(R_D\) factors through
\(R_D/R'\simeq\mathcal U\). The continuous quadratic character
space of \(\mathcal U\simeq\mathbf Z_2\) is generated by \(\tau\).
Subtracting either zero or \(\tau|_N\) therefore gives a character
vanishing on \(R_D\). The known kernel there is precisely the
Kummer span plus \(\eta_3\).
Conversely all displayed characters vanish on \(R'\).
The new summand is independent because \(\tau(w)=1\), whereas
all the older kernel characters vanish on \(R_D\).

If \(\chi\) vanishes on \(V'\), its restriction to \(R'\) factors
through the pro-odd quotient \(R'/V'\), hence is zero.
Thus restriction to \(V'\) has the same kernel.
\end{proof}

\subsection{The arithmetic class of the new quadratic lift}

We identify \(\tau|_{G_{K_2}}\) in the explicit arithmetic target
of Theorem~\ref{thm:fr-Y2-complete}.
Let \(a,b,c_3\) denote the quadratic cyclotomic characters of
\(\sqrt2,\sqrt{-1},\sqrt{-3}\), respectively, and let
\(\delta_2=b^2\), with \(\delta_3\) the cyclic lift obstruction
at the odd prime three. Products below are cup products.
The target normalization is
\[
 \operatorname{tr}(\kappa_{2,3})
  =\operatorname{Inf}(a c_3+\delta_2+\delta_3),\qquad
 \operatorname{tr}(\psi_3)
  =\operatorname{Inf}(b c_3+\delta_3).
\]

\begin{proposition}\label{prop:fr-cubic-two-Y2-class}
One has the exact equality of arithmetic characters
\[
             \tau|_{G_{K_2}}=\kappa_{2,3}+\psi_3.
\]
Consequently
\[
                \vartheta_{2,3}:=\tau|_N+\psi_3\in X_2
\]
is an explicit invariant lift of the corrected cross-cup target
\(\kappa_{2,3}\).
The last summand in Corollary~\ref{cor:fr-cubic-two-joint-gate}
can equivalently be written \(\mathbf F_2\vartheta_{2,3}\).
Neither character is asserted to be arithmetic-trivial.
\end{proposition}

\begin{proof}
Put \(L_v=\mathbf Q(i,v)\), where \(v=(2+\sqrt3)^{1/4}\).
The number
\[
 u=v^2=(\sqrt6+\sqrt2)/2
\]
has degree four over \(\mathbf Q\), since
\(u-u^{-1}=\sqrt2\) and \(u^2-2=\sqrt3\).
As \(v\notin k\), adjoining its square root to
\(\mathbf Q(u)\subset k\) has degree two. Thus
\([\mathbf Q(v):\mathbf Q]=8\), and this real field does not
contain \(i\), giving \([L_v:\mathbf Q]=16\).
All conjugates of \(v\) are \(i^jv^{\pm1}\), so \(L_v/\mathbf Q\)
is Galois.

Its automorphisms have the unique forms
\[
 g=(j,e,b):\qquad
 v\longmapsto i^jv^{(-1)^e},\quad
 i\longmapsto i^{(-1)^b},
 \qquad j\in\mathbf Z/4,\quad e,b\in\mathbf F_2.
\]
Their multiplication is
\[
 (j,e,b)(j',e',b')
   =\bigl((-1)^b j'+(-1)^{e'}j,\ e+e',\ b+b'\bigr).
\]
The element \(z=(2,0,0)\) is central and negates \(v\).
Its fixed field is
\[
                   L_0=\mathbf Q(i,\sqrt2,\sqrt3),
\]
of degree eight. The characters of this quotient on the three
square roots \(\sqrt2,\sqrt{-1},\sqrt{-3}\) have values
\[
                         A=j+e,\qquad b,\qquad C=b+e
                         \quad\text{in }\mathbf F_2.
\]
In particular \(e=b+C\), and \(j\bmod2=\alpha=A+b+C\).

Choose the section of this elementary abelian quotient with
\(j=\alpha\in\{0,1\}\).
The central carry, obtained from the multiplication law, is
\[
 \frac{(-1)^b\alpha'+(-1)^{e'}\alpha
                    -(\alpha+\alpha'\bmod2)}2
  =\alpha\alpha'+b\alpha'+\alpha e'\pmod2.
\]
Consequently the class of the central extension is
\[
                         \alpha^2+b\alpha+\alpha e
                  =a^2+b^2+a c_3+b c_3
\]
in the cohomology of the three-character quotient.
For an explicit check of the cup-order change, expansion gives
the cocycle \(AA'+CA'+bb'+bC'\); replacing \(CA'\) by \(AC'\)
changes it by the coboundary of the function \(AC\).

Pull this extension back to the full cyclotomic unit group
\(U=\widehat{\mathbf Z}^{\times}\).
The character \(a\) of \(\sqrt2\) lies on the procyclic factor
\(1+4\mathbf Z_2\) and lifts to a \(\mathbf Z/4\)-valued
character. Hence \(a^2=0\) in \(H^2(U,\mathbf F_2)\).
The pulled-back class is therefore
\[
                            b^2+a c_3+b c_3.
\]
The quadratic character of \(G_k\) defined by this central
extension is its action on \(v\), hence is \(\tau\) by the
eight-division comparison. Passing to \(G_{K_2}\) inflates the
same class from \(U\) to
\(\operatorname{Gal}(K_2/\mathbf Q)\).
The transgression used in the arithmetic-target theorem is
injective. Finally,
\[
 (a c_3+\delta_2+\delta_3)+(b c_3+\delta_3)
                          =a c_3+b c_3+b^2.
\]
This proves the claimed character identity, with no uncomputed
diagonal or cubic-edge correction.
The stated lift and the equivalent kernel expression follow.
\end{proof}

\paragraph{Scope of the normalization.}
The eliminated group is the coefficient-two level-four norm-one
torus \(\mathcal U\). The finite two-Sylow subgroups at odd coefficient
primes are not included in this assertion. The proof establishes
the full arithmetic matching on \(R_D\) before passing to \(V'\).
It neither identifies the remaining relative quotient with zero
nor asserts \(G=V'H\).

\section{Simultaneous removal of the lifted quadratic families}
\label{sec:simultaneous-quadratic-normalization}

The individual lifts can be normalized simultaneously, because their
arithmetic restrictions remain independent after the preceding
normalizations. We first prove the precise character comparison that
makes this assertion possible.

\begin{lemma}\label{lem:arithmetic-normalization-annihilator}
Let \(G=NH\) be a profinite group, where \(N\) is closed and normal
and \(H\) is closed. Put \(K=N\cap H\). Let \(R\subseteq N\)
be closed and normal in \(G\), with \(G=RH\), and put
\(J=R\cap H\). For any finite abelian group \(C\) with trivial
action, restriction induces an isomorphism
\[
 \ker\bigl(\operatorname{Hom}_{\rm cts}(N,C)^G
                 \to\operatorname{Hom}_{\rm cts}(R,C)\bigr)
 \xrightarrow{\ \sim\ }
 \ker\bigl(\operatorname{Hom}_{\rm cts}(K,C)^H
                 \to\operatorname{Hom}_{\rm cts}(J,C)\bigr).
\]
\end{lemma}
\begin{proof}
The factorization \(G=RH\) gives \(N=RK\): if \(n=rh\in N\)
and \(r\in R\subseteq N\), then \(h\in N\cap H=K\).
Consequently inclusion induces a topological isomorphism
\[
                         K/J\xrightarrow{\ \sim\ }N/R.
\]
It is continuous and bijective between compact Hausdorff groups.
An \(H\)-invariant character of \(K\) killing \(J\) therefore
defines a continuous character of \(N/R\). Conjugation by \(R\)
is trivial on this quotient, and \(G=RH\); thus \(H\)-invariance
is precisely the required \(G\)-invariance. Conversely every
\(G\)-invariant character of \(N\) killing \(R\) restricts to
one of the stated characters of \(K\). The quotient isomorphism
makes these constructions inverse.
\end{proof}

Return to \(G=\GT\), \(H=G_{\Q}\) and the preceding notation.
Let
\[
 \mathcal I=\{(\ell,q):\ell<q\text{ are odd primes}\},\qquad
 \mathcal P=\prod_{(\ell,q)\in\mathcal I}\F_2,
\]
and use the ordered lift \(\xi_{\ell,q}\) constructed from the
cyclic degree-\(\ell\) cover with coefficients modulo \(q\).
Write \(\mathfrak t_2\) for the entire coefficient-two cubic CM
homomorphism, with target
\[
 \mathcal U=\{z\in1+4\Z_2[\omega]:z\bar z=1\},
 \qquad \omega^2+\omega+1=0.
\]
Put \(\delta=1+2\omega\), and write
\[
 \tau=\frac{\log\mathfrak t_2}{4\delta}\pmod2,
 \qquad
 \Sigma_0=\F_2\eta_3\oplus
                 \bigoplus_{p\text{ odd}}\F_2\psi_p.
\]
Thus \(\tau\in X_2\), and its arithmetic restriction is
\(\kappa_{2,3}+\psi_3\), by
Proposition~\ref{prop:fr-cubic-two-Y2-class}.
Define the closed normal subgroups
\[
 \begin{split}
 R^*&=R_D\cap\ker\mathfrak t_2
                   \cap\bigcap_{(\ell,q)\in\mathcal I}\ker\xi_{\ell,q},\\
 V^*&=V_D\cap R^*.
 \end{split}
\]

\begin{theorem}\label{thm:simultaneous-quadratic-normalization}
The joint homomorphism
\[
 (\mathfrak t_2,(\xi_{\ell,q})_{\mathcal I})
                      :R_D\longrightarrow\mathcal U\times\mathcal P
\]
is surjective already on \(R_D\cap H\), and remains surjective
on \(V_D\cap H\). Consequently
\[
 R_D/R^*\simeq\mathcal U\times\mathcal P,
 \qquad V_D/V^*\simeq\mathcal U\times\mathcal P.
\]
Both quotient descriptions hold on the arithmetic intersections as
well. Moreover
\[
 G=R^*H,\qquad
 B\simeq R^*/\overline{[G,R^*](R^*\cap H)},
\]
and
\[
 B_{\max2}\simeq
 \left(V^*/\overline{[G,V^*](V^*\cap H)}\right)_{\max2}.
\]
Restriction from \(X_2\) to either \(R^*\) or \(V^*\) has kernel
\[
 \Sigma_0\oplus\F_2\tau\oplus
                 \bigoplus_{(\ell,q)\in\mathcal I}\F_2\xi_{\ell,q}.
\]
Arithmetic restriction is injective on this kernel. In particular
\(V^*\) is a proper subgroup of infinite index in \(V_D\),
already on the arithmetic intersection.
\end{theorem}
\begin{proof}
Put \(R'=R_D\cap\ker\mathfrak t_2\) and
\(V'=V_D\cap\ker\mathfrak t_2\). Theorem~\ref{thm:fr-cubic-two-exact-image}
and Corollary~\ref{cor:fr-cubic-two-joint-gate} give
\[
 G=R'H,\qquad R'/V'\text{ is pro-odd},\qquad
 \ker(X_2\to H^1(R',\F_2))=\Sigma_0\oplus\F_2\tau.
\]
The arithmetic restrictions of this last space are exactly the span
of \(\eta_3\), all \(\psi_p\), and \(\kappa_{2,3}\).
Theorem~\ref{thm:fr-Y2-complete} says that the characters
\(\kappa_{\ell,q}\), for \((\ell,q)\in\mathcal I\), are linearly
independent modulo that span. Applying
Lemma~\ref{lem:arithmetic-normalization-annihilator} with \(R=R'\)
shows that their restrictions to \(R'\cap H\) are linearly independent.
They are the arithmetic restrictions of the \(\xi_{\ell,q}\).

The image of the continuous map
\(R'\cap H\to\mathcal P\) is compact and hence closed. If it
were proper, some finite-coordinate quotient would also have proper
image: a point outside a closed subset has a basic open neighbourhood
disjoint from it. A proper vector subspace of that finite product
has a nonzero linear annihilator. This would give a nonzero finite
linear combination of the \(\kappa_{\ell,q}\) vanishing on
\(R'\cap H\), a contradiction. Therefore this map is onto.

The same independence holds on \(V'\cap H\). Indeed,
\[
 (R'\cap H)/(V'\cap H)\hookrightarrow R'/V'
\]
is pro-odd. A quadratic character vanishing on \(V'\cap H\)
would factor through this pro-odd quotient and so would vanish on
\(R'\cap H\). The preceding compact-image argument now gives
surjectivity \(V'\cap H\to\mathcal P\).

The coefficient-two theorem gives
\(\mathfrak t_2(R_D\cap H)=\mathfrak t_2(V_D\cap H)=\mathcal U\).
Given \((u,e)\in\mathcal U\times\mathcal P\), first choose an
arithmetic element with coefficient-two value \(u\). Multiplication
by an element of \(R'\cap H\), or of \(V'\cap H\), prescribes
its whole \(\mathcal P\)-coordinate without changing \(u\).
This proves both joint surjectivity assertions, including their full
group versions. Their kernels give the asserted quotient isomorphisms.

For \(r\in R_D\), choose \(h\in R_D\cap H\) with the same joint
image. Then \(rh^{-1}\in R^*\), so
\(R_D=R^*(R_D\cap H)\) and \(G=R^*H\).
The invariant-character extension argument for a closed normal
factor multiplying with \(H\) now gives the stated full relative
abelian quotient. More explicitly, for every finite abelian group
\(C\), a \(G\)-invariant homomorphism \(R^*\to C\) vanishing
on \(R^*\cap H\) extends uniquely to \(G\) by the rule
\(rh\mapsto\chi(r)\); compactness gives continuity. Taking all
finite targets identifies the profinite relative abelian quotients.

The natural injection
\(R^*/V^*\hookrightarrow R_D/V_D\) has pro-odd image.
Lemma~\ref{lem:relative-proodd-restriction}, applied with the proved
factorization \(G=R^*H\), identifies the maximal pro-two relative
quotient on \(V^*\).

Finally, a character in \(X_2\) vanishing on \(R^*\) restricts
to a continuous quadratic character of
\(R_D/R^*\simeq\mathcal U\times\mathcal P\).
Since \(\mathcal U\simeq\Z_2\), its only nonzero quadratic
character is \(\tau\). A continuous character of \(\mathcal P\)
depends on finitely many coordinates. Subtracting the corresponding
finite sum of \(\tau\) and the \(\xi_{\ell,q}\) leaves a character
vanishing on \(R_D\), hence one in \(\Sigma_0\).
Conversely all these characters vanish on \(R^*\).
The sum is direct, and has injective arithmetic restriction, by the
explicit basis of \(Y_2\). A character vanishing on \(V^*\)
vanishes on \(R^*\), since \(R^*/V^*\) is pro-odd. This proves
the identical kernel on \(V^*\) and completes the proof.
\end{proof}

This removes the cubic coefficient-two image and all odd-pair
characters simultaneously. The next construction lifts the remaining
dyadic-pair arithmetic characters; its application gives a further
normalization of the relative kernel.

\section{Lifting the remaining cross-cup axes at coefficient two}
\label{sec:coefficient-two-cross-cups}

Write \(G=\widehat{\mathrm{GT}}\), \(H=G_{\mathbf Q}\),
\(M=\ker\lambda\), and \(N=\ker(\lambda,\rho_2^I)\).
The arithmetic field fixed by \(N\cap H\) is
\[
 K_2=\mathbf Q(\mu_\infty,2^{1/\infty}),\qquad
 A=\operatorname{Gal}(K_2/\mathbf Q)=G/N.
\]
For an odd prime \(p\), let \(c_p\) be the quadratic
character of \(p^*=(-1)^{(p-1)/2}p\), and let
\(\delta_p\) be the degree-two generator of the cyclic
two-primary quotient of \(\mathbf Z_p^\times\), inflated
to the cyclotomic group and then to \(A\). Put
\[
 \begin{gathered}
 b(u)=\frac{u_2-1}{2}\pmod2,\qquad
 a(u)=\frac{u_2^2-1}{8}\pmod2,\\
 \delta_2=b^2,\qquad A_p=\frac{p^2-1}{8}\pmod2.
 \end{gathered}
\]
All cohomology in this section has trivial coefficient
module \(\mathbf F_2\).

\begin{theorem}\label{thm:coefficient-two-cross-cups}
For every odd prime \(p\), there is a continuous
\(G\)-invariant homomorphism
\[
                     \xi_{2,p}:M\longrightarrow\mathbf F_2
\]
whose restriction to \(N\) transgresses to
\begin{equation}
 e_{a,c_p}=a c_p+A_p(\delta_2+\delta_p)
                        \quad\hbox{in }H^2(A,\mathbf F_2).
                                                        \label{eq:coefficient-two-cross-1}
\end{equation}
In particular \eqref{eq:coefficient-two-cross-1} has zero inflation to \(G\), and
\begin{equation}
                       \xi_{2,p}|_{G_{K_2}}=\kappa_{2,p},
                                                        \label{eq:coefficient-two-cross-2}
\end{equation}
where \(\kappa_{2,p}\) is the unique arithmetic target
character transgressing to \eqref{eq:coefficient-two-cross-1}.

The character is obtained from the action on compact
homology, modulo eight, of the degree-\(p\) cyclic cover
of the three-punctured sphere with monodromies
\(1,1,-2\). The same construction is defined and invariant
on the multiplier-one subgroup of any group of continuous
positive peripheral automorphisms of \(\widehat F(x,y)\).
The transgression assertion above concerns specifically
the full quotient \(G/N=A\).
\end{theorem}

The proof also constructs a companion character whose
arithmetic restriction is \(\psi_p\). Its correction
terms provide a useful check on the normalization.

\subsection{Integral homology and the multiplier}

Set
\[
 U_p=\ker\bigl(\widehat F(x,y)\longrightarrow C_p,
                         \quad x,y\longmapsto1\bigr),
 \qquad m=p-1.
\]
The cover has three punctures and compact genus
\((p-1)/2\). Let \(V\) be its first homology over
\(\mathbf Z_2\), modulo the peripheral homology submodule.
It is free of rank \(m\), with its unimodular alternating
intersection form. The deck generator \(D\) has each
nontrivial \(p\)-th-root character once.
Indeed, on nontrivial character summands the cellular
boundary map
\[
 \mathbf Z_2[C_p]^2\longrightarrow\mathbf Z_2[C_p],
          \qquad (r,s)\longmapsto(D-1)(r+s)
\]
has a rank-one kernel. Its two trivial summands are
exactly the peripheral homology. The latter assertion
is integral: \(p\) is a unit in \(\mathbf Z_2\).

A positive peripheral automorphism with common multiplier
\(\lambda\) preserves \(U_p\) and gives an integral
similitude
\begin{equation}
 \rho_{p,2}(g)\in\operatorname{GSp}(V),\qquad
 \rho_{p,2}(g)D\rho_{p,2}(g)^{-1}
       =D^{\lambda_g\bmod p},\qquad
 \operatorname{sim}(\rho_{p,2}(g))=\lambda_{g,2}.
                                                        \label{eq:coefficient-two-cross-3}
\end{equation}
We verify the multiplier at two without passing to an
exponent-two group. A surface presentation is
\begin{equation}
 \prod_{j=1}^{(p-1)/2}[a_j,b_j]c_1c_2c_3=1.
                                                        \label{eq:coefficient-two-cross-4}
\end{equation}
There is one puncture above each original puncture.
The conjugator of a peripheral power can therefore be
moved into \(U_p\) by multiplication by a power of the
corresponding original peripheral generator. Consequently
the automorphism sends each \(c_i\) to a
\(U_p\)-conjugate of \(c_i^\lambda\).

Take the torsion-free class-two pro-two quotient of the
free group on \(a_j,b_j,c_1,c_2\). Its collected
commutator coordinates form a free \(\mathbf Z_2\)-module.
After tensoring with \(\mathbf Q_2\), use the class-two
logarithm and project degree-two coordinates to
\(\bigwedge^2 V\). The logarithms of \(c_1,c_2\), and
their images, have zero projected central coordinate:
conjugation adds a wedge containing a peripheral class.
Equation \eqref{eq:coefficient-two-cross-4} gives projected central coordinate
\(-\sum_j a_j\wedge b_j\) for \(\log c_3\).
The image of \(c_3\) multiplies that coordinate by
\(\lambda_2\). Thus
\[
     (\bigwedge^2\rho_{p,2}(g))
          \left(\sum_j a_j\wedge b_j\right)
       =\lambda_{g,2}\sum_j a_j\wedge b_j.
\]
Both sides are integral. This is the alternating
similitude identity, proving \eqref{eq:coefficient-two-cross-3}. This argument uses
only the stated peripheral action and the presentation
of the cover.

Let \(E\) be the normalizer of \(\langle D\rangle\)
in \(\operatorname{GSp}(V)\). Its quotient map records
deck exponent and similitude multiplier:
\begin{equation}
 \pi:E\longrightarrow Q_p,
       \qquad Q_p=C_m\times\mathbf Z_2^\times.
                                                        \label{eq:coefficient-two-cross-5}
\end{equation}
We construct two \(E\)-invariant characters of
\(T=\ker\pi\), denoted \(\chi_a,\chi_b\), and compute
their central extensions exactly.

\subsection{A norm-one coinvariant calculation}

\begin{lemma}\label{lem:two-norm-one-coinvariants}
Let \(L/\mathbf Q_2\) be unramified of degree \(2s\),
with integer ring \(O\), arithmetic Frobenius \(\sigma\),
and involution \(\bar z=\sigma^s(z)\). Let \(O_0\)
be the fixed integer ring, and put
\[
             T_s=\{z\in O^\times:z\bar z=1\}.
\]
There is a canonical isomorphism of finite groups
\begin{equation}
 \mathcal D_s:
 \mathbf Z_2^\times/(\mathbf Z_2^\times)^2
       \xrightarrow{\ \sim\ } T_s/(\sigma-1)T_s,
 \qquad
 t\longmapsto\left[\frac{\sigma(v)}v\right],
 \qquad v\bar v=t.
                                                        \label{eq:coefficient-two-cross-6}
\end{equation}
Here \((\sigma-1)T_s=\{\sigma(z)/z:z\in T_s\}\).
Define
\[
        \alpha_s=a\circ\mathcal D_s^{-1},\qquad
        \beta_s=b\circ\mathcal D_s^{-1}.
\]
For \(z\in T_s\cap(1+4O)\),
\begin{equation}
 \alpha_s(z)=
 \operatorname{Tr}_{\mathbf F_{2^s}/\mathbf F_2}
       \left(\frac{\log z}{4}\bmod2\right),
 \qquad \beta_s(z)=0.                                   \label{eq:coefficient-two-cross-7}
\end{equation}
Moreover
\begin{equation}
                  \alpha_s(-1)=1,
       \qquad \beta_s(-1)=s\pmod2.                       \label{eq:coefficient-two-cross-8}
\end{equation}
If \(v\bar v=-1\), then
\begin{equation}
                   \alpha_s(v^2)=1,
              \qquad \beta_s(v^2)=0.                    \label{eq:coefficient-two-cross-9}
\end{equation}
In \eqref{eq:coefficient-two-cross-9}, the square belongs to \(T_s\); the values are
independent of the choice of \(v\).
\end{lemma}

\begin{proof}
The norm on the units of an unramified extension is
surjective, and cyclic norm-one elements are Hilbert
90 quotients; see \cite[II, Corollary~1.23, p.~67;
III, Proposition~1.2, p.~98]{MilneCFT}.
For completeness, surjectivity on units follows by
first lifting the surjective residue-field norm and
then successively correcting each principal-unit
level using the surjective residue-field trace.

Choose \(v\) with \(v\bar v=t\). Replacing \(v\) by
\(vz\), \(z\in T_s\), changes its quotient in \eqref{eq:coefficient-two-cross-6}
by \(\sigma(z)/z\). Products of choices show that \eqref{eq:coefficient-two-cross-6}
is a homomorphism. A square \(t=w^2\), with
\(w\in\mathbf Z_2^\times\), has the choice \(v=w\),
so squares map to one.

Conversely, let \(z\in T_s\). Its full norm to
\(\mathbf Q_2\) is one. Hilbert 90 gives
\(z=\sigma(v)/v\) with \(v\in L^\times\). Multiplying
by a rational power of two makes \(v\) a unit because
the extension is unramified. The relative norm
\(v\bar v\) is fixed by \(\sigma\), hence lies in
\(\mathbf Z_2^\times\). This proves surjectivity.
If \(\sigma(v)/v=\sigma(z)/z\) for \(z\in T_s\),
then \(v/z\) is a rational unit, so \(v\bar v\) is
a square. This proves injectivity. The image of
\(\sigma-1\) is closed by compactness; the quotient
and \eqref{eq:coefficient-two-cross-6} are continuous.

To prove \eqref{eq:coefficient-two-cross-7}, choose \(v\) with
\(z=\sigma(v)/v\). Since \(z\equiv1\pmod4\), the
class of \(v\) modulo four is fixed by Frobenius.
Indeed its residue lies in \(\mathbf F_2^\times=\{1\}\);
write \(v=1+2y\pmod4\). Frobenius invariance then gives
\(y^2=y\pmod2\), so \(v\equiv1\) or \(3\pmod4\). Scaling by a rational
unit therefore makes \(v\equiv1\pmod4\).
Write \(v=\exp(4x)\). Then
\[
 t=v\bar v=\exp(4(x+\bar x))\in1+4\mathbf Z_2,
 \qquad \frac{\log z}{4}=\sigma x-x.
\]
Thus \(b(t)=0\) and \(a(t)=x+\bar x\pmod2\).
The reduction of \(\sigma x-x\) is \(x^2+x\), and
it belongs to \(\mathbf F_{2^s}\) because its relative
trace is zero. Its absolute trace telescopes:
\[
 \operatorname{Tr}_{\mathbf F_{2^s}/\mathbf F_2}(x^2+x)
                       =x^{2^s}+x=\bar x+x.
\]
This proves \eqref{eq:coefficient-two-cross-7}.

For \eqref{eq:coefficient-two-cross-8}, the unramified quadratic subfield contains
\(\delta=1+2\omega\), where \(\omega^2+\omega+1=0\).
It satisfies \(\delta^2=-3\) and
\(\sigma\delta=-\delta\). Its relative norm is
\(3\) for odd \(s\), and \(-3\) for even \(s\).
Hence \(\mathcal D_s(3)=[-1]\) or
\(\mathcal D_s(-3)=[-1]\), respectively.
Now \(a(3)=a(-3)=1\), while
\(b(3)=1\) and \(b(-3)=0\).

For \eqref{eq:coefficient-two-cross-9}, remove the odd-order Teichm\"uller factor of
\(v\). Its relative norm is one, so this operation
preserves the norm \(-1\); its square has zero value
under either quadratic character. We may thus take
\(v=1+2x\). The norm equation gives
\[
                    x+\bar x=1\pmod2,
       \qquad v^2=1+4(x+x^2)\in T_s\cap(1+4O).
\]
Formula \eqref{eq:coefficient-two-cross-7} and the same telescoping trace give
\(\alpha_s(v^2)=1\) and \(\beta_s(v^2)=0\).
Replacing \(v\) by another norm-minus-one unit
multiplies \(v^2\) by a square in \(T_s\), so the
values are independent of the choice.
\end{proof}

There is also a direct unit formula for the inverse square class.
If \(z=w/\bar w\in T_s\), then
\begin{equation}\label{eq:two-hilbert-norm-squareclass}
 \mathcal D_s^{-1}([z])=[N_{L/\mathbf Q_2}(w)].
\end{equation}
Changing \(w\) by a unit of the fixed field changes its absolute
norm by a rational square, so the right side is well-defined.
To prove the formula, write \(z=\sigma(v)/v\) and
\(v\bar v=t\in\mathbf Z_2^\times\), as above. Put
\[
 P=\prod_{j=0}^{s-1}\sigma^j(w),\qquad Z=vP.
\]
Then \(\sigma(Z)/Z=(\sigma(v)/v)(\bar w/w)=1\), so
\(Z\) is a rational unit. Therefore
\[
 N_{L/\mathbf Q_2}(w)=P\bar P=Z^2/t.
\]
The square classes of \(t^{-1}\) and \(t\) coincide, proving
\eqref{eq:two-hilbert-norm-squareclass}. In particular
\[
 \alpha_s(z)=a(N_{L/\mathbf Q_2}w),\qquad
 \beta_s(z)=b(N_{L/\mathbf Q_2}w).
\]
These formulas agree on the entire torus, including its first
principal-unit layer.

\subsection{Even Frobenius order}

Suppose \(d=\operatorname{ord}_p(2)=2s\), and let
\(r=m/d\). Every Frobenius orbit is self-dual. With
\(O\) as in Lemma~\ref{lem:two-norm-one-coinvariants},
we may write \(V=O^r\). On each summand the form is
\[
       (z,w)\longmapsto
       \operatorname{Tr}_{L/\mathbf Q_2}(d_i z\bar w),
       \qquad \bar d_i=-d_i,\quad d_i\in O^\times.
\]
The coefficients are units by unimodularity. Their
ratios lie in \(O_0^\times\); surjectivity of the
relative norm permits a rescaling making all of
them equal. Consequently
\begin{equation}
 T=T_s^r,
 \qquad \chi_a(z_1,\ldots,z_r)=\sum_i\alpha_s(z_i),
 \qquad \chi_b(z_1,\ldots,z_r)=\sum_i\beta_s(z_i).
                                                        \label{eq:coefficient-two-cross-10}
\end{equation}

Fix a primitive residue \(u\bmod p\). Choose the
Frobenius-orbit representatives
\(1,u,\ldots,u^{r-1}\). A lift of deck exponent
\(u\) cycles these summands, with Frobenius twist
\(\sigma^k\) on the last passage, where
\[
                  u^r=2^k\pmod p,
              \qquad \gcd(k,d)=1.
\]
In particular \(k\) is odd. The ratio between the
form coefficients before and after this last twist
lies in \(O_0^\times\); multiplying that summand by
a unit of prescribed relative norm restores the
form. This gives an integral symplectic lift \(J\).

For every \(t\in\mathbf Z_2^\times\), choose a unit
\(v\in O^\times\) with \(v\bar v=t\). Uniform scalar
multiplication by \(v\) gives a lift \(L(t)\) of
\((1,t)\) in \eqref{eq:coefficient-two-cross-5}. Thus \eqref{eq:coefficient-two-cross-5} is onto. These scalar
lifts commute with one another and with \(T\).
Conjugation by \(J\) only permutes factors and applies
Frobenius powers. Formula \eqref{eq:coefficient-two-cross-6} makes both characters
in \eqref{eq:coefficient-two-cross-10} invariant. Hence both are invariant under
all of \(E\).

The scalar correction used to define \(J\) commutes
with \(L(t)\). Their commutator has one nontrivial
component, \(\sigma^k(v)/v\), or its inverse according
to the choice of permutation convention. In the
coinvariant quotient of \eqref{eq:coefficient-two-cross-6}, this component is
\(\mathcal D_s(t)^k=\mathcal D_s(t)\). Therefore
\begin{equation}
             \chi_a([J,L(t)])=a(t),\qquad
             \chi_b([J,L(t)])=b(t).                       \label{eq:coefficient-two-cross-11}
\end{equation}

Over an algebraic closure, \(J\) cycles all \(m\)
deck eigenlines in one cycle. The product of its
scalar entries is \(-1\): the cycle has sign \(-1\)
and \(\det J=1\). Thus, exactly,
\begin{equation}
                             J^m=-I.                    \label{eq:coefficient-two-cross-12}
\end{equation}
Using \eqref{eq:coefficient-two-cross-8},
\begin{equation}
                    \chi_a(J^m)=r\pmod2,
           \qquad \chi_b(J^m)=rs\pmod2.
                                                        \label{eq:coefficient-two-cross-13}
\end{equation}
For the cyclic subgroup generated by the multiplier
\(-1\), \eqref{eq:coefficient-two-cross-9} gives
\begin{equation}
                   \chi_a(L(-1)^2)=r\pmod2,
              \qquad\chi_b(L(-1)^2)=0.                  \label{eq:coefficient-two-cross-14}
\end{equation}
In particular the square of the lift of \(-1\) cannot
be discarded in this calculation.

\subsection{Odd Frobenius order}

Suppose \(d=\operatorname{ord}_p(2)\) is odd, and let
\(r=m/(2d)\). Let \(O\) be the integer ring in the
unramified extension of \(\mathbf Q_2\) of degree
\(d\). Frobenius orbits occur in distinct dual pairs.
The lattice is a sum of \(r\) hyperbolic \(O\)-pairs,
with form
\[
 \sum_i\operatorname{Tr}_{O/\mathbf Z_2}
                         (x_i y_i'-y_i x_i').
\]
Its torus and two characters are
\begin{equation}
 \begin{split}
 T&=(O^\times)^r,
       \quad z_i:(x_i,y_i)\longmapsto(z_i x_i,z_i^{-1}y_i),\\
 \chi_a(z)&=\sum_i a(N_{O/\mathbf Z_2}z_i),
       \qquad
 \chi_b(z)=\sum_i b(N_{O/\mathbf Z_2}z_i).
 \end{split}                                             \label{eq:coefficient-two-cross-15}
\end{equation}
They are invariant under Frobenius, inversion, and
permutation of pairs.

For a primitive residue \(u\), the quotient by
\(\langle2,-1\rangle\) has order \(r\), and
\(u^r=-2^k\) for some \(k\). A symplectic lift \(J\)
cycles the \(r\) pairs and on the final passage
exchanges reciprocal coordinates with a Frobenius
twist and one minus sign. A multiplicative similitude
section is
\[
           L(t):(x_i,y_i)\longmapsto(x_i,t y_i)
              \qquad(t\in\mathbf Z_2^\times).
\]
This proves surjectivity of \eqref{eq:coefficient-two-cross-5} and invariance of \eqref{eq:coefficient-two-cross-15}.
The commutator of \(J\) with \(L(t)\) has torus
parameter \(t\) or \(t^{-1}\) on one pair and one
elsewhere. Its norm is \(t^{\pm d}\). Since \(d\)
is odd, \eqref{eq:coefficient-two-cross-11} holds again. The eigenline-cycle proof
of \eqref{eq:coefficient-two-cross-12} also applies, and now
\begin{equation}
               \chi_a(J^m)=0,\qquad \chi_b(J^m)=r\pmod2.
                                                        \label{eq:coefficient-two-cross-16}
\end{equation}
The multiplicative section gives
\begin{equation}
                     \chi_a(L(-1)^2)
                         =\chi_b(L(-1)^2)=0.              \label{eq:coefficient-two-cross-17}
\end{equation}

\subsection{All central-extension coefficients}

For \(\chi=\chi_a\) or \(\chi_b\), the quotient
\begin{equation}
              1\longrightarrow C_2\longrightarrow
                 E/\ker\chi\longrightarrow Q_p
                    \longrightarrow1                   \label{eq:coefficient-two-cross-18}
\end{equation}
is a continuous central extension. Both characters
are nonzero: this follows from \eqref{eq:coefficient-two-cross-6} in even degree
and surjectivity of the unramified norm in odd degree.

The degree-two cohomology of
\(Q_p=C_m\times\mathbf Z_2^\times\), with
\(\mathbf Z_2^\times=\langle-1\rangle\times
\overline{\langle5\rangle}\), has basis
\begin{equation}
                 \delta_p,\quad b^2,\quad ba,
                         \quad b c_p,\quad a c_p.
                                                        \label{eq:coefficient-two-cross-19}
\end{equation}
This follows from the cyclic cohomology calculation
and the product formula; the procyclic factor
\(\overline{\langle5\rangle}\simeq\mathbf Z_2\)
has cohomological dimension one.

Equation \eqref{eq:coefficient-two-cross-11} gives mixed term \(a c_p\) for
\(\chi_a\) and \(b c_p\) for \(\chi_b\).
The preimage in \(E\) of the multiplier factor
\(\mathbf Z_2^\times\) is abelian in both descriptions,
so the coefficient of \(ba\) is zero.
Equations \eqref{eq:coefficient-two-cross-13}--\eqref{eq:coefficient-two-cross-17} give every other coefficient:
\begin{equation}
 \begin{array}{c|cc|cc}
 &\multicolumn{2}{c|}{\chi_a}
 &\multicolumn{2}{c}{\chi_b}\\
 &\delta_p&b^2&\delta_p&b^2\\ \hline
 d=2s &r&r&rs&0\\
 d\text{ odd}&0&0&r&0
 \end{array}                                             \label{eq:coefficient-two-cross-20}
\end{equation}
All entries are reduced modulo two. In both cases
the \(\chi_b\) entry at \(\delta_p\) is
\((p-1)/2\pmod2\).

The \(\chi_a\) entries equal \(A_p\). Indeed, if
\(d\) is odd, two is a square modulo \(p\), so
\(A_p=0\). If \(d\) is even, write
\(2=u^{rh}\), with \(\gcd(h,d)=1\); then \(h\)
is odd and the Legendre symbol \((2/p)=(-1)^r\).
The quadratic reciprocity supplement identifies
this sign with \((-1)^{A_p}\).
Consequently the exact extension classes are
\begin{equation}
 \begin{split}
 [E/\ker\chi_a]&=a c_p+A_p(b^2+\delta_p),\\
 [E/\ker\chi_b]&=b c_p+\frac{p-1}{2}\delta_p.
 \end{split}                                             \label{eq:coefficient-two-cross-21}
\end{equation}

These are exactly the arithmetic Hilbert corrections.
For \((2,p^*)\), the only possible finite ramification
places are \(2,p\), and both Hilbert signs are
\((-1)^{A_p}\). For \((-1,p^*)\), the only finite
sign is at \(p\), with exponent \((p-1)/2\).
The dyadic assertion uses \(p^*\equiv1\pmod4\).
These also follow directly from the explicit dyadic
Hilbert formula in \cite[VIII, \S5, p.~248]{MilneCFT}.
No coefficient in \eqref{eq:coefficient-two-cross-21} is inferred merely from
arithmetic Brauer vanishing.

\subsection{Characters, transgression, and the full target}

For \(g\in M\), \eqref{eq:coefficient-two-cross-3} puts \(\rho_{p,2}(g)\) in
\(T\). Define
\begin{equation}
               \xi_{2,p}(g)=\chi_a(\rho_{p,2}(g)),
       \qquad \widetilde\psi_p(g)=\chi_b(\rho_{p,2}(g)).
                                                        \label{eq:coefficient-two-cross-22}
\end{equation}
They are homomorphisms, and invariance of the torus
characters makes them invariant under conjugation
by every positive peripheral automorphism in the
ambient group. They are continuous. More precisely,
\eqref{eq:coefficient-two-cross-7} kills \(T_s\cap(1+8O)\) for \(\alpha_s\) and
\(T_s\cap(1+4O)\) for \(\beta_s\). In odd degree
the same assertions follow from the norm formulas.
Thus the two characters in \eqref{eq:coefficient-two-cross-22} depend only on
compact homology modulo eight and four, respectively.

The representation \(G\to E\), followed by
either quotient in \eqref{eq:coefficient-two-cross-18}, lifts the cyclotomic map
\(G\to Q_p\). Both classes in \eqref{eq:coefficient-two-cross-21} therefore inflate
to zero on \(G\). Naturality of the five-term
sequence for \(1\to N\to G\to A\to1\) says that
the transgressions of the restrictions of \eqref{eq:coefficient-two-cross-22} to
\(N\) are exactly the pullbacks of \eqref{eq:coefficient-two-cross-21} to \(A\).
There is no additional affine component: this
statement comes from the displayed morphism of
extensions and its torus character.

On the arithmetic subgroup, the transgression
\[
 H^1(G_{K_2},\mathbf F_2)^H\longrightarrow H^2(A,\mathbf F_2)
\]
is injective, as in the complete arithmetic target
theorem. Thus \eqref{eq:coefficient-two-cross-21} identifies the restrictions as
\[
           \xi_{2,p}|_{G_{K_2}}=\kappa_{2,p},
       \qquad \widetilde\psi_p|_{G_{K_2}}=\psi_p|_{G_{K_2}}.
\]
This proves Theorem~\ref{thm:coefficient-two-cross-cups}.

\begin{corollary}\label{cor:all-arithmetic-quadratic-targets-lift}
The arithmetic restriction map is surjective:
\[
           X_2=H^1(N,\mathbf F_2)^G
              \longrightarrow Y_2=H^1(G_{K_2},\mathbf F_2)^H.
\]
Consequently the established structural sequence becomes
\begin{equation}
 0\longrightarrow\operatorname{Hom}_{\mathrm{cont}}
                         (B,\mathbf F_2)
   \longrightarrow X_2\longrightarrow Y_2\longrightarrow0.
                                                        \label{eq:coefficient-two-cross-23}
\end{equation}
This does not determine the left-hand term.
\end{corollary}

\begin{proof}
The complete target basis consists of \(\eta_3|_{G_{K_2}}\),
the \(\psi_p|_{G_{K_2}}\), the odd-pair characters
\(\kappa_{\ell,q}\), and the characters
\(\kappa_{2,p}\). The first two kinds already lift,
the odd-pair construction lifts the third kind, and
Theorem~\ref{thm:coefficient-two-cross-cups} lifts the
last kind. Every target character is a finite sum of
these basis elements. The kernel identification in
\eqref{eq:coefficient-two-cross-23} is the established relative character theorem
for \(G=NH\).
\end{proof}

\begin{remark}
The construction does not assert an equality on all
of \(N\) between distinct lifts of the same arithmetic
character. For \(p=3\), the earlier cubic frame
character \(\tau+\psi_3\) and \(\xi_{2,3}\) have the
same arithmetic restriction. Their difference, if
considered, is a character of the relative term in
\eqref{eq:coefficient-two-cross-23}; no vanishing or nonvanishing of that difference
is required here. Likewise, the companion character
\(\widetilde\psi_p\) is identified with \(\psi_p\)
arithmetically, without assuming their equality
on the full group.
\end{remark}

\section{The companion character is the divided Kummer character}
\label{sec:companion-kummer-equality}

Use the integral compact-homology representation of the degree-\(p\)
cyclic cover and its invariant character \(\widetilde\psi_p\) obtained
from \(b=\epsilon_{-1}\). The possible difference between this character
and the divided Kummer character can be eliminated on full GT.

\begin{theorem}\label{thm:companion-kummer-equality}
For every odd prime \(p\), on the entire multiplier-one subgroup
\(M\) of full profinite GT one has
\begin{equation}
             \boxed{\quad
             \widetilde\psi_p(g)=\frac{\rho_p^I(g)}2\pmod2.
             \quad}                                      
\label{eq:companion-kummer-24}
\end{equation}
More generally, the equality holds for standard positive peripheral
automorphisms \(x\mapsto x\), \(y\mapsto f^{-1}yf\), with
\(f\in F'\), satisfying the two-cycle relation and the integral
Anderson duplication identity. No duplication identity is assumed
for a larger gentle group unless it has been separately established.
\end{theorem}

Use the commutator convention \([x,y]=xyx^{-1}y^{-1}\), and write
\begin{equation}
 [f]=R_g(X,Y)[x,y]\quad\hbox{in }F'/F'',\qquad
 B_g(X,Y)=1+(X-1)(Y-1)R_g(X,Y).
\label{eq:companion-kummer-25}
\end{equation}
Write \(\psi_p=\rho_p^I/2\pmod2\) on \(M\), extending the
earlier notation on \(N\). The coefficient variables record conjugation by \(x,y\). The module
coefficients below are the coefficient-two components. We write
\(\rho_{n,2}^I\) when the integral cocycle is viewed in \(\mathbf Z_2\);
exponents of roots of unity retain the full profinite cocycle.
Relation \eqref{eq:companion-kummer-25} agrees with the normalized Alexander coefficient on \(M\).

\begin{lemma}\label{lem:companion-first-layer}
Let \(O\) be unramified of degree \(2s\) over \(\mathbf Z_2\),
with Frobenius \(\sigma\), involution \(\bar z=\sigma^s z\),
and norm-one group \(T_s\). The companion character is
\(\beta_s(z)=b(v\bar v)\) when \(z=\sigma(v)/v\);
Hilbert 90 permits a unit \(v\), and its relative norm is rational.
Another such unit differs by a rational unit, changing this norm
by a square, so the value is well-defined.
On the first principal-unit layer it is
\begin{equation}
 \beta_s(z)=\operatorname{Tr}_{\mathbf F_{2^s}/\mathbf F_2}
                   \left(\frac{z-1}{2}\bmod2\right)
       \qquad(z\in T_s\cap(1+2O)).                       
\label{eq:companion-kummer-26}
\end{equation}
For an unramified extension of odd degree \(d\),
\begin{equation}
 b(N_{O/\mathbf Z_2}z)
       =\operatorname{Tr}_{\mathbf F_{2^d}/\mathbf F_2}
                   \left(\frac{z-1}{2}\bmod2\right)
       \qquad(z\in1+2O).                                 
\label{eq:companion-kummer-27}
\end{equation}
\end{lemma}

\begin{proof}
For \eqref{eq:companion-kummer-26}, Hilbert 90 gives \(z=\sigma(v)/v\) with \(v\in O^\times\).
Since \(z\equiv1\pmod2\), the residue of \(v\) is fixed by Frobenius
and therefore equals one. Write \(v=1+2x\), and put
\(t=v\bar v\in\mathbf Z_2^\times\). Then
\[
 b(t)=x+\bar x\pmod2,
 \qquad \frac{z-1}{2}=x^2+x\pmod2.
\]
The norm-one condition places the latter value in
\(\mathbf F_{2^s}\), and its trace telescopes to
\(x^{2^s}+x=\bar x+x\). This is \eqref{eq:companion-kummer-26}, by the definition
of \(\beta_s\) through the coinvariant isomorphism.
For \eqref{eq:companion-kummer-27}, expansion of the product of conjugates gives
\(N(1+2x)\equiv1+2\operatorname{Tr}(x)\pmod4\).
\end{proof}

\begin{lemma}\label{lem:cyclic-kummer-fourier}
Let \(g\in M\) be a standard positive peripheral automorphism as
in Theorem~\ref{thm:companion-kummer-equality}. In
\(\mathbf Z_2[C_p]=\mathbf Z_2[X]/(X^p-1)\), put
\[
        C_{p,g}(X)=(X-1)R_g(X,1)=\sum_{j=0}^{p-1}c_jX^j.
\]
Then
\begin{equation}
 C_{p,g}(1)=0,\qquad C_{p,g}(X^{-1})=C_{p,g}(X),
 \qquad \rho_{p,2}^I(g)=-c_0.                            
\label{eq:companion-kummer-28}
\end{equation}
If \(\mathcal S\) contains one representative from each reciprocal
pair of nontrivial \(p\)-th roots of unity, then the sum
\[
                S_{p,g}=\sum_{\zeta\in\mathcal S}C_{p,g}(\zeta)
\]
belongs to \(\mathbf Z_2\), independently of those representatives,
and
\begin{equation}
                 \frac{\rho_{p,2}^I(g)}2=-\frac{S_{p,g}}p.
\label{eq:companion-kummer-29}
\end{equation}
\end{lemma}

\begin{proof}
Use the cyclic \(x\)-cover
\(H_p=\ker(F\to C_p)\), with free basis
\(X_0=x^p\), \(Y_j=x^jyx^{-j}\). The class of \([x,y]\)
in its abelianization is \([Y_1]-[Y_0]\). Conjugation by \(x\)
shifts the indices, and conjugation by \(y\) is trivial.
Consequently \eqref{eq:companion-kummer-25} gives
\[
                    [f]=\sum_j c_j[Y_j]
                  \quad\hbox{in }H_p^{\mathrm{ab}}\otimes\mathbf Z_2.
\]
The zero \(x\)- and \(y\)-exponents of \(f\) give no
\(X_0\)-coordinate and \(\sum_jc_j=0\). By definition
\(\rho_{p,2}^I=-\ell_p(f)_2=-c_0\).

Here is the exact boundary symmetry needed for the second identity
in \eqref{eq:companion-kummer-28}. Write \(f_i=x^ifx^{-i}\). Then
\(\phi(Y_i)=f_i^{-1}Y_if_i\), and
\([f_i]=\sum_jc_{j-i}[Y_j]\). The word
\[
                      Y_0Y_1\cdots Y_{p-1}X_0=(yx)^p
\]
is preserved up to conjugation. Its conjugator can be moved into
\(H_p\) by multiplying by a power of \(yx\). In the free class-two
quotient of \(H_p\), comparison of the \([X_0,Y_i]\) coefficients
makes all coordinates of the boundary conjugator equal. The
remaining coefficient at \([Y_i,Y_j]\) is therefore
\(c_{j-i}-c_{i-j}=0\). This proves \(c_j=c_{-j}\) integrally.
It is the cyclic-cover symmetry, and uses no distribution relation.

Fourier inversion over the unramified splitting ring gives
\[
       p c_0=\sum_{\zeta^p=1}C_{p,g}(\zeta)=2S_{p,g}.
\]
The first equality in \eqref{eq:companion-kummer-28} removes the trivial root, and the
reciprocal symmetry gives the second equality. Frobenius permutes
the reciprocal pairs, so \(S_{p,g}\) is an integral rational
two-adic number. Formula \eqref{eq:companion-kummer-29} follows. In particular \(\rho_p^I\)
is even here without an additional parity assumption.
\end{proof}

\begin{proof}[Proof of Theorem~\ref{thm:companion-kummer-equality}]
The two-cycle relation gives \(R_g(X,Y)=R_g(Y,X)\), and hence
\(B_g(X,Y)=B_g(Y,X)\). The exact duplication identity, in the
completed-group-algebra coordinates used here, is
\begin{equation}
                  B_g(X,X)=X^{2\rho_2^I(g)}B_g(-1,X).
\label{eq:companion-kummer-30}
\end{equation}
This is the integral translation displayed on printed p.~102 in
\cite{Furusho2003}: its exponent is \(-2\Psi_2^{(0)}\), and
\(\rho_2^I=-\Psi_2^{(0)}\). That paper's Theorem~3.1 proves
duplication from (I), (II), (IV), and \((K_2)\). The established
full-GT versions of (IV) and \((K_2)\) therefore supply \eqref{eq:companion-kummer-30}.
We use this explicit integral identity, rather than replacing
the coefficient-two value at \(-1\) by its residue at one.

For a nontrivial \(p\)-th root \(\zeta\), the factor
\(\zeta^{2\rho_2^I(g)}\) has odd order, so every torus character
to \(\mathbf F_2\) kills it. We may compute the companion
character on
\begin{equation}
 z_\zeta=B_g(\zeta,-1)
           =1-2(\zeta-1)R_g(\zeta,-1)
           \equiv1\pmod2.                               
\label{eq:companion-kummer-31}
\end{equation}
This value lies in the required torus: by \eqref{eq:companion-kummer-30} it differs by an
odd-order norm-one root from the compact-homology eigenvalue
\(B_g(\zeta,\zeta)\). Reduction modulo two identifies the
evaluation maps at \(Y=-1\) and \(Y=1\). Thus
\begin{equation}
          \frac{z_\zeta-1}{2}
                   =-C_{p,g}(\zeta)\pmod2.               
\label{eq:companion-kummer-32}
\end{equation}

If \(d=\operatorname{ord}_p(2)=2s\), choose one \(\zeta\)
in each Frobenius orbit. Equation \eqref{eq:companion-kummer-28} places
\(C_{p,g}(\zeta)\) in the degree-\(s\) fixed subfield, and
\[
 S_{p,g}=\sum_{\text{Frobenius orbits}}
       \operatorname{Tr}_{O_0/\mathbf Z_2}C_{p,g}(\zeta).
\]
Equations \eqref{eq:companion-kummer-26} and \eqref{eq:companion-kummer-32} show that
\(\widetilde\psi_p(g)=-S_{p,g}\pmod2\).
If \(d\) is odd, choose one orbit from each reciprocal pair.
Now
\[
 S_{p,g}=\sum_{\text{orbit pairs}}
       \operatorname{Tr}_{O/\mathbf Z_2}C_{p,g}(\zeta),
\]
and \eqref{eq:companion-kummer-27} gives the same conclusion. Finally \(p\) is odd,
so \eqref{eq:companion-kummer-29} gives
\[
          -S_{p,g}\equiv-\frac{S_{p,g}}p
                    =\frac{\rho_p^I(g)}2\pmod2.
\]
This proves \eqref{eq:companion-kummer-24} on all of \(M\).
\end{proof}

\begin{corollary}\label{cor:cubic-two-lifts-equal}
Let \(\tau\) be the earlier coefficient-two cubic frame character.
Then on all of \(M\),
\begin{equation}
                         \xi_{2,3}=\tau+\psi_3.
\label{eq:companion-kummer-33}
\end{equation}
Thus the difference between these two lifts of \(\kappa_{2,3}\)
vanishes on full GT.
\end{corollary}

\begin{proof}
First take \(g\in N\). Proposition~\ref{prop:lr-integral-cubic-beta},
whose scope is \(N\), identifies its coefficient-two
value with
\(T=\varepsilon B_g(\omega,\omega)^{\pm1}\), where
\(\varepsilon\in\{1,-1\}\), \(O=\mathbf Z_2[\omega]\), and
\(T\in1+4O\) has norm one. Both \(\alpha_1\) and \(\beta_1\)
take value one at \(-1\), while
\[
 \alpha_1(T)=\frac{\log T}{4(1+2\omega)}\pmod2=\tau(g),
                     \qquad \beta_1(T)=0.
\]
The possible conjugation or inversion does not change a quadratic
character. Hence \(\xi_{2,3}+\widetilde\psi_3=\tau\) on \(N\).
Theorem~\ref{thm:companion-kummer-equality} proves \eqref{eq:companion-kummer-33} on this subgroup.

We establish the remaining arithmetic values separately. The
compact degree-three cover is
\[
                    C:\quad w^3=t(t-1).
\]
The coordinates \(X=4w\), \(Y=4(2t-1)\) identify its
compactification with
\[
                    E_{16}:\quad Y^2=X^3+16.
\]
The isomorphism with \(E_3:Y^2=X^3+1\), over
\(\mathbf Q(2^{1/3})\), is
\[
                 (X,Y)\longmapsto
                    (X/2^{4/3},Y/4).
\]
Its Galois twisting automorphisms have order dividing three:
changing \(2^{1/3}\) by \(\omega\) multiplies the first
coordinate by \(\omega^{-1}\) and fixes the second.
Likewise, the difference between the restriction of the based
arithmetic action to the cyclic-cover subgroup and the action
using a chosen base point on \(C\) is a deck transformation.
Indeed the two choices differ by conjugation by an element of
\(F\); conjugation by its \(U_3\)-part is trivial on
\(H_1(U_3)\), leaving only \(F/U_3=C_3\).
Both quadratic torus characters kill these order-three
corrections. Consequently, on
\(H^{\mathrm{cyc}}=H\cap M\),
\[
             \xi_{2,3}(\sigma)
                    =\alpha_1(\rho_{E_3,2}(\sigma)).
\]
Here \(\rho_{E_3,2}\) is the Tate representation;
changes between the two CM embeddings do not affect \(\alpha_1\).

Proposition~\ref{prop:adelic-signed-path-comparison} applies on
\(H^{\mathrm{cyc}}\), and gives
\[
 T_\sigma=\kappa(\sigma)\rho_{E_3,2}(\sigma),\qquad
 \kappa(\sigma)=\frac{\sigma(3^{1/4})}{3^{1/4}}
                         =(-1)^{\psi_3(\sigma)}.
\]
Since \(T_\sigma\in1+4O\), applying \(\alpha_1\) gives
\(\tau(\sigma)=\psi_3(\sigma)+\xi_{2,3}(\sigma)\).
Thus \eqref{eq:companion-kummer-33} holds on all of \(H^{\mathrm{cyc}}\), including
its possibly nonzero mass coordinate \(\rho_2^I/2\bmod2\).
Finally \(G=NH\) and \(\lambda(N)=1\) imply
\(M=N H^{\mathrm{cyc}}\). All three terms in \eqref{eq:companion-kummer-33} are
homomorphisms on \(M\), so their equality on these two
factors proves it on \(M\). No full-\(M\) geometric
cubic-frame comparison is used.
\end{proof}

\subsection{The harmonic candidate as an exact comparison of two lifts}

The previous theorem eliminates the companion differences. It does
not identify the coefficient-two character with every other geometric
lift of the same arithmetic target. The harmonic residue supplies
a precise remaining comparison.

For an odd prime \(p\), choose \(i^2=-1\) in
\(\overline{\mathbf F}_p\), put
\(h_p=(p-\chi_{-4}(p))/2\), and define on \(N\)
\[
 (-1)^{\vartheta_{p,4}(g)}=B_{g,p}(i,-1)^{h_p}.
\]
This is the established harmonic residue character. It is the
quadratic torus character of compact homology modulo \(p\) of
the degree-four cyclic cover with monodromies \(1,2,1\).

\begin{proposition}\label{prop:harmonic-two-geometric-lifts}
The transgression of \(\vartheta_{p,4}\) on \(G/N=A\) is
\begin{equation}
                  e_{b,c_p}=b c_p+\frac{p-1}{2}\delta_p.
\label{eq:companion-kummer-34}
\end{equation}
Thus the existing harmonic relative character is exactly
\begin{equation}
                  \Xi_{p,4}
                     =\vartheta_{p,4}-\widetilde\psi_p
                         \quad\hbox{on }N.               
\label{eq:companion-kummer-35}
\end{equation}
Its vanishing on full GT is equivalent to equality of these two
geometric lifts of \eqref{eq:companion-kummer-34}. No vanishing or nonvanishing assertion
for \eqref{eq:companion-kummer-35} is made here.
\end{proposition}

\begin{proof}
The degree-four cover has four punctures and compact genus one.
Its compact deck eigenvalues are \(i,-i\); the order-two
character is peripheral. The positive peripheral action gives
a symplectic similitude normalizer with quotient
\[
                         C_2\times\mathbf F_p^\times.
\]
The two quotient coordinates are the deck exponent modulo four
and the similitude multiplier modulo \(p\), giving characters
\(b,c_p\). The kernel is the split torus \(\mathbf F_p^\times\)
when \(p\equiv1\pmod4\), and the norm-one torus of
\(\mathbf F_{p^2}\) when \(p\equiv3\pmod4\). In both cases
use its unique nonzero quadratic character \(\chi\).

Choose a symplectic lift \(J\) exchanging the two deck eigenlines.
In the split case one may use the matrix with rows \((0,1)\),
\((-1,0)\); in the inert case use \(z\mapsto c\bar z\)
with \(c\bar c=-1\), after normalizing the alternating trace form.
Then \(J^2=-I\). The torus has order \(p-\chi_{-4}(p)\),
which is divisible by four, so \(\chi(J^2)=0\).
Thus the \(b^2\) coefficient vanishes.

In the split case take \(L=\operatorname{diag}(1,t)\) for a
generator \(t\) of \(\mathbf F_p^\times\). Its commutator
with \(J\) has torus parameter \(t\) or \(t^{-1}\), hence
quadratic value one. Also \(L^{p-1}=1\).
In the inert case choose a generator \(v\) of
\(\mathbf F_{p^2}^\times\), and let \(L\) be multiplication
by \(v\). Its norm generates \(\mathbf F_p^\times\).
The commutator has parameter \(v^{p-1}\), or its inverse;
both it and \(L^{p-1}\) have quadratic sign
\(v^{(p^2-1)/2}=-1\) in the norm-one torus.
Hence the mixed coefficient is always one,
and the \(\delta_p\) coefficient is zero in the split case
and one in the inert case. These are all three degree-two
coefficients, proving \eqref{eq:companion-kummer-34} by the homology representation
and naturality of transgression.

The arithmetic transgression is injective, so \eqref{eq:companion-kummer-34} also proves
\(\vartheta_{p,4}|_{K}=\psi_p|_{K}\) directly, without
using an odd Anderson multiplication identity. The difference
therefore extends to a character of \(G\) killing \(H\).
Equation \eqref{eq:companion-kummer-35} follows from \eqref{eq:companion-kummer-24}. Equality of transgression
classes identifies their restrictions arithmetically; it does
not identify the two lifts on \(G\), which is precisely the
remaining relative-kernel issue.
\end{proof}

Equivalently, the unresolved finite compatibility is
\begin{equation}
 B_{g,p}(i,-1)^{(p-\chi_{-4}(p))/2}
       =(-1)^{\chi_b(\rho_{p,2}(g))}
                   \qquad(g\in N).                       
\label{eq:companion-kummer-36}
\end{equation}
The left side uses coefficient \(p\) and the degree-four cover;
the right side uses coefficient two and the degree-\(p\) cover.
The present proof establishes the right side's divided-Kummer
formula exactly. It supplies no additional implication forcing
\eqref{eq:companion-kummer-36}. In particular the exact nonabelian cyclic-filling theorem
cannot be replaced by an endomorphism of the original Alexander
module in an attempt to infer this comparison.

\section{The full quadratic restriction image and its residual kernel}
\label{sec:complete-quadratic-restriction}

The two cyclic-cover constructions lift every basis character of the
arithmetic quadratic target. This determines the image of the original
restriction map, while isolating its still undetermined kernel.
Write \(\nu_p=\xi_{2,p}|_N\in X_2\) for the coefficient-two
cyclic-cover lift of \(\kappa_{2,p}\), for each odd prime \(p\),
given by Theorem~\ref{thm:coefficient-two-cross-cups}.

\begin{theorem}\label{thm:complete-quadratic-restriction}
Restriction \(X_2\to Y_2\) is surjective. The exact sequence
\[
 0\longrightarrow\operatorname{Hom}_{\rm cts}(B,\F_2)
 \longrightarrow X_2\xrightarrow{\ \mathrm{res}\ }Y_2
 \longrightarrow0
\]
has an explicit \(\F_2\)-linear section \(s\), given on the
arithmetic basis by
\[
 \begin{aligned}
 s(\eta_3)&=\eta_3,& s(\psi_p)&=\psi_p,\\
 s(\kappa_{2,3})&=\tau|_N+\psi_3,&
 s(\kappa_{2,p})&=\nu_p\quad(p\ne3),\\
 s(\kappa_{\ell,q})&=\xi_{\ell,q}\quad(\ell<q\text{ odd}).
 \end{aligned}
\]
Thus, putting \(\mathcal S=s(Y_2)\),
\[
 X_2=\operatorname{Hom}_{\rm cts}(B,\F_2)\oplus\mathcal S,
 \qquad \mathcal S\xrightarrow{\ \mathrm{res}\ }Y_2
                 \text{ is an isomorphism}.
\]
This does not determine \(\operatorname{Hom}_{\rm cts}(B,\F_2)\).
\end{theorem}
\begin{proof}
The complete basis theorem for \(Y_2\) lists exactly the generators
in the displayed prescription. The arithmetic identities for
\(\eta_3\), \(\psi_p\), \(\tau\), the odd-pair lifts and the
coefficient-two cyclic-cover lifts show that every prescribed image
restricts to its basis element. Every element of \(Y_2\) is a finite
linear combination, so the prescription defines a linear section.
Surjectivity follows. The previously proved identification of the
restriction kernel with \(\operatorname{Hom}_{\rm cts}(B,\F_2)\)
gives exactness and the direct sum. Explicitly, for \(\chi\in X_2\),
\(\chi-s(\mathrm{res}\chi)\) is the unique component in that kernel.
No assertion about its vanishing is used.
\end{proof}

\begin{corollary}\label{cor:full-quadratic-inflation-kernel}
The two quadratic inflation kernels on the arithmetic quotient agree:
\[
 \ker\bigl(H^2(A,\F_2)\to H^2(G,\F_2)\bigr)
 =\ker\bigl(H^2(A,\F_2)\to\operatorname{Br}(\Q)[2]\bigr).
\]
\end{corollary}
\begin{proof}
The two five-term exact sequences identify these kernels with the
transgression images of \(X_2\) and \(Y_2\), respectively.
Naturality identifies the arithmetic transgression of a restriction
with the same class in \(H^2(A,\F_2)\). Surjectivity of restriction
therefore identifies the two images.
\end{proof}

The same lifts give a further simultaneous arithmetic normalization.
Let \(\mathcal J\) be the set of odd primes other than three, and put
\[
 \begin{split}
 R_{\rm quad}&=R^*\cap\bigcap_{p\in\mathcal J}\ker\nu_p,\\
 V_{\rm quad}&=V^*\cap R_{\rm quad},\qquad
 \mathcal Q=\prod_{p\in\mathcal J}\F_2.
 \end{split}
\]
Corollary~\ref{cor:cubic-two-lifts-equal} proves
\(\nu_3=\tau+\psi_3\) on \(N\). Writing the lift in this
form displays its relation to the entire cubic logarithmic coordinate.

\begin{theorem}\label{thm:complete-quadratic-normalization}
The joint map \((\nu_p)_{p\in\mathcal J}:R^*\to\mathcal Q\)
is onto already on \(R^*\cap H\), and also on \(V^*\cap H\).
Its full and arithmetic kernels are the corresponding intersections
with \(R_{\rm quad}\). In particular
\[
 R^*/R_{\rm quad}\simeq V^*/V_{\rm quad}\simeq\mathcal Q,
 \qquad G=R_{\rm quad}H.
\]
Moreover
\[
 \begin{aligned}
 B&\simeq R_{\rm quad}/
            \overline{[G,R_{\rm quad}](R_{\rm quad}\cap H)},\\
 B_{\max2}&\simeq
 \left(V_{\rm quad}/
            \overline{[G,V_{\rm quad}](V_{\rm quad}\cap H)}\right)_{\max2}.
 \end{aligned}
\]
The full joint quotients, also on arithmetic intersections, are
\[
 R_D/R_{\rm quad}\simeq V_D/V_{\rm quad}
       \simeq\mathcal U\times\mathcal P\times\mathcal Q.
\]
The kernel of restriction from \(X_2\) to either residual subgroup
is exactly \(\mathcal S=s(Y_2)\). Its image is exactly the space
of invariant quadratic characters on that subgroup vanishing on
its arithmetic intersection, canonically identified with
\(\operatorname{Hom}_{\rm cts}(B,\F_2)\).
\end{theorem}
\begin{proof}
The preceding simultaneous theorem gives \(G=R^*H\), with original
restriction kernel
\[
 \Sigma_0\oplus\F_2\tau\oplus
                       \bigoplus_{\mathcal I}\F_2\xi_{\ell,q}.
\]
By Lemma~\ref{lem:arithmetic-normalization-annihilator}, its arithmetic
image is exactly the kernel of restriction from \(Y_2\) to
\(R^*\cap H\). The complete basis of \(Y_2\) therefore makes the
remaining \(\kappa_{2,p}\), \(p\in\mathcal J\), independent on
this arithmetic subgroup. Finite linear character separation and
compactness, exactly as in the preceding theorem, prove
\[
               (\nu_p)_{\mathcal J}(R^*\cap H)=\mathcal Q.
\]
The quotient \((R^*\cap H)/(V^*\cap H)\) embeds in the pro-odd
group \(R^*/V^*\). Hence its quotient in the pro-two image
\(\mathcal Q\) is trivial, proving the full arithmetic image on
\(V^*\cap H\). Matching these values proves
\(G=R_{\rm quad}H\), the two first quotient descriptions, and
preservation of the full relative abelian quotient.

To prescribe all three factors of
\(\mathcal U\times\mathcal P\times\mathcal Q\), first choose an
arithmetic element with the specified \(\mathcal U\times\mathcal P\)
coordinates using Theorem~\ref{thm:simultaneous-quadratic-normalization}.
An element of \(R^*\cap H\), or of \(V^*\cap H\), corrects its
\(\mathcal Q\)-coordinate without changing the first two. This proves
the full joint quotient assertion and its arithmetic version.
The quotient \(R_{\rm quad}/V_{\rm quad}\) embeds in the pro-odd
quotient \(R_D/V_D\). Applying the relative pro-odd theorem with
\(G=R_{\rm quad}H\) gives the stated maximal pro-two quotient.

An original character vanishing on \(R_{\rm quad}\) restricts to
a continuous character of the product
\(R_D/R_{\rm quad}\). Its quadratic character space is the direct
sum of the \(\tau\)-line and finitely supported coordinates of
\(\mathcal P\) and \(\mathcal Q\). Subtracting these lifts leaves
exactly \(\Sigma_0\). This is precisely \(\mathcal S\), since
replacing \(\tau\) by \(\tau+\psi_3\) preserves the span.
The pro-odd quotient gives the same kernel on \(V_{\rm quad}\).

For any \(\chi\in X_2\), subtract \(s(\mathrm{res}\chi)\).
The subtracted character vanishes on both residual subgroups, and
the difference vanishes on \(K\). Thus every restricted character
is invariant and zero on the arithmetic intersection. Conversely
such a character on \(R_{\rm quad}\) extends by zero on \(H\)
using \(G=R_{\rm quad}H\). On \(V_{\rm quad}\), the relative
pro-odd theorem first gives that same extension. Restriction of the
resulting character of \(G\) to \(N\) lies in \(X_2\).
This proves the exact image assertion and its canonical
identification with the relative kernel.
\end{proof}

\section{An integral CM-type criterion at coefficient two}
\label{sec:cm-type-unit-criterion}

Let \(p\) be an odd prime, let \(n=p-1=2h\), and choose a
primitive residue \(g\) modulo \(p\). Write
\[
 F=\mathbf Q(\zeta_p),\qquad
 G_p=\operatorname{Gal}(F/\mathbf Q)=\langle s\rangle,
 \qquad s(\zeta_p)=\zeta_p^g,\qquad c=s^h.
\]
For \(a\in\mathbf F_p^\times\), denote the corresponding automorphism
by \(\sigma_a\). The half-interval CM type and its inverse are
\begin{equation}
 \Phi=\sum_{a=1}^{h}\sigma_a,
 \qquad \Phi^*=\sum_{a=1}^{h}\sigma_a^{-1}.
                                                        \label{eq:cm-type-1}
\end{equation}
This section gives the integral criterion used in the simultaneous CM
normalization. The exact arithmetic image is proved in
Section~\ref{sec:prime-degree-local-cm}. Every local factor above two
is included in the coefficient algebra.

\subsection{The explicit group-ring factor}

Define a periodic sequence, with period \(n\), by
\[
 \varphi_k=
 \begin{cases}
 1,&\text{if the least positive residue of }g^{-k}\text{ is at most }h,\\
 0,&\text{otherwise}.
 \end{cases}
\]
Thus \(\varphi_{k+h}=1-\varphi_k\), and
\(\Phi^*=\sum_{k=0}^{n-1}\varphi_k s^k\). Put
\begin{equation}
 Q_p(s)=\sum_{k=0}^{h-1}(\varphi_k-\varphi_{k-1})s^k,
 \qquad
 R_p=\mathbf Z_2[G_p]/(1+c)
       =\mathbf Z_2[s]/(s^h+1).
                                                        \label{eq:cm-type-2}
\end{equation}
The index \(-1\) in this formula means \(n-1\).

\begin{proposition}\label{prop:explicit-cm-type-factor}
One has the exact identity
\begin{equation}
                 (s-1)\Phi^*=(c-1)Q_p
                  \quad\text{in }\mathbf Z[G_p].        \label{eq:cm-type-3}
\end{equation}
The class of \(Q_p\) modulo \(1+c\) is uniquely determined by \eqref{eq:cm-type-3}.
It is a unit of \(R_p\) if and only if
\begin{equation}
        \gcd\bigl(\overline Q_p(X),X^h+1\bigr)=1
                          \quad\text{in }\mathbf F_2[X].\label{eq:cm-type-4}
\end{equation}
If \(h=2^v m\) with \(m\) odd, one may replace \(X^h+1\)
in \eqref{eq:cm-type-4} by \(X^m-1\). Moreover
\begin{equation}
                         Q_p(1)=\pm1.                  \label{eq:cm-type-5}
\end{equation}
Consequently every prime with \(p-1\) a power of two satisfies
\eqref{eq:cm-type-4}. The unit condition is independent of the choice of primitive
generator of \(G_p\), of translation of the CM type, and of replacing
the type by its complement.
\end{proposition}

\begin{proof}
The coefficient of \(s^k\) in \((s-1)\Phi^*\) is
\(\varphi_{k-1}-\varphi_k\). The relation
\(\varphi_{k+h}=1-\varphi_k\) makes these coefficients
anti-invariant under \(c\). For \(0\leq k<h\), the coefficients
of \((c-1)Q_p\) at \(s^k,s^{k+h}\) are respectively
\(-\varphi_k+\varphi_{k-1}\) and
\(\varphi_k-\varphi_{k-1}\). This proves \eqref{eq:cm-type-3}.

For any integral cyclic group ring of even order, the image of
\(c-1\) is exactly its anti-invariant subgroup: an anti-invariant
element has coefficient pairs \((a_k,-a_k)\). The kernel of
\(c-1\) is exactly \((1+c)\mathbf Z[G_p]\), because its
coefficient pairs are equal. This proves uniqueness with the stated
modulus, without division by two.

The ring \(R_p\) is finite free over \(\mathbf Z_2\). An element
is invertible if and only if its reduction modulo two is invertible:
one direction is immediate, and in the other direction an inverse
modulo two is corrected by the convergent geometric series for an
element of \(1+2R_p\). This gives \eqref{eq:cm-type-4}. In characteristic two,
\(X^h+1=(X^m-1)^{2^v}\), proving the reduced test.
Telescoping and the half-period relation give
\[
 Q_p(1)=\varphi_{h-1}-\varphi_{n-1}
                    =2\varphi_{h-1}-1.
\]
This proves \eqref{eq:cm-type-5}, including the asserted case \(m=1\).

If a different generator is \(s^u\), take \(u>0\) coprime to
\(n\). Its factor in the same ring is
\[
       Q_p\,(1+s+\cdots+s^{u-1}).
\]
The second factor is already a unit in
\(\mathbf Z[s]/(s^h+1)\). Indeed its determinant as a
\(\mathbf Z\)-linear multiplication operator is
\[
 \prod_{\zeta^h=-1}\frac{\zeta^u-1}{\zeta-1}=1:
\]
the map \(\zeta\mapsto\zeta^u\) permutes these roots and none
of them is one. A determinant-one integral operator has an integral
inverse. Translation of the type multiplies \(Q_p\) by a group
element, and complementation changes its sign, because
\((s-1)\sum_{\gamma\in G_p}\gamma=0\). These operations preserve
the unit condition.
\end{proof}

For completeness, the half-interval type has trivial multiplicative
stabilizer. If \(a\{1,\ldots,h\}=\{1,\ldots,h\}\) modulo \(p\),
then summing gives
\[
 a\frac{h(h+1)}2=\frac{h(h+1)}2=-\frac18\neq0
                              \quad\text{in }\mathbf F_p.
\]
Hence \(a=1\). This is the group-theoretic assertion needed when
identifying the reflex field of this CM type with \(F\); it does
not by itself identify any arithmetic inertia image.

\subsection{The full local norm-one torus}

Set
\begin{equation}
 A=\mathcal O_F\otimes\mathbf Z_2,\qquad
 U=1+2A,\qquad
 T=\{z\in U:z\,c(z)=1\}.                                \label{eq:cm-type-6}
\end{equation}
The algebra \(A\) is a product of unramified discrete valuation
rings. Their residue multiplicative groups have odd order. Thus
\(U\) is the maximal pro-two subgroup of \(A^\times\), and
\(T\) is the maximal pro-two subgroup of its norm-one torus.
For a group-ring element, exponent notation denotes its natural
multiplicative action on these abelian groups; for example
\((s-1)z=s(z)/z\).

\begin{theorem}\label{thm:cm-type-torus-unit}
The action of \(\mathbf Z_2[G_p]\) on \(T\) factors through
\(R_p\). The following conditions are equivalent:
\begin{enumerate}
\item \(Q_p\) is a unit of \(R_p\);
\item \(Q_p:T\longrightarrow T\) is an automorphism;
\item \((s-1)\Phi^*(U)=T\).
\end{enumerate}
For every \(p\), without a unit hypothesis, the exact identity is
\begin{equation}
                     (s-1)\Phi^*(U)=Q_pT.               \label{eq:cm-type-7}
\end{equation}
If \(\mathcal D_p\) is the determinant of multiplication by
\(Q_p\) on the free rank-\(h\) module \(R_p\), then the
unit criterion is equivalently \(\mathcal D_p\) odd. In particular
it is an integral determinant criterion, with no class-number
assumption.
\end{theorem}

\begin{proof}
The relation \((1+c)T=0\) is the definition of \(T\). Continuous
\(\mathbf Z_2\)-exponents on a pro-two abelian group define the
claimed ring action. If \(Q_p\) is a unit, its inverse in \(R_p\)
acts as its continuous inverse on \(T\).

The relative unit version of Hilbert 90 gives
\begin{equation}
                             (c-1)U=T.                  \label{eq:cm-type-8}
\end{equation}
Here is the argument for the product algebra in \eqref{eq:cm-type-6}. If \(c\)
interchanges two field factors, a norm-one pair has the form
\((z,z^{-1})\), and is the quotient of a principal-unit pair
by its \(c\)-conjugate. If \(c\) fixes a factor, it acts as the
nontrivial automorphism of an unramified quadratic extension of
the fixed field. Ordinary Hilbert 90 writes a norm-one principal
unit as \(c(v)/v\). Removing a power of two makes \(v\) a unit.
Its residue is fixed by \(c\), since the quotient has residue one.
Dividing by the Teichmuller lift of that fixed residue makes
\(v\) a principal unit. This proves \eqref{eq:cm-type-8} in every orbit of
factors. For the ordinary Hilbert 90 and unramified unit norm
statements, see \cite[II, Proposition~1.22 and Corollary~1.23,
p.~67; III, Propositions~1.1--1.2, pp.~97--98]{MilneCFT}.
Combining \eqref{eq:cm-type-8} with \eqref{eq:cm-type-3} proves \eqref{eq:cm-type-7}.

It remains to justify the converse, including integral rather than
merely rational invertibility. The open subgroup
\(T(4)=T\cap(1+4A)\) is identified by the two-adic logarithm with
\begin{equation}
      4A^-,\qquad A^-=\{a\in A:c(a)=-a\}.               \label{eq:cm-type-9}
\end{equation}
The elements \(\zeta_p^a\), \(1\leq a\leq p-1\), form an
integral basis of \(\mathcal O_F\): they replace the element one
of the usual power basis using
\(1=-\sum_{a=1}^{p-1}\zeta_p^a\), with an integral inverse
change of basis. The group \(G_p\) permutes this basis regularly.
Consequently \(A\simeq\mathbf Z_2[G_p]\), and the elementary
anti-invariant calculation used in Proposition
\ref{prop:explicit-cm-type-factor} gives \(A^-\simeq R_p\).
Thus \eqref{eq:cm-type-9} is a free rank-one \(R_p\)-module.

The group \(T\) is a finitely generated \(\mathbf Z_2\)-module
with finite torsion, and its torsion-free quotient has the same
rational representation as \eqref{eq:cm-type-9}. If \(Q_p\) is onto \(T\),
it is onto that free quotient, so its determinant there is a
two-adic unit. That determinant is \(\mathcal D_p\), because
determinants are independent of the chosen integral lattice in
the rational representation. Multiplication by \(Q_p\) on
\(R_p\) is therefore invertible. This proves both the converse
and the equivalence with condition three.
\end{proof}

\subsection{The exact coinvariants and their quadratic characters}

For \(t\in\mathbf Z_2^\times\), put
\[
 a(t)=\frac{t^2-1}{8}\pmod2,
 \qquad b(t)=\frac{t-1}{2}\pmod2.
\]
These two characters identify
\(\mathbf Z_2^\times/(\mathbf Z_2^\times)^2\) with
\(\mathbf F_2^2\).

\begin{theorem}\label{thm:full-cyclotomic-torus-coinvariants}
There is a canonical isomorphism
\begin{equation}
 \Delta_s:\mathbf Z_2^\times/(\mathbf Z_2^\times)^2
       \xrightarrow{\ \sim\ }T/(s-1)T,
 \qquad
 t\longmapsto[s(v)/v],\quad v\in U,\quad v\,c(v)=t.
                                                        \label{eq:cm-type-10}
\end{equation}
In particular
\begin{equation}
 T^{G_p}=\{1,-1\},\qquad
 T/(s-1)T\simeq C_2^2,
 \qquad [T:(s-1)T]=4.                                  \label{eq:cm-type-11}
\end{equation}
Define
\(\chi_a=a\circ\Delta_s^{-1}\) and
\(\chi_b=b\circ\Delta_s^{-1}\) on \(T\). Then
\begin{equation}
             (s-1)T=\ker\chi_a\cap\ker\chi_b.           \label{eq:cm-type-12}
\end{equation}
They are exactly the two invariant torus characters of the
coefficient-two cyclic-cover construction: \(\chi_a\) defines
\(\xi_{2,p}\), and \(\chi_b\) defines its companion
\(\widetilde\psi_p\). Thus the companion equality
\(\widetilde\psi_p=\rho_p^I/2\bmod2\) identifies which known
Kummer coordinate is removed in \eqref{eq:cm-type-12}.
\end{theorem}

\begin{proof}
The relative norm \(U\to U^c\) is onto. For paired factors this
is immediate. For an internal unramified quadratic extension,
surjectivity of the unit norm follows first on residue units and
then on successive principal-unit quotients by the surjectivity
of the residue trace. If the target is a principal unit, the
residue of an initial lift has norm one, and dividing by its
Teichmuller lift gives a principal-unit lift. This proves the
existence of \(v\) in \eqref{eq:cm-type-10}.

The quotient \(s(v)/v\) has relative norm one, because \(t\)
is rational. Replacing \(v\) by another norm lift changes this
quotient by an element of \((s-1)T\). A rational square \(t=w^2\)
has the lift \(v=w\), so the indicated map is well defined on
square classes and is a homomorphism.

We spell out the Hilbert 90 step for the full product algebra.
The algebra \(E=F\otimes\mathbf Q_2\), with its \(G_p\)-action,
is a cyclic Galois algebra over \(\mathbf Q_2\). Its factors are
permuted transitively by \(s\); the stabilizer of a factor is the
decomposition group \(\langle\sigma_2\rangle\). Cyclic Hilbert
90 for this algebra follows directly by solving the equations
successively around the cycle of factors and then applying
ordinary Hilbert 90 to the resulting equation in one field
factor. The consistency condition is precisely the full norm
equal to one. This argument does not require \(E\) to be a field.

For \(z\in T\), the full norm to \(\mathbf Q_2\) is one,
since \(N_{G_p}=(1+s+\cdots+s^{h-1})(1+c)\). Thus
\(z=s(v)/v\) for \(v\in E^\times\). Since \(z\) is a unit,
the valuation vector of \(v\) is \(G_p\)-fixed. Transitivity
and unramifiedness make all valuations equal integers; dividing
by one rational power of two makes \(v\) a unit in every factor.
Since \(z\equiv1\pmod{2A}\), its residue satisfies \(s(v)=v\).
The invariant subalgebra of \(A/2A\) is \(\mathbf F_2\):
invariants in one factor are the Frobenius-fixed elements, and
the other factors are determined by transitivity. Its only
nonzero unit is one. Consequently \(v\in U\).

Now \(s(v\,c(v))=v\,c(v)\), so its relative norm is in
\(U^{G_p}=\mathbf Z_2^\times\). This proves surjectivity of
\eqref{eq:cm-type-10}. If \(s(v)/v=s(z)/z\) for some \(z\in T\), then
\(v/z\) is a rational unit, and \(v\,c(v)\) is its square.
This proves injectivity. All the groups are compact and the
images of the continuous maps are closed, so the isomorphism
is continuous. The fixed subgroup assertion follows from
\(U^{G_p}=\mathbf Z_2^\times\): relative norm one means \(t^2=1\).
This proves \eqref{eq:cm-type-11} and \eqref{eq:cm-type-12}.

We compare the characters with their local-factor definitions.
Let \(d=\operatorname{ord}_p(2)\). If \(d=2e\) is even,
there are \(r=n/d\) self-dual factors, each an unramified
extension \(L/\mathbf Q_2\) of degree \(2e\), and
\(c\) acts there as \(\sigma^e\), with \(\sigma\) Frobenius.
The local norm-one coinvariant map is
\[
 \Delta_{\sigma}(t)=[\sigma(v)/v],\qquad v\sigma^e(v)=t.
\]
The established local characters are the sums, over the \(r\)
factors, of \(a\circ\Delta_{\sigma}^{-1}\) and
\(b\circ\Delta_{\sigma}^{-1}\). Choose factor representatives
\(1,s,\ldots,s^{r-1}\). The last passage around them is
\(\sigma^k\), where \(s^r=\sigma_2^k\) and
\(\gcd(k,d)=1\); in particular \(k\) is odd. Choose the same
norm lift of \(t\) in these coordinates. All quotients cancel
except the last, \(\sigma^k(v)/v\). Its local coinvariant is
\(\Delta_{\sigma}(t)^k=\Delta_{\sigma}(t)\). Thus the two
sums take the values \(a(t),b(t)\) on \eqref{eq:cm-type-10}.

If \(d\) is odd, the factors occur in \(r=n/(2d)\) reciprocal
pairs. Write their norm-one coordinates as \((z_i,z_i^{-1})\)
with \(z_i\) in an unramified extension of degree \(d\).
The established characters are
\[
        \sum_i a(N_{L/\mathbf Q_2}z_i),\qquad
        \sum_i b(N_{L/\mathbf Q_2}z_i).
\]
Choose the norm lift \((1,t)\) on every pair. A primitive
generator cycles the pairs, and the last passage exchanges
the reciprocal factors, possibly with a Frobenius twist.
Its quotient has a single torus coordinate \(t\) or \(t^{-1}\).
The norm is \(t^{\pm d}\); since \(d\) is odd, its two
characters are again \(a(t),b(t)\). This proves the claimed
identification in both cases.
\end{proof}

The group \((s-1)T\) has \(\mathbf Z_2\)-rank \(h\), because
it is open in \(T\). It need not be torsion-free: when the local
algebra has several orbits under complex conjugation, differences
of their signs can give nontrivial torsion in this image. No
identification of this entire compact group with a free
\(\mathbf Z_2\)-module is intended.

In particular, if an arithmetic construction establishes a first
commutator image equal to \(T\), its further commutator by an
element acting as \(s\) has exact image \((s-1)T\). Formula
\eqref{eq:cm-type-12} describes the two quadratic torus conditions on that
second image. This implication is purely algebraic; it does not
assert that a single inertia group supplies all the local factors.

\subsection{Exact indices and illustrative primes}

There is also a finite formula for the possible failure of the
unit condition. Put
\[
 d=\operatorname{ord}_p(2),\qquad
 r=\begin{cases}
       n/d,& d\text{ even},\\
       n/(2d),&d\text{ odd}.
    \end{cases}
\]
Each unramified factor has only \(\{1,-1\}\) as pro-two
torsion. The relative norm identifies the two signs in a paired
orbit. Thus
\begin{equation}
 T_{\mathrm{tors}}\simeq(\mathbf F_2)^r
              \simeq\mathbf F_2[s]/(s^r-1),              \label{eq:cm-type-13}
\end{equation}
where \(s\) permutes the \(r\) orbits of complex conjugation
transitively. Define
\[
 e_p=\deg\gcd(\overline Q_p(X),X^r-1).
\]
Then \(\ker(Q_p|T_{\mathrm{tors}})\) has order \(2^{e_p}\).
If \(\mathcal D_p\neq0\), the exact image index is
\begin{equation}
                  [T:Q_pT]=2^{v_2(\mathcal D_p)+e_p}.    \label{eq:cm-type-14}
\end{equation}
Indeed, on the free quotient of \(T\) the kernel is zero and
the cokernel has order \(2^{v_2(\mathcal D_p)}\). The exact
sequence for the torsion subgroup and its free quotient then
multiplies this index by the order of the torsion cokernel,
which equals the order of its kernel. If \(\mathcal D_p=0\),
the image has infinite index. No universal nonvanishing of this
determinant is needed for the criterion or asserted here.

Using the least primitive positive residue as \(g\), the first
examples are as follows. Each displayed determinant is the
integer resultant \(\operatorname{Res}_X(X^h+1,Q_p(X))\).
\[
\begin{array}{c|c|l|r|c}
 p&g&Q_p(X)&\mathcal D_p&[T:Q_pT]\\\hline
 5&2&-X&1&1\\
 7&3&-X&1&1\\
 11&2&-X+X^2-X^3&3&1\\
 13&2&-X+X^3-X^4&5&1\\
 17&3&-X^3+X^4-X^7&17&1\\
 19&2&-X+X^2-X^3+X^4-X^6&27&1\\
 23&5&-X+X^3-X^5+X^6-X^7&267&1\\
 29&2&-X+X^3-X^4+X^5-X^7+X^{10}-X^{11}&4520&8
\end{array}
\]
For \(p=29\), the obstruction is visible without a large
determinant:
\[
 \gcd(\overline Q_{29}(X),X^{14}+1)=X^3+X^2+1.
\]
Here \(d=28\), so \(r=1\), and \(Q_{29}(1)=-1\) gives
\(e_{29}=0\). Since \(4520=8\cdot565\), \eqref{eq:cm-type-14} gives the
stated exact index. Thus the unit criterion excludes
some primes. Further exact failures occur at \(113\), with
the same cubic greatest common divisor, and at \(163\), with
greatest common divisor \(X^2+X+1\); their indices are respectively
eight and four. These are finite arithmetic computations of
\eqref{eq:cm-type-2}--\eqref{eq:cm-type-4}, not an extrapolation or a class-number argument.

The case \(p=7\) illustrates why the product algebra matters.
Here \(d=3\), so \(A\) has two unramified cubic factors and
complex conjugation interchanges them. Nevertheless \(Q_7=-s\)
is an integral unit, \(T\) is the principal-unit group of one
cubic factor, and \eqref{eq:cm-type-10}--\eqref{eq:cm-type-12} still give the full two-dimensional
quadratic quotient. Establishing that the relevant arithmetic
inertia groups realize this full product-unit calculation is a
separate CM reciprocity step.

\begin{proposition}\label{prop:cm-commutator-torus-torsion}
With the notation of Theorem
\ref{thm:full-cyclotomic-torus-coinvariants}, let \(P=(s-1)T\)
and let \(r\) be the number of orbits of complex conjugation on
the primes above two. Then
\[
 P_{\mathrm{tors}}=(s-1)T_{\mathrm{tors}}
       \simeq C_2^{r-1}.
\]
In particular, as compact abelian groups, noncanonically,
\[
       T\simeq\mathbf Z_2^h\times C_2^r,
       \qquad P\simeq\mathbf Z_2^h\times C_2^{r-1}.
\]
For example, \(p=17\) has \(r=2\), so \(P\) has one
nontrivial element of order two, even though the unit criterion
holds.
\end{proposition}

\begin{proof}
Under \eqref{eq:cm-type-13}, the image \((s-1)T_{\mathrm{tors}}\) is the
even-parity sign subgroup and has dimension \(r-1\) over
\(\mathbf F_2\). We show that the restriction of
\((\chi_a,\chi_b)\) to \(T_{\mathrm{tors}}\) has kernel
exactly this subgroup. This will prove equality with
\(P\cap T_{\mathrm{tors}}\), which is the torsion subgroup
of \(P\).

If \(d=\operatorname{ord}_p(2)\) is odd, a sign in one paired
factor has odd-degree norm \(-1\). Its character values are
\(a(-1)=0\) and \(b(-1)=1\). Thus the restriction is
\((0,\sum_i\epsilon_i)\) on the sign vector
\((\epsilon_1,\ldots,\epsilon_r)\).

If \(d=2e\) is even, let \(\sigma\) be Frobenius in one
unramified factor. Its unramified quadratic subfield contains
\(\delta=1+2\omega\), with \(\omega^2+\omega+1=0\).
Then \(\delta^2=-3\), \(\sigma(\delta)=-\delta\), and
\[
 \delta\,\sigma^e(\delta)=
 \begin{cases}3,&e\text{ odd},\\-3,&e\text{ even}.
 \end{cases}
\]
Hence the local coinvariant of \(-1\) corresponds to the
square class of \(3\) or \(-3\), respectively. Its two
character values are \((1,e\bmod2)\). On all factors the
restriction is therefore
\((\sum_i\epsilon_i,e\sum_i\epsilon_i)\). In both cases
its kernel is the even-parity sign subgroup, as required.

Finally, a finitely generated compact \(\mathbf Z_2\)-module
splits as its finite torsion subgroup and a free module.
Both \(T\) and its open subgroup \(P\) have rank \(h\),
which gives the two displayed decompositions.
\end{proof}

\subsection{The classical class-number interpretation of the local criterion}
\label{subsec:cm-class-number}
The determinant in the local criterion is governed by the classical
relative class-number formula.  We record the normalization explicitly
to distinguish invertibility of this integral operator from the later
question of global arithmetic saturation.

Let \(p\geq5\) be prime, put \(h=(p-1)/2\), choose a primitive residue
\(g\) modulo \(p\), and let \(s=[g]\), \(c=s^h\).  As above,
\[
 \Phi^\ast=\sum_{a=1}^{h}[a]^{-1},\qquad
 (s-1)\Phi^\ast=(c-1)Q_p,\qquad
 R_p=\mathbf Z[s]/(s^h+1).
\]
Let \(h_p^-\) be the relative class number of
\(\mathbf Q(\zeta_p)\), let \(d=\operatorname{ord}_p(2)\), and put
\(r=(p-1)/d\).

\begin{proposition}\label{prop:cm-class-number-determinant}
Multiplication by \(Q_p\) on \(R_p\) has positive determinant
\[
 \det(Q_p\mid R_p)=\frac{h_p^-}{p}
 \begin{cases}
 (2^{d/2}+1)^r,&d\ \text{even},\\
 (2^d-1)^{r/2},&d\ \text{odd}.
 \end{cases}
\]
Consequently the determinant is nonzero for every \(p\), and
\[
 v_2\bigl(\det(Q_p\mid R_p)\bigr)=v_2(h_p^-).
\]
The operator \(Q_p\) is a unit in
\(\mathbf Z_2[s]/(s^h+1)\) if and only if \(h_p^-\) is odd,
equivalently if and only if the class number of
\(\mathbf Q(\zeta_p)\) is odd.
\end{proposition}
\begin{proof}
For an odd character \(\psi\) of
\((\mathbf Z/p)^\times\), set
\(B_{1,\psi}=p^{-1}\sum_{a=1}^{p-1}a\psi(a)\).
Permutation of the nonzero residues by multiplication by \(2\),
using the least positive representatives and
\(\sum_{a>h}\psi(a)=-\sum_{a\leq h}\psi(a)\), gives
\[
 \sum_{a=1}^{h}\psi(a)
       =(\psi(2)^{-1}-2)B_{1,\psi}.
\]
Evaluation of the defining identity at an odd character \(\chi\)
therefore gives
\[
 Q_p(\chi)=
 \frac{(1-\chi(g))(\chi(2)-2)B_{1,\chi^{-1}}}{2}.
\]
The odd values of \(\chi(g)\) are the roots of \(X^h+1\), so
their factors \(1-\chi(g)\) have product \(2\).  Kummer's
relative class-number formula, in the normalization
\cite[Eq.~(1.2), p.~1]{Agoh2012}, is
\[
 h_p^-=2p\prod_{\chi\ {\rm odd}}\left(-\frac{B_{1,\chi}}2\right).
\]
Multiplying the evaluated identities yields
\[
 \det(Q_p\mid R_p)
     =\frac{h_p^-}{p}\prod_{\chi\ {\rm odd}}(2-\chi(2)).
\]

If \(d\) is even, then \(2^{d/2}=-1\pmod p\), and the values
\(\chi(2)\) for odd \(\chi\) are all roots of \(X^{d/2}+1\),
each with multiplicity \(r\).  If \(d\) is odd, they are all
roots of \(X^d-1\), each with multiplicity \(r/2\).  This gives
the displayed positive formula.  Its factors other than
\(h_p^-\) are odd in \(\mathbf Z_2^\times\), proving the valuation
identity.  An element of a finite free \(\mathbf Z_2\)-algebra
is a unit precisely when its multiplication determinant is
a unit: the inverse linear map supplies a multiplicative inverse
by its value on \(1\).  Finally the equivalence of odd relative
and odd full cyclotomic class number is
\cite[Proposition~1, equivalence (1)--(2), p.~207]{Dummit2020}.
\end{proof}

Thus the prime set \(\Pi\) in the simultaneous local
normalization theorem can equivalently be written
\[
 \Pi=\{p\geq5:\ h(\mathbf Q(\zeta_p))\text{ is odd}\}.
\]
This criterion concerns the local reflex-norm operator.
The degree-\(29\) and degree-\(113\) global arguments show that
an even class number need not prevent arithmetic saturation.
When the local index is computed on the unit group itself,
the finite sign subgroup in the integral index formula must
still be included.

\section{An exact test for the local first-layer hypothesis}
\label{sec:cm-exact-first-layer}
Let \(p\ge5\), \(h=(p-1)/2\), and let \(g\) be a primitive residue
modulo \(p\).  Put
\[
 \begin{gathered}
 s=[g],\quad c=s^h,\quad O=\mathbb Z[\zeta_p]\otimes\mathbb Z_2,\\
 T=\{u\in1+2O:u\,c(u)=1\},\quad P=(s-1)T,\quad
 \mathscr M=(O/2O)^c.
 \end{gathered}
\]
Operators on unit groups are written additively, so
\((s-1)u=s(u)u^{-1}\).
Let \(Q\) be the integral CM operator, assume \(Q(1)\) is odd, and
assume
\[
 D=\det\bigl(Q\mid\mathbb Z_2[s]/(s^h+1)\bigr)\ne0.
\]
Write
\[
 d=\operatorname{ord}_p(2),\quad r=(p-1)/d,\quad
 e=\begin{cases}r,&d\text{ even},\\ r/2,&d\text{ odd},\end{cases}
\]
and
\[
 b=\deg\gcd(\overline Q,X^h-1),\qquad
 t=\deg\gcd(\overline Q,X^e-1).
\]

\begin{proposition}\label{prop:cm-exact-first-layer}
With every finite sign subgroup included,
\[
 [P:QP]=[T:QT]=2^{v_2(D)+t},\qquad v_2(D)\ge b.
\]
The natural reduction
\[
 \alpha:T\longrightarrow \mathscr M,\qquad
 \alpha(u)=\frac{u-1}{2}\pmod{2O}
\]
induces a surjection \(T/QT\to \mathscr M/\overline QM\).  It is an isomorphism
if and only if \(t=0\) and \(v_2(D)=b\).
The logarithmic convention
\[
 \alpha_{\log}(u)=\frac{\log u}{2}\pmod{2O}
\]
induces a map between the same quotients and has exactly the same
isomorphism criterion.
\end{proposition}

\begin{proof}
The algebra \(O\) is a product of \(r\) unramified degree-\(d\) integer
rings, because the irreducible factors of \(\Phi_p\) modulo \(2\)
correspond to the orbits of multiplication by \(2\) on
\((\mathbb Z/p)^\times\).
If \(d\) is even, \(-1\in\langle2\rangle\), so \(c\) fixes every factor
and acts nontrivially within it.  If \(d\) is odd, \(c\) exchanges the
factors in pairs.  In either case there are exactly \(e\) conjugation
orbits, and \(s\) permutes those orbits cyclically.

The logarithm converges on \(1+2O\), and it is inverse to the exponential
on \(1+4O\) and \(4O\).
If \(\log u=0\), then \(u^2\in1+4O\) has zero logarithm and hence
\(u^2=1\).  Conversely every component sign has logarithm zero.
The norm-one condition says that the signs agree on each conjugation
orbit.  Thus
\[
 \mathcal T:=\operatorname{Tor}(T)=(C_2)^e,
 \qquad \mathcal T\simeq\mathbb F_2[X]/(X^e-1)
\]
as an \(s\)-module.  In particular
\[
 |\mathcal T/Q\mathcal T|=2^t.
\]

The image \(\Lambda=\log T\) is a full lattice in
\(O^-=\{a\in O:c(a)=-a\}\): it is contained in \(2O^-\) and contains
\(4O^-\), the latter by exponentiation.
There is an exact sequence of \(Q\)-modules
\[
 0\longrightarrow\mathcal T\longrightarrow T
       \xrightarrow{\log}\Lambda\longrightarrow0.
\]
The elements \(\zeta_p^{g^i}\), \(0\le i<2h\), form a
\(\mathbb Z_2\)-basis of \(O\).
Indeed, a polynomial of degree at most \(p-1\) with zero constant
coefficient cannot be a nonzero scalar multiple of \(\Phi_p\).
The differences
\(\zeta_p^{g^i}-\zeta_p^{-g^i}\), \(0\le i<h\), consequently form a
basis of \(O^-\), with \(s\) acting by the signed cyclic shift.
Thus
\[
 O^-\simeq\mathbb Z_2[s]/(s^h+1).
\]
The determinant of \(Q\) on the commensurable lattice \(\Lambda\)
is also \(D\); hence \(Q\) is injective there and
\(|\Lambda/Q\Lambda|=2^{v_2(D)}\).
The exact sequence above gives
\[
 0\longrightarrow\mathcal T/Q\mathcal T
 \longrightarrow T/QT
 \longrightarrow\Lambda/Q\Lambda\longrightarrow0.
\]
This proves \([T:QT]=2^{v_2(D)+t}\); the torsion contribution is
essential.

The kernel of \(s-1:T\to P\) is \(T^s=\{1,-1\}\), because the
\(s\)-fixed subring of \(O\) is \(\mathbb Z_2\).
The operator \(Q\) is an automorphism on this kernel, since \(Q(1)\)
is odd.  Therefore \(s-1\) induces
\[
 T/QT\xrightarrow{\sim}P/QP.
\]
Explicitly, if \((s-1)x=(s-1)Qy\), then
\(x-Qy\in T^s=QT^s\), so \(x\in QT\); surjectivity follows from
the definition of \(P\).

Modulo \(2\), the signed cyclic module becomes
\(\mathbb F_2[X]/(X^h-1)\).  The nullity of multiplication by
\(\overline Q\) on this module is \(b\).
In the Smith form over \(\mathbb Z_2\), exactly \(b\) diagonal entries
of \(Q\) have positive valuation.  Their valuations sum to \(v_2(D)\);
therefore \(v_2(D)\ge b\).

If \(u=1+2a\in T\), its norm-one equation gives
\(a+c(a)+2a\,c(a)=0\).  Thus \(\overline a\in \mathscr M\), and multiplication
of units shows that \(\alpha\) is an equivariant homomorphism.
The elements
\[
 \beta_i=\overline\zeta_p^{\,g^i}
             +\overline\zeta_p^{-g^i},\qquad0\le i<h,
\]
form a basis of \(\mathscr M\), with cyclic \(s\)-action.
For any \(\overline a\in O/2O\), choose a lift \(a\in O\), put
\(v=1+2a\), and set \(u=v/c(v)\).  Then
\(\alpha(u)=\overline a+c(\overline a)\).
This trace map hits every \(\beta_i\), proving that \(\alpha:T\to \mathscr M\)
is onto.
It follows that the induced map on \(Q\)-quotients is onto and
\[
 |\mathscr M/\overline QM|=2^b.
\]
Its source has order \(2^{v_2(D)+t}\), so it is an isomorphism
exactly when \(v_2(D)+t=b\).
The inequality already proved and \(t\ge0\) make this equivalent to
\(t=0\) and \(v_2(D)=b\).

For completeness, the logarithmic convention does not silently lose
the sign contribution.  The logarithm series gives
\[
 \alpha_{\log}(u)=\alpha(u)+\alpha(u)^2=(1+[2])\alpha(u).
\]
Write \(2=g^m\pmod p\).
On \(\mathscr M\), the operator \([2]\) is \(s^m\), and
\(\gcd(h,m)=e\): the order of \(2\) modulo signs is \(d/2\) for even
\(d\), and \(d\) for odd \(d\).
Consequently
\[
 \gcd(X^h-1,X^m-1)=X^e-1.
\]
Multiplication by \(1+X^m\) on \(\mathscr M/\overline QM\) is invertible
exactly when \(t=0\); more precisely its cokernel has dimension \(t\).
When \(t=0\) and \(v_2(D)=b\), both this operator and the map induced
by \(\alpha\) are isomorphisms.  Conversely an isomorphism induced by
\(\alpha_{\log}\) forces equality of source and target orders, and therefore
the same two conditions.
\end{proof}

The class-number identity
\(v_2(D)=v_2(h_p^-)\) makes these conditions an effective test for the
local hypothesis of the global unit criterion.  This corollary
asserts no global arithmetic saturation.


\section{Arithmetic realization in split cyclotomic coefficient algebras}
\label{sec:prime-degree-local-cm}

Let \(p\ge5\) be prime, put \(g=(p-1)/2\), and consider
\[
 C_p:\ t^2-t=w^p,\qquad J_p=\operatorname{Jac}(C_p),
 \qquad F_p=\mathbf Q(\zeta_p).
\]
Write
\[
 \Gamma_p=(\mathbf Z/p)^\times,\qquad
 [a](\zeta_p)=\zeta_p^a,\qquad
 \mathcal A_p=O_{F_p}\otimes\mathbf Z_2,\qquad
 \mathcal U_p=1+2\mathcal A_p.
\]
Choose a primitive residue \(\gamma\bmod p\), put
\(s=[\gamma]\), and let \(c=[-1]=s^g\).
The groups relevant to the coefficient-two action are
\[
 \mathcal T_p=\{u\in\mathcal U_p:u\,c(u)=1\},
 \qquad \mathcal P_p=(s-1)\mathcal T_p.
\]
For an integral group-ring element \(A=\sum n_a[a]\),
the notation \(A(u)\) means \(\prod[a](u)^{n_a}\).
On pro-two unit groups this extends continuously to
coefficients in \(\mathbf Z_2\).

Let \(\rho_{p,2}\) be the established full integral
compact-homology representation of the cyclic cover with
monodromies \(1,1,-2\). Its scalar restriction to
\(M=\ker\lambda\) is denoted by \(\mathfrak h_p\).
The full duplication identity gives
\(\mathfrak h_p(N)\subset\mathcal T_p\).
We use the established full companion equality and
torus coinvariant identification in the form
\begin{equation}\label{eq:pdl-upper-input}
 \mathfrak h_p(R_{\rm quad})\subset\mathcal P_p.
\end{equation}
Explicitly, the two torus characters are the chosen
\(\nu_p\) and \(\psi_p\); their common kernel in
\(\mathcal T_p\) is \((s-1)\mathcal T_p\).
The arithmetic realization below is independent of
any assumption that two is primitive modulo \(p\).

\subsection{CM type and all local factors}

\begin{lemma}\label{lem:pdl-reflex-and-local}
The CM type of \(J_p\) is \(\{[1],\ldots,[g]\}\).
It is primitive, its reflex field is \(F_p\), and its
reflex norm on coefficient-two units is the group-ring map
\[
                           \Phi_p^*=\sum_{a=1}^{g}[a^{-1}].
\]
For every prime \(v\mid2\) of \(F_p\), let
\(\iota_v:O_{F_p,v}^\times\to\mathcal A_p^\times\)
put a unit in the \(v\)-component and one in every other
component. A corresponding inertia subgroup \(I_v\subset H\)
satisfies the exact arithmetic identity
\begin{equation}\label{eq:pdl-one-inertia}
                \rho_{J_p,2}(I_v)
                    =\Phi_p^*(\iota_v O_{F_p,v}^\times).
\end{equation}
Consequently, products of elements of these conjugate
inertia subgroups realize \(\Phi_p^*(u)\) for every
\(u\in\mathcal U_p\).
\end{lemma}

\begin{proof}
The regular differentials
\[
                   \frac{w^{j-1}\,dw}{2t-1},\qquad 1\le j\le g,
\]
have deck eigenvalues \(\zeta_p^j\). The integral deck
action embeds \(O_{F_p}\) in \(\operatorname{End}_{F_p}(J_p)\),
and its degree is \(p-1=2\dim J_p\). This proves the CM
assertion and the displayed type.
Its stabilizer in \(\Gamma_p\) is trivial: the sum of the
residues in the type is
\[
             \sum_{j=1}^{g}j=\frac{p^2-1}{8}
                                      \equiv-\frac18\pmod p,
\]
which is nonzero. If multiplication by \(a\) preserved
the type, comparison of sums would force \(a=1\).
The reflex field is therefore \(F_p\). An induced type
from a proper CM subfield would have a nontrivial stabilizer,
so the type is primitive. The inverse-embedding reflex norm
is precisely \(\Phi_p^*\), by
\cite[Example 1.19, p.~14; Proposition 1.26, p.~17]{MilneCM2020}.

The curve has good reduction at two. Besides its displayed
affine chart, use \(u=w^{-1}\), \(v=t/w^{g+1}\), giving
\[
                          v^2-u^{g+1}v=u.
\]
The two monic quadratic charts glue to a finite flat cover
of \(\mathbf P^1_{\mathbf Z_2}\). The first has derivative
\(2t-1\), a unit modulo two. In the second chart, vanishing
of the derivative in \(v\) modulo two forces \(u=0\),
where the derivative in \(u\) is \(-1\). Thus the proper
model is smooth, and its Jacobian has good reduction.

The field \(F_p\) is unramified at two. Each \(I_v\)
is consequently an inertia group of \(\mathbf Q\) at two
contained in \(G_{F_p}\). Local and global reciprocity
identify its abelian image with the global Artin image
of ideles supported on \(O_{F_p,v}^\times\).
For such an idele \(s_v\) and a corresponding \(b_v\in I_v\),
Milne's main theorem
\cite[Theorem 9.10, p.~78]{MilneCM2020} gives
\[
                   \rho_{J_p,f}(b_v)=A_v N_\Phi(s_v),
\]
where \(A_v:J_p\to b_vJ_p\) is an \(F_p\)-linear
quasi-isogeny. Here \(b_vJ_p=J_p\), since the Jacobian
is defined over \(\mathbf Q\), and \(b_v\) fixes the CM
field. At every auxiliary odd Tate coefficient, good
reduction at two makes the left side one, while the
reflex norm of \(s_v\) also has component one.
Faithfulness of the Tate action on quasi-endomorphisms
therefore gives \(A_v=1\). This proves
\eqref{eq:pdl-one-inertia}, on all units and without a
finite local correction. We use Milne's Artin convention
from printed p.~76; its reversal inverts the parameter
and does not change any image.

Every \(u\in\mathcal U_p\) is a product of the separate
\(\iota_v(u_v)\). Choose a corresponding \(b_v\) for each
factor and take their product \(b(u)\). The factors fix
\(F_p\), so their Tate images commute and multiply to
\(\Phi_p^*(u)\). No assertion that one local inertia
group supplies all factors is used.
\end{proof}

Here is an explicit description of the local product.
Let \(d=\operatorname{ord}_p(2)\) and \(r=(p-1)/d\).
Then \(\mathcal A_p\) is a product of \(r\) copies of
the unramified degree-\(d\) integer ring \(O_d\).
Choose representatives \(e_i\) for
\(\Gamma_p/\langle2\rangle\), and let \(\sigma\)
be Frobenius on \(O_d\). In these coordinates,
\begin{equation}\label{eq:pdl-frobenius-products}
 (\Phi_p^*(u))_i
    =\prod_{a=1}^{g}\sigma^{k(i,a)}(u_{j(i,a)}),
 \qquad e_i a^{-1}=2^{k(i,a)}e_{j(i,a)}\pmod p.
\end{equation}
This formula keeps every Frobenius and permutation factor.
If \(d\) is even, complex conjugation acts in each
component as \(\sigma^{d/2}\); \(\mathcal T_p\) is a
product of \(r\) relative norm-one principal-unit groups.
If \(d\) is odd, conjugation pairs components, and
\(\mathcal T_p\simeq(1+2O_d)^{r/2}\).

\subsection{An exact inertia-generated arithmetic subgroup}

\begin{proposition}\label{prop:pdl-arithmetic-realization}
For every \(u\in\mathcal U_p\), there is an element
\(w_p(u)\in V_{\rm quad}\cap H\) such that
\begin{equation}\label{eq:pdl-witness-values}
 \mathfrak h_p(w_p(u))=(s-1)^2\Phi_p^*(u),\qquad
 \mathfrak h_q(w_p(u))=1
        \quad\text{for every odd prime }q\ne p.
\end{equation}
Every original \(\chi\in X_2\) vanishes on this
entire family. In particular the image
\((s-1)^2\Phi_p^*(\mathcal U_p)\) is realized
arithmetically with all other prime-degree coordinates zero.
\end{proposition}

\begin{proof}
Choose \(a\in I_p\) with \(\lambda(a)\bmod p=\gamma\).
It has \(\lambda_2(a)=1\). Obtain \(b(u)\) as in
Lemma~\ref{lem:pdl-reflex-and-local}, and put
\[
                  z=[a,b(u)],\qquad z_2=z^{e_2},
                         \qquad w_p(u)=[a,z_2].
\]
The idempotent exponent \(e_2\in\widehat{\mathbf Z}\)
is one at two and zero elsewhere; powers are taken in
the procyclic closure of the element. Both \(z\) and
\(z_2\) have full multiplier one.

For the arithmetic Jacobian, conjugation by \(a\) acts
on the coefficient algebra by \(s\). Consequently
\[
                    \rho_{J_p,2}(z)=(s-1)\Phi_p^*(u).
\]
This value is pro-two because \(u\in\mathcal U_p\).
The based-cover action and the Jacobian action may differ
by a deck \(C_p\) factor. Indeed two lifts of the same
outer action differ by conjugation in the base geometric
group; modulo an inner automorphism of its covering
subgroup, this is exactly a deck transformation. Inner
automorphisms have trivial action on compact homology.
The discrepancy is a crossed map to \(\mu_p\), and
on the cyclotomic kernel it is a homomorphism.
It therefore vanishes on the pro-two element \(z_2\).
Thus
\[
                    \mathfrak h_p(z_2)=(s-1)\Phi_p^*(u).
\]
The full cover representation now computes its second
commutator by \(s-1\), proving the first equality in
\eqref{eq:pdl-witness-values}.

For membership in \(N\), use the coefficient-two Kummer
cocycle of two. It vanishes on \(a\in I_p\), since
\(p\ne2\). The cocycle commutator formula
\[
 \rho_q^I([h,k])=(1-\lambda_k)\rho_q^I(h)
                              +(\lambda_h-1)\rho_q^I(k)
\]
therefore gives \(\rho_{2,2}^I(z)=0\), because
\(\lambda_2(a)=1\). Taking \(e_2\) kills all other
coefficient components, so \(z_2\in N\cap H\).
Every rational Kummer coordinate of \(z_2\) is supported
at two. Hence
\[
                \rho_q^I(w_p(u))
                   =(\lambda(a)-1)\rho_q^I(z_2)=0
                           \quad\text{for every }q,
\]
and \(w_p(u)\in D\cap H\).

The based pro-\(\ell\) action of an inertia group
at two is trivial for \(\ell\ne2\), and the same is
true of every conjugate of that inertia group.
Thus \(\Phi_p(b(u))=1\). For \(\ell\ne p\), one
has \(\Phi_\ell(a)=1\). This is the automorphism-level
unramifiedness assertion with the rational tangential
base point in \cite[\S8.5, pp.~151--152]{HainMatsumoto2003}.
Consequently every \(\Phi_\ell\) kills \(z,z_2,w_p(u)\).

Since \(p\ge5\), inertia at \(p\) fixes both
\(\mathbf Q(i)\) and \(\mathbf Q(\sqrt{-3})\).
The element \(z_2\) fixes these fields as well, because
it has full multiplier one. Their harmonic and cubic
elliptic Tate actions lie in commutative CM algebras,
so their Tate commutators are trivial. The cubic frame's
quartic Kummer correction of three is trivial on \(D\).
Hence both complete corrected frames of \(w_p(u)\)
are one. The cubic radical--CM defect \(d_{(3)}\)
vanishes on the arithmetic subgroup under consideration.
This proves all the word, radical, defect, and full
coefficient-two/odd-frame conditions.

For any \(\chi\in X_2\), conjugation invariance gives
\[
                  \chi(w_p(u))
                    =\chi(a z_2a^{-1})-\chi(z_2)=0.
\]
This kills every chosen quadratic coordinate, including
the whole odd-pair family. Thus \(w_p(u)\in V_{\rm quad}\cap H\).

Finally fix an odd prime \(q\ne p\). Both \(a\) and
\(z_2\) have cyclotomic value one modulo \(q\): for
\(a\) this follows because \(\mathbf Q(\zeta_q)\) is
unramified at \(p\). Their degree-\(q\) homology actions
therefore commute with the deck generator. Its integral
centralizer is the commutative algebra
\(O_{F_q}\otimes\mathbf Z_2\): the nontrivial deck
eigencharacters occur once, and compact homology is
free of rank one over this algebra. The two actions
commute, proving the second equality in
\eqref{eq:pdl-witness-values}. This is an identity of the
full based-cover representations; no further
arithmetic or tangential comparison is needed for it.
\end{proof}

\subsection{The split prime seven}

For \(p=7\), let \(O_3\) be unramified cubic over
\(\mathbf Z_2\), with Frobenius \(\sigma\). Use the two
embeddings represented by \(1,-1\). Then
\[
 \mathcal A_7=O_3\times O_3,\qquad
 c(a,b)=(b,a),\qquad
 s(a,b)=(\sigma^2 b,\sigma^2 a),\quad \gamma=3.
\]
In particular \(s^2=\sigma\) componentwise and \(s^3=c\).
The reflex norm is
\begin{equation}\label{eq:pdl-seven-reflex}
 \Phi_7^*(a,b)
       =(a\sigma^2(a)\sigma(b),\ b\sigma^2(b)\sigma(a)).
\end{equation}
For a unit in the first factor, this gives
\[
 (s-1)\Phi_7^*(u,1)
                         =((\sigma^2u)^{-1},\sigma^2u).
\]
As \(u\) ranges over \(1+2O_3\), these are every
element of \(\mathcal T_7\). This verifies the exact
split local image directly, including its first unit layer.

Identifying \(\mathcal T_7\) with \(1+2O_3\) by
\(t\mapsto(t,t^{-1})\), one has
\[
 \mathcal P_7=(1+\sigma^2)(1+2O_3)
  =\{t\in1+2O_3:N_{O_3/\mathbf Z_2}(t)
                                  \in(\mathbf Z_2^\times)^2\}.
\]
To verify the last equality, the norm of
\(t\sigma^2(t)\) is a square. Conversely suppose
\(N(t)=v^2\), and choose the unique rational unit \(a\)
with \(a^3=v\). Then \(t/a^2\) has norm one. On this
norm-one subgroup, \(1+\sigma^2=-\sigma\), an
automorphism; and \(a^2=(1+\sigma^2)(a)\).
This proves the equality on the entire principal-unit group.
The two existing torus characters here are
\(a(Nt)\) and \(b(Nt)\), where
\(a(v)=(v^2-1)/8\pmod2\) and \(b(v)=(v-1)/2\pmod2\).
Thus their common kernel is exactly \(\mathcal P_7\),
also verifying \eqref{eq:pdl-upper-input} directly in this case.
The group \(\mathcal P_7\) has rank three and no torsion:
the only principal-unit torsion is \(-1\), whose cubic
norm is \(-1\), not a square. Hence
\(\mathcal P_7\simeq\mathbf Z_2^3\).

It follows from Proposition~\ref{prop:pdl-arithmetic-realization}
that all four images on \(R_{\rm quad},V_{\rm quad}\)
and their arithmetic intersections are \(\mathcal P_7\).
The same is true after imposing the degree-five kernel,
because the degree-seven witnesses have degree-five value one.
Thus adjoining the degree-seven kernel after degree five
gives a further strict reduction, with exact quotient
\(\mathbf Z_2^3\), already arithmetically.

\subsection{Simultaneous consequence of the integral unit criterion}

We formulate the arithmetic consequence of the explicit
CM-type unit calculation. Suppose
\[
 (s-1)\Phi_p^*=(c-1)Q_p
       \quad\text{in }\mathbf Z[\Gamma_p],
\]
and the image of \(Q_p\) is a unit in
\(\mathbf Z_2[s]/(s^g+1)\). Product-algebra unit Hilbert
90 gives \((c-1)\mathcal U_p=\mathcal T_p\): in a
self-conjugate unramified factor use ordinary unit
Hilbert 90 and remove its invariant Teichmuller part;
in a conjugate pair use \((1,t)\).
The operator \(Q_p\) is an automorphism on
\(\mathcal T_p\), since \(c\) acts there by inversion.
It follows that
\begin{equation}\label{eq:pdl-unit-surjectivity}
                     (s-1)\Phi_p^*(\mathcal U_p)=\mathcal T_p.
\end{equation}
Combining \eqref{eq:pdl-unit-surjectivity},
\eqref{eq:pdl-upper-input}, and
Proposition~\ref{prop:pdl-arithmetic-realization} gives
all four exact images \(\mathcal P_p\).
For seven the group-ring identity is
\[
 \Phi_7^*=1+s^4+s^5,\qquad
 (s-1)\Phi_7^*=(s^3-1)(-s),
\]
so its unit criterion requires no numerical certificate.

\begin{theorem}\label{thm:simultaneous-cyclotomic-cm}
Let \(\Pi\) be any set of primes at least five satisfying the equivalent
unit conditions of Theorem~\ref{thm:cm-type-torus-unit}, and put
\[
 \mathcal R_\Pi=R_{\rm quad}\cap\bigcap_{p\in\Pi}\ker\mathfrak h_p,
 \qquad \mathcal V_\Pi=V_{\rm quad}\cap\mathcal R_\Pi.
\]
These subgroups are closed and normal in \(G\). The joint image of the
old subgroups, and of either arithmetic intersection, is exactly
\begin{equation}\label{eq:pdl-simultaneous}
       \prod_{p\in\Pi}\mathcal P_p.
\end{equation}
Both old-to-new quotients are this product, also on arithmetic
intersections. Moreover \(G=\mathcal R_\Pi H\); full \(B\) is preserved
on \(\mathcal R_\Pi\), and \(B_{\max2}\) is preserved on
\(\mathcal V_\Pi\). The exact original restriction kernel on either
subgroup is \(\mathcal S=s(Y_2)\). Its complete image is the invariant
quadratic characters killing that subgroup's arithmetic intersection,
canonically \(\operatorname{Hom}_{\rm cts}(B,\mathbf F_2)\).
No nonzero continuous quadratic character of the displayed product,
pulled back to \(R_{\rm quad}\), extends to a member of \(X_2\).
\end{theorem}

\begin{proof}
The normalizer acts on each torus through \(\Gamma_p\), preserving its
identity. Thus the two kernels are closed and normal. The upper bound
and the axis witnesses have just been proved.
For every finite set of coordinates, multiply the arithmetic
witnesses from Proposition~\ref{prop:pdl-arithmetic-realization}:
each has arbitrary prescribed value on its own coordinate
and value one on all the others. This gives every finite
projection of \eqref{eq:pdl-simultaneous}. The image of
the compact arithmetic subgroup is closed, so these finite
surjections give the whole infinite product.

Arithmetic correction now gives
\[
 G=\mathcal R_{\Pi}H,\qquad
 B\simeq \mathcal R_{\Pi}/
 \overline{[G,\mathcal R_{\Pi}](\mathcal R_{\Pi}\cap H)}.
\]
The quotient \(\mathcal R_{\Pi}/\mathcal V_{\Pi}\)
embeds in the pro-odd group \(R_{\rm quad}/V_{\rm quad}\).
The central coprime extension lemma therefore preserves
\(B_{\max2}\) on \(\mathcal V_{\Pi}\).
The original restriction kernel on either subgroup remains
\(\mathcal S\). Indeed an original character vanishing on the new
kernel factors through \eqref{eq:pdl-simultaneous} on the
old subgroup. It vanishes on every arithmetic axis witness.
The finite-support subgroup is dense in the product, so
continuity forces that factored character to be zero.
The complete restricted images are again precisely the
invariant quadratic characters killing the arithmetic
intersection, canonically \(\operatorname{Hom}(B,\mathbf F_2)\),
by the complete quadratic splitting, ordinary gluing on
\(\mathcal R_{\Pi}\), and pro-odd extension for
\(\mathcal V_{\Pi}\).

A nonzero continuous quadratic character of the product is nonzero on
some arithmetic axis combination, whereas every original \(X_2\)
character vanishes on all those witnesses. This proves the nonextension
assertion.
\end{proof}

\subsection{Two limitations of a field-only generalization}

First, one inertia factor need not fill the torus. For
\(p=31\), one has \(d=5\), \(r=6\), and
\(\operatorname{rank}_{\mathbf Z_2}\mathcal T_{31}=15\).
One inertia image is a quotient of the unit group of an
unramified degree-five local field and therefore has
pro-two rank at most five. Applying fixed group-ring
operators cannot increase that rank. Thus replacing all
local factors by one inertia group cannot yield the
rank-fifteen target. The product construction is essential.

Second, \(\mathcal P_p\) need not be torsion-free. Set
\[
 e_p=\begin{cases}r,&d\text{ even},\\r/2,&d\text{ odd}.\end{cases}
\]
The sign torsion of \(\mathcal T_p\) is \((C_2)^{e_p}\),
and \(s\) cycles these sign coordinates. Its difference
has image \((C_2)^{e_p-1}\). This is all torsion in
\(\mathcal P_p\): if \((s-1)t\) has torsion image,
then \((s-1)\log t=0\). The invariant vector \(\log t\)
is diagonal and rational, while \((1+c)\log t=0\),
so it is zero. Thus \(t\) itself is sign torsion.
The rational anti-conjugation space has dimension \(g\),
and \(s-1\) is invertible on it. Hence, noncanonically,
\[
                  \mathcal P_p\simeq
                   \mathbf Z_2^g\times(C_2)^{e_p-1}.
\]
For example, at \(p=17\), \(d=8\) and \(r=2\),
so the target has a nontrivial sign-torsion factor.
All simultaneous assertions above retain the exact group
\(\mathcal P_p\); they do not discard this torsion.

\section{The degree-five example and a single inertia group}
\label{sec:genus-two-coefficient-two}
The degree-five case exhibits the integral normalization in the smallest
higher-genus curve. Put \(O=\mathbf Z_2[\zeta_5]\), let
\(\sigma(\zeta_5)=\zeta_5^2\), and write \(\bar z=\sigma^2z\).
The CM type is \(\{1,\sigma\}\), its inverse is
\(1+\sigma^{-1}\), and the integral factor is \(Q_5=-\sigma\).
Thus five satisfies the unit criterion. Set
\[
 T_5^{(2)}=\{z\in1+2O:z\bar z=1\},\qquad
 P_5=(\sigma-1)T_5^{(2)},\qquad
 R_{{\rm quad},5}=R_{\rm quad}\cap\ker\mathfrak h_5,
\]
\[
 V_{{\rm quad},5}=V_{\rm quad}\cap\ker\mathfrak h_5.
\]
There is only one prime above two, so
Proposition~\ref{prop:cm-commutator-torus-torsion} gives
\(P_5\simeq\mathbf Z_2^2\).

\begin{corollary}\label{thm:genus-two-exact-image}
The degree-five homology scalar has exact image \(P_5\) on
\(R_{\rm quad}\), \(V_{\rm quad}\), and both arithmetic intersections.
The four corresponding quotients by its kernel are \(\mathbf Z_2^2\).
Moreover \(G=R_{{\rm quad},5}H\); the full relative quotient on
\(R_{{\rm quad},5}\) is \(B\), and its maximal pro-two version on
\(V_{{\rm quad},5}\) is \(B_{\max2}\). Both original restriction
kernels are exactly \(\mathcal S\), and both images are the full
relative invariant character space \(\operatorname{Hom}_{\rm cts}(B,\F_2)\).
\end{corollary}
\begin{proof}
Apply Theorem~\ref{thm:simultaneous-cyclotomic-cm} with \(\Pi=\{5\}\).
All its hypotheses were verified explicitly above. The torus has free
rank two and no sign-torsion image, giving the stated quotient.
\end{proof}

The arithmetic realization can be made inside one inertia group before
taking the final commutator.
\begin{lemma}\label{lem:genus-two-inertia-M}
Every \(v\in T_5^{(2)}\) is \(\mathfrak h_5(j)\) for a pro-two
arithmetic element \(j\in I_2\cap M\).
\end{lemma}
\begin{proof}
Let \(J\) be the Jacobian of \(t^2-t=w^5\). By
Lemma~\ref{lem:pdl-reflex-and-local}, the exact local CM formula is
\(r(u)=u\sigma^{-1}(u)\) on \(O^\times\). To prescribe \(v\),
choose a unit \(z\) with \(\bar z/z=v\), by unit Hilbert 90.
Set \(u=\sigma^{-1}(z)/z\). Then
\[
 N_{O/\mathbf Z_2}(u)=1,\qquad
 u\sigma^{-1}(u)=\bar z/z=v.
\]
An inertia element \(b\) representing this local unit therefore has
Jacobian value \(v\). The Weil pairing gives
\(\lambda_2(b)=v\bar v=1\); all odd cyclotomic components are one
on inertia at two. Thus \(b\in M\). The based homology value may
have an additional deck factor in \(\mu_5\). Taking
\(j=b^{e_2}\), with the two-primary idempotent exponent, kills this
factor and leaves \(v\) unchanged. It preserves inertia and \(M\),
which proves the assertion.
\end{proof}

For inertia \(a\) at five with deck exponent two, the commutator
\([a,j]\) has scalar \(\sigma(v)/v\). Every Kummer coordinate is
zero because \(j\) is pro-two and \(\lambda_2(a)=1\).
Its primary word actions are trivial: the coefficient-two action kills
\(a\), and every odd-primary IA action kills the pro-two element
\(j\). Both entries fix the two elliptic CM fields, so their corrected
CM commutators are one; the cubic Kummer correction is zero on the
rational-radical kernel. Finally all chosen cup characters extend
invariantly to \(M\), and therefore kill \([a,j]\). This gives a
single-commutator realization in \(V_{\rm quad}\cap H\).
The degree-seven example in Section~\ref{sec:prime-degree-local-cm}
gives a further quotient \(\mathbf Z_2^3\) after imposing this kernel.

\section{An exact global unit criterion}
\label{sec:cm-global-unit-criterion}

Fix a prime \(p\geq5\), write \(h=(p-1)/2\), choose a primitive
residue \(g\) modulo \(p\), and put
\[
 F=\mathbb Q(\zeta_p),\quad F^+=\mathbb Q(\zeta_p+\zeta_p^{-1}),
 \quad \mathcal O=\mathbb Z[\zeta_p]\otimes\mathbb Z_2,\quad
 s=[g],\quad c=s^h.
\]
The algebra \(\mathcal O\) may be a product. All its factors are
included in the following groups:
\[
 T=\{t\in1+2\mathcal O:t\,c(t)=1\},\qquad P=(s-1)T.
\]
We use additive notation for operators on these multiplicative
abelian groups. Let \(Q=Q_p(s)\) be the integral reflex operator and
assume \(Q(1)\) is odd. Set
\[
 \mathscr M=(\mathcal O/2\mathcal O)^c,\qquad
 \beta_k=\overline\zeta_p^{\,g^k}+
                \overline\zeta_p^{-g^k}\quad(0\leq k<h).
\]
These elements form an \(\mathbb F_2\)-basis on which \(s\) acts cyclically. The hypothesis
needed here is the natural finite isomorphism
\begin{equation}
 \theta:T/QT\xrightarrow{\ \sim\ }
               \mathscr M/\overline Q\mathscr M,\qquad
 \theta(t)=(t-1)/2\bmod2\mathcal O.
 \label{eq:uniform-natural-layer}
\end{equation}
An equality of rational ranks is insufficient for this hypothesis.

\begin{lemma}
\label{lem:uniform-operator-inverses}
Under \eqref{eq:uniform-natural-layer}, inclusion identifies
\(P/QP\) with \(T/QT\). On this quotient both \(s-1\) and \(1+[2]\)
are automorphisms.
\end{lemma}
\begin{proof}
The operator \(Q\) acts on \(T/P\) by the odd scalar \(Q(1)\).
It is therefore an automorphism of this pro-\(2\) group. Hence
\(T=P+QT\), and \(Qt\in P\) implies \(t\in P\); consequently
\(QT\cap P=QP\).

The map \(\theta:T\to\mathscr M\) is surjective. Indeed, on each
unramified factor complex conjugation is either the nontrivial
automorphism of an unramified quadratic extension or exchanges
two factors. Thus \(1+c:\mathcal O/2\mathcal O\to\mathscr M\) is
surjective. If \(w+c(w)\equiv v\pmod2\), then
\((1+2w)/(1+2c(w))\) belongs to \(T\) and has image \(v\).
Also \(T[2]\) consists exactly of componentwise signs, equal on
the factors exchanged by \(c\). Since \([2]\) on
\(\mathcal O/2\mathcal O\) is absolute Frobenius, this gives
\begin{equation}
 T[2]\xrightarrow{\ \theta\ }\ker(1+[2]:\mathscr M\to\mathscr M)
 \quad\text{isomorphically}.
 \label{eq:uniform-sign-kernel}
\end{equation}
Let \(b=\dim_{\mathbb F_2}(\mathscr M/\overline Q\mathscr M)\).
The integral anti-invariant lattice \(\mathcal O^-\) has basis
\(\zeta_p^{g^k}-\zeta_p^{-g^k}\), and is
\(\mathbb Z_2[s]/(s^h+1)\). Its reduction is \(\mathscr M\).
The logarithm identifies \(T\otimes\mathbb Q_2\) with
\(\mathcal O^-\otimes\mathbb Q_2\). The Smith normal form therefore
gives
\[
 e:=v_2\det(Q\mid T\otimes\mathbb Q_2)\geq b.
\]
The exact index on the finitely generated abelian pro-\(2\) group
\(T\) is
\begin{equation}
 [T:QT]=|\ker(Q:T[2]\to T[2])|\,2^e.
 \label{eq:uniform-sign-index}
\end{equation}
To see this formula, pass to the torsion-free quotient of \(T\).
The induced endomorphism is injective, its cokernel has order
\(2^e\), and the remaining cokernel is the cokernel of \(Q\)
on the finite torsion group. Kernel and cokernel there have equal
order. Hypothesis \eqref{eq:uniform-natural-layer} forces \(e=b\)
and forces \(Q\) to be invertible on \(T[2]\).

Put \(D=1+[2]\). In the Fitting decomposition
\(\mathscr M=\mathscr M_0\oplus\mathscr M_1\) for \(D\),
the operator \(D\) is nilpotent on the first summand and invertible
on the second. The operator \(\overline Q\) commutes with \(D\).
A nonzero kernel vector of \(\overline Q\) in \(\mathscr M_0\)
would have a last nonzero \(D\)-iterate in
\(\ker D\cap\ker\overline Q\), contrary to
\eqref{eq:uniform-sign-kernel}. Thus the first summand contributes
nothing to the cokernel of \(\overline Q\), and \(D\) is invertible
on that cokernel. Finally, \(\overline Q(1)=1\), so a Bezout
identity for \(X-1\) and \(\overline Q(X)\) makes \(s-1\)
invertible there as well. No semisimplicity is used.
\end{proof}

Write
\[
 \mathscr C=P/QP=T/QT=\mathscr M/\overline Q\mathscr M,
 \qquad \pi:\mathscr M\to\mathscr C.
\]
Let \(\mathscr E=\mathcal O_{F^+}^{\times}\) be the full real
unit group. Define the linear square-class map
\begin{equation}
 \mathscr K:\mathscr C^\vee\longrightarrow
       \mathscr E/(\{\pm1\}\mathscr E^2),\qquad
 \ell\longmapsto[u_\ell],\qquad
 u_\ell=\prod_{k=0}^{h-1}
  (\zeta_p^{g^k}+\zeta_p^{-g^k})^{\ell(\pi\beta_k)}.
 \label{eq:uniform-real-unit-map}
\end{equation}
Each factor is a unit: up to a root of unity it is
\(1+\zeta_p^{2g^k}\), whose norm is \(\Phi_p(-1)=1\).
Square roots in \eqref{eq:uniform-real-unit-map} may be arbitrary
units of \(F^+\).

\begin{theorem}
\label{thm:uniform-arithmetic-cm-image}
Let \(R_\Lambda\) and \(V_\Lambda\) be obtained from
\(R_{\rm quad}\) and \(V_{\rm quad}\) by imposing any collection
of the scalar kernels \(h_q\) with \(q\ne p\).
Under \eqref{eq:uniform-natural-layer}, their arithmetic images are
\begin{equation}
 h_p(R_\Lambda\cap H)=h_p(V_\Lambda\cap H)
  =\{v\in P:\ell(v\bmod QP)=0\text{ for every }
                         \ell\in\ker\mathscr K\}.
 \label{eq:uniform-arithmetic-image}
\end{equation}
In particular,
\[
 [P:h_p(R_\Lambda\cap H)]=2^{\dim\ker\mathscr K},\qquad
 \dim_{\mathbb F_2}\bigl(h_p(R_\Lambda\cap H)/QP\bigr)
                      =\dim\operatorname{im}\mathscr K.
\]
Every value in \eqref{eq:uniform-arithmetic-image} has a witness
in \(V_{\rm quad}\cap H\) with all other prime scalar coordinates
equal to \(1\).
\end{theorem}

\begin{proof}[Proof of Theorem~\ref{thm:uniform-arithmetic-cm-image}]
First, a real unit \(u\in\mathscr E\) is a square in \(F(i)\)
if and only if \(u\in\{\pm1\}\mathscr E^2\).
Indeed \(F(i)/F^+\) is biquadratic: \(F\) is unramified at \(2\),
so it does not contain \(i\). Of its three quadratic subfields,
\(F\) and the third, real subfield are ramified at the prime
above \(p\); \(F^+(i)\) is unramified there. Adjoining a square
root of a unit is unramified at every odd prime. Thus a nontrivial
\(F^+(\sqrt u)\) inside \(F(i)\) must be \(F^+(i)\).
Writing a square root as \(a+bi\) gives \(ab=0\), so \(u=a^2\)
or \(u=-b^2\). The square root is a unit, by valuations.

Let \(\rho:G_F\to\mathcal O^\times\) be the scalar of the compact-homology
Tate representation, and remove the odd Teichmuller factor in
every component to obtain its principal projection
\(\operatorname{pr}_U\rho\). For \(b\in G_{F(i)}\), set
\[
 v_b=\bigl(\operatorname{pr}_U\rho(b)-1\bigr)/2\bmod2\mathcal O.
\]
Polarization and \(\lambda_2(b)\equiv1\pmod4\) imply
\(v_b\in\mathscr M\). This is a continuous additive character.
For \(\ell\in\mathscr C^\vee\), define
\[
 \mathscr D(\ell)=
 \left[\prod_{k=0}^{h-1}
   (1-\zeta_p^{g^k})^{\ell(\pi\beta_k)}\right]
          \in F(i)^\times/(F(i)^\times)^2.
\]
We claim
\begin{equation}
 \ell(\pi v_b)=\kappa_{\mathscr D(\ell)}(b),
 \label{eq:uniform-jacobi-kummer}
\end{equation}
where \(\kappa\) is the quadratic Kummer character.

It suffices to check Frobenius elements at primes of \(F(i)\)
away from \(2p\). Let \(q\) be the full residue-field cardinality,
\(m=(q-1)/p\), and let \(\chi\) be the residue character with
\(\chi(x)\equiv x^m\) at that prime and \(\chi(0)=0\).
Then \(q\equiv1\pmod{4p}\) and \(m\) is even. The
character comparison in Theorem~8.1, Remark~8.3, and
Corollary~8.4 of \cite{Otsubo2016} applies over the number field
\(F(i)\), with auxiliary residue degree \(1\), and gives the
negative Jacobi sum
\[
 j=-\sum_{x\in\mathbb F_q}\chi(x)\chi(1-x).
\]
The substitution \(y=2x-1\) gives
\[
 \chi(4)j=-\sum_{y\in\mathbb F_q}\chi(1-y^2)
        =-1-2\sum_{y\in\mathbb F_q^\times/\{\pm1\}}
                                      \chi(1-y^2).
\]
This element is \(1\) modulo \(2\), so \(\chi(4)^{-1}\) is
exactly the odd Teichmuller factor of \(j\). Therefore
\[
 v_{\operatorname{Frob}}\equiv
 1+\sum_{y\in\mathbb F_q^\times/\{\pm1\}}\chi(1-y^2)
                                                   \pmod2.
\]
For \(a\ne0\), let \(n_a\) count the pairs contributing
\(\zeta_p^a\). Equivalently, it counts roots of \(z^m=\zeta_p^a\)
for which \(1-z\) is a nonzero square. All \(m\) roots lie in
\(\mathbb F_q\), none is \(1\), and
\[
 \prod_{z^m=\zeta_p^a}(1-z)=1-\zeta_p^a.
\]
The quadratic residue sign of this product is
\((-1)^{m-n_a}=(-1)^{n_a}\). Since \(-1\) and \(\zeta_p\)
are squares in \(\mathbb F_q\), opposite exponents have the
same parity. Thus
\[
 v_{\operatorname{Frob}}=
 \epsilon\cdot1+\sum_{k=0}^{h-1}
     \kappa_{1-\zeta_p^{g^k}}(\operatorname{Frob})\beta_k
       \quad\text{in }\mathscr M.
\]
The constant term is killed by \(\pi\), because
\(\overline Q(1)=1\). Applying \(\ell\pi\) proves
\eqref{eq:uniform-jacobi-kummer} on Frobenius elements;
Chebotarev proves it on all of \(G_{F(i)}\).

Put \(d=[2]\) as an operator, and \(L_0=d^{-2}(1+d)\).
The identity
\[
 \zeta_p^a+\zeta_p^{-a}
  =\zeta_p^{-a}\frac{1-\zeta_p^{4a}}{1-\zeta_p^{2a}}
\]
gives, in square classes of \(F(i)\),
\begin{equation}
 [u_\ell]=\mathscr D(\ell\circ L_0).
 \label{eq:uniform-dual-reindexing}
\end{equation}
The inverse exponents in this formula are essential:
the coefficient at \(1-\zeta_p^a\) is
\(\ell(\pi\beta_{a/4})+\ell(\pi\beta_{a/2})\).
Replacing an exponent by its negative changes the factor
only by a square in \(F(i)\).

The kernel of \(\mathscr D\) is stable under precomposition
by \(s\), since
\(\mathscr D(\ell\circ s)=s^{-1}\mathscr D(\ell)\).
It is therefore stable under every polynomial in \(s\).
Lemma~\ref{lem:uniform-operator-inverses} makes \(L_0\)
and \(s-1\) invertible on \(\mathscr C\). Their dual actions
preserve this finite kernel and are bijective on it.
The real-unit square test and \eqref{eq:uniform-dual-reindexing}
consequently give
\[
 \ker\mathscr K=\ker\mathscr D.
\]
Finite-dimensional duality in \eqref{eq:uniform-jacobi-kummer}
now shows that
\begin{equation}
 \{\pi v_b:b\in G_{F(i)}\}
       =W:=(\ker\mathscr K)^\perp\subseteq\mathscr C.
 \label{eq:uniform-image-annihilator}
\end{equation}
The subspace \(W\) is stable under \(s\) and under
\((s-1)^{-1}\).

Every element of \(R_\Lambda\cap H\) or \(V_\Lambda\cap H\)
fixes \(F(i)\), and its principal Tate scalar is \(h_p\).
This proves the upper bound in
\eqref{eq:uniform-arithmetic-image}. For the reverse bound,
use the established mixed-inertia construction. Choose primitive
\(p\)-inertia \(a\), and take \(b\in G_{F(i)}\). The word
\[
 w_b=[a,[a,b]^{e_2}]
\]
belongs to \(V_{\rm quad}\cap H\), and its scalar coordinates
in every degree \(q\ne p\) are \(1\). In degree \(p\),
the scalar of \([a,b]\) is \((s-1)\rho(b)\).
The pro-\(2\) power removes its odd part; the outer commutator
applies \(s-1\) once more. Hence
\[
 h_p(w_b)\bmod QP=(s-1)^2\pi v_b.
\]
By \eqref{eq:uniform-image-annihilator}, these words realize
all of \(W\) modulo \(QP\). The local norm construction realizes
every element of \(QP\), again with every other coordinate
equal to \(1\). Multiplication gives every value in the claimed
preimage, with all other coordinates equal to \(1\).
Such witnesses survive any collection, finite or infinite,
of other-degree scalar kernels. This proves the exact image
and both dimension formulas.
\end{proof}


\section{A global Frobenius fills the degree-twenty-nine deficit}
\label{sec:global-frobenius-29}

Put \(\zeta=\zeta_{29}\), \(O=\mathbf Z_2[\zeta]\), and let
\(s(\zeta)=\zeta^2\). Since two has order twenty-eight modulo
\(29\), the ring \(O\) is unramified of degree twenty-eight.
Write \(c=s^{14}\), and put
\[
 U=1+2O,\qquad T=\{u\in U:u\,c(u)=1\},\qquad P=(s-1)T.
\]
The explicit CM-type operator is
\[
 Q=-s+s^3-s^4+s^5-s^7+s^{10}-s^{11}.
\]
The preceding integral calculation gives
\[
 (s-1)\Phi^*=(c-1)Q,\qquad Q(1)=-1,\qquad [T:QT]=8.
\]
Let \(\mathcal R_\Pi,\mathcal V_\Pi\) denote the established simultaneous kernels
for the primes satisfying the integral unit criterion. In particular
\(29\notin\Pi\), \(G=\mathcal R_\Pi H\), and \(\mathcal R_\Pi/\mathcal V_\Pi\) is
pro-odd. The compact-homology scalar in degree twenty-nine is
\(\mathfrak h_{29}\).

\begin{theorem}\label{thm:global-frobenius-29}
All four images are the whole target:
\[
 \mathfrak h_{29}(\mathcal R_\Pi)=\mathfrak h_{29}(\mathcal R_\Pi\cap H)
 =\mathfrak h_{29}(\mathcal V_\Pi)=\mathfrak h_{29}(\mathcal V_\Pi\cap H)=P.
\]
In this inert case \(P\simeq\mathbf Z_2^{14}\).
For
\[
 \mathcal R_{\Pi,29}=\mathcal R_\Pi\cap\ker\mathfrak h_{29},\qquad
 \mathcal V_{\Pi,29}=\mathcal V_\Pi\cap\ker\mathfrak h_{29},
\]
one has \(G=\mathcal R_{\Pi,29}H\). Both new kernels have infinite
index, with exact quotient \(P\), already on their arithmetic
intersections. They preserve respectively the full relative quotient
\(B\) and its maximal pro-two quotient. The original invariant
quadratic-character restriction kernel remains exactly \(\mathcal S\), and
its image on either new kernel is exactly the invariant quadratic
characters killing the arithmetic intersection, canonically
\(\operatorname{Hom}_{\rm cts}(B,\mathbf F_2)\).
\end{theorem}

\subsection{The exact missing finite quotient}

The torus coinvariant calculation gives \(T/P\simeq C_2^2\),
and \(s\) acts trivially there. Thus \(Q\) induces multiplication
by \(Q(1)=-1\), an automorphism. Inclusion induces an isomorphism
\begin{equation}\label{eq:g29-cokernel-identification}
                         P/QP\xrightarrow{\sim}T/QT.
\end{equation}
Indeed, if \(Qt\in P\), invertibility on \(T/P\) forces
\(t\in P\); this proves injectivity. For any \(t\in T\),
choose \(u\) with \(Q[u]=[t]\) in \(T/P\), and then
\(t/(Qu)\in P\), proving surjectivity.

Set \(T(4)=T\cap(1+4O)\). The first principal-unit layer is
\[
 \theta:T\longrightarrow \mathbf F_{2^{14}},\qquad
             t\longmapsto\frac{t-1}{2}\pmod2,
 \qquad \ker\theta=T(4).
\]
The norm condition puts the image in the fixed field of \(c\).
This map is onto: lift any element of that fixed field to \(a\in O\).
The norm of \(1+2a\) lies in \(1+4O^c\). The norm on
\(1+4O\) is onto \(1+4O^c\), by logarithm and surjectivity
of the trace in the unramified quadratic extension. Correcting by
an element of \(1+4O\) therefore produces a norm-one lift with
the specified first layer.

The fourteen elements
\[
                   \beta_j=\zeta^j+\zeta^{-j},\qquad 1\le j\le14,
\]
are an \(\mathbf F_2\)-basis of this fixed field. They are independent
because \(\zeta,\ldots,\zeta^{28}\) are a basis of
\(O/2O\), and these sums have disjoint pairs of supports.
Frobenius cycles the normal basis \(\beta_{2^k}\), \(0\le k<14\),
where indices are read modulo twenty-nine and up to sign. Thus, as an
\(\mathbf F_2[s]\)-module,
\[
                    \mathbf F_{2^{14}}\simeq
                         \mathbf F_2[s]/(s^{14}+1).
\]
The exact polynomial calculation is
\[
 \gcd(\overline Q,s^{14}+1)=f(s),\qquad f(s)=s^3+s^2+1.
\]
Consequently
\(\mathbf F_{2^{14}}/\overline Q\mathbf F_{2^{14}}
\simeq\mathbf F_2[s]/(f)\), of order eight.
The induced map from \(T/QT\) is surjective; since its domain also
has order eight, it is an isomorphism. Together with
\eqref{eq:g29-cokernel-identification}, this identifies the missing
quotient exactly as
\begin{equation}\label{eq:g29-layer-certificate}
                         P/QP\simeq\mathbf F_2[s]/(s^3+s^2+1).
\end{equation}
In particular it has exponent two. The cubic polynomial is irreducible,
so the \(s\)-orbit of every nonzero element spans the quotient.

\subsection{The arithmetic word allows arbitrary CM Frobenius}

Choose \(a\in I_{29}\) with \(\lambda(a)\bmod29=2\).
For any \(b\in G_{\mathbf Q(\zeta)}\), put
\[
                       z=[a,b],\qquad z_2=z^{e_2},\qquad w=[a,z_2],
\]
where \(e_2\in\widehat{\mathbf Z}\) is one at two and zero at
all other primes, with exponentiation in the procyclic closure.
The element \(b\) is not required to be an inertia element.

Both \(z\) and \(z_2\) have full multiplier one. The
coefficient-two Kummer cocycle of two vanishes on \(a\in I_{29}\),
and \(\lambda_2(a)=1\). The cocycle commutator formula therefore
gives \(\rho_{2,2}^I(z)=0\). Primary projection kills every other
coefficient, so \(z_2\in N\cap H\). Every rational Kummer
coordinate of \(z_2\) is supported at two, and hence
\[
                 \rho_q^I(w)=(\lambda(a)-1)\rho_q^I(z_2)=0
                         \quad\text{for all rational primes }q.
\]
Thus \(w\in D\cap H\).

At coefficient two the based action of \(a\) is trivial, by
unramifiedness outside two, so \(\Phi_2(z_2)=\Phi_2(w)=1\).
For every odd prime \(\ell\), the normalized image of a
multiplier-one element lies in the pro-\(\ell\) kernel of
\(\operatorname{Aut}(\widehat F_2^{(\ell)})\to
\operatorname{GL}_2(\mathbf F_\ell)\). The image of the pro-two
element \(z_2\) in that group is therefore trivial.
For completeness, the pro-\(\ell\) assertion follows on finite
characteristic quotients from the following elementary fact. The
Frattini kernel in the automorphism group of a finite \(\ell\)-group
acts freely on the set of generating tuples lifting a fixed basis of
the Frattini quotient. That nonempty set has \(\ell\)-power cardinality,
so the automorphism kernel is an \(\ell\)-group. Passing to
characteristic quotients proves the assertion.
Hence all primary based word actions kill \(w\).

The element \(a\) fixes both elliptic CM fields
\(\mathbf Q(i)\) and \(\mathbf Q(\sqrt{-3})\), as does \(z_2\).
Their elliptic Tate actions therefore commute. The corrected harmonic
frame of \(w\) is one, and so is the corrected cubic frame, since
its quartic Kummer correction of three vanishes on \(D\).
The cubic radical--CM defect vanishes arithmetically. Finally
\[
                 \chi(w)=\chi(a z_2a^{-1})-\chi(z_2)=0
                              \quad(\chi\in X_2).
\]
All the quadratic coordinates defining \(V_{\rm quad}\) vanish.
Thus \(w\in V_{\rm quad}\cap H\).

For an odd prime \(q\ne29\), both \(a\) and \(z_2\) have
trivial degree-\(q\) deck exponent. Their compact-homology actions
lie in the commutative deck centralizer. Consequently
\(\mathfrak h_q(w)=1\) for every such \(q\), and
\begin{equation}\label{eq:g29-word-membership}
                               w\in\mathcal V_\Pi\cap H.
\end{equation}
This verifies all previous kernels for arbitrary \(b\) fixing
\(\mathbf Q(\zeta)\).

If \(\alpha=\rho_{J_{29},2}(b)\), the Jacobian value of the first
commutator is \((s-1)\alpha\). The based-cover/Jacobian comparison
may have a deck \(C_{29}\) discrepancy; it vanishes after pro-two
projection. The full cover representation then gives exactly
\begin{equation}\label{eq:g29-word-value}
                  \mathfrak h_{29}(w)=\bigl((s-1)^2\alpha\bigr)^{e_2}.
\end{equation}

\subsection{The Frobenius at fifty-nine and its exact normalization}

The Fermat curve \(u^{29}+v^{29}=1\) maps to \(C_{29}\) by
\[
                             (u,v)\longmapsto(t=u^{29},w=-uv).
\]
It is the quotient by \((u,v)\mapsto(\zeta u,\zeta^{-1}v)\).
At coefficient two, averaging over this order-twenty-nine subgroup
identifies compact homology of the quotient with its invariants,
integrally. The surviving Fermat character components are exactly
\((a,a)\), \(a\ne0\). The degree-twenty-nine deck action on the
quotient is induced by \((u,v)\mapsto(\zeta u,v)\), so the
\((1,1)\) component has deck eigenvalue \(\zeta\).

By \cite[Theorem 8.1 and Remark 8.3, pp.~638--639]{Otsubo2016},
arithmetic Frobenius acts on this homology component, equivalently
on the Jacobian Tate module, by
\[
 j_{29}^{1,1}(\mathfrak l)
       =-\sum_{x\in\mathbf F_{59}\setminus\{0,1\}}
                             \chi_{\mathfrak l}(x)\chi_{\mathfrak l}(1-x).
\]
The cited remark explicitly distinguishes geometric Frobenius,
cohomology, and homology; the scalar here is the displayed Jacobi
sum, not its inverse. Choose \(\mathfrak l\mid59\) with
\(\zeta\equiv4\pmod{\mathfrak l}\). Since two is primitive
modulo fifty-nine, the residue character satisfies
\(\chi_{\mathfrak l}(2)=\zeta\). Choose \(b\) lifting arithmetic
Frobenius at this prime. It fixes \(\mathbf Q(\zeta)\), and its
coefficient-two scalar is exactly
\begin{equation}\label{eq:g29-jacobi}
 \alpha=\zeta^{27}
       +2\sum_{j\in\{4,5,9,13,14,24,26\}}\zeta^j
       -2\sum_{j\in\{7,10,17,19,22,23,25\}}\zeta^j.
\end{equation}
For a direct finite verification, the counts of the exponents
\(\log_2x+\log_2(1-x)\pmod{29}\), in order from zero to
twenty-eight, are
\[
\begin{gathered}
(2,2,2,2,0,0,2,4,2,0,4,2,2,0,0,\\
2,2,4,2,4,2,2,4,4,0,4,0,1,2).
\end{gathered}
\]
Subtracting the final coefficient with the relation
\(\sum_{j=0}^{28}\zeta^j=0\), and taking the minus sign,
gives \eqref{eq:g29-jacobi}. Exact multiplication also gives
\(\alpha\bar\alpha=59\).

Set \(\delta=\alpha\zeta^2\). Then \(\delta\in1+2O\), and
\[
              (s-1)^2\alpha=\zeta^{-2}(s-1)^2\delta.
\]
The first factor has odd order, while the second is pro-two.
Equation \eqref{eq:g29-word-value} consequently becomes the exact
identity \(\mathfrak h_{29}(w)=(s-1)^2\delta\).
Writing \(\delta=1+2r\), its first layer is
\[
 \theta(\mathfrak h_{29}(w))=r+s^2r
                   =\beta_2+\beta_{11}+\beta_{12}+\beta_{13}.
\]
The last equality follows by substituting
\eqref{eq:g29-jacobi} and reducing modulo four. In the cyclic
basis \(\beta_{2^k}\), it is the polynomial
\[
                   A(s)=s+s^4+s^7+s^{11}.
\]
Its remainder modulo \(s^3+s^2+1\) is \(s+1\), which is
nonzero. Thus this arithmetic Frobenius word has a
nonzero image in \(P/QP\). Its first three conjugates by \(a\)
have classes
\[
                        s+1,\qquad s^2+s,\qquad 1,
\]
which form a basis of that three-dimensional quotient.

\subsection{Saturation and the new relative kernel}

The local product-inertia construction already realizes \(QP\)
on \(\mathcal V_\Pi\cap H\): its exact value is
\((s-1)^2\Phi^*(U)=(s-1)QT=QP\), and its other
prime-degree coordinates are one. The Frobenius word just
constructed and its conjugates supply every class in \(P/QP\).
Therefore \(\mathfrak h_{29}(\mathcal V_\Pi\cap H)=P\).
The full-group upper bound \(\mathfrak h_{29}(\mathcal R_\Pi)\subset P\)
proves all four image equalities in the theorem.

Arithmetic correction in this common image proves
\(G=\mathcal R_{\Pi,29}H\), and ordinary invariant-character gluing
preserves the full relative quotient \(B\). The quotient
\(\mathcal R_{\Pi,29}/\mathcal V_{\Pi,29}\) embeds in the old pro-odd quotient;
central coprime extension therefore preserves \(B_{\max2}\).
Every arithmetic element used to fill \(P\) is killed by every
original invariant quadratic character. Thus an original character
vanishing on a new kernel, and hence factoring through \(P\)
on the old subgroup, already vanishes on that old subgroup.
Its original restriction kernel is exactly \(\mathcal S\).
The complete quadratic splitting, gluing on \(\mathcal R_{\Pi,29}\), and
pro-odd extension for \(\mathcal V_{\Pi,29}\) give the stated complete
restricted-image description. Finally \(P\) is torsion-free of
rank fourteen in this single inert factor, so both reductions are
strict and have infinite index.

\section{Global saturation in degree \(113\)}
\label{sec:global-cm-113}

Put \(p=113\), \(h=56\), and choose \(s=[3]\).
The order of \(2\) modulo \(113\) is \(28\), with \(2=3^{12}\).
Thus
\(\mathcal O=\mathbb Z[\zeta_{113}]\otimes\mathbb Z_2\)
is a product of four unramified degree-\(28\) integer rings.
Complex conjugation \(c=s^{56}\) acts within each factor
as its degree-\(14\) Frobenius. Write
\[
 T=\{u\in1+2\mathcal O:u\,c(u)=1\},\qquad P=(s-1)T.
\]
Then \(T\) has rank \(56\) and torsion \((C_2)^4\).
The generator \(s\) permutes the four signs cyclically.
Since \(s-1\) is invertible on the rational anti-invariant
space, a torsion element in its image must come from a
torsion element. Consequently
\begin{equation}
 P\simeq\mathbb Z_2^{56}\times(C_2)^3.
 \label{eq:113-full-target}
\end{equation}

The integral reflex operator in this primitive coordinate is
\[
\begin{split}
 Q={}&-s^2+s^5-s^6+s^{13}-s^{15}+s^{16}-s^{19}
       +s^{21}-s^{22}+s^{23}\\
    &-s^{24}+s^{26}-s^{30}+s^{37}-s^{39}
       +s^{45}-s^{47}+s^{52}-s^{53}.
\end{split}
\]
In particular \(Q(1)=-1\). On the four sign factors its
reduction is
\begin{equation}
 \overline Q=1+s^2+s^3\pmod{s^4-1},\qquad
 (1+s^2+s^3)^{-1}=1+s+s^2.
 \label{eq:113-sign-inverse}
\end{equation}
All signs therefore belong to the local image.
On the signed cyclic lattice \(\mathbb Z_2[s]/(s^{56}+1)\),
exact fraction-free elimination gives
\[
 \det Q=1028233560187751537651440165000,\qquad
 v_2(\det Q)=3.
\]
Also the Euclidean algorithm over \(\mathbb F_2[s]\) gives
\begin{equation}
 \gcd(\overline Q,s^{56}+1)=s^3+s^2+1.
 \label{eq:113-residual-polynomial}
\end{equation}
The sign calculation and the torsion/free index formula imply
\[
 [T:QT]=[P:QP]=8.
\]
The natural first-layer map is surjective and its target has
dimension three by \eqref{eq:113-residual-polynomial};
hence it is an isomorphism. Inclusion identifies \(P/QP\)
with \(T/QT\), because \(Q\) is multiplication by \(-1\)
on \(T/P\). Thus the quotient is the additive group
\begin{equation}
 \mathscr C=\mathbb F_2[\xi]/(\xi^3+\xi^2+1),
 \qquad s\text{ acts by multiplication by }\xi.
 \label{eq:113-residual-field}
\end{equation}

For the Frobenius calculation it is convenient to use
\[
 L(u)=\log(u)/2\bmod2\mathcal O
     =a+a^2\quad\text{if }u=1+2a.
\]
On \(T\) this is \((1+[2])\theta=(1+s^{12})\theta\).
Since \(\xi\) has order seven, \(1+\xi^{12}\ne0\).
Consequently \(L\), followed by the cyclic quotient
\[
 \beta_k=\overline\zeta_{113}^{\,3^k}
            +\overline\zeta_{113}^{-3^k}\longmapsto\xi^k,
 \qquad0\leq k<56,
\]
also gives an isomorphism \(P/QP\to\mathscr C\).
This logarithmic detector does not remove an unresolved
sign quotient: \eqref{eq:113-sign-inverse} has already
included every finite sign.

\begin{theorem}
\label{thm:113-global-saturation}
After any collection of the other-prime scalar kernels,
the degree-\(113\) images of \(R_{\rm quad}\), \(V_{\rm quad}\),
and their arithmetic intersections are all \(P\).
In particular degree \(113\) can be added to the previously
normalized primes \(29\) and \(\Pi\), preserving the full
relative quotient on the \(R\)-group, its maximal pro-\(2\)
quotient on the \(V\)-group, and the original quadratic
restriction kernel \(\mathcal S\) and full relative image.
The new scalar kernel has infinite index on both groups.
\end{theorem}

\begin{proof}
The prime \(227=2\cdot113+1\) splits completely in
\(\mathbb Q(\zeta_{113})\). The residue \(2\) is primitive
modulo \(227\). Choose \(\chi:\mathbb F_{227}^{\times}
\to\mu_{113}\) by \(\chi(2)=\zeta_{113}\), and extend it
by \(\chi(0)=0\). Put
\[
 J=-\sum_{x\in\mathbb F_{227}}\chi(x)\chi(1-x).
\]
The positive residue-character convention corresponds to
the prime at which \(\zeta_{113}\) reduces to \(4\).
The character-by-character Frobenius comparison
\cite[Theorem~8.1, Remark~8.3, Corollary~8.4]{Otsubo2016}
identifies this Jacobi sum with the degree-\(113\) Tate scalar,
with all four dyadic factors labelled compatibly.
In the inverse-character convention of
\cite[definition on p.~201; Lemma~5.1 and its proof]{AnglesNuccio},
the corresponding residue is \(4^{-1}=57\); this changes
the common label by conjugation and does not change the
image argument.

For an explicit integral certificate, write
\(J=-\sum_{j=0}^{112}n_j\zeta_{113}^j\). The counts
\(n_j=\#\{x\ne0,1:\chi(x(1-x))=\zeta_{113}^j\}\)
are obtained by listing the powers of \(2\) in
\(\mathbb F_{227}^{\times}\). Their sum is \(225\).
Their cyclic autocorrelation has constant coefficient
\(673\) and every other coefficient \(446\), giving
\(Jc(J)=673-446=227\).
The unique odd \(n_j\) is \(n_{111}=1\).
Thus
\[
 \delta=\zeta_{113}^{\,2}J\in1+2\mathcal O.
\]
Its norm is the rational unit \(227\). For a rational
\(2\)-adic unit \(v\), \(\log(v)/2\equiv0\pmod2\);
this follows directly from \(a+a^2=0\) for
\(a\in\mathbb F_2\). Therefore \(L(\delta)\) is
\(c\)-invariant, although the raw first layer of \(\delta\)
need not be.

In the cyclic basis \(\beta_k\), the support of
\(L(\delta)\) is the following subset of \(\{0,\ldots,55\}\):
\[
\begin{split}
 \mathcal I=\{&
 0,1,2,5,7,12,14,15,16,17,18,20,22,24,\\
 &26,28,33,34,38,39,40,41,43,45,46,49,50,55\}.
\end{split}
\]
Reduction by \(\xi^3+\xi^2+1\) consequently gives
\begin{equation}
 \pi L(\delta)=\sum_{k\in\mathcal I}\xi^k
             =1+\xi+\xi^2.
 \label{eq:113-jacobi-detector}
\end{equation}
This certificate uses only arithmetic in the two finite
rings \(\mathbb F_{227}\) and
\(\mathbb F_2[\xi]/(\xi^3+\xi^2+1)\).

Choose \(a\) in \(113\)-inertia inducing \(s\), and let \(b\)
be a Frobenius lift at the indicated prime over \(227\).
Set
\[
 z=[a,b]^{e_2},\qquad w=[a,z].
\]
The established mixed-inertia construction places \(w\)
in \(V_{\rm quad}\cap H\); it also gives \(h_q(w)=1\)
for every \(q\ne113\). These assertions are simultaneous.
For example, at another prime degree the two elements
in the outer commutator have linear rank-one scalar
actions, which commute. The construction also kills the
earlier based-word, radical, cubic, and invariant
quadratic-character coordinates.

In degree \(113\), the inner commutator applies \(s-1\)
to the Frobenius scalar, the pro-\(2\) power removes its
odd part, and the outer commutator applies \(s-1\) again.
Equation \eqref{eq:113-jacobi-detector} yields
\[
 \pi L(h_{113}(w))
    =(\xi-1)^2(1+\xi+\xi^2)=1.
\]
The corresponding values on \(awa^{-1}\) and
\(a^2wa^{-2}\) are \(\xi\) and \(\xi^2\).
These three elements span the residual quotient.
The local norm axes already give \(QP\), including the
signs in \eqref{eq:113-sign-inverse}. Multiplication
therefore gives every element of \(P\), by arithmetic
elements with all other scalar coordinates equal to \(1\).
The four image equalities follow.

On the preceding normalized groups \(R_0,V_0\), write
\(R_1=R_0\cap\ker h_{113}\) and
\(V_1=V_0\cap\ker h_{113}\). Matching a scalar value
by an element of \(V_0\cap H\) gives
\[
 R_0=R_1(V_0\cap H),\qquad
 R_1/V_1\simeq R_0/V_0,\qquad G=R_1H.
\]
The same matching preserves the arithmetic image in the
pro-odd quotient. Extension of invariant relative
characters along \(G=R_1H\), and coprime extension along
\(R_1/V_1\), give the stated relative quotient and original
restriction consequences. Finally
\eqref{eq:113-full-target} has positive rank, so the
new kernels have infinite index.
\end{proof}

\section{The degree-\(163\) arithmetic image}
\label{sec:arithmetic-cm-163}

Put \(p=163\), \(s=[2]\), \(h=81\), and \(c=s^{81}\).
The residue \(2\) has order \(162\) modulo \(163\);
therefore
\(\mathcal O=\mathbb Z_2[\zeta_{163}]\) is unramified of
degree \(162\). The group
\[
 T=\{u\in1+2\mathcal O:u\,c(u)=1\}
\]
has rank \(81\) and torsion \(\{\pm1\}\).
Since \(s-1\) kills this torsion and is invertible on the
rational anti-invariant space,
\[
 P=(s-1)T\simeq\mathbb Z_2^{81}.
\]
For a reproducible description of the reflex polynomial,
write \(Q=\sum_{j\in I_+}s^j-\sum_{j\in I_-}s^j\), where
\[
\begin{split}
 I_+={}&\{2,4,9,13,16,18,20,25,27,31,35,42,44,49,
                                     52,54,57,67,70,73\},\\
 I_-={}&\{1,3,5,11,15,17,19,22,26,28,32,38,43,47,50,
                                  53,56,62,68,71,75\}.
\end{split}
\]
Then \(Q(1)=-1\), and exact elimination on the signed
cyclic lattice gives
\[
 \det Q=40176902897933369663096255063282475371772362696316,
 \qquad v_2(\det Q)=2.
\]
The corresponding mod-two calculation is
\[
 \gcd(\overline Q,s^{81}+1)=s^2+s+1.
\]
Because \(Q\) is an automorphism on the only sign factor,
\begin{equation}
 [T:QT]=[P:QP]=4,\qquad
 T/QT\xrightarrow{\ \theta\ }
      (\mathcal O/2\mathcal O)^c/\overline Q
                       (\mathcal O/2\mathcal O)^c
 \quad\text{is an isomorphism}.
 \label{eq:163-first-layer}
\end{equation}
The last assertion follows from surjectivity of the
first-layer map and equality of the two orders.
Thus
\[
 \mathscr C=P/QP\simeq
     \mathbb F_2[s]/(s^2+s+1),
\]
with the irreducible cyclotomic \(C_3\)-action.

\begin{lemma}
\label{lem:163-period-unit-square}
Let \(H_0=\langle8\rangle\subset(\mathbb Z/163)^\times\),
and put
\[
 \theta_i=\sum_{a\in2^iH_0}\zeta_{163}^a,\qquad i=0,1,2.
\]
Then
\begin{equation}
 \prod_{j=0}^{26}
      (\zeta_{163}^{8^j}+\zeta_{163}^{-8^j})
            =-(5+\theta_1)^2.
 \label{eq:163-period-square}
\end{equation}
The element \(5+\theta_1\) is a real unit.
The two conjugate identities hold for the other two cosets.
\end{lemma}
\begin{proof}
The subgroup \(H_0\) has order \(54\), contains \(-1\),
and its three cosets partition the nonzero residues.
Thus the periods are real and
\(\theta_0+\theta_1+\theta_2=-1\).
Counting sums of two residues in these cosets gives
\begin{equation}
\begin{aligned}
 \theta_0^2&=54+20\theta_0+17\theta_1+16\theta_2,\\
 \theta_1^2&=54+16\theta_0+20\theta_1+17\theta_2,\\
 \theta_0\theta_1&=17\theta_0+16\theta_1+21\theta_2.
\end{aligned}
 \label{eq:163-period-table}
\end{equation}
Here every coefficient is an integer residue count.
For the product in \eqref{eq:163-period-square}, one can
work entirely in \(\mathbb Z[X]/(X^{163}-1)\), starting
with \(C_0(X)=1\) and using
\[
 C_{j+1}(X)=C_j(X)(X^{8^j}+X^{-8^j})
                     \pmod{X^{163}-1},\qquad0\leq j<27.
\]
The final coefficient is constant on each of
\(\{0\},H_0,2H_0,4H_0\), with respective values
\[
 823364,\qquad823427,\qquad823413,\qquad823426.
\]
These integers give
\[
 C_{27}(\zeta_{163})=-62+\theta_0-13\theta_1.
\]
The second line of \eqref{eq:163-period-table}, together
with the sum of the periods, gives
\[
 -(5+\theta_1)^2
  =-25-10\theta_1-\theta_1^2
  =-62+\theta_0-13\theta_1.
\]
This proves the identity. Every factor on its left is
a unit, since it is a root of unity times
\(1+\zeta_{163}^{2a}\) and \(\Phi_{163}(-1)=1\).
The algebraic integer \(5+\theta_1\) is consequently a
unit as well. Applying \(s\) and \(s^2\) proves the
conjugate identities.
\end{proof}

\begin{theorem}
\label{thm:163-exact-arithmetic-image}
After imposing any collection of the scalar kernels
\(h_q\) with \(q\ne163\), the degree-\(163\) arithmetic
images of the resulting \(R\)- and \(V\)-groups are
both exactly
\[
 Q_{163}P_{163},\qquad [P_{163}:Q_{163}P_{163}]=4.
\]
This is an unconditional integral image statement.
\end{theorem}
\begin{proof}
The exact first-layer hypothesis for
Theorem~\ref{thm:uniform-arithmetic-cm-image} holds by
\eqref{eq:163-first-layer}.
Write
\(\xi=s\bmod(s^2+s+1)\) in \(\mathscr C\).
For \(\ell\in\mathscr C^\vee\), the exponents
\(\ell(\xi^k)\) in the real-unit map are periodic of
period three. Put
\[
 U_i=\prod_{j=0}^{26}
  \bigl(\zeta_{163}^{2^{i+3j}}
                   +\zeta_{163}^{-2^{i+3j}}\bigr),
 \qquad i=0,1,2.
\]
The unit assigned to \(\ell\) is
\[
 u_\ell=U_0^{\ell(1)}U_1^{\ell(\xi)}U_2^{\ell(\xi^2)}.
\]
Lemma~\ref{lem:163-period-unit-square} and its conjugates
make every \(U_i\) a negative square of a real unit.
Thus \(\mathscr K_{163}\) is the zero map into real units
modulo signs and squares. Its kernel is all of
\(\mathscr C^\vee\), whose annihilator is zero.
The exact uniform arithmetic-image theorem now gives
the claimed image \(QP\). That theorem also provides
witnesses for every element of \(QP\) with all other
prime scalar coordinates equal to \(1\), so the equality
persists under any collection of the other-prime kernels.
\end{proof}

On a normalized pair \(R_0,V_0\) with pro-odd
\(R_0/V_0\), the full images of \(h_{163}\) coincide:
their quotient is both pro-odd and pro-\(2\), hence is
trivial. The common image contains \(QP\).
Its quotient in \(P/QP\) is stable under the full
cyclotomic action and this module is irreducible of
dimension two. Thus the common full image is either
\(QP\) or \(P\). The arithmetic equality alone does not
select either alternative.

\begin{lemma}
\label{lem:163-integral-pair-decomposition}
There is a cyclotomic-stable decomposition
\[
 P=B_2\oplus L_{79},\qquad QP=2B_2\oplus L_{79},
 \qquad
 \operatorname{rank}_{\mathbb Z_2}B_2=2,\quad
 \operatorname{rank}_{\mathbb Z_2}L_{79}=79.
\]
Moreover, if \(t=-s\) acts on the additive
\(\mathbb Z_2\)-lattice \(P\), then \(P\) is a free
rank-one \(\mathbb Z_2[C_{81}]\)-module.
\end{lemma}
\begin{proof}
Complex conjugation acts on the norm-one torus as
inversion, so \(s^{81}=-1\) on the additive lattice.
Consequently \(t=-s\) has order \(81\).
The rational representation is the signed cyclic
representation for \(s\), hence the regular cyclic
representation for \(t\), with characteristic
polynomial \(X^{81}-1\).

Since \(81\) is odd, \(\mathbb Z_2[C_{81}]\) is finite
etale over \(\mathbb Z_2\), a product of unramified
discrete valuation rings. The lattice \(P\) decomposes
under its integral idempotents. Each component is
torsion-free of rank one over its corresponding
valuation ring, because the rational representation
is regular. Each component is therefore free of
rank one, proving the module assertion.

Modulo \(2\), the only factor on which \(Q\) is
noninvertible is \(X^2+X+1\). It defines an unramified
degree-two component \(B_2\).
The exact determinant valuation is \(2\), so the
valuation of \(Q\) on this component is \(1\);
it is zero on every other component.
Therefore \(Q B_2=2B_2\), while \(Q\) is invertible
on the complementary rank-\(79\) lattice.
This proves the stated decomposition.
\end{proof}

\section{Simultaneous global arithmetic normalization}
\label{sec:simultaneous-global-cm}

\begin{theorem}
\label{thm:simultaneous-global-cm-normalization}
Let \(\Pi_{\rm global}\) be the set of primes \(p\geq5\) such that
\(Q_p(1)\) is odd, the natural first-layer map
\eqref{eq:uniform-natural-layer} is a finite isomorphism, and
\(\mathscr K_p\) is injective. Define
\[
 R_{\rm global}=R_{\rm quad}\cap
       \bigcap_{p\in\Pi_{\rm global}}\ker h_p,\qquad
 V_{\rm global}=V_{\rm quad}\cap R_{\rm global}.
\]
Then the joint images of \(R_{\rm quad}\), \(V_{\rm quad}\),
\(R_{\rm quad}\cap H\), and \(V_{\rm quad}\cap H\) are all
\[
 \mathcal P=\prod_{p\in\Pi_{\rm global}}P_p,
\]
including every finite sign subgroup. Consequently,
\begin{equation}
 R_{\rm quad}/R_{\rm global}\simeq\mathcal P,\qquad
 V_{\rm quad}/V_{\rm global}\simeq\mathcal P,\qquad
 G=R_{\rm global}H,
 \label{eq:uniform-global-factorization}
\end{equation}
and inclusion induces an isomorphism
\begin{equation}
 R_{\rm global}/V_{\rm global}
       \xrightarrow{\ \sim\ }R_{\rm quad}/V_{\rm quad}
 \label{eq:uniform-proodd-identification}
\end{equation}
which preserves the image of the arithmetic intersections.
In particular the quotient in \eqref{eq:uniform-proodd-identification}
is pro-odd. The relative quotients satisfy
\begin{equation}
 B(R_{\rm global})\simeq B,\qquad
 B(V_{\rm global})_{\max2}\simeq B_{\max2}.
 \label{eq:uniform-relative-quotients}
\end{equation}
The original restrictions of \(X_2\) to both new groups have
kernel exactly \(\mathcal S\), and image all continuous
\(G\)-invariant quadratic characters which vanish on the
corresponding arithmetic intersection.

The older unit set \(\Pi\) and the primes \(29,113\) belong to
\(\Pi_{\rm global}\), whereas \(163\) does not. No infinitude
assertion about \(\Pi_{\rm global}\) is needed.
\end{theorem}
\begin{proof}
For each \(p\in\Pi_{\rm global}\),
Theorem~\ref{thm:uniform-arithmetic-cm-image} gives the complete
\(P_p\)-axis inside the joint image of \(V_{\rm quad}\cap H\):
the witness is \(1\) in every other prime coordinate.
Finite products of such axes are dense in \(\mathcal P\).
The image of the compact group \(V_{\rm quad}\cap H\) is
closed, hence equals \(\mathcal P\). All four stated images
therefore equal \(\mathcal P\).

Matching every scalar coordinate by an element of
\(V_{\rm quad}\cap H\) gives
\[
 R_{\rm quad}=R_{\rm global}(V_{\rm quad}\cap H).
\]
The known factorization \(G=R_{\rm quad}H\) now gives
\eqref{eq:uniform-global-factorization}. The same equality,
and \(V_{\rm quad}\cap R_{\rm global}=V_{\rm global}\),
prove \eqref{eq:uniform-proodd-identification}.
If \(r\in R_{\rm quad}\cap H\), match its coordinates by
\(v\in V_{\rm quad}\cap H\). Then
\(rv^{-1}\in R_{\rm global}\cap H\), with the same class
in \(R_{\rm quad}/V_{\rm quad}\). This proves the assertion
about the arithmetic image.

For completeness, if a closed normal \(L\subseteq N\)
satisfies \(G=LH\), every continuous \(G\)-invariant
homomorphism \(\varphi:L\to E\), with \(E\) finite abelian
and \(\varphi(L\cap H)=0\), extends uniquely by
\[
 \widetilde\varphi(lh)=\varphi(l).
\]
The intersection condition proves well-definedness,
invariance proves multiplicativity, and the compact
quotient map \(L\times H\to G\) proves continuity.
The same description holds for \(N\), since \(G=NH\).
This identifies all finite abelian relative duals and
therefore gives the first isomorphism in
\eqref{eq:uniform-relative-quotients}.

If \(Z\subset L\) are closed normal and \(L/Z\) is pro-odd,
continuous coprime cohomology gives, for every finite
abelian \(2\)-group \(E\),
\[
 H^1(L,E)\xrightarrow{\ \sim\ }H^1(Z,E)^{L/Z}.
\]
Indeed \(H^1(L/Z,E)=H^2(L/Z,E)=0\).
A \(G\)-invariant character on \(Z\) extends uniquely to
a \(G\)-invariant character on \(L\). If it vanishes on
\(Z\cap H\), its restriction to \(L\cap H\) factors through
the pro-odd group \((L\cap H)/(Z\cap H)\), and is zero.
Applying this with \(L=R_{\rm global}\) and
\(Z=V_{\rm global}\) proves the second isomorphism in
\eqref{eq:uniform-relative-quotients}. It does not require
a factorization \(G=V_{\rm global}H\).
The accepted splitting
\(X_2=\operatorname{Hom}(B,\mathbb F_2)\oplus\mathcal S\),
together with the vanishing of \(\mathcal S\) on
\(R_{\rm quad}\), proves the original restriction assertions.

For \(p\in\Pi\), one has \(T_p=Q_pT_p\). Surjectivity of
the first-layer map makes its cokernel zero, so the
hypothesis holds and \(\mathscr K_p\) is injective.
At \(29\) and \(113\), the exact indices
\([T_p:Q_pT_p]=8\) and the three-dimensional first-layer
cokernels prove the natural isomorphism. The established
full arithmetic images \(P_p\), together with
\eqref{eq:uniform-arithmetic-image}, force
\(\ker\mathscr K_p=0\).
At \(163\), the corresponding index is \(4\), the
first-layer cokernel has dimension \(2\), and the
arithmetic image is exactly \(Q_pP_p\). Thus
\(\mathscr K_{163}=0\) on a two-dimensional domain,
and \(163\notin\Pi_{\rm global}\).
\end{proof}

For any one added prime satisfying this criterion, the newly
imposed kernel has infinite index on each preceding normalized
\(R\)- or \(V\)-group. Indeed \(P_p\) has
\(\mathbb Z_2\)-rank \((p-1)/2\): the rational logarithm
identifies its span with the anti-invariant coefficient space,
and \(s-1\) is invertible there because its characteristic
polynomial is \(X^{(p-1)/2}+1\). This also proves that the
simultaneous normalization is a strict, infinite-index
normalization of \(R_{\rm quad}\) and \(V_{\rm quad}\).
If \(\mathscr K_p\) is not injective, the uniform theorem
determines the arithmetic image but does not determine the
full GT image or the relative quadratic quotient.

\section{Restriction after every prime scalar is killed}
\label{sec:all-cm-original-restriction}

Define
\[
 R_{\rm all}=R_{\rm quad}\cap\bigcap_{p\geq5}\ker h_p,\qquad
 V_{\rm all}=V_{\rm quad}\cap R_{\rm all}.
\]
For any closed normal subgroup \(L\subseteq N\), write
\[
 B(L)=L/\overline{[G,L](L\cap H)}.
\]
These are intrinsic relative quotients of the indicated
groups. Their comparison with \(B\) requires attention
to which characters extend.

\begin{theorem}
\label{thm:all-cm-original-restriction}
The original quadratic restriction maps satisfy
\begin{equation}
 \ker\bigl(X_2\to H^1(R_{\rm all},\mathbb F_2)^G\bigr)
 =\ker\bigl(X_2\to H^1(V_{\rm all},\mathbb F_2)^G\bigr)
 =\mathcal S.
 \label{eq:all-cm-original-kernel}
\end{equation}
There are canonical maps
\begin{equation}
 B(V_{\rm all})_{\max2}
       \xrightarrow{\ \sim\ }B(R_{\rm all})_{\max2}
       \twoheadrightarrow B_{\max2}.
 \label{eq:all-cm-relative-surjection}
\end{equation}
The last map is a surjection; an isomorphism is not
asserted. Moreover,
\[
 R_{\rm all}/V_{\rm all}
       \simeq R_{\rm quad}/V_{\rm quad},
\]
and this isomorphism preserves the image of the
arithmetic intersections.
\end{theorem}
\begin{proof}
Let \(C\) be the image of the joint scalar homomorphism
\[
 h:R_{\rm quad}\longrightarrow\prod_{p\geq5}P_p.
\]
It is a compact abelian pro-\(2\) group.
The isolated local arithmetic axes and compactness give
\begin{equation}
 A_0:=\prod_{p\geq5}Q_pP_p
                   \subseteq h(V_{\rm quad}\cap H)\subseteq C.
 \label{eq:all-cm-arithmetic-axes}
\end{equation}
The image of \(V_{\rm quad}\) equals \(C\): the quotient
of these two images is both a quotient of the pro-odd
group \(R_{\rm quad}/V_{\rm quad}\) and a pro-\(2\)
group, so it is trivial.
The same argument applied to
\((R_{\rm quad}\cap H)/(V_{\rm quad}\cap H)\)
shows that their arithmetic images under \(h\) also agree.

We claim
\begin{equation}
 C=A_0+\overline{[G,C]}.
 \label{eq:all-cm-projection-identity}
\end{equation}
Here additive notation is used in the compact scalar group.
For each \(p\), let \(S_p\) be conjugation by a primitive
arithmetic \(p\)-inertia element. Equivariance makes
\(S_p\) act as \(s_p\) on the \(p\)-coordinate and as
the identity on every other prime coordinate.
The reflex polynomial has augmentation
\(\epsilon_p=Q_p(1)\in\{1,-1\}\).
Write
\[
 Q_p(s)=\epsilon_p+(s-1)F_p(s),\qquad F_p\in\mathbb Z[s].
\]
For \(c=(c_q)_q\in C\), let \(i_p\) denote the
single-coordinate inclusion into the ambient product.
There is an identity in that product:
\begin{equation}
 \epsilon_p i_p(c_p)
   =i_p(Q_pc_p)-(S_p-1)F_p(S_p)c.
 \label{eq:all-cm-coordinate-projection}
\end{equation}
The first term on the right belongs to \(A_0\),
and the second to \([G,C]\).
Thus \(i_p(c_p)\) belongs to the right side of
\eqref{eq:all-cm-projection-identity}, as well as to \(C\).
Finite sums of these coordinate projections converge
to \(c\). The right side of
\eqref{eq:all-cm-projection-identity} is closed, being
the sum of two compact subgroups. This proves the claim.
In particular every continuous \(G\)-invariant
character of \(C\) which vanishes on its arithmetic
image is zero.

Use the established decomposition
\(X_2=\operatorname{Hom}(B,\mathbb F_2)\oplus\mathcal S\).
The summand \(\mathcal S\) vanishes on \(R_{\rm quad}\).
If a relative character from \(\operatorname{Hom}(B,\mathbb F_2)\)
vanishes on \(R_{\rm all}\), its restriction to
\(R_{\rm quad}\) factors through \(C\), is invariant,
and kills the arithmetic image containing \(A_0\).
It is zero by \eqref{eq:all-cm-projection-identity};
the known factorization \(G=R_{\rm quad}H\) makes the
original character zero. The same proof on \(V_{\rm quad}\)
works because its joint image is \(C\) and its original
relative-character restriction is an isomorphism.
This proves \eqref{eq:all-cm-original-kernel}.

Equality of the full joint images gives
\(R_{\rm quad}=R_{\rm all}V_{\rm quad}\), hence
\[
 R_{\rm all}/V_{\rm all}\simeq R_{\rm quad}/V_{\rm quad}.
\]
Equality of the joint arithmetic images allows the same
matching using \(V_{\rm quad}\cap H\), proving the
assertion about arithmetic images in this quotient.
It is pro-odd. Coprime extension of invariant
characters with finite abelian \(2\)-group values,
including their zero arithmetic restriction, therefore
gives the first isomorphism in
\eqref{eq:all-cm-relative-surjection}.

Finally, the canonical image of
\(B(R_{\rm all})_{\max2}\) in \(B_{\max2}\) is closed.
If it were proper, the resulting nonzero abelian
pro-\(2\) quotient would have a nonzero continuous
character to \(\mathbb F_2\). Its pullback would be
an original relative quadratic character vanishing
on \(R_{\rm all}\), contrary to the injectivity just
proved. The last arrow is consequently surjective.
\end{proof}

For clarity, the augmentation used above is integral.
Represent the inverse CM type as
\[
 \Phi_p^*=(1+s+\cdots+s^{h-1})+(c-1)A(s),
 \qquad\deg A<h,
\]
where each coefficient of \(A\) is \(0\) or \(1\).
The identity
\((s-1)\Phi_p^*=(c-1)(1+(s-1)A(s))\)
gives \(Q_p\) after reduction by \(s^h+1\).
Its augmentation is \(1-2a_{h-1}\), hence is \(1\)
or \(-1\). This also explains why no determinant
or irreducibility hypothesis is needed in
\eqref{eq:all-cm-coordinate-projection}.

\subsection{The intrinsic extension problem}

Let \(R_0=R_{\rm quad}\), let \(C=h(R_0)\), and put
\(D_0=h(R_0\cap H)\). Thus \(R_{\rm all}=\ker h\),
\(G=R_0H\), and
\[
 C=D_0+\overline{[G,C]}
\]
by \eqref{eq:all-cm-projection-identity}.
The following criterion separates the central extension class
from the conjugation action on its possible splittings.
Theorem~\ref{thm:all-cm-central-extension} will eliminate
the second condition for this CM quotient.

\begin{proposition}
\label{prop:all-cm-intrinsic-extension}
Let \(\kappa:R_{\rm all}\to\mathbb F_2\) be continuous,
\(G\)-invariant, and zero on \(R_{\rm all}\cap H\).
There is a central relative obstruction
\[
 o_2(\kappa)\in H^2(C,D_0;\mathbb F_2).
\]
If it vanishes, there is a conjugation obstruction
\[
 o_1(\kappa)\in
 H^1_{\rm cont}\!\left(G/R_0,\
    \operatorname{Hom}_{\rm cont}(C/D_0,\mathbb F_2)\right).
\]
The character extends to a continuous character of
\(G\) which is zero on \(H\) if and only if both
obstructions vanish. When it exists, the extension
is unique. Characters on \(V_{\rm all}\) satisfy
the same criterion after their unique coprime
extension to \(R_{\rm all}\).
\end{proposition}
\begin{proof}
For nonzero \(\kappa\), form the central extension
\(R_0/\ker\kappa\to C\) with kernel \(\mathbb F_2\).
Its kernel is central because \(\kappa\) is
\(G\)-invariant. The image of \(R_0\cap H\) gives
a prescribed continuous splitting over \(D_0\);
its kernel is \(R_{\rm all}\cap H\), on which
\(\kappa\) is zero. Define \(o_2(\kappa)\) to be
the relative extension class, and set \(o_2(0)=0\).
This class vanishes exactly when the extension has
a continuous splitting agreeing with the prescribed
one. Projection onto the central factor is then
a character of \(R_0\) extending \(\kappa\) and
zero on \(R_0\cap H\). Conversely the kernel of
such a character gives the required splitting.

These relative classes use continuous cochains
without an extra topological hypothesis.
A finite-kernel profinite extension has a continuous
section: choose an open subgroup meeting the kernel
trivially and choose representatives for its finitely
many image cosets. The section can be made to agree
with a prescribed section over a closed subgroup,
because a continuous finite-valued function on a
closed subset of a profinite space extends by
clopen separation. This realizes the extension
and its splitting in continuous relative cohomology.

When \(o_2(\kappa)=0\), the possible extensions to
\(R_0\) form a torsor under
\(\mathscr U=\operatorname{Hom}_{\rm cont}(C/D_0,\mathbb F_2)\).
Conjugation by \(G\) preserves this torsor: the subgroup
\(H\) preserves \(R_0\cap H\), and inner conjugation
by \(R_0\) acts trivially on characters of \(R_0\).
The action factors through \(G/R_0\).

Give these character spaces the discrete topology.
Every continuous character \(\lambda:R_0\to\mathbb F_2\)
has open stabilizer: choose an open normal
\(W\triangleleft G\) with
\(W\cap R_0\subseteq\ker\lambda\); then
\([W,R_0]\subseteq W\cap R_0\), so \(W\) fixes
\(\lambda\). This proves continuity even when
\(\mathscr U\) is infinite.
Choosing one extension defines the continuous
affine cocycle of the torsor, whose cohomology
class is \(o_1(\kappa)\). It vanishes exactly when
there is a \(G\)-invariant extension to \(R_0\).

Finally a \(G\)-invariant character of \(R_0\)
zero on \(R_0\cap H\) extends along \(G=R_0H\)
by \(\chi(rh)=\chi(r)\), continuously by compact
descent. The difference of two possible invariant
extensions belongs to \(\mathscr U^{G/R_0}\).
This fixed space is zero because
\(C=D_0+\overline{[G,C]}\).
Thus the extension is unique.

The pro-odd quotient
\(R_{\rm all}/V_{\rm all}\) gives the asserted
unique extension from \(V_{\rm all}\), including
the zero arithmetic restriction, by continuous
coprime cohomology.
\end{proof}

For an isomorphism of the relative pro-\(2\) groups
in \eqref{eq:all-cm-relative-surjection}, the extension
property must hold for invariant relative characters
with every finite abelian \(2\)-group as target.
The proposition gives the exact first case, with
target \(\mathbb F_2\). Theorem~\ref{thm:all-cm-central-extension}
proves that its conjugation obstruction always vanishes.
The relative central class remains to be determined;
the factorization \(G=R_{\rm all}H\) is not asserted.

\subsection{The two possible full images at degree \(163\)}
Write \(\mathcal R=R_{\rm global}\) and
\(\mathcal V=V_{\rm global}\), let \(P=P_{163}\), and put
\(A=Q_{163}P\).  Thus \(P/A\) is the two-dimensional irreducible
\(\mathbf F_2[C_3]\)-module.  The cyclic group \(C_3\) is the quotient
of \(\operatorname{Gal}(\mathbf Q(\zeta_{163})/\mathbf Q)\) through
which this module is acted on.

\begin{proposition}\label{prop:cm163-full-image-dichotomy}
The full images \(h_{163}(\mathcal R)\) and
\(h_{163}(\mathcal V)\) coincide.  There are precisely two
possibilities.
\begin{enumerate}
\item Their common image is \(A\).  The groups
\(\mathcal R'=\mathcal R\cap\ker h_{163}\) and
\(\mathcal V'=\mathcal V\cap\ker h_{163}\) then give a further
arithmetic normalization \(G=\mathcal R'H\).  The full relative
abelian quotient on \(\mathcal R'\), its maximal pro-\(2\) quotient
on \(\mathcal V'\), and the exact original quadratic restriction
kernel and image are unchanged.
\item Their common image is \(P\).  There is a continuous
surjection \(G\longrightarrow A_4\) whose restriction to \(H\)
has image a subgroup of order three.
\end{enumerate}
In the second case the displayed finite quotient has zero relative
quadratic character space.
\end{proposition}
\begin{proof}
Since \(\mathcal R/\mathcal V\) is pro-odd, the quotient
\(h_{163}(\mathcal R)/h_{163}(\mathcal V)\) is both pro-odd and
pro-\(2\), and is trivial.  The exact arithmetic image is \(A\).
The common full image modulo \(A\) is an invariant subspace of
the irreducible module \(C=P/A\), and hence is zero or \(C\).

If the common image is \(A\), its equality with the arithmetic
image on both arithmetic intersections proves
\(\mathcal R=\mathcal R'(\mathcal V\cap H)\), and
therefore \(G=\mathcal R'H\).  The relative gluing and pro-odd
comparison arguments apply as before.  The exact original
restriction statements follow from the decomposition into the
relative character space and the fixed lift space \(\mathcal S\).

Suppose instead that the common full image is \(P\), and let
\(t:\mathcal R\to C\) be \(h_{163}\) followed by reduction modulo
\(A\).  It is surjective and vanishes on \(\mathcal R\cap H\).
Write \(\gamma:G\to C_3\) for the cyclotomic quotient acting on
\(C\).  For \(g=rh\), \(r\in\mathcal R\), \(h\in H\), set
\(z(g)=t(r)\).  This is well defined because \(t\) vanishes on
the overlap.  It is continuous: the continuous function
\((r,h)\mapsto t(r)\) on the compact space \(\mathcal R\times H\)
descends through its quotient map onto \(G\).  Equivariance and
the multiplication rule for \(rh\) give
\[
 z(g_1g_2)=z(g_1)+\gamma(g_1)z(g_2).
\]
Here \(\gamma(\mathcal R)=1\), as \(\mathcal R\subset\ker\lambda\).
It follows that
\[
 g\longmapsto (z(g),\gamma(g))
\]
is a continuous homomorphism to \(C\rtimes C_3\).  Its image
contains all translations \(C\) and the entire complement \(C_3\);
the latter follows from \(G=\mathcal RH\) and
\(\gamma(\mathcal R)=1\).  It is therefore surjective.
The faithful irreducible action identifies
\(C\rtimes C_3\) with \(A_4\), for example through its affine
action on the four elements of \(C\).  The image of \(H\) is
exactly the complement.

Finally \((C_3-1)C=C\), since the irreducible nontrivial
\(C_3\)-module \(C\) has no trivial constituent.  Every invariant
homomorphism \(C\to\mathbf F_2\) is consequently zero.  This
proves the last assertion without deciding the original relative
quadratic character space.
\end{proof}

\section{The integral degree-\(163\) extension obstruction}
\label{sec:163-relative-extension}

Let \(P=B_2\oplus L_{79}\) and \(A=2B_2\oplus L_{79}\)
be the pair from Lemma~\ref{lem:163-integral-pair-decomposition}.
Put \(E=P/A\) and \(U=E^\vee\).
Relative cohomology uses continuous cochains. In degree two
it classifies central extensions by \(\mathbb F_2\) together
with a specified continuous splitting over \(A\).

\begin{proposition}
\label{prop:163-relative-cohomology}
There are natural identifications, relative to the displayed
stable decomposition,
\[
 H^1(P,A;\mathbb F_2)=U,
\]
\begin{equation}
 H^2(P,A;\mathbb F_2)
  \simeq\operatorname{Quad}(E,\mathbb F_2)
       \oplus
       \operatorname{Bil}(E\times(L_{79}/2L_{79}),\mathbb F_2).
 \label{eq:163-relative-forms}
\end{equation}
The latter space has dimension \(161\).
Inversion acts trivially on it, and
\begin{equation}
 H^2(P,A;\mathbb F_2)^{C_{162}}
       =\mathbb F_2\,\mathfrak q.
 \label{eq:163-quaternion-line}
\end{equation}
Here \(\mathfrak q\) is the pullback of
\(Q_8\to E\simeq\mathbb F_4\) along \(P\to E\), with its
canonical splitting over \(A\). Its quadratic form is
\(q(z)=z^3\). This class is fixed by every automorphism
of the pair \((P,A)\).
\end{proposition}
\begin{proof}
The relative degree-one group is the kernel of restriction
\(\operatorname{Hom}(P,\mathbb F_2)\to
\operatorname{Hom}(A,\mathbb F_2)\), hence is \(E^\vee\).

Consider a central extension \(\widetilde P\) with splitting
\(\sigma:A\to\widetilde P\), and identify its central kernel
with \(\mathbb F_2\). For \(b\in B_2\), choose a lift
\(\widetilde b\), and set
\[
 q(\overline b)=\widetilde b^{\,2}\sigma(2b)^{-1}.
\]
Changing the central lift changes no square.
Changing \(b\) by \(2b'\) also changes nothing: commutators
are bilinear with values of exponent two, and hence the
commutator against \(2b'\) is zero. Thus \(q\) is
well-defined on \(E\). The square of a product shows that
its polarization is the commutator pairing on \(B_2\).
For \(l\in L_{79}\), define the mixed pairing
\[
 m(\overline b,\overline l)
             =[\widetilde b,\sigma(l)].
\]
It is bilinear and factors through \(E\) and \(L_{79}/2L_{79}\).
The remaining commutators vanish because \(\sigma\) is
a homomorphism on \(A\).

Conversely, for any quadratic form \(q\) on \(E\), choose
a bilinear cocycle \(c_q:E\times E\to\mathbb F_2\) with
diagonal \(c_q(e,e)=q(e)\). Such a cocycle exists by
writing the quadratic polynomial in a basis of \(E\).
Given also \(m\), define a central extension on
\(P\times\mathbb F_2\) by the cocycle
\[
 ((b,l),(b',l'))\longmapsto
       c_q(\overline b,\overline b')
              +m(\overline b,\overline l').
\]
It is continuous and zero on \(A\times A\), so the zero
central coordinate gives the required splitting.
Its forms are exactly \(q,m\).

If \(q=m=0\), all commutators vanish. Choose lifts of
a \(\mathbb Z_2\)-basis of \(B_2\); their squares equal
the prescribed lifts of twice those basis vectors.
Their continuous \(\mathbb Z_2\)-powers, together with
\(\sigma\) on \(L_{79}\), then define a splitting on
all of \(P\) extending \(\sigma\). This proves both
injectivity and surjectivity of
\eqref{eq:163-relative-forms}. The dimension is
\(3+2\cdot79=161\). Inversion leaves both the quadratic
and mixed forms unchanged.

The action may therefore be computed using \(t=-s\),
of odd order \(81\). Its regular representation on
\(P/2P\) contains exactly one copy of the irreducible
two-dimensional module \(E\); the complement \(L_{79}/2L_{79}\)
contains no such copy. This module is self-dual.
Hence the invariant mixed-pairing space is zero.
Quadratic forms on \(E\) are the functions on its three
nonzero vectors, and the cyclotomic \(C_3\) permutes
those vectors transitively. The invariant quadratic
forms thus have dimension one, generated by the form
which is \(1\) on every nonzero vector.
In the field model for \(E\) it is \(q(z)=z^3\).
This proves \eqref{eq:163-quaternion-line}.

Writing \(z=a+b\omega\), where
\(\omega^2+\omega+1=0\), gives
\[
 q(z)=a+b+ab,\qquad
 q(z+w)+q(z)+q(w)=zw^2+z^2w.
\]
The cocycle
\(aa'+bb'+ba'\) realizes the associated central extension
as \(Q_8\): every nonzero quotient vector has lifts of
order four, and the commutator pairing is nonzero.
Pullback to \(P\) splits canonically over \(A\).
The form \(q\) is invariant under every permutation
of the three nonzero vectors, hence under
\(\operatorname{GL}_2(\mathbb F_2)\).
Its relative pullback class is consequently fixed by
every automorphism of \((P,A)\).
\end{proof}

One can also identify the entire cyclotomic module.
Let \(\mathcal R=\mathbb F_2[C_{81}]\) and let
\(\mathcal I\) be its augmentation summand.
Semisimplicity gives
\[
 H^2(P,A;\mathbb F_2)\simeq\mathcal R\oplus\mathcal I
\]
noncanonically, with inversion acting trivially.
Indeed \(\operatorname{Quad}(E)=\mathbb F_2\oplus E\),
\(E\otimes E\simeq\mathbb F_2^2\oplus E\), and
\(E\otimes\mathcal R\simeq\mathcal R^2\); substitution
in \eqref{eq:163-relative-forms} gives
\(2\mathcal R-\mathbb F_2=\mathcal R\oplus\mathcal I\)
in the semisimple module category.

\subsection{The exact two-stage extension criterion}

Suppose \(R=R_{\rm global}\) has full degree-\(163\) image
\(P\), and put \(J=\ker(h_{163}:R\to P)\).
We know \(G=RH\) and \(h_{163}(R\cap H)=A\).
Consider the intrinsic relative character space
\[
 \mathscr X(J)=
 \{\kappa:J\to\mathbb F_2:
   \kappa\text{ is continuous and }G\text{-invariant},
   \ \kappa(J\cap H)=0\}.
\]

\begin{proposition}
\label{prop:163-two-stage-extension}
For \(\kappa\in\mathscr X(J)\) there is a linear central
obstruction
\[
 o_2(\kappa)\in H^2(P,A;\mathbb F_2)^{C_{162}}
                       =\mathbb F_2\,\mathfrak q.
\]
If it vanishes, there is a second obstruction
\[
 o_1(\kappa)\in H^1_{\rm cont}(\Gamma,U),
 \qquad\Gamma=G/R,\quad U=(P/A)^\vee.
\]
A continuous character of \(G\), zero on \(H\), extends
\(\kappa\) if and only if both obstructions vanish.
When it exists, this extension is unique.
Thus restriction from the original relative quadratic
character space to \(\mathscr X(J)\) is injective,
with image precisely the characters passing these two tests.
\end{proposition}
\begin{proof}
For nonzero \(\kappa\), the quotient \(R/\ker\kappa\)
is a central extension of \(P\) by \(\mathbb F_2\).
Centrality follows from \(G\)-invariance of \(\kappa\).
The image of \(R\cap H\) maps isomorphically to \(A\):
its kernel is the image of \(J\cap H\), which
\(\kappa\) kills. It therefore gives a specified
continuous splitting over \(A\).
Define \(o_2(\kappa)\) to be this relative extension
class, and define \(o_2(0)=0\).
Pushout and the Baer-sum construction show linearity.
Conjugation by \(H\) preserves the extension and
arithmetic splitting, so its class is cyclotomically
invariant. Proposition~\ref{prop:163-relative-cohomology}
therefore places it on the displayed line.

Vanishing of this class means that the central
extension splits by a section extending the prescribed
section on \(A\). Projection onto the central factor
then gives a character of \(R\) extending \(\kappa\)
and killing \(R\cap H\). Conversely, the kernel of
any such character supplies a section with the required
restriction. Thus \(o_2(\kappa)=0\) is exactly the
first extension condition.

The possible extensions to \(R\), zero on \(R\cap H\),
form a torsor under \(U\): their difference factors
through \(P\), and vanishes on \(A\).
Conjugation by \(G\) acts on this torsor.
To check preservation of the zero condition, write
\(g=rh\). Inner conjugation by \(r\in R\) acts trivially
on characters of \(R\), and conjugation by \(h\in H\)
preserves \(R\cap H\). The action therefore factors
through \(\Gamma=G/R\).
Choose an extension \(\widetilde\kappa\). The differences
\[
 \gamma\longmapsto
       \gamma\widetilde\kappa-\widetilde\kappa
\]
form a continuous \(U\)-valued cocycle; changing the
chosen extension changes it by a coboundary.
Its class \(o_1(\kappa)\) vanishes precisely when
the torsor has a fixed point.

A fixed point is a \(G\)-invariant character of \(R\)
zero on \(R\cap H\). The formula
\(\chi(rh)=\widetilde\kappa(r)\) then extends it to
\(G\), continuously and with zero restriction to \(H\).
Conversely a character of \(G\) with this restriction
gives a fixed point. The difference of two fixed points
belongs to \(U^\Gamma\), which is zero because the
cyclotomic \(C_3\)-action is irreducible.
This proves the criterion and uniqueness.
If an original relative character is zero on \(J\),
its restriction to \(R\) factors through the same
fixed space \(U^\Gamma\), so it is zero on \(R\),
and hence on \(G=RH\). This proves injectivity.
\end{proof}

The cyclic calculation alone does not settle the second obstruction;
Theorem~\ref{thm:163-concrete-odd-inertia} will settle it using the
arithmetic inertia and exact defining intersection of \(R\). To see why
those further inputs are needed, if
\(\Gamma_0=\ker(\Gamma\to C_3)\), then continuous
inflation--restriction and odd-order averaging give
\begin{equation}
 H^1(\Gamma,U)
   \simeq\operatorname{Hom}_{\rm cont}(\Gamma_0,U)^{C_3}.
 \label{eq:163-full-conjugation-obstruction}
\end{equation}
The action on \(U\) is trivial on \(\Gamma_0\);
both \(H^1(C_3,U)\) and \(H^2(C_3,U)\) vanish.
Although all positive-degree cohomology
\(H^j(C_{162},U)\) vanishes, the right side of
\eqref{eq:163-full-conjugation-obstruction} need not vanish for an
arbitrary profinite group with the same coefficient action.
For example, for the abstract group
\(\Gamma=E\rtimes C_3\) it has dimension two over
\(\mathbb F_2\), since
\(\operatorname{Hom}_{C_3}(E,U)\simeq\mathbb F_4\).
These continuous-cohomology statements use the
finite-quotient formalism and spectral sequence in
\cite[II, Remark~1.35 and Proposition~4.2]{MilneCFT}.

\begin{proposition}
\label{prop:163-central-obstruction-model}
The quaternion class can obstruct extension in a
profinite group with the stated integral pair and
cyclotomic action.
\end{proposition}
\begin{proof}
Form
\[
 \widetilde P=P\times_E Q_8,
\]
with central kernel \(\mathbb F_2\) and canonical
splitting \(\sigma:A\to\widetilde P\).
The order-three automorphism of \(Q_8\) cycling its
three quaternion directions lifts the action on \(E\).
Together with the given action on \(P\), it defines
a \(C_{162}\)-action on \(\widetilde P\); its
order-two factor acts on \(P\) by inversion and
trivially on \(Q_8\).
Set
\[
 G_0=\widetilde P\rtimes C_{162},\qquad
 H_0=\sigma(A)\rtimes C_{162}.
\]
The identity character of the central kernel is
\(G_0\)-invariant and zero on its intersection with
\(H_0\). It cannot extend to \(\widetilde P\),
because lifts of an \(\mathbb F_2\)-basis of \(E\)
have nontrivial central commutator. Every abelian
character kills that commutator.
Its relative extension class is exactly
\(\mathfrak q\). This is a group-theoretic model;
no identification with a quotient or subgroup
of \(G\) is asserted.
\end{proof}

\subsection{An odd-inertia vanishing criterion}
\label{sec:odd-inertia-criterion}

\begin{proposition}\label{prop:odd-inertia-vanishing}
Let \(R\subseteq L\) be closed normal subgroups of a profinite group
\(G\), and put \(\Gamma=G/R\).
Assume \(C=L/R\) is a central pro-\(2\) subgroup of \(\Gamma\).
Let \(a\in G\) have pro-odd procyclic closure, and assume its closed
normal closure in \(G/L\) is pro-odd.
If \(U\) is a finite continuous \(2\)-primary \(\Gamma\)-module with
\(U^a=0\), then
\[
 H^j_{\rm cont}(\Gamma,U)=0\qquad\text{for every }j\ge0.
\]
\end{proposition}

\begin{proof}
We first show that the closed normal closure \(N\) of the image of
\(a\) in \(\Gamma\) is itself pro-odd.

We use an elementary finite-group fact.
If a finite group \(E_0\) is a central extension of an odd-order group
\(D_0\) by a finite \(2\)-group \(C_0\), then it has a unique odd-order
complement \(J_0\) and \(E_0=C_0\times J_0\).
Indeed \(C_0\) is abelian.  For a central-extension cocycle \(f\),
put
\[
 b(x)=|D_0|^{-1}\sum_{z\in D_0}f(x,z)\in C_0.
\]
Multiplication by \(|D_0|\) is invertible on \(C_0\).
Summing the cocycle identity over \(z\) gives
\(f(x,y)=b(x)+b(y)-b(xy)\), so the extension splits.
Two homomorphic splittings differ by a homomorphism \(D_0\to C_0\),
which is zero by coprimality.
The complement is therefore unique and characteristic, and every
odd-order element of \(E_0\) lies in it.

Now take any finite quotient \(F\) of \(\Gamma\), and let \(C_F\)
be the image of \(C\).  In \(F/C_F\), the image of the assumed
pro-odd normal closure in \(G/L\) is an odd-order normal subgroup.
Its preimage \(E_F\) in \(F\) is a central extension by the \(2\)-group
\(C_F\).  The finite-group fact gives its unique odd-order complement
\(J_F\).  Since \(E_F\) is normal in \(F\), uniqueness makes \(J_F\)
normal in \(F\).
The image of \(a\) has odd order and belongs to \(E_F\), hence to
\(J_F\).  All its \(F\)-conjugates belong to \(J_F\).
Thus the image of \(N\) in every finite quotient of \(\Gamma\) is
odd-order.  Ambient finite quotients are cofinal for the topology
on the closed subgroup \(N\), so \(N\) is pro-odd.

For a pro-odd group and a finite \(2\)-primary continuous module,
averaging in finite quotients gives
\(H^j(N,U)=0\) for every \(j>0\).
Also \(U^N\subseteq U^a=0\), since \(a\) belongs to \(N\).
All rows of the continuous Hochschild--Serre spectral sequence
\[
 H^i(\Gamma/N,H^j(N,U))
       \ \Longrightarrow\ H^{i+j}(\Gamma,U)
\]
are therefore zero.  This proves the assertion, including degree zero.
The spectral sequence and its profinite interpretation are recalled in
\cite[II, Remark~1.35(b), p.~72; Proposition~4.2, p.~87]{MilneCFT}.
\end{proof}

The following application verifies both the central pro-\(2\) layer
and the full normal closure for the exact defining intersection of
\(R_{\rm global}\).

\section{Odd inertia and the full CM extension obstruction}
\label{sec:general-odd-inertia}

The scalar coordinates admit exact product descriptions even when
their individual full images are unknown. An inertia argument then
eliminates the conjugation obstruction for the entire relative
character module. The relative central extension class remains
the condition for extending an intrinsic character.

Put \(R_0=R_{\rm quad}\), \(M=\ker\lambda\), and \(D=N_{\rm rad}\).
The definitions in
Theorems~\ref{thm:simultaneous-quadratic-normalization}
and~\ref{thm:complete-quadratic-normalization} give
\begin{equation}
 \begin{split}
 R_0={}&D\cap\ker\Phi_2\cap\ker\mathfrak t_2\\
 &\cap\bigcap_{\ell<q,\ \ell,q\ {\rm odd}}\ker\xi_{\ell,q}
 \cap\bigcap_{p\ {\rm odd},\ p\ne3}\ker\nu_p,
 \qquad G=R_0H.
 \end{split}
 \label{eq:general-cm-residual-intersection}
\end{equation}
Here \(\Phi_2\) is the full based free-pro-two action and
\(\mathfrak t_2\) is the entire corrected cubic coefficient-two
homomorphism.

For \(q\ge5\), write \(\rho^{\rm bas}_{q,2}\) for the full
compact-homology representation of the degree-\(q\) cyclic cover,
and \(\Delta_q\) for its deck generator.
Equation~\eqref{eq:coefficient-two-cross-3} gives
\[
 \rho^{\rm bas}_{q,2}(g)\Delta_q
       \rho^{\rm bas}_{q,2}(g)^{-1}
       =\Delta_q^{\lambda(g)\bmod q}.
\]
The integral centralizer of \(\Delta_q\) is the commutative
algebra \(\mathcal O_q=\mathbb Z[\zeta_q]\otimes\mathbb Z_2\),
since every nontrivial deck eigencharacter occurs once.
On \(M\), the representation is scalar in this algebra.
Its principal pro-two projection is \(h_q\).
Projection to principal units commutes with the cyclotomic
automorphisms, including permutations of the local factors.
Consequently
\begin{equation}
 h_q(gmg^{-1})=[\lambda(g)\bmod q]\bigl(h_q(m)\bigr)
                              \qquad(g\in G,\ m\in M).
 \label{eq:general-cm-scalar-equivariance}
\end{equation}
In particular every \(\ker h_q\) is closed and normal in \(G\).

\subsection{A prime outside the imposed scalar kernels}

For a set \(\Lambda\) of primes at least five define
\[
 R_\Lambda=R_0\cap\bigcap_{q\in\Lambda}\ker h_q.
\]
The next theorem assumes arithmetic factorization for this
particular intersection.

\begin{theorem}\label{thm:omitted-prime-vanishing}
Suppose \(G=R_\Lambda H\), and let \(p\ge5\) be a prime
outside \(\Lambda\).
Let \(W\) be a finite abelian \(2\)-primary
\(G/R_\Lambda\)-module whose action factors through
\[
 G/R_\Lambda\xrightarrow{\lambda\bmod p}
                  (\mathbb Z/p)^\times=\langle s_p\rangle.
\]
If \(s_p-1\) is bijective on \(W\), then
\begin{equation}
                  H^j_{\rm cont}(G/R_\Lambda,W)=0
                                     \qquad(j\ge0).
 \label{eq:omitted-prime-cohomology}
\end{equation}
\end{theorem}

\begin{proof}
We apply Proposition~\ref{prop:odd-inertia-vanishing}.
Define
\begin{equation}
 L_\Lambda=D\cap\ker\Phi_2\cap\ker\mathfrak t_2
                    \cap\bigcap_{q\in\Lambda}\ker h_q.
 \label{eq:general-cm-intermediate-subgroup}
\end{equation}
Every kernel here is closed and \(G\)-normal.
For \(D\), use the full affine Kummer homomorphism of
Theorem~\ref{thm:rational-radical-normalization}.
The map \(\Phi_2\) is a representation on \(G\).
For the corrected cubic frame, the integral conjugation
identity \eqref{eq:adelic-cubic-kernel-conjugation}
preserves its identity fibre in \(M\).
The remaining kernels are normal by
\eqref{eq:general-cm-scalar-equivariance}.

The two quadratic families in
\eqref{eq:general-cm-residual-intersection} have joint
kernel \(R_\Lambda\) on \(L_\Lambda\).
Their invariance follows from
Theorems~\ref{thm:all-odd-cross-cup-lifts}
and~\ref{thm:coefficient-two-cross-cups}.
Thus \(L_\Lambda/R_\Lambda\) is a closed subgroup of
a product of copies of \(\mathbb F_2\), and
\([G,L_\Lambda]\subseteq R_\Lambda\).
In particular it is central pro-two in \(G/R_\Lambda\).
The intersection \eqref{eq:general-cm-intermediate-subgroup}
includes the kernel of the entire cubic coefficient-two map.

\paragraph{The odd part of arithmetic inertia.}
Write \(p-1=2^e m\), with \(m\) odd.
Let \(Z_p\subset\widehat{\mathbb Z}^{\times}\) be the
finite order-\(m\) subgroup in the \(p\)-component,
with all other components equal to \(1\).
Cyclotomic inertia at \(p\) maps onto
\(\mathbb Z_p^\times\) and is trivial at every other
coefficient prime: the \(p\)-power cyclotomic tower
is totally ramified locally, whereas roots of unity
of order prime to \(p\) are unramified.
Choose \(b\in I_p\subset H\) whose multiplier generates
\(Z_p\).
Let \(e_{\rm odd}\in\widehat{\mathbb Z}\) be zero at
two and one at every odd prime, and put
\(a=b^{e_{\rm odd}}\), with powers taken in the
procyclic closure of \(b\).
Then
\[
 \overline{\langle a\rangle}\text{ is pro-odd},\qquad
 \lambda(a)\text{ generates }Z_p,\qquad
 \lambda_\ell(a)=1\quad(\ell\ne p).
\]

We have \(W^a=0\).
Indeed \(a\) acts through a generator of the odd-order
subgroup of \(\langle s_p\rangle\).
If \(V=W^{\langle s_p^{2^e}\rangle}\) were nonzero,
then \(V[2]\) would be a nonzero finite
\(\mathbb F_2\)-vector space.
The action of \(s_p\) on it has order dividing \(2^e\),
so \((s_p-1)^{2^e}=0\).
This contradicts the injectivity of \(s_p-1\).
Thus \(V=0\), including the case \(m=1\), when the
hypothesis forces \(W=0\).

Let \(\mathscr N_H\) be the closed normal closure of
\(a\) in \(H\).
We claim
\begin{equation}
                       \mathscr N_H\cap D\subseteq L_\Lambda.
 \label{eq:general-cm-arithmetic-kernel-comparison}
\end{equation}

\paragraph{The based action and the corrected cubic frame.}
The full based image \(\Phi_2(G)\) is pro-two.
The peripheral formulas
\(x\mapsto x^\lambda\), \(y\mapsto f^{-1}y^\lambda f\)
act trivially on the mod-two Frattini quotient of
the free pro-two group, because \(\lambda\) is odd.
The kernel of the corresponding automorphism action
is pro-two, as seen on finite characteristic
two-group quotients; this is the argument in the
proof of Theorem~\ref{thm:full-dyadic-free-product}.
It follows that \(\Phi_2(a)=1\), and hence
\(\Phi_2(\mathscr N_H)=1\).
This is an equality of full based automorphisms.

The curve \(E_3:Y^2=X^3+1\) has discriminant \(-432\),
so it has good reduction at \(p\ge5\).
Its two-adic Tate representation is unramified there.
The Galois field
\(\mathbb Q(\mu_4,3^{1/4})\) is also unramified at \(p\):
the discriminant of \(X^4-3\) has prime divisors only
\(2,3\), and adjoining \(\mu_4\) adds ramification
only at two.
Both arithmetic kernels contain \(\mathscr N_H\).

For \(\sigma\in\mathscr N_H\cap M\), the signed comparison
of Proposition~\ref{prop:adelic-signed-path-comparison}
therefore gives
\[
 \mathfrak t_2(\sigma)
   =\frac{\sigma(3^{1/4})}{3^{1/4}}\rho_{E_3,2}(\sigma)=1.
\]
This comparison is applied only on its domain \(H\cap M\),
with its quartic Kummer correction included.

\paragraph{Good reduction and the odd deck factors.}
Fix \(q\in\Lambda\), so \(q\ne p\).
The curve \(w^q=t(t-1)\) has good reduction at \(p\).
For a direct check, write \(q=2g+1\), set \(y=2t-1\),
and use the charts
\[
 y^2=1+4w^q,\qquad
 v^2=u^{q+1}+4u,\qquad
 u=w^{-1},\quad v=y/w^{g+1}.
\]
The first branch polynomial is separable modulo \(p\),
since \(p\nmid2q\).
At the point \(u=v=0\) of the second chart the derivative
with respect to \(u\) is \(-4\), a unit.
The two monic quadratic charts give a finite flat cover
of \(\mathbb P^1_{\mathbb Z_p}\) which is proper and smooth.
Its Jacobian is an abelian scheme.
Its \(2^n\)-torsion is finite etale for every \(n\), hence
constant over the strict henselization.
This proves unramifiedness of the entire two-adic
Tate representation, both here and for \(E_3\).

The arithmetic Tate representation and the based
compact-homology representation differ only by a
deck element, as in the proof of
Proposition~\ref{prop:pdl-arithmetic-realization}.
Indeed two lifts of the same outer arithmetic action
differ by an inner automorphism of the base geometric
group. The covering-subgroup part acts trivially on
compact homology, leaving its deck class.
Since inertia at \(p\) fixes \(\mathbb Q(\zeta_q)\),
good reduction gives
\[
                 \rho^{\rm bas}_{q,2}(a)\in
                              \langle\Delta_q\rangle\simeq C_q.
\]
This deck element need not be trivial.
The full normalizer identity shows that
\(\langle\Delta_q\rangle\) is normalized by
\(\rho^{\rm bas}_{q,2}(H)\).
Consequently
\[
                 \rho^{\rm bas}_{q,2}(\mathscr N_H)
                                  \subseteq\langle\Delta_q\rangle.
\]
For \(\sigma\in\mathscr N_H\cap M\), its principal
pro-two scalar \(h_q(\sigma)\) is therefore one,
because that projection kills the odd-order deck group.
This holds for every imposed coordinate, including
an infinite collection.
Together with \(D\subseteq M\), it proves
\eqref{eq:general-cm-arithmetic-kernel-comparison}.

\paragraph{The full affine radical quotient.}
By Theorem~\ref{thm:rational-radical-normalization},
\[
 G/D\simeq A_{\rm rad}
  =\left\{(u,(t_r)_r):
      u\in\widehat{\mathbb Z}^{\times},\
      t_r\in\widehat{\mathbb Z},\
      t_r\bmod2=\epsilon_r(u)\right\},
 \qquad (u,t)(v,z)=(uv,t+uz),
\]
where \(r\) ranges over all rational primes and
\(\epsilon_r\) is the character of \(\mathbb Q(\sqrt r)\).
The subgroup
\[
 \mathscr O_p
   =\{(u,t)\in A_{\rm rad}:u\in Z_p,\
                    t_{r,2}=0\text{ for every }r\}
\]
is closed and pro-odd.
Every \(\epsilon_r\) vanishes on \(Z_p\), so it is
the semidirect product
\[
 \left(\prod_{r\ {\rm prime}}\prod_{\ell\ {\rm odd}}\mathbb Z_\ell\right)
                                      \rtimes Z_p.
\]
It is normal in the entire affine group.
Indeed conjugation by \((v,z)\) sends \((u,t)\) to
\[
                            (u,vt+(1-u)z);
\]
every dyadic translation stays zero because \(u_2=1\).

The image of \(a\) belongs to \(\mathscr O_p\).
Its multiplier belongs to \(Z_p\), and every
dyadic Kummer cocycle restricts on
\(\overline{\langle a\rangle}\) to a homomorphism
to the additive group \(\mathbb Z_2\).
That homomorphism is zero by coprimality.
The full image of \(\mathscr N_H\) is therefore
contained in \(\mathscr O_p\), and
\[
                        \mathscr N_H/(\mathscr N_H\cap D)
                                         \text{ is pro-odd}.
\]
This argument includes all rational-prime translations
and all odd coefficient components.

Finally the assumed factorization \(G=R_\Lambda H\)
implies \(G=L_\Lambda H\), and compactness gives
\[
                      G/L_\Lambda\simeq H/(H\cap L_\Lambda).
\]
The closed normal closure of \(a\) in this quotient
is exactly the image of \(\mathscr N_H\).
Equation~\eqref{eq:general-cm-arithmetic-kernel-comparison}
and \(L_\Lambda\subseteq D\) give
\[
                   \mathscr N_H\cap L_\Lambda=\mathscr N_H\cap D.
\]
Its normal closure in \(G/L_\Lambda\) is therefore
pro-odd. We have verified the central pro-two layer,
the pro-odd procyclic closure of \(a\), its pro-odd
normal closure in \(G/L_\Lambda\), and \(W^a=0\).
Proposition~\ref{prop:odd-inertia-vanishing} proves
\eqref{eq:omitted-prime-cohomology}.
\end{proof}

\subsection{Exact products of the scalar images}

Let \(\mathcal I\) be any set of primes at least five.
Define the full and arithmetic images by
\begin{equation}
 \begin{aligned}
 C_p&=h_p(R_0),&
 D_p&=h_p(R_0\cap H),\\
 C&=(h_p)_{p\in\mathcal I}(R_0),&
 D_0&=(h_p)_{p\in\mathcal I}(R_0\cap H).
 \end{aligned}
 \label{eq:all-cm-scalar-images}
\end{equation}
All these groups are compact abelian pro-two groups.
The group \(C\) is \(G\)-stable.
The group \(D_0\) is \(H\)-stable and is also \(G\)-stable:
conjugation by \(R_0\) is trivial on \(C\), and \(G=R_0H\).
Their conjugation actions therefore factor through
\(\Gamma_0=G/R_0\).

Use additive operator notation on the scalar unit groups.
For each \(p\), retain
\[
 \mathcal O_p=\mathbb Z[\zeta_p]\otimes\mathbb Z_2,\qquad
 \mathcal U_p=1+2\mathcal O_p,\qquad
 T_p=\{t\in\mathcal U_p:t\,c(t)=1\},\qquad P_p=(s_p-1)T_p.
\]
Here \(s_p\) is a primitive cyclotomic generator and
\(c=s_p^{(p-1)/2}\).
Proposition~\ref{prop:explicit-cm-type-factor} and
Theorem~\ref{thm:cm-type-torus-unit} give the exact identities
\begin{equation}
 (s_p-1)\Phi_p^*=(c-1)Q_p,\qquad
 (c-1)\mathcal U_p=T_p,\qquad
 Q_p(1)=\varepsilon_p\in\{1,-1\}.
 \label{eq:all-cm-integral-operator-inputs}
\end{equation}
The upper bound \eqref{eq:pdl-upper-input} gives
\(C_p\subseteq P_p\).

\begin{lemma}\label{lem:all-cm-isolated-axes}
If \(i_p:P_p\to\prod_{q\in\mathcal I}P_q\)
denotes inclusion in the \(p\)-coordinate, then
\[
                           i_p(Q_pP_p)\subseteq D_0\subseteq C.
\]
In particular \(Q_pP_p\subseteq D_p\) and \(Q_pC_p\subseteq D_p\).
\end{lemma}

\begin{proof}
The witnesses of Proposition~\ref{prop:pdl-arithmetic-realization}
lie in \(V_{\rm quad}\cap H\subseteq R_0\cap H\),
have every other prime coordinate equal to one, and realize
\((s_p-1)^2\Phi_p^*(\mathcal U_p)\) in degree \(p\).
By \eqref{eq:all-cm-integral-operator-inputs}, this image is
\[
 (s_p-1)(c-1)Q_p(\mathcal U_p)
     =Q_p(s_p-1)T_p=Q_pP_p.
\]
These are equalities of the full compact groups,
including their finite sign subgroups.
\end{proof}

\begin{proposition}\label{prop:all-cm-exact-products}
For every set \(\mathcal I\) of primes at least five,
\begin{equation}
                         C=\prod_{p\in\mathcal I}C_p,\qquad
                         D_0=\prod_{p\in\mathcal I}D_p.
 \label{eq:all-cm-exact-products}
\end{equation}
Consequently
\[
                         C/D_0\simeq\prod_{p\in\mathcal I}C_p/D_p.
\]
These statements require no CM-unit or finite-index hypothesis.
\end{proposition}

\begin{proof}
Choose an arithmetic inertia element at \(p\) whose
cyclotomic reduction modulo \(p\) is primitive.
Its induced operator \(S_p\) on the ambient scalar product
acts by \(s_p\) in the \(p\)-coordinate and by the identity
in every other coordinate, since cyclotomic inertia at \(p\)
is trivial modulo every \(q\ne p\).
Both \(C\) and \(D_0\) are stable under \(S_p\).

Divide the integral polynomial \(Q_p(X)-\varepsilon_p\)
by \(X-1\), and write
\begin{equation}
                     Q_p(X)=\varepsilon_p+(X-1)F_p(X),
                     \qquad F_p\in\mathbb Z[X].
 \label{eq:all-cm-polynomial-division}
\end{equation}
For \(x\in C\), the following identity holds in the whole
ambient product:
\begin{equation}
 \varepsilon_p i_p(x_p)
       =i_p(Q_px_p)-(S_p-1)F_p(S_p)x.
 \label{eq:all-cm-exact-coordinate-projection}
\end{equation}
At \(p\), it is \eqref{eq:all-cm-polynomial-division};
at every other coordinate the second term is zero.
The first term belongs to \(D_0\) by
Lemma~\ref{lem:all-cm-isolated-axes}, and the second
belongs to \(C\) by stability.
Since \(\varepsilon_p=\pm1\), it follows that
\(i_p(C_p)\subseteq C\).

For \(x\in D_0\), both terms on the right of
\eqref{eq:all-cm-exact-coordinate-projection}
belong to \(D_0\), so \(i_p(D_p)\subseteq D_0\).
Finite sums of coordinate axes are dense in the
corresponding products. The images \(C,D_0\) are
compact and hence closed, proving
\eqref{eq:all-cm-exact-products}.
This also covers infinite and empty index sets.
The quotient isomorphism is induced by the coordinate
quotient maps and is a continuous bijection between
compact Hausdorff groups.
\end{proof}

The proposition describes the full and arithmetic images \(C_p,D_p\).
It does not assert that either of them equals \(P_p\).

\subsection{The full continuous conjugation module}

Set
\[
 E_p=C_p/D_p,\qquad
 U_p=\operatorname{Hom}_{\rm cont}(E_p,\mathbb F_2),\qquad
 U=\operatorname{Hom}_{\rm cont}(C/D_0,\mathbb F_2).
\]
Give the character groups the discrete topology and the
left conjugation action
\((g\chi)(x)=\chi(g^{-1}xg)\).
Proposition~\ref{prop:all-cm-exact-products} gives
\begin{equation}
 C/D_0\simeq\prod_{p\in\mathcal I}E_p,\qquad
 U\simeq\bigoplus_{p\in\mathcal I}U_p.
 \label{eq:all-cm-character-decomposition}
\end{equation}
Indeed the open kernel of a continuous character on the
product contains the entire factors outside a finite set.
Conversely every finite sum of coordinate characters is
continuous. The decomposition is equivariant, and the
action on \(U_p\) factors through \(\lambda\bmod p\).

\begin{lemma}\label{lem:all-cm-coordinate-module}
Each \(U_p\) is finite, and \(s_p-1\) is bijective on it.
\end{lemma}

\begin{proof}
The principal-unit group \(1+2\mathcal O_p\) has the
finite-index subgroup \(1+4\mathcal O_p\), which logarithm
identifies with \(4\mathcal O_p\simeq\mathbb Z_2^{p-1}\).
It is therefore a finitely generated abelian pro-two group.
A closed subgroup has the same finiteness property:
its intersection with the logarithmic subgroup is a
closed submodule of a finite free \(\mathbb Z_2\)-module,
and its remaining quotient is finite.
Thus \(C_p\) and its quotient \(E_p\) are finitely generated,
so \(U_p\) is finite. This argument retains all finite torsion.

By Lemma~\ref{lem:all-cm-isolated-axes}, \(Q_p\) annihilates
\(E_p\). Equation~\eqref{eq:all-cm-polynomial-division}
therefore gives
\[
 (s_p-1)(-\varepsilon_p F_p(s_p))
       =1=(-\varepsilon_p F_p(s_p))(s_p-1)
                         \quad\text{on }E_p.
\]
Hence \(s_p-1\) is a continuous automorphism of \(E_p\).
The dual action is precomposition with \(s_p^{-1}\),
and \(s_p^{-1}-1=-s_p^{-1}(s_p-1)\) is also invertible.
Duality proves the assertion on \(U_p\).
\end{proof}

\begin{corollary}\label{cor:all-cm-conjugation-vanishing}
With its discrete conjugation action, \(U\) is a continuous
\(\Gamma_0=G/R_0\)-module, and
\begin{equation}
                       H^j_{\rm cont}(\Gamma_0,U)=0
                                      \qquad(j\ge0).
 \label{eq:all-cm-conjugation-cohomology}
\end{equation}
\end{corollary}

\begin{proof}
Each element of \(U\) has finite coordinate support.
The action on each coordinate factors through a finite
cyclotomic group, so each element has an open stabilizer.
This proves continuity for the discrete topology.

Apply Theorem~\ref{thm:omitted-prime-vanishing} with
\(\Lambda=\varnothing\) and \(R_\Lambda=R_0\).
The required factorization is
\eqref{eq:general-cm-residual-intersection}, and
Lemma~\ref{lem:all-cm-coordinate-module} supplies the
coefficient hypotheses. Thus
\(H^j(\Gamma_0,U_p)=0\) for every \(p,j\).
The same holds for every finite sum
\(U_F=\bigoplus_{p\in F}U_p\).

For every \(n\ge0\), a continuous cochain
\(\Gamma_0^n\to U\) has finite image because its domain
is compact and \(U\) is discrete.
The union of the supports of those finitely many values
is finite. Consequently, as cochain complexes,
\begin{equation}
 C^\bullet_{\rm cont}(\Gamma_0,U)
      =\bigcup_{\substack{F\subseteq\mathcal I\\F\ {\rm finite}}}
                            C^\bullet_{\rm cont}(\Gamma_0,U_F).
 \label{eq:all-cm-continuous-cochain-union}
\end{equation}
Every differential preserves these subcomplexes.
A positive-degree cocycle therefore belongs to a finite
summand complex and is a coboundary there.
In degree zero an invariant element likewise belongs
to a finite summand with zero invariant space.
This proves \eqref{eq:all-cm-conjugation-cohomology}
for the full, possibly infinite module.
\end{proof}

\subsection{The remaining relative central extension class}

Let
\[
 J=\ker\bigl((h_p)_{p\in\mathcal I}:R_0\to C\bigr),\qquad
 H_0=R_0\cap H,
\]
and let
\[
 \mathscr X(J)=
 \{\kappa:J\to\mathbb F_2:
   \kappa\text{ is continuous and }G\text{-invariant},\
   \kappa(J\cap H)=0\}.
\]
For \(\mathcal I\) equal to all primes at least five,
one has \(J=R_{\rm all}\).

\begin{theorem}\label{thm:all-cm-central-extension}
For \(\kappa\in\mathscr X(J)\), the central extension
\(R_0/\ker\kappa\to C\), with its prescribed arithmetic
splitting over \(D_0\), defines a continuous relative class
\[
                         o_2(\kappa)\in
                                H^2_{\rm cont}(C,D_0;\mathbb F_2)
\]
when \(\kappa\ne0\); set \(o_2(0)=0\).
The assignment \(o_2\) is \(\mathbb F_2\)-linear.
The character \(\kappa\) extends to a continuous character
of \(G\) vanishing on \(H\) if and only if \(o_2(\kappa)=0\).
When it exists, this extension is unique.
Equivalently, restriction identifies
\begin{equation}
 \operatorname{Hom}_{\rm cont}(B,\mathbb F_2)
   \simeq
 \ker\bigl(o_2:\mathscr X(J)\longrightarrow
                       H^2_{\rm cont}(C,D_0;\mathbb F_2)\bigr).
 \label{eq:all-cm-central-kernel-identification}
\end{equation}
\end{theorem}

\begin{proof}
For \(\kappa\ne0\), \(G\)-invariance makes \(\ker\kappa\)
normal in \(G\) and gives a central extension
\[
 1\longrightarrow\mathbb F_2
   \longrightarrow E:=R_0/\ker\kappa
   \xrightarrow{\pi}C\longrightarrow1.
\]
The image of \(H_0\) in \(E\) maps isomorphically onto
\(D_0\): its kernel comes from \(J\cap H\), where
\(\kappa\) is zero.
Compactness gives a prescribed continuous splitting
\(\sigma_D:D_0\to E\).

We make the continuous relative class explicit.
A finite-kernel profinite extension admits a continuous
set-theoretic section.
Choose an open normal subgroup of \(E\) meeting the finite
kernel trivially. It maps homeomorphically onto its open
image in \(C\), and finitely many translated copies give
a section on all of \(C\).
The discrepancy between this section on \(D_0\) and
\(\sigma_D\) is a continuous finite-valued function.
A finite clopen partition of a closed subset of a
profinite space extends to a finite clopen partition
of the whole space, by disjoint clopen separation.
Thus the discrepancy extends to \(C\), and adjusting
the section gives a continuous section \(s:C\to E\)
agreeing with \(\sigma_D\).

Its continuous cocycle
\(s(x)s(y)s(xy)^{-1}\) has values in \(\mathbb F_2\)
and vanishes on \(D_0\times D_0\).
Changing \(s\) while keeping this prescribed restriction
changes the cocycle by the coboundary of a continuous
one-cochain zero on \(D_0\).
This describes its relative class.
It also agrees with the relative mapping-cone definition:
restriction of continuous finite-valued cochains from
\(C^n\) to \(D_0^n\) is surjective by the same clopen
extension argument, so the kernel relative complex
computes the relative cohomology.
The class is zero exactly when a continuous
one-cochain zero on \(D_0\) changes \(s\) into a
homomorphic splitting extending \(\sigma_D\).

A compatible splitting identifies
\(E\simeq\mathbb F_2\times C\).
Projection onto the first factor, pulled back to \(R_0\),
gives a continuous character \(\ell:R_0\to\mathbb F_2\)
extending \(\kappa\) and zero on \(H_0\).
Conversely such a character factors through \(E\),
is the identity on the central \(\mathbb F_2\),
and kills \(\sigma_D(D_0)\).
Its kernel maps homeomorphically onto \(C\), giving
the required compatible splitting.
For \(\kappa=0\), the zero character provides such
an extension and \(o_2(0)=0\).
Pushout and Baer sum preserve the prescribed arithmetic
splitting. They identify the class associated with
\(\kappa+\kappa'\) with \(o_2(\kappa)+o_2(\kappa')\),
so \(o_2\) is \(\mathbb F_2\)-linear.

Suppose \(o_2(\kappa)=0\).
The possible characters \(\ell\) on \(R_0\) form a torsor
under the full module
\(U=\operatorname{Hom}_{\rm cont}(C/D_0,\mathbb F_2)\):
the difference of two extensions kills \(J\), factors
continuously through \(C\), and kills \(D_0\).
Conjugation by \(G\) preserves this torsor.
The subgroup \(H\) preserves \(H_0\) and the specified
restriction on \(J\), while inner conjugation by \(R_0\)
acts trivially on characters of \(R_0\).
Since \(G=R_0H\), the action factors through \(\Gamma_0\).

This is a continuous action on a discrete torsor, even
when \(U\) is infinite.
For any continuous \(\ell:R_0\to\mathbb F_2\), choose an
open normal subgroup \(V\triangleleft G\) with
\(V\cap R_0\subseteq\ker\ell\).
Then \([V,R_0]\subseteq V\cap R_0\), so \(V\) fixes \(\ell\).
The stabilizer, and its image in \(\Gamma_0\), is open.

Choose one extension \(\ell_0\).
The differences
\[
                         z(\gamma)=\gamma\ell_0-\ell_0
                                      \qquad(\gamma\in\Gamma_0)
\]
form a continuous \(U\)-valued one-cocycle.
Corollary~\ref{cor:all-cm-conjugation-vanishing} gives
\(H^1(\Gamma_0,U)=0\), so
\(z(\gamma)=\gamma u-u\) for some \(u\in U\).
The extension \(\ell_0-u\) is \(G\)-invariant.
It is the unique invariant extension, because the same
corollary gives \(U^{\Gamma_0}=0\).

Finally define
\[
                         \chi(rh)=\ell(r)
                            \qquad(r\in R_0,\ h\in H)
\]
using this invariant extension.
The condition \(\ell(H_0)=0\) proves well-definedness;
\(H\)-invariance proves the homomorphism identity.
Its pullback along the compact surjection
\(R_0\times H\to G\) is continuous, so \(\chi\) is continuous.
It extends \(\kappa\) and vanishes on \(H\).
Conversely a character of \(G\) with these properties
restricts to a compatible character on \(R_0\), forcing
\(o_2(\kappa)=0\).
The difference of two such characters is zero on \(R_0\)
by \(U^{\Gamma_0}=0\), and zero on \(H\); hence it is zero
on \(G=R_0H\).
Characters of \(G\) vanishing on \(H\) are the continuous
characters of \(B\). The criterion and uniqueness therefore
give \eqref{eq:all-cm-central-kernel-identification}.
\end{proof}

For the collection of all prime degrees, this theorem
eliminates the conjugation obstruction in
Proposition~\ref{prop:all-cm-intrinsic-extension}.
The relative central class remains a condition.
No arithmetic factorization \(G=JH\), equality of the
individual full and arithmetic scalar images, or
vanishing of \(\operatorname{Hom}_{\rm cont}(B,\mathbb F_2)\)
follows from these results.


\subsection{Relative central classes of product pairs}
\label{sec:cm-cross-prime-central}

Let \(\{(C_i,D_i)\}_{i\in I}\) be abelian profinite pairs, with
\(D_i\) closed in \(C_i\), and put
\[
 C=\prod_{i\in I}C_i,\qquad D=\prod_{i\in I}D_i,
 \qquad k=\mathbf F_2.
\]
The coefficient action of all these groups on \(k\) is trivial.
Write
\[
 W_i=\operatorname{Hom}_{\rm cont}(C_i,k),\quad
 V_i=\operatorname{Hom}_{\rm cont}(D_i,k),\quad
 U_i=\ker(W_i\longrightarrow V_i).
\]
For distinct indices let
\[
 K_{ij}=\ker\bigl(W_i\otimes_k W_j
                  \longrightarrow V_i\otimes_k V_j\bigr).
\]
Equivalently, \(K_{ij}\) is the space of continuous bilinear maps
\(C_i\times C_j\to k\) that vanish on \(D_i\times D_j\).
Every such bilinear map factors through finite elementary abelian
quotients, which justifies its identification with a tensor of the
continuous character spaces. Linear algebra over the field \(k\)
gives
\begin{equation}
 K_{ij}=U_i\otimes W_j+W_i\otimes U_j
               \quad\hbox{inside }W_i\otimes W_j.       \label{eq:cm-mixed-central-module}
\end{equation}
No surjectivity of \(W_i\to V_i\) is asserted or needed.

Fix an ordering of \(I\) to write each unordered pair once.

\begin{proposition}[Relative product decomposition]
\label{prop:cm-relative-product-cohomology}
There is a natural isomorphism
\begin{equation}
 H^2_{\rm cont}(C,D;k)
 \simeq
 \bigoplus_{i\in I}H^2_{\rm cont}(C_i,D_i;k)
 \ \oplus\!\!\bigoplus_{i<j}K_{ij}.                    \label{eq:cm-relative-product-cohomology}
\end{equation}
It is equivariant for every group acting continuously and
coordinatewise on the pairs. The first coordinates are restrictions
to individual pairs. The mixed coordinate of a relative central
extension is the commutator pairing between its two indicated
coordinate subgroups. All sums are algebraic direct sums, so a class
has finite coordinate support.
\end{proposition}

\begin{proof}
First suppose \(I\) is finite. We use the standard elementary
description of continuous relative degree--two classes as central
extensions by \(k\), equipped with a prescribed homomorphic splitting
over \(D\), up to equivalence preserving that splitting. It follows
directly from normalized continuous cocycles: a cocycle which is zero
on \(D\times D\) defines the group law on \(k\times C\), and
relative coboundaries are changes of section that are zero on \(D\).
Conversely a finite-kernel profinite extension has a continuous
set-theoretic section, and this section can be chosen to agree with
the prescribed one on the closed subgroup \(D\).

For completeness, choose an open subgroup of the extension meeting
the finite kernel trivially. Its map to its open image in \(C\)
is a homeomorphism; choosing finitely many coset representatives
constructs a continuous section. The discrepancy from the prescribed
section on \(D\) is a continuous finite-valued function. Such a
function on a closed subset of a profinite space extends to the
whole space by clopen separation, so the section can be corrected.

Let \(\widetilde C\to C\) represent a relative central class.
Because \(C\) is abelian, all commutators in \(\widetilde C\)
lie in its central kernel. Consequently the commutator of lifts of
\(x_i\in C_i\) and \(y_j\in C_j\) is independent of their
lifts and is a continuous bilinear function \(b_{ij}(x_i,y_j)\).
The prescribed splitting is a homomorphism on \(D\), so
\(b_{ij}\) vanishes on \(D_i\times D_j\). This gives the map
in \eqref{eq:cm-relative-product-cohomology}, together with the individual restrictions. The construction
respects Baer sum.

The map is surjective. An individual relative extension can be
pulled back by \(C\to C_i\), using the corresponding prescribed
splitting on \(D\); this has the stated individual component and
zero other components. For \(b\in K_{ij}\), the continuous
normalized cocycle
\[
 \alpha_b(x,y)=b(x_i,y_j)
\]
defines a relative extension because it is zero on
\(D\times D\). Its individual restrictions vanish, and its only
mixed commutator pairing is \(b\). Finite Baer sums realize all
the prescribed coordinates.

For injectivity, suppose all individual and mixed components
vanish. The restriction over each \(C_i\) then has a continuous
homomorphic section \(\sigma_i\) agreeing with the prescribed
section on \(D_i\). Vanishing of every mixed commutator says that
these sections have mutually commuting images. Their product
\[
 \sigma((x_i)_i)=\prod_i\sigma_i(x_i)
\]
is therefore a continuous homomorphic section over \(C\).
It agrees with the prescribed section on \(D\), since that
section is itself a homomorphism and agrees on each factor.
The original relative class is zero. This proves \eqref{eq:cm-relative-product-cohomology} for finite
products. Each step is coordinatewise natural, giving equivariance.

Now allow \(I\) to be infinite. In every degree, restriction of
continuous \(k\)-valued cochains from \(C\) to \(D\) is
surjective by the same clopen extension argument. Thus the kernel
of restriction is a relative cochain model. A continuous function
\(C^n\to k\) depends on finitely many coordinates: choose
cylinder neighborhoods on which it is constant, take a finite
subcover by compactness, and collect the finitely many coordinate
sets used by that subcover.

If such a cochain is zero on \(D^n\) and depends on a finite set
\(F\subset I\), the induced cochain on \(C_F^n\) is zero on
\(D_F^n\), where \(C_F=\prod_{i\in F}C_i\) and similarly
for \(D_F\). Hence the relative cochain complex is the filtered
union of the finite-coordinate relative complexes. Filtered unions
of vector spaces are exact, so its cohomology is the corresponding
filtered colimit. Pullback from a finite product is injective because
coordinate inclusion is a retraction as a map of pairs. Applying the
finite-product calculation proves \eqref{eq:cm-relative-product-cohomology}, with finite support as
stated.
\end{proof}

\subsection{Eliminating the mixed invariant summands}

Let a group \(\mathcal G\) act coordinatewise on the pairs.
Suppose that for every \(i\) there is an operator \(S_i\) induced
by an element of \(\mathcal G\), acting trivially at every
coordinate other than \(i\), such that its action on \(C_i\)
has finite order dividing \(2^{a_i}d_i\), with \(d_i\) odd.
Assume further that \(S_i-1\) acts bijectively on \(U_i\).
Here and below the action on character spaces is the dual action.

\begin{theorem}[Absence of mixed invariant central classes]
\label{thm:cm-cross-prime-central}
Under these hypotheses,
\begin{equation}
 H^2_{\rm cont}(C,D;k)^{\mathcal G}
 \simeq
 \bigoplus_{i\in I}
       H^2_{\rm cont}(C_i,D_i;k)^{\mathcal G}.         \label{eq:cm-prime-central-invariants}
\end{equation}
The isomorphism is restriction; its inverse is the finite sum of
coordinate pullbacks. No assertion that the individual invariant
groups vanish is part of the theorem.
\end{theorem}

\begin{proof}
Set \(T_i=S_i^{2^{a_i}}\), so its action on the \(i\)-th
coordinate has odd order dividing \(d_i\). In characteristic two,
\[
 T_i-1=(S_i-1)^{2^{a_i}}
                    \quad\hbox{on }U_i.
\]
This is invertible by hypothesis. The averaging operator
\[
 N_i=\sum_{a=0}^{d_i-1}T_i^a
\]
is the projection onto the \(T_i\)-fixed subspace: the usual
normalizing factor \(d_i^{-1}\) is one in \(k\). In particular
\(N_iU_i=0\). Its action on every other coordinate is the identity.

Let \(b\in K_{ij}^{\mathcal G}\). Since \(T_i,T_j\) come
from elements of \(\mathcal G\), invariance gives
\(N_iN_jb=b\). On the other hand \eqref{eq:cm-mixed-central-module} expresses \(b\) as
a sum from \(U_i\otimes W_j\) and \(W_i\otimes U_j\).
The commuting operator product \(N_iN_j\) kills the first
summand by \(N_iU_i=0\), and the second by \(N_jU_j=0\).
Thus \(b=0\).

The decomposition \eqref{eq:cm-relative-product-cohomology} is equivariant and every class has finite
support. Taking invariants therefore gives precisely \eqref{eq:cm-prime-central-invariants}.
\end{proof}

\subsection{The individual-prime central obstructions}

Use the notation of Proposition~\ref{prop:all-cm-exact-products}:
\[
 C=\prod_{p\in\mathcal I}C_p,\qquad
 D_0=\prod_{p\in\mathcal I}D_p,\qquad
 U_p=\operatorname{Hom}_{\rm cont}(C_p/D_p,\mathbf F_2).
\]
These are the full and arithmetic images; neither is replaced
by an ambient torus or a lower bound for the arithmetic image.

\begin{corollary}\label{cor:cm-individual-central-obstructions}
Restriction induces an isomorphism
\[
 H^2_{\rm cont}(C,D_0;\mathbf F_2)^G
 \simeq
 \bigoplus_{p\in\mathcal I}
       H^2_{\rm cont}(C_p,D_p;\mathbf F_2)^G.
\]
Every intrinsic relative central obstruction in
Theorem~\ref{thm:all-cm-central-extension} has only finitely many
nonzero individual-prime components and has no mixed component.
The corresponding character extends to a continuous character of
\(G\) vanishing on \(H\) if and only if all its individual-prime
components vanish. Such an extension is unique.
\end{corollary}

\begin{proof}
A primitive arithmetic \(p\)-inertia element induces \(s_p\)
at \(p\) and the identity at every other coordinate. Its coordinate
action has order dividing \(p-1\). By
Lemma~\ref{lem:all-cm-coordinate-module}, \(s_p-1\) is invertible
on \(U_p\). Theorem~\ref{thm:cm-cross-prime-central} therefore
applies to the exact products of
Proposition~\ref{prop:all-cm-exact-products}.

Each intrinsic relative central obstruction is invariant. Conjugation
by the arithmetic subgroup preserves its prescribed arithmetic
splitting, and the arithmetic subgroup supplies the same action on
\(C\) as \(G=R_0H\), since \(R_0\) acts trivially on the
abelian scalar image. The decomposition and finite-support assertion
follow. The extension assertion and uniqueness now follow from
Theorem~\ref{thm:all-cm-central-extension}.
\end{proof}

This result leaves the individual relative central groups to be
determined. In particular, it does not give a finite list of all
nonzero individual groups as the primes vary. At degree \(163\),
when \((C_{163},D_{163})=(P,A)\), the quaternion line of
Proposition~\ref{prop:163-relative-cohomology} is one of the
individual invariant summands. Its vanishing for full GT is not
decided by the product decomposition.

\subsection{Vanishing of the degree-\(163\) conjugation obstruction}
\label{sec:163-concrete-odd-inertia}

Put
\[
 R=R_{\rm global},\qquad \Gamma=G/R,\qquad
 P=P_{163},\qquad A=Q_{163}P,\qquad U=(P/A)^\vee.
\]
The module \(U\) is the nontrivial irreducible two-dimensional
\(\mathbb F_2\)-module for the cyclotomic quotient
\[
 \Gamma\xrightarrow{\lambda\bmod163}C_{162}\longrightarrow C_3;
\]
see Lemma~\ref{lem:163-integral-pair-decomposition}.
The omitted-prime theorem applies to the exact defining
intersection of \(R\), independently of its full degree-\(163\) image.

\begin{theorem}\label{thm:163-concrete-odd-inertia}
For every \(j\ge0\),
\begin{equation}
                         H^j_{\rm cont}(\Gamma,U)=0.
 \label{eq:163-concrete-cohomology-vanishing}
\end{equation}
\end{theorem}

\begin{proof}
Theorem~\ref{thm:simultaneous-global-cm-normalization} gives
\[
 R=R_{\rm quad}\cap
       \bigcap_{q\in\Pi_{\rm global}}\ker h_q,\qquad
 G=RH,\qquad 163\notin\Pi_{\rm global}.
\]
A primitive generator \(s\) of \(C_{162}\) acts on \(U\)
through a nontrivial element of \(C_3\).
Its minimal polynomial is \(X^2+X+1\), so \(s-1\) is
bijective on \(U\).
Theorem~\ref{thm:omitted-prime-vanishing}, with
\(\Lambda=\Pi_{\rm global}\), \(p=163\), and \(W=U\),
therefore proves
\eqref{eq:163-concrete-cohomology-vanishing}.
\end{proof}

\begin{corollary}\label{cor:163-single-extension-obstruction}
Suppose \(h_{163}(R)=P\), and put
\(J=\ker(h_{163}:R\to P)\).
For \(\kappa\in\mathscr X(J)\), the second obstruction
\(o_1(\kappa)\) of
Proposition~\ref{prop:163-two-stage-extension} is always
zero whenever it is defined.
The character \(\kappa\) extends to a continuous
character of \(G\) vanishing on \(H\) if and only if
\[
                   o_2(\kappa)=0
                       \quad\text{in }\mathbb F_2\,\mathfrak q.
\]
When it exists, this extension is unique.
Equivalently, restriction identifies
\begin{equation}
 \operatorname{Hom}_{\rm cts}(B,\mathbb F_2)
   \simeq
 \ker\bigl(o_2:\mathscr X(J)\longrightarrow
                               \mathbb F_2\,\mathfrak q\bigr).
 \label{eq:163-remaining-relative-obstruction}
\end{equation}
\end{corollary}

\begin{proof}
Theorem~\ref{thm:163-concrete-odd-inertia} gives
\(H^1(\Gamma,U)=0\), so the conjugation obstruction in
Proposition~\ref{prop:163-two-stage-extension} vanishes.
That proposition supplies the remaining condition,
uniqueness, and injectivity of restriction.
Characters of \(G\) vanishing on \(H\) are the
continuous characters of the relative abelian quotient
\(B\), giving \eqref{eq:163-remaining-relative-obstruction}.
\end{proof}

The theorem is independent of which full-image
alternative in
Proposition~\ref{prop:cm163-full-image-dichotomy} holds.
It does not decide between \(h_{163}(R)=A\) and
\(h_{163}(R)=P\), and does not show that the relative
quaternion obstruction \(o_2\) vanishes.
Thus the original question concerning
\(\operatorname{Hom}_{\rm cts}(B,\mathbb F_2)\) remains open.

\subsection{A binary tetrahedral lifting criterion}
\label{sec:163-binary-tetrahedral}

The remaining central obstruction has an exact finite lifting
interpretation. The criterion below is conditional on the larger
scalar image and does not establish the existence of the lift.

Write
\[
 G=\widehat{\mathrm{GT}},\quad H=G_{\mathbf Q},\quad
 R=R_{\rm global},\quad \Gamma=G/R,
 \qquad G=RH.
\]
Suppose that the larger degree-163 image in
Proposition~\ref{prop:cm163-full-image-dichotomy} occurs.  Thus
\[
 h:R\twoheadrightarrow P,\qquad h(R\cap H)=A,\qquad
 E=P/A\simeq\mathbf F_2^2,
 \qquad J=\ker h.
\]
The cyclotomic action on \(E\) factors through a surjection
\(\gamma:\Gamma\to C_3\) and is faithful and irreducible.
Proposition~\ref{prop:163-relative-cohomology} identifies the invariant
relative degree-two group with the line generated by the pullback of
\(Q_8\to E\). Theorem~\ref{thm:163-concrete-odd-inertia} gives
\[
 H^1_{\rm cont}(\Gamma,E^\vee)=0.
\]
The polarization of \(q(z)=z^3\), namely
\(\langle z,w\rangle=zw^2+z^2w\), identifies \(E\) with
\(E^\vee\) equivariantly.  Therefore
\begin{equation}
 H^1_{\rm cont}(\Gamma,E)=0.                 \label{eq:163-selfdual-h1-vanishing}
\end{equation}

The normalized relative degree--163 cocycle and its finite quotient are
\[
 z(rh_0)=h(r)\bmod A,\qquad
 \pi(g)=(z(g),\gamma(g)):
 G\twoheadrightarrow E\rtimes C_3\simeq A_4.
\]
Here \(h_0\in H\); the symbol \(h\) without a subscript denotes
the scalar homomorphism.  In particular \(z|_H=0\).
Let
\[
 \widetilde A_4=Q_8\rtimes C_3,
 \qquad p:\widetilde A_4\to E\rtimes C_3
\]
where the chosen order-three automorphism cycles the three quaternion
directions.  This is the binary tetrahedral group
\(\mathrm{SL}_2(\mathbf F_3)\). Explicitly, in the latter group
\[
 i=\begin{pmatrix}0&1\\-1&0\end{pmatrix},\qquad
 j=\begin{pmatrix}1&1\\1&-1\end{pmatrix}
\]
satisfy \(i^2=j^2=-I\) and \(ij=-ji\), so they generate
\(Q_8\). An element of order three acts cyclically on its three
noncentral pairs. This gives the semidirect product of order
\(8\cdot3=24=|\mathrm{SL}_2(\mathbf F_3)|\).  Its center is the kernel
\(C_2\), and its displayed \(C_3\) maps isomorphically to the
displayed complement in \(A_4\).

\begin{proposition}\label{prop:163-binary-tetrahedral-lift}
The following conditions are equivalent.
\begin{enumerate}
\item Some continuous, \(G\)-invariant
\(\kappa:J\to\mathbf F_2\), zero on \(J\cap H\), has
nonzero central relative obstruction \(o_2(\kappa)\).
\item There is a continuous homomorphism
\(\widetilde\pi:G\to\widetilde A_4\) such that
\(p\widetilde\pi=\pi\) and
\(\widetilde\pi(H)\) lies in the displayed \(C_3\).
\item The pullback of
\(\widetilde A_4\to A_4\) by \(\pi\), equipped with its
canonical splitting over \(H\), is split as a relative central
extension of \((G,H)\).
\end{enumerate}
Every homomorphism in condition two is surjective and has
\(\widetilde\pi(H)=C_3\).  When such homomorphisms exist, their
set is a torsor under the continuous quadratic characters of \(G\)
which vanish on \(H\).
\end{proposition}

\begin{proof}
Assume condition one and put \(K=\ker\kappa\).  Invariance makes
\(K\) a closed normal subgroup of \(G\).  Since
\(o_2(\kappa)\ne0\), Proposition~\ref{prop:163-relative-cohomology}
gives an equivalence of central extensions, preserving the specified
arithmetic splitting,
\[
 R/K\simeq P\times_E Q_8.
\]
The image of \(R\cap H\) under this equivalence is
\[
 L=\{(a,1):a\in A\}.
\]
It is central in \(R/K\).  It is also preserved by the image of
\(H\), since it is exactly the image of \(R\cap H\).
As \(G=RH\), it follows that \(L\) is normal in \(G/K\).
All groups and images involved are compact, so \(L\) is closed.

Quotienting \(G/K\) by \(L\) now gives a profinite group
\(G_1\), with subgroups \(R_1\) and \(H_1\) which are the
images of \(R\) and \(H\).  They satisfy
\[
 R_1\simeq Q_8,\qquad R_1\cap H_1=1,\qquad G_1=R_1H_1.
\]
For the intersection, the image of \(H\) in \(G/K\) meets
\(R/K\) in the image of \(R\cap H\), because \(K\subset R\).
That intersection is precisely \(L\), which has just been killed.
Moreover \(H_1\simeq H/(H\cap R)\simeq\Gamma\).  Thus
\[
 G_1\simeq Q_8\rtimes_\alpha\Gamma
\]
for a continuous action \(\alpha:\Gamma\to\operatorname{Aut}(Q_8)\).

The induced action on \(Q_8/Z(Q_8)=E\) is the specified
\(C_3\)-action.  There is an exact sequence
\[
 1\longrightarrow\operatorname{Inn}(Q_8)\simeq E
 \longrightarrow\operatorname{Aut}(Q_8)
 \longrightarrow\operatorname{GL}_2(\mathbf F_2)
 \longrightarrow1.
\]
Indeed an automorphism can send the first two quaternion generators
to arbitrary signed independent directions; there are \(6\cdot4=24\)
choices.  Those inducing the identity on the three directions are
the four inner automorphisms.  Let
\(\theta:C_3\to\operatorname{Aut}(Q_8)\) be the displayed
order-three lift.  The difference between \(\alpha\) and
\(\theta\gamma\) is a continuous cocycle
\(v:\Gamma\to E\), through
\[
 \alpha(\delta)=\operatorname{Inn}(v(\delta))
                         \theta(\gamma(\delta)).
\]
Equation~\eqref{eq:163-selfdual-h1-vanishing} says that \(v\) is a coboundary.  Changing the
identification of \(Q_8\) by one inner automorphism therefore makes
\(\alpha=\theta\gamma\).  This change acts trivially on \(E\).
The map
\[
 Q_8\rtimes_\alpha\Gamma\longrightarrow Q_8\rtimes C_3,
 \qquad(q,\delta)\longmapsto(q,\gamma(\delta))
\]
then gives the required lift of the original \(\pi\), with
the required arithmetic complement.  It is onto.

Conversely, suppose condition two holds.  The image of \(R\)
lies in \(Q_8\), and maps onto \(E\), because \(z(R)=E\).
No proper subgroup of \(Q_8\) maps onto \(E\): every proper
subgroup is cyclic or central.  Hence
\(\widetilde\pi(R)=Q_8\).  Its commutator subgroup is the
central \(C_2\), and \([R,R]\subset J\).  Therefore the
restriction of \(\widetilde\pi\) to \(J\) is a nonzero
character
\[
 \kappa:J\longrightarrow Z(Q_8)=\mathbf F_2.
\]
It is \(G\)-invariant, since this center is central in
\(\widetilde A_4\), and it vanishes on \(J\cap H\), since
the displayed \(C_3\) meets the center trivially.

The homomorphism
\[
 R\longrightarrow P\times_E Q_8,
 \qquad r\longmapsto(h(r),\widetilde\pi(r))
\]
is surjective.  It surjects to \(P\), and its kernel above
\(0\in P\) maps onto the central \(C_2\), by the preceding
nonvanishing of \(\kappa\).  Hence every two-element fiber is
covered.  Its kernel is \(\ker\kappa\), and it sends the
arithmetic intersection to \(\{(a,1):a\in A\}\).
Consequently its relative class is exactly the quaternion class.
This proves condition one.

The equivalence of conditions two and three is the usual elementary
pullback description: a relative splitting is a homomorphic section
\(g\mapsto(g,\widetilde\pi(g))\), whose prescribed value on
\(H\) is the canonical lift through its \(C_3\)-image.
If two such lifts exist, their pointwise ratio lies in the common
central kernel and is a continuous character \(G\to\mathbf F_2\)
vanishing on \(H\).  Conversely multiplying a lift by any such
character gives another lift.  Finally the image of \(H\) is all
\(C_3\), since \(\gamma(H)=C_3\).
\end{proof}

The proposition identifies the finite quotient which would witness
a nonzero central obstruction.  No exclusion or realization of this
binary tetrahedral quotient for full GT has been proved here.  In
particular the result is not a proof that \(o_2\) is zero or nonzero,
does not decide \(h(R)=A\) versus \(P\), and does not decide
\(\operatorname{Hom}_{\rm cont}(B,\mathbf F_2)\).

\paragraph{Commutator evaluation.}
There is also a direct evaluation test.  If \(r_1,r_2\in R\) map to
an \(\mathbf F_2\)-basis of \(E\), then
\[
 o_2(\kappa)\ne0
 \quad\Longleftrightarrow\quad
 \kappa([r_1,r_2])=1.
\]
Indeed the unique invariant nonzero quadratic form has nonzero
polarization on every independent pair.  This test still needs a
character \(\kappa\) or a full-GT identity; it does not
produce either one.

\subsection{A fixed finite word for the degree-\(163\) obstruction}
\label{sec:163-finite-word}

Use \(R=R_{\rm global}\), \(P=P_{163}\), \(A=Q_{163}P\),
and \(E=P/A\). Put
\[
 \mathcal O=\mathbb Z_2[\zeta_{163}],\quad
 s(\zeta_{163})=\zeta_{163}^2,\quad c=s^{81},\quad
 W=\ker\bigl(\operatorname{Tr}_{\mathbb F_8/\mathbb F_2}\bigr).
\]
The exact detector proved above identifies \(E\) with the
two-dimensional additive space \(W\):
\begin{equation}
 \Psi(t)=\sum_{j=0}^{26}s^{3j}
             \left((t-1)/2\bmod2\mathcal O\right),\qquad
 \ker(\Psi:P\to W)=A.
 \label{eq:163-trace-detector}
\end{equation}
Here \(W\) is a vector space inside \(\mathbb F_8\), not a
subfield of order four.

Let \(\mathcal O_4=\mathcal O/4\mathcal O\), and define the fixed
finite group
\[
 \mathscr D=(\mathcal O_4)_{\rm add}\rtimes\langle t\rangle,
 \qquad t^{163}=1,\qquad tvt^{-1}=\zeta_{163}v.
\]
Write \(d=\zeta_{163}-1\), a unit, and put
\[
 X=(0,t),\qquad Y=(d^{-1},t),\qquad n=2^{81}+1.
\]
For \(0\leq j\leq26\), let \(a_j\) be the least positive residue
of \(2^{3j}\) modulo \(163\).

\begin{proposition}
\label{prop:163-finite-word-detector}
The elements \(X,Y\) generate \(\mathscr D\), so they define
one finite quotient of \(\widehat F_2\), of order \(163\cdot4^{162}\).
For every \(r\in R\), use the conjugator \(f_{r^n}\) of the
group power \(r^n\), in the convention
\[
 \phi_f(x)=x,\qquad \phi_f(y)=f^{-1}yf,\qquad
 [x,y]=xyx^{-1}y^{-1}.
\]
Then
\begin{equation}
 \prod_{j=0}^{26}f_{r^n}(X^{a_j},Y^{a_j})
       =\bigl(2d^{-2}\Psi(h_{163}(r)),1\bigr)
       \quad\text{in }\mathscr D.
 \label{eq:163-exact-word-value}
\end{equation}
In particular,
\begin{equation}
 \prod_{j=0}^{26}f_{r^n}(X^{a_j},Y^{a_j})=1
       \quad\Longleftrightarrow\quad h_{163}(r)\in A.
 \label{eq:163-word-equivalence}
\end{equation}
The word value lies in \(2d^{-2}W\), a subgroup of order four.
\end{proposition}
\begin{proof}
The multiplication is
\((u,t^i)(v,t^j)=(u+\zeta_{163}^iv,t^{i+j})\).
It gives \([X,Y]=(1,1)\).
Conjugating this translation by powers of \(X\) gives
\((\zeta_{163}^i,1)\). These elements span \(\mathcal O_4\)
over \(\mathbb Z/4\mathbb Z\), and \(X\) generates the quotient
of order \(163\). Hence \(X,Y\) generate \(\mathscr D\).
Its derived subgroup is exactly the translation subgroup,
because \([X,(u,1)]=(du,1)\) and \(d\) is a unit.

Let \(A_f(U,V)\) be the Alexander coefficient of \(f\) relative
to the generator \([x,y]\) of the metabelian commutator module.
In that module the convention in the statement gives
\[
 [x,f^{-1}yf]
   =\bigl(1+(U-1)(V-1)A_f(U,V)\bigr)[x,y].
\]
The established diagonal compact-homology identification
therefore gives the full coefficient-two scalar
\begin{equation}
 b_r=1+d^2A_{f_r}(\zeta_{163},\zeta_{163}).
 \label{eq:163-alexander-scalar}
\end{equation}
This identity uses the representation on the full group \(G\).
On \(R\) it is linear over \(\mathcal O\), so
\(b_{r^n}=b_r^n\).

For a nonzero residue \(a\) modulo \(163\), write
\(d_a=\zeta_{163}^a-1\) and
\(\sigma_a(\zeta_{163})=\zeta_{163}^a\).
The power formula gives
\[
 X^a=(0,t^a),\qquad Y^a=(d^{-2}d_a,t^a),\qquad
 [X^a,Y^a]=(d^{-2}d_a^2,1).
\]
Consequently
\begin{equation}
 f(X^a,Y^a)
   =\bigl(d^{-2}d_a^2A_f(\zeta_{163}^a,\zeta_{163}^a),1\bigr).
 \label{eq:163-substitution}
\end{equation}
These evaluations all factor through the same finite quotient.
Explicitly,
\[
 \beta_a(v,t^i)=
       \bigl(d_a^2d^{-2}\sigma_a(v),t^{ai}\bigr)
\]
is an automorphism, \(\beta_a\beta_b=\beta_{ab}\), and
\(\beta_a(X)=X^a,\ \beta_a(Y)=Y^a\).

The scalar \(b_r\) need not be principal. Its norm is one,
and its odd Teichmuller factor has order dividing
\(2^{81}+1\), because \(c\) acts on that factor by the
\(2^{81}\)-power map. The chosen odd power thus gives
\begin{equation}
 b_{r^n}=b_r^n=h_{163}(r)^n
                     \equiv h_{163}(r)\pmod{4\mathcal O}.
 \label{eq:163-odd-power}
\end{equation}
It also preserves the image in \(E\), which has exponent two.
No identity between \(f_{r^n}\) and the word power \(f_r^n\)
is used.

Apply \eqref{eq:163-substitution} and
\eqref{eq:163-alexander-scalar} to \(f_{r^n}\).
The factors are translations, so their product has additive
coordinate
\[
 d^{-2}\sum_{j=0}^{26}s^{3j}(b_{r^n}-1)
   =2d^{-2}\sum_{j=0}^{26}s^{3j}
             \left((h_{163}(r)-1)/2\bmod2\mathcal O\right).
\]
This proves \eqref{eq:163-exact-word-value}. Multiplication
by \(2\) here denotes the canonical injection
\(\mathcal O/2\mathcal O\to2\mathcal O/4\mathcal O\);
there is no division by \(2\) in \(\mathcal O_4\).
That injection and multiplication by the unit \(d^{-2}\)
are injective. Equation \eqref{eq:163-trace-detector}
therefore proves \eqref{eq:163-word-equivalence}.
\end{proof}

The associated relative class has a unique normalized
representative:
\begin{equation}
 z(rh)=h_{163}(r)\bmod A,\qquad
 [z]\in\ker\bigl(H^1_{\rm cont}(G,E)\to
                         H^1_{\rm cont}(H,E)\bigr),\qquad z|_H=0.
 \label{eq:163-relative-twisted-class}
\end{equation}
Well-definedness follows from \(h_{163}(R\cap H)=A\);
equivariance gives the cocycle identity, and compact descent
from \(R\times H\) gives continuity.
If two cohomologous cocycles vanish on \(H\), their difference
is the coboundary of an element of \(E^H\). This fixed space
is zero, because the arithmetic cyclotomic \(C_3\)-action
is irreducible. Thus \([z]=0\) if and only if \(z=0\).
The common full image is \(A\) precisely when
\eqref{eq:163-word-equivalence} vanishes for every \(r\in R\).
Otherwise the quotient \(G\twoheadrightarrow E\rtimes C_3\simeq A_4\)
described above follows. This class uses \(G\) and the nontrivial
coefficient module \(E\); the character-extension torsor uses
\(G/R\) and \(E^\vee\).

The defining cycle relations alone have not yet been shown here
to force this universal word identity. Their diagonal
two-cycle specialization is automatic in \(\mathscr D\):
swapping \(X,Y\) acts as minus one on its commutator subgroup.
The three-cycle relation at multiplier one gives the necessary
Alexander identity
\[
 A_f(U,V)+V^{-1}A_f((UV)^{-1},U)
       +(UV)^{-1}A_f(V,(UV)^{-1})=0.
\]
Indeed \([z,x]=V^{-1}[x,y]\) and
\([y,z]=(UV)^{-1}[x,y]\) for \(z=(xy)^{-1}\).
Using two-cycle symmetry and then setting \(U=V=\zeta_{163}\)
gives
\[
 A_f(\zeta_{163},\zeta_{163})
   +(\zeta_{163}^{-1}+\zeta_{163}^{-2})
                 A_f(\zeta_{163},\zeta_{163}^{-2})=0.
\]
The second coefficient is a unit, so this equation relates
the diagonal and mixed values without making the displayed
trace vanish. The usual forgetful projection of the five-cycle
relation gives the two-cycle relation
\cite[Section~1, Proposition~1.3]{Furusho2003}.
A different use of the five-cycle relation may give more
information. No solution of a finite projected system is
being used as a realization of an element of \(G\).

Finite-field rigidity does not evaluate this detector. Guillot proves
that \(\mathrm{GT}_1(\mathrm{PSL}_2(2^s))\) is trivial, using
rigidity of generating triples with prescribed traces
\cite[Theorem~1.1, Proposition~3.2 and Corollary~3.3, pp.~1--3]{Guillot2016}.
Here the residual scalar lies in \(1+2O\), so its reduction modulo
two is already one. The test takes place modulo four. Applying
a finite-field rigidity theorem would therefore require an additional
integral lifting argument, which is not established here.

\section{Further identities at degree \(163\)}
\label{sec:163-additional-identities}

The arithmetic square identity of
Lemma~\ref{lem:163-period-unit-square} has three useful refinements.
The first identifies precisely the index-four enlargement of the
circular-unit lattice needed to contain its square roots. The second
shows that a mixed beta ratio retains the same residual scalar
coordinate. The third places the square roots in an explicit
three-point orbit of a rational map. These identities do not select
an alternative in Proposition~\ref{prop:cm163-full-image-dichotomy}.

\subsection{The circular-unit saturation defect}

Put \(p=163\), \(\zeta=\zeta_p\), \(F^+=\mathbf Q(\zeta+\zeta^{-1})\),
and let \(\mathcal E\) be its real unit group.  Let \(C\subset\mathcal E\)
be the real circular-unit group, including \(-1\).  For nonzero
\(a\in\mathbf F_p\), use the standard real circular unit
\[
 \eta_a=\zeta^{(1-a)/2}\frac{1-\zeta^a}{1-\zeta},
\]
where the half exponent is interpreted modulo \(p\).
Then \(\eta_1=1\) and \(\eta_{-a}=-\eta_a\).  The group \(C\)
is generated by \(-1\) and the \(\eta_a\), taking one representative
from each pair \(\{a,-a\}\).
The circular-unit finite-index theorem gives
\(\operatorname{rank}(C/\{\pm1\})=80\); see
\cite[p.~206, the definition and index of \(C\)]{Dummit2020}.

Write
\[
 \epsilon_k=\zeta^{2^k}+\zeta^{-2^k}\qquad(0\leq k<81),
 \qquad
 U_i=\prod_{j=0}^{26}\epsilon_{i+3j}\quad(i=0,1,2).
\]
The equality \(2^{81}=-1\pmod {163}\) makes the \(\epsilon_k\)
periodic with period \(81\).  Let \(s(\zeta)=\zeta^2\).

\begin{proposition}[Circular-unit saturation defect]
\label{prop:163-circular-unit-saturation}
There is an isomorphism
\[
 C/\{\pm1\}\simeq
 \mathbf Z^{81}/\mathbf Z(1,\ldots,1)
\]
which sends the \(k\)-th standard vector to \(\epsilon_k\).
The classes of \(U_0,U_1,U_2\) in \(C/(\{\pm1\}C^2)\)
span the irreducible two-dimensional \(\mathbf F_2[C_3]\)-module;
their sole linear relation is
\[
 [U_0]+[U_1]+[U_2]=0.
\]
In particular none of the \(U_i\) is a square modulo sign of a
circular unit.

Let \(v_i=s^i(5+\theta_1)\), with the period notation of Lemma~\ref{lem:163-period-unit-square}.
Then \(U_i=-v_i^2\), and
\[
 C^*=\langle C,v_0,v_1,v_2\rangle\subset\mathcal E,
 \qquad C^*/C\simeq\mathbf F_2^2.
\]
The quotient has the same nontrivial irreducible \(C_3\)-action.
No assertion that \(C^*=\mathcal E\) is made.
\end{proposition}

\begin{proof}
The elementary identity
\[
 \frac{\eta_{4b}}{\eta_{2b}}
 =\zeta^{-b}\frac{1-\zeta^{4b}}{1-\zeta^{2b}}
 =\zeta^b+\zeta^{-b}
\]
gives
\(\epsilon_k=\eta_{2^{k+2}}/\eta_{2^{k+1}}\).
Modulo signs, the \(81\) classes \(2^k\) enumerate all nonzero
residues modulo sign.  Their consecutive differences therefore
generate all \(\eta_a\) modulo signs.  The \(80\) units
\(\eta_{2^k}\), \(1\leq k\leq80\), generate the free group
\(C/\{\pm1\}\) of rank \(80\), so they form a basis.
It follows that the only integral relation among the \(81\)
consecutive ratios, modulo signs, is their total product relation.
In fact telescoping gives the exact value
\[
 \prod_{k=0}^{80}\epsilon_k
 =\frac{\eta_{2^{82}}}{\eta_2}=-1.
\]
This proves the displayed presentation.

Modulo two, let \(w_i\in\mathbf F_2^{81}\) be the characteristic
vector of the indices congruent to \(i\pmod3\).  The image of
\(w_i\) is \([U_i]\).  The vector
\(a_0w_0+a_1w_1+a_2w_2\) is constant precisely when
\(a_0=a_1=a_2\), since the supports of the \(w_i\) are disjoint
and cover all \(81\) indices.  Their image therefore has dimension
two with exactly the claimed relation.  The automorphism \(s\)
cycles the three nonzero classes.  This is the irreducible nontrivial
\(C_3\)-module over \(\mathbf F_2\).

Lemma~\ref{lem:163-period-unit-square} gives \(v_i^2=-U_i\in C\).
Suppose \(v_0^{a_0}v_1^{a_1}v_2^{a_2}\in C\), with
\(a_i\in\{0,1\}\).  Squaring shows
\(U_0^{a_0}U_1^{a_1}U_2^{a_2}\in\{\pm1\}C^2\), so the
previous paragraph forces all \(a_i\) equal.  Conversely
\[
 (v_0v_1v_2)^2=(-1)^3U_0U_1U_2=1,
\]
and hence \(v_0v_1v_2=\pm1\in C\).  Thus the quotient is exactly
the stated group of order four.  The action is again the cyclic
permutation of its three nonzero directions.
\end{proof}

\paragraph{A concrete character separating the two notions of square.}
The preceding presentation permits a homomorphism
\(\ell:C\to\mathbf F_2\) with
\[
 \ell(-1)=0,\quad \ell(\epsilon_0)=\ell(\epsilon_1)=1,
 \quad\ell(\epsilon_k)=0\quad(2\leq k<81).
\]
The sum of these values is zero, as required by the sole relation
modulo signs.  This is a homomorphism on the circular-unit
group, so it respects every multiplicative identity among its
elements.  Nevertheless
\[
 (\ell(U_0),\ell(U_1),\ell(U_2))=(1,1,0).
\]
It cannot extend to a homomorphism \(C^*\to\mathbf F_2\):
an extension would send \(-v_0^2\) to zero.

Consequently an argument which only treats the \(\epsilon_k\)
as generators of the circular-unit group, respects their
multiplicative identities, and kills signs and circular squares
does not force the period products to vanish.  It must establish
the additional extension across the explicit index-four lattice
\(C\subset C^*\), or use other full-GT information.
This conclusion concerns the specified conductor and specified
multiplicative data.  It is not a construction of a GT cocycle,
not a counterexample to the full pentagon, and not a claim that
all higher-conductor distribution relations preserve this character.

\paragraph{Unit signatures.}
Dummit's Proposition~2 identifies a defect of order four between
totally positive circular units and circular squares at conductor
\(163\), while the full real unit group realizes all signatures
\cite[Proposition~2, pp.~209--210]{Dummit2020}.
The explicit lattice argument above isolates the three period products
and their square roots without importing a value for the full real
class number.  The equality of that class number to four is conditional on GRH
in the cited proposition and is not used here.

\subsection{A mixed beta ratio with the same residual coordinate}

Use the full coefficient-two beta values of a multiplier-one GT element
\(r\), and set \(p=163\), \(O=\mathbf Z_2[\zeta_p]\),
\(s(\zeta_p)=\zeta_p^2\).  For a unit of \(O\), write its
principal factor with a superscript \(+\); the complementary
Teichmuller factor has odd order.  All products below are in the
commutative unit group.

\begin{lemma}[Canonical principal gamma factor on an odd cyclic group]
\label{lem:finite-cyclic-principal-gamma}
Let \(p\) be odd and let
\(b:\mathbf F_p\times\mathbf F_p\to1+2O\) be a normalized
two-cocycle for the additive group \(\mathbf F_p\), with trivial
action on the coefficient group.  There is a unique function
\(\gamma:\mathbf F_p\to1+2O\), with \(\gamma(0)=1\), such that
\[
 b(a,b)=\frac{\gamma(a)\gamma(b)}{\gamma(a+b)}.
\]
Explicitly,
\[
 \gamma(a)=\left(\prod_{c\in\mathbf F_p}b(a,c)\right)^{1/p},
\]
where the \(p\)-th root is the unique root in the pro-two group
\(1+2O\).
\end{lemma}

\begin{proof}
Write \(T(a)=\prod_c b(a,c)\).  Multiplying the cocycle identity
\(b(a,b)b(a+b,c)=b(b,c)b(a,b+c)\) over \(c\) gives
\[
 b(a,b)^p T(a+b)=T(a)T(b).
\]
The \(p\)-power map is an automorphism of the abelian pro-two
group \(1+2O\), including its signs.  Taking its inverse proves
the claimed formula.  The ratio of two solutions would be a
homomorphism from a group of odd order to a pro-two group, hence
would be trivial.
\end{proof}

Apply the lemma to
\[
 b_r(a,b)=B_r(\zeta_p^a,\zeta_p^b)^+.
\]
The full-GT beta cocycle gives its hypothesis
\cite[Proposition~1.6.1]{Ihara1999},
\cite[Section~2.3, p.~101]{Furusho2003}. Frobenius
equivariance and beta unitarity give, by uniqueness,
\[
 \gamma_r(2a)=s(\gamma_r(a)),\qquad
 \gamma_r(-a)=\gamma_r(a)^{-1}.
\]
Thus \(u_r=\gamma_r(1)\) lies in the principal norm-one torus
\(T\).  If
\[
 a_r=B_r(\zeta_p,\zeta_p),\quad
 b_r=B_r(\zeta_p^2,\zeta_p),\quad
 k_r=a_r/b_r,
\]
then, writing group-ring actions additively,
\begin{equation}
 a_r^+=(2-s)u_r,
 \qquad k_r^+=(1+[3]-2s)u_r.                 \label{eq:163-gamma-beta-ratios}
\end{equation}

\begin{proposition}[Equality of the residual coordinates]
\label{prop:163-mixed-beta-residual}
For \(r\in R_{\rm global}\), the principal ratio \(k_r^+\)
belongs to \(P=P_{163}\), and
\[
 k_r^+\bmod QP=h_{163}(r)\bmod QP.
\]
Consequently the mixed ratio retains exactly the original
degree--163 alternative.  With \(n=2^{81}+1\) and the
first-layer detector \(\Psi:P\to W\) of
\eqref{eq:163-trace-detector},
\begin{equation}
 \Psi(k_r^n)=\Psi(h_{163}(r)).               \label{eq:163-mixed-ratio-detector}
\end{equation}
The expression on the left is defined because \(k_r^n\) is
principal and belongs to \(P\).
\end{proposition}

\begin{proof}
Equation~\eqref{eq:163-alexander-scalar} gives \(a_r^+=h_{163}(r)\in P\).
The operator \(2-s\) is an automorphism of \(T\), with inverse
\[
 (2-s)^{-1}=-\sum_{m\geq0}2^m s^{-m-1}.
\]
The series converges in the endomorphisms of the compact pro-two
group.  It preserves the closed stable subgroup \(P\).  Therefore
\(u_r\in P\), and \eqref{eq:163-gamma-beta-ratios} also places \(k_r^+\) in \(P\).

Direct modular arithmetic gives \(2^{20}=-3\pmod {163}\) and
\(2^{81}=-1\pmod {163}\), so
\[
 [3]=s^{101}.
\]
On \(E=P/QP\), the relations are \(2=0\), \(s^3=1\), and
\(s^2+s+1=0\).  Consequently
\[
 1+s^{101}-2s=1+s^2=s=2-s
                 \quad\text{on }E.
\]
Equation \eqref{eq:163-gamma-beta-ratios} proves the required equality in \(E\).

The full beta values in the ratio have norm one, by coefficient
equivariance and beta unitarity.  Hence the odd Teichmuller factor
of \(k_r\) has order dividing \(2^{81}+1=n\).  Thus
\(k_r^n=(k_r^+)^n\), and this lies in \(P\).  Since \(n\) is
odd, \((k_r^+)^n\equiv k_r^+\pmod {4O}\).  The detector in \eqref{eq:163-trace-detector}
factors through this layer and kills exactly \(QP\).
This proves \eqref{eq:163-mixed-ratio-detector}.
\end{proof}

\paragraph{Scope.}
This proposition uses the full-GT beta cocycle, its coefficient
equivariance and unitarity, and the degree--163 scalar and
first-layer identifications.  It does not use an arbitrary choice of
an adelic gamma function or an extension to motivic Galois groups.
The canonical factor is constructed directly on the finite cyclic
group.  The equality says that a particular off-diagonal ratio does
not lose the unknown residual; it supplies no new reason for that
residual to vanish.

\subsection{A three-point M\"obius orbit of the square roots}

Retain the periods \(\theta_i\) of
Lemma~\ref{lem:163-period-unit-square}, numbered by multiplication
of their cosets by two, and put
\[
 v_0=5+\theta_1,\qquad v_1=5+\theta_2,\qquad
 v_2=5+\theta_0.
\]
The period identity is \(U_i=-v_i^2\).

\begin{lemma}[An explicit special fiber]
\label{lem:163-mobius-period-fiber}
The three units satisfy
\[
 v_iv_{i+1}=v_i-1,\qquad
 v_{i+1}=1-v_i^{-1}\quad(i\bmod3),
\]
and are the roots of
\[
 X^3-14X^2+11X+1.
\]
In particular they lie in the fiber above \(14\) of the
degree-three quotient map
\[
 J(v)=\frac{v^3-3v+1}{v^2-v}
\]
for the order-three transformation \(v\mapsto1-v^{-1}\).
The discriminant of the polynomial is \(163^2\).
\end{lemma}

\begin{proof}
Apply \(s\) to \eqref{eq:163-period-table} to get
\[
 \theta_1\theta_2=17\theta_1+16\theta_2+21\theta_0.
\]
Using \(\theta_0+\theta_1+\theta_2=-1\), one obtains
\[
 (5+\theta_1)(5+\theta_2)
 =25+22\theta_1+21\theta_2+21\theta_0
 =4+\theta_1=v_0-1.
\]
Its two conjugates give the other equalities.  Division is allowed
because all three \(v_i\) are units.  Their sum is \(14\).
Summing the pair products gives \(11\), and multiplying the
Mobius relations, or using \(U_0U_1U_2=-1\), gives their product
\(-1\).  More directly the first two Mobius relations give
\(v_2=-1/(v_0-1)\), so \(v_0v_1v_2=-1\).
These elementary symmetric functions give the cubic polynomial.

The transformation \(\varphi(v)=1-v^{-1}\) has order three, and
\[
 v+\varphi(v)+\varphi^2(v)
   =v+\frac{v-1}{v}-\frac1{v-1}
   =\frac{v^3-3v+1}{v^2-v}.
\]
Thus \(J\) is invariant under \(\varphi\), has degree three,
and is its quotient map.  Setting \(J(v)=14\) gives exactly
the cubic above.  The elementary cubic discriminant formula gives
\(26569=163^2\).
\end{proof}

\paragraph{What additional compatibility would be needed.}
The relations express the square roots by the standard order-three
permutation of the three punctures of the thrice-punctured line.
This is a concrete geometric structure beyond an unspecified real
unit.  Applying it requires preservation of the particular fiber
\(J=14\).  An arithmetic identification
of that fiber with the cubic period field does not itself identify
the two actions of full GT on its relevant finite covers.

A sufficient next input would be a proved full-GT comparison between
the Kummer covers of this \(J=14\) fiber and the cyclotomic
\(163\)-cover defining the three \(U_i\), including the chosen
base paths and their signs.  Such a comparison would permit the
relations \(U_i=-v_i^2\) to be applied to the full-GT first-layer
coordinates.  The two- and three-cycle equations supply the Mobius
symmetry, but have not here been shown to supply preservation of
this particular fiber or its arithmetic identification.  No such comparison has been established here.

\section{The scalar determinant of the acyclotomic twist}
\label{sec:twist-scalar-determinant}

Lochak--Schneps express an acyclotomic GT element through four
block-substitution actions on a free profinite group of rank three
\cite[Theorem~2.1, equation~(8), pp.~121--122]{LS2015Twists}.
The definitions of the four actions are equations~(9)--(12),
and their peripheral product preservation is Lemma~3.1
\cite[pp.~122--124]{LS2015Twists}. We compute what their scalar
Fox determinants retain.

There is a printed error in equation~(10): the first conjugator in
the displayed image of \(z\) has first argument \(y\).
The block-substitution definition gives instead
\[
 F^{2,3}(z)=f(z,(yz)^{-1})\,z\,f((yz)^{-1},z).
\]
We use this block-substitution action. It is also the one compatible
with the two-generator action and the product assertion of
Lemma~3.1 in the cited source.

Use \([x,y]=xyx^{-1}y^{-1}\), and write
\[
 [f]=R_f(X,Y)[x,y],\qquad
 B_f(X,Y)=1+(X-1)(Y-1)R_f(X,Y).
\]
All Fox derivatives are abelianized. They can be computed in the
completed coefficient-two group algebra and then specialized to
finite abelian-cover quotients. For an IA automorphism \(\alpha\)
of a free profinite group, let \(J_\alpha\) be the matrix of
abelianized Fox derivatives, with the image of the \(i\)-th
generator in row \(i\). Its determinant is multiplicative on IA
automorphisms, since their action on the coefficient algebra is
trivial.

\begin{proposition}\label{prop:twist-determinant-beta}
For the four block-substitution actions of an acyclotomic element
\((1,f)\), the Fox determinants are respectively
\[
 B_f(X,Y),\qquad B_f(Y,Z),\qquad
 B_f(XY,Z),\qquad B_f(X,YZ)
\]
for \(F^{1,2},F^{2,3},F^{12,3},F^{1,23}\).
The action on the rank-three fiber of the pure-braid proword
\(f(x_{12},x_{23})\) has determinant one. Consequently, at
the determinant level the twist relation gives exactly
\begin{equation}
 B_f(X,Y)B_f(XY,Z)=B_f(X,YZ)B_f(Y,Z).
 \label{eq:twist-determinant-beta}
\end{equation}
This determinant operation supplies no scalar equation beyond the
beta two-cocycle identity.
\end{proposition}

\begin{proof}
First consider \(\phi_f(x)=x\), \(\phi_f(y)=f^{-1}yf\).
The abelianized derivative vector of \(f\) is
\[
 (\partial_x f,\partial_y f)
          =R_f(X,Y)(1-Y,X-1).
\]
Thus
\[
 J_{\phi_f}=
 \begin{pmatrix}
 1&0\\
 -(Y-1)^2R_f&1+(X-1)(Y-1)R_f
 \end{pmatrix},
 \qquad \det J_{\phi_f}=B_f(X,Y).
\]
If \(w\) belongs to the commutator subgroup of a free group of
rank \(d\), and \(q_j=\partial_j w\), then
\[
 J_{\operatorname{Inn}(w)}
   =I+\bigl(1-X_i\bigr)_i(q_j)_j.
\]
The determinant lemma and the Fox fundamental identity give
\[
 \det J_{\operatorname{Inn}(w)}
 =1+\sum_j(1-X_j)q_j=1,
\]
because the abelianization of \(w\) is one. An inner adjustment
by a commutator proword therefore does not change the determinant.

The peripheral action used in the block substitutions is based at
infinity. For \(z=(xy)^{-1}\), put \(a=f(x,z)\),
\(b=f(y,z)\), and \(c=f(x,y)\). The two- and three-cycle
relations give \(b=ac^{-1}\). Consequently this action is
\(\operatorname{Inn}(a)\phi_c\), and its determinant is
\(B_f(X,Y)\).

The action \(F^{1,2}\) adds a fixed third generator, so its
determinant is \(B_f(X,Y)\). Similarly \(F^{2,3}\) has
determinant \(B_f(Y,Z)\). For \(F^{12,3}\), use the free
Nielsen basis \(x,a=xy,z\). The images of \(a,z\) depend only
on \(a,z\). The image of \(x\) is its conjugate by a word in
\(a,z\) with trivial abelianization. Its first Jacobian diagonal
entry is therefore one, and the remaining lower block is the
two-generator Jacobian in variables \(XY,Z\). Its determinant
is \(B_f(XY,Z)\). The change-of-basis factors cancel because
the automorphism is IA, so changing back does not alter this
determinant. The same argument gives \(B_f(X,YZ)\) for
\(F^{1,23}\).

The two pure braids \(x_{12},x_{23}\) act trivially on the
abelianization of the rank-three fiber. The determinant of their
Fox representation is consequently a continuous homomorphism from
the free base group \(\widehat F(x_{12},x_{23})\) to a
commutative unit group. It kills the commutator proword \(f\),
so the pure-braid side has determinant one. This conclusion also
holds if the inner action is inverted.

The determinant of the fourfold block product is
\[
 \frac{B_f(X,Y)B_f(XY,Z)}{B_f(X,YZ)B_f(Y,Z)}.
\]
Equating it to the determinant-one pure-braid side proves
\eqref{eq:twist-determinant-beta}. All calculations commute with
finite coefficient reduction and inverse limits, so they apply to
the full continuous actions.
\end{proof}

The four factors are strand-substitution actions on the rank-three
fiber. This description does not give a commutator factorization in
GT. The proposition identifies only the information captured by the
scalar determinant; it makes no assertion about the remaining
matrix entries.

\section{The oriented twist and a finite matrix comparison}
\label{sec:twist-matrix-comparison}

We determine the conjugating proword directly from two pentagon
identities, then compute the complete two-dimensional homology
matrix of the twist. The matrix gives the same degree--163 scalar
test as before and explains a limitation for central commutators.

Use the conventions
\[
 \operatorname{Inn}(g)(h)=ghg^{-1},\qquad
 x_{ij}=\sigma_{j-1}\cdots\sigma_{i+1}\sigma_i^2
                \sigma_{i+1}^{-1}\cdots\sigma_{j-1}^{-1}
                    \quad(i<j).
\]
Work in the pure five-point sphere mapping class group and its
profinite completion, and abbreviate
\[
 u=x_{12},\quad v=x_{23},\quad
 x=x_{14},\quad y=x_{24},\quad z=x_{34},\quad
 w=x_{45},\quad t=x_{51}.
\]
Then \(xyzw=1\). Set
\[
 c=x'_{35}=\sigma_4^{-1}\sigma_3^2\sigma_4.
\]
The two sphere identities used below are
\begin{equation}
 c=uxy,\qquad t=vyz.                                  \label{eq:twist-sphere-complements}
\end{equation}
They identify the boundary twist of three indicated punctures with
the twist of the complementary two punctures. In these conventions
they are equation~(18) and a cyclic instance of equation~(16) in
\cite{LS2015Twists}; the braid-word definition is its equation~(14).
These are relations in the sphere mapping class group and do not
use any GT action formula.

\subsection{Pure-braid conjugation on the fiber}

\begin{lemma}\label{lem:twist-direct-pure-braid}
On the fiber \(\widehat F_3=\langle x,y,z\rangle\),
\begin{align*}
 \operatorname{Inn}(u)(x)&=(xy)^{-1}x(xy)=y^{-1}xy,\\
 \operatorname{Inn}(u)(y)&=(xy)^{-1}y(xy)=y^{-1}x^{-1}yxy,\\
 \operatorname{Inn}(u)(z)&=z,
\end{align*}
and
\begin{align*}
 \operatorname{Inn}(v)(x)&=x,\\
 \operatorname{Inn}(v)(y)&=(yz)^{-1}y(yz)=z^{-1}yz,\\
 \operatorname{Inn}(v)(z)&=(yz)^{-1}z(yz)=z^{-1}y^{-1}zyz.
\end{align*}
In particular both base generators fix \(xyz\), and commute
with \(w=(xyz)^{-1}\).
\end{lemma}

\begin{proof}
In the three-strand braid group put \(a=\sigma_1\),
\(b=\sigma_2\). The pure elements are
\(a^2,ba^2b^{-1},b^2\). Their product is
\[
 a^2(ba^2b^{-1})b^2=a^2ba^2b
      =a(aba)ab=a(bab)ab=(ab)^3,
\]
where the middle equality is the braid relation \(aba=bab\).
This is central: if \(\Delta=aba=bab\), then
\(\Delta a=b\Delta\) and \(\Delta b=a\Delta\), and
\((ab)^3=\Delta^2\). Hence conjugation by \(a^2\) on the
other two pure generators is conjugation by the inverse of their
product.

For the strands \(1,2,4\), conjugate this three-strand calculation
by \(\sigma_3\). The base \(\sigma_1^2=u\) is fixed,
because \(\sigma_1\) and \(\sigma_3\) commute. The other two
pure elements become \(x\) and \(y\). Thus \(uxy\) commutes
with \(u,x,y\), proving the first two formulas. The last follows
from \([\sigma_1,\sigma_3]=1\).

The same three-strand argument with \(a=\sigma_2,b=\sigma_3\)
shows that \(vyz\) is central in \(\langle v,y,z\rangle\),
proving the last two formulas for \(v\). To verify \([v,x]=1\)
directly, put \(b=\sigma_2\), \(d=\sigma_3\), and
\(q=db\). The braid relation gives
\[
 d^{-1}bd=bdb^{-1},\qquad q^{-1}bq=d.
\]
Since \(x=q\sigma_1^2q^{-1}\) and \(d\) commutes with
\(\sigma_1\), it follows that \(b\), and hence \(v=b^2\),
commutes with \(x\). These identities hold already in the
discrete braid groups, so they persist in the indicated quotient
and completion. The product assertion follows immediately.
\end{proof}

\subsection{Two pentagons determine the orientation}

Let \((1,f)\) be a full GT element. Use the peripheral action
based at infinity and the four block-substitution actions
\(F^{1,2},F^{2,3},F^{12,3},F^{1,23}\) of equations
(9)--(12) of \cite{LS2015Twists}, with the disclosed correction of the first
argument in \(F^{2,3}(z)\) from \(y\) to \(z\).
Explicitly, with \(a=f(x,(xy)^{-1})\),
\(b=f(y,(xy)^{-1})\), the first action sends
\((x,y,z)\) to \((axa^{-1},byb^{-1},z)\). The second
is the same two-generator construction on \((y,z)\), fixing
\(x\). The third conjugates \(x,y\) by \(f(xy,w)\)
and \(z\) by \(f(z,w)\); the fourth conjugates \(x\)
by \(f(x,w)\) and \(y,z\) by \(f(yz,w)\).
They preserve \(xyz\); moreover \(F^{1,2}\) fixes \(xy,z\),
and \(F^{2,3}\) fixes \(x,yz\). These are the elementary
product identities of Lemma~3.1 of \cite{LS2015Twists}, derived from the two- and
three-cycle relations, rather than from the twist theorem.

Set
\[
 L=F^{1,2}F^{12,3},\qquad R=F^{2,3}F^{1,23},
 \qquad d=f(v,u)=f(u,v)^{-1}.
\]
Products of automorphisms mean composition, with the rightmost
factor applied first.

\begin{proposition}\label{prop:twist-oriented-pentagon}
With these conventions,
\begin{equation}
 L=\operatorname{Inn}(f(v,u))R,
 \qquad
 \operatorname{Inn}(f(x_{23},x_{12}))
 =F^{1,2}F^{12,3}(F^{1,23})^{-1}(F^{2,3})^{-1}.
                                                               \label{eq:twist-oriented}
\end{equation}
The equalities are automorphisms of the rank-three fiber.
\end{proposition}

\begin{proof}
We will use the elementary central-factor rule: if \(a,b,e\)
are group elements with \(e\) commuting with \(a,b\), then
\[
 f(ae,b)=f(a,b)=f(a,be).
\]
For discrete commutator words this follows by collecting the powers
of the central factor; their exponent sum is zero. Continuity gives
the rule for a proword in the closed commutator subgroup.

The block definitions and their product preservation give
\[
 L(x)=\operatorname{Inn}\bigl(f(xy,w)f(x,(xy)^{-1})\bigr)(x),
 \qquad L(z)=\operatorname{Inn}(f(z,w))(z),
\]
and
\[
 R(x)=\operatorname{Inn}(f(x,w))(x),\qquad
 R(z)=\operatorname{Inn}\bigl(f(yz,w)f(z,(yz)^{-1})\bigr)(z).
\]
Because \(u\) commutes with \(xy,w\), and \(c=uxy\)
commutes with \(u,x\), the first conjugator for \(L(x)\) is
\[
 a=f(c,w)f(x,u).
\]
Likewise \(v\) commutes with \(yz,w\), while \(t=vyz\)
commutes with \(v,z\), so the conjugator for \(R(z)\) is
\[
 a_z=f(t,w)f(z,v).
\]

The defining pentagon is
\begin{equation}
 f(u,v)f(z,w)f(t,u)f(v,z)f(w,t)=1.                    \label{eq:twist-pentagon}
\end{equation}
Conjugate this identity by \(\sigma_4^{-1}\). The braid-word
definitions show that \(u,v,w\) are fixed, \(z\) goes to
\(c=\sigma_4^{-1}\sigma_3^2\sigma_4\), and \(t\) goes
to \(x\), since \(t=x_{15}=\sigma_4x_{14}\sigma_4^{-1}\).
Thus
\begin{equation}
 f(u,v)f(c,w)f(x,u)f(v,c)f(w,x)=1.                    \label{eq:twist-conjugated-pentagon}
\end{equation}
Put \(b=f(x,w)\). By duality, \eqref{eq:twist-conjugated-pentagon} says
\[
 a=d\,b\,f(c,v).
\]
Both \(c\) and \(v\) commute with \(x\), so the final
factor commutes with \(x\). Hence
\[
 L(x)=\operatorname{Inn}(d)R(x).
\]

Next put \(e=f(z,w)\) and \(q=f(t,u)\). Since
\(a_z^{-1}=f(v,z)f(w,t)\), \eqref{eq:twist-pentagon} gives
\[
 e\,q=d\,a_z.
\]
Both \(t\) and \(u\) commute with \(z\), so \(q\)
commutes with \(z\). It follows that
\[
 L(z)=\operatorname{Inn}(d)R(z).
\]
Finally \(L\), \(R\), and \(\operatorname{Inn}(d)\)
all fix \(xyz\). Equality on \(x,z,xyz\) therefore gives
equality on \(y\) as well. This proves \eqref{eq:twist-oriented}.
\end{proof}

\paragraph{Comparison of conventions.}
With \(\operatorname{Inn}(g)(h)=ghg^{-1}\), the proposition
has \(f(x_{23},x_{12})\) on the left. This reverses the two
arguments in the printed equation~(8) of \cite{LS2015Twists}.
The proof derives the orientation directly from the defining
pentagon and the explicit block actions. It does not use the
standard-action display~(26) or Propositions~2.2--2.3 of that
source.

\subsection{The exact finite matrix identity}

Let \(O=\mathbf Z_2[\zeta_{163}]\), choose a primitive
\(163\)-rd root \(t\) in \(O/4O\), and put \(q=1+t\).
For units \(u\), define
\[
 C_t(u)=\begin{pmatrix}u&-(u-1)/q\\0&1\end{pmatrix},\qquad
 D_t(u)=\begin{pmatrix}1&0\\-t(u-1)/q&u\end{pmatrix}.
\]
For the local system sending each of \(x,y,z\) to \(t\),
take the homology basis
\[
 v_1=(1,-1,0),\qquad v_2=(0,1,-1).
\]
The direct conjugation lemma gives
\begin{equation}
 A_t=C_t(t^{-2}),\qquad B_t=D_t(t^{-2})                \label{eq:twist-base-matrices}
\end{equation}
for \(\operatorname{Inn}(u),\operatorname{Inn}(v)\).
Indeed the standard Artin half-twists have matrices
\[
 S_1=\begin{pmatrix}-t&1\\0&1\end{pmatrix},\qquad
 S_2=\begin{pmatrix}1&0\\t&-t\end{pmatrix},
\]
and the lemma identifies the pure-braid actions with their inverse
squares. The equalities \(S_1^2=C_t(t^2)\),
\(S_2^2=D_t(t^2)\), and the multiplicativity of \(C_t,D_t\)
give \eqref{eq:twist-base-matrices}.

Write
\[
 a=B_f(t,t),\qquad b=B_f(t^2,t)=B_f(t,t^2),\qquad k=a/b.
\]
The four block matrices, in the same column convention, are
\[
 M_{1,2}=C_t(a),\quad M_{2,3}=D_t(a),\quad
 M_{12,3}=bC_t(b^{-1}),\quad
 M_{1,23}=bD_t(b^{-1}).
\]
For completeness, a two-generator peripheral action fixes the product
chain \((1,t)\) and acts by its beta scalar on the cycle
\((1,-1)\). These form a basis since \(1+t\) is a unit.
This yields the first two matrices. The third action fixes \(v_1\)
and acts by \(b\) on \(v_1+(1+t)v_2\), arising from the
subgroup generated by \(xy,z\). The fourth fixes \(v_2\)
and acts by \(b\) on \((1+t)v_1+t v_2\), arising from
\(x,yz\). These properties give the last two matrices.

Proposition~\ref{prop:twist-oriented-pentagon} therefore gives the fully normalized identity
\begin{equation}
 f(B_t,A_t)=C_t(k)D_t(k)^{-1},\qquad
 f(A_t,B_t)=D_t(k)C_t(k)^{-1}.                        \label{eq:twist-finite-word}
\end{equation}
The second equality is duality. Every proword evaluation here is
in one fixed finite matrix group.

The first matrix on the right is
\[
 \begin{pmatrix}
 k-\dfrac{t(k-1)^2}{q^2k}&-\dfrac{k-1}{qk}\\[4pt]
 \dfrac{t(k-1)}{qk}&k^{-1}
 \end{pmatrix}.
\]
If \(k=1+2u\pmod{4O}\), it equals
\begin{equation}
 I+2u\begin{pmatrix}1&-q^{-1}\\tq^{-1}&-1\end{pmatrix}.
                                                               \label{eq:twist-first-layer}
\end{equation}
Its inverse has the same value modulo four. In particular its
upper-left entry is \(1+2u\), so this coefficient is not lost.

For \(r\in R_{\rm global}\), let \(n=2^{81}+1\) and use
the proword of the GT group power \(r^n\), not the ordinary word
power of \(f_r\). The beta ratio of \(r^n\) is \(k_r^n\)
and is principal. The gamma-ratio identity of Proposition~\ref{prop:163-mixed-beta-residual} gives
\[
 \Psi(k_r^n)=\Psi(h_{163}(r)).
\]
Consequently \eqref{eq:twist-finite-word}--\eqref{eq:twist-first-layer}
are an exact finite matrix representation of the same unresolved
degree--163 scalar coordinate.

\begin{corollary}\label{cor:twist-finite-scalar-and-commutator}
For an acyclotomic full GT element, the finite matrices satisfy
\eqref{eq:twist-finite-word}. For \(r\in R_{\rm global}\),
the first layer after the GT group power \(2^{81}+1\) determines
\(\Psi(h_{163}(r))\). For any two multiplier-one elements
\(r_1,r_2\), one has
\[
 f_{[r_1,r_2]}(A_t,B_t)=1.
\]
Thus this specialization cannot detect the commutator value of the
relative quaternion obstruction.
\end{corollary}
\begin{proof}
The matrix calculation and gamma comparison just proved give the
first two assertions. On the multiplier-one subgroup, each beta evaluation, and hence the
ratio \(k_r\), is multiplicative in the GT group law. Thus
\(k_{[r_1,r_2]}=1\), and \eqref{eq:twist-finite-word} gives
\[
 f_{[r_1,r_2]}(A_t,B_t)=1
\]
for multiplier-one elements \(r_1,r_2\). This proves the last assertion.
\end{proof}

The matrix identity remains an equivalent test for the full scalar
image if an additional full-GT identity controls these proword
values. Such an identity is not proved here.

\section{The scope of a common-cover reciprocity argument}
\label{sec:common-cover-reciprocity-limits}

The remaining harmonic comparison is
\begin{equation}
 \overline{B_{g,p}(i,-1)}^{(p-\chi_{-4}(p))/2}
      =(-1)^{\widetilde\psi_p(g)}
      =(-1)^{\rho_p^I(g)/2},\qquad g\in N.               \label{eq:common-cover-reciprocity-question}
\end{equation}
Here the bar denotes reduction in the residue norm-one torus.
The companion equality proves the second equality on all of \(M\),
and hence on \(N\).
The first equality compares two different coefficient primes.
The following proposition specifies what an integral common cover and
the componentwise beta identities do not, by themselves, prove.
It does not construct a counterexample in GT.

\begin{proposition}\label{prop:common-cover-homology-countermodel}
For every odd prime \(p\), the degree-four cyclic cover with
monodromies \((1,2)\) and the degree-\(p\) cyclic cover with
monodromies \((1,1)\) have a common degree-\(4p\) cover.
Its compact genus is \(2p-1\), and it has four punctures.

There is an automorphism of its full profinite integral first
homology which commutes with the deck group, preserves the
intersection form, fixes every peripheral homology class,
and is compatible with pushforward and transfer to both smaller
covers. It induces the identity on the full integral homology
of the degree-\(p\) cover, but has nontrivial harmonic quadratic
sign at coefficient \(p\) on the degree-four cover. In particular
it is the identity at coefficient two.
\end{proposition}

\begin{proof}
The common deck group is \(Q=C_4\times C_p\), with the
monodromy images
\[
                  x=(1,1),\qquad y=(2,1),\qquad z=(1,-2).
\]
Their orders are \(4p,2p,4p\), respectively. Thus the three
original punctures have \(1,2,1\) preimages. The Euler
characteristic of the open cover is \(-4p\), so
\(2-2g-4=-4p\) gives \(g=2p-1\).

Let \(a=(1,0)\in Q\), of order four, and let \(V_p\)
be the integral first homology of the open common cover over
\(\mathbf Z_p\). Since \(p\) is odd, the following decomposition
is integral and etale:
\begin{equation}
 \mathbf Z_p[C_4]
      \simeq\mathbf Z_p\times\mathbf Z_p\times\mathbf Z_p[i].
                                                        \label{eq:common-cover-c4-algebra}
\end{equation}
The factors correspond to the trivial, order-two, and primitive
order-four characters. The involution induced by \(a\mapsto a^{-1}\)
is trivial on the first two factors and complex conjugation on
the last one. This notation includes the split algebra when
\(p\equiv1\pmod4\).

The residue norm-one group in the last factor is cyclic of order
\(p-\chi_{-4}(p)\). Choose an element with nontrivial quadratic
sign, and lift it to a norm-one unit \(v\in\mathbf Z_p[i]^\times\).
Such a lift can be taken among the prime-to-\(p\) roots of unity:
in the split case lift the two reciprocal coordinates, and in
the inert case lift a root of unity of order \(p+1\).
Let \(u\in\mathbf Z_p[C_4]^\times\) be the element with
components \((1,1,v)\) in \eqref{eq:common-cover-c4-algebra}. Then
\[
                              u^\iota u=1.
\]
Let \(u(a)\) act on \(V_p\). It is an automorphism and
commutes with the entire abelian deck group. Deck transformations
preserve the intersection form, so the adjoint of \(a\) is
\(a^{-1}\). Hence \(u^\iota u=1\) proves that \(u(a)\)
preserves this form, including its induced nondegenerate form
on compact homology.

The boundary representation is the sum of the permutation
representations of the four punctures, modulo their one sum
relation. The two punctures over \(y\) have stabilizer
\(\langle(2,1)\rangle\), so their quotient deck action is
\(C_2\); the other punctures are fixed. Consequently every
boundary class lies in the first two factors of \eqref{eq:common-cover-c4-algebra}.
The element \(u\) fixes all of them, integrally.

Under the map to the degree-four cover, \(a\) becomes its
deck generator. Thus \(u(a)\) induces multiplication by
\(v\) on its compact primitive order-four summand, giving
the prescribed nontrivial harmonic sign modulo \(p\).
Under the map to the degree-\(p\) cover, \(a\) becomes one;
the induced operator is the augmentation of \(u\), which
is one. Pushforward and transfer commute with these deck
actions. A polynomial in \(a\), with \(\mathbf Z_p\)
coefficients, therefore commutes with both maps in the
corresponding equivariant sense.

Finally choose this operator at coefficient \(p\) and the
identity at every other coefficient prime. The product is
an automorphism of full profinite integral homology. It
has all the stated properties. No denominators divisible
by \(p\) occur in \eqref{eq:common-cover-c4-algebra}, so the construction does not discard
the integral structure at the coefficient prime where the
\(C_p\)-deck action is ramified.
\end{proof}

This proposition concerns homology operators. It does not lift
the operator to an automorphism of the ambient free profinite
group, to a positive peripheral automorphism, or to a solution
of the GT pentagon. It shows that an argument using only
integral homology, intersection pairings, the indicated deck
action, and the common-cover transfer maps leaves both signs
possible in \eqref{eq:common-cover-reciprocity-question}.


The componentwise formal construction in
Proposition~\ref{prop:harmonic-formal-residue-gap} gives a separate
limitation: beta associativity, unitarity, integral dyadic duplication
and the logarithmic Anderson identities do not force the first equality
in \eqref{eq:common-cover-reciprocity-question}. The companion theorem
now identifies its other side with the divided Kummer character.
Neither construction gives a full-GT element with nonzero discrepancy.
A decision requires further nonabelian compatibility or a GT witness.

\section{Odd arithmetic targets invisible to the primewise and CM realizations}

Write \(G=\widehat{\mathrm{GT}}\), \(H=G_{\mathbf Q}\),
\(M=\ker\lambda\), \(N=\ker(\lambda,\rho_2)\), and
\(K=N\cap H\).  Let \(D\subset N\) be the common kernel of
all rational-prime Kummer cocycles.  For every prime \(r\), let
\(\Phi_r\) be the based standard action on the free pro-\(r\)
fundamental group, with its rational tangential base point.
Let \(T,W\) be the signed cubic and harmonic frames on \(M\).
Define the closed subgroup
\begin{equation}\label{eq:odd-mixed-common-kernel}
 J=D\cap\bigcap_r\ker\Phi_r\cap\ker T\cap\ker W,
 \qquad J_H=J\cap H.
\end{equation}
All intersections and kernels in this formula are taken inside \(N\).

Fix an odd prime \(p\).  Retain the arithmetic target basis
\(\kappa_{ij}\in Y_p=H^1(K,\mathbf F_p)^H\), indexed by
\(i<j\) in \(\{0\}\cup\mathcal S_p\), obtained from
\[
 \zeta_{ij}=\gamma_i\smile\gamma_j-D_{ij}
                      \in H^2(U,\mathbf F_p),
 \qquad U=\widehat{\mathbf Z}^{\times}.
\]
Here \(D_{ij}\) is a finite sum of the cyclic-extension classes.
We use the pushout convention for transgression: an invariant
character is sent to the class of its central pushout extension.
Changing the convention changes every \(\kappa_{ij}\) by the
same sign and has no effect on the assertions of independence.

\begin{theorem}\label{thm:odd-target-survives-all-standard-kernels}
For each pair \(i<j\), there is an arithmetic element
\(w_{ij}\in J_H\) such that, for every \(a<b\),
\begin{equation}\label{eq:odd-mixed-evaluation}
                  \kappa_{ab}(w_{ij})=2\,\delta_{ai}\delta_{bj}.
\end{equation}
Consequently restriction
\[
                         Y_p\longrightarrow H^1(J_H,\mathbf F_p)
\]
is injective.  In particular, if \(\chi\in H^1(N,\mathbf F_p)^G\)
factors through the joint rational Kummer, primewise word, and
two-frame realization of \(N\), then \(\chi|_K=0\).

Thus no nonzero odd arithmetic target character can be obtained
solely from those realizations.  This conclusion does not assert
that a target character fails to extend to full GT.
\end{theorem}

\subsection{Central extensions and alternating evaluation}

We first explain why an odd target character can be evaluated on
the entire cyclotomic kernel \(H^{\mathrm{cyc}}=H\cap M\).
Arithmetic inflation of \(\zeta_{ij}\) is zero.  Represent this
class by a continuous central extension
\[
             1\longrightarrow\mathbf F_p\longrightarrow E_{ij}
                         \longrightarrow U\longrightarrow1.
\]
Vanishing of its arithmetic inflation is equivalent to the
existence of a continuous homomorphism
\(\widetilde\lambda:H\to E_{ij}\) lifting \(\lambda\).
Its restriction to \(H^{\mathrm{cyc}}\) is an \(H\)-invariant
character \(\widetilde\kappa_{ij}\).
Two lifts differ by a character of \(H\), and every such
character factors through \(H^{\rm ab}=U\) by
Kronecker--Weber.  Therefore \(\widetilde\kappa_{ij}\) is
independent of the chosen lift.  Its restriction to \(K\) is
\(\kappa_{ij}\), by the naturality and injectivity of arithmetic
transgression and the isomorphism
\(H^2(U,\mathbf F_p)\simeq H^2(A,\mathbf F_p)\).

If \(c\) represents a central-extension class on an abelian
group, the commutator of lifts of \(u,v\) has central coordinate
\(c(u,v)-c(v,u)\).  The cocycle
\(\gamma_a(u)\gamma_b(v)\) represents
\(\gamma_a\smile\gamma_b\).  Every cyclic-extension class
has zero alternating form, because it is pulled back from an
extension of abelian groups.  Hence, for any \(h,h'\in H\),
\begin{equation}\label{eq:odd-mixed-alternation}
 \widetilde\kappa_{ab}([h,h'])
   =\gamma_a(\lambda(h))\gamma_b(\lambda(h'))
      -\gamma_a(\lambda(h'))\gamma_b(\lambda(h)).
\end{equation}
We use \([h,h']=hh'h^{-1}h'^{-1}\).  The correction
\(D_{ab}\), although essential for arithmetic inflation to
vanish, does not affect this alternating formula.

\subsection{Inertia choices}

Fix decomposition groups and inertia groups \(I_r\subset H\)
for rational primes \(r\).  Inertia at \(r\) acts trivially on
roots of unity of order prime to \(r\), and its \(r\)-adic
cyclotomic image is all of \(\mathbf Z_r^\times\).
The latter assertion follows, for example, from total
ramification of the local cyclotomic tower; its reciprocity
description is \cite[I, Theorem~1.1, p.~20]{MilneCFT}.

Choose \(h_0\in I_p\) whose multiplier has \(p\)-component
in \(1+p\mathbf Z_p\) and satisfies \(\gamma_0(h_0)=1\).
For each \(\ell\in\mathcal S_p\), choose
\(h_\ell\in I_\ell\) with \(\gamma_\ell(h_\ell)=1\).
Every other character among the \(\gamma_i\) is zero on
the chosen element.  In concise notation,
\begin{equation}\label{eq:odd-mixed-inertia-coordinates}
                         \gamma_a(\lambda(h_i))=\delta_{ai}.
\end{equation}
The inertia primes attached to distinct indices are distinct:
index zero uses \(p\), while \(p\notin\mathcal S_p\).

All these elements fix \(\mathbf Q(i,\sqrt{-3})\).
Indeed, every \(\ell\in\mathcal S_p\) is different from
two and three.  For \(p\ne3\), inertia at \(p\) also fixes
this field, which is unramified outside two and three.
For \(p=3\), the imposed condition
\(\lambda_3(h_0)\in1+3\mathbf Z_3\) makes \(h_0\) fix
\(\mu_3\), and inertia at three fixes \(\mathbf Q(i)\).
Thus the assertion holds in this case too.

Put \(z_{ij}=[h_i,h_j]\).  The full multiplier of \(z_{ij}\)
is one.  The tangential pro-\(r\) action of \(H\) on the
thrice-punctured line is unramified outside \(r\):
\begin{equation}\label{eq:odd-mixed-unramified-word}
                              \Phi_r(I_s)=1\qquad(s\ne r).
\end{equation}
The primary statement in \cite[\S8.5, pp.~151--152]{HainMatsumoto2003}
is explicitly an action into \(\operatorname{Aut}\) of the
pro-\(r\) group with tangent vector \(\overrightarrow{01}\);
the subsequent passage to the unipotent completion is separate.
Thus \eqref{eq:odd-mixed-unramified-word} holds for the based
action, not merely its outer or unipotent image.
At least one member of the pair \(h_i,h_j\) has trivial
\(\Phi_r\)-image, for every \(r\).  Therefore
\begin{equation}\label{eq:odd-mixed-z-word}
                              \Phi_r(z_{ij})=1\quad\hbox{for all }r.
\end{equation}

Let \(E_4:y^2=x^3-x\) and \(E_3:y^2=x^3+1\).
Over their CM fields their adelic Tate representations have
abelian images: the Tate modules are rank one over the
completed CM orders, whose unit groups are abelian.
The choices \(h_i,h_j\) fix both CM fields.  Consequently
\begin{equation}\label{eq:odd-mixed-z-CM}
                 \rho_{E_4}(z_{ij})=\rho_{E_3}(z_{ij})=1.
\end{equation}
No good-reduction or local surjectivity assertion at the
coefficient prime is required for this step.

\subsection{Kummer cancellation and detection}

Let \(c\in H\) be complex conjugation, with full multiplier
\(\lambda(c)=-1\).  Define
\begin{equation}\label{eq:odd-mixed-witness}
                         w_{ij}=z_{ij}\,c z_{ij}c^{-1}.
\end{equation}
Each rational Kummer cocycle is additive on \(M\) and obeys
\[
                  \rho_q(aga^{-1})=\lambda(a)\rho_q(g)
                                  \qquad(g\in M).
\]
Since \(z_{ij}\in M\), the two summands in
\eqref{eq:odd-mixed-witness} have opposite Kummer values.
Thus \(\rho_q(w_{ij})=0\) for every rational prime \(q\),
and \(w_{ij}\in D\cap H\).  Normality of the word kernels
and \eqref{eq:odd-mixed-z-word} give
\(\Phi_r(w_{ij})=1\) for every \(r\).

The two Tate representations are homomorphisms, so
\eqref{eq:odd-mixed-z-CM} also holds for \(w_{ij}\).
The established signed arithmetic comparison on
\(H^{\mathrm{cyc}}\) is
\[
 W_g=\rho_{E_4}(g),\qquad
 T_g=\kappa_{3^{1/4}}(g)\rho_{E_3}(g).
\]
The quartic Kummer factor is one on \(D\cap H\).
It follows that \(T(w_{ij})=W(w_{ij})=1\), proving
\(w_{ij}\in J_H\).

Finally, \(\widetilde\kappa_{ab}\) is an \(H\)-invariant
character of \(H^{\mathrm{cyc}}\).  Equations
\eqref{eq:odd-mixed-alternation} and
\eqref{eq:odd-mixed-inertia-coordinates} give
\[
 \kappa_{ab}(w_{ij})
  =2\widetilde\kappa_{ab}(z_{ij})
  =2(\delta_{ai}\delta_{bj}-\delta_{aj}\delta_{bi}).
\]
For ordered pairs \(a<b\) and \(i<j\), the last negative
term is zero.  This proves \eqref{eq:odd-mixed-evaluation}.
Because \(p\) is odd, its diagonal entry is nonzero.
Every element of \(Y_p\) is a finite linear combination
of these basis elements, so their evaluations on the
\(w_{ij}\) prove injectivity of restriction to \(J_H\).

The joint realization of \(N\) in the theorem has kernel
exactly \(J\).  Any character factoring through it is zero
on \(J\), and hence its arithmetic restriction is zero
on \(J_H\).  The proved injectivity makes that arithmetic
restriction zero on all of \(K\), as claimed.

For \(p=3\), this applies in particular to both
\(\kappa_{0,7}\) and \(\kappa_{7,13}\).  Each has an
arithmetic witness with trivial rational Kummer data,
trivial action on every free pro-prime completion, and
trivial full signed cubic and harmonic frames.  A full GT
lift of either class, if one exists, must detect this
common kernel and therefore requires additional information.

\section{The odd witnesses also kill every abelian-cover homology action}

Let \(F=\widehat F(x,y)\), and let \(M\) be the multiplier-one
subgroup of full profinite GT.  Use the normalized full adelic
Alexander unit \(B_g(X,Y)\), with its semilinear cocycle law
\[
 B_{gh}(X,Y)=B_g(X,Y)B_h(X^{\lambda_g},Y^{\lambda_g}).
\]
This is the full completed group-ring coefficient, not only its
augmentation expansion at one coefficient prime.  The convention
is \cite[\S2.3, p.~101]{Furusho2003}.  On multiplier one it is
the based coefficient on the Alexander module.

\begin{lemma}\label{lem:full-integral-alexander-unitarity}
For every \(g\in M\),
\begin{equation}\label{eq:full-alexander-unitarity}
                    B_g(X,Y)B_g(X^{-1},Y^{-1})=1
\end{equation}
in the full completed adelic group algebra.  Consequently,
for complex conjugation \(c\) and any \(z\in M\),
\begin{equation}\label{eq:full-alexander-symmetrization}
                           B_{z c z c^{-1}}=1.
\end{equation}
\end{lemma}

\begin{proof}
Fix a positive integer \(d\), a coefficient prime \(r\), and a
character \((a,b)\) of \((\mathbf Z/d)^2\) with values in a
finite extension \(L/\mathbf Q_r\).  The cover corresponding to
\[
                         F_d=\ker(F\to(\mathbf Z/d)^2)
\]
is the punctured Fermat cover.  The automorphism of \(F\)
associated with \(g\) preserves \(F_d\), acts trivially on its
deck group, and preserves every positive peripheral class.
For every nontrivial deck character its eigenspace in
\(H_1(F_d,L)\) has dimension one, and its scalar action is
\(B_g(a,b)\).  This follows either by specializing the
Alexander module or by the cellular complex with differential
\((a-1,b-1):L^2\to L\).

We specify the boundary characters and their scalars.  In
the usual normalization \(x\) is fixed and \(y\) is conjugated
by an element of \(F'\).  Write \(z_0=(xy)^{-1}\) and
\(\phi_g(z_0)=k^{-1}z_0k\).  If \([k]=(u,v)\) in
\(F^{\rm ab}\), then \(u=v\): modulo \(\gamma_3F\),
the automorphism is the identity, so the commutator of
\([k]\) with \((-1,-1)\) is zero.  The degree-two free
commutator coordinate gives \(u-v=0\) in
\(\widehat{\mathbf Z}\).

The boundary subspace is the sum of the permutation modules
induced from the three cyclic peripheral subgroups, with one
trivial relation.  Its nontrivial characters satisfy
\(a=1\), \(b=1\), or \(ab=1\).
The conjugator at \(x\) is trivial, the one at \(y\) has
trivial deck image, and the one at \(z_0\) has diagonal deck
image.  Their scalars on the respective boundary characters
are therefore one; in the last case it is
\(a^u b^v=(ab)^u=1\).
Hence \(B_g(a,b)=1\) at every boundary character.
At \((a,b)=(1,1)\), the scalar is also one because the
class-two commutator is fixed.

The homology modulo boundary classes has the nondegenerate
alternating form of the compactified cover.  The induced
automorphism preserves this form with multiplier one.
Algebraically, in a surface presentation, preservation of
each positive boundary word preserves the degree-two
fundamental relation \(\sum_j[A_j,B_j]\) after quotienting
by the boundary classes.  Conjugating a boundary word adds
only wedges involving its boundary class, which are zero in
that quotient.  This gives preservation of the associated
alternating form without asserting that the profinite
automorphism is realized by a discrete homeomorphism.

The deck action preserves the same form.  Thus the
\((a,b)\) and \((a^{-1},b^{-1})\) eigenspaces pair
nondegenerately whenever \(a,b,ab\ne1\), and
\[
                         B_g(a,b)B_g(a^{-1},b^{-1})=1.
\]
The boundary calculation proves this identity in all
remaining cases as well.  It is valid also when \(r\mid d\),
since the character decomposition is being taken over the
characteristic-zero field \(L\).

For a fixed finite deck group, its characters over a splitting
field separate the characteristic-zero group algebra.
The integral algebra \(\mathbf Z_r[(\mathbf Z/d)^2]\)
embeds in that algebra.  The identity therefore holds in
this integral algebra for every \(r,d\).  Passing to the
inverse limit proves \eqref{eq:full-alexander-unitarity}
in the full adelic completed algebra.

Because this algebra is commutative, the cocycle law gives
\[
 B_{a z a^{-1}}(X,Y)=B_z(X^{\lambda_a},Y^{\lambda_a})
                                            \qquad(z\in M).
\]
Take \(a=c\), whose full multiplier is \(-1\).
Equation \eqref{eq:full-alexander-unitarity}, followed by
the multiplier-one product law, proves
\eqref{eq:full-alexander-symmetrization}.
\end{proof}

\begin{corollary}\label{cor:odd-target-invisible-abelian-covers}
Let \(J\) be the common radical, primewise-word, and signed-CM
kernel from the preceding theorem, and put
\[
                              J_B=J\cap\ker B.
\]
For every odd prime \(p\), restriction
\[
                  Y_p\longrightarrow H^1(J_B\cap H,\mathbf F_p)
\]
is injective.  Every witness \(w_{ij}\) from that theorem
lies in \(J_B\cap H\), with the same diagonal evaluation
\(\kappa_{ab}(w_{ij})=2\delta_{ai}\delta_{bj}\).
Each witness acts trivially on the full adelic first
homology of every finite abelian cover of the punctured line.
In particular, Fermat-cover homology alone cannot provide
a full GT lift of a nonzero odd target character.
\end{corollary}

\begin{proof}
The witnesses have the form \(z c z c^{-1}\) with \(z\in M\),
so the lemma adds \(B_w=1\) to their previously proved
properties.  Their evaluations prove the asserted
injectivity exactly as before.

For the square abelian cover \(F_d\), every nontrivial
character line has scalar one.  The deck-trivial subspace
over \(\mathbf Q_r\) is the transfer of \(H_1(F,\mathbf Q_r)\),
on which multiplier one acts as the identity.  Thus the
whole \(H_1(F_d,\mathbf Q_r)\) is fixed.  Its integral
\(\mathbf Z_r\)-lattice is torsion-free, so it too is fixed.
Any finite abelian cover is dominated by one of these
square covers.  Transfer is injective over \(\mathbf Q_r\)
and is equivariant, so its homology is likewise fixed,
first rationally and then integrally.  This proves the
assertion for every coefficient prime.

Alternatively, on the image of \(B:M\to\Lambda^\times\),
complex conjugation acts by inversion.  An invariant
\(\mathbf F_p\)-character of this image is consequently
its own negative and is zero for odd \(p\).
The witness statement also excludes obtaining a nonzero
arithmetic target by jointly using these homology actions
with the earlier Kummer, word, and CM realizations.
It imposes no restriction on a genuinely nonmetabelian
mixed finite quotient.
\end{proof}

\subsection{An explicit inertia witness at seven}
The same strict reduction also has a local inertia construction.
\begin{proposition}\label{prop:seven-inertia-mixed-prime-witness}
There is an arithmetic element \(w\in T_D\) for which every
single-prime word action is trivial and
\[
 W_7(w)=\left(\frac{1+7i}{1-7i}\right)^{-2}.
\]
This value has infinite order.
\end{proposition}
\begin{proof}
Fix decomposition groups at the primes two and seven,
with inertia groups \(I_2,I_7\subset H\).  The normalized
tangential action on the free pro-\(p\) fundamental group of
\(\mathbf P^1-\{0,1,\infty\}\) is unramified outside \(p\).
This is a statement about the tangential action, not just the
outer action; see \cite[\S8.5]{HainMatsumoto2003}.  Thus
\begin{equation}\label{eq:disjoint-inertia-word-actions}
                  \Phi_p(I_r)=1\qquad(r\ne p).
\end{equation}

Here is the local CM input in an explicit form.  Let
\(K=\mathbf Q_7(i)\) and \(O=\mathbf Z_7[i]\).
This is the unramified quadratic extension of \(\mathbf Q_7\).
The curve \(E_4\) has good reduction at seven, and a direct
count gives \(\#E_4(\mathbf F_7)=8\): the numbers of affine
points above \(x=0,\ldots,6\) are respectively
\(1,1,0,0,2,2,1\), with the point at infinity added.
Consequently the seven-power Frobenius \(\pi\) on the
reduction satisfies \(\pi^2=[-7]\).
In the formal parameter \(t=-x/y\), the endomorphism
\([-7]\) of \(\widehat E_4/O\) therefore satisfies
\[
 [-7](t)\equiv-7t\pmod{t^2},\qquad
 [-7](t)\equiv t^{49}\pmod7.
\]
The CM automorphism \((x,y)\mapsto(-x,iy)\) acts on this
parameter by \(t\mapsto it\).  Its action and the
\(\mathbf Z_7\)-action give the formal \(O\)-module structure.
The preceding two congruences identify this formal module
with the Lubin--Tate module for the uniformizer \(-7\).
Since \([7]\) is purely inseparable on the reduction, every
point of \(E_4[7^n]\) reduces to the identity.  Thus its
Tate module is the rank-one \(O\)-module of this formal
group.  The Lubin--Tate reciprocity theorem says that the
inertia image on it is all of \(O^\times\); see
\cite[p.~380, the displayed reciprocity formula~(2)]{LubinTate1965}.
Indeed the finite torsion fields are totally ramified with
groups \((O/7^nO)^\times\), and inverse limits give the
assertion.  As \(K/\mathbf Q_7\) is unramified, its inertia
is \(I_7\).

Choose
\[
             a=\frac{1+7i}{1-7i}\in1+7O.
\]
Then \(a\bar a=1\), and \(a\) has infinite order.  For
example its seven-adic logarithm has first term \(14i\),
of valuation one, while all subsequent terms have larger
valuation.  Choose \(j\in I_7\) whose action on \(T_7(E_4)\)
is multiplication by \(a\).  The Weil pairing gives
\(\lambda_7(j)=\det(a)=a\bar a=1\).  Inertia at seven
fixes all roots of unity of order prime to seven, so in fact
\begin{equation}\label{eq:inertia-seven-full-multiplier}
                              \lambda(j)=1.
\end{equation}

The local cyclotomic extension at two gives an element
\(h\in I_2\) with \(\lambda_2(h)=-1\).
All other components of \(\lambda(h)\) are one.
Since \(h(i)=-i\), its action on the seven-adic CM module
is semilinear for the involution of \(O\).  Whatever its
unit factor, it therefore conjugates multiplication by
\(a\) to multiplication by \(\bar a\).

Use the convention \([h,j]=h j h^{-1}j^{-1}\), and put
\(z=[h,j]\).  Its full multiplier is one.  For each prime
\(p\), at least one of \(\Phi_p(h),\Phi_p(j)\) is the
identity by \eqref{eq:disjoint-inertia-word-actions}, so
\begin{equation}\label{eq:commutator-all-word-kernels}
                           \Phi_p(z)=1\quad\hbox{for every }p.
\end{equation}
For any rational-prime Kummer cocycle \(\rho_q\), the
crossed-homomorphism law and \(\lambda(j)=1\) give
\[
              \rho_q(z)=(\lambda(h)-1)\rho_q(j).
\]
This adelic vector is supported only at its two-component.
Let \(e_7\in\widehat{\mathbf Z}=\prod_p\mathbf Z_p\)
have seven-component one and all other components zero,
and define \(w=z^{e_7}\).  The exponent is taken through
the continuous map
\(\widehat{\mathbf Z}\to\overline{\langle z\rangle}\);
it is a profinite group power.  On the multiplier-one
closed cyclic group generated by \(z\), all Kummer
cocycles are additive.  Hence
\[
                  \rho_q(w)=e_7\rho_q(z)=0\quad\hbox{for all }q.
\]
Thus \(w\in D\cap H\), and
\eqref{eq:commutator-all-word-kernels} implies
\(w\in P_D\cap H\).  The cubic residue defect is zero
on arithmetic elements, so \(w\in T_D\cap H\).

Finally, in the signed CM identification,
\[
                    W_7(z)=\bar a/a=a^{-2}.
\]
This element belongs to the pro-seven group \(1+7O\),
so exponentiation by \(e_7\) leaves it unchanged:
\[
                         W_7(w)=a^{-2}.
\]
Projection from the local norm-one torus to its pro-seven
Sylow factor annihilates every pro-two element.  It
therefore factors through \(N_4(4)_{\rm odd}\).
The displayed value has infinite order, so the image of
the closed cyclic subgroup generated by \(w\) in
\(T_D/U_D\) is infinite.  This proves both strictness
assertions.
\end{proof}

\subsection{An explicit cubic inertia witness}

The following construction gives a local arithmetic witness for the
strictness of the second odd CM restriction.

\begin{proposition}\label{prop:five-inertia-odd-cubic}
There is an element \(w\in U_D\cap H\) whose full harmonic
frame is one and whose cubic frame at five is
\[
 T_5(w)=\left(\frac{1+5\omega}{1+5\omega^2}\right)^{-2}.
\]
This value has infinite pro-five order.
\end{proposition}
\begin{proof}
Use the CM curve \(E_3:y^2=x^3+1\).
It has good reduction at five.  The affine point counts
above \(x=0,\ldots,4\) are \(2,0,2,0,1\), so
\(\#E_3(\mathbf F_5)=6\).  The five-power Frobenius
therefore satisfies \(\pi^2=[-5]\) on the reduction.
Over the unramified quadratic extension
\(K=\mathbf Q_5(\omega)\), with integers
\(O=\mathbf Z_5[\omega]\), the formal CM parameter
\(t=-x/y\) has \(\omega\)-action \(t\mapsto\omega t\).
As in the preceding local computation,
\[
 [-5](t)\equiv-5t\pmod{t^2},\qquad
 [-5](t)\equiv t^{25}\pmod5.
\]
The formal group is the Lubin--Tate module for the
uniformizer \(-5\), and all five-power torsion lies in
it.  Hence inertia \(I_5\) has image \(O^\times\) on
its rank-one \(O\)-Tate module
\cite[Theorem~2(b),(d),(e), pp.~383--384]{LubinTate1965}.

Choose
\[
                       a=\frac{1+5\omega}{1+5\omega^2}.
\]
It has norm one and belongs to \(1+5O\).  Its logarithm
has leading term \(5(\omega-\omega^2)\), of valuation
one, so it has infinite order.  Choose \(j\in I_5\)
whose cubic Tate action is multiplication by \(a\).
The determinant and the Weil pairing give
\(\lambda_5(j)=1\), and prime-to-five roots of unity
are unramified.  Thus \(\lambda(j)=1\) at every prime.

Choose \(h\in I_3\) with \(\lambda_3(h)=-1\);
all its other cyclotomic components are one.  Such an
element exists in the local cyclotomic inertia image.
It sends \(\omega\) to \(\omega^2\).  Accordingly its
cubic Tate action conjugates multiplication by \(a\)
to multiplication by \(\bar a\).
Put \(z=[h,j]=h j h^{-1}j^{-1}\).
The tangential pro-\(p\) action itself is unramified
outside \(p\), as in \cite[\S8.5]{HainMatsumoto2003}.
For every \(p\), one of the two factors has trivial
pro-\(p\) action.  Therefore
\[
                    \lambda(z)=1,\qquad \Phi_p(z)=1
                                      \quad\hbox{for every }p.
\]
The rational Kummer cocycle formula gives
\[
                \rho_q(z)=(\lambda(h)-1)\rho_q(j)
                                      \quad\hbox{for every }q.
\]
Each such value is supported only at the coefficient
prime three.  Let \(e_5\in\widehat{\mathbf Z}\)
be the idempotent supported at five, and put
\(w=z^{e_5}\), using the continuous power map on
\(\overline{\langle z\rangle}\).
Additivity on multiplier one gives \(\rho_q(w)=0\)
for every rational prime \(q\).  Thus
\(w\in P_D\cap H\), and arithmetic vanishing of the
cubic residue defect gives \(w\in T_D\cap H\).

We check that the harmonic image is already exactly
trivial.  Both inertia groups \(I_3\) and \(I_5\)
fix \(\mathbf Q(i)\), which is unramified at these
primes.  The Gaussian Tate representation restricted
to \(G_{\mathbf Q(i)}\) is abelian: all CM
endomorphisms are defined there, so the image lies in
the commutative unit group
\((\mathbf Z[i]\otimes\widehat{\mathbf Z})^\times\).
Consequently \(\rho_{E_4}(z)=1\), and likewise
\(\rho_{E_4}(w)=1\).  The signed harmonic identification
on the cyclotomic arithmetic subgroup yields
\[
                               W(w)=1.
\]
In particular \(w\in U_D\cap H\).

On the cubic Tate module at five,
\(\rho_{E_3,5}(z)=\bar a/a=a^{-2}\).
This value is pro-five, so the \(e_5\) power preserves it.
The signed cubic frame on arithmetic multiplier one
is \(T=\kappa\rho_{E_3}\), where
\(\kappa\) is the fourth-root Kummer character of
three.  Since \(w\in D\), it fixes that fourth root,
and \(\kappa(w)=1\).  It follows that
\[
                              T_5(w)=a^{-2}.
\]
The logarithm calculation above proves its infinite pro-five order.
\end{proof}

\subsection{What the finite Gaussian translation quotient does not determine}

There is a concrete finite obstruction to concluding the harmonic
square test solely from the translation group, its symplectic
action and the three orbifold inertia classes. This obstruction
is at the level of a finite peripheral model; it is not a GT
counterexample.

\begin{proposition}\label{prop:gaussian-finite-peripheral-scope}
For every odd prime \(\ell\), there is a finite quotient of the
three-punctured sphere group with peripheral orders \(4,2,4\),
and an automorphism that fixes its maximal pro-\(p\) quotient
for every prime \(p\), preserves each positive peripheral
conjugacy class, acts symplectically on its translation subgroup,
and has nontrivial Gaussian square sign.
\end{proposition}
\begin{proof}
Use the finite etale algebra
\[
 A=\mathbf F_\ell[i]=\mathbf F_\ell[t]/(t^2+1),\qquad
 J=A_{\rm add}\rtimes\langle i\rangle.
\]
The algebra has dimension two, including in the split case.
Its involution sends \(i\) to \(-i\). Write
\((v,r)(w,s)=(v+rw,rs)\), and choose
\[
 x=(0,i),\qquad y=(1,-1),\qquad z=(xy)^{-1}=(1,i).
\]
These elements have orders \(4,2,4\), respectively, and \(xyz=1\).
Their commutator is \([x,y]=(i-1,1)\).
Since \(i-1\) is invertible and conjugation by \(x\) multiplies
translations by \(i\), they generate \(J\).
Thus they define the specified quotient.

The norm-one subgroup
\[
 A^1=\{a\in A^\times:a\overline a=1\}
\]
is cyclic of order \(\ell-\chi_{-4}(\ell)\), which is divisible
by four. Choose a nonsquare \(a\in A^1\), and define
\(\phi_a(v,r)=(av,r)\). This is an automorphism.
Its determinant on the two-dimensional translation space is
\(a\overline a=1\), hence it preserves its alternating form.
It fixes \(x\). If
\[
 f=((1-a)/2,1),
\]
then \(\phi_a(y)=f^{-1}yf\). The two elements \(z=(1,i)\)
and \(\phi_a(z)=(a,i)\) are conjugate by a translation, because
translation conjugacy changes the first coordinate by
\((1-i)v\), and \(1-i\) is invertible.
Thus all three positive peripheral classes are preserved.

The maximal pro-two quotient of \(J\) is \(C_4\), and
\(\phi_a\) fixes it. Its maximal pro-\(\ell\) quotient is
trivial: killing \(i\), whose order is four, also kills every
translation by the relation \(iv=v\), since \(i-1\) is invertible.
All other single-prime quotients are trivial. Nevertheless
\[
             a^{(\ell-\chi_{-4}(\ell))/2}=-1.
\]
The affine Alexander calculation identifies its beta scalar
with \(a\): since \([x,y]=i-1\), the conjugator has coefficient
\(R=(1-a)/(2(i-1))\), and
\[
             1+(i-1)(-2)R=a.
\]
\end{proof}

This construction specifies only an automorphism of the displayed
finite quotient. It does not lift that automorphism to full GT,
impose the pentagon, or prescribe the simultaneous rational
Kummer coordinates of a lift. In particular it does not assert
that an element of \(P_D\) has nontrivial harmonic residue.
It shows why triviality of all single-prime word actions does
not by itself settle a finite quotient with translation and
rotation of different prime orders.

For a geometric attempt using the full \(\ell\)-division
translation group on a torus, the multiplication-by-\(\ell\)
cover is characteristic on closed torus homology. Its liftable
mapping-class subgroup is therefore the full mapping-class
group. This removes the particular normalizer-surjectivity
obstruction of a cyclic isogeny. On first homology, however,
the resulting intertwining is initially just commutation with
the scalar \(\ell I\). The finite proposition shows that the
translation and orbifold inertia data still permit both
Gaussian square classes. A proof of the desired relation
must additionally identify the tangential or rotation
correction with the divided Kummer coordinate
\(\rho_\ell^I/2\bmod2\). That identification remains open here.

\subsection{An explicit division-point norm for the harmonic sign}

Let \(E:y^2=x^3-x\), with the order-four automorphism
\(\iota(x,y)=(-x,iy)\). Fix an odd prime \(\ell\), and put
\(\chi=\chi_{-4}(\ell)\).
For \(Q\in E[\ell]\setminus\{0\}\), define
\[
        F(Q)=\frac{x(Q)^2+1}{2y(Q)}.
\]
Both numerator and denominator are nonzero: a zero denominator
would make \(Q\) two-torsion, and the duplication formula shows
that a zero numerator would make \(2Q=(0,0)\).
The same duplication formula gives
\begin{equation}\label{eq:gaussian-norm-function}
       F(Q)^2=x(2Q),\qquad F(\iota Q)=-iF(Q).
\end{equation}
Choose one representative of each \(\langle\iota\rangle\)-orbit
on \(E[\ell]\setminus\{0\}\), and let
\[
       C_\ell=\prod_{Q\ {\rm representative}}F(Q).
\]
Changing the representatives multiplies \(C_\ell\) by a fourth
root of unity.

\begin{proposition}\label{prop:gaussian-division-norm}
One has the exact identity
\[
                         C_\ell^4=\frac{\chi}{\ell}.
\]
For the action of the residue norm-one torus
\[
 T_\ell=
 \{a\in(\mathbf F_\ell[i])^\times:a\overline a=1\}
\]
on the \(\ell\)-torsion points, the multiplier of \(C_\ell\)
is its quadratic character:
\[
        \frac{a(C_\ell)}{C_\ell}
                    =a^{(\ell-\chi)/2}\in\{1,-1\}.
\]
The notation on the left means the monomial action on the factors,
using the second identity in
\eqref{eq:gaussian-norm-function}; it does not presume an
arithmetic realization of \(a\).
\end{proposition}
\begin{proof}
We recall the division-polynomial computation, including its
normalization. On \(y^2=x^3-x\), the standard addition-law
recursions begin with
\[
 \psi_0=0,\quad \psi_1=1,\quad \psi_2=2y,\quad
 \psi_3=3x^4-6x^2-1,\quad
 \psi_4=4y(x^6-5x^4-5x^2+1)
\]
and are
\[
 \psi_{2m+1}=\psi_{m+2}\psi_m^3-\psi_{m-1}\psi_{m+1}^3,
 \qquad
 \psi_{2m}=\frac{\psi_m}{2y}
       \bigl(\psi_{m+2}\psi_{m-1}^2
                       -\psi_{m-2}\psi_{m+1}^2\bigr).
\]
The apparent denominator cancels in the even recursion.
These are the recursions obtained from the chord-and-tangent
addition formulas. In particular \(\psi_\ell\) has leading
coefficient \(\ell\), degree \((\ell^2-1)/2\), and simple roots
at the abscissae of nonzero \(\ell\)-torsion points, modulo sign.

At \((0,0)\), every even division polynomial is zero.
The odd recursion proves by induction that
\[
                \psi_n(0)=(-1)^{(n-1)/2}\qquad(n\text{ odd}).
\]
Indeed for \(n=2m+1\) with \(m=2k\) the surviving term is
\(-\psi_{2k-1}\psi_{2k+1}^3\), of sign \(+1\);
with \(m=2k+1\) it is
\(\psi_{2k+3}\psi_{2k+1}^3\), of sign \(-1\).
Let \(M=(\ell^2-1)/4\), which is even.
Pairing the abscissae under \(\iota\) gives
\[
       \psi_\ell(x)=
          \ell\prod_{Q\ {\rm representative}}
                       (x^2-x(Q)^2).
\]
Evaluating at zero, and using that \(M\) is even, proves
\[
            \prod_{Q\ {\rm representative}}x(Q)^2
                         =\frac{\chi}{\ell}.
\]
Multiplication by two permutes the orbit set. Equation
\eqref{eq:gaussian-norm-function} now gives the claimed fourth
power of \(C_\ell\).

For the second assertion, identify \(E[\ell]\) with the
two-dimensional \(\mathbf F_\ell\)-space \(\mathbf F_\ell[i]\).
The torus \(T_\ell\) is cyclic of order \(n=\ell-\chi\),
and it acts freely on the nonzero vectors. In the nonsplit
case this follows because every nonzero vector is a field unit.
In the split case a nonidentity element is
\((a,a^{-1})\) with \(a\ne1\); neither eigenvalue is one,
so it also fixes no nonzero vector. The number of torus orbits
is
\[
                \frac{\ell^2-1}{\ell-\chi}=\ell+\chi,
                                      \qquad \ell+\chi\equiv2\pmod4.
\]
Choose a generator \(b\) of \(T_\ell\) with \(b^{n/4}=\iota\).
On each regular \(T_\ell\)-orbit, choose representatives
\(v,bv,\ldots,b^{n/4-1}v\) for its \(\langle\iota\rangle\)-orbits.
Multiplication by \(b\) permutes all these factors except the
last, which acquires the multiplier \(-i\).
Over all \(\ell+\chi\) torus orbits the multiplier is therefore
\[
                       (-i)^{\ell+\chi}=-1.
\]
For \(a=b^r\), it is \((-1)^r=a^{n/2}\), as asserted.
Changing the chosen representatives does not affect the result,
because \(T_\ell\) commutes with \(\iota\).
\end{proof}

For an arithmetic element \(\sigma\in H^{\rm cyc}\), its action
on \(E[\ell]\) is in \(T_\ell\).
The first assertion places \(C_\ell\) in
\(\mu_8\ell^{-1/4}\). Since \(\sigma\) fixes all roots of unity,
\[
        \frac{\sigma(C_\ell)}{C_\ell}
                =i^{-\rho_\ell^I(\sigma)}
                =(-1)^{\rho_\ell^I(\sigma)/2}.
\]
Thus the proposition recovers the arithmetic Gaussian square
relation directly from an explicit division-point product.
The evenness of \(\rho_\ell^I\) on the cyclotomic kernel is
used in the last equality.

For a full-GT element, the finite-cover and harmonic-frame
comparison supplies the torus action and hence the monomial
character in the proposition. It does not by itself identify
that action on the product with the Ihara Kummer action on
\(\ell^{-1/4}\). The outstanding comparison can therefore be
stated on this particular norm tensor: compatibility must be
proved for the product of the functions \(F(Q)\), retaining
the fourth-root normalization. The division-polynomial identity
and the finite torus calculation alone do not establish that
compatibility for full GT.

\subsection{The cover carrying the Gaussian norm tensor}

Let \(E:y^2=x^3-x\), let \(T_0=(0,0)\),
\(T_+=(1,0)\), \(T_-=(-1,0)\), and put
\[
                         F(Q)=\frac{x(Q)^2+1}{2y(Q)}.
\]
The notation \(\tau_T\) means translation by \(T\).

\begin{lemma}\label{lem:gaussian-norm-v4-cover}
The morphism \(F:E\to\mathbf P^1\) is a Galois cover
of degree four with deck group
\[
 \{1,\tau_{T_0},\tau_{T_+}[-1],\tau_{T_-}[-1]\}
                              \simeq(\mathbf Z/2)^2.
\]
Its branch values are \(1,-1,i,-i\).  It restricts to
an unramified cover
\[
 E\setminus E[4]\longrightarrow
       \mathbf P^1\setminus\{0,\infty,1,-1,i,-i\}.
\]
For every odd prime \(\ell\), the values \(F(Q)^4\),
with one \(Q\) in each order-four CM orbit on
\(E[\ell]\setminus\{0\}\), form a configuration of
\(M=(\ell^2-1)/4\) distinct points of
\(\mathbf P^1\setminus\{0,1,\infty\}\).
\end{lemma}
\begin{proof}
The addition formulas give
\[
 \tau_{T_0}(x,y)=(-1/x,y/x^2),\qquad
 \tau_{T_+}(x,y)=((x+1)/(x-1),-2y/(x-1)^2).
\]
Thus \(F\circ\tau_{T_0}=F\),
\(F\circ\tau_{T_+}=-F\), and \(F\circ[-1]=-F\).
The four displayed automorphisms preserve \(F\) and form
a four-group.  The poles of \(F\) are the four points
of \(E[2]\), each simple, so its degree is four.  This
proves the deck-group assertion.

The nontrivial translation has no fixed points.  The fixed
points of \(\tau_{T_+}[-1]\), respectively
\(\tau_{T_-}[-1]\), satisfy \(2Q=T_+\), respectively
\(2Q=T_-\).  Since \(F(Q)^2=x(2Q)\), their images are
respectively \(\{1,-1\}\) and \(\{i,-i\}\).
The fibers over zero and infinity are respectively
\([2]^{-1}(T_0)\) and \(E[2]\).  Together with the four
branch fibers, these are exactly \(E[4]\), proving the
unramified assertion.

The function \(F(Q)^4=x(2Q)^2\) is the order-four CM
quotient of \(2Q\).  Multiplication by two is a bijection
on \(E[\ell]\), commuting with the CM automorphism.
Thus distinct CM orbits give distinct values.  None of
these values is zero, one, or infinity, since the omitted
fibers consist of four-torsion points.
\end{proof}

This describes a second, explicit cover which carries the
product tensor.  Let
\(\mathcal C_M=\operatorname{Conf}_M
(\mathbf P^1\setminus\{0,1,\infty\})\) be the space of
unordered configurations of \(M\) distinct points.
The product of their affine coordinates is a regular unit
on \(\mathcal C_M\).  The equation
\begin{equation}\label{eq:gaussian-norm-configuration-cover}
                     c^4=\prod_{j=1}^{M}z_j
\end{equation}
therefore defines a finite etale cyclic cover of degree
four over \(\mathbf Q(i)\).  Its local monodromy is one
modulo four when one point travels around zero, minus one
around infinity, and zero around one or a collision
divisor.  These statements follow by the respective
orders of the product unit; in particular a braid
exchanging two points introduces no additional sign.

At the configuration in Lemma~\ref{lem:gaussian-norm-v4-cover},
the fourth-root fiber is exactly the norm tensor
\(C_\ell=\prod_Q F(Q)\).  The division-polynomial identity
identifies its fourth power as \(\chi_{-4}(\ell)/\ell\).
Consequently this is a specified finite cover and a
specified fiber in which to formulate the outstanding
full-GT comparison.

Preservation of the geometric monodromy class of
\eqref{eq:gaussian-norm-configuration-cover} does not yet
specify the action on its fiber at that configuration.
The missing assertion identifies the action transported
from the CM division configuration, including its
tangential normalization, with the Ihara Kummer action
on \(\ell^{-1/4}\).  Neither the four-group deck description
nor a normalizer statement for the full translation
group supplies that identification by itself.

One can also see the same missing descent in gamma
notation.  The adelic gamma function takes values in
Witt vectors over algebraic closures of the residue
fields, not necessarily in the finite residue field
containing a beta value; see
\cite[\S2.3, p.~101]{Furusho2003}.
An expression for a beta value as a square of such a
gamma value therefore does not prove that it is a square
in the rational norm-one torus.  The Frobenius action on
that square root is part of the remaining harmonic
residue condition.

\subsection{What survives a deformation of the Gaussian tensor}

The fourth-root tensor has a useful quadratic extension away
from the CM point.  This also specifies a limitation of a
direct argument by degeneration to the Tate cusp.

\begin{proposition}\label{prop:generic-elliptic-division-norm}
Let
\[
 E_{A,B}:y^2=x^3+Ax^2+Bx,\qquad B(A^2-4B)\ne0,
\]
over a field of characteristic zero.  Thus \((0,0)\) is a
specified nonzero two-torsion point.  For an odd prime
\(\ell\), put \(M=(\ell^2-1)/4\), which is even, and
choose one point in each pair \(\{Q,-Q\}\) of nonzero
\(\ell\)-torsion points.  Define
\[
 F_B(Q)=\frac{x(Q)^2-B}{2y(Q)},\qquad
 D_\ell=B^{-M/2}\prod_{Q\bmod\{\pm1\}}F_B(Q).
\]
Then
\begin{equation}\label{eq:generic-elliptic-norm-square}
                          D_\ell^2=\frac{\chi_{-4}(\ell)}{\ell}.
\end{equation}
Changing representatives changes \(D_\ell\) only by a sign.
It is unchanged by the coordinate scaling
\((x,y,A,B)\mapsto(u^2x,u^3y,u^2A,u^4B)\).

For \(A=0,B=-1\), group the representatives in pairs
\(Q,\iota Q\), one such pair for each order-four CM orbit,
where \(\iota(x,y)=(-x,iy)\).  Then the tensor
\(C_\ell=\prod_{Q\bmod\langle\iota\rangle}F_{-1}(Q)\)
satisfies, with those representative choices,
\begin{equation}\label{eq:quadratic-norm-is-square-of-cm-norm}
                                  D_\ell=C_\ell^2.
\end{equation}
\end{proposition}
\begin{proof}
The addition law gives the duplication identity
\[
                         x(2Q)=\frac{(x(Q)^2-B)^2}{4y(Q)^2}.
\]
The denominator and numerator of \(F_B(Q)\) are nonzero
for nonzero odd-order \(Q\): otherwise \(Q\) is
two-torsion, or \(2Q=(0,0)\), respectively.

The odd division polynomial has the form
\[
                \psi_\ell(x)=\ell
                      \prod_{Q\bmod\{\pm1\}}(x-x(Q)).
\]
We check its constant coefficient.  The initial values at
\((0,0)\) are \(\psi_1=1\), \(\psi_3=-B^2\), and
every even division polynomial vanishes there.  The usual
addition-law recursion
\[
             \psi_{2m+1}=\psi_{m+2}\psi_m^3
                              -\psi_{m-1}\psi_{m+1}^3
\]
then gives, for every odd \(n\),
\[
                         \psi_n(0)=(-1)^{(n-1)/2}B^{(n^2-1)/4}.
\]
For clarity, when \(n=4k+1\), the only surviving term is
\(-\psi_{2k-1}\psi_{2k+1}^3\).  Its sign is positive
and its exponent of \(B\) is
\((k^2-k)+3(k^2+k)=4k^2+2k\).
When \(n=4k+3\), the surviving term is
\(\psi_{2k+3}\psi_{2k+1}^3\).  Its sign is negative
and its exponent is
\((k^2+3k+2)+3(k^2+k)=4k^2+6k+2\).
These are precisely the required exponents.

The number of factors in the displayed product for
\(\psi_\ell\) is \(2M\), which is even.  Evaluation
at zero therefore gives
\[
                    \prod_{Q\bmod\{\pm1\}}x(Q)
                         =\frac{\chi_{-4}(\ell)B^M}{\ell}.
\]
Multiplication by two permutes the nonzero torsion pairs.
Squaring the product defining \(D_\ell\), and applying
the duplication identity, proves
\eqref{eq:generic-elliptic-norm-square}.
The sign assertion follows from \(F_B(-Q)=-F_B(Q)\).
Under the coordinate scaling, each factor of \(F_B\)
is multiplied by \(u\); there are \(2M\) factors.
The normalizing term \(B^{M/2}\) is also multiplied
by \(u^{2M}\), proving invariance.

At the CM point,
\(F_{-1}(\iota Q)=-iF_{-1}(Q)\).  Thus
\[
 D_\ell=(-1)^{-M/2}(-i)^M C_\ell^2=C_\ell^2,
\]
because \(M\) is even.  This proves
\eqref{eq:quadratic-norm-is-square-of-cm-norm}.
\end{proof}

The generic identity controls a square root of
\(\chi_{-4}(\ell)/\ell\).  At the Gaussian point,
the desired tensor is one further square root.
After the value of \(D_\ell\) has been fixed, the two
possible values \(C_\ell\) and \(-C_\ell\) remain.
The involution \(C_\ell\mapsto-C_\ell\) preserves
both \eqref{eq:generic-elliptic-norm-square} and
\eqref{eq:quadratic-norm-is-square-of-cm-norm}.
Consequently a compatibility proved only for the generic
quadratic norm does not establish the fourth-root
compatibility needed for the harmonic residue character.

\begin{proposition}\label{prop:no-cm-preserving-tate-degeneration}
A family of smooth elliptic curves in characteristic zero,
equipped with an origin-preserving automorphism of exact
order four, has constant \(j\)-invariant \(1728\).
It cannot have a degeneration to a multiplicative Tate
cusp while retaining that automorphism on the generic
elliptic curve.  If additional marked points are allowed
to collide in a stable compactification, forgetting
those points and stabilizing still has a smooth
elliptic core of \(j\)-invariant \(1728\).
\end{proposition}
\begin{proof}
Over an algebraic closure of the generic field, put the
curve in short Weierstrass form
\(y^2=x^3+ax+b\).  An automorphism fixing the origin
has the form \((x,y)\mapsto(u^2x,u^3y)\).
Exact order four gives \(u^2=-1\).  Preservation of
the equation then forces \(b=0\); smoothness forces
\(a\ne0\).  The formula
\[
                  j=1728\frac{4a^3}{4a^3+27b^2}
\]
gives \(j=1728\).

For a multiplicative degeneration over a trait, a minimal
semistable equation has nonzero reduction of \(c_4\)
and vanishing discriminant.  Hence
\(j=c_4^3/\Delta\) has a pole.  This is incompatible
with the constant value \(1728\).  Equivalently, in
the compactified moduli of one-pointed elliptic curves,
the boundary has \(j=\infty\), whereas the constant
map has image \(1728\).  The same conclusion after
forgetting and stabilizing additional marks proves the
last assertion.
\end{proof}

Thus a direct Tate-cusp comparison that keeps the
order-four orbit decomposition throughout cannot supply
the desired normalization: that fixed locus has no such
cusp.  If the CM automorphism is forgotten before
deforming, the norm naturally displayed above is
\(D_\ell=C_\ell^2\), and an additional comparison is
needed to recover the remaining sign.  This is an
obstruction to that particular degeneration argument,
not a proof that another full-GT comparison is impossible.

\section{The two first degeneracy correspondences}
The first degeneracy correspondence is already governed by relation (IV).
Put $B=\operatorname{PSL}_2(\mathbb Z)$ and write
\[
 u=\begin{pmatrix}1&1\\0&1\end{pmatrix},\qquad
 v=\begin{pmatrix}1&0\\-1&1\end{pmatrix}
\]
for their classes in $B$. The first map from $\Gamma_0(N)$ to $B$ is inclusion.
The second is conjugation by $D_N=\operatorname{diag}(N,1)$; denote it by $d_N$.
Thus $d_N(u)=u^N$ and $d_N(v^N)=v$. These maps extend to full profinite completions,
since each of their images has finite index in $B$.

\begin{proposition}\label{prop:degeneracy-two}
Let $F=(1,f)$ be a normalized element of $\widehat{\operatorname{GT}}$.
Let $\Phi_F$ be its standard action on $\widehat B$, so that
\[
 \Phi_F(u)=u,\qquad
 \Phi_F(v)=q^{-1}vq,\qquad q=f(u^2,v^2).
\]
Then $\Phi_F$ preserves the closure of $\Gamma_0(2)$, and its restriction
commutes with both degeneracy maps to $\widehat B$.
\end{proposition}
\begin{proof}
The standard index-three subgroup $\Gamma_0(2)$ is generated by $u,v^2$;
its presentation is $\langle u,w\mid(uw)^2=1\rangle$ with $w=v^2$.
This follows, for example, by applying the Schreier method to the action of
$B$ on $\mathbb P^1(\mathbb F_2)$: the stabilizer of infinity is $\Gamma_0(2)$,
and the two resulting generators can be chosen as $u,v^2$.
Because $q$ belongs to $\overline{\langle u^2,v^2\rangle}$, the formulas for
$\Phi_F$ take both generators into $\overline{\Gamma_0(2)}$.
The same argument for $F^{-1}$ gives equality of this subgroup with its image.
Commutation with inclusion is now its definition. For the other map,
\[
 d_2(q)=f(u^4,v)=f(u^2,v^2)=q.
\]
The middle equality follows by applying the interchange of $u,v$ to normalized
relation (IV), then applying $f(b,a)=f(a,b)^{-1}$ to both sides.
On the generator $u$ the two compositions both have value $u^2$.
On $v^2$ they both have value $q^{-1}vq$. Equality on these topological
generators proves the assertion.
\end{proof}

\section{Normalized congruence triviality and the level-three gate}
\label{sec:normalized-congruence}

The level-two identity has a consequence at every congruence level.
Use the signed representation
$\widetilde q:\widehat B_3\to\operatorname{SL}_2(\widehat{\mathbb Z})$
with $u\mapsto U=\left(\begin{smallmatrix}1&1\\0&1\end{smallmatrix}\right)$
and $v\mapsto V=\left(\begin{smallmatrix}1&0\\-1&1\end{smallmatrix}\right)$.

\begin{proposition}\label{prop:normalized-congruence-triviality}
Let $g=(1,f)\in\widehat{\mathrm{GT}}$ satisfy $\rho_2(g)=0$.
Then
\[
 f(U^2,V^2)=I
       \qquad\text{in }\operatorname{SL}_2(\widehat{\mathbb Z}).
\]
Consequently $\widetilde q\Phi_g=\widetilde q$ on $\widehat B_3$.
\end{proposition}
\begin{proof}
Normalized relation (IV), together with interchange of the two braid
generators and the two-cycle relation, gives
\[
 f(u^4,v)=f(u^2,v^2)
       \qquad\text{in }\widehat B_3.
\]
Fix a prime $\ell$ and put $Q=f(U^2,V^2)$ in
$\operatorname{SL}_2(\mathbb Z_\ell)$.
Conjugation by $D=\operatorname{diag}(2,1)$ in
$\operatorname{GL}_2(\mathbb Q_\ell)$ sends $U^2$ to $U^4$ and
$V^2$ to $V$. Therefore the displayed signed identity gives
$DQD^{-1}=Q$, with no scalar ambiguity. Since the two diagonal
entries of $D$ are distinct, $Q=\operatorname{diag}(a,a^{-1})$
for a unit $a\in\mathbb Z_\ell^\times$.

The images of the braid generators under
$\widetilde q_\ell\Phi_g$ are $U$ and
\[
 M=Q^{-1}VQ=\begin{pmatrix}1&0\\-a^2&1\end{pmatrix}.
\]
Their braid relation $UMU=MUM$, in its lower-left entry, is
$\gamma(1+\gamma)=0$ with $\gamma=-a^2$.
The coefficient ring is a domain and $a$ is a unit, so $a^2=1$.
Thus $Q=\pm I$ and $M=V$.

It remains to determine the sign using the three-cycle relation.
Put $A=U^2$, $B=V^2$, $Z=(AB)^{-1}$, and $R=UV$.
Direct multiplication gives
\[
 RAR^{-1}=B,\qquad RBR^{-1}=-Z,\qquad RZR^{-1}=-A.
\]
A commutator word is unchanged when one of its arguments is multiplied
by a central scalar. Since $Q$ is central, conjugating its defining
expression by $R$ and $R^2$ consequently gives
\[
 f(A,B)=f(B,Z)=f(Z,A)=Q.
\]
The multiplier-one three-cycle relation, with $ABZ=1$, now implies
$Q^3=I$. Together with $Q^2=I$, this forces $Q=I$.
This holds for every $\ell$, proving the signed adelic assertion.
Both braid generators now have unchanged images, which proves the
intertwining identity.
\end{proof}

In the following statements the ambient modular group is the
centerless group $B=\operatorname{PSL}_2(\mathbb Z)$ and its full
profinite completion, rather than the central braid extension. Put
\[
 \mathcal C=\ker\bigl(\widehat B\longrightarrow
      \operatorname{SL}_2(\widehat{\mathbb Z})/\{\pm I\}\bigr),
\]
where the scalar sign is the common sign in all components.

\begin{lemma}\label{lem:parabolic-centralizer}
The centralizer of $v$ in $\widehat B$ is
$\overline{\langle v\rangle}\simeq\widehat{\mathbb Z}$, and
it has trivial intersection with $\mathcal C$.
\end{lemma}
\begin{proof}
The free pure subgroup is
$\widehat F_2=\overline{\langle x\rangle}\amalg
\overline{\langle y\rangle}$, where $x=u^2,y=v^2$.
The centralizer theorem for a nonidentity element of a factor in a
free profinite product gives
$C_{\widehat F_2}(y)=\overline{\langle y\rangle}$
\cite[Theorem B, p.~155; Theorem B$'$, \S4, p.~159]{HerfortRibes1985}.
This theorem has no torsion assumption on the element of the factor.
If $c\in C_{\widehat B}(v)$, its image in $S_3$ centralizes the
transposition represented by $v$. Thus $cv^{-e}$ is pure for
some $e\in\{0,1\}$ and centralizes $y$. It lies in
$\overline{\langle y\rangle}$, proving the centralizer assertion.
Finally $v^t$ has congruence image represented by
$\left(\begin{smallmatrix}1&0\\-t&1\end{smallmatrix}\right)$.
It is a scalar sign only if $t=0$. This proves injectivity of the
procyclic subgroup in the congruence quotient and the last assertion.
\end{proof}

\begin{corollary}\label{cor:level-three-exact-gate}
For a normalized $g=(1,f)\in\widehat{\mathrm{GT}}$, put
$q=f(u^2,v^2)\in\widehat B$. The element $q$ lies in
$\mathcal C$, hence in $\overline{\Gamma_0(3)}$.
The automorphism $\Phi_g$ preserves that subgroup. For
$d_3:\overline{\Gamma_0(3)}\to\widehat B$ induced by
conjugation by $\operatorname{diag}(3,1)$, set
\[
 h_3=d_3(q)q^{-1}.
\]
Then $h_3\in\mathcal C$, and strict compatibility with both
level-three degeneracy maps is equivalent to the exact identity
\[
 d_3(q)=q.
\]
\end{corollary}
\begin{proof}
The preceding proposition shows that $q$ and the action of $g$
are trivial at every congruence level. Therefore $q$ belongs to
$\mathcal C$ and $\Phi_g$ preserves the level-three subgroup.
On its congruence image, $d_3$ is conjugation by the same diagonal
matrix: the equality holds on the discrete subgroup and extends
by continuity. Consequently $d_3(q)$ also has trivial congruence
image, and $h_3\in\mathcal C$.

The subgroup $\Gamma_0(3)$ is generated by $u,v^3$, and
$d_3(u)=u^3$, $d_3(v^3)=v$. The two proposed compositions agree
on $u$. On $v^3$ they agree exactly when
$h_3\in C_{\widehat B}(v)$. The preceding lemma and
$h_3\in\mathcal C$ make this equivalent to $h_3=1$.
\end{proof}

The corollary identifies the remaining full profinite identity. It does
not prove it: every congruence representation already kills $h_3$, so
congruence matrices alone cannot decide whether this element is trivial
in $\widehat B$.


\section{One local Brauer class for the two frame signs}
The two frame signs admit one local Brauer invariant which can be defined
from the traces of the two matrices. No choices of centralizer generators
or common matrix basis enter its definition.

\begin{proposition}\label{prop:local-brauer-frame}
Let
\[
 \mathcal D=\{(T,W)\in\operatorname{SL}_2(\mathbb Z_3)^2:
       \operatorname{tr}(T)\operatorname{tr}(W^2)\in\mathbb Z_3^\times\}.
\]
On this domain define
\[
 \mathfrak b(T,W)=
 \big[(\operatorname{tr}(T)\operatorname{tr}(W^2),3)_{\mathbb Q_3}\big]
       \in\operatorname{Br}(\mathbb Q_3)[2].
\]
Here brackets denote a quaternion algebra's Brauer class.
This invariant is unchanged by independent conjugation of $T$ and $W$.
Suppose their reductions have the form
\[
 \bar T\in(-I)^{\epsilon_3}\langle R_0\rangle,
 \qquad \bar W\in S_0^{\epsilon_2}\langle-I\rangle,
\]
where $R_0$ is the standard nonidentity unipotent matrix of order three
in $\operatorname{SL}_2(\mathbb F_3)$ and $S_0^2=-I$.
Then $(T,W)\in\mathcal D$ and
\[
 \mathfrak b(T,W)=(\epsilon_2+\epsilon_3)
                 [(-1,-3)_{\mathbb Q_3}].
\]
Its local Hilbert sign is $(-1)^{\epsilon_2+\epsilon_3}$.
\end{proposition}
\begin{proof}
Every power of $R_0$ has trace $2$ in $\mathbb F_3$, including its
zeroth power. Also $\bar W^2=(-I)^{\epsilon_2}$. Consequently
\[
 \operatorname{tr}(T)\equiv2(-1)^{\epsilon_3},\qquad
 \operatorname{tr}(W^2)\equiv2(-1)^{\epsilon_2}\pmod3.
\]
Both traces are units and their product has residue
$(-1)^{\epsilon_2+\epsilon_3}$. By Hensel's lemma, a unit of
$\mathbb Z_3$ is a square in $\mathbb Q_3$ precisely when its residue
is a square in $\mathbb F_3$. The odd-prime Hilbert formula gives
$(e,3)_3=(e/3)$ for a unit $e$. In particular $(-1,-3)_3=-1$,
so this quaternion algebra represents the unique nonzero class of
$\operatorname{Br}(\mathbb Q_3)[2]$. The formula follows.
Both traces are unchanged by independent conjugation, which proves
its invariance on the entire stated domain.
\end{proof}

For comparison with centralizer coordinates, suppose
\[
 T=aI+bR,\quad W=cI+dS,\qquad
 R^2+R+I=0,\quad S^2=-I,
 \qquad a^2-ab+b^2=c^2+d^2=1,
\]
where $R,S$ are the standard cubic and harmonic matrices over $\mathbb Z_3$.
The trace identities give
\[
 \operatorname{tr}(T)=2a-b\equiv2(a+b),\qquad
 \operatorname{tr}(W^2)=2(c^2-d^2)\pmod3.
\]
The norm equations imply that $a+b$ and $c^2-d^2$ are units. Thus
\[
 \mathfrak b(T,W)=\big[((a+b)(c^2-d^2),3)_{\mathbb Q_3}\big].
\]
This recovers the product of the two previously given local symbols,
while the trace expression makes the definition intrinsic to the ordered
pair of matrices.

The connection with the modular frame character is exact on its stated
domain. For $F\in H_0$, take the signed three-adic images of $A_F$ and
$\beta_F$ as $T$ and $W$. Their reductions satisfy the displayed conditions
with $\epsilon_3=a_3(F)$ and $\epsilon_2=a_2(F)$, by the definition of the
frame character. Therefore
\[
 \mathfrak b(T,W)=\delta(F)[(-1,-3)_{\mathbb Q_3}].
\]
Here $H_0=\{F:\lambda_F=1,\ t_F=0\pmod{12}\}$; this conclusion uses only
the finite frame conditions and does not require the uncorrected
three-adic matrices themselves to centralize the elliptic matrices.
For corrected frame pairs satisfying the centralizer norm forms above,
the same formula holds with their specified signed frame bits.

The class can be nonzero on the matrix domain: $T=-I,W=I$ gives its
nonzero value. This matrix calculation alone does not prove that such a
pair comes from any designated subgroup. On normalized elements of
$\widehat{\operatorname{GT}}$, the class vanishes by \cite[Theorem~3]{H17}. The algebra $(-1,-3)$ over $\mathbb Q$
has ramification set $\{3,\infty\}$, whereas $(2,3)$ has ramification set
$\{2,3\}$; these rational algebras cannot be identified merely because
their completions at three are isomorphic.

\section{The quartic Kummer layer over the Gaussian field}
The mod-four Kummer extension is a classical dihedral extension.
The following computation fixes its Hecke conductors and specifies the
additional analytic hypothesis that a square-root error term would use.
Let $K=\mathbb Q(i)$, let $\mathfrak p=(1+i)$, and put
$L=K(\alpha)$ with $\alpha=2^{1/4}>0$.
For a prime ideal $\mathfrak q\nmid2$ of $K$, use arithmetic Frobenius and define
\[
 \psi(\mathfrak q)=\frac{\operatorname{Frob}_{\mathfrak q}(\alpha)}{\alpha}
  =\left(\frac{2}{\mathfrak q}\right)_4.
\]
The last expression is the unique fourth root of unity whose reduction modulo
$\mathfrak q$ is $2^{(N\mathfrak q-1)/4}$. Extend multiplicatively to ideals
prime to two.

\begin{proposition}\label{prop:quartic-conductor}
The extension $L/K$ is cyclic of degree four. Its two faithful finite-order
Hecke characters are $\psi,\psi^{-1}$ and
\[
 \mathfrak f(\psi)=\mathfrak f(\psi^{-1})=\mathfrak p^6,
 \qquad \mathfrak f(\psi^2)=\mathfrak p^4.
\]
All three characters have trivial infinity type. The absolute discriminant of
$L$ is $2^{24}$. The representation
$\operatorname{Ind}_{G_K}^{G_{\mathbb Q}}\psi$ is odd and irreducible of
dimension two, has conductor $256$, and has determinant $\chi_{-4}$.
\end{proposition}
\begin{proof}
The polynomial $X^4-2$ is Eisenstein at two. Its real root field has degree four
and does not contain $i$, so $[L:\mathbb Q]=8$ and $[L:K]=4$.
The automorphism $r(\alpha)=i\alpha$, $r(i)=i$ generates
$\operatorname{Gal}(L/K)$; complex conjugation $s$ satisfies $srs=r^{-1}$.
This proves the dihedral description. Artin reciprocity identifies the
characters of this finite cyclic quotient with finite-order Hecke characters.
The displayed Frobenius formula follows by reducing $\alpha^{N\mathfrak q-1}$.
There is no ramification away from two, since the polynomial discriminant
of $X^4-2$ is a power of two and $K/\mathbb Q$ is unramified away from two.

We compute the different at two directly. Let $M=\mathbb Q(\zeta_8)=K(\sqrt2)$,
let $\varpi=1-\zeta_8$, and normalize the local valuation on $M$ by
$v_M(\varpi)=1$. Then $v_M(2)=4$ and $v_M(\sqrt2)=2$.
Put $D=\sqrt2/\varpi^2$. From
\[
 \sqrt2=\zeta_8+\zeta_8^{-1}
        =\frac{2-2\varpi+\varpi^2}{1-\varpi}
\]
we obtain $D\equiv1+\varpi\pmod{\varpi^2}$. For
$\beta=\alpha/\varpi$ and $\gamma=\beta-1$ we have
\[
 \gamma^2+2\gamma+(1-D)=0.
\]
This polynomial is Eisenstein over $M$ because $v_M(1-D)=1$.
Thus $L/M$ is totally ramified quadratic at two, $\gamma$ is a uniformizer,
and its different is generated by the derivative $2\gamma+2=2\beta$.
Since $\beta$ is a unit and $v_L(2)=8$, the different exponent is eight.
The discriminant of $M/\mathbb Q$ is $2^8$, from the cyclotomic polynomial
$X^4+1$. The discriminant formula in a tower therefore gives
$|d_L|=(2^8)^2\,2^8=2^{24}$.

Since $d_K=-4$, the relative discriminant ideals are
\[
 \mathfrak d_{M/K}=\mathfrak p^4,
 \qquad \mathfrak d_{L/K}=\mathfrak p^{16}.
\]
Indeed their norms are respectively $2^8/4^2=2^4$ and
$2^{24}/4^4=2^{16}$. The conductor-discriminant formula
\cite[Chapter V, Theorem~3.27, p.~165]{MilneCFT}
gives the conductor of the quadratic character $\psi^2$ as $\mathfrak p^4$.
Writing the common conductor exponent of the inverse faithful characters as
$a$, the same formula gives $2a+4=16$, hence $a=6$.

Conjugation by $s$ sends $\psi$ to $\psi^{-1}\ne\psi$, so induction is
irreducible. On $r$ its two eigenvalues are $i,-i$, and on $s$ they are
$1,-1$. Its determinant is consequently the character of $K/\mathbb Q$,
and it is odd. The induction formula for conductors gives
$|d_K|N_{K/\mathbb Q}\mathfrak f(\psi)=4\cdot2^6=256$.
\end{proof}

Hecke's theta construction associates to $\psi$ the weight-one form
\[
 \theta_\psi(z)=\sum_{\substack{\mathfrak a\subset\mathbb Z[i]\\
                     (\mathfrak a,2)=1}}
                 \psi(\mathfrak a)e^{2\pi i N\mathfrak a z}.
\]
It is the classical newform of level $256$ and character $\chi_{-4}$, recorded
as LMFDB $256.1.c.a$ \cite{LMFDB256}. Its first coefficients are
$q+q^9-2q^{17}-q^{25}+O(q^{33})$.

For an unramified prime ideal $\mathfrak q$ of $K$, specifying the value
$\psi(\mathfrak q)=i^a$ is a cyclic Chebotarev condition. A square-root
error term obtained from the explicit formula uses the named hypothesis
that every nontrivial zero of $L_K(s,\psi^j)$, $j=0,1,2,3$, has real part
$1/2$, where $j=0$ means $\zeta_K(s)$. The factorization
\[
 L_K(s,\psi^2)=L_{\mathbb Q}(s,\chi_8)
                    L_{\mathbb Q}(s,\chi_{-8}),\qquad
 \zeta_K(s)=\zeta_{\mathbb Q}(s)L_{\mathbb Q}(s,\chi_{-4})
\]
shows the new analytic input at this layer to be the quartic Hecke
$L$-function (the inverse character has the same $L$-function).
If extra cyclotomic congruences are imposed, the corresponding finite family
of Hecke twists $\psi^j(\chi\circ N_{K/\mathbb Q})$ must be included.
The existence of appropriate primes, and a terminating search for one, does
not require GRH. Nor does the algebraic decomposition itself require it.
A quantitative assertion must state its requested bound before a stronger
analytic hypothesis is described as necessary.


\section{Questions remaining}
The full relative abelian group \(B\) and its primary quotients remain
undetermined, so equality \(G=G_{\Q}\) is not decided. The arithmetic
targets \(Y_p\) are completely computed for every prime. At prime two,
restriction is surjective onto the complete target, with the explicit
section in Theorem~\ref{thm:complete-quadratic-restriction}.
The relative kernel is a different question: the ordered differences
\(\varepsilon_{\ell,q}\), the harmonic residue discrepancies and
the finite-cover determinant give explicit candidates, but none is
proved nonzero.

The full relative abelian obstruction is preserved on \(R_{\rm global}\),
with \(G=R_{\rm global}H\). Its maximal pro-two quotient is preserved
on \(V_{\rm global}\). On the latter group all rational Kummer
coordinates, all single-prime word images, the entire coefficient-two
cubic image, both odd CM images, the chosen lifts of every corrected
cup-product character, and the scalars \(h_p\) for
\(p\in\Pi_{\rm global}\) vanish.
The prime set is governed by the exact real-unit square-class criterion
of Theorem~\ref{thm:main-cyclotomic-cm-normalization}. It includes
\(29\) and \(113\), whose full arithmetic images fill their local
index-eight deficits, and excludes \(163\).
The exact original restriction kernel is \(\mathcal S=s(Y_2)\),
and the restricted image is precisely the remaining relative kernel.
Determining the invariant quadratic characters of \(V_{\rm global}\)
that kill \(V_{\rm global}\cap H\) would settle the original quadratic
relative question. No equality \(G=V_{\rm global}H\) is asserted.

At degree \(163\), put \(R=R_{\rm global}\), \(P=P_{163}\), and
\(A=Q_{163}P\). It remains unknown whether \(h_{163}(R)=A\) or
\(h_{163}(R)=P\). The second alternative would give a quotient
\(A_4\) with arithmetic image of order three. The finite-word detector
in Theorem~\ref{thm:main-163-finite-word} distinguishes the alternatives
but does not choose one.

For an arbitrary collection of CM scalar kernels, including the
infinite collection, the conjugation obstruction has been eliminated.
The relative central class \(o_2\) in
Theorem~\ref{thm:main-all-cm-central-obstruction} is now the sole
obstruction to extending an intrinsic invariant quadratic character
that vanishes on the arithmetic intersection. Its vanishing is not
proved in general. At degree \(163\), conditional on the full-image
alternative, the possible invariant relative class is the
one-dimensional quaternion line of
Theorem~\ref{thm:main-163-odd-inertia}; this remaining central
obstruction is also undecided.

At odd primes the entire target survives after the rational,
single-prime, complete CM and finite-abelian-cover homology actions
are all killed, by the mixed-inertia construction above. Accordingly
these abelian-cover and primary realizations cannot furnish a
nonzero arithmetic restriction of an odd invariant character.
A construction using further nonabelian finite quotients, or an
obstruction in the full continuous cohomology of \(G\), is still
needed to determine the odd-prime image and relative kernel.

The radical and complete CM normalizations also have their own exact
joint-image question. Its remaining conditions are the tame cubic
and harmonic quadratic tests; the wild cubic, complete prime-two and
exceptional prime-three conditions are proved. The harmonic tests
are the relative characters \(\Xi_{\ell,4}\), equivalently the
adelic beta square tests on \(D\). Their geometric division-point
norm gives an explicit correspondence, but its identification with
the rational Kummer torsor for every element of full GT remains
unproved. The generic norm calculation determines only its square,
and the order-four CM locus cannot degenerate to a Tate cusp.

The determinant character \(\psi_{17,11}\) of
Corollary~\ref{cor:concrete-radical-anderson-character} likewise has
undecided vanishing. Strict level-three compatibility is reduced to
\(d_3(q)=q\) in the full profinite modular group; equality of all its
congruence images does not decide that word identity.
Comparison with another completed-tower construction still requires
its action formulas or an appropriate uniqueness theorem for its
objects and morphisms. For punctured surfaces, passing from
nonseparating twist inertia to all twist inertia additionally requires
recognition of multicurves with connected complement. Commutation of
individual twists recognizes disjointness and does not supply that
recognition.

\section*{Method and records}
The mathematics and text were produced by GPT-6 Astra Pro
(ChatGPT Pro mode). Separate same-model workers reconstructed main
statements before reading their authored proofs; complete comparisons
then checked the written arguments, conventions and cited hypotheses.
The accompanying statement-level records distinguish complete
reconstructions, partial attempts, proof-informed checks and independently
implemented calculations. They also disclose permitted background
results, source pointers and normalization hints. Earlier partial
attempts have not been relabelled after later successful checks.

The exact cubic frame identity and its finite-character consequences
have complete fresh independent reconstructions. The frame proof includes
the full profinite topology, identification of the descended parameter,
based signs and final central correction; its general nonzero-Kummer
variant is not asserted. The full harmonic beta comparison and its
residue-character assertion also have a complete fresh independent
reconstruction, accepted after a full comparison with the authored proof.
The earlier harmonic attempt remains partial at the geometric comparison.
The exact scope of each check is recorded with its statement and file
digest.

Finite computations use exact arithmetic and explicit reconstruction
programs. They support the stated finite calculations; they do not replace
the arguments passing to full profinite groups or infinite products.
The three stipulated background results are cited rather than reproved.
Current independent workers used the same model, with effort setting not
exposed. The cited background notes record checks by Claude Fable 5.2
and Claude Opus 5.5 at effort max. No human has checked this paper.
The text and original programs are dedicated to the public domain under
CC0 1.0. The workbench is \url{\workbenchurl}; the permanent records are
at \url{\recordurl}.

\begin{thebibliography}{99}
\bibitem{BBR2026} L.~Basualdo Bonatto and M.~Robertson,
Galois actions on surfaces and a higher genus
Grothendieck--Teichm\"uller group,
\href{https://arxiv.org/abs/2606.01466v2}{arXiv:2606.01466v2},
8 October 2026, Corollary~4.4, p.~39; Definition~4.8 and
Theorem~4.9, pp.~40--41; Corollary~6.13 and
Proposition~6.14, pp.~57--58.
\bibitem{Guillot2016} P.~Guillot, The Grothendieck--Teichm\"uller
group of $\mathrm{PSL}(2,q)$,
\href{https://arxiv.org/abs/1604.04415v1}{arXiv:1604.04415v1},
15 April 2016, Theorem~1.1, Proposition~3.2 and Corollary~3.3,
pp.~1--3.
\bibitem{BCL2022} F.~M.~Bleher, T.~Chinburg and A.~Lubotzky,
On representations of $\operatorname{Gal}(\overline\Q/\Q)$,
$\widehat{GT}$ and $\operatorname{Aut}(\widehat F_2)$,
J. Algebra \textbf{607} (2022), 134--159,
\href{https://arxiv.org/abs/2004.08860v3}{arXiv:2004.08860v3},
Theorems 1.1--1.3.
\bibitem{BDHP2026} I.~Bortnovskyi, V.~A.~Dolgushev, B.~Holikov
and V.~Pashkovskyi, Accessing non-abelian quotients of the
Grothendieck--Teichmueller group via elementary tools,
\href{https://arxiv.org/abs/2405.11725v3}{arXiv:2405.11725v3},
6 August 2026, Theorem 5.3, Corollary 5.4, Theorem 5.7
and Proposition 5.8.
\bibitem{Drinfeld1990} V. G. Drinfeld, On quasitriangular quasi-Hopf algebras and on a group that is closely connected with $\operatorname{Gal}(\overline{\Q}/\Q)$, Leningrad Math. J. \textbf{2} (1991), 829--860.
\bibitem{Ihara1999} Y. Ihara, On beta and gamma functions associated with the Grothendieck--Teichm\"uller group, in \emph{Aspects of Galois Theory}, London Math. Soc. Lecture Note Ser. 256, Cambridge University Press, 1999, 144--179, Proposition~1.6.1 and \S1.8.
\bibitem{IchimuraKaneko1989} H. Ichimura and M. Kaneko, On the universal power series for Jacobi sums and the Vandiver conjecture, J. Number Theory \textbf{31} (1989), 312--334, Section~1.2 and Theorem~1, pp.~314--315; \href{https://www2.math.kyushu-u.ac.jp/~mkaneko/papers/04universal_power.pdf}{author's scan}.
\bibitem{Anderson1989} G. W. Anderson, The hyperadelic gamma function, Invent. Math. \textbf{95} (1989), 63--131, Corollary 8.6.3; the beta formulation used here is stated in \cite[Remark 2.4, p.~101]{Furusho2003}.
\bibitem{HerfortRibes1985} W. Herfort and L. Ribes, Torsion elements and centralizers in free products of profinite groups, J. reine angew. Math. \textbf{358} (1985), 155--161, Theorem B, p.~155, and Theorem B$'$ in \S4, p.~159.
\bibitem{Furusho2003} H. Furusho, Geometric and arithmetic subgroups of the Grothendieck--Teichm\"uller group, Math. Res. Lett. \textbf{10} (2003), 97--108, \S2.2, p.~100.
\bibitem{Enriquez2007} B. Enriquez, Quasi-reflection algebras and cyclotomic associators, Selecta Math. (N.S.) \textbf{13} (2007), 391--463, \href{https://arxiv.org/abs/math/0408035v2}{arXiv:math/0408035v2}, \S8, Theorem 8.1, equation (50) and Remark 8.6, author pp.~30--32.
\bibitem{MoreeStevenhagen} P. Moree and P. Stevenhagen, Computing higher rank primitive root densities, Acta Arith. \textbf{163} (2014), 15--32, \href{https://arxiv.org/abs/1203.4313v2}{arXiv:1203.4313v2}, \S2, Theorem 2.5, author pp.~4--6.
\bibitem{NS2000} H. Nakamura and L. Schneps, On a subgroup of the Grothendieck--Teichm\"uller group acting on the tower of profinite Teichm\"uller modular groups, Invent. Math. \textbf{141} (2000), 503--560, Theorems 1.3--1.4 and Definition 5.1.
\bibitem{HLS} A. Hatcher, P. Lochak and L. Schneps, On the Teichm\"uller tower of mapping class groups, J. reine angew. Math. \textbf{521} (2000), 1--24.
\bibitem{LS2006} P. Lochak and L. Schneps, Open problems in Grothendieck--Teichm\"uller theory, in \emph{Problems on Mapping Class Groups and Related Topics}, Proc. Sympos. Pure Math. 74, American Mathematical Society, 2006, 165--186, \S1.
\bibitem{Lochak2025} P. Lochak, Automorphisms of profinite and procongruence curve complexes and the Grothendieck--Teichm\"uller group, Algebra i Analiz \textbf{37} (2025), no.~2, author version, Definition 5.4 and Theorem 5.10, pp.~50, 60--62.
\bibitem{LNS} P. Lochak, H. Nakamura and L. Schneps, The Grothendieck--Teichm\"uller group acts on the profinite mapping class group of genus $g$, \href{https://arxiv.org/abs/2602.12462}{arXiv:2602.12462}, Theorem 1.1.
\bibitem{Boggi} M. Boggi, Automorphisms of profinite mapping class groups, \href{https://arxiv.org/abs/2011.15075v6}{arXiv:2011.15075v6}, 22 April 2026, Theorem 3.12, Corollary 3.15, Remark 3.11, Proposition 6.9 and Lemma 6.16.
\bibitem{H17} \emph{Relations $(\mathrm{III}')$ and $(\mathrm{IV})$ of Nakamura and Schneps hold on the profinite Grothendieck--Teichm\"uller group, with an application to the two elliptic points of the modular curve}, version 2, 24 November 2026, \href{https://hexagonmath.org/2611.00940v2}{hexagon:2611.00940v2}, Theorems 1--3 and 20, Corollaries 21--22, Lemmas 5, 7 and 15, Propositions 11 and 18.
\bibitem{Monoid} \emph{The profinite Grothendieck--Teichm\"uller monoid is a group, on a question of Lochak and Schneps}, 25 November 2026, \href{https://hexagonmath.org/2611.01024}{hexagon:2611.01024}, Theorems A and E.
\bibitem{MilneCFT} J. S. Milne, \emph{Class Field Theory}, version 4.03, 2020, Chapter V, Theorem 3.27, p.~165, \url{https://www.jmilne.org/math/CourseNotes/CFT.pdf}.
\bibitem{LMFDB256} The LMFDB Collaboration, Classical modular form $256.1.c.a$, \url{https://www.lmfdb.org/ModularForm/GL2/Q/holomorphic/256/1/c/a/}, accessed 25 November 2026.
\bibitem{SchmidtWingberg} A. Schmidt and K. Wingberg, Shafarevich's theorem on solvable groups as Galois groups, \href{https://arxiv.org/abs/math/9809211}{arXiv:math/9809211}, Theorem 1, p.~1.
\bibitem{MilneMF110}
J. S. Milne, \emph{Modular Functions and Modular Forms},
version 1.10, 1997, Remark~12.12, Proposition~12.13,
and Theorem~12.27, printed pp.~119 and~124,
\url{https://www.jmilne.org/math/CourseNotes/MF110.pdf}.

\bibitem{PapikianTate}
M. Papikian, \emph{Rigid analytic geometry and abelian varieties},
notes dated 13 June 2004, \S2, Theorem~2.2 and pp.~3--4,
\url{https://math.stanford.edu/~vakil/snowbird/mihranjun21.pdf}.

\bibitem{Marin2008}
I. Marin, Characters of the Grothendieck--Teichm\"uller group through
rigidity of the Burau representation, Ann. Fac. Sci. Toulouse Math.
(6) \textbf{17} (2008), 765--780, \S6.3, p.~780,
\url{https://www.numdam.org/item/10.5802/afst.1202.pdf}.

\bibitem{NT2002}
H. Tsunogai and H. Nakamura, Harmonic equations in the
Grothendieck--Teichm\"uller group, RIMS K\^oky\^uroku
\textbf{1267} (2002), 197--208, Theorem~2, pp.~199--200,
\url{https://www.kurims.kyoto-u.ac.jp/~kyodo/kokyuroku/contents/pdf/1267-18.pdf}.


\bibitem{KKSV2016}
K. Katz, M. Katz, M. M. Schein and U. Vishne,
\emph{Bolza quaternion order and asymptotics of systoles along
congruence subgroups}, Experimental Mathematics \textbf{25}
(2016), no.~4, 399--415,
doi:10.1080/10586458.2015.1073642.
Author version: \url{https://u.math.biu.ac.il/~scheinm/text61.pdf};
arXiv:1405.5454.
\bibitem{KRRBolza}
V. Koziarz, C. Rito and X. Roulleau,
\emph{The Bolza curve and some orbifold ball quotient surfaces},
author PDF, \S4, equation (4.1),
\url{https://www.math.u-bordeaux.fr/~vkoziarz/Bolza.pdf}.

\bibitem{CubicHMM} Y.~Hoshi, A.~Minamide and S.~Mochizuki,
\emph{Group-theoreticity of numerical invariants and distinguished
subgroups of configuration space groups}, Kodai Math. J.
\textbf{45} (2022), 295--348. Definition~2.7 and
Corollary~2.8(i), author pp.~31--32.
\href{https://www.kurims.kyoto-u.ac.jp/~motizuki/Numerical\%20invariants\%20and\%20distinguished\%20subgroups\%20of\%20configuration\%20spaces.pdf}{Author PDF}.
\bibitem{CubicBZ} M.~Boggi and P.~Zalesskii,
\emph{Finite subgroups of the profinite completion of good groups},
arXiv:2406.08639v2, 28 October 2024. Theorem~A(i),(iii),
p.~3; Section~6.1 and Theorem~6.2(i),(iii), pp.~16--17.
\href{https://arxiv.org/pdf/2406.08639v2}{Version 2 PDF}.
\bibitem{CubicAsada} M.~Asada,
\emph{The faithfulness of the monodromy representations associated
with certain families of algebraic curves}, RIMS preprint 1179
(1997). Introduction Theorem, p.~3, parts (i)--(ii), summarizing
Theorems~2 and~3A; the identification with full profinite
mapping class groups is given there.
\href{https://www.kurims.kyoto-u.ac.jp/preprint/file/RIMS1179.ps.gz}{Original preprint}.
\bibitem{CubicBH} J.~S.~Birman and H.~M.~Hilden,
\emph{On isotopies of homeomorphisms of Riemann surfaces},
Ann. of Math. \textbf{97} (1973), 424--439.
Theorem~1, p.~424, and the extremal-map proof on p.~433.
\href{https://www.math.columbia.edu/~jb/bir-hilden-annals.pdf}{Author-hosted scan}.
\bibitem{CubicBHErr} J.~S.~Birman and H.~M.~Hilden,
\emph{Erratum}, 2017, one page, concerning Lemma~5.1 and
Theorem~5 of the preceding article; Theorems~1--4 are unaffected.
\href{https://www.math.columbia.edu/~jb/bh-erratum.pdf}{Author-hosted erratum}.
\bibitem{CubicKerckhoff} S.~P.~Kerckhoff,
\emph{The Nielsen realization problem}, Ann. of Math.
\textbf{117} (1983), 235--265. Section~4, pp.~255--256,
extends the realization argument to finite-type groups with cusps.
\href{https://people.math.harvard.edu/~ctm/home/text/others/kerckhoff/nielsen/nielsen.pdf}{Primary article scan}.
\bibitem{CubicGH} F.~P.~Gardiner and J.~Hu,
\emph{A short course on Teichm\"uller's theorem},
in \emph{Proceedings of the Year on Teichm\"uller Theory},
HRI, Allahabad, vol.~10 (2009), 195--228.
Section~5, Theorem~1 and its punctured-surface hypothesis,
author pp.~14--15; Section~8, Theorem~6, author p.~22.
\href{https://userhome.brooklyn.cuny.edu/gardiner/A\%20short\%20course\%20on\%20Teichmuller\%27s\%20theorem.pdf}{Author PDF}.
\bibitem{Schneps2006Special} L.~Schneps,
Automorphisms of curves and their role in Grothendieck--Teichm\"uller
theory, Math. Nachr. \textbf{279} (2006), 656--671,
Definition~1.2, p.~658, Theorem~3.1, p.~665, and
Theorem~3.4 with equations~(3.10)--(3.11), p.~667.
\href{https://webusers.imj-prg.fr/~leila.schneps/automorphisms.pdf}{Author-hosted journal PDF}.
\bibitem{HMT2025} Y.~Hoshi, S.~Mochizuki and S.~Tsujimura,
Combinatorial construction of the absolute Galois group of the field of
rational numbers, J. Math. Sci. Univ. Tokyo \textbf{32} (2025), 1--125,
Theorem~C, p.~5, and Definition~5.12, pp.~87--90;
\href{https://www.ms.u-tokyo.ac.jp/journal/32-1-1.pdf}{journal text}.
\bibitem{SchnepsSpecialLoci} L.~Schneps, Special loci in moduli spaces
of curves, in \emph{Galois Groups and Fundamental Groups}, MSRI
Publications \textbf{41}, Cambridge University Press, 2003, 217--275,
Section~4.3, Theorems~4.3.1--4.3.2;
\href{https://webusers.imj-prg.fr/~leila.schneps/MSRI.pdf}{author's text},
pp.~54--58.
\bibitem{Furusho2010} H.~Furusho, Pentagon and hexagon equations,
Ann. of Math. \textbf{171} (2010), 545--556,
Theorem~3, pp.~547--548, and \S3, p.~554;
\href{https://annals.math.princeton.edu/wp-content/uploads/annals-v171-n1-p13-p.pdf}{journal text}.
\bibitem{Willwacher} T.~Willwacher, The Deligne--Drinfeld conjecture
after OpenAI, \href{https://arxiv.org/abs/2610.15592v4}{arXiv:2610.15592v4}
(28 October 2026), equations~(1)--(6), p.~1, Theorems~1.1 and~1.3 and
\S1.1, p.~2; \S4.3, p.~9.
\bibitem{OpenAI} OpenAI, The Deligne--Drinfeld conjecture, 23 September
2026, Theorem 1.1, p.~2; \href{https://github.com/openai/math/blob/main/preprints/The-Deligne-Drinfeld-conjecture-September-23-2026/paper.pdf}{OpenAI Math Release}.
\bibitem{Markarian} N.~Markarian, A remark on the AI proof of the
Deligne--Drinfeld conjecture,
\href{https://arxiv.org/abs/2611.01745}{arXiv:2611.01745v2}
(18 November 2026), Theorem 4.8, \S5 and Appendix C.
\bibitem{HainMatsumoto2003} R.~Hain and M.~Matsumoto,
Weighted completion of Galois groups and Galois actions on the
fundamental group of $\mathbf P^1-\{0,1,\infty\}$,
Compositio Math. \textbf{139} (2003), 119--167,
\S8.5 and Remark~8.6, p.~151; \href{https://doi.org/10.1023/B:COMP.0000005077.42732.93}{journal article}.
\bibitem{Neumann1977} O.~Neumann, On maximal $p$-extensions, class numbers
and unit signatures, Ast\'erisque \textbf{41--42} (1977), 239--246,
Theorem~3.7, p.~244;
\href{https://www.numdam.org/article/AST_1977__41-42__239_0.pdf}{Numdam archive}.
\bibitem{LubinTate1965} J.~Lubin and J.~Tate,
Formal complex multiplication in local fields,
Annals of Math. (2) \textbf{81} (1965), 380--387,
p.~380, conditions~(1) and reciprocity formula~(2), the formal-group
criterion on p.~383, and Theorem~2(b),(d),(e), pp.~383--384;
\href{https://www.sas.rochester.edu/mth/sites/doug-ravenel/otherpapers/lubin-tate.pdf}{primary scan}.
\bibitem{StevenhagenRedei2022} P. Stevenhagen, R\'edei reciprocity, governing fields, and negative Pell, Math. Proc. Cambridge Philos. Soc. \textbf{172} (2022), 627--654, \href{https://arxiv.org/abs/1806.06250v3}{arXiv:1806.06250v3}, Lemma~5.1 and Corollary~5.2, author pp.~9--10.
\bibitem{HochschildSerre1953}
G. Hochschild and J.-P. Serre, Cohomology of group extensions,
Trans. Amer. Math. Soc. \textbf{74} (1953), 110--134,
Chapter~II, pp.~118--119; Chapter~III, Theorem~2, pp.~129--130,
and \S6, Theorems~3--4, pp.~132--133;
\href{https://www.sas.rochester.edu/mth/sites/doug-ravenel/otherpapers/hochschild-serre.pdf}{primary scan}.
\bibitem{Otsubo2016} N. Otsubo, Homology of the Fermat tower and universal measures for Jacobi sums, Canad. Math. Bull. \textbf{59} (2016), no.~3, 624--640, \href{https://doi.org/10.4153/CMB-2016-012-0}{doi:10.4153/CMB-2016-012-0}, \S7, p.~637, (7.1)--(7.2).
\bibitem{MilneCM2020} J.~S. Milne, \emph{Complex Multiplication},
version 0.10, 14 July 2020, Example~1.19, p.~14,
Proposition~1.26, p.~17, Corollary~6.11, p.~51, and
Theorem~9.10, p.~78;
\href{https://www.jmilne.org/math/CourseNotes/CM.pdf}{author's course notes}.
\bibitem{Agoh2012} T.~Agoh, On the relative class number of special
cyclotomic fields, Math.\ Appl.\ \textbf{1} (2012), 1--12,
eq.~(1.2), p.~1, \href{https://doi.org/10.13164/ma.2012.01}{doi:10.13164/ma.2012.01}.
\bibitem{Dummit2020} D.~S.~Dummit, A note on the equivalence of the parity
of class numbers and the signature ranks of units in cyclotomic fields,
Nagoya Math.\ J.\ \textbf{238} (2020), 206--214, Proposition~1,
(1)--(2), p.~207, \href{https://doi.org/10.1017/nmj.2018.42}{doi:10.1017/nmj.2018.42}.
\bibitem{AnglesNuccio} B.~Angl\`es and F.~A.~E.~Nuccio, On Jacobi sums
in \(\mathbb Q(\zeta_p)\), Acta Arith.\ \textbf{142} (2010), no.~3,
199--218, definition p.~201; Lemma~5.1 and proof, p.~214,
\href{https://doi.org/10.4064/aa142-3-1}{doi:10.4064/aa142-3-1}.
\bibitem{LS2015Twists} P.~Lochak and L.~Schneps, Every acyclotomic element of the
profinite Grothendieck--Teichm\"uller group is a twist,
Rev. Roumaine Math. Pures Appl. \textbf{60} (2015), no.~2, 117--128;
Theorem~2.1, equation~(8), pp.~121--122; equations~(9)--(12),
Lemma~3.1, pp.~122--124; \href{https://imar.ro/journals/Revue_Mathematique/pdfs/2015/2/3.pdf}{journal text}.
\end{thebibliography}
% Generated after the source commit; excluded from the source manifest.
\par\bigskip\begingroup\footnotesize
\noindent\textbf{Version identifiers.} UNNUMBERED-AQ/v1. The following SHA-256 identifies the source manifest, whose scope excludes the PDF and generated identifiers. The complete payload manifest is available separately.\par
\noindent Source manifest SHA-256:\par\noindent\texttt{4fdcfd31037fec93\allowbreak{}99c174a85ec05d5e\allowbreak{}c57b36be2d02854d\allowbreak{}eba75423674431b5}.\par
\noindent Local source commit:\par\noindent\texttt{d6e646cedccfd6aa\allowbreak{}897b9acc2fbbdfd9\allowbreak{}355c3716}.\par
This commit precedes generation of these identifiers and the PDF; the final release commit is recorded by the version tag.\par\endgroup

\end{document}
